% dossier.tex — source dossier on the word Tilegnelse.
% GENERATED by dossier.py; do not edit by hand. Regenerate with:
%   python3 dossier.py 'tilegn*' --authors sibbern,kierkegaard,brochner,hoeffding --title 'Tilegnelse' -o <this file>
% Each passage is followed by a visible \dcite{...} and by sks-search's own
% citation string as a % comment, in the shape sks.el's `w' key produces.
% Compile with lualatex (Makefile) or xelatex.
\documentclass[11pt]{article}
\usepackage[letterpaper,margin=1in]{geometry}
\usepackage{fontspec}
\IfFileExists{libertinus.sty}{\usepackage{libertinus}}{\setmainfont{DejaVu Serif}}
\IfFileExists{danish.ldf}{\usepackage[english,danish]{babel}}{\usepackage[english]{babel}}
\usepackage{microtype}
\usepackage{xcolor,booktabs,longtable}
\usepackage[hidelinks,bookmarksnumbered]{hyperref}
\setcounter{secnumdepth}{0}
\setcounter{tocdepth}{2}
\makeatletter\renewcommand{\@pnumwidth}{2.2em}\makeatother
\emergencystretch=5em
\tolerance=1500
\newcommand{\dcite}[1]{{\nopagebreak\raggedleft\small #1\par}\medskip}
\newcommand{\tlg}[1]{{\setlength{\fboxsep}{0.8pt}\colorbox{yellow}{#1}}}
\IfFontExistsTF{SBL Hebrew}{\newfontfamily\hebrewfont{SBL Hebrew}}{%
  \IfFontExistsTF{Arial Hebrew}{\newfontfamily\hebrewfont{Arial Hebrew}}{%
    \newfontfamily\hebrewfont{DejaVu Sans}}}
\ifdefined\XeTeXversion \TeXXeTstate=1
  \newcommand{\hebr}[1]{{\hebrewfont\beginR #1\endR}}
\else
  \newcommand{\hebr}[1]{{\hebrewfont\textdir TRT #1}}
\fi
\title{\emph{Tilegnelse} hos Sibbern, Kierkegaard, Brøchner og Høffding\\[0.5ex]\large Kildedossier}
\author{Danish Philosophical Texts}
\date{26 September 2026}
\begin{document}
\maketitle

\begin{otherlanguage}{english}
\section{About this dossier}

\noindent Every passage in which a form of the word occurs (query \texttt{tilegn*},
prefix-matched, so compounds are included), in 4 bodies of text:

\begin{enumerate}
\item \textbf{Sibbern:} 35 passages from 5 works transcribed in this collection. Diplomatic transcriptions; page references are the printed pages recorded in the transcription, and ``p.~109+'' means the passage begins on p.~109 and runs on.
\item \textbf{Kierkegaard:} \emph{Søren Kierkegaards Skrifter}, TEI edition (Det Kgl.\ Bibliotek), reading text only: no deletions, variants or editorial commentary. Published works are ordered by SKS volume, the journals, notebooks and papers by year of composition. Each passage is the whole SKS unit in which the word occurs, which for a published work can be a full page. The citation strings are sks-search's, verbatim.
\item \textbf{Brøchner:} 69 passages from 8 works transcribed in this collection. Diplomatic transcriptions; page references are the printed pages recorded in the transcription, and ``p.~109+'' means the passage begins on p.~109 and runs on.
\item \textbf{Høffding:} 19 passages from 8 works transcribed in this collection. Diplomatic transcriptions; page references are the printed pages recorded in the transcription, and ``p.~109+'' means the passage begins on p.~109 and runs on.
\end{enumerate}

\noindent Caveats:

\begin{enumerate}
\item The search is lexical. It catches every sense of the word, alongside the epistemic and religious uses; nothing has been filtered by sense. Every word form containing \emph{tilegn} is \tlg{highlighted} (386 in all).
\item Much of Kierkegaard's \emph{Notesbøger} (the \texttt{Not} sigla) records lectures and reading. Check the SKS commentary before attributing a notebook passage to him.
\item The quotations are machine-assembled. An entry without a page break of its own carries the page it begins on; for a quotation taken from late in a long entry, check the printed page.
\item 57 further passages in e-book reissues of Høffding are not reproduced: those editions are in copyright.
\item The collection also holds passages by Nielsen (224), Treschow (3) and Møller (1), not included here.
\end{enumerate}

\subsection{Contents by work}
\begin{longtable}{@{}p{0.17\textwidth}p{0.66\textwidth}r@{}}
\toprule
Author & Text & Passages\\
\midrule
\endhead
Sibbern & Om Elskov eller Kjærlighed imellem Mand og Qvinde (1819) & 11 \\
Sibbern & Om Erkjendelse og Granskning (1822) & 9 \\
Sibbern & Udaf Gabrielis' Breve til og fra Hjemmet (1826) & 1 \\
Sibbern & Bemærkninger og Undersøgelser, fornemmelig betreffende Hegels Philosophie (1838) & 11 \\
Sibbern & Om Philosophiens Begreb, Natur og Væsen (1843) & 3 \\
Kierkegaard & BI (SKS 1) & 3 \\
Kierkegaard & EE1 (SKS 2) & 6 \\
Kierkegaard & EE2 (SKS 3) & 6 \\
Kierkegaard & G (SKS 4) & 1 \\
Kierkegaard & PS (SKS 4) & 2 \\
Kierkegaard & BA (SKS 4) & 4 \\
Kierkegaard & 2T43 (SKS 5) & 2 \\
Kierkegaard & 3T43 (SKS 5) & 2 \\
Kierkegaard & 4T43 (SKS 5) & 2 \\
Kierkegaard & 3T44 (SKS 5) & 2 \\
Kierkegaard & TTL (SKS 5) & 6 \\
Kierkegaard & SLV (SKS 6) & 9 \\
Kierkegaard & AE (SKS 7) & 33 \\
Kierkegaard & LA (SKS 8) & 4 \\
Kierkegaard & OTA (SKS 8) & 5 \\
Kierkegaard & KG (SKS 9) & 9 \\
Kierkegaard & CT (SKS 10) & 4 \\
Kierkegaard & TSA (SKS 11) & 2 \\
Kierkegaard & IC (SKS 12) & 5 \\
Kierkegaard & OI6 (SKS 13) & 2 \\
Kierkegaard & OIA (SKS 15) & 1 \\
Kierkegaard & JC (SKS 15) & 3 \\
Kierkegaard & BOA (SKS 15) & 7 \\
Kierkegaard & SFV (SKS 16) & 2 \\
Kierkegaard & Journals etc., 1833 & 1 \\
Kierkegaard & Journals etc., 1835 & 1 \\
Kierkegaard & Journals etc., 1836 & 3 \\
Kierkegaard & Journals etc., 1837 & 4 \\
Kierkegaard & Journals etc., 1838 & 3 \\
Kierkegaard & Journals etc., 1839 & 2 \\
Kierkegaard & Journals etc., 1840 & 2 \\
Kierkegaard & Journals etc., 1841 & 3 \\
Kierkegaard & Journals etc., 1843 & 1 \\
Kierkegaard & Journals etc., 1845 & 1 \\
Kierkegaard & Journals etc., 1846 & 5 \\
Kierkegaard & Journals etc., 1847 & 2 \\
Kierkegaard & Journals etc., 1848 & 2 \\
Kierkegaard & Journals etc., 1849 & 5 \\
Kierkegaard & Journals etc., 1850 & 4 \\
Kierkegaard & Journals etc., 1852 & 1 \\
Kierkegaard & Journals etc., 1853 & 1 \\
Kierkegaard & Journals etc., 1854 & 5 \\
Kierkegaard & Journals etc., Undated & 2 \\
Brøchner & Om Søren Kierkegaards Virksomhed som religieus Forfatter (1855) & 2 \\
Brøchner & Benedict Spinoza. En Monographie (1857) & 5 \\
Brøchner & Problemet om Tro og Viden (1868) & 17 \\
Brøchner & Et Svar til Professor R. Nielsen (1868) & 2 \\
Brøchner & Om det Religiøse i dets Enhed med det Humane (1869) & 13 \\
Brøchner & Bidrag til Opfattelsen af Philosophiens historiske Udvikling (1869) & 13 \\
Brøchner & Den græske Philosophies Historie i Grundrids (1873–1874) & 9 \\
Brøchner & Den nyere Philosophies Historie i Grundrids (1873–1874) & 8 \\
Høffding & Filosofien og Darwinismen (1874) & 3 \\
Høffding & Om Realisme i Videnskab og Tro (1882 / 1884) & 1 \\
Høffding & Forholdet mellem Tro og Viden i dets historiske Udvikling (1885) & 2 \\
Høffding & Kontinuiteten i Kants filosofiske Udviklingsgang (1893) & 2 \\
Høffding & Religionsfilosofi (1901) & 1 \\
Høffding & Vort Hjem. En Indledning (1903) & 2 \\
Høffding & Den menneskelige Tanke (1910) & 5 \\
Høffding & Erindringer — Uddrag af Kapitel II (1928) & 3 \\
\midrule
& Total & 293\\
\bottomrule
\end{longtable}

\end{otherlanguage}

\tableofcontents
\clearpage

\part*{I.\quad F.C. Sibbern}
\addcontentsline{toc}{part}{I.\quad F.C. Sibbern}

\subsection{\textit{Om Elskov eller Kjærlighed imellem Mand og Qvinde} (1819)}

\begin{quote}
Første Bog. — Med uudsigelige Følelser fylder den opvaagnende Elskov enhver inderligt
elskende Sjæl, og i et nyt, herligt og rigt Liv udfolder sig den i
lykkelig Elskov gjenfødte Tilværelse. Men kun altfor ofte skjænkede
Naturen Mennesket denne jordiske Velsignelses rigeste Fylde ved en
Lykke, den har bestemt for Enhver, uden at den Lykkelige lod Det, der
tildeeltes ham, saaledes gaae ind i sig og blive sit, at det kunde
vorde ham til en ægte Skat i Livet og til et Glædens bestandig rindende
Væld. Kun alt for ofte gaaer Elskovs skjønne Blomstretid forbi som en
Drøm, og bliver ikke, hvad den kunde blive, Indgangen til en
uomskiftelig Tilværelse. Og dog er ingen jordisk Vei til det sande Liv
skjønnere, end denne, og intet jordiskt Lys lyser klarere for den ægte
Stræben, end Elskovs kraftigt vækkende og inderligt qvægende Lys. Med
indtrængende Stemme forkynder Naturen i denne Følelse Livets indre
Betydning, og kalder ved den ud af det Synliges vidt udbredte
Mangfoldighed ind i det Usynliges dybere Helligdom. For Den, der vil
trænge ind i Livets og Naturens indre Sandhed og Væsenhed, er den
skjønneste Vielse beredet her. Men skal den vorde ham til en saadan
Indvielse, skal af den en varende Glæde ved Tilværelsen udvikle sig hos
ham, saa maa han fatte den i dens indre Væsen og \tlg{tilegne} sig den af
Hjertet ved en dyb Erkjendelse. Ikke at han, eftersporende og sondrende
Følelsernes vexlende Række, saaledes som de maaskee paa en egen og
mærkelig Maade synes ham at fremstille sig i ham, skulde, grublende
over sig selv, falde ud af Livets fulde Fornemmelse, og nu
sønderskille sig det Liv, han forhen levede heelt og som af eet Stykke.
Men fatte og erkjende maa han, hvad der er skjænket ham, for netop at
kunne leve og nyde det inderligen og tilfulde, idet han føler dets
indre i alle vexlende Momenter sig rørende Dybde, og lader denne komme
til fuld Klarhed og til en gjennemtrængende Udbredelse i hans hele
Tilvær, hvilket ikke skeer uden en Erkjendelse. Ved denne maa Elskoven
hæves frem til at vise sin Magt som et ægte Livet besjælende og
dannende Princip, og ligeledes maa den derved finde sin faste Grund,
hvori den kan rodfæstes, og sit høiere Hjem, hvori den kan vorde
forklaret, i det Eviges Sphære. Uden at finde denne Grundvold og at
naae til denne Forklarelse, vil Elskoven blomstre af, som den
blomstrede frem. Menneskets hele aandelige Liv har sin Rod i Kjærlighed
til den Evige og i Troen paa hans uendelige Væsen, og denne Tro og
denne Kjærlighed maa i enhver Henseende udgjøre Hjertet, Livsfølelsernes
indre varmende og bevægende Middelpunct, i alle det aandelige Livs
Lemmer og Dele. — Derfor, baade for at naae til denne Erkjendelse, og
tillige for overhovedet at fornemme den Naturens Røst, der taler til os
i denne den skjønneste blandt de til en blot jordisk Velsignelse
virkende Livets Magter, ville vi gaae ind i en Eftertænkning over
Elskovs Natur og Indhold. Vi følge kun Naturens Vink, naar vi i
Elskoven, saaledes som den lever i lykkelige, rene og kraftige Sjæle,
og udfolder sig i deres Samliv, søge et Billede paa det ægte Liv. I
denne dens skjønneste og heldigste Skikkelse ville vi da betragte den
først.
\end{quote}
\dcite{Sibbern, Om Elskov eller Kjærlighed imellem Mand og Qvinde, pp. 1–3  [sibbern/om-elskov:11]}
% sibbern/om-elskov:11  (Om Elskov eller Kjærlighed imellem Mand og Qvinde; pp. 1–3)

\begin{quote}
Elskeren er i sin Kjærlighed en Seer, der formaaer igjennem det ydre
Dække at skue ind til Hjertet i det andet Væsen. Et lykkeligt Seerblik
hører til at vække den ægte Elskov. Men enhver ungdommelig bevæget Sjæl
har lettelig saadanne geniale Øieblikke. Dog kan der stundom virkelig
høre Held til, at et andet Væsens Indre kommer til at oplade sig for os
og en Straale fra den Himmel, som deri ligger beredet for os, treffer
vort Øie i en Stund, da vort Sind er frit og aabent nok til at modtage
et reent Indtryk med Reenhed. Ei sjeldent levede Tvende længe ved Siden
af hinanden, inden et lykkeligt Øieblik bragte dem til at oplade deres
inderste Væsen for hinanden og gjensidigen at erkjende hinandens Værd,
saa at nu begge følte sig som rørte af noget Guddommeligt. Stundom laae
vel ogsaa i et Menneskehjerte den himmelske Gehalt tilsluttet og skjult
som en Diamant i Klippen, indtil en lykkelig Finder erkjendte dens Værd
og kaldte den frem til Dagen, for at den nu i Kjærlighedens Sol kunde
udstraale, det er gjenstraale, en uudsigelig Herlighed af Lys. Ogsaa
skeer det ofte, at et Individ, der syntes Andre, og bliver ved at synes
dem, lidet betydeligt, elskes med en Hjertelighed og et Velbehag og en
vedvarende Tilfredshed og Lyst, der for meget beviser Elskovens indre
Ægthed og Dybde, til at der ikke til Grund for samme skulde ligge noget
særdeles Ægte, men hvilket ikkun Elskeren, der lykkeligen har fattet og
skuet den dybere Grundskikkelse i sin Elskede, har været i Stand til at
erkjende og \tlg{tilegne} sig saa hjerteligen. Vistnok kan Kjærlighed, uagtet
den efter sin indre Natur og Væsen gjør seende, dog undertiden bringe
et Individ til at hildes i Vildfarelse og Blindhed, ligesom en
philosophisk Grunden og Stræben ei sjeldent førte til Sværmerie eller
Ideeforvirring, uagtet Philosophie og philosophisk Grunden dog netop er
Noget, der oplader Aandens høiere Øie og gjør ægte seende. Men
ingenlunde er altid den Elskov fremkaldet ved et Blændværk, der synes
Andre saa, fordi de i den Elskede ikke see noget saa Elskværdigt, at
det kan synes dem i Stand til at yde saa megen Lyst og Fyldestgjørelse.
En menneskelig Sjæl gjemmer ofte meget Mere i sit indre Dyb, end hvad
der lader sig tilsyne for Andre, end Den, hvem Elskoven gjør til Seer.
Det Guddommelige lever jo dog virkeligen i ethvert godt Menneske og
giver Mennesket Væsenhed, være det end overskygget eller trykket
tilbage formedelst jordiske Mangler, hvilke dog letteligen svinde hen
for Den, der erkjender og skuer hiint Guddommelige i dets indre
Reenhed.
\end{quote}
\dcite{Sibbern, Om Elskov eller Kjærlighed imellem Mand og Qvinde, pp. 13–15  [sibbern/om-elskov:19]}
% sibbern/om-elskov:19  (Om Elskov eller Kjærlighed imellem Mand og Qvinde; pp. 13–15)

\begin{quote}
Dog for Den, der med inderlig Lyst \tlg{tilegner} sig sin Elskov og sin
Elskede, er denne altid skjøn. Og det er aldeles naturligt, at den
elskede Skikkelse udøver Skjønhedens Magt over ham. Al Skjønhed beroer
paa, at Noget, som udspringer af Livets høiere Kilde og bærer dens Præg
i sig, træder fuldeligen frem i en ydre sandselig eller timelig Form,
udfyldende og gjennemtrængende denne, saa at denne bærer det fulde
Udtryk af hiint. Men hvad Elskeren skuer som Elskværdigt, tilmed
Himmelsk, i den Elskedes indre Væsen, dette Samme netop seer han med
inderlig Glæde udtrykt og udtalt i den hele elskede Individualitet og
enhver af dens Yttringer. Og saaledes vil den Elskede ogsaa i
Henseende til sit Ydre gjøre Skjønhedens Indtryk hos Elskeren, og hans
Glæde og Velbehag vil være af samme Art som den, Skjønheden fremkalder.
Er det ikke ogsaa en Straale fra den sande Himmel, som falder i hans
Sjæl igjennem hans Elskede, og vækker denne ikke den uendelige Følelse
af det evige Skjønnes himmelske Herlighed? Men hvad der er i Stand til
at vække hos os denne Følelse, staaer derved selv forklaret for os, og
endog hvor den individuelle Fuldendthed og Fylde fattes, prise vi det
Væsen som skjønt, ved hvis Beskuelse en Mindelse om den evige Skjønhed
bestandigen lever op igjen hos os. Derfor udvikler sig da af Elskov et
Velbehag og en Lyst ei blot ved det i sig Fortrinlige og Væsentlige,
men ogsaa ved det Ubetydeligste i den Elskedes Yttringer og Bevægelser,
fordi den Elskende deri ideligen møder den elskede Individualitet, og
forlyster sig ved at see den udpræget indtil den yderste Overflade af
Formen. Dog fornemmes denne rene Glæde fuldest i det totale Indtryk, i
hvilket alt Enkelt og Mangfoldigt smelter sammen i en eneste rolig og
heel Anskuelse, hvori den Elskedes hele Væsen udfylder vor Sjæl.
\end{quote}
\dcite{Sibbern, Om Elskov eller Kjærlighed imellem Mand og Qvinde, pp. 16–17  [sibbern/om-elskov:21]}
% sibbern/om-elskov:21  (Om Elskov eller Kjærlighed imellem Mand og Qvinde; pp. 16–17)

\begin{quote}
Men som Elskov er en ubegrændset Hengivenhed, saa er den ogsaa en
ubegrændset Attraa. Ligesaa inderlig Elskeren taber sig i sin Elskede
og hænger ved denne, lige saa indtrængende staaer hans Hu til at møde
samme Hengivenhed og ganske at kunne \tlg{tilegne} sig det tilbedte Væsen og
tage det i Besiddelse. Begge Dele, Hengivenheden og Attraaen efter
udelukkende Forening, ligge i den af Elskovens Inderste fremgaaende
Drivt til Sammensmeltning og fuldkommen Eenhed. Som Trang og Længsel og
Begjær er Elskoven lige saa udelukkende og færdig til at fortrænge alt
Fremmedt, ja ligesaa uslukkelig, som dens sig selv hengivende og
hentabende Lyst er uudtømmelig. Baade som en Lyst til i sig at optage,
ja ligesom fortære den Elskede, saa og som en Lyst til at løses op og
forgaae, ja ligesom at døe hen i den Anden, yttrer sig Elskovens
guddommelige Afsindighed (μανία θεία) i tusinde endog mimiske Udbrud.
Ved denne udelukkende Attraa kan nu ogsaa det Heftige, Brændende,
Urolige, i Fryd og Qvide Vexlende, komme ind i Elskovens Følelser, og
den Elskende komme til at prøve, hvad Platon (Symposion, 403) skildrer
som Eros's Vilkaar: At han, født af Poros og Penia, Overflod og Armod,
som et Mellemvæsen mellem Guderne og de Dødelige, bærer Præg af sin
dobbelte Herkomst. Først — saa hedder det om ham — er Eros „altid
fattig og langt fra at være fiin og smuk, som de Fleste troe; derimod
er han raa, af forvildet Udseende, løber om uden Skoe og uden Hjem,
ligger paa den bare Jord uden Dække, sover ved Dørene og paa Veiene
under aaben Himmel; han har sin Moders Natur, boer altid sammen med
Trangen. Men atter efter Faderen efterstræber han det Gode og Skjønne,
er mandig, kjæk, drivtig, en drabelig Jæger, altid pønsende paa
Konstgreb, begjærlig efter Indsigt og fuld af Paafund; han
philosopherer hele Livet igjennem, er stærk i Trolddom og i at blande
Trylledrikke, en vældig Sophist; han er hverken som en Udødelig eller
som en Dødelig af Natur; men snart paa selvsamme Dag blomstrer og lever
han, naar det gaaer ham vel, snart døer han hen, men lever atter op
igjen, ifølge Faderens Natur; hvad han erhverver, rinder bestandig bort
igjen, saa at Eros hverken er rig eller fattig, og ligeledes staaer
midt imellem Viisdom og Uforstand.“
\end{quote}
\dcite{Sibbern, Om Elskov eller Kjærlighed imellem Mand og Qvinde, pp. 27–29  [sibbern/om-elskov:30]}
% sibbern/om-elskov:30  (Om Elskov eller Kjærlighed imellem Mand og Qvinde; pp. 27–29)

\begin{quote}
En dyb om Individets eget Jeg sig dreiende Følelse, en dyb jegisk
Følelse altsaa, gaaer i Elskoven i Eet med en dyb sympathisk. Allerede
i blot organisk Henseende er Driften i sit Udspring en individuel og
jegisk, thi den er en umiddelbar Fortvirkning af den til bestandig
Selvuddannelse og individuel Formation virkende organiske Tendents. Og
siden sættes ved denne Drivt ikke blot Sympathien i Bevægelse, idet den
drager det ene Individ til det andet, men dernæst ogsaa reiser sig
Følelsen af dybe sjælelige Savn. Her viser sig en inderlig stærk Trang,
som kræver Fyldestgjørelse. Bevæget af inderlig Glæde over det Skjønne
og Guddommelige, det Tiltrækkende og Qvægende, der skues i et andet
Individ, føler Mennesket sig snart som Den, hvis isolerede Liv kun er et
halvt, altsaa som Den, der, for at komme til fuldelig Existents, maa
voxe sammen med et andet Væsen, for nu i et uafbrudt Samliv med dette,
delende Alt med samme, at nyde Livet som en uafbrudt Elsken. Derfor
attraaer han en fast og inderlig og udelukkende Forening; ikke i denne
eller hiin Henseende alene, men i alle Henseender. Han begjærer nu af
al Magt, ikke blot at hengive sig selv, men ogsaa at \tlg{tilegne} sig den
Elskede og udelukkende at besidde denne; ei blot at elske, men og at
elskes igjen med samme Inderlighed.
\end{quote}
\dcite{Sibbern, Om Elskov eller Kjærlighed imellem Mand og Qvinde, pp. 29–30  [sibbern/om-elskov:31]}
% sibbern/om-elskov:31  (Om Elskov eller Kjærlighed imellem Mand og Qvinde; pp. 29–30)

\begin{quote}
Men hvad den Elskende attraaer, er, saa himmelsk end Elskovens Salighed
monne være, dog altid et jordiskt Gode. Thi det er et jordiskt Væsens
Gjenkjærlighed og en jordisk Forening. Det er overhovedet en Lykke, som
ei altid undes selv Den, der længes meest derefter; den kan ogsaa gives
og fortabes igjen. Hvor afhængig er Elskovs Lykke af jordiske
Sammenstød? Derfor kan det ved lykkelig Elskovs dybe Følelser opløste
og hensvundne Selv ikke blive ved at være fortabt i et høiere Livs
udeelte og ustandsede Henstrømmen. Jeghedsfølelsen kan ikke hvile. Paa
tusinde Maader ægges den op igjen, og under en idelig Stræben efter at
\tlg{tilegne} sig, holde fast og bevare det \tlg{Tilegnede}, eller forvisse sig om
den udeelte Bevaren af dets Besiddelse, optændes nu letteligen alt det
Hede og Bryndefulde i Sjælen, der ligger i al jegisk Existents, naar
Individet kommer til skarpt at fornemme, at det har at gjøre med en
muligen modstridende Omverden. Ved den store Opæggelighed
(Irritabilitet) hos det ved Elskov blødnede Væsen reiser sig da altfor
let midt under den himmelske Glæde en fra Skinsygens frygtelige Helvede 
kommende Fornemmelse. Ved al jordisk Besiddelse er Fare. Tilmed er
Frygten tilrede med de fjendtligste Følelser. Allerede ved sin Tilvær
alene kan Omverdenen da bringe heftige Gemyter i lidenskabelig
Bevægelse; og endog om herfra intet Fiendtligt ægges op i Sjælen, kan
letteligen den Individualiteternes Kamp, som maa opstaae, naar tvende
allerede til en vis Grad fixerede Organisationer skulle smelte hen og
løsnes, for at begge kunne voxe sammen til et nyt begge forenigende
Liv, fremkalde urolig Gjæring og optænde Harme. Ja det kan maaskee være
godt, om denne Gjæring indtræder tidlig under Følelsernes endnu rigeste
Fylde og dybeste Inderlighed, naar Elskoven endnu er frodigst til at
fortrænge al Standsning og at udfolde sit indre Væsen til ydre
Skjønhed. Saaledes som Menneskets Vilkaar nu eengang ere paa Jorden,
kan den rene, i sin indre Lyst heelt hensjunkne Elskov neppe holde ud,
endog i de uskyldigste Hjerter, uden at de engang komme til at fornemme
og erkjende, at det skumle Mørkes glædefiendske Vælde ogsaa kan gribe
ind i den jordiske Lyksaligheds meest himmelske Sphære, naar Mennesket
ei vogter sit Hjerte. Men hvad et dybtfølende Gemyt lider, naar det
første Gang føler Muligheden af et Brud paa Samlivets rene Heelhed, er
sikkert ofte langt Mere, end Nogen troer. Dog, vi ville siden finde
Stedet til videre at oplyse Elskovs, om end stundom ganske naturlige,
dog atter igjen ofte saa høist unaturlige, Lidenskabelighed. Endnu ville
vi blive ved med at betragte Elskoven i den skjønne Skikkelse, hvori
den lever i tvende lykkeligt elskende Væseners Samliv.
\end{quote}
\dcite{Sibbern, Om Elskov eller Kjærlighed imellem Mand og Qvinde, pp. 30–32  [sibbern/om-elskov:32]}
% sibbern/om-elskov:32  (Om Elskov eller Kjærlighed imellem Mand og Qvinde; pp. 30–32)

\begin{quote}
Men vistnok maa Naturen her være naturlig. Den er det aldrig, naar den
ei er reen og sædelig; thi den er det ei, naar den i vilde Udbrud
sønderbryder sin egen rolige Bestaaen og indre Fred. Al ægte Elskov er
af Hjertet reen og derfor ogsaa paalidelig, naar den kun er i sig selv
nogenlunde klar. Men det er den ei altid, og saaledes kan der være
Fare, hvor der mindst skulde være den. Dog, da det ellers Naturlige saa
er saa lidet naturligt og lidet menneskeligt, er ogsaa, naar Elskov kun
ikke er alt for ubesindig eller uægte, den selv i høi Grad værnende. I
alle gode Øieblikke føles Elskov som en uendelig Omhu og Baretægt, og i
dens fulde, skjønne Henrykkelse løser enhver Begjær sig op i den
aandeligste. Saadanne Øieblikkes Kraft maa da formedelst indre
Betænkning og Bekræftelse gjennemtrængende udbrede sig til de øvrige. Jo
mere Elskovs dybe Følelser komme til Klarhed, jo fuldeligere den Attraa
er, hvormed Elskeren \tlg{tilegner} sig sin Elskovs hele Fylde og paa alle
Maader knytter sig til den og lever i den, desto større vil Elskovs
indre Reenhed være. Den bevare da ogsaa selv under den nøieste
Forbindelses retlige Frihed denne Reenhed, der holder al ἀκολασία
borte! Thi her er endog da en Grændse, det for Elskovs egen Tryghed og
Ro er farligt at overskride.
\end{quote}
\dcite{Sibbern, Om Elskov eller Kjærlighed imellem Mand og Qvinde, p. 48+  [sibbern/om-elskov:45]}
% sibbern/om-elskov:45  (Om Elskov eller Kjærlighed imellem Mand og Qvinde; p. 48+)

\begin{quote}
Men hvad der saaledes af sig selv er udsprunget i Mennesket, og, som en
vældig Naturindflydelse af høiere Art, bevæger ham i de stærkeste
Sindsrørelser, dette komme nu til inderlig Klarhed i ham, og han \tlg{tilegne}
sig det af Hjertet og tilfulde! En klar Besindighed træde her, som
overalt, til det Geniale, for at det ved en høiere Beaanding givne Eie
ikke skal besiddes med Blindhed, eller Henrevenhedens skjønne Opløftelse
gaae forbi som en Beruusning. Men det Geniale (τὸ δαιμόνιον, som Platon
kalder det) \tlg{tilegner} sig den Besindige, deels erkjendende, deels
villende.
\end{quote}
\dcite{Sibbern, Om Elskov eller Kjærlighed imellem Mand og Qvinde, pp. 58–59  [sibbern/om-elskov:56]}
% sibbern/om-elskov:56  (Om Elskov eller Kjærlighed imellem Mand og Qvinde; pp. 58–59)

\begin{quote}
Men ligesom paa den ene Side ved en Erkjendelse, saaledes maae vi paa
den anden Side ved en fast Beslutning \tlg{tilegne} os det nye Liv, der gaaer
op for os i vor Elskov. Vi elske, fordi vi maae elske; men hvad vi saa
uimodstaaeligen maae og føle os drevne til, dette maae vi nu ogsaa ville.
Den Elskende vil sin Kjærlighed, saa ville han den da nu ogsaa af
Hjertens Grund, med et sig selv besiddende Sinds indre Frihed. Vi leve
os ind i Elskoven. Men den inderlige udelukkende Forbindelse for hele
Livet skal udtrykkeligen indgaaes. Hvad der er skeet af sig selv, skal
Mennesket med fri Selvbestemmelse bekræfte, derfor ogsaa først prøve sig
selv, om han ogsaa af Hjertens Grund kan ville og virkeligen vil, hvad
han nu selv skal beslutte. Forholdet blive da til et ægte moralskt
Forhold, stiftet ved en fri Villies og Beslutnings faste Bestemmelse! I
denne faste hellige Beslutning, nu ogsaa af Hjertet at ville \tlg{tilegne} sig
den vundne Lykke og holde den fast og ganske give sig hen til den, og
derved at fixere sig selv og den let bevægede Lyst og Drivt, har
Troskaben sin Rod og sit Sæde.
\end{quote}
\dcite{Sibbern, Om Elskov eller Kjærlighed imellem Mand og Qvinde, pp. 59–60  [sibbern/om-elskov:58]}
% sibbern/om-elskov:58  (Om Elskov eller Kjærlighed imellem Mand og Qvinde; pp. 59–60)

\begin{quote}
At den Forening, hvortil Elskov fører, ei skal være en Forening i denne
eller hiin Henseende alene, men i alle Henseender, og at den ei skal
medføre et partielt, men et totalt Samliv, Dette føler i Ungdommens
skjønne Periode enhver lykkelig Elskende, og medens han selv stræber
efter at \tlg{tilegne} sig den Erkjendelsens og Sjælenydelsens Fylde, der har
opladt sig for ham, og hvorved et opvakt og ædelt ungdommeligt Sind
glæder sig saa meget, attraaer han tillige af al Magt at leve og nyde
det rige Liv med sin Elskede. Begge finde da tusinde Berøringspuncter,
og tusindfoldig Meddelelse finder Sted imellem dem, saafremt de kun
ellers høre sammen. Men endog De, hvis Elskov faldt i denne skjønne
ungdommelige Tid, og endnu oftere De, der først senere fandt Elskovs
Lykke, glemme saa ofte, naar den fulde Forening er indtraadt og Livets
øvrige Krav og Kald indfinde sig igjen, den lykkelige Skikkelse, deres
ungdommelige Elskov havde. Alt for ofte bliver endog Elskov taget som en
Overgang: de tvende forbundne Individer træde efterhaanden i sjælelig
Henseende tilbage i deres forrige Maade at existere paa, og komme Tid
efter anden saavidt, at begge tilsidst leve, hver for sig, deres
isolerede Liv ved Siden af hinanden, istedetfor at leve et, begges hele
Liv omfattende, Samliv med hinanden. Istedetfor at de skulle stræbe at
gaae ind i hinandens Individualiteter og knytte sig under mangfoldig
Meddelelse til hinanden i mange Puncter, indskrænker sig Samlivet til
nogle faa Puncter, hvori begge føle sig som sammenvoxede; udenfor disse
er det, som om de ei hørte sammen. Elskende ere paa mange Maader og i
mange Henseender Incitament for hinanden. Men til hvad Side Elskoven
fornemmelig kommer til at vende sig, beroer paa, hvorledes de tage den.
Og saa skeer det da, at ofte stor Eensidighed faaer Overhaand, og at
hele Regioner i deres Samliv komme til at ligge øde. Ofte indtræder
denne Tilstand nu vistnok, fordi begge i sig selv virkelig ikke hørte
sammen, fordi den ene ei kan gaae ind i den andens Ideekreds; stundom
derimod af en vis overstor passiv Tilfredshed i Begyndelsen, og af
Ladhed eller dog Mangel paa Opmærksomhed og rask Aandskraft siden: en
Ladhed eller Inertie, der endog ofte fører til at bornere Omdømmet og
det frie Blik paa Livet, saa at de Elskende endog tabe Sandsen for det
mangfoldige Herlige, der, forskjelligt fra Det, de skue i hinanden,
bevæger sig omkring dem, og nu derfor tabe Ansporelsen til egen
Fremadstræben og Udvikling. Imidlertid give de sig da dog gjærne hen til
en vis Art af de Bevægelser, hvori Elskov sætter. Der gives Mange, hvis
hele Elskov synes at blive staaende ved den Lyst, hvormed de beskue
hinanden og gjøre af hinanden; og nu, for dog at bringe en vis
Mangfoldighed og Afvexling ind i Forholdet og fordrive en
Kjedsommelighedsfølelse, hvilken Kjærtegn ei altid kunne forjage, søge
de idelige nye Opvækkelser, der dreie sig om det samme Punct, og Aanden
kan derover ei komme op. Heri forsynde sig mange Elskende og glemme, at
selv Eros's Tjeneste ei maa vanhelliges, glemme, at Elskov skal føre til
en Sjælemeddelelse og et aandeligt Forhold, og tabe nu Livets uendelige
Nydelse over en høist indsnævrende, der nu ei engang ret bliver en
Nydelse af Hjertens Grund, og let kan blive en aandsfortærende. Eller og
skeer det, at man lader Livets Torne og Tidsler faae saa megen Overhaand
i det aandelige Samliv, at de tilsidst underkue det fremspirende Frø.
Undertiden var vel ogsaa et af den i saa mange Henseender ødelæggende
Egoismus udspringende Overmod, en selvraadig og herrisk Lyst, til
Fordærvelse her, naar Den, der havde givet sig ganske hen, dog nu var
for stolt til at ville give sig ganske hen. Men ei sjeldent lider endog
et fra først af frodigt Samliv, og falder eller visner bort
efterhaanden, fordi der fattes en tilstrækkelig Næring, indtil tilsidst
Opmærksomheden paa den Næring, der endda kunde frembyde sig, ogsaa
svækkes. Omgivelserne bidrage da stundom ei lidet til at forstørre
Kløften mellem Begge. Der vises ofte for ringe Opmærksomhed eller for
liden Sands for Egteskabets Hellighed, og betænkes ei altid hvad hver En
ogsaa i denne Henseende skylder Andre. De kjæreste Mennesker ere os ei
altid kjærkomne, og kunne ei være det, naar de træde forstyrrende ind i
vor Kreds. For at nære et ægte Samqvem med hinanden, udfordres blandt
Andet ogsaa undertiden Eensomhedens Ro. I en Kreds, hvori et udbredt
selskabeligt Samliv let danner sig, lide ikke sjeldent de
fortreffeligste Mennesker meest i denne Henseende, fordi de
fortreffeligste snarest indeslutte saadanne Følelser dybt ind i sig
selv, der ved at komme frem til Dagen let kunde profaneres.
\end{quote}
\dcite{Sibbern, Om Elskov eller Kjærlighed imellem Mand og Qvinde, pp. 98–101  [sibbern/om-elskov:82]}
% sibbern/om-elskov:82  (Om Elskov eller Kjærlighed imellem Mand og Qvinde; pp. 98–101)

\begin{quote}
Endog Bitterhed og Harme fremægges let ved den ulykkelige Elskov, ikke
blot mod en mere lykkelig Tredie, men ogsaa mod det elskede Væsen selv.
Forbittrelse og Vrede og en vis, endog stundom til grusomme Tanker
stigende, Lyst til fiendtligen at antaste, ja ligesom at sønderrive det
Væsen, der ikke kan bevæges ved saa megen Tilbedelse, griber let om sig
ved Siden af Elskoven, fordi, naar den stærke Begjæring og Higen kommer
op, saa er det Egoistiske i Bevægelse, men dette er efter sin Natur lige
saa nær til at yttre sig fiendtligt og bortstødende, som til at vise
Tilbøielighed. Det heftige Incitament pirrer og ægger, men irriterer og
forbittrer derfor ogsaa, naar det, fordi Tilfredsstillelsen ei kommer,
bliver piinligt. Saaledes træder da let i en ny Skikkelse Jalousiens Hede
frem, ei mindre fordærvelig, end i dens forhen betragtede Skikkelser.
Dette er det Værste, hvortil det kan komme, og rører noget Saadant sig op
i Sindet, da er det vistnok høi Tid at betænke, at det er bedre at flye,
end at lide saadan Brynde. Elskoven viser atter her sin dobbelte Natur,
hvorledes den ikke blot er af guddommelig, men ogsaa af en aldeles
jordisk Herkomst. Lysten til udelukkende \tlg{Tilegnelse}, til at besidde og
kalde sit, hvad der er Gjenstand for Velbehag, fordærver atter her det
skjønne Eie, der, naar der elskedes med mere philosophisk Kraft og Dybde,
kunde nydes og bevares. Derfor tænke den i den haabløse Elskov fængslede
Sjæl allerede tidligen paa at bevare sig det skjønne, fulde Billede, den
har optaget i sig. At bære et reent og herligt Billede af en elsket
Individualitets Skjønhed og Fylde og hele Eiendommelighed i sig, og netop
at bære det i sig med en reen og total Hengivenhed, fri for Begjær, er
noget usigeligt Qvægende, Oplivende og Sjælen inderligt Nærende. Just af
en saadan Natur skal ogsaa den Hvile være, hvortil den lykkeligt Elskende
skal komme ved en lykkelig, al Attraa stillende Forbindelse; og for den
haabløs Elskende skal saaledes af Kjærlighedens egen Fylde den høiere
Erkjendelse og Glæde ved det Skjønne og Guddommelige gaae frem, der kan
føre den en Tid lang i en snæver Kreds hildede og fangne Sjæl tilbage
igjen til en friere Tilvær og vorde den frugtbare Spire til et nyt i
høiere Regioner sig med ny Kraft bevægende Liv. Saa at det igjen maa gaae
i Opfyldelse, at Den, der veed at lade det Guddommelige komme til
Herredømme i sit Indre, ham maa Alt komme tilgode. Hvortil iøvrigt en vis
bestemt ret dybt i Livet indtrængende Virksomhed, enten det saa er som
Studium og Ideers Udvikling og Betragtning, eller paa en mere praktisk, i
det ydre Menneskeliv mere umiddelbart indgribende Maade, er en nødvendig
Betingelse. Især i Stand til at frigjøre Sindet og fremme den nye Livets
Kraft er en saadan Virksomhed, der staaer i nær Forbindelse med
Kjærligheden til det Evige og frembyder mangfoldig Anledning til, at høie
Ideer og Følelser for det Store og Skjønne sættes i Bevægelse paa en
Maade, der river Individet udaf sit individuelle Livs Tankekreds midt ind
i den store Natur og det almene overalt rige og betydningsfulde Liv.
\end{quote}
\dcite{Sibbern, Om Elskov eller Kjærlighed imellem Mand og Qvinde, pp. 162–164  [sibbern/om-elskov:117]}
% sibbern/om-elskov:117  (Om Elskov eller Kjærlighed imellem Mand og Qvinde; pp. 162–164)

\subsection{\textit{Om Erkjendelse og Granskning} (1822)}

\begin{quote}
Men overvejes nærmere, hvad der hører derhen og hvad ikke, saa vise sig
Vanskeligheder, hvilke ogsaa Forfatteren af det nærværende Forsøg ej har
kunnet andet end fornemme. For at oplyse det videnskabelige Studiums Natur og
Væsen maa man gaae ind i Betragtninger, hvilke dog andre videnskabelige
Udviklinger \tlg{tilegne} sig som sig tilhørende, og Spørsmaalet bliver nu,
hvorledes en Grændse drages mellem disse og hiin. Det være tilladt her at
omtale denne Sag med nogle Ord, for at det Forhold, hvori denne Fremstilling
er tænkt at skulle staae til de andre videnskabelige Udviklinger, med hvilke
den paa flere Steder mødes, kan blive oplyst.
\end{quote}
\dcite{Sibbern, Om Erkjendelse og Granskning  [sibbern/erkjendelse:4]}
% sibbern/erkjendelse:4  (Om Erkjendelse og Granskning)

\begin{quote}
En Fremstilling af denne Art kan betragtes som et Slags Monographie. Af
saadanne for sig fremtrædende videnskabelige Fremstillinger af visse enkelte,
eller til større videnskabelige Systemer henhørende, Gjenstandes hele Natur
og Beskaffenhed, kan man tænke sig to Arter. De sædvanlige naturhistoriske,
og de, der ligne dem, behandle deres Gjenstand ganske saaledes, som den
ellers behandles paa sit Sted i den Lærebygning, hvortil den hører, og ere da
kun at ansee for udhævede, blot maaskee mere end sædvanligen udførte, Afsnit
eller Stykker af det større videnskabelige Hele. Af en anden Beskaffenhed
derimod er en Monographie, som den nærværende. Erkjendelser, der ellers høre
hjemme paa de forskjelligste Steder i andre Videnskaber, skulle her samle sig
om eet Punct, for tilligemed, og i eet med, Betragtningerne over det særegne,
der er ejendommeligt for Gjenstanden, at oplyse denne tilfulde og fra alle
Sider. En saadan Monographie, mere vanskelig og indviklet end en af første
Art, har ogsaa større Adkomst til at træde op som en for sig bestaaende
videnskabelig Betragtning, eller som et mindre Hele i Videnskabernes store
Hele, da der her er et ejendommeligt Middelpunct for en egen videnskabelig
Organisation, hvorved et Indbegreb af sammenhængende Erkjendelser paa en
særegen Maade voxe sammen til Eenhed; og denne Organisation, eller
Fremstillingens indre Grundform, dens Deles indre Forbindelsesmaade og Følge,
eller Udviklingens Gang, bestemmes ved Gjenstandens ejendommelige Grundnatur,
ikke ved en fra en anden Videnskab hentet og ved dennes Gang og Methode
bestemt Norm og Fremstillingsmaade. En Monographie af denne Art forholder sig
til de Videnskaber, til hvis Grundbetragtninger den knytter sig, og til
hvilke meget er at føre tilbage af det, hvad den, skjønt altid paa en ved
Gjenstanden bestemmet ejendommelig Maade og i et eget Lys, fremstiller og
udfører, ligesom et Talent forholder sig til de almindelige
Aandsvirksomheder. Drag bort af Talentet, hvad der hører Forstand og
Judicium, Sands og Phantasie, Fornuft og Følelse til, der vil intet blive
tilbage, som kunde kaldes Talent, og dog er Talentet noget andet, end hine
almenere Aandsvirksomheder. Hvad er Talent da? Den er den ejendommelige
Combination af dem alle i visse bestemte Retninger, den ejendommelige Aandens
Organisation i Henseende til alles forenede Samvirkning, der gaaer ud paa en
vis Art af Productivitet. Ligesom nu et vist Talents grundige psychologiske
Skildring ikke bestaaer i en Skildring af hine almenere Aandsvirksomheder i
Almindelighed, men maa forudsætte dem som bekjendte, medens netop de selv dog
ere det, den overalt maa agte paa og have for Øje, saaledes kan en
Monographie, som den nærværende, have at see hen til mange enkelte Puncter,
som i og for sig tagne høre hen til bestemte Steder i bestemte Videnskaber,
og kan have at fremstille meget, hvilket stykkeviis andre Videnskaber kunne
\tlg{tilegne} sig og drage til sig, og kan dog være, og haandhæve sig som, et
ejendommeligt videnskabeligt Hele for sig.
\end{quote}
\dcite{Sibbern, Om Erkjendelse og Granskning  [sibbern/erkjendelse:7]}
% sibbern/erkjendelse:7  (Om Erkjendelse og Granskning)

\begin{quote}
hvori en Person lader sig til Syne og træder frem som saadan, idet den ganske
knytter sig til, og giver sig hen til, hiin Virken. I den rene Modtagen og Skuen
saaes endnu kun det Sande som fremstillende sig selv, saa at kun dets egen
(universelle) Fremtræden deri viiste sig, ingen personlig Leven. Modtagelsen var
blot som en Afspejling eller Gjenspejling, liig den, hvormed Vandet gjengiver
Lyset og Farven, ikke en personlig Optagelse, liig den, hvorved Ædelstenen i sin
Farve har en Lysets Kraft og Virkning optaget i sig til Eenhed med sin egen
Natur og Tilvær. Men i Overbeviisningen bliver den til en inderlig \tlg{Tilegnelse},
hvorved det sig i og for en Individualitet fremstillende Sande nu ogsaa kommer
til kraftigen at boe i denne Individualitet, som i et Væsen, i hvilket det har
taget Sæde, og hvilket nu selv er og bliver noget ved Optagelsen. I denne det
Sandes og Individets Sammenknytning til Eenhed er ogsaa den Nødvendighed, som
vistnok er deri, saaledes fortabt, at den ikke særligen fornemmes eller kommer
til Syne. Her er en udeelt, over Nødvendighedens Følelse opløftet, fri
Leven i det Sande. Kun naar det individuelle, rivende sig løs fra det almene,
bevæger sig i en særlig personlig Stræben, kommer Modsætningen og tilligemed
samme Nødvendigheden for Lyset.
\end{quote}
\dcite{Sibbern, Om Erkjendelse og Granskning, pp. 9–10  [sibbern/erkjendelse:49]}
% sibbern/erkjendelse:49  (Om Erkjendelse og Granskning; pp. 9–10)

\begin{quote}
Sandheden viser sig herved som den, der ved sin egen Magt og Kraft søger
at tage Skikkelse i Individet, medens dette aldeles giver sig hen til at lade
den virke hos sig, hvad den selv vil. Og ligesom dette gjelder om alt, hvad der
i hvilkensomhelst Videnskab eller Kundskab viser sig som det Sande, paa hvis
Erkjendelse Studiet gaaer ud, saaledes gjelder ogsaa alt, hvad der her er
fremsat, om enhver videnskabelig Stræben, selv om den, der endnu staaer saa
langt fra de højere Ideer, at den kun sjeldent føler sig beaandet deraf. Under
ethvert Studium vil den trolige Hengivenhed til den Erkjendelse, som
former sig i os, bringe den Fylde, Kjerlighed giver. Synes i Granskningens
lavere Kredse hvad vi erkjende, ikke at bringe os det Evige nærmere, saa vil dog
det, at det virker fuldkommen Erkjendelse og levende Overbeviisning hos os,
gjøre, at vi føle os som opfyldte af en Guddom. (Som da ogsaa, selv i de laveste
Sphærer, Erkjendelsen i sig selv virkelig kommer fra det højere Udspring, om den
end ikke kommer til Bevidsthed som saadan. See længere hen §. 11 i Anm. 2). Saa
at ethvert Studium, hvori troelig og klar Granskning hersker, har sin
Muse. Individet træder deri lige over for noget, som det vil \tlg{tilegne} sig,
trængende dertil, som til noget, uden hvilket det selv ej kan komme til fuldelig
Leven. Men hvad det saaledes vil \tlg{tilegne} sig, er ikke noget, det kan behandle
som Gjenstand for en sandselig Fornødenheds Tilfredsstillelse, fortære eller
tage som et blot Vehikel eller Middel. Det bliver altid staaende lige over for
Individet, som noget, der bestaaer i og for sig i sin Betydenhed og Værd, hvad
enten det kommer til Individet eller ej, ja som noget, der netop kun forsaavidt
har Værd for Individet, som det har Værd og Betydenhed (det er paa dette Sted
har Sandhed) i og for sig. Og saaledes er overalt, hvor Sandhed søges, i hvilken
Sphære det saa er, Forholdet mellem Individet og det Sande, det søger, et
saadant Tilknytningsforhold, der kan, og under en ægte Stræben ogsaa vil, blive
et Kjerlighedsforhold.
\end{quote}
\dcite{Sibbern, Om Erkjendelse og Granskning, pp. 12–14  [sibbern/erkjendelse:51]}
% sibbern/erkjendelse:51  (Om Erkjendelse og Granskning; pp. 12–14)

\begin{quote}
4. Formedelst den frie Tænkning skeer en højere og inderligere \tlg{Tilegnelse} af den
til at erkjendes sig frembydende Erfaringssphære. Den opfattende Personlighed,
der derved ligesom assimilerer sig den og optager den til inderlig Eenhed med sin
egen (fornuftige) Leven, bliver derved ligesom hjemme deri og raadig derover, saa
at den bevæger sig og gaaer omkring deri som i sin egen Ejendom, hvori intet mere
er den fremmed. Men i og ved den samme \tlg{Sigtilegnen} bliver denne Personlighed mere
og mere taget i Besiddelse af den erkjendte Sandhed, saa at den personlige Leven
gaaer op i denne Sandheds egen Sigfremstillen, og mere og mere røres og bevæger
sig i Sandheden og efter Sandheden.
\end{quote}
\dcite{Sibbern, Om Erkjendelse og Granskning, p. 71  [sibbern/erkjendelse:120]}
% sibbern/erkjendelse:120  (Om Erkjendelse og Granskning; p. 71)

\begin{quote}
2. Al ægte og levende Kjerlighed til Ideen og Glæde ved samme føder af sig en
Interesse og Glæde ved Detaillet, en Lyst til at forfølge det Enkelte
indtil dets mindste Dele eller særligste Modificationer, for ret tilfulde at skue
og \tlg{tilegne} sig det intelligente igjennem det empiriske, det almene igjennem det
specielle; ja denne Lyst til at forfølge det Enkelte hører væsentligen med til at
vidne om et sandt Kald til en vis Videnskab; hvorimod det at ringeagte Detaillet
vidner om Mangel af ægte Sands og Følelse for Ideen, til hvis Væsen det netop
hører at træde frem som det altbestemmende i sin Mangfoldighed, og hvis Kraft kun
lader sig tilsyne i dens Virkninger, saa at vi kun kunne knytte os til den og skue
den, ved at knytte os til det, som bærer dens Præg og Udtryk, og skue den i dette.
Naar vi spørge os, hvorledes dog en genial Maler af den nederlandske Skole har
kunnet finde Glæde ved saa troligen at iagttage og studere Lysets Reflectioner paa
en Stol eller Kobberkjædel under visse Belysninger, for troligen at kunne gjengive
den i sit Malerie, saa ligger Svaret nær: netop Lyset, det almene alt oplivende
Lys, glædede det ham at skue og forfølge i de utallige Nuancer af dets
Reflectioner paa Gjenstandene; det almene Lys var det, der oplivede ham i dette
Studium, og var det, hvilket han vilde fremstille i nogle af dets utallige
Virkninger. Og naar saaledes et tænkende Hoved med Sands for det højeste ej
sjeldent har indskrænket sig til en speciel Gjenstands (et Insects eller en
Plantearts) nøjagtige og trolige Studium, saa var det, fordi det Hele var ham
magtpaaliggende og fordi han saae, at det for dette Heles Skyld var nødvendigt
ogsaa at studere denne enkelte Gjenstand. Og naar han studerede den enkelte
Gjenstand med den Aand og Tendents deri at levere et Bidrag til det højeste og
almeneste, kom han derved ogsaa selv til at studere det almeneste i den enkelte
Gjenstand og derunder at føle sig i levende Berøring med den hele almene,
altomfattende, alt i sig bærende Natur, og kunde nu ogsaa i den enkelte Gjenstands
grundige Fremstilling levere noget, der, ligesom det ej kunde komme i Stand uden
et mangfoldigt Henblik til den hele øvrige Natur, nu ogsaa kunde kaste et Lys hen
over en vidtomfattende Kreds. Ogsaa det meest sig begrændsende Studium, som vi i
Anm. til §. 1. allerede have sagt, har sin Muse.
\end{quote}
\dcite{Sibbern, Om Erkjendelse og Granskning, pp. 151–153  [sibbern/erkjendelse:199]}
% sibbern/erkjendelse:199  (Om Erkjendelse og Granskning; pp. 151–153)

\begin{quote}
Hvad i Særdeleshed den til al ægte Videnskabelighed fornødne philosophiske Grunden
angaaer, saa har enhver, længe førend han kan være moden til philosophiskt
Studium, allerede sat sig fast i en Maade at tage og skue Livet paa, hvilken
Philosophien ej kan lade blive staaende, han har bortgivet Realitetens Navn til
det, der ingen sand Realitet har. En anden Forestillingsmaade end den
philosophiske holder sig ogsaa hele Livet igjennem fast ved Siden af denne,
omtrent som vi blive ved at tale, som om Solen gik omkring Jorden, saa vel vi end
vide, at dette ej er saa, og saa let den øvede kan i ethvert Øjeblik sætte sig over
paa det rette Standpunct. Allerede derfor er det intet ringe, der skal skee, naar
Philosophien skal komme til et Individ og blive levende i dette; thi dets hele
Forestillingskreds skal løsne sig og lige som smelte hen, for at undergaae en ny
Dannelsesproces og at gjenfødes. Dog har det, naar Tiden dertil kommer, ikke saa
stor Vanskelighed i denne Henseende hos den lærelystne Studerende, thi netop en
Længsel efter Ideen vaagner i den Tid, da han er moden til det akademiske Studium,
lige med hans Stræben efter Selvstændighed, og han er let beredt til at aabne sig
for det, der først skal bringe det rette Lys og det rette Liv ind i alt, hvad han
hidtil lærte. En Tid lang giver han sig villig hen til Ideens Betragtning. Under
den Form, den bydes ham, tager han aabent imod den, og den højere Selvtænkning
gaaer i eet med en højere Læren. Det ham saaledes paa en livfuld Maade, igjennem
en fri Selvtænkning til hans egen Selvtænkning, overleverede tænker han sig let
ind i, \tlg{tilegner} sig det, og glæder sig i det og ved det; fører det ogsaa i mangt
et enkelt Punct videre paa en selvstændig Maade, og knytter dertil mangen egen
Besvarelse af selvopkastede Spørsmaal; det synes nu ganske at være assimileret.
Men lever en dybere Drivt til Philosophie og ægte Videnskab i ham, saa vil det med
Tiden vise sig, at \tlg{Tilegnelsen} dermed ej var til Ende. Det saaledes optagne vil
efterhaanden komme op igjen, gjennemarbejdes endnu engang for nu først under en
egentlig og fuldkommen Gjenfødelse at assimileres tilfulde og, (for at bruge endnu
et physiologiskt Ord) optages i, eller gaae over i, Constitution; alle eller dog en
Deel af de Hovedspørsmaal, der syntes vel besvarede og bragte til Ro, især det til
Grunden for dem alle liggende Grundspørsmaal, komme nu først saaledes op, at
Braaden deri eller Veerne derved (ὠδίς, for at tage et Udtryk af det andet af de
platoniske Breve) blive fornumne; og nu spørges der først egentligen ret af
Hjertens Grund. Hvad vi forhen fattede og oversaae i dets hele Sammenhæng og
derfor erkjendte i denne dets Sammenhæng og Conseqvents, og hvori vi saaledes
fandt det, vi søgte, Ideen og Grunderkjendelsen, at vi villig optoge det, dette see
vi nu i et andet Lys. Hvad den allerede modnede Tænker maaskee strax havde fundet,
naar en saadan Lære var bleven bragt ham, at den, sin maaskee gode, velgrundede og
meget omfattende Sammenhæng uagtet, dog enten aldeles ikke eller ej tilstrækkeligen
havde iagttaget dette eller hiint, der ellers var tænkt i denne Videnskab eller dog
kunde være bleven tænkt paa, dette vil nu og under sit Studiums Fremgang den
Studerende kunne bemærke om det, han tidligere beredvillig lod finde Indgang hos
sig. Fuldkomment er jo intet, der som videnskabeligt System overleveres, lige saa
lidet hvad et alsidigt Hensyn og Omblik angaaer, som ellers. Endog en vis
umiddelbar, vistnok ej altid ubedragelig, Sandhedsfølelse, som, uden at
ugjendrevne Indvendinger endnu have dannet sig, kan stræbe det modtagne eller dog
den Maade, hvorpaa det er udviklet, imod, kommer nu i Betragtning med. Hos den, hos
hvem en saadan Følelse eller Fornemmelse modstræber, er den noget reelt, der, lige
saa vel som bestemte Indvendinger, skal klares, og ikke tør tilsidesættes. Idet
man skal føres ind i et ellers dybtgaaende og vel sammenhængende Tankehele, maa
man vistnok ikke lade deslige modstræbende Fornemmelser, ja ej engang mange
bestemte, fra den hidtil havte Forestillingsmaade hentede eller øjenblikkeligen
opstaaende, Indvendinger faae en betydelig eller endog forstyrrende Indflydelse.
Der er mangen ny Lære, mangen ny Tanke eller ny Fremstillingsmaade, hvorom der med
Føje kan siges: Prøver den practiskt, optager og sætter Eder ind i den, og I
skulle vel forfare, at den er af Sandheden. Men har nogen just optaget den saaledes
og levet med den, og er han netop derved bleven modnet til at stille sig imod den
(som man da overhovedet kun ved at erkjende, bliver ret i Stand til paa den rette
Maade at bestride, det, der ellers har nogen sand Gehalt), saa kan han ogsaa give
en modstræbende Følelse Magt; maaskee deri noget anmelder sig, der igjennem ham
vil bryde en ny Bane, aabne Indgangen til en ny Erkjendelse.
\end{quote}
\dcite{Sibbern, Om Erkjendelse og Granskning, pp. 184–188  [sibbern/erkjendelse:239]}
% sibbern/erkjendelse:239  (Om Erkjendelse og Granskning; pp. 184–188)

\begin{quote}
4. Men naar vi da nu saaledes see den Fare eller vovelige Prøve, man i det
philosophiske Studium har at gjennemgaae, saa tør vi ej fordølge for os, at det
just ikke kan være raadeligt for enhver at føres ind i Philosophiens Dyb, og at vi
ingenlunde tør undre os over, om selv Mænd af fortrinlig Aandsgave ej sjeldent have
Utilbøjelighed til at indlade sig med den i den abstracte Form, der er den
Philosophien egentligen tilkommende Form. Et udmærket Talent til Konst eller til en
vis Art af practisk Virksomhed, ja til visse Arter af videnskabelige Granskninger,
kan, naar Reflection og indtrængende Grunden ikke naae Dybden af det, man i sit
Talent besidder, letteligen forvirres og forstyrres eller standses ved
philosophiske Speculationer. Philosophering i videnskabelig Form er ej heller
hverken den eneste eller den højeste Maade, paa hvilken vi kunne stræbe at \tlg{tilegne}
os og besidde Ideen og den aandelige Gehalt, Ideen giver Livet. Fromhed,
Kjerlighed og Dyd ere højere end Philosophie, og (Poesien kan stille sig ved Siden
af den). Eet er det, under en levende Receptivitet og Productivitet at besjæles af
Idee, et andet, at stræbe at udgrunde Beskaffenheden af det, der saaledes ligger
til Grund for den kraftige Virksomhed, til hvis Udøvelse vi føle Kald og Gave. Med
Talent til Virksomheden selv er just Talent til at udgrunde dens indre Natur og
Væsen ej altid, ja ej letteligen forbunden. Men saa kan den umiddelbare kraftige
Udøvelse let gaae tabt over Bestræbelsen efter, philosopherende at trænge ind i
det, der constituerer den. Bekjendte ere Bacos Ord, at Philosophien, naar man blot
kommer halvt ind i den, fører fra Gud, og kun, naar den grundig studeres, fører
tilbage igjen til Gud. Og saaledes vil man og ved et ikke noksom til Grunden
gaaende Studium af sin egen Levens ideelle Gehalt komme bort fra denne. Ideen, man
formedelst abstract Tænkning vil stræbe at skue klarere, som den i og for sig er,
kan ikke vedblive at leve i os, som den forhen levede i os, som den ikke særligen
til Syne kommende Sjæl i vor practiske Virksomhed, og det gjelder nu om den
Tænkning, der har gjort den til Gjenstand for indtrængende Reflectioner, kan
igjenvinde den under den philosophiske Form, og paa ny bringe den til den
umiddelbare kraftige Leven, eller om denne midt under hiin Tænkning kan bevare sig
ved Siden deraf.
\end{quote}
\dcite{Sibbern, Om Erkjendelse og Granskning, pp. 192–193  [sibbern/erkjendelse:245]}
% sibbern/erkjendelse:245  (Om Erkjendelse og Granskning; pp. 192–193)

\begin{quote}
3. Hensynet til en Praxis, hvortil man vil gjøre sig duelig er vistnok allerede
under Studiet forsaavidt af megen Vigtighed, som det kan gjøre opmærksom paa alt,
hvad der engang i Tiden vil kunne udfordres dertil og nu allerede maa omfattes med
af Studiet; som da i alle Hovedfag egentligen Ideen om et bestemt Kald, en bestemt
i Livet indgribende Virken, leder til at samle en Mængde ellers adskilte
Granskninger omkring et Middelpunct. Men ikke maa dette Hensyn hilde saaledes, at
noget, som Videnskabens Idee fordrer med til sin totale Udførelse til et
systematisk Hele, skulde tilsidesættes, fordi det ej synes vigtigt for det senere
practiske Liv, eller at overhovedet det Spørsmaal, hvorvidt det, Videnskabens Gang
fører til, er fornødent for det practiske Liv, skulde komme frem med Hensyn til de
enkelte Puncter. Den bedste Beredelse til Praxis er et grundigt Studium af
Videnskaben selv saaledes, som den er; thi til en solid Praxis udfordres at være
aldeles hjemme i Videnskaben og at have \tlg{tilegnet} sig dens Aand, dens Principer og
Methoder, men dertil kommer kun den, som studerer Videnskaben i dens Heelhed,
tilmed dens Principer og Methoder i deres hele Udførelse og forfølger dem i alle
deres Hovedgrene eller Retninger. Men ligesom Hensynet til Praxis ej bør betinge
eller hilde, saa bør det ej heller bekymre eller ængste, saa at vi ved ethvert
enkelt Punct skulde sørge for, hvorledes vi skulde komme udaf det dermed, naar
engang Regnskabens Dag kommer. Det første er at trænge ind til Aanden og de
Grunderkjendelser, hvorpaa alt kommer an; ere vi grebne og opfyldte heraf og have
ganske \tlg{tilegnet} os disse, saa vil ogsaa Detaillets mangfoldige Erkjendelse komme
efterhaanden. Ogsaa i Studiet maae vi have for Øjne, hvad der staaer skrevet: Søger
først Guds Rige og hans Retfærdighed, saa skulle og alle disse Ting tillægges
Eder. Endog hvad den practiske Udøvelse angaaer, saa er vistnok nøjagtig Kundskab
om et mangfoldigt Detail fornøden, og det væsentligen fornøden dertil, men det
rette practiske er dog Besiddelsen af det, hvorfra dette Details sande Erkjendelse
gaaer ud. Hvor lidet gives der vel af den egentlige, dybere gaaende Philosophie,
som kan anvendes ligefrem i en ydre Virkekreds? Og dog er, næst Religionen, intet
mere practisk end Philosophien, der, idet den bringer det, hvorpaa hele Livet
hviler, til fuldkommen Erkjendelse og Klarhed, danner til overalt at lægge de rette
Principer, den rette Anskuelsesmaade, til Grund for sine Domme, og bestandigen at
have Hensyn til det, hvorpaa alt kommer an.
\end{quote}
\dcite{Sibbern, Om Erkjendelse og Granskning, pp. 199–201  [sibbern/erkjendelse:253]}
% sibbern/erkjendelse:253  (Om Erkjendelse og Granskning; pp. 199–201)

\subsection{\textit{Udaf Gabrielis' Breve til og fra Hjemmet} (1826)}

\begin{quote}
„Jeg tænkte paa de Dage fra fordum og de Aar, som ere henrundne“ —
hvor ofte ere ikke disse Assaphs Ord vendte tilbage i min Erindring,
fra den Tid af, da jeg i Aaret 1812 paa Aarets sidste Dag i
Aftendæmringen hørte en Prædiken over dem i den store Domkirke i
Breslau, hvor jeg dengang opholdt mig, for at være i Nærheden af
Steffens. Men især vendte jeg ofte tilbage til disse Ord fra den Tid
af, da jeg begyndte at erholde og samle de Breve af Gabrielis, som jeg
her meddeler, og tænkte paa, at, naar jeg en Gang kom til at udgive
dem, maatte jeg og vilde jeg lade dem ledsages af en \tlg{Tilegnelse} til
Dem. Det er nu 25 Aar siden, at jeg sendte Dem „Gabrielis' efterladte
Breve“ i Haandskrivten, og fra Dem med stor Glæde modtog den Recension
deraf, som blev trykket foran dem. Da vare Tiderne ganske andre, end
de nu ere, men ogsaa da stredes og kæmpedes der i Gemyterne og af
Gemyterne i Forvirring og Uro; det var i kirkelig Henseende, at der
blev stridt, og jeg udgav netop hine Breve (som da allerede, paa nogle
faa nær, havde ligget hos mig i over 9 Aar), fordi jeg forestillede
mig, at de maaskee kunde virke til, at Tankerne hos os i de offentlige
Omtalers Gebeet, kunde vende sig fra Spørgsmaalet om de rette første
Kilder til Livets helligste Vande, hen til Omsorgen for den rette
Gjenfødelse og Qvægelse ved disse Vande selv. Hine Breve lade os see
Gabrielis i en Krise, og det er kun denne Krise, de føre frem for os,
ladende Enhver i egne Tanker forfølge den til dens Opløsning eller
Klarelse; men Tilværelsens hellige Vande rinde og raade deri. Jeg
havde efter hiin Tid ofte det Haab, engang at skulle kunne knytte
tvende Samlinger af Breve til hiin første, og at deraf den ene skulde
kunne kaldes „Gabrielis' Breve til Hjemmet“, den anden, hans „Breve fra
Hjemmet“. Men kun fra nogle faa Uger, men vistnok for ham indholdsrige
og betydningsfulde Uger, fra en Tid, som falder 10 Aar efter hiin Tid i
hans Liv, har jeg kunnet komme i Besiddelse af Breve fra hans Haand, og
dem meddeler jeg nu.
\end{quote}
\dcite{Sibbern, Udaf Gabrielis' Breve til og fra Hjemmet  [sibbern/gabrielis:3]}
% sibbern/gabrielis:3  (Udaf Gabrielis' Breve til og fra Hjemmet)

\subsection{\textit{Bemærkninger og Undersøgelser, fornemmelig betreffende Hegels Philosophie} (1838)}

\begin{quote}
Man seer let, at det med Udødeligheden i et heelt, gjennemført philosophiskt
System, til hvis Verdensanskuelse den hører væsentligen med, ikke kan
forholde sig saa, at den mueligen kunde være bleven forglemt af dets
Ophavsmand, og da kunde af Tilhængere føies til og knyttes an eet eller
andet Sted, som Noget, hiin ei var kommen til at omtale. Den udgjør —
især naar en saadan Philosophie vil sætte sig i et bestemt Forhold til
Christendommen, ja vil indoptage dettes væsentligste Grundlag i sig — et
ei blot saa væsentligt, men og saa indgribende Punct, at den allerede i et
saadant Systems Anthropologie, end sige i dets Psychologie, og fremdeles i
dets hele Moralphilosophie, maa bestemt komme tilsyne, naar alt Dette skal
ville \tlg{tilegne} sig Navn af speculativt, og naar det endog tager saadanne
Lærdomme med, som den om den dyriske Magnetisme, for at haandhæve, og ved
speculative Anskuelser at oplyse, Gyldigheden af at antage en saadan. Men
især gjelder Dette om en Philosophie, som baade i sin naturphilosophiske og
i sin aandsphilosophiske Deel just især gaaer Tilværelsens opadstigende
Trin og Stadier igjennem, derover endog deels forglemmende, deels
tilsidesættende de allervigtigste Betragtninger med Hensyn til Det, der paa
collateral Maade og i collateral Henseende gjør sig gjeldende i Liv og
Tilvær. Thi, naar i en saadan Philosophie Erkjendelsen af Udødelighedens 
Realitet virkelig har Plads, hvorledes kan da den trinmæssige Betragtning
aldeles herske frem deri, uden at Menneskelivet her paa Jorden bliver
fremstillet som et Livets Trin, der viser ud over sig selv til et høiere
Trin paa Individualitetens og Personlighedens, saavelsom paa det i
Personlighederne opgaaende og sig deri indlevende aandelige Riges,
Udviklings-Scala? Hvorledes vil det da kunne undgaae Opmærksomheden, at
Livsrækken i Mennesket først i Udødeligheden finder sin Integration? Dog vi
maae rigtignok bemærke, at, saa meget end den opadstigende Idee-Udviklings-gang
hos Hegel hersker, er det egentlige Historiske i Liv og Tilværelse dog hos
ham ikke kommen til at blive gjeldende, saaledes som f. Ex. hos Steffens,
der seer hele Livet historiskt. — Som det hos et menneskeligt Individ maa
gjøre en Hovedforskjel, hvad enten det tager sit Liv som det, hvis hele
Fylde, Formaal og Sphære er allene at søge paa Jorden (saa at det
forsaavidt om samme maa hedde: qvod petis hic est), eller om det
tager Livet som det, der her paa Jorden skal modnes en anden Verden imøde,
saa at det veed sig her paa Jorden endnu kun at befinde sig i Propylæerne,
ei i Templet, saaledes maa det og i et philosophiskt System strax i dets
hele Totalitet vise sig, om Udødeligheden er indoptaget deri. Hos Leibnitz,
Kant, Jacobi, ogsaa Fichte den Ældre, behøve vi ei at spørge. At
Udødelighedstroen er heelt gjennemgribende i Christendommen, er bekjendt
nok. Den ligger ogsaa til Grund for den hele romantiske Poesie. Det hører
til Christendommens store Virkninger, at den har ladet den Tanke blive
mægtig i Bevidsthederne, at vi her endnu kun bære det Uforkrænkelige i
Forkrænkelighed, men at det i os liggende aandelige Foster skal finde den
Dag, da det bryder ud af dette Leie, river sin Navlestreng over, med det
hele derhen hørende Apparat af Organer for Ernæring og Respiration, og
gjennem saadant Brud og saadan Mellembød kommer ud til det Dagens Lys, til
det Element og den friere Verden, som det er bestemt for, hvor det skal
aande og inddrage Næring gjennem andre Organer og i andre Sphærer, end
hidtil. Til de stærkeste, af Aanden meest opfyldte, varmeste og
henrykkelsesfuldeste Steder i det nye Testament høre Pauli begeistrende og
dybt indlysende Ord herom i 1 Cor. 15, hvilket jeg ei skulde bemærke,
dersom jeg ei vidste, hvor lidet kjendt det nye Testament, som den hele
Bibel, i vor Tid er hos mange Dannede. At nu alt Herhenhørende i den
Hegelske Lære enten reent fattes, eller træder saa tilbage, at man har
kunnet strides om, om Noget deraf er der, eller ikke, det maa man dog
tilstaae, at være i høieste Grad paafaldende, naar man betænker, at denne
samme Philosophie skal gjelde for at have ført Philosophien lige op til
Christendommens Standpunct, at have gjort Christendommen speculativ, og at
have givet den christelige Tænkning en Grundvold, ved hvilken endog
Theologer anmodes om at holde fast.
\end{quote}
\dcite{Sibbern, Bemærkninger og Undersøgelser, fornemmelig betreffende Hegels Philosophie, pp. 12–14  [sibbern/bemaerkninger:16]}
% sibbern/bemaerkninger:16  (Bemærkninger og Undersøgelser, fornemmelig betreffende Hegels Philosophie; pp. 12–14)

\begin{quote}
Jeg har her berørt adskillige Hovedpuncter hos Hegel, som noksom maae vise,
hvorledes det staaer sig med Hegels Forhold til Christendommen, og tilmed
til de Fordringer, der i en Tid, i hvilken der er kommen nyt friskt Røre i
den christelige Ideekredses Dybder, maae gjøres til den Philosophie, der
skal kunne tilfredsstille Tiden. Men jeg berører her ogsaa Noget, hvori jeg
ei har kunnet Andet, end see talende Beviser for, at Hegel selv, hvis han
nu kunde træde frem i sin Manddoms bedste Kraft, og da havde sin egen hele
Philosophie for sig, umueligen kunde ville blive staaende ved den, men
maatte føle sig dreven til at gaae ud over den paa mange Maader, især om han
tillige kunde træde frem friet for denne fatale og kjedsommelige Polemik, i
hvis Vold vi see ham, og som ei blot har havt en høist forstyrrende
Indvirkning paa hans Systems hele Eurythmie, men endog, ved Siden af hans
haardnakne Fastholden ved visse, hans Dannelsestid tilhørende,
Eensidigheder, synes at have grebet forvirrende ind i hans hele Systems
inderste Constitution. Men nu ei at ville gaae ud over denne Philosophie
— nu, da den ligger heelt for os, og da der har været Tid nok til
Sundelse og til friere Omblik i Verden — hvad skal man dertil sige? Paa
mange Maader er det mig klart, at denne hele Lære af en grundig Hegelianer
maatte kunne tages saaledes, at Hegel vilde komme til at staae med
Betydningen af at danne Mellemleddet i et Slags stor Hegelsk Trilogie,
uagtet hans Philosophie i sig selv gaaer ud paa at afslutte en Æra, nemlig
den hele Schellingske Periodes Æra*) [note: Jeg undlader ei at bemærke,
hvad der vel ellers maa forstaae sig af sig selv, at jeg, naar jeg her
nævner Schelling, bestandig kun har Schelling, som han da var, for Øie.
Schelling, som han nu er, har jeg her slet ikke at tage Hensyn til, og jeg
har ei heller her kunnet tage Hensyn til ham, som han nu er. Selv har han jo
i mange Aar saa godt som Intet leveret os.]. Thi det maa vel erkjendes, at
den Hegelske Philosophie i Grunden kun gaaer ud paa en Udførelse,
Gjennemførelse og Fuldendelse — og det endda kun i een særdeles Retning
— af Det, hvad der i den rige Schellingske Tids Foraarsperiode begyndte at
udfolde sig. I denne ene Retning har nu forresten Hegel gjennemført det med
stor Originalitet, og saaledes, at man vilde lade det Hegelske Genie
vederfares alt for liden Ret, naar man ikke heri vilde erkjende den
grandiose Tænkers mægtige Selvstændighed og store Tankestyrke, samt rige
Fylde af gehaltfuld Bemærkning. Hvortil da kommer, at en Ontologie for denne
Philosophie har Hegel først skabt. Blandt de Lignelser man, deels har
anvendt, deels kan anvende, for at betegne Hegels Forhold til hans nærmeste
Forgjængere (til hvilke Schelling ei ene hører, men den hele Schellingske
Periodes Mænd), turde den være den meest treffende, at han staaer i samme
Forhold til dem, som Spinoza til Cartesius og Cartesianerne, medens der ved
Siden af Spinoza var Plads og Rum nok for en Malebranche, og snart efter
maatte opstaae en Leibnitz, kun at Hegels Værk er saa meget rigere og mere
speculativt tilfredsstillende, som vor Tid er rigere og mere speculativ, end
Spinozas. Men allerede deraf maa, for ei at tale om hine store, Hegels
Person tilhørende Eensidigheder, kunne skjønnes, at Hegel ei kan komme til,
saa stærkt at gribe ind i al Videnskabeligheds Rige, som Schelling gjorde
det, eller som Kant tidligere greb ind, skjønt Hegel har grebet meget mere
ind, end Spinoza gjorde det i sin Tid, eller end Fichte den Ældre i sin Tid.
Philosophien vil især kun i de Perioder kunne gribe bevægende ind i Livet, i
hvilke den undergaaer sine Metamorphoser, og udfolder sig i nye Skikkelser
under en Stræben efter nye Løsninger af de philosophiske Grundproblemer,
idet der paa Ny kommer Røre i disse. Der kom et saadant nyt Røre i dem i den
kantisk-fichtisk-schellingske Tid; men egentligen er intet af dem paa en
gjennemgribende Maade discuteret paa Ny af Hegel, der altfor umiddelbart og
thetiskt har lagt sine Grundanskuelser til Grund, og forresten har ladet
saadanne Problemer, som Spørgsmaalet om Sjælens Forhold til
Legemet*) [note: Hvor han berører dette Forhold (i Anm. til § 389 i
Encyklop.), siger han selv Intet til dets Opklaring, og hvad han herved
siger med Hensyn til de Philosopher, han her nævner, maa, især hvad
Cartesius og Leibnitz angaaer, forundre hver Den, som kjender disse
Tænkeres Philosophie. Men det er af flere Steder hos ham klart, at han ei
noksom kjender dertil.] , om Sjælens Udødelighed, om Guds Personlighed
o. fl., næsten urørte, ja ei engang har ladet Spørgsmaalet om
Philosophiens reale Mulighed, dens Natur og Væsen, samt om dens Forhold til
Empirien paa den ene, til Troen paa den anden Side, komme til nogen sand
indtrængende og gjennemtrængende Discussion. Naar en vis Tids Philosophie
begynder at sætte sig, for at udorganisere sig efter en vis bestemt
Grundtypus, og derover imidlertid en heel Menneskealder gaaer hen, vil dens
Bevægelsesperiode være forbi. Impulserne have da allerede virket, og
imidlertid har ogsaa Observationsaanden grebet om sig i Livets og i de
observerende Videnskabers Gebeet (men observerende ere de i Grunde alle),
saa at Forraad samle sig til Bedste for en ny Metamorphose i Philosophiens
Rige. Hos Hegel selv er et stort Forraad af slige Observationer at finde,
skjønt vistnok gjennemblandede med mange Konstlerier, og hans nærmeste
Virken for Fremtiden turde beroe paa, hvad Fremtidens philosophiske
Observatorer ville finde at \tlg{tilegne} sig og vide at føre sig til Nytte af
hans rige Forraad.
\end{quote}
\dcite{Sibbern, Bemærkninger og Undersøgelser, fornemmelig betreffende Hegels Philosophie, pp. 15–18  [sibbern/bemaerkninger:19]}
% sibbern/bemaerkninger:19  (Bemærkninger og Undersøgelser, fornemmelig betreffende Hegels Philosophie; pp. 15–18)

\begin{quote}
Hvad nu især den Hegelske Philosophies Forhold til Christendommen angaaer,
— og derpaa maae vi da her især have Opmærksomheden henvendt — saa er
vel nu allerede det Tidspunct indtraadt, da den Erkjendelse kan siges
noksom at have gjort sig gjeldende hos Mange, som ellers have følt sig
hentrukne til Hegel, at den Hegelske Christendomsphilosophie ikke svarer til
den ved Hegel selv til al Lære gjorte Hovedfordring, at den skal være
immanent, eller at den skal være Sagens egen sig udviklende Gang; hvoraf da
dog vel maa blive en Følge, at udaf Sagens eget Indre Videnskabens bevægende
Grundidee og det hele deri sig rørende Begreb maa gaae frem. Kun ved en sand
levende Vexelvirkning imellem den philosophiske Speculations Aand paa den
ene Side, og det givne Indbegreb, den givne Fylde af christelig Tro og
Erkjendelse paa den anden Side, kan Grundlaget for, og det sig heelt
igjennem Bevægende i, en Christendomsphilosophie hæve og klare sig frem. Kun
da kommer Christendomsphilosophien frem som et nyt Væsen, der har baade
Fader og Moder, og gjenudtrykker begges Naturer i en inderlig og væsentlig
Eenhed. Men den Hegelske Christendomsphilosophie, som vistnok er aprioriskt,
er ingenlunde immanent paa denne Maade. Den har ingenlunde udviklet sig udaf
en grundig philosophisk Fordybelse i selve Christendommen. Dens Forhold til
Læren om Udødeligheden, dens Trinitetslæres egentlige Beskaffenhed, hvorom
længere hen, saavelsom dens Miskjendelse af Følelsens og Troens egentlige
Væsen, samt det Suprematie, den, ligesom Fichte den Ældre, vil tillægge den
philosophiske Erkjendelse over Religionen, kan noksom lære Dette. Ikke at
tale om, at hos Hegel det hele Moralske ei bygges paa Religion, et Punct,
hvori vi atter see hele Tiden kommen ud over det Hegelske Stadium, skjønt
den hele Hegelske Philosophies besynderlige inverse Gang, som jeg pleier
kalde den, her har gjort Sit til at føre til denne store Ulempe. Som en
allerede færdig Philosophie træder den Hegelske Philosophie hen til
Christendommen, for nu og at \tlg{tilegne} sig den, der da maa lade sig arrangere
efter Principer, ved hvis Udvikling og Etablering den ikke retteligen og
fuldeligen har været paa Raad med. Christendommen har ei fra først af været
tilstede ved, og virksomt grebet ind i den Debat, hvoraf den hegelske
Philosophie har hævet sig frem. Spinozistiske Ideer have ganske anderledes
— Enhver, som kjender til den ældre Schellingske Philosophie, og seer sig
om hos Hegel, maa kunne see det — været i Bevægelse, da den Hegelske
Philosophie constituerede sig sit Grundlag. Med sine færdige
Grundanskuelser træder da Hegelianismen hen til det Christelige. Paa saadan
Maade kommer da Christendommen kun til at tjene den Hegelske Philosophie til
Føde; en paa Hegels Viis ontologiseret eller logisticeret Christendomslære
bliver da Resultatet, ingen sand speculativ christelig Theologie. At mange
enkelte Bemærkninger over det Christelige af immanent Udspring hos Hegel
kunne forekomme, skal ikke nægtes; jeg har allerede sagt, at der hos Hegel
er ikke Lidet, som peger hen ud over ham; men et forunderligt exoteriskt
Udseende faae ofte slige Yttringer hos ham. Men det maa herved ei heller
glemmes, at Hegels Philosophie hos ham selv allerede ganske havde
consolideret sig, eller fastnet sig, da de nye Bevægelser, det nye Røre fra
Grunden af, i de christne Erkjendelsers Verden brød frem med al Styrke.
\end{quote}
\dcite{Sibbern, Bemærkninger og Undersøgelser, fornemmelig betreffende Hegels Philosophie, pp. 18–20  [sibbern/bemaerkninger:20]}
% sibbern/bemaerkninger:20  (Bemærkninger og Undersøgelser, fornemmelig betreffende Hegels Philosophie; pp. 18–20)

\begin{quote}
Selve den Hegelske Philosophie indeholder Nok til at være en Stagnation i
denne Henseende imod; den indeholder Nok, der peger ud over dens eget
Stadium, Nok for at afholde fra den Tanke, at her kunde man nu indtil videre
slaae sig til Ro, og give sig hen i Hvilens Fylde. At man ret snart bliver
Dette almindeligen vaer, er ogsaa for de Hegelske Bestræbelsers egen Skyld
saare ønskeligt. Man har, hvad Virkningen af den Hegelske Philosophies
Studium i Tydskland angaaer, villet bemærke, at det synes at Hegel blandt
sine Modstandere besidder nogle af sine bedste Læsere, hos hvilke hans
Philosophie meest har baaret Frugt. Dette kan ei forundre, og det er for
dens egen Skyld kun at ønske, at det underlige unaturlige Forhold maa
ophøre, ifølge hvilket der har dannet sig en Tilstand, at det seer ud, som
om det kom an paa, enten at bekæmpe den, som en Philosophie, man aldeles
maatte lade falde, eller at forfægte den, som en Philosophie, der i Eet og
Alt var saa tilfredsstillende, at det kun kom an paa at udbrede den og gjøre
den gjeldende overalt. Saa synes at blive vaer, at, for retteligen at nytte
den her meddeelte store Gehalt, ja jeg kan sige, for tilbørligen at lade den
vederfares Ret, ved at lade den vorde saa frugtbringende, som den kan blive
det, maa man netop gaae ud over den. Istedetfor at optage den som et stort,
høist priisværdigt Bidrag til den philosophiske Stræbens og det
philosophiske Værks Fremme, dele sig de Philosopherende for en stor Deel,
forsaavidt de paa en livfuld Maade indlade sig med Hegel, i Modstandere, som
blot forholde sig angribende, Tilhængere, som blot forfægte og yttre sig
paaholdende, og Saadanne, som ønske at nyde Godt af hvad her er givet, uden
at ville slaae sig til noget Partie, men som dog ei kunne undgaae at drages
ind med i denne Antagonismus. Til de sidste turde nærværende Recensent vel
være at henregne. Saa at, tvertimod den hele Hegelske Philosophies egen
Synsmaade i Henseende til Det, som i Philosophiens Historie er traadt stærkt
indgribende frem, — hvorefter man ei skal gaae ud fra, at et philosophiskt
System netop enten skal antages for gyldigt, eller gjendrives som ugyldigt,
— trækkes nu Hegel selv med ind under et saadant aut-aut, maaskee
til Straf, fordi han selv med Hensyn til Andre, især Jacobi, saa lidet har
havt for Øie, hvad han (Logik III, strax i Beg., Pag. 8 i den første Udg.)
fordrer med særdeles Hensyn til Spinoza.*) [note: „Die wahrhafte
Widerlegung“ — saa hedder det her — „muß in die Kraft des Gegners
eingehen und sich in den Umkreis seiner Stärke stellen; ihn ausserhalb
seiner selbst angreifen, und da Recht zu behalten, wo er nicht ist, fördert
die Sache nicht.“ — „Die einzige Widerlegung“ — lægges der til — „des
Spinozismus kann daher nur darin bestehen, daß sein Standpunct zuerst als
wesentlich und nothwendig anerkannt werde, daß aber zweitens dieser
Standpunct aus sich selbst auf den höhern gehoben werde.“ — At saadan Art
af Modvirkning allene skulde være den rette, kan iøvrigt vist Ingen
indrømme; thi megen Modvirkning maa dog have sit Grundlag i, at man friere
og i en større Omkreds seer sig om i Livet. Men Hegel selv kunde indbyde
til, at Nogen netop behandlede ham, som han her har fordret Spinoza
behandlet, og selv har handlet med Hensyn til Spinoza (hvorved denne
forresten vistnok alt for meget er bleven ham hans egentlige Grundlag).] 
Dertil have da de altfor eensidige, blot og bart Hegelske Tilhængere, samt
tillige Hegels egne store Eensidigheder, saavelsom hans slemme,
sædvanligviis haanende, absolutistiske Maade at polemisere paa, bidraget
meget, især da denne Polemik har smittet mange af Tilhængerne, deels med den
Fornemhed, deels med den Sufficence, som Eensidighed dobbelt kommer til,
naar den, under en stor Auctoritets Beskyttelse, troer sig berettiget til
haanlig at afvise. Hvoraf Følgen atter er bleven, at Modstanderne dobbelt
ere blevne æggede, og have viist, at ei blot Kjærlighed, men ogsaa
Indignation gjør skarpseende, kun paa forskjellig Maade. — Men ligesom
neppe Den, der allene betragter Hegels Philosophie med et mod samme indtaget
Sinds Øie, vil kunne blive fuldkomment vaer, hvorledes man rettest fører ud
over den, og derved bedst overvinder den, saaledes vil ogsaa neppe Den med
nogen sand Frihed, altsaa paa den rette frugtbringende Maade, \tlg{tilegne} sig
det Meget, som her er givet, der ikke tillige kan see, at Meget af det
Betydeligste i denne Philosophie staaer i et aabenbart Misforhold, deels til
Philosophiens egentlige Væsen, deels til Livets Indhold, ogsaa
til Christendommen og hvad der i dennes Regioner i vor Tid atter
med ny Kraft saa mægtigen har gjort sig gjeldende, og bliver ved
at gjøre sig gjeldende.
\end{quote}
\dcite{Sibbern, Bemærkninger og Undersøgelser, fornemmelig betreffende Hegels Philosophie, pp. 29–32  [sibbern/bemaerkninger:30]}
% sibbern/bemaerkninger:30  (Bemærkninger og Undersøgelser, fornemmelig betreffende Hegels Philosophie; pp. 29–32)

\begin{quote}
Professor Heiberg gjør (Perseus I, Pag. 34), hvor han kommer til at
afhandle Spørgsmaalet om Philosophiens Begyndelse med Licentiat
Martensen, den Lignelse, at, ligesom Mennesker vistnok maae have
Forældre, men dog, naar de træde frem i Verden, maae leve deres
selvstændige Liv, uafhængigt af deres Forældre, saaledes maa ogsaa
Philosophien, skjønt den ei kan komme frem, uden at forudsætte Meget
foran sig, dog, naar den først træder frem, føre sit selvstændige Liv
ud af sig selv allene [note: Jeg vil hensætte Ordene, som de hos
Prof. Heiberg lyde Pag. 34 nederst: „Hvo vilde philosophere, dersom
han ikke iforveien var umiddelbart overbeviist om sin egen, om den ydre
Verdens og om Guds Tilværelse? Man kan altsaa vel indrømme, at
Philosophien forudsætter noget umiddelbart Vist og Givet, men heraf
følger ikke, at denne Omstændighed maa forandre Videnskabens Skikkelse,
eller frembringe et nyt System, der gaaer qvalitativt videre; thi hiint
Givne er ikke paa anden Maade Forudsætning for Philosophien, end f.
Ex. Forældrene for et Barn; skal dette blive Menneske, da maa det
engang begynde sit Liv saa selvstændigt, som om det stod ene i Verden;
og paa samme Maade maa Philosophien, i det Øieblik, den begynder sin
selvstændige Udvikling, begynde med at ignorere alle de Forudsætninger,
hvorfra den ikke desmindre har sin Oprindelse.“] . Jeg vil forfølge
denne Lignelse. Spørgsmaalet kan ikke være, om Børnene, det er: de nye
selvstændige Væsener, som fremkomme, skulle leve deres Liv, uafhængigt
af deres Forældre, men om de skulle leve det uafhængigt af den dem
gjennem deres Forældre, og tillige gjennem mange Andre, bibragte
Opdragelse, Dannelse og hele Tilværelsesmaade. Man kan og maa hertil
svare baade Nei og Ja. Hiint, nemlig Nei, synes simplest, da de dog ei
kunne fare ud af deres hele nu saaledes bestemte og tildannede Tilvær,
som om det gjaldt, at de skulde fare ud af deres egen Hud, for at
begynde en aldeles selvstændig Tilvær, reent ud fra en reen og bar,
endnu tom Væren. Men dog ogsaa Ja er at svare paa hiint Spørgsmaal. Thi
de skulle i Eet og Alt, for at naae sand Personlighed og
Selvstændighed, gjennem deres egen Udviklingsproces, under deres egen
Genius's Indflydelse og Styrelse, selv bekræfte sig i alt Det, hvad der
af hiint Bibragte skal vedblive, og skulle derimod hos sig selv benægte,
hvad der ei skal blive. De skulle da, idet de saaledes \tlg{tilegne} sig
deres Livs Indhold og tage sig selv i Besiddelse, tillige bestemme selv
deres hele Tilvær, og komme til at leve deres eget Liv paa Basis af
deres egen Natur og egen Frihed. De maa da saaledes gjennemgaae deres
egen selvstændige Udviklingsproces. Netop saaledes maa nu og
Philosophien det. Og derfor maa den have hiin sin Indledning, som
væsentligen første Stadium i dens Organisations Gang. Men vilde den
begynde med at ignorere alle de Forudsætninger, hvorfra dens
Udviklingsgang har taget sin Begyndelse, saa maatte den jo ogsaa
ignorere Det, den gjennem denne Udviklingsgang selv har constitueret
sig til at være — thi Dette gaaer i Eet med denne Constitutionsproces
selv. Den vilde da ligne det Menneske, der, for ret at begynde sit Liv
selvstændigt, virkelig vilde begynde det, „som om det stod ene i
Verden“, og, istedetfor at tage sin selvstændige Udviklingsproces paa
den Maade, jeg nys omtalte, virkelig vilde ignorere alt det ved
Opdragelse og Dannelse Bibragte, tilmed ignorere sit Livs hele Indhold,
for nu reent og bart ud af sig selv allene at begynde sin Leven. Jeg
gad med Mephistopheles om et saadant Menneske sige:
\end{quote}
\dcite{Sibbern, Bemærkninger og Undersøgelser, fornemmelig betreffende Hegels Philosophie, pp. 49–51  [sibbern/bemaerkninger:45]}
% sibbern/bemaerkninger:45  (Bemærkninger og Undersøgelser, fornemmelig betreffende Hegels Philosophie; pp. 49–51)

\begin{quote}
Hvad dernæst angaaer „at sige Noget om Veien, som fører dertil,“ nemlig
til den attraaede Verdensanskuelse, da kan det ei være saa vanskeligt.
Reisen er at gjøre paa samme Maade, som dog Hegel maa have gjort sin.
Det gjelder her aandeligen at \tlg{tilegne} sig det Land, hvori man allerede
er og lever, men med det Samme at udvide sin Synskreds deri, saa at man
fuldeligere og i større Omfang kommer til at røre sig i de Egne, hvori
man allerede existerer, hvorved da ogsaa mangt et Punct i dem nu først
kan vorde opdaget, og mangen Udsigt udover dem kan aabne sig. Hegel har
dog vel ogsaa bereist Terrainet frem og tilbage, til Høire og Venstre,
saa vidt han kunde naae, eller saavidt hans Hu drog ham, forend han gik
til at give os sit Kort og sin Geographie deraf. Han selv peger i
Begyndelsen af sin Logik hen til sin tidligere „Phänomenologie des
Geistes“ i saa Henseende. Hans philosophiske System, hvor han giver os
det heelt, nemlig i hans Encyklopædie, indledes rigtignok ikke ved
noget Uddrag af af denne Phænomenologie,
saa at den Tanke, han først synes at have havt, at hans System deri
skulde finde sin Indledning, som Henledning til sit egentlige
Udgangspunct, maa antages for opgivet; ligesom hans Erklæringer i
Begyndelsen af hans Logik (begge Udgaver) vise, at han i hiint Skrivt
ei vil have givet en saadan Indledning til sit System, hvortil dette i
sin Begynden behøvede at støtte sig, eller skulde søge Hjemmelen for
sin Begynden. Jeg har forhen talt herom; men det har der maattet vise
sig, at han dog virkelig ei har kunnet komme til sin Philosophie ad
anden Vei, end gjennem en foreløbig Gjennemgaaelse af Alt, hvad der
skulde optages deri, eller igjennem et Slags forudgaaende almindeligt
Overslag. — Fichte den Ældre var i saa Henseende klar nok over sig
selv til at erklære, — i hans første Tid nemlig — at han havde
forvisset sig om sin Veis Retning lige hen til Maalet forud, men at han
kun kunde indbyde til at følge med ham, da Udfaldet nok vilde lære, om
han havde begyndt retteligen. Det er Det, Hegel kalder at begynde
„bittweise“; og saaledes vil Hegel ei begynde, men absolut, hvilket dog
i Grunden bliver at begynde postulatviis. Fichte indsaae, at sin fulde
Retfærdiggjørelse kunde hans hele Philosophie i hvert af dens Puncter,
tilmed og i Henseende til dens Begynden og dens Gang, ikkun faae ved
sin Gjennemførelse og den Samstemming, som da vilde vise sig med det i
Erfaringen Givne. Senere gik Fichte saaledes tilværks, at han
erklærede, at, som enhver Videnskab havde sin Gjenstand, den søgte at
opklare, saaledes havde ogsaa Philosophien sin; denne var den hele
menneskelige Viden, som Viden, i sin Heelhed — das Wißthum, som han
han kaldte den,
— eller das Wissen, als Wissenschaft. Deraf det Navn Videnskabslære
(Wissenschaftslehre, Wißthumslehre), som han gav Philosophien. Men nu
begyndte han da med, først at exponere eller beskrive denne Gjenstand,
saaledes som den viser sig for Iagttageren, forend han gik til at
udgrunde den i dens egentlige Grund og Væsen. Saaledes kom han til at
adskille en blot explicativ Deel af sin Philosophie, som den blot
forberedende Deel, hvilken han kaldte „die Lehre von den Thatsachen des
Bewußtseyns“, fra den egentlige Philosophie, eller den egentlige
Wissenschaftslehre. — Den Gang, Philosopheringen har at gaae, for at
komme til et Grundlag, der nu retteligen kan gjøre sig gjeldende som
Grundlag, seer man saaledes hos Fichte anerkjendt. Og Hegel er i
Gjerningen paa sin Viis gaaet en lignende Gang.
\end{quote}
\dcite{Sibbern, Bemærkninger og Undersøgelser, fornemmelig betreffende Hegels Philosophie, pp. 56–57  [sibbern/bemaerkninger:56]}
% sibbern/bemaerkninger:56  (Bemærkninger og Undersøgelser, fornemmelig betreffende Hegels Philosophie; pp. 56–57)

\begin{quote}
Herigjennem komme vi til en dybere \tlg{Tilegnelse} og, saa at sige, Assimilation
af det Sande; det bliver herved inderligere og dybere optaget, idet med det
Samme dets Indhold herunder først oplader og udfolder sig for os i al sin
Rigdom og al sin Gehalt. Men om vi nu end begjære paa saadan Maade at
trænge ind i og i os at indoptage [note: Jeg vilde ønske, at det danske
Sprog her frembød et Udtryk, svarende til det tydske Innewerden.] Alt, hvad
der paa hiin første umiddelbare Maade gjør sig gjeldende ved sin
umiddelbare Nærværelses Magt, og derfor at forvandle den ligefrem anskuende
Erkjenden til en intelligent eller indseende Erkjenden, saa begjære vi dog,
omvendt, ogsaa, at Alt, hvad der under denne Begriben og begrebsmæssige
Indtrængen oplader sig for os, maa komme til at træde saaledes ud i Eenhed
og Heelhed for os, at det kan gaae tilbage i hiin Anskuelsens
Umiddelbarhed, og blive os Noget, som vi i umiddelbar Nærværelse have heelt
og udeelt for os. Herhen hører, at philosophisk Erkjendelse og Gehalt søger
at træde frem i umiddelbar Anskuelighed i Poesien, og der gives derfor lige
saa vel en Poesie, der ligger paa den anden Side af en Philosophie, og hvori
en Tidsalders philosophiske Erkjendelser gjøre sig gjeldende, som der gives
den Poesie, der gaaer forud for Philosophien. Philosophiens intermediære
Natur viser sig her. Hvad vi begrebsmæssigen have indseet, begjære vi
overhovedet med Rette, saavidt mueligt, at see umiddelbart for vore Øine.
Man kalder det, at man fordrer det sandseliggjort; men man seer her
Betydningen af denne Fordring og dens Gyldighed. Exempler ere Tilnærmelser
hertil. De fordres med fuld Ret, deels ifølge den begrebsmæssige
Erkjendelses eget Krav og Medfør, fordi Begrebet dog først derved
tilstrækkeligen godtgjør, at det formaaer at præstere hvad det skal
præstere, eller at det formaaer at realisere sig og gjøre sig gjeldende i det
Concrete, deels fordi vi forlange at see det Sande gjort gjeldende saa
umiddelbart, som mueligt, ved dets egen umiddelbare Magt og „Præsents.“ —
Experimenter til Beviislig- og Anskueliggjørelse høre herhen. —
Forvisningens Styrkelse ved megen Omgang med Begreberne i visse Sphærer, hvor
de kunne komme til praktisk eller in usu at gjøre sig gjeldende, og
at staae deres Prøve, idet de lade vidne om sig, og selv bære Vidnesbyrd om
sig selv, hører ogsaa herhen.
\end{quote}
\dcite{Sibbern, Bemærkninger og Undersøgelser, fornemmelig betreffende Hegels Philosophie, pp. 102–103  [sibbern/bemaerkninger:107]}
% sibbern/bemaerkninger:107  (Bemærkninger og Undersøgelser, fornemmelig betreffende Hegels Philosophie; pp. 102–103)

\begin{quote}
Det er, inden vi gaae videre, ikke at forglemme, som Noget, der
udtrykkeligen bør bemærkes, at vistnok Alt, hvad vi kunne henregne til den
reelle Anskuelses Gebeet, kun kommer til at blive anskueligt formedelst en
medierende og \tlg{tilegnende} Tænkning, og at denne da paa ethvert Punct, som en
uvitterlig Reflexion, er tilstede i Anskuelsen, da der i denne overalt er en
Sammenfatten af et Mangfoldigt til en sammenhængende Erkjenden, hvori da
overalt er Noget af Det, i hvis Iagttagelse og Forfølgelse den begribende
Erkjenden viser sig. Gjorde dette sig gjeldende, som saadant (og det gjør det,
saasnart Nogen i en uvant Sphære skal orientere sig frem, for at komme til at
see eller høre retteligen), saa vilde det vise sig, at der ogsaa paa ethvert
Punct var noget Anerkjendende, noget Tilfredsstillende, noget Indlysende,
hvori Tænkningens Maa bestandigen indtraadte. Saaledes ere da fra først af de
to Arter af Erkjendelse forbundne i Eet, men derpaa, idet Erkjendelsen gjør
Fremgang og vil arbeide sig hen til Andet og Mere, hvorved da Reflexionen
træder frem som saadan, træde de to Arter af Erkjenden ud fra hinanden, for
atter i høiere Erkjendelsesregioner, eller paa høiere Erkjendelsens
Standpuncter, at gaae sammen igjen til Eenhed. — Fremdeles bemærker jeg her
Følgende: den intelligente eller begribende Erkjenden gaaer især og
væsentligt ud paa en Indtrængen i det Givnes Natur og Væsen, og gaaer derved
ud paa en ganske anden Erkjendelsesart, og fører ind i en ganske anden
Erkjendelsens Region, end den første umiddelbare anskuende, eller sandselige
Erkjenden. Men den kan ogsaa vise sig som den, der blot træder i dennes
Tjeneste. Det ad den umiddelbare Affectionsnødvendigheds Vei Givne kan
underkastes en Randsagelse, der blot gaaer ud paa, enten, at vi gjøre os Rede
for, hvad det er, som fremstiller sig for os — da have vi med det
Sandseliges Explication blot i sandselig Henseende at gjøre; eller ogsaa at vi
spørge om Gyldigheden af saadan sandselig Anskuen, da jo Illusioner kunne
indtræde — da ville vi ved den prøvende og til Indsigt førende Erkjenden see
os bestemte til, og tilmed berettigede til, at følge den sandselige Anskuen,
eller vel og til at forkaste denne. Forsaavidt kan den sandselige Skuen
bestyrkes i sin Gyldighed ved den begribende Erkjenden, men den kan og
underkjendes ved denne. Omvendt synes det ogsaa, at denne atter kan ved hiin
snart bestyrkes, snart bringes til at anerkjende sig i sin
Ugyldighed [note: Jeg maa herved gjøre følgende Anmærkning: Ordet
Erkjendelse tages deels i en objectiv Betydning, og da ere Sandhed
og Erkjendelse correlata, idet Erkjendelsen da er en Bevidsthedens
Indtrængen i og Viden af det Sande — deels i en subjectiv
Betydning, og da er Erkjendelse det Samme som forvisningsfuld Antagen. En
saadan er imidlertid ikke altid i sig selv (objectiv) gyldig, og
forsaavidt kan der tales om en Erkjendelses Ugyldighed.] . Herhen synes Det at
høre, hvad vi nys angaaende Exempler, Experimenter o. s. v. have sagt. Men
egentligen er det her altid den begribende Erkjenden, som derigjennem enten
selv bestyrker sig selv, eller selv anerkjender sine Antagelsers Ugyldighed.
Hvad enten Noget betragtes in abstracto, eller in concreto,
er det altid den intelligente Erkjenden, paa hvis Skuen det beroer; thi det er
det Intellectuelle, der i begge Tilfælde skal betragtes, den ene Gang i og for
sig, den anden Gang i sin Udtræden i det bestemte Enkelte eller Concrete. At
prøve hører ei den umiddelbare, sandselige Skuen til; thi Prøvelse er strax
Tænkningens Sag; men Bestyrkelse, saavelsom Underkjendelse forudsætter
Prøvelse.
\end{quote}
\dcite{Sibbern, Bemærkninger og Undersøgelser, fornemmelig betreffende Hegels Philosophie, pp. 103–106  [sibbern/bemaerkninger:109]}
% sibbern/bemaerkninger:109  (Bemærkninger og Undersøgelser, fornemmelig betreffende Hegels Philosophie; pp. 103–106)

\begin{quote}
Det er vel at mærke, at her, i den Erkjendelsens Region, hvori vi nu ere
traadte ind, hører hiint Sympathiske, hiin Grebethed og hiint Hengivelsesfulde,
ei blot væsentligen med til den fulde \tlg{Tilegnelse}, men og til den virkelige
Erkjendelse og Anerkjendelse selv, som væsentligt Grundmoment i denne.
\end{quote}
\dcite{Sibbern, Bemærkninger og Undersøgelser, fornemmelig betreffende Hegels Philosophie, p. 109+  [sibbern/bemaerkninger:116]}
% sibbern/bemaerkninger:116  (Bemærkninger og Undersøgelser, fornemmelig betreffende Hegels Philosophie; p. 109+)

\begin{quote}
At til \tlg{Tilegnelsen} det sympatethiskt Bevægende hører med, er erkjendt i den
almindelige Sætning, at det overalt i al Art af Virken, altsaa og Erkjenden,
er Lysten, der driver Værket, som man siger. Den hører til at levendegjøre
Erkjendelsen. I al Erkjendelse besidde vi dog egentligen kun det Erkjendte
paa den fulde, fyldige Maade, forsaavidt vi deri besidde det Sande, som
personliggjort i os, som Noget, hvori vi villigen og gjærne røre os, hvori vi
have og nyde vor egen personlige Tilvær. Jo mere dette Pathos taber sig, desto
mere staaer Erkjendelsen som et Skyggebillede i os, en Reminiscents; og det er
ogsaa af høieste Grad af Vigtighed, i Opdragelsen at agte paa, at ikke
Kundskaberne bibringes paa saadan Maade, at de blive til saadanne
Skyggebilleder i Individerne, uden at der er sandt Liv, sand forvisningsfuld og
tillige drivtfuld Existents i dem. Den matte Skygge af en Reminiscents,
hvortil megen af den christelige Overbeviisning, blandt Andet og Troen paa
Udødelighed, er bleven hos Utallige i vor Tid, ret som om her hiint af Goethe
ofte nævnte Dæmoniske havde været mægtigere, end al Underviisning og Lærdom,
maa tjene til stor Advarsel for os, at vi ei over Omsorgen for Badet og
Vandet til Badet lade Barnet forkomme.
\end{quote}
\dcite{Sibbern, Bemærkninger og Undersøgelser, fornemmelig betreffende Hegels Philosophie, pp. 109–111  [sibbern/bemaerkninger:117]}
% sibbern/bemaerkninger:117  (Bemærkninger og Undersøgelser, fornemmelig betreffende Hegels Philosophie; pp. 109–111)

\begin{quote}
Men ei blot i \tlg{Tilegnelsen} i enhver Sphære, men og i selve den egentlige
Erkjenden, som saadan, i den Sphære, vi her have for os, er hiint Sympathiske
et væsentligt Moment. Den rationale eller intelligente Erkjenden forholder sig
til denne sympathiske, omtrent som til den ligefrem anskuende. Ligesom hvad
der for den Tænkning, der gaaer ud fra Anskuelsen, udfolder sig, maa gaae
tilbage igjen i den Anskuelse, hvorfra det gik ud, og her maa gjenskues; og
ligesom det her er noget Andet: begrebsmæssigt at erkjende og indsee, og noget
Andet: ligefrem anskuende at see umiddelbart med egne Øine; og ligesom nu her
begge Dele, det Anskuende og det Begribende, maae gaae sammen til Eet,
saaledes er det og i den sympathiske Erkjendens Gebeet. Ved denne gives os
først Sagen selv i den Region, hvor denne Art af Erkjenden hører hen. Den
intelligente Indtrængen i Sagen kan ei sætte sig i Sagens Sted, men maa gaae
tilbage til Eenhed med den i denne Henseende mere umiddelbare sympathiske
Erkjenden. Ved det Skjønne see vi det saa klart, at Velbehaget er et
væsentligt Moment i Erkjendelsen; thi kun derigjennem gaaer den hele
Erkjendelses-Region op for os, og deri maa atter Alt, hvad der her gjennem
intelligent Udvikling træder ud, gaae tilbage igjen, for at blive Erkjendelse
af Skjønhed, som Skjønhed. Men paa lignende Maade forholder det sig med den
paa en Basis af Tillid og Pietet hvilende Erkjenden, at ogsaa her disse
Følelser, baade ligge til Grund, og maae samle Alt, hvad der her udvikler sig,
tilbage i deres Fylde igjen. Det følende Hjerte, den „det Erkjendte og meget
Mere, end dette, som i en Livskjærne indsluttende Tro“ [note: See ovenfor
Pag. 29.] er her, hvad i den anskuende Erkjendens Region den umiddelbare
Seen, det Alt tilsammenfattende Øie er, eller hvad i Musikens Rige Gehøret og
den musikalske Sands er, der baade fra først af giver, og heelt igjennem bærer
og holder og væsentligen constituerer den hele Erkjendelseskreds.
\end{quote}
\dcite{Sibbern, Bemærkninger og Undersøgelser, fornemmelig betreffende Hegels Philosophie, pp. 111–112  [sibbern/bemaerkninger:118]}
% sibbern/bemaerkninger:118  (Bemærkninger og Undersøgelser, fornemmelig betreffende Hegels Philosophie; pp. 111–112)

\subsection{\textit{Om Philosophiens Begreb, Natur og Væsen} (1843)}

\begin{quote}
Anm. 5. Men denne Tredeling maae vi endnu i en anden Henseende her
tage i Betragtning. Vi ere komne til at see, hvorledes Erkjendelsen fra sit
første Hovedtrin, Umiddelbarhedens, eller den første umiddelbare Erkjendens
Trin, gaaer over til at antage en ny Charakteer, ved at træde frem som den ved
Randsagelsesaanden fremkaldte middelbare, men tillige intelligente og derved
høiere Erkjenden. Men nu maa da det Spørgsmaal opstaae, hvorledes denne forholder
sig til hiin sympathetiske Erkjenden, der ogsaa viser sig som en saaledes
medieret, at ogsaa ved samme Subjectiviteten træder frem som saadan ligeoverfor
det Objective. Vi bemærke da, at denne her bliver at henregne med til hiin første
umiddelbare Erkjenden, da den — om den end i Sammenligning med den umiddelbart
skuende Erkjenden viser sig at være en mere middelbar, der forsaavidt bliver at
sammenstille med den intelligente — dog, ligesaavel som den umiddelbart
skuende, danner sig forud for al Randsagelse og Prøvelse, og afgiver Noget som
ved den intelligente Erkjenden først skal gjøres til Gjenstand for Spørgsmaal om
dets Gyldighed og dets Indhold. — Ligeoverfor den intelligente Erkjenden, der
ogsaa, som vi længere hen nærmere ville see, kan kaldes den rationale, finde vi
en anden, til hvilken vi ikke komme gjennem en rational Stræben, men som danner
sig forud for denne, og har sit eget, af det Rationale, hvorved vi mene det i sin
Rationalitet fremtrædende Rationale, uafhængige Grundlag. Denne anden Erkjenden
kunne vi da kalde den extrarationale eller præterrationale (fordi vi ikke have
den ifølge Randsagelse og Regnskabsaflæggelse eller ifølge Grunde, der have
fremstillet sig for os som saadanne, men faae den paa anden Maade). Men til
saadan extrarational Erkjenden hører nu baade den, der allene skylder den
udvortes Sandsning sin Tilvær, samt hvad der bliver at stille sammen hermed
(hvorhen Rums- og Tidsanskuelsen henhører), saa og den hele sympathetiske
Erkjenden, der ogsaa danner og haandhæver sig udenfor hiin Rationalitet. Al
saadan extrarational Erkjenden er ogsaa, i Henseende til sin Opstaaen, ifølge sin
Natur mere beroende paa umiddelbar Paavirkning, Affection, end den intelligente.
Iøvrigt finde vi ogsaa i dennes Sphære Meget, der danner sig paa Umiddelbarhedens
Maade, og maa det overhovedet bemærkes, at Menneskeheden, især i de dannede
Nationer, nu paa Umiddelbarhedens Maade, eller af sig selv, d. e. uden
udtrykkelig Overveielse og Betænkning, \tlg{tilegner} sig Meget, som man fra først af
kun langsomt har arbeidet sig hen til. Dette er, hvad især de Gamles Musik
angaaer, bleven oplyst af Hr. Pastor Bindesbøl i hans Anmældelse af Dr.
Bojesens Doctordissertation de harmonica scientia Græcorum (1833) i
Maanedskr. f. Liter. 1837, 8de Hefte; cfr. Indl. til Dr. Bojesens
Skoleprogram 1843 de tonis s. harmoniis Græcorum, Pag. 6.
\end{quote}
\dcite{Sibbern, Om Philosophiens Begreb, Natur og Væsen, p. 18  [sibbern/philosophiens-begreb:46]}
% sibbern/philosophiens-begreb:46  (Om Philosophiens Begreb, Natur og Væsen; p. 18)

\begin{quote}
Den hele Rationalitet i vor Erkjenden gaaer nemlig ud paa Mere, end blot
Erkjendelsernes Grundethed og Gyldighed. Den gaaer ud paa selve Gjenstandenes
Grundforhold, og altsaa tillige ud paa disses egne Grunde, idet vi stræbe efter,
at Det, hvori hvert Punct i selve Tilværelsen er grundet, maa vise sig for os. Thi
vor hele Erkjendelsesbestræbelse gaaer jo ud paa, at vi i og gjennem Erkjendelse
kunne fuldeligen \tlg{tilegne} os og indleve os i Objectiviteten; men da maae vi gaae
ind i Sammenhængen og Forholdene i alt det Givne, for at Alt heri kan vise sig for
os, hvert i sin Betydning i Totaliteten. Men saaledes komme vi da til, ei blot at
kræve os selv til Regnskab med Hensyn til vor Opfattelse af det Objective, men
komme ogsaa til at kræve selve det Objective paa ethvert Punct, eller kræve selve
Tingene til Regnskab, for at de kunne ligesom gjøre os Rede for, hvori de selv ere
grundede, og hvorledes det staaer sig med deres Grundforhold. — Men idet vi nu
ogsaa herved tragte efter at gaae til sidste Grunde og Grundforhold, komme vi,
idet hele Tilværelsen gjør sig gjeldende for os som den, der danner en eneste
Heelhed, til at tragte efter en hertil svarende altomfattende Grunderkjendelse,
ikke som et Aggregat af flere, men med indre Eenhed og Heelhed i sig.
\end{quote}
\dcite{Sibbern, Om Philosophiens Begreb, Natur og Væsen, pp. 27–28  [sibbern/philosophiens-begreb:67]}
% sibbern/philosophiens-begreb:67  (Om Philosophiens Begreb, Natur og Væsen; pp. 27–28)

\begin{quote}
Anm. 1. Af ganske anden Art end den \tlg{Tilegnelse} og Indoptagelse, som vi
her omtale, er a) den, som skeer ved den Art af ideel Skuen, der lader os see det
Algyldiges og Alconstitutives indre Momenter træde ud og røre sig som det inderste
til Grund Liggende i den reelle Virkeligheds aldeles individuelle Skikkelser og
Livsbevægelser; hvilket deels skeer i en ideel Empiries deels i Poesiens og den
øvrige ideelle Konsts Gebeet, hvori da det til Grund Liggende viser sig i sin
magtfulde Nærværelse i det Concrete, paa ganske individuel Maade; fremdeles ogsaa b)
den, som paa sympathetisk Maade skeer ved vor egen interessefulde og villiefulde,
kjærlige og religiøse Tilknytning til og Leven i og for Tilværelsens Totalitet. Vi
have ogsaa i alt Dette Noget, hvorved eller hvorigjennem det Sande lige saa fuldt,
som gjennem Philosophering, kun paa en ganske anden Maade, kan sande sig for os; saa
at vel den Bemærkning, at det kan, ja maa det, bliver en philosophisk, men ei denne
Sandelse selv. Paa en saadan Sandelse er det, at Christus viser hen ved de Ord
(Joh. 7, 17): „Min Lærdom er ikke min, men hans, som mig udsendte. Dersom Nogen vil
gjøre hans Villie, han skal kjende, om Lærdommen er af Gud, eller jeg taler af mig
selv."
\end{quote}
\dcite{Sibbern, Om Philosophiens Begreb, Natur og Væsen, p. 48  [sibbern/philosophiens-begreb:112]}
% sibbern/philosophiens-begreb:112  (Om Philosophiens Begreb, Natur og Væsen; p. 48)

\part*{II.\quad Søren Kierkegaard}
\addcontentsline{toc}{part}{II.\quad Søren Kierkegaard}

\section{Published works}

\subsection{BI \normalfont(SKS 1)}

\begin{quote}
Socrates indledede sit Foredrag ved en Ironi, men dette var blot, om jeg saa maa sige, en ironisk Figur, og han vilde i Sandhed ikke fortjene Navn af Ironiker, hvis det blot var den Færdighed, hvormed han talte ironisk, ligesom Andre tale Kragemaal, hvorved han udmærkede sig. De foregaaende Talere havde sagt saare Meget om Kjærlighed, hvoraf vel Meget ikke vedkom denne Gjenstand, men den Forudsætning blev dog tilbage, at der var en Mængde at sige om Kjærlighed. Nu udvikler Socrates det for dem. Og see, Kjærlighed er Attraa, er Trang og saa videre Men Attraa, Trang og saa videre ere Intet. Vi see nu Methoden. Kjærligheden frigjøres bestandig mere og mere fra den tilfældige Concretion, hvori den i de foregaaende Taler viste sig, og føres tilbage til sin abstracteste Bestemmelse, hvor den viser sig ikke som Kjærlighed i Dette eller Hiint eller til Dette eller Hiint, men som Kjærlighed til et Noget, som den ikke har det vil sige som Attraa, Længsel. Dette er nu i en vis Forstand meget sandt; men Kjærlighed er derhos ogsaa den uendelige Kjærlighed. Naar vi sige, at Gud er Kjærlighed, da sige vi jo derved, at han er den uendeligt sig meddelende; naar vi tale om at forblive i Kjærlighed, da tale vi jo om en Deelagtighed i en Fylde. Dette er det Substantielle i Kjærligheden. Det Attraaende, det Længselsfulde er det Negative i Kjærligheden, det vil sige den immanente Negativitet. Attraa, Trang, Længsel og saa videre er Kjærlighedens uendelige Subjectivering, for at bruge et af Hegels Udtryk, der her netop erindrer om, hvad det skal erindre om. Denne Bestemmelse er nu tillige den meest abstracte, eller rettere, den er selve det Abstracte, ikke i Betydning af det Ontologiske, men i Betydning af det sit Indhold Manglende. Man kan naturligviis enten opfatte det Abstracte som det Alt Constituerende, forfølge det i sine egne tause Bevægelser, og lade det bestemme sig hen til det Concrete, udfolde sig deri. Eller man kan gaae ud fra det Concrete og med det Abstracte in mente finde det i og ud af det Concrete. Men ingen af Delene er Tilfældet med Socrates. Det er ikke til Kategorier, han kalder Forholdet tilbage. Hans Abstracte er en aldeles indholdsløs Betegnelse. Han gaaer ud fra det Concrete og kommer til det meest Abstracte, og der hvor nu Undersøgelsen skulde begynde, der hører han op. Det Resultat, han kommer til, er egentlig den rene Værens ubestemmelige Bestemmelse: Kjærlighed er; thi det Tillæg, den er Længsel, Attraa, er ingen Bestemmelse, da det blot er et Forhold til et Noget, som ikke er givet. Paa samme Maade kunde man ogsaa føre Erkjendelse tilbage paa et aldeles negativt Begreb, ved at bestemme den som \tlg{Tilegnelse}, Erhvervelse, thi dette er jo aabenbart Erkjendelsens ene Forhold til det Erkjendte, men paa den anden Side er den og Besiddelse. Men som nu det Abstracte i Betydning af det Ontologiske har sin Gyldighed i det Speculative, saaledes har det Abstracte som det Negative sin Sandhed i det Ironiske.
\end{quote}
\dcite{BI:107  (SKS 1, 107)}
% BI:107  (SKS 1, 107)

\begin{quote}
Forinden jeg imidlertid gaaer over til denne Documentation, maa jeg endnu dvæle ved lidt nærmere at betragte den Længsel efter Døden, der her i Phædon tillægges Philosophen. Den Anskuelse nemlig, at Livet egentlig bestaaer i at afdøe, kan opfattes deels moralsk, deels intellectuelt. Moralsk er den bleven opfattet i Christendommen, der da heller ikke er bleven staaende ved det blot Negative; thi i samme Grad som Mennesket afdøer fra det Ene, voxer det Andet en guddommelig Væxt, og naar dette Andet har ligesom indsuget og \tlg{tilegnet} sig og derved forædlet al den Spirekraft, der var i det Syndens Legeme, der skulde dødes, da er dette ogsaa efterhaanden skrumpet sammen, udtørret, og idet det brister og hensmuldrer, da hæver det fuldvoxne Guds-Menneske sig ud deraf, der er skabt efter Gud i Sandheds Retfærdighed og Hellighed. Forsaavidt Christendommen ogsaa antager en fuldkomnere Erkjendelse forbunden med denne Gjenfødelse, saa er dette dog kun secundairt og fornemmelig kun forsaavidt, som Erkjendelsen tilforn ogsaa har været befængt med Syndens Smitte. Intellectuelt er denne Afdøen opfattet i Græciteten, det vil sige reent intellectuelt, og man kjender ogsaa strax herpaa Hedenskabets sorgløse Pelagianisme. Man tillade mig at udhæve for en Forsigtigheds Skyld, hvad forøvrigt formodentlig de fleste Læsere indsee. Paa den ene Side er i Christendommen det, der skal afdøes fra, opfattet i sin Positivitet, som Synd, som et Rige, der kun altfor overbevisende forkynder sin Gyldighed for enhver, der sukker under dets Love; paa den anden Side er det, der skal fødes og opstaae, opfattet ligesaa positivt. Ved den intellectuelle Afdøen er det, der skal afdøes fra, noget Indifferent, det der skal voxe derunder noget Abstrakt. Forholdet mellem disse to store Anskuelser er omtrent følgende. Den Ene siger, man skal afholde sig fra usunde Spiser, beherske Lysten, saa skal Sundheden udfolde sig, den Anden siger, man skal lade være at spise og drikke, saa kan man nok have Haab om efterhaanden at blive til Intet. Man seer saaledes, at Grækeren er større Rigorist end den Christne, men derfor er hans Anskuelse ogsaa usand. Efter den christelige Anskuelse er da ogsaa Dødsøieblikket den sidste Strid mellem Dag og Nat, Døden er, som den i Kirken skjønt er bleven kaldet, Fødselen. Med andre Ord, den Christne dvæler ikke ved Striden, ved Tvivlen, ved Smerten, ved det Negative, men fryder sig ved Seiren, ved Visheden, ved Saligheden, ved det Positive. Platonismen vil, at man skal afdøe fra den sandselige Erkjenden for ved Døden at opløses i det Udødelighedens Rige, hvor det i og for sig Lige, det i og for sig Skjønne og saa videre leve i Dødsstilhed. Dette udtrykkes endnu stærkere ved en Yttring af Socrates, at Philosophens Attraa er efter at døe og at være død. Men at attraae Døden saaledes i og for sig, kan ikke være grundet i Begeistring, forsaavidt man vil agte dette Ord og ikke for Exempel bruge det om det Vanvid, hvormed man vel stundom seer Mennesket at ønske at blive til Intet, det maa være grundet i en Art Livslede. Saalænge der endnu ikke kan være Tale om, at man ret indseer, hvad der ligger i dette Ønske, saalænge kan Begeistringen endnu være der, hvorimod, naar Ønsket har sin Grund i en vis Dorskhed, eller den Ønskende er sig bevidst, hvad han ønsker, da er Livslede det Fremherskende. I mine Øine indeholder den bekjendte Gravskrift af Wessel: »tilsidst gad han ei heller leve,« Ironiens Anskuelse af Døden. Men den der døer, fordi han ikke gider leve, han ønsker visselig heller ikke et nyt Liv, thi det vilde jo være en Modsigelse. Denne Mathed, der i den Forstand ønsker Døden, er nu aabenbart en fornem Sygdom, der kun har hjemme i de høieste Cirkler, og er i sin fuldkomne Ublandethed ligesaa stor som den Begeistring, der i Døden seer Livets Forklarelse. Det er de to Poler, mellem hvilke det almindelig menneskelige Liv bevæger sig døsigt og uklart. Ironien er nemlig en Sundhed, forsaavidt den frelser Sjælen ud af det Relatives Besnærelser, den er en Sygdom, forsaavidt den ikke kan bære det Absolute uden i Form af Intet, men denne Sygdom er en Climatfeber, som kun faa Individer paadrage sig, færre overstaae.
\end{quote}
\dcite{BI:135  (SKS 1, 135)}
% BI:135  (SKS 1, 135)

\begin{quote}
Naar man derfor stundom hører En og Anden med stor Fornemhed at tale om Ironien i den uendelige Stræben, hvori den løber løbsk, saa kan man vel gjerne give ham Ret, men saafremt han ikke kjender den Uendelighed, der rører sig i Ironien, saa staaer han ikke over men under Ironien. Saaledes gaaer det overalt, hvor man overseer Livets Dialectik. Der hører Mod til, ikke at give efter for det Fortvivlelsens kløgtige eller medlidende Raad, der tillader En, at udslette sig selv af de Levendes Tal; men deraf følger ingenlunde, at enhver Spekhøker, der fedes og næres i Selvtilfredshed, derfor har mere Mod end den, der gav efter for Fortvivlelsen. Der hører Mod til, naar Sorgen vil bedaare En, naar den vil lære at forfalske al Glæde til Veemod, al Længsel til Savn, ethvert Haab til Erindring, der hører Mod til da at ville være glad; men deraf følger ingenlunde, at ethvert voxent gammelt Barn med sit vamle Smiil, sine glædedrukne Øine har mere Mod end den, der bøiede sig i Sorgen, og som glemte at smile. Saaledes ogsaa med Ironien. Maa man derfor end advare mod Ironien som mod en Forfører, saa maa man ogsaa anprise den som en Veileder. Netop i vor Tid maa man anprise den. Videnskaben er saaledes i vor Tid kommen i Besiddelse af et saa uhyre Resultat, at det neppe kan hænge rigtigt sammen dermed; Indsigt ikke blot i Menneskeslægtens, men selv i Guddommens Hemmeligheder falbydes for saa godt Kjøb, at det seer ganske betænkeligt ud. Man har i vor Tid af Glæde over Resultatet glemt, at et Resultat dog ingen Værdi har, naar det ikke er erhvervet. Men vee den, der ikke kan taale, at Ironien vil gjøre Regnskabet op. Ironien er som det Negative Veien; ikke Sandheden, men Veien. Enhver, som har et Resultat som saadant, han eier det ikke; thi han har ikke Veien. Naar Ironien nu træder til, da bringer den Veien, men ikke den Vei, hvorved den, der bilder sig ind at have Resultatet, kommer til at besidde det; men den Vei, ad hvilken Resultatet forlader ham. Dertil kommer, at det vel nærmest maa ansees for vor Tids Opgave, at sætte Videnskabens Resultater over i det personlige Liv, personligt at \tlg{tilegne} sig disse. Naar saaledes Videnskaben lærer, at Virkeligheden har absolut Gyldighed, saa gjelder det i Sandhed om, at den faaer Gyldighed, og man kan dog ikke nægte, at det vilde være saare latterligt, om En, der i sin Ungdom lærte og maaskee endog lærte Andre, at Virkeligheden havde absolut Betydning, blev gammel og døde, uden at Virkeligheden havde havt anden Gyldighed end den, at han baade i Tide og Utide havde forkyndt den Viisdom, at Virkeligheden havde Gyldighed. Naar Videnskaben medierer alle Modsætninger, saa gjelder det om at denne fyldige Virkelighed i Sandhed kommer tilsyne. Der er i en anden Retning i vor Tid en uhyre Begeistring og besynderligt nok, det, der begeistrer den, synes at være uhyre Lidet. Hvor velgjørende kan her ikke Ironien være. Der er en Utaalmodighed, der vil høste, førend den saaer, lad kun Ironien tugte den. Der er i ethvert personligt Liv saa Meget, der maa afvises, saa mange vilde Skud, der maae bortskjæres: her er Ironien atter en fortrinlig Operateur; thi, som sagt, naar Ironien er bleven behersket, da er dens Function af yderste Vigtighed, for at det personlige Liv kan faae Sundhed og Sandhed.
\end{quote}
\dcite{BI:356  (SKS 1, 356)}
% BI:356  (SKS 1, 356)

\subsection{EE1 \normalfont(SKS 2)}

\begin{quote}
Det er maaskee dog stundom faldet Dig ind, kjære Læser, at tvivle en Smule om Rigtigheden af den bekjendte philosophiske Sætning, at det Udvortes er det Indvortes, det Indvortes det Udvortes. Du har maaskee selv gjemt en Hemmelighed, om hvilken Du følte, at den, i sin Glæde eller i sin Smerte, var Dig for kjær til, at Du kunde indvie Andre i den. Dit Liv har maaskee bragt Dig i Berøring med Mennesker, om hvilke Du ahnede, at noget Saadant var Tilfælde, uden at dog Din Magt eller Din Besnærelse var istand til at bringe det Forborgne til Aabenbarelse. Maaskee passer intet af Tilfældene paa Dig og Dit Liv, og dog er Du ikke ukjendt med hiin Tvivl; den har nu og da som en flygtig Skikkelse svævet Din Tanke forbi. En saadan Tvivl kommer og gaaer, og Ingen veed, hvorfra den kommer eller hvor den farer hen. Jeg for min Part har altid været noget kjættersk sindet paa dette Punkt af Philosophien, og har derfor tidlig vant mig til, saa godt som muligt, selv at anstille Iagttagelser og Efterforskninger; jeg har søgt Veiledning hos de Forfattere, hvis Anskuelse jeg i denne Henseende deelte, kort, jeg har gjort hvad der stod i min Magt for at bøde paa det Savn, de philosophiske Skrifter efterlode. Efterhaanden blev da Hørelsen mig den kjæreste Sands; thi ligesom Stemmen er Aabenbarelsen af den for det Ydre incommensurable Inderlighed, saaledes er Øret det Redskab, ved hvilket denne Inderlighed opfattes, Hørelsen den Sands, ved hvilken den \tlg{tilegnes}. Hver Gang jeg da fandt en Modsigelse mellem hvad jeg saae, og hvad jeg hørte, da fandt jeg min Tvivl bestyrket, og min Iagttagelseslyst blev forøget. En Skriftefader er ved et Gitter adskilt fra den Skriftende, han seer ikke, han hører blot. Efterhaanden som han hører, danner han et Udvortes, der svarer hertil; han kommer altsaa ikke i Modsigelse. Anderledes derimod, naar man paa eengang seer og hører, og dog seer et Gitter mellem sig og den Talende. Mine Bestræbelser for i denne Retning at anstille Iagttagelser have, hvad Resultatet angaaer, været saare forskjellige. Snart har jeg havt Lykken med mig, snart ikke, og Lykke hører der altid til for paa disse Veie at vinde noget Udbytte. Imidlertid har jeg dog aldrig tabt Lysten til at fortsætte mine Efterforskninger. Har jeg end en enkelt Gang været nær ved at fortryde min Vedholdenhed, saa har ogsaa en enkelt Gang et uventet Held kronet mine Bestræbelser. Et saadant uventet Held var det, der paa en høist besynderlig Maade satte mig i Besiddelse af de Papirer, jeg herved har den Ære at forelægge det læsende Publicum. I disse Papirer fik jeg Leilighed til at gjøre et Indblik i tvende Menneskers Liv, der bestyrkede min Tvivl om, at det Udvortes ikke er det Indvortes. Dette gjælder i Særdeleshed om den Ene af dem. Hans Udvortes har været i fuldkommen Modsigelse med hans Indvortes. Ogsaa om den Anden gjælder det til en vis Grad, forsaavidt han under et ubetydeligere Udvortes har skjult et betydeligere Indre.
\end{quote}
\dcite{EE1:12  (SKS 2, 12)}
% EE1:12  (SKS 2, 12)

\begin{quote}
Alle classiske Frembringelser staae, som ovenfor blev bemærket, lige høit, fordi enhver staaer uendelig høit. Vil man altsaa desuagtet gjøre et Forsøg paa at indføre en vis Orden i denne Procession, saa følger det af sig selv, at denne ikke kan være begrundet paa noget Væsentligt; thi deraf vilde jo følge, at der var en væsentlig Forskjel, hvoraf atter vilde følge, at Ordet »classisk« med Urette prædiceredes om dem Alle. Vilde man saaledes begrunde en Classification paa Stoffets forskjellige Beskaffenhed, saa vilde man strax herved indvikle sig i en Misforstaaelse, der i sin vidtløftigere Udbredelse vilde ende med at ophæve det hele Begreb af det Classiske. Stoffet er nemlig et væsentligt Moment, forsaavidt det er den ene Faktor, men det er ikke det Absolute, da det kun er det ene Moment. Man kunde saaledes gjøre opmærksom paa, at der ved visse Arter af classiske Frembringelser paa en Maade intet Stof er til, medens derimod ved andre Stoffet spiller en saa betydelig Rolle. Det Første er Tilfældet med de Værker, vi beundre som classiske i Architektur, Sculptur, Musik, Maleri, især med de tre første, saaledes, at selv hvad Maleriet angaaer, forsaavidt der er Tale om Stoffet, dette dog nærmest kun har Betydning af at være Anledning. Det Andet gjælder om Poesien, dette Ord taget i sin videste Betydning, hvorved det betegner al kunstnerisk Frembringelse, der er baseret paa Sproget og den historiske Bevidsthed. Denne Bemærkning er i sig selv ganske rigtig; men vil man paa den begrunde en Classification, ved at ansee Mangelen paa Stof eller dettes Tilstedeværelse for en Fordeel eller en Besvær for det producerende Subject, saa løber man vild. Stricte taget kommer man nemlig til at urgere det Modsatte af det man egentlig vilde, saaledes som det altid gaaer, naar man abstrakt bevæger sig i dialektiske Bestemmelser, hvor det da ikke blot gjælder, at man siger Eet, mener et Andet, men man siger det Andet; det man troer man siger, siger man ikke, men siger det Modsatte. Saaledes hvor man gjør Stoffet gjældende som Inddelings-Princip. Idet man taler derom, taler man om noget ganske Andet, nemlig om den formende Virksomhed. Vil man derimod gaae ud fra den formende Virksomhed og alene udhæve den, saa har man samme Skjæbne. Idet man her vil gjøre Forskjellen gjældende, og saaledes udhæve, at i nogle Retninger den formende Virksomhed i den Grad er skabende, at den skaber Stoffet med, hvorimod den i andre modtager Stoffet, saa taler man her igjen, skjøndt man troer at tale om den formende Virksomhed, egentlig om Stoffet, og begrunder egentlig Classificationen paa Inddelingen af Stoffet. Om den formende Virksomhed som Udgangspunkt for en saadan Classification gjælder aldeles det Samme som om Stoffet. Til at begrunde en Rangfølge kan altsaa den enkelte Side aldrig bruges; thi den er bestandig for væsentlig til at være tilfældig nok, for tilfældig til at begrunde en væsentlig Anordnen. Men denne absolute gjensidige Gjennemtrængen, der gjør, at naar man vil tale distinct, man ligesaa godt kan sige, at Stoffet gjennemtrænger Formen, som at Formen gjennemtrænger Stoffet, denne gjensidige Gjennemtrængen, dette Lige for Lige i det Classiskes udødelige Venskab, kan tjene til at belyse det Classiske fra en ny Side og begrændse det saaledes, at det ikke bliver for ampelt. De Æsthetikere nemlig, der eensidig urgerede den digteriske Virksomhed, have udvidet dette Begreb saa meget, at hiint Pantheon blev i den Grad beriget, ja overlæsset med classiske Snurrepiberier og Bagateller, at den naturlige Forestilling om en kølig Halle med enkelte bestemte store Skikkelser aldeles forsvandt og hiint Pantheon snarere blev et Pulterkammer. Enhver i kunstnerisk Henseende fuldendt lille Nethed er efter denne Æsthetik et classisk Værk, der var sikkret absolut Udødelighed; ja i dette Hocus-Pocus indrømmede man slige Smaating allermest; skjøndt man ellers hadede Paradoxer, frygtede man dog ikke for det Paradox, at det Mindste egentlig var Kunsten. Det Usande ligger i, at man eensidig udhævede den formelle Virksomhed. En saadan Æsthetik kunde derfor kun holde sig en bestemt Tid, saa længe man nemlig ikke blev opmærksom paa, at Tiden spottede den og dens classiske Værker. Paa Æsthetikens Gebeet var denne Anskuelse en Form af den Radikalisme, der paa en tilsvarende Maade har yttret sig paa saa mange Gebeter, den var en Yttring af det tøilesløse Subject i dets ligesaa tøilesløse Indholdsløshed. Denne Stræben som saa mange andre fandt imidlertid sin Betvinger i Hegel. Det er overhovedet en sørgelig Sandhed med Hensyn til den hegelske Philosophi, at den ingenlunde har faaet den Betydning hverken for en forgangen eller en nærværede Tid, som den vilde have faaet, hvis den forgangne ikke havde havt saa travlt med at skræmme Folk ind i den, men derimod lidt mere præsentisk Ro i at \tlg{tilegne} sig den, den nærværende ikke været saa utrættelig virksom for at jage Folk ud over den. Hegel indsatte igjen Stoffet, Ideen i dens Rettigheder, og fordrev derved disse flygtige classiske Værker, disse lette Væsener, Tusmørke-Sværmere fra Classicitetens Hvælvinger. Det kan ingenlunde være vor Hensigt at frakjende disse Værker det dem tilkommende Værd, men det gjælder om at vaage over, at ikke her som saa mange andre Steder Sproget forvirres, Begreberne enerveres. En vis Evighed kan man gjerne tillægge dem, og dette er det Fortjenstfulde ved dem; men denne Evighed er dog egentlig blot det evige Øieblik, som enhver sand kunstnerisk Frembringelse har, ikke den fyldige Evighed midt i Tidernes vexlende Omskiftelser. Det, hine Frembringelser manglede, var Ideer, og jo mere fuldendte de vare i formel Henseende, desto hurtigere forbrændte de i dem selv, jo mere den techniske Færdighed udvikledes til den høieste Grad af Virtuositet, desto flygtigere blev denne, og havde ikke Mod og Kraft eller Holdning til at modstaae Tidens Pust, idet den fornemmere og fornemmere bestandig gjorde større Fordring paa at være den mest rectificerede Spiritus. Kun hvor Ideen er bragt til Hvile og Gjennemsigtighed i en bestemt Form, kun der kan være Tale om et classisk Værk; men det vil da ogsaa være istand til at modstaae Tiderne. Denne Eenhed, denne gjensidige Inderlighed i hinanden, har ethvert classisk Værk, og man seer saaledes let, at ethvert Attentat paa en Classification af de forskjellige classiske Værker, der til sit Udgangspunkt har en Adskillelse af Stof og Form eller Idee og Form, eo ipso er forfeilet.
\end{quote}
\dcite{EE1:60  (SKS 2, 60)}
% EE1:60  (SKS 2, 60)

\begin{quote}
Labdakos's Slægt er altsaa Gjenstand for de vrede Guders Forbittrelse, Oedip har dræbt Sphinx, befriet Theben, Oedip har myrdet sin Fader, ægtet sin Moder og Antigone er Frugten af dette Ægteskab. Saaledes i den græske Tragedie. Her afviger jeg. Alt forholder sig hos mig ligesaa, og dog er Alt anderledes. At han har dræbt Sphinx og befriet Theben, det er Alle bekjendt, og Oedip lever hædret og beundret, lykkelig i sit Ægteskab med Jokaste. Det Øvrige er skjult for Menneskenes Øine, og ingen Ahnelse har nogensinde kaldet denne rædsomme Drøm til Virkelighed. Kun Antigone veed det. Hvorledes hun har faaet det at vide, ligger udenfor den tragiske Interesse, og Enhver kan i den Henseende overlade sig til sin egen Combination. I en tidlig Alder, før hun endnu var fuldkommen udviklet, have dunkle Hentydninger paa denne rædsomme Hemmelighed momentviis grebet hendes Sjæl, indtil Visheden med eet Slag kaster hende i Angstens Arme. Her har jeg nu strax en Bestemmelse af det moderne Tragiske. Angst er nemlig en Reflexion og er forsaavidt væsentlig forskjellig fra Sorg. Angst er det Organ, hvorved Subjectet \tlg{tilegner} sig Sorgen og assimilerer sig den. Angst er den Bevægelses Kraft, ved hvilken Sorgen borer sig ind i Eens Hjerte. Men Bevægelsen er ikke hurtig som Pilens, den er successiv, den er ikke eengang for alle, men den vorder bestandig. Som et lidenskabeligt erotisk Øiekast attraaer sin Gjenstand, saaledes seer Angsten paa Sorgen for at attraae den. Som et stille uforkrænkeligt Kjærligheds Blik sysler med den elskede Gjenstand, saaledes er Angsten Selvbeskæftigelse med Sorgen. Men Angsten har et Moment mere i sig, som gjør, at den endnu stærkere holder fast ved sin Gjenstand, thi den baade elsker og frygter den. Angst har en dobbelt Function, deels er den den opdagende Bevægelse, der bestandig berører og ved denne Tasten opdager Sorgen, idet den gaaer rundt om Sorgen. Eller Angsten er pludselig, sætter i et eneste Nu hele Sorgen, dog saaledes, at dette Nu øieblikkeligt opløser sig i Succession. Angst i denne Betydning er en ægte tragisk Bestemmelse, og det gamle Ord: quem deus vult perdere, primum dementat, lader sig her ret med Sandhed anvende. At Angst er en Reflexions-Bestemmelse, det viser Sproget selv; thi jeg siger altid: at ængstes for Noget, hvorved jeg adskiller Angsten fra det, hvorfor jeg ængstes, og jeg kan aldrig bruge Angst i objectiv Forstand, medens jeg derimod, naar jeg siger: min Sorg, ligesaa meget kan udtrykke det, jeg sørger over, som min Sørgen derover. Dertil kommer, at Angst altid i sig indeholder en Reflexion paa Tid, thi jeg kan ikke ængstes over det Nærværende, men kun for det Forbigangne eller Tilkommende, men det Forbigangne og det Tilkommende, saaledes holdte imod hinanden, at det Præsentiske forsvinder, er Reflexions-Bestemmelser. Den græske Sorg derimod er, som hele det græske Liv, præsentisk, og derfor Sorgen dybere men Smerten mindre. Angsten hører derfor væsentlig med til det moderne Tragiske. Derfor er Hamlet saa tragisk, fordi han ahner Moderens Forbrydelse. Robertle diable spørger, hvoraf det dog kunde komme, at han gjør saa meget Ondt. Høgne, hvem Moderen havde avlet med en Trold, kommer tilfældigviis til at see sit Billede i Vandet og spørger nu Moderen, hvorfra hans Legeme har faaet en saadan Form.
\end{quote}
\dcite{EE1:154  (SKS 2, 154)}
% EE1:154  (SKS 2, 154)

\begin{quote}
Dog, vi vende tilbage til Marie Beaumarchais. Det Eiendommelige ved hendes Sorg er, som ovenfor bemærket, den Uro, der forhindrer hende i at finde Sorgens Gjenstand. Hendes Smerte kan ikke faae Stilhed, hun mangler den Fred, der er nødvendig for ethvert Liv, naar det skal kunne \tlg{tilegne} sig sin Næring og qvæges af den; ingen Illusion overskygger hende med sin stille Kølighed, medens hun indsuger Smerten. Hun har tabt Barndommens Illusion, da hun vandt Elskovens, hun har tabt Elskovens, da Clavigo bedrog hende; hvis det var muligt for hende at vinde Sorgens Illusion, da var hun hjulpen. Da vilde hendes Sorg naae Mands Modenhed og hun have Vederlag for den Tabte. Men hendes Sorg trives ikke, thi hun har ikke tabt Clavigo, han har bedraget hende, den bliver altid et spædt Barn med sit Skrig, et fader- og moderløst Barn; thi hvis Clavigo var bleven hende berøvet, da vilde det i Erindringen om hans Trofasthed og Elskelighed have havt Faderen og i Marias Sværmeri Moderen; og hun har Intet, hvoraf hun kan opfostre det; thi det Oplevede var vel skjønt, men havde dog ikke Betydning i og for sig, men som Forsmag paa det Tilkommende; og haabe kan hun ikke, at dette Smertens Barn skal forvandles til en Glædens Søn, haabe kan hun ikke, at Clavigo skulde vende tilbage; thi hun vil ikke have Styrke til at bære en Fremtid, hun har tabt den glade Tillid, med hvilken hun var fulgt ham uforfærdet i Afgrunden, og hun har faaet hundrede Betænkeligheder isteden, hun vilde i det Høieste kun være istand til med ham at opleve det Forbigangne endnu engang. Der laae en Fremtid for hende, da Clavigo forlod hende, en Fremtid saa skjøn, saa fortryllende, at den næsten forvirrede hendes Tanker, den udøvede dunkelt sin Magt over hende, hendes Metamorphose var allerede begyndt, da blev Udviklingen afbrudt, hendes Forvandling standset; et nyt Liv havde hun ahnet, havde fornummet dets Kræfter røre sig i hende, da blev det brudt og hun stødt tilbage, og der er intet Vederlag for hende, hverken i denne eller i den tilkommende Verden. Det, der skulde komme, smilede hende saa rigt imøde og speilede sig i hendes Elskovs Illusion, og dog var Alt saa naturligt og saa ligefrem; nu har en afmægtig Reflexion maaskee stundom malet hende en afmægtig Illusion, der ikke virker fristende paa hende selv, men vel i et Øieblik dulmende. Saaledes vil Tiden gaae hen for hende, indtil hun har fortæret Gjenstanden selv for hendes Sorg, der ikke var identisk med hendes Sorg, men Anledning til at hun bestandig søgte en Gjenstand for Sorgen. Hvis et Menneske eiede et Brev, hvorom han vidste eller troede, at det indeholdt Oplysning om hvad han maatte ansee for sit Livs Salighed, men Skrifttegnene vare fine og blege, Haandskriften næsten ulæselig, da vilde han vel med Angst og Uro, med al Lidenskab læse og atter læse, og i eet Øieblik faae een Mening ud, i det næste en anden, i Forhold til som han, naar han troede med Bestemthed at have læst et Ord, vilde forklare Alt efter dette; men han vilde aldrig komme videre end til den samme Uvished, med hvilken han begyndte. Han vilde stirre, ængsteligere og ængsteligere, men jo mere han stirrede, jo mindre saae han; hans Øie vilde stundom fyldes med Taarer, men jo oftere det hændte ham, jo mindre saae han; i Tidens Løb blev Skriften blegere og utydeligere, tilsidst hensmuldrede Papiret selv, og han beholdt intet Andet tilbage end et taareblændet Øie.
\end{quote}
\dcite{EE1:186  (SKS 2, 186)}
% EE1:186  (SKS 2, 186)

\begin{quote}
Alt dette har Goethe fuldkommen indseet, og derfor er Margrete en borgerlig lille Pige, en Pige, man næsten kunde fristes til at kalde ubetydelig. Vi ville nu, da det er af Vigtighed med Hensyn til Margretes Sorg, lidt nærmere overveie, hvorledes Faust vel maa have virket paa hende. De enkelte Træk, Goethe har fremhævet, have naturligviis stor Værdi; men dog troer jeg, at man for Fuldstændighedens Skyld maa tænke en lille Modification. I sin uskyldige Enfoldighed mærker Margrete snart, at det ikke hænger rigtigt sammen med Faust i Henseende til Troen. Hos Goethe fremtræder dette i en lille Katechisationsscene, der unegtelig er en fortræffelig Opfindelse af Digteren. Spørgsmaalet bliver nu, hvilke Følger denne Examination kan have paa deres Forhold til hinanden. Faust viser sig som Tvivleren, og det synes, at Goethe, da han intet Nærmere antyder i denne Henseende, har villet lade Faust vedblive at være Tvivler ogsaa ligeoverfor Margrete. Han har stræbt at bortlede hendes Opmærksomhed fra alle slige Undersøgelser og at fæste den ene og alene paa Kjærlighedens Realitet. Men deels troer jeg, at dette vilde falde Faust vanskeligt, da Problemet engang er kommet frem, deels troer jeg, at det ikke er psychologisk rigtigt. For Fausts Skyld skal jeg ikke nærmere dvæle ved dette Punkt, men vel for Margretes Skyld; thi dersom han ikke har viist sig som Tvivler for hende, har hendes Sorg et Moment mere. Faust er altsaa Tvivler, men han er ingen forfængelig Nar, der vil gjøre sig betydningsfuld ved at tvivle om det, Andre troe paa; hans Tvivl har en objectiv Grund i ham. Dette være sagt til Ære for Faust. Saasnart han derimod vil gjøre sin Tvivl gjældende mod Andre, kan der let indblande sig en ureen Lidenskab. Saasnart Tvivlen gjøres gjældende mod Andre, ligger der en Misundelse deri, der glæder sig ved at fravriste dem, hvad de ansaae for sikkert. Men for at denne Misundelsens Lidenskab skal vækkes hos Tvivleren, maa der kunne være Tale om en Modstand hos det paagjeldende Individ. Hvor der enten ikke kan være Tale derom eller hvor det endog vilde være uskjønt at tænke den, der ophører Fristelsen. Det Sidste er Tilfælde med en ung Pige. Ligeoverfor hende befinder en Tvivler sig altid i en Forlegenhed. At fravriste hende hendes Tro er ingen Opgave for ham; tvertimod føler han, at det kun er ved den hun er det Store, hun er. Han føler sig ydmyget; thi der ligger i hende en naturlig Fordring til ham at være hendes Forsørger, forsaavidt hun selv er bleven vaklende. Ja, en Stymper af en Tvivler, en halvstuderet Røver, han vil vel kunne finde en Tilfredsstillelse i at frarive en ung Pige hendes Tro, en Glæde i at skrække Koner og Børn, da han ikke kan forfærde Mænd. Men dette gjælder ikke om Faust; dertil er han for stor. Med Goethe kan man altsaa være enig i, at Faust første Gang forraader sin Tvivl, men derimod troer jeg neppe, det vil hænde ham anden Gang. Dette er af stor Vigtighed med Hensyn til Opfattelsen af Margrete. Faust seer lettelig, at Margretes hele Betydning beroer paa hendes uskyldige Enfoldighed; tages denne fra hende, da er hun Intet i sig, Intet for ham. Den maa altsaa bevares. Han er Tvivler, men som saadan har han alle det Positives Momenter i sig, thi ellers er han en daarlig Tvivler. Han mangler Slutningspunktet; herved blive alle Momenter til negative Momenter. Hun derimod har Slutningspunktet, har Barnligheden og Uskyldigheden. Intet er altsaa lettere for ham end at udstyre hende. Hans Praxis i Livet har ofte nok lært ham, at hvad han foredrog som Tvivl har paa Andre virket som positiv Sandhed. Nu finder han da sin Glæde i at berige hende med en Anskuelses rige Indhold, han tager den umiddelbare Troes hele Pynt frem, han finder en Glæde i at smykke hende dermed, fordi det anstaaer hende vel, og hun derved bliver skjønnere i hans Øine. Deraf drager han tillige den Fordeel, at hendes Sjæl knytter sig fastere og fastere til hans. Hun forstaaer ham egentlig slet ikke; som et Barn knytter hun sig fast til ham, hvad der for ham er Tvivl er for hende usvigelig Sandhed. Men medens han saaledes opbygger hendes Tro, undergraver han den paa samme Tid, thi han bliver tilsidst selv en Troesgjenstand for hende, en Gud, ikke et Menneske. Kun maa jeg her stræbe at forebygge en Misforstaaelse. Det kunde synes, at jeg gjør Faust til en nedrig Hykler. Dette er ingenlunde Tilfældet. Grete er selv den, der har bragt den hele Sag paa Bane; med et halvt Øie overskuer han den Herlighed, hun troer at eie, og seer, at den ikke kan bestaae for hans Tvivl, men han nænner ikke at tilintetgjøre den, og nu er det endog af en vis Godmodighed, han bærer sig saaledes ad. Hendes Kjærlighed giver hende Betydning for ham, og dog bliver hun næsten et Barn; han nedlader sig til hendes Barnlighed og har sin Glæde i at see, hvorledes hun \tlg{tilegner} sig Alt. For Margretes Fremtid har dette imidlertid de sørgeligste Følger. Havde Faust viist sig for hende som Tvivler, da havde hun maaskee senere kunnet frelse sin Tro, hun havde da i al Ydmyghed erkjendt, at hans høitflyvende og dristige Tanker ikke vare for hende, hun havde holdt fast ved hvad hun havde. Nu derimod skylder hun ham Troens Indhold, og dog indseer hun, da han har forladt hende, at han ikke selv har troet derpaa. Saalænge han var hos hende, opdagede hun ikke Tvivlen, nu han er borte, forandrer Alt sig for hende, og hun seer Tvivl i Alt, en Tvivl, hun ikke kan beherske, da hun bestandig tænker den Omstændighed med, at Faust selv har ikke kunnet magte den.
\end{quote}
\dcite{EE1:204  (SKS 2, 204)}
% EE1:204  (SKS 2, 204)

\begin{quote}
Hvis det lod sig gjøre, at staae bagved Cordelia, idet hun modtager et Brev fra mig, kunde det interessere mig meget. Jeg vilde da lettelig overbevise mig om, hvorvidt hun i egentligste Forstand erotisk \tlg{tilegner} sig dem. I det Hele er og bliver Breve altid et ubetaleligt Middel til at gjøre Indtryk paa en ung Pige; det døde Bogstav har ofte langt større Indflydelse end det levende Ord. Et Brev er en hemmelighedsfuld Communication; man er Herre over Situationen, føler intet Tryk af nogen Tilstedeværende, og sit Ideal, troer jeg, vil en ung Pige helst være ganske ene med, det vil sige i enkelte Øieblikke, og netop i de Øieblikke, hvor det virker stærkest paa hendes Sind. Om end hendes Ideal har fundet et nok saa fuldstændigt Udtryk i en bestemt, elsket Gjenstand, saa er der dog Momenter, i hvilke hun føler, at der er en Overvætteshed i Idealet, som Virkeligheden ikke har. Disse store Forsonings-Fester maae indrømmes hende; kun maa man passe paa, at benytte dem rigtigt, at hun ikke fra dem vender udmattet tilbage til Virkeligheden, men styrket. Dertil hjælpe Breve, der bevirke, at man usynlig er aandelig tilstede i disse hellige Indvielsens Øieblikke, medens Forestillingen om, at den virkelige Person er Forfatteren af Brevet, danner en naturlig og let Overgang til Virkeligheden.
\end{quote}
\dcite{EE1:402  (SKS 2, 402)}
% EE1:402  (SKS 2, 402)

\subsection{EE2 \normalfont(SKS 3)}

\begin{quote}
Vi ville nu see, hvorvidt det er lykkedes den Tid, der har tilintetgjort den romantiske Kjærlighed, at sætte noget bedre istedet. Først vil jeg dog angive Kjendetegnene paa den romantiske Kjærlighed. Man kunde med eet Ord sige, den er umiddelbar; at see hende og elske hende var eet, eller uagtet hun kun gjennem en Ridse af det tillukkede Jomfrubuurs Vindue saae ham en eneste Gang, saa elskede hun ham dog fra dette Øieblik, ham ene i hele Verden. Jeg burde nu vel egentlig efter Aftale her give Plads for nogle polemiske Udgydelser, for derved hos Dig at befordre den Secretion af Galde, der er en nødvendig Betingelse for en sund og gavnlig \tlg{Tilegnelse} af hvad jeg har at sige. Imidlertid kan jeg dog ikke beslutte mig dertil, og det af to Grunde, deels fordi det i vor Tid er temmelig forslidt, og oprigtigt talt er det ubegribeligt, at Du i denne Henseende vil være med Strømmen, da Du ellers altid er imod; deels fordi jeg virkelig har bevaret en vis Tro paa Sandheden deraf, en vis Ærbødighed derfor, en vis Vemod derved. Jeg nævner derfor blot Løsnet til Din Polemik i denne Retning, Overskriften over en lille Opsats af Dig, empfindsame og ubegribelige Sympathier eller to Hjerters harmonia præstabilita. Hvad Goethe med saa megen Kunst i hans Wahlverwandschaften først har ladet os ahne i Naturens Billedsprog for siden at realisere i Aandens Verden, det er det vi her tale om, kun at Goethe har bestræbt sig for at motivere denne Tiltrækning gjennem en Succession af Momenter (maaskee for at vise Forskjellen mellem Aandens Liv og Natur-Livet) og ikke udhævet den Hurtighed, den forelskede Utaalmodighed og Bestemthed, hvormed det Sammenhørende søger sammen. Og er det dog ikke skjønt saaledes at tænke sig, at to Væsener ere bestemte for hinanden! Hvor ofte har man dog ikke en Trang til at gaae ud over den historiske Bevidsthed, en Længsel, en Hjemvee efter den Urskov, der ligger bag ved os, og faaer denne Længsel ikke dobbelt Betydning, naar dertil knytter sig Forestillingen om et andet Væsen, der ogsaa har sit Hjem i hine Egne. Ethvert Ægteskab, endog det, der er indgaaet efter besindigt Overlæg, har derfor en Trang til, idetmindste i enkelte Øieblikke at tænke sig en slig Forgrund. Og hvor skjønt er det ikke, at den Gud, der er Aand, tillige elsker den jordiske Kjærlighed. At der nu blandt Ægtefolk lyves meget i denne Henseende, det vil jeg gjerne indrømme Dig, og at Dine Observationer i denne Retning ofte have moret mig, men det Sande deri bør man ikke glemme. Maaskee tænker En og Anden, at det dog er bedre at have fuldkommen Raadighed i Valget af »sit Livs Ledsagerinde«, men en saadan Yttring røber en høi Grad af Bornerethed og taabelig Forstands Vigtighed, og ahner ikke, at den romantiske Kjærlighed er fri i sin Genialitet, og at netop denne Genialitet er det Store i den.
\end{quote}
\dcite{EE2:29  (SKS 3, 29)}
% EE2:29  (SKS 3, 29)

\begin{quote}
Dette staaer altsaa fast mellem os: som Moment betragtet er den ægteskabelige Kjærlighed ikke blot ligesaa skjøn som den første men endnu skjønnere, fordi den i sin Umiddelbarhed indeholder en Eenhed i flere Modsætninger. Det er da ikke saa: Ægteskabet er en høist respektabel men kjedelig moralsk Person, Elskoven Poesi; nei Ægteskabet er ret egentlig det Poetiske. Og naar Verden saa ofte med Smerte har seet, at en første Kjærlighed ikke lod sig gjennemføre, saa vil jeg gjerne sørge med, men tillige erindre om, at Feilen ikke saa meget laae i det Senere, som i at man ikke begyndte rigtigt. Den første Kjærlighed mangler nemlig det andet æsthetiske Ideal, det romantiske. Den har ikke Bevægelses-Loven i sig. Dersom jeg tog Troen i det personlige Liv ligesaa umiddelbar, saa vilde til den første Kjærlighed svare en Tro, som i Kraft af Forjættelsen troede sig istand til at flytte Bjerge, og som nu vilde drage om og gjøre Undergjerninger. Maaskee vilde det lykkes den, men denne Tro havde ingen Historie; thi Ramsen paa alle dens Undergjerninger er ikke dens Historie, derimod er Troens \tlg{Tilegnelse} i det personlige Liv Troens Historie. Denne Bevægelse har den ægteskabelige Kjærlighed; thi i Forsættet er Bevægelsen vendt ind efter. I det Religiøse lader den ligesom Gud sørge for hele Verden, i Forsættet vil den i Forening med Gud stride for sig selv, erhverve sig selv i Taalmodighed. I Syndens Bevidsthed er der optaget en Forestilling om den menneskelige Skrøbelighed, men i Forsættet er den seet overvundet. Dette kan jeg ikke noksom indskærpe med Hensyn til den ægteskabelige Kjærlighed. Den første Kjærlighed har jeg visselig ladet vederfares al Ret, og jeg troer, jeg er en bedre Lovpriser af den end Du, men dens Feil ligger i dens abstrakte Charakteer.
\end{quote}
\dcite{EE2:99  (SKS 3, 99)}
% EE2:99  (SKS 3, 99)

\begin{quote}
Det Første Du vil nævne er »Vane, den uundgaaelige Vane, denne rædsomme Monotoni, det evindelige Einerlei i det ægteskabelige Huslivs ængstende Stilleben. Jeg elsker Natur, men er en Hader af den anden Natur.« Du veed, det maa man lade Dig, med forførerisk Varme og Veemod at skildre den lykkelige Tid, da man endnu gjør Opdagelser, med Angst og Forfærdelse at afmale den Tid, da det er forbi; Du veed til det Latterlige og Modbydelige at udpensle en ægteskabelig Eensformighed, som end ikke Naturen har Mage til; »thi her er dog, som allerede Leibnitz har viist, Intet ganske eens, en saadan Eensformighed er kun forbeholdt de fornuftige Skabninger, enten som Frugt af deres Søvnighed eller af deres Pedanteri.« Jeg har nu ingenlunde i Sinde at negte Dig, at det er en skjøn Tid, en evig uforglemmelig Tid, (læg vel Mærke til, i hvilken Betydning jeg kan sige dette) da Individet i Elskovens Verden forbauses ved og lyksaliggjøres i det, som vel for længe siden har været opdaget, som han vel ofte har hørt og læst om, men som han dog nu først \tlg{tilegner} sig med Overraskelsens hele Enthusiasme, med Inderlighedens hele Dybde; det er en skjøn Tid lige fra Kjærlighedens første Ahnelse, det første Syn, den første Forsvinden af den elskede Gjenstand, den første Accord af denne Stemme, det første Blik, det første Haandtryk, det første Kys lige til den første fuldkomne Forvisning om sin Besiddelse; det er en skjøn Tid, den første Uro, den første Længsel, den første Smerte, fordi hun udeblev, den første Glæde, fordi hun kom uventet, men hermed skal ingenlunde være sagt, at det Paafølgende ikke er ligesaa smukt. Du, der indbilder Dig at have en saa ridderlig Tænkemaade, prøv Dig selv. Naar Du siger det første Kys er det smukkeste, det sødeste, da fornærmer Du den Elskede; thi det, der altsaa giver Kysset absolut Værdi, er Tiden og dens Bestemmelse.
\end{quote}
\dcite{EE2:126  (SKS 3, 126)}
% EE2:126  (SKS 3, 126)

\begin{quote}
Dog, jeg vender tilbage til Betragtningen af hele Din aandelige Habitus. Du siger, at Du er en erobrende Natur, der ikke kan besidde. Naar Du siger dette, da mener Du vel ikke at have sagt Noget til Din egen Forringelse, tvertimod, Du føler Dig snarere større end Andre. Lad os betragte det lidt nærmere. Hvad hører der størst Kraft til, at gaae op ad Bakke eller ned ad Bakke? Naar Bakken er lige steil, hører der aabenbart mest Kraft til det Sidste. En Tilbøielighed til at gaae op ad Bakke er næsten ethvert Menneske medfødt, hvorimod de Fleste have en vis Angst for at gaae ned ad Bakke. Saaledes troer jeg ogsaa, at der er langt flere erobrende Naturer end besiddende, og naar Du føler Din Overlegenhed ligeoverfor en Mængde Ægtefolk og »deres dumme dyriske Tilfredshed«, saa kan det vel til en vis Grad være sandt, men Du skal jo heller ikke lære af dem, der staae under Dig. Den sande Kunst gaaer i Almindelighed den modsatte Vei af den, Naturen gaaer, uden at den derfor tilintetgjør denne, og saaledes vil ogsaa den sande Kunst vise sig i at besidde, ikke i at erobre; Besidden er nemlig en baglænds Erobren. I disse Udtryk seer Du allerede, hvorvidt Kunst og Natur modstræbe hinanden. Den, der besidder, han har jo ogsaa Noget, der er erobret, ja naar man vil være stræng i sit Udtryk, kan man sige, at først den, der besidder, han erobrer. Nu mener Du vel og at besidde; thi Du har jo Besiddelsens Øieblik, men det er ingen Besiddelse; thi det er ingen dybere \tlg{Tilegnelse}. Naar jeg saaledes vilde tænke mig en Erobrer, der undertvang Riger og Lande, saa besad han jo ogsaa disse undertrykte Provindser, han havde store Besiddelser, og dog kalder man jo en saadan Fyrste en erobrende og ikke en besiddende. Først naar han med Viisdom ledede disse Lande til deres eget Bedste, først da besad han dem. Dette er nu meget sjeldent med en erobrende Natur, i Almindelighed vil han mangle den Ydmyghed, den Religiøsitet, den sande Humanitet, der hører til at besidde. Seer Du, derfor var det, at jeg ved Udviklingen af Ægteskabets Forhold til den første Kjærlighed netop fremhævede det religiøse Moment, fordi dette vil dethronisere Erobreren og lade Besidderen komme tilsyne; derfor anpriste jeg det, at Ægteskabets Anlæg netop var beregnet paa det Høieste, paa den varige Besiddelse. Her kan jeg erindre Dig om et Ord, Du ofte nok slaaer om Dig med: »ikke det Oprindelige er det Store, men det Erhvervede«; thi det Erobrende i et Menneske, og det, at han gjør Erobring, det er egentlig det Oprindelige, men det, at han besidder og vil besidde, det er det Erhvervede. Til at erobre hører der Stolthed, til at besidde Ydmyghed, til at erobre hører der Hæftighed, til at besidde Taalmod, til at erobre – Begjærlighed; til at besidde – Nøisomhed; til at erobre hører der Spise og Drikke, til at besidde Bøn og Fasten. Men alle de Prædikater, jeg her, og dog vel med Rette, har brugt for at charakterisere den erobrende Natur, lade sig alle anvende paa og passe absolut paa det naturlige Menneske; men det naturlige Menneske er ikke det Høieste. En Besiddelse er nemlig ikke et aandeligt dødt og ugyldigt om end juridisk kraftigt »Schein«, men en stadig Erhvervelse. Her seer Du atter, at den besiddende Natur har den erobrende i sig; han erobrer nemlig som en Landmand, der ikke sætter sig i Spidsen for sine Karle, og fordriver sin Nabo, men han erobrer ved at grave i Jorden. Det sande Store er altsaa ikke at erobre, men at besidde. Naar Du nu her vil sige: »jeg vil ikke afgjøre, hvad der er det Største, men jeg vil gjerne indrømme, at det er de to store Formationer af Mennesker; Enhver maa nu afgjøre med sig selv, hvilken han tilhører, og nu vogte sig for, ikke at lade sig omkalfatre af en eller anden Omvendelsens-Apostel«, saa føler jeg nok, at Du med det sidste Udtryk saa smaat har mig i Kikkerten. Jeg vil imidlertid svare, det Ene er ikke blot større end det Andet, men det Ene er der Mening i, det er der ikke i det Andet. Det Ene har baade Forsætning og Eftersætning, det Andet er blot Forsætning, og istedetfor Eftersætning en betænkelig Tankestreg, hvis Betydning jeg en anden Gang skal forklare Dig, hvis Du ikke allerede veed den.
\end{quote}
\dcite{EE2:131  (SKS 3, 131)}
% EE2:131  (SKS 3, 131)

\begin{quote}
Du giver Dig jo ogsaa af med at lovprise Elskoven. Jeg vil da ikke berøve Dig, hvad der vel ikke er Din Eiendom, thi det er jo Digterens, men hvad Du dog har \tlg{tilegnet} Dig; men da ogsaa jeg har \tlg{tilegnet} mig det, saa lad os dele, Du faaer det hele Vers, jeg det sidste Ord: og saa Uskyldighed.
\end{quote}
\dcite{EE2:143  (SKS 3, 143)}
% EE2:143  (SKS 3, 143)

\begin{quote}
Da Mystikeren vælger sig selv abstrakt, saa er hans Ulykke, at han har saa vanskeligt ved at komme i Bevægelse, eller rettere sagt, at det er ham en Umulighed. Som det gaaer Dig med Din jordiske første Kjærlighed, saaledes gaaer det Mystikeren med hans religiøse første Kjærlighed. Han har smagt dens hele Salighed, og har nu Intet at gjøre uden at vente paa, om den vil komme igjen i ligesaa megen Herlighed, og herom kan han let fristes til at nære en Tvivl, der er den saa ofte af mig paapegede, at Udvikling er Tilbagegang, er en Sætten af. For en Mystiker er Virkeligheden en Forsinkelse, ja af saa betænkelig Art, at han næsten løber Fare for, at Livet berøver ham, hvad han engang har eiet. Vilde man derfor spørge en Mystiker, hvad er Livets Betydning, saa vilde han maaskee svare: det er Livets Betydning at lære Gud at kjende, at forelske sig i ham. Dette er imidlertid ikke Svar paa Spørgsmaalet; thi her er Livets Betydning opfattet som Moment, ikke som Succession. Naar jeg derfor vilde spørge ham, hvilken Betydning har det for Livet, at Livet har havt denne Betydning, eller med andre Ord, hvilken er Timelighedens Betydning, saa har han ikke Stort at svare, i ethvert Tilfælde ikke meget Glædeligt. Siger han, at Timeligheden er en Fjende, der skal overvindes, saa maatte man da nærmere spørge, om det ingen Betydning skulde have, at denne Fjende blev overvunden. Dette mener Mystikeren egentlig ikke, og helst vilde han dog være færdig med Timeligheden. Som han derfor miskjendte Virkeligheden og fik den metaphysisk opfattet som Forfængelighed, saa miskjender han nu det Historiske og faaer det metaphysisk opfattet som unyttig Møie. Den høieste Betydning, han da kan tillægge Timeligheden, er, at den er en Prøvetid, i hvilken man atter og atter gjør Prøve, uden at der dog egentlig resulterer Noget deraf, eller at man er kommen videre, end man var i Begyndelsen. Dette er imidlertid en Miskjendelse af Timeligheden, thi vel beholder den altid Noget af en ecclesia pressa i sig, men den er tillige Muligheden af den endelige Aands Forherligelse. Det er netop det Skjønne ved Timeligheden, at den uendelige og den endelige Aand i den skilles ad, og det er netop den endelige Aands Storhed, at Timeligheden er den anviist. Timeligheden er derfor ikke til, om jeg saa tør sige, for Guds Skyld, for at han i den, at jeg skal tale mystisk, kan prøve og forsøge den Elskende, men den er til for Menneskets Skyld og er den største Naadegave af alle. Deri ligger nemlig et Menneskes evige Værdighed, at han kan faae Historie, deri ligger det Guddommelige i ham, at han selv, hvis han vil, kan give denne Historie Continuitet; thi det faaer den først, naar den ikke er Indbegrebet af hvad der er skeet eller hændt mig, men min egen Gjerning, saaledes, at selv det, der er hændt mig, ved mig er forvandlet og overført fra Nødvendighed til Frihed. Dette er det Misundelsesværdige ved et Menneskeliv, at man kan komme Guddommen til Hjælp, kan forstaae ham, og det er igjen den eneste et Menneske værdige Maade at forstaae ham paa, at man i Frihed \tlg{tilegner} sig Alt, hvad der tilkommer Een, baade det Glade og det Sørgelige. Eller synes det Dig ikke saa? Mig forekommer det saaledes, ja, jeg synes, at man blot behøvede at sige det høit til et Menneske, for at gjøre ham misundelig paa sig selv.
\end{quote}
\dcite{EE2:239  (SKS 3, 239)}
% EE2:239  (SKS 3, 239)

\subsection{G \normalfont(SKS 4)}

\begin{quote}
Dersom jeg ikke havde Job! Det er umuligt at beskrive og nuancere, hvilken Betydning og hvor mangfoldig en Betydning han har for mig. Jeg læser ham ikke som man læser en anden Bog med Øiet, men jeg lægger Bogen ligesom paa mit Hjerte, og med Hjertets Øie læser jeg den, forstaaer i en clairvoyance det Enkelte paa den forskjelligste Maade. Som Barnet lægger Lærebogen under sit Hoved, for at være sikker paa, at han ikke har glemt sin Lectie, naar han vaagner om Morgenen, saaledes tager jeg Bogen med mig om Natten i Sengen. Hvert Ord af ham er Føde og Klæde og Lægedom for min elendige Sjæl. Snart vækker et Ord af ham mig af min Lethargi, saa jeg vaagner til ny Uro, snart stiller det den ufrugtbare Rasen i mit Indre, ender det Rædsomme i Lidenskabens stumme Qvalmen. De har dog læst Job? Læs den, læs den atter og atter. Jeg nænner ikke engang at skrive et eneste Udbrud af i et Brev til Dem, skjøndt jeg har min Glæde af at tage ny og ny Afskrift af Alt, hvad han har sagt, snart med danske, snart med latinske Bogstaver, snart i een Format snart i en anden. Enhver saadan Afskrift bliver lagt som et Gudshaandsplaster paa mit syge Hjerte. Og paa hvem laae vel Guds Haand som paa Job! Men citere ham – det kan jeg ikke. Det var at ville give min Besyv med, det var at ville gjøre hans Ord til mine i en Andens Nærværelse. Naar jeg er alene, da gjør jeg det, \tlg{tilegner} mig Alt; men saasnart Nogen er tilstede, da veed jeg vel, hvad et ungt Menneske har at gjøre, naar gamle Folk tale.
\end{quote}
\dcite{G:72.3  (SKS 4, 72)}
% G:72.3  (SKS 4, 72)

\subsection{PS \normalfont(SKS 4)}

\begin{quote}
Fordelen af Conseqventserne synes at ligge i at hiint Faktum lidt efter lidt skulde være blevet naturaliseret. Er dette Tilfældet (det vil sige lader dette sig tænke), da er den senere Generation endog ligefrem i Fordelen for den samtidige (og det Menneske maatte da være meget dumt, der kunde tale om Conseqventsen i denne Forstand og dog phantasere om det Lykkelige i at være Samtidig med hiint Faktum), og kan \tlg{tilegne} sig hiint Faktum høist ugeneert uden at formærke Noget til den Opmærksomhedens Tvetydighed, af hvilken Forargelsen kan fremgaae saavel som Troen. Hiint Faktum respecterer imidlertid ikke nogen Dressur, er for stolt til at ønske en Discipel, der vil være med, i Kraft af det heldige Udfald som Sagen fik, forsmaaer at naturaliseres under en Konges eller en Professors Protection; det er og bliver Paradoxet og lader sig ikke tilspeculere. Hiint Faktum er kun for Troen. Troen kan nu vel blive til anden Natur i et Menneske, men det Menneske, hvis anden Natur den bliver, han maa dog vel have havt en første, siden Troen blev den anden. Skal hiint Faktum naturaliseres, da lader dette sig i Forhold til Individet udtrykke saaledes, at Individet fødes med Troen, det vil sige med sin anden Natur. Dersom vi begynde saaledes vor Udvikling, saa begynder samtidigen alt muligt Galimathias at jubilere; thi nu er det jo slaaet løs og ikke mere til at standse. Dette Galimathias maa naturligviis være opfundet ved at gaae videre, thi i Socrates' Anskuelse var der sandeligen god Mening, om vi end forlode det for at opdage det tidligere Projekterede, og et saadant Galimathias vilde vel føle det som en dyb Fornærmelse, at det ikke skulde være langt videre end det Socratiske. Selv i en Sjelevandring er der dog Mening, men at fødes med sin anden Natur, en anden Natur der refererer sig til et i Tiden givet historisk Faktum, det er et sandt non plus ultra af Galenskab. Socratisk forstaaet, har Individet været til før det blev til, og erindrer sig selv, saa Erindringen er Præexistentsen (ikke Erindring om Præexistentsen); Naturen (den ene; thi her er ikke Tale om en første og anden Natur) er bestemmet i Continuitet med sig selv. Her derimod er Alt forlænds og historisk, saaledes at det at fødes med Troen i Grunden er ligesaa plausibelt som at fødes 24 Aar gammel. Skulde man virkelig kunne paavise et Individ, der var født med Troen, da var dette et mere seeværdigt Monstrum end det Barberen i den Stundesløse fortæller at være født in den neuen Buden, om det end vilde forekomme Barberer og Stundesløse at være et allerkjæreste lille Væsen, Speculationens høieste Triumph. – Eller fødes maaskee Individet med begge Naturer paa eengang, vel at mærke ikke to saadanne Naturer der høre sammen for at danne den almindelige menneskelige Natur, men med to hele menneskelige Naturer, hvoraf den ene forudsætter et mellemliggende Historisk. I saa Fald er Alt forvirret hvad vi projekterede i Capitel I, vi staae ei heller ved det Socratiske, men ved en Confusion, som end ikke Socrates vilde være istand til at hæve. Det bliver en Forvirring forlænds, der har meget tilfælleds med den af Apollonius af Tyana opfundne baglænds. Han nøiedes nemlig ikke som Socrates med at erindre sig selv som værende før han blev til (Bevidsthedens Evighed og Continuitet er det Dybsindige og Tanken i det Socratiske), men var hurtig til at gaae videre, han erindrede nemlig, hvem han havde været før han blev sig selv. Dersom hiint Faktum er blevet naturaliseret, saa er Fødselen ikke mere Fødselen, men den er tillige Igjenfødelsen, saaledes at den, der aldrig har været, han fødes igjen – idet han fødes. – I det individuelle Liv udtrykkes det saaledes, at Individet fødes med Troen; i Slægten maa det Samme udtrykkes saaledes, at Slægten efter hiint Faktums Indtræden er bleven en ganske anden, og desuagtet er bestemmet i Continuitet med den første. I saa Fald burde Slægten antage et nyt Navn; thi Troen, saaledes som vi har projekteret den, er vel intet Umenneskeligt, som en Fødsel indenfor en Fødsel (Igjenfødelsen), men vel derimod blev den et eventyrligt Monstrum, hvis det var som vi have ladet Indvendingen ville det.
\end{quote}
\dcite{PS:293  (SKS 4, 293)}
% PS:293  (SKS 4, 293)

\begin{quote}
Den første Generation af secundaire Disciple har da den Fordeel, at Vanskeligheden er der; thi det er altid en Fordeel, en Lettelse, naar det er det Vanskelige, jeg skal \tlg{tilegne} mig, at det bliver gjort mig vanskeligt. Skulde den sidste Generation, naar den betragter den første og seer den næsten segnende under Forfærdelsen, falde paa at sige: det er ubegribeligt; thi det Hele er ikke sværere end at man kan tage og løbe med det; da var der vel Den, som vilde svare: »vær saa god, løb Du kun, men Du skulde dog see efter, om Det, Du løber med, virkelig er Det, hvorom Talen er; og derom disputere vi jo ikke, at en Vind er let nok at løbe med«.
\end{quote}
\dcite{PS:296.2  (SKS 4, 296)}
% PS:296.2  (SKS 4, 296)

\subsection{BA \normalfont(SKS 4)}

\begin{quote}
Egentlig hører Synden slet ikke hjemme i nogen Videnskab. Den er Prædikenens Gjenstand, hvor den Enkelte taler som den Enkelte til den Enkelte. I vor Tid har den videnskabelige Vigtighed faaet narret Præsterne til at være et Slags Professor-Degne, der ogsaa tjene Videnskaben og finde det under deres Værdighed at prædike. Forsaavidt er det da intet Under, at det at prædike er blevet anseet for en meget fattig Kunst. At prædike er imidlertid den vanskeligste af alle Kunster og er egentlig den Kunst, som Socrates anpriser: at kunne samtale. Det følger af sig selv, at derfor behøver der ingenlunde Een at svare i Menigheden, eller at det skulde hjælpe bestandig at indføre Een talende. Det, Socrates egentlig dadlede hos Sophisterne under den Distinction: at de vel kunde tale, men ikke samtale, var, at de om enhver Ting kunde sige meget, men manglede \tlg{Tilegnelsens} Moment. \tlg{Tilegnelsen} er netop Samtalens Hemmelighed.
\end{quote}
\dcite{BA:323.2  (SKS 4, 323)}
% BA:323.2  (SKS 4, 323)

\begin{quote}
Med Syndigheden blev Sexualiteten sat. I samme Øieblik begynder Slægtens Historie. Som da Syndigheden i Slægten bevæger sig i quantitative Bestemmelser, saa gjør Angesten det ogsaa. Arvesyndens Følge eller Arvesyndens Tilstedeværen i den Enkelte er Angest, der kun er quantitativt forskjellig fra Adams. I Uskyldighedens Tilstand, og en saadan maa der jo ogsaa kunne være Tale om i det senere Menneske, maa Arvesynden have den dialektiske Tvetydighed, ud af hvilken Skylden bryder i det qualitative Spring. Derimod vil Angesten i et senere Individ kunne være mere reflekteret, end i Adam, fordi den quantitative Tilvæxt, som Slægten tilbagelægger, nu gjør sig gjeldende i ham. Angest bliver dog her saa lidt som nogensinde en Ufuldkommenhed hos Mennesket, og man maa tvertimod sige, jo oprindeligere et Menneske, jo dybere Angest, fordi den Syndighedens Forudsætning, som hans individuelle Liv supponerer, da han jo kommer ind i Slægtens Historie, maa \tlg{tilegnes}. Forsaavidt har Syndigheden faaet en større Magt, og Arvesynden er voxende. At der gives Mennesker, der slet ingen Angest mærke, maa forstaaes ligesom, at Adam ingen vilde have fornummet, hvis han havde været blot Dyr.
\end{quote}
\dcite{BA:357.4  (SKS 4, 357)}
% BA:357.4  (SKS 4, 357)

\begin{quote}
Skulde jeg her i een eneste Sætning udtrykke det Mere, der er for ethvert senere Individ i Forhold til Adam, da vilde jeg sige det er: at Sandseligheden kan betyde Syndighed, det vil sige, denne dunkle Viden derom, samt en dunkel Viden om hvad Synden forøvrigt kan betyde, samt en misforstaaet historisk \tlg{Tilegnelse} af det Historiske de te fabula narratur, hvorved Pointen, den individuelle Oprindelighed udelades, og Individet uden videre forvexler sig selv med Slægten og dens Historie. Vi sige ikke, at Sandselighed er Syndighed, men at Synden gjør den dertil. Naar vi nu tænke os det senere Individ, da har ethvert saadant en historisk Omgivelse, i hvilken det kan vise sig, at Sandseligheden kan betyde Syndighed. For Individet selv betyder den det ikke, men denne Viden giver Angesten et Mere. Aanden er da ikke blot sat i Forhold til Sandselighedens Modsætning, men til Syndighedens. Det følger af sig selv, at det uskyldige Individ endnu ikke forstaaer denne Viden; thi den forstaaes først qualitativt, men denne Viden er dog igjen en ny Mulighed, saaledes at Friheden i sin Mulighed forholdende sig til det Sandselige bliver end mere Angest.
\end{quote}
\dcite{BA:377.5  (SKS 4, 377)}
% BA:377.5  (SKS 4, 377)

\begin{quote}
At ogsaa Videnskaben ligesaa fuldt som Poesie og Kunst forudsætter Stemning baade hos den Producerende og den Reciperende, at en Feil i Modulationen er ligesaa forstyrrende som en Feil i Tankens Udvikling, har man aldeles glemt i vor Tid, hvor man aldeles har glemt Inderligheden og \tlg{Tilegnelsens} Bestemmelse af Glæde over al den Herlighed, man meente at eie, eller har i Begjerlighed renonceret paa den ligesom Hunden, der foretrak Skyggen. Dog enhver Feil føder sin egen Fjende. Tænkningens Feil har Dialektiken uden for sig, Stemningens Udebliven eller Forfalskning har det Comiske udenfor sig som Fjende.
\end{quote}
\dcite{BA:461.4  (SKS 4, 461)}
% BA:461.4  (SKS 4, 461)

\subsection{2T43 \normalfont(SKS 5)}

\begin{quote}
»Al god og al fuldkommen Gave er ovenfra, og kommer ned fra Lysenes Fader, hos hvilken er ikke Forandring eller Skygge af Omskiftelse.« Disse Ord ere saa husvalende og lindrende, og dog hvor Mange vare de, der ret forstode at suge Trøstens rige Næring af dem, ret forstode at \tlg{tilegne} sig denne! De Bekymrede, de hvem Livet ikke tillod at ældes og døe som Børn, hvem det ikke diede ved Medgangs Melk, men som tidlig bleve vendte fra; de Sørgende, hvis Tanke søgte at trænge gjennem det Vexlende til det Blivende, de fornam Apostelens Ord og agtede paa dem. Jo mere de formaaede at sænke deres Sjæl i dem, at glemme Alt over dem, desto mere følte de sig styrkede og tillidsfulde. Dog viste det sig snart, at denne Styrke var en Skuffelse; hvor megen Tillid de end vandt, de vandt dog ikke Kraft til at gjennemtrænge Livet, snart søgte igjen det bekymrede Sind og den raadvilde Tanke hen til hiin rige Trøst, snart fornam de igjen Modsigelsen. Tilsidst forekom det dem maaskee, at disse Ord næsten vare farlige for deres Fred, de vakte en Tillid i dem, som bestandig blev skuffet, de gave dem Vinger, som vel kunde hæve dem til Gud, men ikke hjælpe dem i deres Gang gjennem Livet; de negtede ikke Ordenes uudtømmelige Trøst, men de frygtede den næsten, om de end priste den. Hvis et Menneske eiede et herligt Smykke, og der kom ingen Tid, da han vilde negte, at det var herligt, saa tog han det vel af og til frem, glædede sig derover, men snart sagde han: til daglig Brug kan jeg dog ikke pryde mig dermed, og den festlige Anledning, ved hvilken det ret skulde faae sin Betydning, den venter jeg forgjeves paa. Saa vilde han vel lægge Smykket til en Side og med Veemod tænke paa, at han eiede et saadant Smykke og at Livet ikke forundte ham Leilighed til ret med Glæde at tage det frem.
\end{quote}
\dcite{2T43:44  (SKS 5, 44)}
% 2T43:44  (SKS 5, 44)

\begin{quote}
Eller, min Tilhører, var der maaskee i Dit Liv ingen Anledning til at finde hine Ord vanskelige? Var Du altid tilfreds med Dig selv, saa tilfreds, at Du maaskee takkede Gud, at Du ikke var som andre Mennesker? Var Du maaskee bleven saa klog, at Du fattede den dybe Mening i den meningsløse Tale, at det var herligt ikke at være som andre Mennesker?.... Hvad var det da, der gjorde Dig dem vanskelige? Dersom et Menneske selv var en god og en fuldkommen Gave, dersom han blot forholdt sig modtagende og modtog Alt af Guds Haand, ja! hvorledes skulde han da kunne modtage Andet end gode og fuldkomne Gaver? Men da Du bøiede Dig under Menneskets almindelige Lod, da tilstod Du jo, at Du hverken var god eller fuldkommen, at Du ikke forholdt Dig blot modtagende, men at der med Alt, hvad Du modtog, foregik en Forvandling. Kan da det Lige forstaaes af Andet end af det Lige, kan det Gode blive godt i Andet end i det Gode; kan den sunde Føde bevare sin Sundhed i den syge Sjæl? Et Menneske forholder sig ikke reent modtagende, han er selv meddelende, og det blev vanskeligt for Dig at forstaae, hvorledes det Usunde, der kom fra Dig, kunde blive til andet end Skade for Andre. Vel forstod Du, at det var kun ved Taksigelse til Gud, at Alt for Dig blev en god og en fuldkommen Gave, Du holdt det fast, at paa samme Maade maatte det andet Menneske ogsaa \tlg{tilegne} sig Alt; men var da selv den Kjærlighed, der fødte Taksigelsen, var den reen, forvandlede den ikke det Modtagne? Kan da et Menneske gjøre mere end at elske? Har Tanken og Sproget noget høiere Udtryk for det at elske end det altid at takke? Ingenlunde, den har et lavere, et ydmygere; thi selv Den, der altid vil takke, elsker dog efter sin Fuldkommenhed, og Gud kan et Menneske i Sandhed kun elske, naar han elsker ham efter sin Ufuldkommenhed. Hvilken er denne Kjærlighed? det er Angerens, der er skjønnere end al anden Kjærlighed; thi i den elsker Du Gud! trofastere og inderligere, end al anden Kjærlighed; thi i Angeren er det Gud, der elsker Dig. I Angeren modtager Du Alt af Gud, selv Taksigelsen, som Du bringer ham, saa at selv denne er hvad Barnets Gave er i Forældrenes Øine, en Spøg, en Modtagen af Noget, man selv har givet. Var det ikke saaledes, min Tilhører? Du vilde altid takke Gud, men selv dette var saa ufuldkomment. Da forstod Du, at Gud er den, som gjør Alt i Dig, og som da forunder Dig den barnlige Glæde over at han betragter Din Taksigelse som en Gave fra Dig. Denne Glæde skænker han Dig, naar Du ikke frygtede Angerens Smerte, ikke den dybe Sorg, i hvilken et Menneske bliver glad som et Barn i Gud, naar Du ikke frygtede at forstaae, at dette er Kjærligheden, ikke at vi elske Gud, men at Gud elsker os.
\end{quote}
\dcite{2T43:54  (SKS 5, 54)}
% 2T43:54  (SKS 5, 54)

\subsection{3T43 \normalfont(SKS 5)}

\begin{quote}
Det beroer da ikke blot paa hvad man seer, men det hvad man seer beroer paa, hvorledes man seer; thi al Betragtning er ikke blot en Modtagen, en Opdagen, men tillige en Frembringen, og forsaavidt den er dette, da bliver det jo afgjørende, hvorledes den Betragtende selv er. Naar da Een seer Eet, en Anden et Andet i det Samme, saa opdager den Ene hvad den Anden skjuler. Forsaavidt det, der er Gjenstand for Betragtningen, tilhører den udvortes Verden, da er det vel ligegyldigere, hvorledes den Betragtende er, eller rettere, da er det til Betragtningen Nødvendige noget hans dybere Væsen Uvedkommende; jo mere derimod Betragtningens Gjenstand tilhører Aandens Verden, desto vigtigere er det, hvorledes han selv er i sit inderste Væsen; thi alt Aandeligt \tlg{tilegnes} kun ved Frihed; men hvad der \tlg{tilegnes} ved Frihed, det frembringes ogsaa. Forskjellen ligger da ikke i det Udvortes men i det Indvortes, og indvortes fra udgaaer Alt det som gjør et Menneske ureent, og hans Betragtning ureen. Det udvortes Øie gjør Intet til Sagen, men »indvortes fra udgaaer et ondt Øie.« Men et ondt Øie opdager meget, som Kjerligheden ikke seer; thi et ondt Øie seer endog, at Herren gjør Uretfærdighed, naar han er god. Naar der i Hjertet boer Ondskab, da seer Øiet Forargelse, men naar der i Hjertet boer Reenhed, da seer Øiet Guds Finger; thi de Rene see altid Gud; »men den, der gjør ondt, seer ikke Gud« (3 Joh. 11).
\end{quote}
\dcite{3T43:70  (SKS 5, 70)}
% 3T43:70  (SKS 5, 70)

\begin{quote}
Den, der har en Lære at anbefale Menneskene og stræber at vinde dem, han har jo et Vidnesbyrd, til hvilket han trøstigt henviser den Enkelte. Men naar dette Vidnesbyrd svigter, da indseer han vel, at Magten er tagen fra ham, og, skjønt det er saare tungt, forliger han sig dog med Gud i sit Hjerte; og sørger vel som Den, hvem Brudgommen forlod og Glæden, men ogsaa som Den, der ikke løb paa det Uvisse, der ikke glemmer, at høiere end at frelse Andre, er det at frelse sin egen Sjæl, at underlægge det urolige Sind under Troens Lydighed, at holde de vildfarende Tanker bundne ved Overbeviisningens Magt i Kjerlighedens Baand. Det er velgjørende at tale derom, og ethvert redeligt Menneske bekjender vel, at det er saligt saaledes at beskikke sit eget Huus, naar man har udtjent ved den større Gjerning, og bliver sat over den mindre. – Nu Paulus! Levede han i de Mægtiges Gunst, at den kunde anbefale hans Lære? Nei! han var fangen. Hyldede de Vise hans Lære, at deres Anseelse kunde borge for Sandheden? Nei den var dem en Daarskab. Formaaede hans Lære hurtigt at sætte den Enkelte i Besiddelse af en overnaturlig Magt, falbød den sig til Menneskene ved Gjøgleværk? Nei! den maatte langsomt erhverves, \tlg{tilegnes} i Prøvelse, der begyndte med Forsagelse af Alt. Havde Paulus da noget Vidnesbyrd? Ja! han havde ethvert menneskeligt Vidnesbyrd mod sig, og dertil havde han endnu den Bekymring, at Menigheden skulde forsage, eller hvad værre var, forarges paa ham; thi Forargelsen ligger vel aldrig nærmere end, naar Sandhed nedkues, naar Uskyld lider, naar Uretfærdighed er tryg med sin Seier, naar Vold har Fremgang, naar Vankundighed end ikke behøver at bruge Magt mod det Gode, men sorgløs og ubekymret forbliver uvidende om, at det er til. Men forsager Paulus forladt af Vidnesbyrdet? Ingenlunde. Da han intet andet Vidnesbyrd havde at beraabe sig paa, da beraaber han sig paa sine Trængsler. Er dette ikke som en Undergjerning? Hvis Paulus ellers ikke havde beviist krafteligen, at han havde Underets Magt, er det ikke et Beviis? At forvandle Trængsler til et Vidnesbyrd for Lærens Sandhed, at forvandle Skjendsel til en Ære for sig og for den troende Menighed; at forvandle den tabte Sag til en Æressag, der har Vidnesbyrdets hele begeistrende Magt, er det ikke som at bringe den Halte til at gaae, og gjøre den Stumme talende!
\end{quote}
\dcite{3T43:90.2  (SKS 5, 90)}
% 3T43:90.2  (SKS 5, 90)

\subsection{4T43 \normalfont(SKS 5)}

\begin{quote}
Ikke blot den kalde vi en Menneskenes Lærer, der ved en særdeles lykkelig Begunstigelse opdagede, eller ved utrættelig Møie og gjennemgribende Vedholdenhed udgrundede en eller anden Sandhed, efterlod det Erhvervede som en Lærdom, hvilken de følgende Slægter stræbe at forstaae og i denne Forstaaelse at \tlg{tilegne} sig; men ogsaa den, maaskee i endnu strængere Forstand, en Menneskehedens Lærer, der ingen Lære havde at overgive Andre, men kun efterlod Slægten sig selv som et Forbillede, sit Liv som en Veiledning for ethvert Menneske, sit Navn som en Betryggelse for de Mange, sin Gjerning som en Opmuntring for de Forsøgte. En saadan Menneskehedens Lærer og Veileder er Job, hvis Betydning ingenlunde ligger i hvad han har sagt, men i hvad han har gjort. Vel har han efterladt sig et Udsagn, som ved sin Korthed og Skjønhed er blevet et Ordsprog, bevaret fra Slægt til Slægt, og Ingen har formastelig føiet Noget til eller taget Noget fra; men Udsagnet selv er ikke det Veiledende, og Jobs Betydning ligger ikke i, at han sagde det, men i at han efterkom det i Gjerning. Ordet selv er vel skjønt og værd at overveie, men dersom en Anden havde sagt det, eller dersom Job havde været en Anden, eller dersom han havde sagt det ved en anden Leilighed, da var ogsaa Ordet selv blevet et andet, betydningsfuldt, hvis det ellers var det, som udtalt, men ikke betydningsfuldt derved, at han handlede ved at udsige det, at Udsigelsen selv var en Handling. Dersom Job havde anvendt hele sit Liv paa at indskærpe dette Ord, dersom han havde betragtet det som Summen og Fuldendelsen af, hvad et Menneske bør lade Livet lære sig, dersom han bestandig blot havde lært det fra sig, men aldrig selv forsøgt det, aldrig selv handlet, idet han udsagde det, saa er Job en Anden, hans Betydning en anden. Da vilde Jobs Navn være glemt, eller det vilde dog være ligegyldigt, om man vidste det, Hovedsagen var Ordets Indhold, den Tankefylde, der laae i det. Hvis Slægten havde annammet Ordet, da var det dette, hvilket den ene Slægt overgav til den anden; medens det nu derimod er Job selv, der ledsager Slægten. Naar da den ene Slægt har udtjent, har fuldkommet sin Gjerning, udkæmpet sin Strid, da har Job ledsaget den, naar den nye Slægt med sine uoverskuelige Rækker og hver Enkelt i disse paa sin Plads staaer færdig til at begynde Vandringen, da er Job atter tilstede, indtager sin Plads, hvilken er Menneskehedens Yderpost. Seer Slægten kun glade Dage i lykkelige Tider, da følger Job trofast med, og hvis den Enkelte dog i Tanken oplever det Forfærdelige, ængstes ved Forestillingen om, hvad Livet kan gjemme af Rædsel og Nød, om, at Ingen veed, naar Fortvivlelsens Time slaaer for ham, da søger hans bekymrede Tanke hen til Job, hviler sig paa ham, beroliges ved ham; thi han følger troligen med og trøster vel ikke saaledes, som havde han engang for alle lidt, hvad aldrig siden skulde lides, men trøster som Den, der vidner, at det Forfærdelige er lidt, at Rædselen er oplevet, at Fortvivlelsens Kamp er stridt, Gud til Ære, ham til Frelse, Andre til Gavn og Glæde. I glade Dage, i lykkelige Tider gaaer Job ved Slægtens Side, og betrygger den sin Glæde, bekæmper den angstfulde Drøm, at en pludselig Rædsel skulde overfalde et Menneske, og have Magt til at myrde hans Sjel som sit visse Bytte. Kun den Letsindige kunde ønske, at Job ikke var med, at hans ærværdige Navn ikke skal minde ham om, hvad han søger at glemme, at der er Forfærdelse til i Livet og Angest; kun den Selvkjerlige kunde ønske, at Job ikke var til, at Forestillingen om hans Lidelse ikke skal med sit strænge Alvor forstyrre hans skrøbelige Glæde, skrække ham ud af hans i Forhærdelse og Fortabelse berusede Tryghed. I stormfulde Tider, naar Tilværelsens Grundvold vakler, naar Øieblikket skjælver i angstfuld Forventning af hvad der maa komme, naar enhver Forklaring forstummer ved Synet af det vilde Oprør, naar Menneskets Inderste vaander sig i Fortvivlelse og »i Sjelens Bitterhed« skriger mod Himlen, da gaaer Job endnu ved Slægtens Side og borger for, at der er en Seier, borger for, at om end den Enkelte taber i Striden, saa er der dog en Gud, der, som han gjør enhver Fristelse menneskelig, selv om et Menneske ikke bestod i Fristelsen, dog skal gjøre dens Udgang, saa vi den kunne bære, ja herligere end nogen menneskelig Forventning. Kun den Trodsige kunde ønske, at Job ikke var til, at han ganske kunde frigjøre sin Sjel fra den sidste Kjerlighed, der dog blev tilbage i Fortvivlelsens Klageskrig, at han kunde klage, ja forbande Livet saaledes, at der end ikke var nogen Medlyd af Tro og Tillid og Ydmyghed i hans Tale, at han i sin Trods kunde qvæle Skriget, at det end ikke skulde synes, som var der Nogen, hvem det udæskede. Kun den Blødagtige kunde ønske, at Job ikke var til, at han jo før jo hellere kunde slappe enhver Tanke, opgive enhver Bevægelse i den modbydeligste Afmagt, udslette sig selv i den elendigste og usleste Glemsel.
\end{quote}
\dcite{4T43:116  (SKS 5, 116)}
% 4T43:116  (SKS 5, 116)

\begin{quote}
»Det er saligere at give end at modtage« det Ord udtrykker Forskjellen, som den er i Jordlivet, ikke speilende sig i eller laanende sit Udtryk fra menneskelige Lidenskaber, at det er bedre, lykkeligere, herligere, ønskeligere at kunne give end at maatte modtage; men Ordet betegner Forskjellen, som denne er i det kjerlige Hjerte. Thi den, der siger, at det er bedre, eller lykkeligere, eller herligere eller ønskeligere, han siger jo ikke hvad Brug han vil gjøre af denne Magt, der forundtes ham, han vidner heller ei om, hvad han haver fornummet ved at bruge den rettelig. Den derimod, der veed, at det er saligere at give end at modtage, han haver fattet Livets Forskjellighed, men ogsaa for sin Deel forsonet den ved at give med Glæde og ved at have sin Salighed i at give. Dog bliver Forskjellen tilbage, hvis ikke Den, der maa modtage, fornemmer en lignende Salighed. Skulde han ikke det? Skulde Den, der blev forsøgt i en vanskeligere Strid, ikke have en lignende Løn, og det er dog i Sandhed større at modtage end at give, det vil sige, hvis begge Dele skee paa den rette Maade. Saa lader os da nu betragte det apostoliske Ord for at see, om dette ikke udtrykker Ligheden. Den, der giver, tilstaaer, at han er ringere end Gaven; thi det var jo denne Bekjendelse, som det apostoliske Ord med Glæde modtog af den Enkelte som frivillig eller afnødte ham som ufrivillig. Den, der modtager, han bekjender, at han er ringere end Gaven; thi det var denne Tilstaaelse, som Ordet gjorde til Betingelse for at opløfte ham, og denne Tilstaaelse lyder jo ogsaa ofte i Verden i al sin Ydmyghed. Men naar Den, der giver, er ringere end Gaven, og Den, der modtager, er ringere end Gaven, saa er jo Ligheden tilveiebragt, Ligheden nemlig i Ringhed, Lighed i Ringhed overfor Gaven, fordi Gaven er ovenfra, og derfor ikke egentlig tilhører eller ligemeget tilhører begge, det vil sige den tilhører Gud. Ufuldkommenheden laa da ikke i, at den Trængende modtog Gaven, men i at den Rige besad, hvilken Ufuldkommenhed han hævede derved, at han bortskjenkede Gaven; nu tilveiebragtes Ligheden, ikke derved, at den Trængende besad Gaven, men derved, at Gaven ikke mere tilhørte den Rige som Eiendom, thi han havde jo skjenket den Trængende den; og heller ikke tilhørte den Trængende som Eiendom, thi han havde jo modtaget den som Gave. Naar den Rige takker Gud for Gaven, og for at der forundtes ham Leilighed til skjønt at bortgive den, saa takker han jo for Gaven og for den Fattige; naar den Fattige takker Giveren for Gaven og Gud for Giveren, saa takker han jo ogsaa Gud for Gaven; altsaa er Ligheden tilstede i Taksigelsen til Gud, Ligheden overfor Gaven i Taksigelsen. Dette er uden Vanskelighed og let at fatte; forsaavidt Vanskeligheden stundom frister et Menneske, da er det fordi de jordiske Gaver, paa hvilke vi destoværre nærmest tænke, naar vi minde om Forskjellen mellem at kunne give og at maatte modtage, i sig ere saa ufuldkomne, og derfor have denne Ufuldkommenhedens Magt til at gjøre Forskjel. Jo fuldkomnere Gaven derimod er, desto tydeligere og uomtvisteligere viser ogsaa strax Ligheden sig. Men det er jo den gode og fuldkomne Gave, om hvilken vi sige, at den er ovenfra, men den eneste gode og fuldkomne Gave, et Menneske kan give, er Kjerlighed, og om den have jo alle Mennesker til alle Tider bekjendt, at den haver sit Hjem i Himlen og kommer ned herovenfra. Vil et Menneske ikke modtage Kjerligheden som en Gave, da see han til, hvad han faaer, men modtager han den som en Gave, da ere Giveren og Modtageren ikke til at adskille i Gaven, begge gjort væsentligen og fuldelig lige overfor Gaven, saaledes, at kun den jordiske Forstand i sin Ufuldkommenhed kan tvetyde, hvad der betyder eet og det Samme. Og jo mere et Menneske afvænner sin Sjel fra at forstaae det Ufuldkomne til at fatte det Fuldkomne, desto mere vil han \tlg{tilegne} sig den Forklaring af Livet, der trøster medens det er Dag, og som bliver hos ham, naar Natten kommer, naar han ligger glemt i sin Grav, og selv har glemt hvad Møl og Rust har fortæret, og menneskelig Kløgt udfundet, og dog skal have en Tanke, der kan udfylde ham den lange Tid, der ikke vil vide af den Forskjellighed, med hvilken han bekymredes, men kun om den Lighed, der er ovenfra, den Lighed i Kjerlighed, som bliver og som er det eneste Blivende, den Lighed, der ikke tillader noget Menneske at være det andets Skyldner, uden, som Paulus siger, i dette Ene at elske hverandre.
\end{quote}
\dcite{4T43:157  (SKS 5, 157)}
% 4T43:157  (SKS 5, 157)

\subsection{3T44 \normalfont(SKS 5)}

\begin{quote}
Der gives en Sandhed, hvis Storhed, hvis Ophøiethed man pleier at prise ved beundrende at udsige om den, at den er ligegyldig, lige gyldig, hvad enten Nogen antager den eller ikke; ligegyldig ved den Enkeltes særlige Tilstand, om han er ung eller gammel, glad eller bedrøvet; ligegyldig ved sit Forhold til ham, om den gavner ham eller skader ham, om den afholder ham fra Andet eller forhjælper ham til det; lige gyldig, om han bifalder den af sin ganske Sjel eller kold og følesløs bekjender den, om han lader sit Liv for den, eller benytter den i slet Vinding; ligegyldig ved, om han har udfundet den, eller han blot eftersiger det Lærte. Og kun Dens Forstaaelse var den sande, Dens Beundring den berettigede, der fattede, at denne Ligegyldighed var det Store, og lod sig danne i Overeensstemmelse med den til Ligegyldighed mod, hvad der angik ham selv eller noget Menneske som Menneske eller særligen som Menneske. Der gives en anden Art Sandhed, eller om dette er beskednere, en anden Art af Sandheder, hvilke man kunde kalde de bekymrede. Deres Liv er ikke i det Ophøiede, allerede af den Grund, at de ligesom beskæmmede ere sig selv bevidste ikke at passe ganske i Almindelighed ved alle Leiligheder, men kun særligen ved enkelte. De ere ikke ligegyldige ved den Enkeltes særlige Tilstand, om han er ung eller gammel, glad eller bedrøvet; thi dette afgjør for dem, om de skulle være Sandheder for ham. De slippe ei heller strax den Enkelte og opgive ham; men vedblive at bekymre sig om ham, indtil han selv løsriver sig ganske, og end ikke dette er dem ligegyldigt, skjøndt han dog ikke formaaer at gjøre dem tvivlsomme om sig selv. En saadan Sandhed er ikke ligegyldig ved, hvorledes den Enkelte annammer den, om han \tlg{tilegner} sig den af ganske Hjerte, eller den vorder ham et Mundsveir, thi netop denne Forskjellighed viser jo, at den er nidkjær paa sig selv; ikke ligegyldig ved, om Sandheden vorder ham til Velsignelse eller til Fortabelse, thi denne modsatte Afgjørelse vidner jo netop mod den lige Gyldighed; om han i Oprigtighed fortrøster sig til den, eller vil bedrage Andre, selv bedraget, thi denne hævnende Vrede viser jo netop, at den ikke er ligegyldig. En saadan bekymret Sandhed er ikke ligegyldig ved, hvo der haver fremsat den, tvertimod vedbliver han bestandigen at være tilstede i den, for igjen at bekymre sig om den Enkelte.
\end{quote}
\dcite{3T44:234  (SKS 5, 234)}
% 3T44:234  (SKS 5, 234)

\begin{quote}
Hvor ofte er ikke saaledes hiint gamle Ord gjentaget i Verden, anbragt paa rette Sted af Betragteren; men hvor sjeldent er det forstaaet itide, og naar da endeligen Forstaaelsen kom og det var for silde, hvorledes lød den? Hvor let er det ikke ogsaa for den Enkelte at unddrage sig Hentydningen, thi Ordet taler saa almindeligt! Men den almindelige Tale om almindelige Sandheder kan vel berige et Menneskes Hukommelse og udvikle hans Forstand, men er kun saare lidet til Gavn for ham, saalidet som det er til Gavn at have et Rustkammer fuldt af Vaaben, som dog ikke passe for En, naar man skal bruge dem. Og fremfor Alt er denne Almindelighed ikke til Opbyggelse; thi man opbygges aldrig i Almindelighed, ligesaalidet som et Huus opføres i Almindelighed. Først naar Ordet siges af den Rette, under de rette Forhold, paa den rette Maade, først da har Udsagnet gjort Alt hvad det formaaer for at veilede den Enkelte til oprigtigt at gjøre, hvad man ellers er snar nok til at gjøre: at henføre Alt til sig selv. Og om det end er forbudet i guddommelige og menneskelige Love, at begjære hvad Næstens er, Næstens Formaning er det aldrig forbudet at begjære, ei heller at benytte dens Veiledning. Saaledes er Alt her i Orden; thi hvad vi have sagt, gjelder i fuld Maade om hiint forelæste Ord af Johannes den Døber, og Ingen gjøre sig Samvittighed af at \tlg{tilegne} sig Ordet.
\end{quote}
\dcite{3T44:271  (SKS 5, 271)}
% 3T44:271  (SKS 5, 271)

\subsection{TTL \normalfont(SKS 5)}

\begin{quote}
Skjøndt denne lille Bog (Leilighedstaler, som den kunde kaldes, hvorvel den ikke har Leiligheden, der gjør Taleren og gjør ham myndig, eller Leiligheden, der gjør Læseren og gjør ham lærende) er uden al Opfordring og saaledes i sin Mangelfuldhed uden al Undskyldning, er uden al Bistand af Omstændigheder og saaledes i sin Udførlighed hjælpeløs, er den dog ikke uden Haab og fremfor Alt ikke uden Frimodighed. Den søger hiin Enkelte, hvem jeg med Glæde og Taknemlighed kalder min Læser, eller den søger ham end ikke. Uvidende om Tiden og Timen venter den i Stilhed, at hiin rette Læser skal komme som Brudgommen og bringe Leiligheden med sig. Hver gjøre Sit, Læseren altsaa meest. Betydningen ligger i \tlg{Tilegnelsen}. Deraf Bogens glade Hengivelse. Her er intet verdsligt Mit og Dit, der adskiller og forbyder at \tlg{tilegne} sig hvad Næstens er. Thi Beundring er dog lidt Misundelse og saaledes en Misforstaaelse, og Dadel i al sin Berettigelse dog lidt Modstand og saaledes en Misforstaaelse, og Gjenkjendelse i Speilet kun et flygtigt Bekjendtskab og saaledes en Misforstaaelse, men at see rigtigt til og ikke at ville glemme hvad Speilets Afmagt ikke formaaer at bevirke, det er \tlg{Tilegnelsen}, og \tlg{Tilegnelsen} er Læserens endnu større, er hans seierrige Hengivelse.
\end{quote}
\dcite{TTL:389.2  (SKS 5, 389)}
% TTL:389.2  (SKS 5, 389)

\begin{quote}
Der var engang i Verden, da Mennesket træt af Forundring, træt af Skjebnen vendte sig bort fra det Udvortes, og fandt, at der var ingen Forundringens Gjenstand, at det Ubekjendte var et Intet og Forundringen et Bedrag. Og hvad der engang var Livets Indhold, det kommer igjen i Slægtens Gjentagelse. Om Nogen vil tykkes sig viis ved at sige, at der ere tilbagelagte Skikkelser, som for Aartusinder siden ere færdige: i Livet er det ikke saaledes. Og Du mener dog vel heller ikke, min Tilhører, at jeg vilde spilde Dig Din Tid ved at fortælle store Begivenheder og nævne snurrige Navne, og blive aandløst vigtig i Betragtning af hele Menneskeheden! Ak nei, hvis det er saa, at Den er bedragen, som kun faaer Lidet at vide, mon da ikke ogsaa Den, som fik saa Meget at vide, at han slet Intet \tlg{tilegnede} sig! Langsomt gaaer Mennesket frem, selv den herligste Viden er dog kun en Forudsætning. Vil man forøge Forudsætningerne mere og mere, da er man jo som den Gjerrige, der dynger Penge sammen, som han ingen Brug har for. Selv hvad der er værd at anslaae høit: en lykkelig Opdragelse, selv den er jo kun en Forudsætning, og der gaaer Tid efter Tid, inden den \tlg{tilegnes}, og et heelt Liv er ikke for meget, hvis man vil \tlg{tilegne} sig den. O, dersom Den var bedragen, hvis Opdragelse forsømtes, mon da ikke ogsaa Den, som blev uvidende om, at den var en Forudsætning, et betroet Gods, en hellig Arv, der skulde erhverves, og som uden videre tog den hen, og tyktes sig at være det, han havde Navn af. Har den Bedre stundom sukket, fordi det Søgte var saa fjernt, saa har Du vel min Tilhører, fattet, at der ogsaa er en anden Vanskelighed, at der er et Videns Blendværk, som bedaarer Sjelen, at der er en Tryghed, hvori man er vidende – og dog bedragen, at der er en Fjernhed fra al Afgjørelse, hvor man, uden at drømme derom, fortabes. Lad Forfærdelsen tage sit Bytte, o, denne Tryghed er et forfærdeligere Uhyre! Lad Savnets Nød hungre ud, er det bedre at omkomme af Overflod? Rystende var det, da Forundringen slap Mennesket, og han fortvivlede over sig selv, men lige saa rystende er det, at man kan være vidende om dette, vide langt mere, og end ikke have oplevet hiint, og meest rystende, at man kan vide Alt og ikke have begyndt paa det Mindste. Og er det saaledes, o, lad mig da begynde forfra; vend tilbage, Du Ungdom med Din Ønsken og Din elskelige Forundring, vend tilbage, Du Ungdommens vilde Stræben med Din Dumdristighed og Din Gysen for hiin Ubekjendte, grib mig, Du Fortvivlelse, der bryder med Forundringen og Ungdommens Forundring, men hurtigt, hurtigt, hvis det er muligt, hvis jeg har spildt min bedste Tid uden at opleve Noget, lær mig dog idetmindste ikke at blive ligegyldig derved, at søge Trøst med Andre i fælles Tab, saa er vel Forfærdelsen over Tabet en Begynden paa min Helbredelse; hvor sildig den end er, det er dog bedre end at leve hen som en Løgner, bedragen ikke af hvad der kunde synes skikket til at bedrage, ak, og derfor rædsomt bedragen – bedragen af megen Viden!
\end{quote}
\dcite{TTL:402  (SKS 5, 402)}
% TTL:402  (SKS 5, 402)

\begin{quote}
Her ender denne Tale – i Syndens Bekjendelse. Men kan det ogsaa være en Ende; skal Glæden nu da ikke seire; skal Synden kun gaae med Sorgen; skal Sjelen da sidde beklemt, men Frydens Harpe ikke stemmes? Du pleier maaskee at faae mere at vide, Du veed vel selv meget mere, saa søg da Feilen hos Talen og den Talende. Er Du virkeligen videre, saa lad Dig ikke sinke; men hvis ikke, o, saa betænk, at man er forfærdelig bedragen, naar man er bedragen af megen Viden. Lad os tænke en Styrmand og antage, at han med Udmærkelse havde bestaaet alle Lære-Prøver, men han havde dog endnu ikke været ude – tænk Dig ham i en Storm: han veed Alt hvad han har at gjøre, men han kjendte ikke den Forfærdelse, som griber den Søfarende, naar Stjernerne blive borte i Nattens Mulm; han kjendte den ikke den Afmagt, med hvilken den Styrende seer, at Roret i hans Haand er et Legetøi for Havet; han vidste ikke, hvorledes Blodet stormer til Hovedet, naar man i et saadant Øieblik skal anstille Beregninger: kort, han havde ingen Forestilling om den Forandring, der foregaaer med den Vidende, naar han skal bruge sin Viden. Hvad Magsveir er for Sømanden, er for det enkelte Menneske det at leve hen i jævn Fart med Andre og med Slægten, men Afgjørelsen, Besindelsens farefulde Øieblik, naar han skal tage sig ud af Omgivelsen og blive ene for Gud og blive en Synder, dette er en Stilhed der forandrer det Sædvanlige ligesom Stormen. Han veed det Alt, veed hvad der skal hænde ham, men han vidste ikke, hvilken Angest der griber, idet han føler sig forladt af det Mangfoldige, hvori han havde sin Sjel; han vidste ikke hvorledes Hjertet banker, naar Hjælpen fra Andre og Veiledningen fra Andre og Maalestokken fra Andre og Adspredelsen ved Andre forsvinder i Stilheden; han vidste ikke, hvorledes det gyser, naar det er for silde at raabe om Menneskenes Hjælp, da Ingen kan høre ham: kort, han havde ingen Forestilling om, hvorledes den Vidende forandres, naar han skal \tlg{tilegne} sig sin Viden. Skulde dette maaskee ogsaa være Dit Tilfælde, min Tilhører? Jeg dømmer jo ikke, jeg spørger Dig kun. Ak, medens De blive flere og flere, som vide saa saare meget, saa blive de heelbefarne Mænd sjeldnere og sjeldnere! Men en Saadan var det jo Du engang ønskede at være. Du har vel ikke glemt, hvad vi talte om Oprigtighed mod sig selv: at man tydeligt husker, hvad man engang vilde være; og selv er Du jo betænkt paa at ville være oprigtig for Gud i Syndernes Bekjendelse. Hvad var det da Du vilde? Du vilde stræbe efter det Høieste, at fatte Sandheden og at være i den; Du vilde ikke spare Tid og Flid; Du vilde forsage Alt, deriblandt dog vel ogsaa ethvert Bedrag. Om Du end ikke greb det Høieste, Du vilde dog sikkre Dig, at Du tydeligen vidste med Dig selv, hvad Du hidtil havde forstaaet ved at naae det. Om dette var nok saa lidet, Du vilde dog hellere være tro over Lidet end utro over Meget; om det var een eneste Tanke, og Du blev den Fattige blandt de Rige, der vide Alt, Du vilde dog hellere være tro som Guld – og det kan jo Enhver, naar han vil det, thi Guldet, det tilhører den Rige, men tro som Guld kan jo den Fattige ogsaa være. Og Den, som var tro over Lidet, tro paa Nødens Dag, naar Regnskabet gjøres op, tro i at forstaae sin Gjeld, tro i Stilheden, hvor ingen Løn vinker men Skylden bliver tydelig, tro i Oprigtigheden, der vedgaaer Alt, endog at denne Oprigtighed dog er mangelfuld, trofast i Angerens Kjærlighed, hiin ydmyge Kjærlighed, hvis Fordring er Selvanklagen: han skal vel ogsaa blive sat over Mere.
\end{quote}
\dcite{TTL:415  (SKS 5, 415)}
% TTL:415  (SKS 5, 415)

\begin{quote}
Var det ikke saaledes Du vilde? Thi, ikke sandt, derom ere vi jo enige: at i Forhold til det Væsentlige er det at kunne væsentligen at kunne gjøre det. Barnet mener det anderledes, og naar den Lille lærer paa, hvad der er ham foresat, og nu maaskee siger til den voxne Søster, om hun vil overhøre ham, men hun har Andet at tage vare og svarer, Nei, min kjære Dreng, nu har jeg ikke Tid, men læs det Foresatte fem Gange eller læs det ti Gange, og sov saa paa det, saa kan Du det udmærket imorgen: saa troer Barnet det, gjør som det bliver ham sagt, og kan det udmærket den næste Dag. Men den Modnere lærer paa en anden Maade. Og om Nogen end vilde tage den hellige Skrift og lære udenad, da kunde det vel være smukt, forsaavidt der var noget Barnligt i hans Adfærd, men væsentligen lærer den Ældre kun ved at \tlg{tilegne} sig, og væsentligen \tlg{tilegner} han sig kun det Væsentlige ved at gjøre det. O, midt i al Nød skjønne Glæde over Livet og over Slægten og over selv at være Menneske; o, under Stilheden skjønne Samdrægtighed med Enhver; o, i Eensomheden skjønne Fællesskab med Alle! Thi det er ikke saaledes, at det ene Menneske ikke har det samme Væsentlige til Opgave, som det andet Menneske, saa lidet som det ene Menneskes Udvortes er væsentligen forskjelligt fra det andets, men det er saaledes, at Hver forstaaer det lidt forskjelligt og paa sin Viis. Og det er ikke saaledes som i Forvirringen, at der er forskjellige Veie og forskjellige Sandheder og nye Sandheder, men det er saaledes, at der er mange Veie, som føre til den ene Sandhed, og Hver gaaer sin. Deraf Eiendommeligheden, naar det Væsentlige bliver den Enkeltes Eiendom, og denne Eiendommelighed betinges ved at gjøre det og opdages derved. Skulde denne Tale være splidagtig? Langtfra den være at nævne hiin Eiendommelighed, om hvilken der strides i Verden, ligesom om andre Lykkens Gaver; nei, Enhver der eier noget Væsentligt ved at have gjort det, han har Eiendom og Eiendommelighed. Og saaledes for at minde om Talens Gjenstand, at forstaae hiin Stilhed er at kunne blive stille. Og hvor skal man blive det? Ja, der er Stedet dertil, men dog ikke udvortes og ligefrem, thi bringer man ikke Stilheden med, saa gavner Stedet Intet. Altsaa der er i en vis Forstand intet Sted; o, dette: i en vis Forstand er det ikke allerede uroligende! Og naar behøver man denne Stilhed meest? Naar man er stærkest bevæget. Er denne Tanke ikke istand til at jage Stilheden bort? Hvor flyer man saa hen for at undflye sig selv? Ja, vil man flye, saa undflyer man netop Stilheden. Saa er der slet Intet at gjøre? Ja, vil man slet Intet gjøre, saa undflyer man atter Stilheden i den aandelige Døds Stilhed. O, mon det er saa let at kunne blive stille! Nu frister en Tryghed, fordi der jo er Tid nok, nu en Utaalmodighed, fordi det er for silde, nu et vinkende Haab, nu en dvælende Erindring, nu en stormende Beslutning, nu en Gjenlyd af Verdens Dom der spottende indhenter Dig, som gik Du ad denne Stilhedens Vei til Bedragets Ørken, hvor den Eensomme omkommer, nu en Gjenlyd fra det Selviske i Dig, der forstyrrer med Selvbeundring, nu en Sammenligning, der adspreder, nu et Overslag, der adspreder, nu lidt Glemsel ved Hjælp af Tankeløshed, nu lidt Forskud ved Hjælp af Selvtillid, nu en eventyrlig Forestilling om Guds Uendelighed, nu Nedstemthed ved at skulle betroe den Alvidende hvad han dog veed, nu et letsindigt Spring, der Intet gavner, nu et tungsindigt Suk der nærer Tungsindet, nu lidt Veemod som bedøver, nu en Klarhed som overrasker, nu en Stilhed ved Planer og Tanker og Forsætsdrømme og Indbildningsbilleder istedenfor ved Skyld og Regnskab og Forsættets Pagt med den tydelige Skyld og den alvidende Gud. O, mon det er saa let at blive stille! At have været ganske nær derved, og dog at have grebet en Skuffelse, og skulle begynde igjen og altsaa med mere Uro! At have fundet Trøst hos en Anden og nu opdage, at det var et Selvbedrag, en forfalsket Stilhed, og altsaa at skulle begynde med mere Uro! At være bleven forstyrret af Verden, af en Fjende, af en Ven, af en falsk Lærer, af en Hykler, en Spotter, og nu opdage at det er et Selvbedrag at ville skyde Skylden paa en Anden, og altsaa skulle begynde med mere Uro! At have stridt, at ville af yderste Evne, og da opdage, at man Intet formaaer, at man dog ikke kan give sig selv denne Stilhed, fordi den tilhører Gud! Vil Nogen sige, at dette er det rette Udtryk, at man ikke kan det, da betænke man vel, om det er Dorskheden, som her taler. Vel er det nemlig sandt, ja selv en Apostel vidner jo derom, men mon dette Vidnesbyrd var et Indfald, en almindelig Bemærkning i Hastværk, eller var det ikke saa svært at forstaae denne menneskelige Intethed og at have sit Bevidsthedsliv deri, at selv han, den Myndige, den for evigt Besluttede ikke var ene derom, men behøvede en Medhjælper, en Satans Engel nemlig, der ved daglige Erfaringer og ved daglig Pine hjalp ham ud af Sandsebedrag, ud af at have sin Viisdom i det udenad Lærte, sin Fred i almindelige Forsikkringer, sin Tillid til Gud i et Mundheld. Eller havde Nogen lært Apostelen dette, saa han kunde tale bagefter? Man har vel hørt før i Verden, at den Vise havde en Engel, der veiledede eller advarede ham; havde Paulus talt derom, saa kunde det være udenadlært, men at den Vise behøver en Satans Engel til daglig Brug, det har visseligen kostet lang Tid at lære.
\end{quote}
\dcite{TTL:417  (SKS 5, 417)}
% TTL:417  (SKS 5, 417)

\begin{quote}
Saa har da Talen, idet den søgte at tydeliggjøre Ægteskabets hellige Beslutning, vel mindet Dig, min Tilhører, om hvad Du selv oftere har betænkt, thi Talen er langtfra at være belærende. Du er da enig med den, og dog siger Du maaskee: Talen er rigtig, men der fordres megen Alvor for at en saadan Tale skal frembringe det rette Indtryk, at den ikke føder Utaalmodighed eller forstyrrer i Forstemthed. Og heri har Du vist ganske Ret, at der fordres megen Alvor; sandeligen, det at være en god Læser eller en god Tilhører er lige saa stort som at være en god Taler, og det er meget stort, naar, som her, Talen er i Ufuldkommenhed og uden Myndighed. Var det ikke saaledes Din Mening, thi Du vilde dog ikke med hiin Indvending skyde Skylden paa den Talende, som havde Du vundet Noget ved at anklage ham? Lader os betragte dette lidt nærmere. Som der til at fatte en Beslutning fordredes en virkelig Forestilling om Livet, saa fordredes der ogsaa, hvad der blev sagt med det Samme, en virkelig Forestilling om sig selv. Der var maaskee Den, som sendte de speidende Tanker ud for at faae mangfoldige Indtryk af Livet, men som ikke kunde tage sig selv tilbage, som gav sig hen – ak, og tabte sig selv. Men Den, som ved Vielsen knytter et andet Menneskes Liv til sit, ved Vielsen indgaaer en Forpligtelse, som ingen Tid skal løse og som hver Dag skal fordre opfyldt, af ham fordres der jo en Beslutning, og i denne Beslutning altsaa en virkelig Forestilling om sig selv. Og denne virkelige Forestilling om sig selv, og denne Forestillings Inderlighed, det er Alvoren. Nu er det vel saaledes, at som hiin Elskov, hvilken Digterne besynge, er en Længsel i ethvert Menneskes Sjel, saa er der ogsaa i Enhver en Længsel, et Ønske der begjerer hvad man maatte kalde en Veileder og Lærer i Livet, den forsøgte Mand, paa hvem man kan forlade sig, den Vise, der veed at raade, den Ædle, der opmuntrer ved sit eget Exempel, den Begavede, der har Veltalenhedens Magt og Overbeviisningens Fynd, den Alvorlige, der betrygger \tlg{Tilegnelsen}. Som Barn gaaer det saa let, man er fritagen for Valgets Besværlighed; ja selv om en Fader ikke var som han burde, den Ærbødighed Barnet har, den ubetingede Lydighed hjælper det stundom til at lære det Gode selv af en saadan Fader. Men saa kommer Ungdommens og Frihedens Tid og den Tid, da han med den Elskede søger en saadan Veileder. Da gjelder det, at Friheden og Valget ikke bliver en Snare. Hiin Længselens søgte Veileder er den Sjeldne. Stundom findes han end ikke i hver Slægt, og selv om Du er samtidig med en saadan Høiærværdig, hvem Du ganske tør hengive Dig til, han er maaskee ikke paa samme Sted som Du, eller han var det, men maa forlade det, eller Du maa forlade det – og saa, ja saa maa Du nøies med Mindre, det er, Du maa see til at hjælpe Dig selv. I Livet er der Forvirring nok; det Forskjelligste forkyndes og anprises og foragtes og gjentages; de forskjelligste Forbilleder vise sig og skuffe og vise sig atter; de forskjelligste Veiledninger tilbydes, og altid er der Reiseselskab; Trøstegrunde og Udflugter og Tilraab og Advarsler og Seierssange og Klageskrig høres forvirret imellem hinanden. Ak, Elskov og Ægteskabet er jo Noget, hvori Enhver forsøges og hvorom altsaa Enhver har en Mening og jo ogsaa med Sandhed kan have det, hvis han i Alvor vil. Ethvert Menneske, ogsaa den Ugifte, skal jo have et blivende Sted, og dog er der maaskee mangt et Ægteskab som ikke har det, men omtumles for enhver Luftning. Nu mener den Unge, vildledet af en tilfældig Erfaring, at naar de ydre Omstændigheder i Velstand og lykkelige Livsforhold begunstige Elskoven, da er den betrygget, og han betænker maaskee ikke, hvorledes det Raaderum som Sjelstilstandene derved faae, kan føde sine Vanskeligheder. Nu har Een en overspændt Forestilling om den forelskede Sjelesorgs vidtløftige Gjerning, og kan ikke nedstemme sig til Næringssorgens simple Arbeide; nu giver en Anden sig for meget hen i Følelsers Overskummen, og en Afsmag følger, nu har den ene Part dog nogen Besindighed og vil bruge den, men den anden Part misforstaaer det og troer det er Kulde og Ligegyldighed; nu vil den ene Part holde til Raade og spare, og den anden Part forstaaer det ikke og troer at det er Mangel paa Sands for noget Høiere. Nu mistrøstes Een, fordi Gjentagelsen uden om ham gjør ham kjed ved sit Eget; nu gjør en Andens første Lykke ham utaalmodig; nu sammenligner han, nu minder han, nu taber han – ja hvo skulde vel blive færdig med at nævne dette; det kan ingen Tale, og dette var jo ogsaa ligegyldigt, men intet Menneske kan det, det er det Forfærdelige; kun een Magt formaaer det, det er Beslutningen, som passer paa itide. Hvor læres da Alvor? I Livet. Ganske vist, og Ægteskabets Gud velbehagelige Stand har jo en sjelden Beleilighed. Saa læres Alvoren – hvis man tager Beslutningen med og i den en virkelig Forestilling om sig selv. Beslutningen selv er Alvoren. Og for at man skal lære Alvor af det man kalder Livets Alvor forudsættes allerede Alvor. Thi Livets Alvor er ikke som en Læremester i Forhold til den Lærende, men i en vis Forstand som en ligegyldig Magt i Forhold til Den, der selv maa være noget af en Læremester i Forhold til sig selv som den Lærende. Ellers kan man endog af Livets Alvor lære Ligegyldighed for Alt. Man ønsker vel en Veiledning, og dog gjelder det ogsaa i Forhold til den, at man selv maa have Alvoren for at hjælpes derved. Eller er det ikke seet, at man selv i Forhold til hiin sjeldne Veileder, naar han stod iblandt os, hittede paa mange Ting for at svække sig Indtrykket af ham, som tabte han derved ikke En selv, som var det Aarenes Viisdom at blive mere og mere kræsen, mere og mere duelig til at vrage, istedenfor at blive mere og mere skjønsom. Og naar der nu ingen Saadan er, hvad saa? Ja, Verden lader det aldrig mangle paa Veiledere. See, nu vil En veilede Alle og kan ikke hjælpe sig selv; nu udraabes En som Viis og Beundringen kjender ham paa, at han ikke engang kan forstaae, hvad den Eenfoldige forstaaer; nu har En Overtalelsens Magt, og leder vild, har Usandhedens kraftige Gjerninger; nu skal det være forældet, hvad man lærte i Barndommen, og man maa lære om igjen. Nu vil En rive Ægtemanden fra Hustruens Side og gjøre ham vigtig ved Deeltagelse i store Foretagender og lære ham at tænke ringe om Ægteskabets hellige Kald; nu frister En Hustruen og lærer hende at sukke under Ægteskabets Aag; nu vil man foregøgle Mand og Hustru Fællesskab i Bedrifter, der gjør det ægteskabelige Forhold ligegyldigt; nu vil man lære Ægtefolkene Nydelse, vil tage Børnene og dermed Sorgerne fra dem, for at Forældrene kunne leve for høiere Beskæftigelser. Saa spændes Forventningen paa noget Overordentligt, en ny Tingenes Orden der skal komme, og vi faae Alle, baade Gifte og Ugifte, fri som Skolebørnene, fordi Skolemesteren skal flytte, fri indtil han kommer i Orden – men vi ere jo dog ikke Skolebørn længere, og Enhver skal gjøre Gud Regnskab for sig, og Ægteskabets hellige Forpligtelse skal jo give hver Dag sin Gjerning og sit Ansvar. Hvor finder man saa Veiledningen, hvis man ikke selv forarbeider sin egen Sjels Salighed med Frygt og Bæven, thi saa bliver man nok alvorlig. Ellers skal det vel vise sig, at den ene Veileder kan man ikke følge, thi om han end mener det saa godt, er han dog svag; og den anden kan man ikke følge, thi vel har han Kraft og Fynd i hvad han siger, men han skal ikke mene det, saa siger man; eller den ene er for gammel til at tilfredsstille Tiden, og den anden er for ung. Ja hvor skulde en Tale vel blive færdig, hvis den vilde skildre denne Livets Forvirring, men hvad tykkes Dig, min Tilhører, om et saadant Menneske, der vel havde Alvorens Skin i at vrage Alt, men ikke Spor af dens Kraft i at besidde det Mindste; hvad tykkes Dig om en Saadan, hvis han var Ægtemand, og altsaa var bleven gift uden Beslutning, og altsaa levede hen ubekymret om den helligste Forpligtelse, det vil sige, uden i Alvor at bekymre sig om den. Havde ikke Skilsmissen ogsaa mærket et saadant Ægteskab, hvor Ægtefolkene vel tilhørte hinanden, men ikke i Alvor!
\end{quote}
\dcite{TTL:433  (SKS 5, 433)}
% TTL:433  (SKS 5, 433)

\begin{quote}
Den umyndige Tale kan ikke saaledes gjøre Alvor deraf, ingen Død venter paa den, for at Alt kan være forbi. Men derfor kan Du jo gjerne agte paa Talen, min Tilhører. Thi Døden selv har jo sin Alvor; det Alvorlige ligger ikke i Begivenheden, ikke i det Udvortes: at der nu igjen er en Mand død, saa lidet som Alvorens Forskjel ligger i, at der var mange Karreter; ja saa lidet som hiin formildede Stemning, der kun vil tale godt om de Døde, er Alvor eller i fjerneste Maade kunde tilfredsstille Den, der alvorligen betænkte sin egen Død. Døden kan netop lære, at Alvor ligger i det Indvortes, i Tanken, lære, at det kun er et Sandsebedrag, naar der letsindigt eller tungsindigt sees paa det Udvortes, eller naar Betragteren dybsindigt over Dødens Tanke glemmer at tænke og betænke sin egen Død. Vil man nævne ret en Gjenstand for Alvor, da nævner man Døden, og »Dødens alvorlige Tanke«; og dog er det jo som laae der en Spøg til Grund for Døden, og denne Spøg, forskjelliggjort i Stemnings og Udtryks Forskjellighed, er det Væsentlige i enhver Betragtning af Døden, hvor Betragteren ikke selv bliver under fire Øine med Døden og tænker sig selv sammen med Døden. En Hedning har allerede sagt, at Døden skal man ikke frygte, »thi naar den er, er jeg ikke, og naar jeg er, er den ikke.« Dette er Spøgen, hvormed den listige Betragter stiller sig selv udenfor; men selv om Betragtningen brugte Rædselens Billeder for at skildre Døden, forfærdede en syg Indbildning, det er dog kun Spøg, hvis han blot tænker Døden, ikke sig selv i Døden, hvis han tænker den som Slægtens Vilkaar, men ikke som sit. Spøgen er, at hiin ubøielige Magt ligesom ikke kan faae Ram paa sit Bytte, at der er en Modsigelse, at Døden ligesom narrer sig selv. Thi Sorgen, hvis Du dermed vil sammenligne Døden, og hvis Du vil kalde den en Bueskytte som Døden er det: Sorgen rammer ikke feil, thi den rammer den Levende, og naar den har rammet ham, saa først begynder Sorgen: men naar Dødens Piil har rammet, saa er det jo forbi. Og Sygdommen, hvis Du dermed vil sammenligne Døden, og hvis Du vil kalde den en Snare, som jo Døden er den Snare, hvori Livet fanges: Sygdommen fanger virkeligt, og naar den har fanget den Karske, saa begynder Sygdommen: men naar Døden trækker Snaren sammen, saa har den jo Intet fanget, thi saa er det forbi. Men heri netop ligger Alvoren, og herved netop er Dødens Alvor forskjellig fra Livets, der saa let lader En bedrage sig selv. Thi naar Nogen gaaer sin Gang bøiet i Modgang, i Lidelser, i Sygdom, i Miskjendelse, i trange Kaar, i kummerlige Udsigter, da slutter han feil, om han ligefrem deraf slutter, at han er alvorlig; thi Alvor er ikke den ligefremme Gjengivelse, men den forædlede, det er, atter her er det det Indvortes og Tanken og \tlg{Tilegnelsen} og Forædlingen, der er Alvoren. Eller naar man er travl i den vidtløftige Bedrift, maaskee har mange Stridsmænd at befale over, maaskee skriver mange Bøger, maaskee er i de høie Stillinger, naar man maaskee har mange Børn, eller ofte maa være med i Livsfare, eller har den alvorlige Gjerning, at klæde Liig, saa slutter man feil, hvis man deraf uden videre slutter, at man er alvorlig, thi Alvoren er i Indtrykket, Alvoren er det indvortes Menneskes, ikke Bestillingens. Døden derimod er ikke i den Forstand noget Virkeligt, og naar man først er død, er det bag efter med at blive alvorlig, og naar man fandt en brad Død, hvilket en alvorligere Tid betragtede som den største Ulykke, hvorfor den ogsaa nævnes i den gamle Bøn, hvilket en nyere Tid anseer for den største Lykke, saa var man jo hjulpen. Livets Alvor er alvorlig, og dog er der ingen Alvor uden Bevidsthedens Forædling af det Udvortes, heri ligger Skuffelsens Mulighed; Dødens Alvor er uden Bedrag, thi det er ikke Døden der er alvorlig, men Tanken om Døden.
\end{quote}
\dcite{TTL:445  (SKS 5, 445)}
% TTL:445  (SKS 5, 445)

\subsection{SLV \normalfont(SKS 6)}

\begin{quote}
Kun om Ægteskabet ønsker jeg at skrive; at overbevise en Enkelt, er mit Haab; at fjerne dem, som tale imod, er min Agt. Saaledes er Ægteskabet mig min ene Stræng, men den er saa sammensat, at jeg, uden just at fortrøste mig til Virtuositet, hvad der jo ellers fordres af Den, der kun har een Stræng, tør vove at lade mig høre, ikke just som en Kunstner for et talrigt Publikum, men hellere som en omvankende Musikant, der staaer udenfor det enkelte Huses Dør, og ikke kalder Nogen bort fra sin Gjerning, om der dog ogsaa i hans Musik, naar den lyder med i Gjerningen, er Tækkelighed. Jeg mener nemlig ingenlunde, at hvad jeg kan have at sige, skulde være uskjønt. Adskilligt skylder jeg min Kone, om jeg end ikke taler ganske saaledes med hende, som jeg her skriver, men hvad der kommer fra hende, har altid en vis Ynde, hvad der er Qvindens Medgift. Ofte har jeg undret mig derover. Som Een, der kun skriver en maadelig Haand, naar han seer sit eget Manuscript udført af en kaligraphisk Kunstner, maa blive forundret, som Den, der sendte et gnidret Ark i Bogtrykkeriet, naar han modtager en nydelig Reentryk, neppe tør vedkjende sig dette som sit Eget, saaledes er det oftere gaaet mig i mit huslige Liv. Hvad der dunkelt bevæger sig i hende, det udtrykker jeg saa godt jeg kan, da forundres hun, at det netop var, hvad hun vilde sige, jeg udtaler det altsaa saagodt jeg kan, da \tlg{tilegner} hun sig det, men nu kommer Touren til mig, naar jeg med Forundring seer, at mine Tanker, mine Ord have faaet en Beaandelse, en Inderlighed, en Ynde, saa jeg med Rette kan sige, at det ikke er mine Tanker. Ulykken er, at denne Ordenes og Tankernes elskelige Pynt meer eller mindre næsten ganske forsvinder, naar jeg vil gjentage dem, og lader sig ikke udtrykke, saa lidet som jeg her paa Papiret kan beskrive hendes Stemme. Imidlertid er hun til en vis Grad Med-Forfatter, og et saadant literairt Firma forekommer mig ikke uskjønt, naar man kun vil skrive om Ægteskabet. Hun billiger, det veed jeg, at jeg gjør Brug af hvad jeg skylder hende, hun tilgiver mig, det veed jeg, at jeg benytter Leiligheden til at sige Eet og Andet om hende, som jeg ellers ikke kan faae sagt uden i Eensomhed, thi hvor meget hun er for mig, kan jeg jo ikke sige ligefrem til hende selv, at min Lovpriisning ikke skulde blive besværende og maaskee næsten forstyrre vor gode Forstaaelse. Som anonym, og som den, der med al Omhyggelighed vil bevare sin Anonymitet, har jeg jo sikkret mig mod, hvad jeg forøvrigt haaber, at Delicatesse vil forbyde Enhver: at mit huslige Liv skulde blive Gjenstand for Nogens Nysgjerrighed.
\end{quote}
\dcite{SLV:94  (SKS 6, 94)}
% SLV:94  (SKS 6, 94)

\begin{quote}
Mig selv bliver der da slet ikke Tid til at tænke paa, og dog er min inderste Existents en saadan, at den kan give nok at betænke. Jeg er egentlig ingen religieus Individualitet, jeg er kun en ordentlig og fuldstændig construeret Mulighed til en saadan. Med Sværdet over Hovedet, i Livsfare opdager jeg de religieuse Kriser med en saadan Primitivitet, som havde jeg slet ikke kjendt dem iforveien, med en saadan Primitivitet, at, hvis de ikke vare opdagede, maatte jeg opdage dem. Dette behøves nu ikke, jeg kan i den Henseende ydmyge og tugte mig selv, som jeg engang trøstede en noget indskrænket Mand, hvem en Anden spottende sagde om: han har ikke opfundet Krudtet, jeg svarede, det behøves jo heller ikke, da det er opfundet. Men eet er at lære af Lærebogen og udenads Lectier, eet at kunne læse for Præsten ja selv for Biskoppen, naar han visiterer, eet at kunne præke som en Præst – et Andet er Oprindeligheden i \tlg{Tilegnelsen}. Godt er det, at jeg ikke skal lære Andre. Jeg betaler med Glæde baade Tavle-Penge og Præste-Penge, Held Den, der er saa sikker paa sig selv, at han tør tage Penge for Underviisningen.
\end{quote}
\dcite{SLV:241  (SKS 6, 241)}
% SLV:241  (SKS 6, 241)

\begin{quote}
Som Mulighed er jeg god nok, men idet jeg i Katastrophen vil \tlg{tilegne} mig de religieuse Forbilleder, møder en philosophisk Tvivl, som jeg ikke vil udtale som saadan ikke til et eneste Menneske. Det, det kommer an paa, er \tlg{Tilegnelses} Momentet. Disponeret som jeg er i den religieuse Katastrophe, griber jeg efter Paradigmet, men see! Paradigmet kan jeg slet ikke forstaae, om jeg end venererer det med en Barndommens Pietet, der ikke vil slippe det. Det ene Paradigma beraaber sig paa Syner, det andet paa Aabenbarelser, det tredie paa Drømme. At tale derom, at gjennemgløde Foredraget med Phantasi og dog beholde Forudsætningen, Forudsætningen, der netop angaaer \tlg{Tilegnelses}-Momentet for den Senere, er nemt nok, men at forstaae det!
\end{quote}
\dcite{SLV:241.2  (SKS 6, 241)}
% SLV:241.2  (SKS 6, 241)

\begin{quote}
Hvad nytter saa Klogskab? Men jeg er heller ikke klog. Er jeg i en vis Forstand den Klogeste iblandt Mine, saa er jeg ogsaa i en anden Forstand den Dummeste maaskee af Alle. I Alt, hvad jeg har hørt og læst, har Intet saaledes truffet og slaaet mig som et Ord, der siges om Periander. Om ham siges der, at han talede som en Viis og handlede som en Afsindig. At Ordet passer netop paa mig bevises deraf, at jeg annammer det med den lidenskabeligste Sympathi, og at det dog ikke har den mindste Indflydelse paa at forandre mig. Denne Maade at \tlg{tilegne} mig det paa, er ganske à la Periander. Inden min Forudsætning er jeg klog, men min Handlens Forudsætning er saa ideel, at denne Forudsætning gjør al min Klogskab til Daarskab. Naar jeg kunde lære at slaae af i min Forudsætning, saa vilde min Klogskab tage sig ud. Kunde jeg saaledes handle klogt, havde jeg været gift for længe siden. At faae sit Ønske opfyldt, at faae Tak til som for en Velgjerning, dernæst at indrette sig som man væsentligen havde sin Frihed, det havde været klogt, og jeg havde været et agtværdigt Menneske, der ikke bryder givent Ord, en seeværdig Ægtemand, der er sin Kone tro, og ærer en Pige; thi min ideelle Forudsætning ærer hende formodentligen ikke, og det at jeg hellere vilde sætte Alt til og Himmel og Jord i Bevægelse, end at skulke mig ind i den Gud velbehagelige Stand og skulke mig gjennem Livet, beviser formodentligen, at jeg ingen Ære har.
\end{quote}
\dcite{SLV:288.2  (SKS 6, 288)}
% SLV:288.2  (SKS 6, 288)

\begin{quote}
Hvad skeer? Store Gud, hun har været paa mit Værelse, medens jeg var ude. Jeg finder en Billet, affattet med en fortvivlet Lidenskab, hun kan ikke leve uden mig, det bliver hendes Død, hvis jeg forlader hende, hun besværger mig for Guds, for min Saligheds Skyld, ved ethvert Minde, der binder mig, ved det hellige Navn, som jeg kun sjeldent nævner, fordi min Tvivl har forhindret mig fra at \tlg{tilegne} mig det, om end netop derfor min Veneration er for det som for intet andet.
\end{quote}
\dcite{SLV:307.5  (SKS 6, 307)}
% SLV:307.5  (SKS 6, 307)

\begin{quote}
Bedst men ogsaa tungest har jeg lært dette ved mit Forhold til hende, hvor den ønskende Sympathi bestandigt har villet gjøre Undtagelse, hvor jeg indtil Fortvivlelse har ønsket at kunne være Alt for hende, indtil jeg i Smerten lærte, at det er uendeligt Høiere at være slet Intet for hende. Det trøster mig, at jeg i mit Forhold til hende aldrig har bildet mig ind at være Lærer, eller kaldet til at sige et Par formanende Ord. Om den Viseste anvendte 6 Timer af Dagen paa et Menneske, om han anvendte 6 andre til at betænke, hvorledes han bedst skulde gjøre det, om han vedblev dermed i 6 Aar, han var en Bedrager, hvis han vovede at sige, at han væsentligen havde gavnet ham. Og denne Tanke er for mig idetmindste Begeistringens dybeste Væld. Sprog, Kunster, Haandfærdigheder og saa videre kan det ene Menneske lære det andet, men i ethisk-religieus Forstand kan den Ene ikke gavne væsentligen den Anden. Og derfor er det skjønt og er det begeistrende at udtrykke dette i Bedragets yderste Anstrængelse, thi et Bedrag under ethisk Ansvar er ingen Magelighed, og kan altid hamle op med de formanende Ord. Det trøster mig nu, da alt det Senere er skeet, at hun intet den Lærendes Forhold har til mig, hvilket kunde forstyrre. Hvad jeg har yttret det har jeg yttret som var det til mig selv jeg talte, og jeg har hverken gjort Gestikulationer eller Applikationer. \tlg{Tilegner} hun sig det, saa gjør hun det ved sig selv, ikke stolende paa »hans Ord og Kjole.« Det er nemt nok at springe i en Omnibus og kjøre omkring og sige et Par formanende Ord, der kan ogsaa ligge noget Smukt i at ville gjøre det, men det er dumt, at kunne lære, at et Menneske formaaer slet Intet, og saa at kunne tillægge et Par formanende Ord en saa uhyre Virkning. For Virkningen tilkommer der Gud Forundringens og Beundringens Taksigelse. Thi hvert Menneske see til sig selv i Livet, i Evigheden er der Tid til at see, hvad Gud bragte ud deraf. Og dette forstaaes ikke om det Iøinefaldende ved enkelte Individer, men om den mindste Brøksdeel af Virkning i Forhold til det ubetydeligste Menneskes Gjerning.
\end{quote}
\dcite{SLV:320  (SKS 6, 320)}
% SLV:320  (SKS 6, 320)

\begin{quote}
Altsaa han seer hende igjen. Men netop fordi han for at tjene hende holder sig selv som en Afdød, har han forhindret sig i at faae et eneste ligefremt Indtryk. Her er han ifærd med at blive normal, Aberrationen ligger i, at han dog med hele sin Lidenskab bekymrer sig derom. Han faaer aldrig nogen Begivenhed i Forhold til hende eller et ligefremt Forhold til en saadan, aldrig nogen Vished, men vedblivende med hele sin Lidenskab at bekymre sig derom \tlg{tilegner} han sig Alt selv det Ubetydeligste, og stort mere faaer han ikke, ved dialektiske Anstrængelser. Dette vil sige, han kommer mere og mere til at fordybe sig i sig selv.
\end{quote}
\dcite{SLV:394  (SKS 6, 394)}
% SLV:394  (SKS 6, 394)

\begin{quote}
For at fatte Idealiteten maa jeg kunne løse det Historiske op i Idealiteten, eller gjøre, hvad man med et fromt Udtryk siger om Gud i Forhold til en Døende: klare op for det. Omvendt kommer jeg ikke ind i Idealiteten ved at gjentage den historiske Ramse. Den, som derfor ikke i Forhold til det Samme fatter ligesaa godt Slutningen ab posse ad esse som ab esse ad posse, han fatter ikke Idealiteten i dette Samme. Han opfostres kun ved Indbildninger. Idealiteten som det besjelende Princip bliver ikke uden videre historisk. Det, der kan overleveres mig, er en Mangfoldighed af Data, som ikke er Idealiteten, og saaledes er det Historiske altid raat Materiale, som nu Den, der \tlg{tilegner} sig det, veed at opløse i et posse og assimilere sig som et esse. Der er derfor intet Taabeligere paa det religieuse Gebeet end at høre det forstandige Spørgsmaal, der spørger, naar der læres Noget: er det nu ogsaa virkeligt saaledes gaaet til, for saa vil man troe det. Det, om det virkeligt saaledes er gaaet til, det er saa idealt som det fremstilles, kan kun controleres ved Idealiteten, men den kan man ikke faae historisk afproppet.
\end{quote}
\dcite{SLV:406.3  (SKS 6, 406)}
% SLV:406.3  (SKS 6, 406)

\begin{quote}
Medens man hundrede Gange læser: Umiddelbarheden er ophævet, seer man ikke en eneste Yttring om, hvorledes et Menneske bærer sig ad med at existere saaledes. Deraf turde man muligen slutte, at de Skrivende gjøre Nar ad Een og selv i al Stilhed existere i Kraft af Umiddelbarheden, og derhos leve af at skrive Bøger om, at den er ophævet. Maaskee er Systemet endda ikke saa vanskeligt at forstaae, men det, der gjør \tlg{Tilegnelsen} saa vanskelig, er, at man har sprunget alle Mellembestemmelser over om, hvorledes Individet pludseligen bliver et metaphysisk Jeg-Jeg, hvorvidt det er gjørligt, hvorvidt det er tilladeligt, hvorvidt ikke hele det Ethiske er tilsidesat, hvorvidt Systemets evige Sandhed ikke til Forudsætning (i Retning af det Existentielle, Psychologiske, Ethiske, Religieuse) har en nødvendig lille Løgn i Mangel af anden Indledning, og Systemets himmelske Text til Forklaring temmelig schofle Noter samt en tvetydig Tradition, der dispenserer de Indviede fra at tænke noget Afgjørende endog ved det mest Afgjørende. Et umiddelbart Geni kan blive Digter, Kunstner, Mathematiker og saa videre, men en Tænkende maa jo dog vide sit Forhold til den menneskelige Tilværelse, at han ikke bliver en Utydske (ved Hjælp af den rene Væren som er en U-Ting) trods alle tydske Skrifter. Han maa jo vide, hvorvidt det er ethisk og religieust forsvarligt at slutte sig metaphysisk inde, ikke at ville respektere Tilværelsens Krav, ikke paa hans lyksaliggjørende mange Tanker, ikke paa hans indbildte Jeg-Jeg, men paa hans menneskelige: Du, hvad enten Tilværelsen kalder til Lyst og Glæde og Nydelse, eller til Forfærdelse og Bæven, thi tankeløst ikke at være bleven opmærksom derpaa er lige misligt. Og kan han tankeløst oversee dette, saa gjør et Experiment med en saadan Tænker, sæt ham hen i Grækenland: han vil blive udleet i hiint udvalgte Land, lykkeligt ved sin deilige Beliggenhed, lykkeligt ved sit rige Sprog, lykkeligt ved sin uopnaaede Kunst, lykkeligt ved Folkets glade Sind, lykkeligt ved sine skjønne Piger, lykkeligt først og sidst ved Tænkere, der søgte og stræbte at forstaae sig selv og sig selv i Tilværelsen, inden de prøvede at forklare hele Tilværelsen.
\end{quote}
\dcite{SLV:454.8  (SKS 6, 454)}
% SLV:454.8  (SKS 6, 454)

\subsection{AE \normalfont(SKS 7)}

\begin{quote}
Det forskende, det spekulerende, det erkjendende Subjekt spørger saaledes vel om Sandheden, men ikke om den subjektive Sandhed, om \tlg{Tilegnelsens} Sandhed. Det forskende Subjekt er saaledes vel interesseret, men ikke uendelig personligt i Lidenskab i Retning af sin evige Salighed interesseret for sit Forhold til denne Sandhed. Langt fra det objektive Subjekt være en saadan Ubeskedenhed, en saadan Forfængelighed.
\end{quote}
\dcite{AE:29.2  (SKS 7, 29)}
% AE:29.2  (SKS 7, 29)

\begin{quote}
Dette blot her strax itide for at gjøre opmærksom paa, hvad der i anden Deel vil blive udført, at Problemet ad denne Vei slet ikke kommer afgjørende frem, det vil sige ikke kommer frem, da Problemet netop ligger i Afgjørelsen. Lad den videnskabeligt Forskende arbeide med rastløs Iver, lad ham endog forkorte sit Liv i Videnskabens begeistrede Tjeneste, lad den Speculerende ikke spare Tid og Flid: de ere dog ikke uendeligt personligt i Lidenskab interesserede, tvertimod ville de end ikke være det. Deres Betragtning vil være objektiv, interesseløs. Hvad Subjektets Forhold til den erkjendte Sandhed angaaer, da antages, at naar blot det objektive Sande er tilveiebragt, saa er det en smal Sag med \tlg{Tilegnelsen}, den følger af sig selv i Kjøbet, og am Ende er det da ligegyldigt med Individet. Deri ligger netop den Forskendes ophøiede Ro, og den Eftersnakkendes comiske Tankeløshed.
\end{quote}
\dcite{AE:30  (SKS 7, 30)}
% AE:30  (SKS 7, 30)

\begin{quote}
Den objektive Tænkning er aldeles ligegyldig mod Subjektiviteten og derved mod Inderligheden og \tlg{Tilegnelsen}; dens Meddelelse er derfor ligefrem. Det følger af sig selv, at den ingenlunde derfor behøver at være let, men den er ligefrem, den har ikke Dobbelt-Reflexionens Svig og Kunst, den har ikke hiin subjektive Tænknings gudfrygtige og humane Omsorg i at meddele sig, den lader sig ligefrem forstaae, den lader sig ramse. Den objektive Tænkning er derfor blot opmærksom paa sig selv, og er derfor ingen Meddelelse, idetmindste ingen kunstnerisk Meddelelse, forsaavidt der dog altid vilde fordres at tænke den Modtagende og at agte paa Meddelelsens Form i Forhold til Modtagerens Misforstaaelse. Den objektive Tænkning er ligesom de fleste Mennesker saa inderlig god og meddeelsom; den meddeler sig uden videre og griber i det Høieste til Forsikkringer om sin Sandhed, Anbefalinger, Forjættelser om at alle Mennesker engang ville antage denne Sandhed – saa sikker er den. Eller maaskee snarere saa usikker, thi Forsikkringen og Anbefalingen og Forjættelsen, som jo rigtignok er for de Andres Skyld, som skulle antage den, turde maaskee ogsaa være for Lærerens Skyld, der trænger til Stemmefleerhedens Sikkerhed og Tilforladelighed. Nægter Samtiden ham den, saa trækker han paa Eftertiden – saa sikker er han. Denne Sikkerhed har noget tilfælleds med den Uafhængighed, der, uafhængig af Verden, behøver Verden som Vidne til sin Uafhængighed, for at være vis paa, at man er uafhængig.
\end{quote}
\dcite{AE:77  (SKS 7, 77)}
% AE:77  (SKS 7, 77)

\begin{quote}
Meddelelsens Form er noget Andet end Meddelelsens Udtryk. Naar Tanken har faaet sit rette Udtryk i Ordet, hvilket naaes ved den første Reflexion, saa kommer den anden Reflexion, der betræffer Meddelelsens eget Forhold til Meddeleren og gjengiver den existerende Meddelers eget Forhold til Ideen. Lad os endnu engang anføre et Par Exempler, vi har jo ogsaa god Tid, thi hvad jeg skriver er ikke den forventede sidste §, hvormed Systemet er færdigt. Sæt altsaa, at En vilde meddele følgende Overbeviisning: Sandheden er Inderligheden; der er objektivt ingen Sandhed; men \tlg{Tilegnelsen} er Sandheden; sæt han havde Iver og Begeistring for at faae dette sagt, thi naar Folk hørte det, saa vare de frelste, sæt han saa sagde det ved alle Leiligheder og bevægede ikke blot de Letsvedende, men ogsaa haardføre Mennesker: hvad saa? Saa vilde der nok findes nogle Arbeidere, der havde staaet ledigt paa Torvet, og først ved dette Kald gik hen at arbeide i Viingaarden – med at forkynde denne Lære for Alle. Og hvad saa? Saa havde han yderligere modsagt sig selv, som han havde gjort det lige fra Begyndelsen, thi allerede Iveren og Begeistringen for at faae det sagt, og faae det hørt, var en Misforstaaelse. Hovedsagen var jo netop at blive forstaaet, og Forstaaelsens Inderlighed vilde jo netop være, at den Enkelte forstod det ved sig selv. Nu havde han endog bragt det til at faae Udraabere; og en Udraaber af Inderlighed er et seeværdigt Dyr. For virkeligen at meddele en saadan Overbeviisning vilde der udfordres Kunst og Selvbeherskelse; Selvbeherskelse nok til i Inderlighed at fatte, at det enkelte Menneskes Guds-Forhold er Hovedsagen og Trediemands Travlhed er Mangel paa Inderlighed og Overflødighed af elskværdig Dumhed; Kunst nok til uudtømmeligt, som Inderligheden var det, at variere Meddelelsens dobbelt-reflekterede Form. Jo mere Kunst, jo mere Inderlighed; ja havde han megen Kunst, kunde han endog gjerne sige, at han brugte den, sikker paa i næste Øieblik at kunne sikkre Meddelelsens Inderlighed, fordi han var uendelig bekymret for at bevare sin egen Inderlighed, hvilken Bekymring frelser den Bekymrede fra al positiv Sluddervorrenhed. – Sæt En vilde meddele, at ikke Sandheden er Sandheden, men Veien er Sandhed, det vil sige at Sandheden kun er i Vordelsen, i \tlg{Tilegnelsens} Proces, at der altsaa intet Resultat er; sæt han var en Menneskeven, der nødvendigviis maatte avertere alle Mennesker derom; sæt han skjød den fortræffelige Gjenvei at meddele det i den ligefremme Form i Adresseavisen, hvorved han vandt en Masse Tilhængere, medens den kunstneriske Vei vilde trods hans yderste Anstrængelse lade det uafgjort, om han havde hjulpet noget Menneske: hvad saa? Ja saa blev hans Udsagn netop et Resultat. – Sæt En vilde meddele, at al Reciperen er en Produceren; sæt han gjentog det saa tidt, at denne Sætning endogsaa blev brugt paa Forskrifter til Skjønskrivning: saa havde han rigtignok faaet sin Sætning bekræftet. – Sæt En vilde meddele den Overbeviisning, at et Menneskes Guds-Forhold er en Hemmelighed; sæt det var ret saadan en rar Mand, der holdt saa meget af andre Mennesker, at det maatte ud af ham; sæt han dog endnu havde saa megen Forstand, at han mærkede lidt af Modsigelsen ved at meddele det ligefrem, og han altsaa meddeelte det under Tausheds Løfte: hvad saa? Saa maatte han enten antage, at Discipelen var visere end Læreren, at Discipelen virkelig kunde tie, hvad Læreren ikke kunde (en ypperlig Satire over det at være Lærer!), eller han maatte blive lyksaliggjort i Galimathias, saa han slet ikke opdagede Modsigelsen. Det er besynderligt med disse rare Mennesker; det er nu saa Rørende, at det maa ud af dem – og det er saa Forfængeligt, at troe, at noget andet Menneske i sit Guds-Forhold trænger til Ens Bistand, som kunde Gud ikke nok hjælpe sig selv og den Paagjeldende. Men det er lidt anstrængende, existerende at fastholde den Tanke, at man for Gud er Intet, at det er en Spøg med al Ens Anstrængelse; lidt tugtende, at ære ethvert Menneske, saa man ikke ligefrem tør blande sig i hans Guds-Forhold, deels fordi man skal have nok med sit eget at gjøre, deels fordi Gud ikke er en Ven af Næsviished.
\end{quote}
\dcite{AE:78  (SKS 7, 78)}
% AE:78  (SKS 7, 78)

\begin{quote}
Overalt hvor det Subjektive er af Vigtighed i Erkjendelsen, altsaa \tlg{Tilegnelsen} er Hovedsagen, der er Meddelelsen et Kunstværk, den er dobbelt-reflekteret, og dens første Form er netop det Underfundige, at Subjektiviteterne maae holdes gudeligt ud fra hinanden, og ikke løbe skjørnende sammen i Objektivitet. Dette er Objektivitetens Afskedsord til Subjektiviteten.
\end{quote}
\dcite{AE:79  (SKS 7, 79)}
% AE:79  (SKS 7, 79)

\begin{quote}
Den almindelige Meddelelse, den objektive Tænkning har ingen Hemmeligheder, først den dobbelt-reflekterede subjektive Tænkning har Hemmeligheder, det vil sige al dens væsentlige Indhold er væsentlig Hemmelighed, fordi det ikke ligefrem lader sig meddele. Dette er Hemmelighedens Betydning. Det, at Erkjendelsen ikke er ligefrem at udsige, fordi det Væsentlige ved Erkjendelsen netop er \tlg{Tilegnelsen}, gjør at den bliver en Hemmelighed for Enhver, der ikke paa samme Maade ved sig selv er dobbelt reflekteret; men det, at det er Sandhedens væsentlige Form, gjør, at denne ikke kan siges paa nogen anden Maade. Naar derfor En vil meddele det ligefrem, saa er han dum; og naar En vil fordre det af ham, saa er han ogsaa dum. Ligeoverfor en saadan svigefuld kunstnerisk Meddelelse vilde den sædvanlige menneskelige Dumhed raabe: det er Egoisme. Naar saa Dumheden seirede, og Meddelelsen blev ligefrem, saa havde Dumheden vundet saa meget, at Meddeleren var bleven ligesaa dum.
\end{quote}
\dcite{AE:79.2  (SKS 7, 79)}
% AE:79.2  (SKS 7, 79)

\begin{quote}
Det gjelder da om, at det existerende Subjekts Tænkning har en Form, hvori det kan gjengive dette. I det ligefremme Udsagn siger det netop noget Usandt, hvis det siger dette; thi i det ligefremme Udsagn er netop Svigen udeladt, saa altsaa Meddelelsens Form forstyrrer, ligesom naar den Epileptiskes Tunge siger det forkeerte Ord, om den Talende maaskee ikke mærker det saa tydeligt som den Epileptiske. Lader os tage et Exempel. Det existerende Subjekt er evigt, men som existerende er det timeligt. Uendelighedens Svig er nu, at Dødens Mulighed er tilstede i ethvert Øieblik. Al positiv Tilforladelighed er saaledes gjort mistænkt. Er jeg mig ikke i ethvert Øieblik dette bevidst, saa er min positive Tillid til Livet Barnagtighed, uagtet den er bleven spekulativ, fornem paa den systematiske Cothurne; men bliver jeg mig dette bevidst, saa er Uendelighedens Tanke saa uendelig, at den ligesom forvandler min Existents til et forsvindende Intet. Hvorledes gjengiver nu det existerende Subjekt denne sin Tanke-Existents? At det er saaledes med det at existere, veed ethvert Menneske, men de Positive veed det positivt, det vil sige de veed det slet ikke – men det forstaaer sig, de have jo ogsaa travlt med hele Verdenshistorien. Engang om Aaret ved en høitidelig Leilighed griber denne Tanke dem, og nu udsige de i Forsikkringens Form, at det er saa. Men dette at de kun engang ved en høitidelig Leilighed bemærke det, forraader tilstrækkeligen, at de ere meget positive, og det at de sige det i Forsikkringens Tilforladelighed, viser, at de, selv idet de udsige det, ikke veed, hvad de sige, hvorfor de ogsaa i næste Øieblik kunne glemme det igjen. I Forhold nemlig til saadanne negative Tanker er en svigefuld Form den eneste adæqvate, fordi den ligefremme Meddelelse ligger i Continueerlighedens Tilforladelighed, den svigefulde Tilværelse derimod, hvor jeg fatter den, isolerer mig. Den, der er opmærksom herpaa, den der, nøiet med at være Menneske, har Styrke og Ørkesløshed nok til ikke at ville bedrages for saa at faae Lov at spreche om hele Verdenshistorien, beundret af Ligesindede, spottet af Tilværelsen, han undgaaer det ligefremme Udsagn. Som bekjendt var Socrates en Dagdriver, der hverken brød sig om Verdenshistorien eller Astronomien (denne opgav han som Diogenes fortæller, og naar han senere stod stille og stirrede hen for sig, kan jeg dog, uden forøvrigt at afgjøre Noget om hvad han egentlig foretog sig, ikke antage, at han kiggede Stjerner), men havde god Tid og Særhed nok til at bekymre sig om det simple Menneskelige, hvilken Bekymring, besynderligt nok, ansees for Særhed hos Mennesker, medens det derimod slet ikke er sært, at have travlt med Verdenshistorien, Astronomien og andet Saadant. Ifølge en udmærket Afhandling i det Fyenske Tidsskrift seer jeg, at Socrates skal have været noget ironisk. Det var virkelig paa den høie Tid, at det blev sagt, og jeg er nu i det Tilfælde at turde beraabe mig paa hiin Afhandling, naar jeg antager noget Lignende. Socrates's Ironi bruger blandt Andet den Form, netop naar han vil have Uendeligheden frem, at han i første Instants taler som en Afsindig. Ligesom Tilværelsen er lumsk, saaledes er hans Tale ogsaa, maaskee (thi jeg er ikke en saadan vis Mand som den Positive i det fyenske Tidsskrift) for at forhindre, at faae en rørt og troende Tilhører, der positivt \tlg{tilegnede} sig Udsagnet om Tilværelsens Negativitet. Denne Afsindighed i første Instants kan da tillige have betydet for Socrates, at han, medens han talte med Menneskene, tillige i det Sagte confererede privatissime med Ideen, hvilket Ingen kan forstaae som kun kan tale ligefrem, og hvilket det ikke kan hjælpe at sige En eengang for alle, da Hemmeligheden netop er, at det altid maa være overalt tilstede i Tanken og i dennes Gjengivelse, ligesom det overalt er tilstede i Tilværelsen. Det er forsaavidt netop det Rigtige ikke at blive forstaaet, thi derved er man jo sikkret mod Misforstaaelse. Naar saaledes Socrates etsteds siger, at det er mærkeligt nok, at Skipperen, der har sat En over fra Grækenland til Italien, naar han er ankommen, gaaer rolig frem og tilbage paa Strandbredden, modtager sin Betaling, som havde han gjort noget Godt, medens han dog ikke kan vide, om han har gavnet Passagererne, eller om det ikke havde været disse bedre at sætte Livet til paa Havet: saa taler han egentlig som en Afsindig. Maaskee har En af de Tilstedeværende virkeligen anseet ham for gal (thi ifølge Plato og Alcibiades skal det have været en almindeligere Mening idetmindste at ansee ham for sær, ατοπος); maaskee har en Anden tænkt, at det var en snurrig Vending i Talen, maaskee. Socrates derimod, han har maaskee paa samme Tid holdt et lille Stævnemøde med sin Idee, med Uvidenheden. Dersom han har fattet Uendeligheden i Uvidenhedens Form, saa maa han jo overalt have denne hos sig. Sligt uleiliger ikke en Privat-Docent, han gjør det eengang om Aaret i § 14 med Pathos, og deri gjør han vel, at han ikke gjør det anderledes, hvis han nemlig har Kone og Børn og Udsigt til et godt Levebrød – men ingen Forstand at tabe.
\end{quote}
\dcite{AE:82  (SKS 7, 82)}
% AE:82  (SKS 7, 82)

\begin{quote}
Som den subjektive existerende Tænker selv existerer saaledes, saa gjengiver hans Fremstilling ogsaa dette, og derfor kan Ingen uden videre \tlg{tilegne} sig hans Pathos. Som de comiske Partier i et romantisk Drama, saaledes slynger ogsaa Comiken sig gjennem Lessings Fremstilling, maaskee stundom paa urette Sted, maaskee, maaskee ikke; med Bestemthed kan jeg ikke sige det. Hauptpastor Goetze er en høist ergötzlich Figur, som Lessing har comisk conserveret for Udødeligheden ved at gjøre ham uadskillelig fra sin Fremstilling. Det forstaaer sig, det er forstyrrende, man kan ikke med den Tilforladelighed hengive sig til Lessing, som til Deres Fremstilling, der med ægte speculativ Alvor gjøre Alt ud af Eet, og saaledes have Alt færdigt.
\end{quote}
\dcite{AE:90  (SKS 7, 90)}
% AE:90  (SKS 7, 90)

\begin{quote}
Christendommen vil jo skjænke den Enkelte en evig Salighed, et Gode der ikke uddeles i større Partier, men kun til Een og Een ad Gangen. Antager den end at Subjektiviteten er, som \tlg{Tilegnelsens} Mulighed, Muligheden af at annamme dette Gode, saa antager den dog ikke at Subjektiviteten uden videre er fix og færdig, uden videre endog blot har en virkelig Forestilling om dette Godes Betydning. Denne Subjektivitetens Udvikling eller Omskabelse, denne dens uendelige Concentration i sig selv ligeoverfor Forestillingen om Uendelighedens høieste Gode, en evig Salighed, er den udviklede Mulighed af Subjektivitetens første Mulighed. Saaledes protesterer Christendommen mod al Objektivitet; den vil, at Subjektet uendeligt skal bekymre sig om sig selv. Det, den spørger om, er Subjektiviteten, først i denne er Christendommens Sandhed, hvis den overhovedet er, objektivt er den slet ikke. Om den kun er i et eneste Subjekt, saa er den kun i ham, og der er større christelig Glæde i Himlen over denne Ene end over Verdenshistorien og Systemet, hvilke, som objektive Magter, ere incommensurable for det Christelige.
\end{quote}
\dcite{AE:122.2  (SKS 7, 122)}
% AE:122.2  (SKS 7, 122)

\begin{quote}
For den existerende Aand qva existerende Aand bliver Spørgsmaalet igjen om Sandheden; thi den abstrakte Besvarelse er kun for det Abstraktum, hvilket en existerende Aand bliver ved at abstrahere fra sig selv qva existerende, hvad den kun momentviis kan gjøre, medens den dog selv i saadanne Øieblikke betaler Tilværelsen sin Gjeld ved alligevel at existere. Altsaa det er en existerende Aand, der spørger om Sandheden, formodentlig for at ville existere i den, men i ethvert Tilfælde er Spørgeren sig bevidst at være et existerende enkelt Menneske. Saaledes troer jeg at kunne gjøre mig forstaaelig for enhver Græker, og for ethvert fornuftigt Menneske. Om en tydsk Philosoph følger sin Lyst til at skabe sig, og først skaber sig om til et overfornuftigt Noget, ligesom Guldmagere og Troldmænd udpynte sig phantastisk, for da paa en yderst tilfredsstillende Maade at besvare Spørgsmaalet om Sandheden: vedkommer ikke mig, saa lidet som hans tilfredsstillende Svar, der vistnok er yderst tilfredsstillende – naar man phantastisk er bleven udklædt. Om derimod en tydsk Philosoph gjør det eller ikke, vil Enhver let forvisse sig om, der begeistret samler sin Sjel paa at ville lade sig veilede af en saadan Viis, uden Critik blot følgsomt benyttende hans Veiledning ved at ville danne sin Existents efter den; netop naar man saaledes begeistret forholder sig som en Lærende til en saadan tydsk Professor, saa realiserer man det fortræffeligste Epigram over ham, thi en saadan Speculant er Intet mindre tjent med end en Lærendes redelige og begeistrede Iver for at udtrykke og realisere, for existerende at \tlg{tilegne} sig hans Viisdom, da den er et Noget, som Hr. Professoren har indbildt sig selv, skrevet Bøger om, men aldrig selv forsøgt, ja det er end ikke faldet ham ind, at det skulde gjøres. Der gives Speculanter, der, ligesom hiin Portskriver skrev hvad han ikke selv kunde læse, i den Formening at det blot var hans Sag at skrive, blot skrive og skrive hvad der, naar det, om jeg saa tør sige, skal læses ved Hjelp af Gjerning, viser sig at være Nonsens, med mindre det maaskee skulde være bestemt for phantastiske Væsener.
\end{quote}
\dcite{AE:176  (SKS 7, 176)}
% AE:176  (SKS 7, 176)

\begin{quote}
Idet der for den existerende Aand qva existerende bliver Spørgsmaal om Sandheden, kommer hiin Sandhedens abstrakte Reduplikation igjen, men Existentsen selv, Existentsen selv i Spørgeren, der jo existerer, holder de tvende Momenter ud fra hinanden, og Reflexionen udviser tvende Forhold. For den objektive Reflexion bliver Sandheden et Objektivt, en Gjenstand, og det gjælder om at see bort fra Subjektet; for den subjektive Reflexion bliver Sandheden \tlg{Tilegnelsen}, Inderligheden, Subjektiviteten, og det gjelder netop om existerende at fordybe sig i Subjektiviteten.
\end{quote}
\dcite{AE:176.2  (SKS 7, 176)}
% AE:176.2  (SKS 7, 176)

\begin{quote}
Naar Subjektiviteten er Sandheden, maa Sandhedens Bestemmelse tillige indeholde i sig et Udtryk for Modsætningen til Objektiviteten, en Erindring fra hiint Veiskille, og dette Udtryk angiver da tillige Inderlighedens Spændstighed. Her er en saadan Definition paa Sandhed: den objektive Uvished, fastholdt i den meest lidenskabelige Inderligheds \tlg{Tilegnelse}, er Sandheden, den høieste Sandhed der er for en Existerende. Der hvor Veien svinger af (og hvor det er, lader sig ikke objektivt sige, da det netop er Subjektiviteten), bliver den objektive Viden sat i Bero. Objektivt har han da kun Uvisheden, men netop dette strammer Inderlighedens uendelige Lidenskab, og Sandheden er netop dette Vovestykke, med Uendelighedens Lidenskab at vælge det objektivt Uvisse. Jeg betragter Naturen for at finde Gud, jeg seer jo ogsaa Almagt og Viisdom, men jeg seer tillige meget Andet som ængster og forstyrrer. Summa summarum heraf bliver den objektive Uvished, men netop derfor er Inderligheden saa stor, fordi Inderligheden omfatter den objektive Uvished med Uendelighedens hele Lidenskab. I Forhold til en mathematisk Sætning for Exempel er Objektiviteten given, men derfor er dens Sandhed ogsaa en ligegyldig Sandhed.
\end{quote}
\dcite{AE:187  (SKS 7, 187)}
% AE:187  (SKS 7, 187)

\begin{quote}
Da jeg ikke er ganske ukjendt med hvad der er sagt og skrevet om Christendommen, kunde jeg vel sige Eet og Andet; her vil jeg det dog ikke, men blot gjentage, at der er Eet, jeg skal vogte mig for at sige om den: at den til en vis Grad er sand. Det var jo dog muligt, at Christendommen var Sandheden, det var jo dog muligt, at der engang kom en Dom, hvor Adskillelsen blev Inderlighedens Forhold til den. Sæt der da traadte et Menneske frem som maatte sige: vel har jeg ikke troet, men saa meget har jeg dog æret Christendommen, at jeg har anvendt hver en Time i mit Liv for at grunde over den; eller der kom et Menneske frem, om hvem Anklageren maatte sige: han har forfulgt de Christne, og den Anklagede sagde: ja, jeg vedgaaer det, Christendommen har opflammet min Sjel, saa jeg ikke har villet Andet end udrydde den af Verden, netop fordi jeg fattede dens forfærdelige Magt; eller sæt der kom et Menneske frem, om hvem Anklageren maatte sige: han har afsvoret Christendommen, og den Anklagede sagde: ja, det er sandt, thi jeg indsaae, at Christendommen var en saadan Magt, at den, hvis jeg gav den een Finger, tog mig ganske, og ganske kunde jeg ikke tilhøre den – men sæt saa at der endeligen kom en vever Privat-Docent med iilsomme og travle Fjed og talte saaledes: jeg er ikke som hine Trende, jeg har ikke blot troet, men endogsaa forklaret Christendommen, viist, at den, som den foredroges af Apostlene og \tlg{tilegnedes} i de første Aarhundreder, kun til en vis Grad er sand, hvorledes den derimod ved Spekulationens Forstaaelse er den sande Sandhed, hvisaarsag jeg maa udbede mig en passende Godtgjørelse for mine Fortjenester af Christendommen: hvilken af disse Fires Stilling var da den forfærdeligste? Det var jo dog muligt, at Christendommen var Sandheden, sæt den nu, da dens utaknemlige Børn vil have den erklæret umyndig under Spekulationens Værgemaal, sæt den nu, ligesom hiin græske Digter, hvis Børn ogsaa fordrede den alderstegne Fader erklæret umyndig, forbausede Dommerne og Folket ved at digte en af sine skjønneste Tragedier til Tegn paa, at han endnu var myndig, sæt den saaledes reiste sig forynget: der var dog Ingen, hvis Stilling vilde blive saa forlegen som Privat-Docenternes. Jeg nægter ikke, at det er fornemt at staae saa høit over Christendommen; jeg nægter ikke, at det er beqvemt at være Christen og dog være fritaget for det Martyrium, som altid bliver, selv om ingen udvortes Forfølgelse hjemsøger, selv om en Christen forblev ubemærket, som havde han slet ikke levet, det Martyrium at troe mod Forstanden, den Livsfare at ligge paa de 70,000 Favne Vand og først der finde Gud. See, den Vadende han føler sig for med Foden, at han ikke gaaer længere ud end han kan bunde: og saaledes føler den Forstandige sig med Forstanden for i Sandsynligheden og finder Gud der, hvor Sandsynligheden slaaer til, og takker ham paa Sandsynlighedens store Høitidsdage, naar han har faaet et rigtigt godt Levebrød, og der ovenikjøbet er Sandsynlighed for snart at avancere; naar han faaer sig baade en smuk og rar Pige til Hustru, og selv Krigsraad Marcussen siger, at det vil blive et lykkeligt Ægteskab, og at Pigen er af den Art Skjønheder, der efter al Sandsynlighed vil holde sig længe, og saaledes bygget, at hun efter al Sandsynlighed vil føde smukke og stærke Børn. At troe mod Forstanden er noget Andet, og at troe med Forstanden lader sig slet ikke gjøre, thi Den der troer med Forstanden taler kun om Levebrød og Hustru og Ager og Øxne og andet Saadant, som slet ikke er Troens Gjenstand, da Troen altid takker Gud, altid i Livsfaren, i hiint Uendelighedens og Endelighedens Sammenstød, der netop er Livsfaren for Den, som er sammensat af begge. Sandsynligheden er derfor den Troende saa lidet kjær, at han frygter den meest af Alt, da han godt veed, at det er, fordi han begynder at tabe Troen. Troen har nemlig tvende Opgaver: at passe paa og i ethvert Øieblik opdage Usandsynligheden, Paradoxet, for da med Inderlighedens Lidenskab at fastholde det. I Almindelighed forestiller man sig det saaledes, at det Usandsynlige, det Paradoxe er Noget, som Troen blot forholder sig lidende til, den maa nu foreløbigt nøies med dette Forhold, og lidt efter lidt bliver det nok bedre, det er der endogsaa Sandsynlighed for. O, vidunderlige Confusionsmagerie i at tale om Troen! Man skal begynde at troe i Tillid til, at der er Sandsynlighed for, at det nok bliver bedre. Paa den Maade faaer man dog Sandsynligheden smuglet ind og sig selv forhindret i at troe; paa den Maade er det let at forstaae, at Frugten af at have troet i længere Tid bliver den, at man ophører at troe, istedenfor at man skulde troe, at Frugten blev, at man troede inderligere. Nei Troen forholder sig selvvirksom til det Usandsynlige og det Paradoxe, selvvirksom i at opdage det og hvert Øieblik fastholde det – for at kunne troe. Der hører allerede hele Uendelighedens Lidenskab og dennes Sluttethed til, for at standse ved det Usandsynlige, thi det Usandsynlige og det Paradoxe er ikke til at naae ved en Forstandens Qvantiteren af det mere og mere Vanskelige. Der hvor Forstanden fortvivler, er allerede Troen med for ret at gjøre Fortvivlelsen afgjørende, at Troens Bevægelse ikke bliver en Omsætning indenfor Forstandens pruttende Omfang. Men at troe mod Forstanden det er et Martyrium, at begynde at faae Forstanden lidt med sig, er Anfægtelse og Tilbagegang. Dette Martyrium er Spekulanten fri for. At han maa studere, især at han maa læse mange af de moderne Bøger, indrømmer jeg gjerne er byrdefuldt, men Troens Martyrium er dog noget Andet.
\end{quote}
\dcite{AE:212  (SKS 7, 212)}
% AE:212  (SKS 7, 212)

\begin{quote}
Da jeg havde fattet dette, blev det mig tillige tydeligt, at dersom jeg vilde meddele Noget desangaaende, maatte det fornemlig gjelde om, at min Fremstilling blev i den indirecte Form. Dersom nemlig Inderligheden er Sandheden, saa er Resultat kun Skramlerie, man ikke skal besvære hinanden med, og det at ville meddele Resultat en unaturlig Omgang mellem Menneske og Menneske, forsaavidt ethvert Menneske er Aand, og Sandheden netop \tlg{Tilegnelsens} Selvvirksomhed, hvilken et Resultat forhindrer. Lad Læreren i Forhold til den væsentlige Sandhed (thi ellers er jo det ligefremme Forhold mellem Lærer og Lærende ganske i sin Orden), som man siger, have megen Inderlighed og gjerne vilde forkynde sin Lære Dag ud og Dag ind: dersom han antager, at der er et ligefremt Forhold mellem ham og den Lærende, saa er hans Inderlighed ikke Inderlighed, men umiddelbar Udgydelse, thi den Ærbødighed for den Lærende, at han netop er i sig selv Inderligheden, er Lærerens Inderlighed. Lad en Lærende være begeistret, og i de stærkeste Udtryk forkynde Lærerens Priis, og saaledes lægge, som man siger, sin Inderlighed for Dagen: hans Inderlighed er ikke Inderlighed, men den umiddelbare Hengivelse, thi den gudfrygtige tause Overeenskomst, ifølge hvilken den Lærende ved sig selv \tlg{tilegner} sig det Lærte, fjernende sig fra Læreren fordi han vender sig ind i sig selv, det er netop Inderligheden. Pathos er vel Inderlighed, men det er umiddelbar Inderlighed, derfor gives den ud, men Pathos i Modsætningens Form er Inderlighed, den bliver hos den Meddelende, skjøndt den gives ud, og den kan ikke ligefrem \tlg{tilegnes} uden gjennem den Andens Selvvirksomhed, og Modsætningens Form er netop Inderlighedens Kraftmaaler. Jo fuldeligere Modsætningens Form er, desto større Inderlighed, og jo mindre den er tilstede indtil Meddelelsen er ligefrem, desto mindre Inderlighed. Det kan være vanskeligt nok for et begeistret Genie, der gjerne vil lyksaliggjøre alle Mennesker og føre dem til Sandheden, at lære saaledes at holde igjen og fatte Reduplikationens NB., fordi Sandheden ikke er som en Circulaire, paa hvilken der samles Underskrifter, men i Inderlighedens valore intrinseco; for en Løsgænger og Letsindig falder dette mere naturligt at forstaae. Saasnart Sandheden, den væsentlige Sandhed, kan antages at være Enhver bekjendt, saa er \tlg{Tilegnelsen} og Inderligheden det for hvilket der maa arbeides, og her kan kun arbeides i indirecte Form. En Apostels Stilling er en anden, thi han har at forkynde Sandheden som er ubekjendt, og derfor kan den ligefremme Meddelelse altid midlertidigt have sin Gyldighed.
\end{quote}
\dcite{AE:221  (SKS 7, 221)}
% AE:221  (SKS 7, 221)

\begin{quote}
Skriftet »Gjentagelsen« blev paa Titelbladet kaldet »psychologisk Experiment«. At dette var en dobbelt reflekteret Meddelelses-Form, blev mig snart tydeligt. Thi derved, at Meddelelsen skeer i Experimentets Form, danner den sig selv en Modstand, og Experimentet befæster et svælgende Dyb mellem Læser og Forfatter, sætter Inderlighedens Skilsmisse mellem dem, saa den ligefremme Forstaaelse er gjort umulig. Experimentet er Meddelelsens bevidste, drillende Tilbagekaldelse, hvilket altid er af Vigtighed for en Existerende der skriver for Existerende, at ikke Forholdet forandres til en Ramsende der skriver for Ramsende. Dersom en Mand vilde staae paa eet Been, eller i en snurrig dandsende Stilling svinge med Hatten, og i denne Attitude foredrage noget Sandt, saa vilde hans faae Tilhørere dele sig i to Classer; og mange fik han ikke, da de fleste strax vilde opgive ham. Den ene Classe vilde sige: hvor skulde det kunne være det Sande han siger, naar han gesticulerer saaledes; den anden vilde sige: ja hvad enten han nu slaaer Entrechats eller han gaaer paa Hovedet, om han saa slog Kulbøtter, det han siger er sandt, det vil jeg \tlg{tilegne} mig og lade ham fare. Saaledes ogsaa med Experimentet. Er det Udsagte den Skrivendes Alvor, saa beholder han Alvoren væsentlig for sig selv; opfatter den Modtagende det som Alvor, gjør han det væsentligen ved sig selv, og dette er netop Alvoren, og allerede i Barne-Underviisningen gjør man jo Forskjel mellem »Udenads-Læsning« og »Forstands-Øvelse«, hvilken Forskjel ofte er paafaldende nok i Forhold til den systematiske »Udenads-Læsning«. Experimentets Mellemværende begunstiger de Tvendes Inderlighed bort fra hinanden i Inderlighed. Denne Form vandt ganske mit Bifald, og jeg troede deri tillige at opdage, at de pseudonyme Forfattere bestandigt holdt Sigte paa det at existere, og saaledes vedligeholdt en indirecte Polemik mod Speculationen. Naar et Menneske veed Alt, men veed det udenad, saa er Experimentets Form et godt Explorations-Middel; man sige ham endog hvad han veed i denne Form: han kjender det ikke. – Senere har en ny Pseudonym, Frater Taciturnus, udviist Experimentets Plads i Forhold til den æsthetiske, ethiske og religieuse Frembringelse. (confer »Stadier paa Livets Vei« P. 340 o. flere følgende § 3).
\end{quote}
\dcite{AE:240  (SKS 7, 240)}
% AE:240  (SKS 7, 240)

\begin{quote}
Og saa kom endeligen mine Smuler; thi Existents-Inderligheden var nu bestemmet saavidt hen, at det Christelig-Religieuse kunde fremtages uden strax at forvexles med Allehaande. Dog Eet endnu. Magister Kierkegaards opbyggelige Taler havde stadigen fulgt med, i mine Øine et Vink om at han fulgte med, og var det mig paafaldende, at de fire sidste antog et omhyggeligt fortonet Anstrøg af det Humoristiske. Saaledes ender vel ogsaa hvad der er at naae ved Immanentsen. Medens det Ethiskes Fordring gjøres gjeldende, medens Livet og Tilværelsen accentueres som en møisommelig Gang, bliver dog Afgjørelsen ikke sat i et Paradox, og den metaphysiske Erindringens Tilbagetagen i det Evige er bestandig mulig, og giver Immanentsen Humorens Anstrøg som en Uendelighedens Tilbagekaldelse af det Hele i Evighedens Afgjorthed bag ved. Existentsens (det vil sige det at existere) paradoxe Udtryk som Synd, den evige Sandhed som Paradoxet ved at være blevet til i Tiden, kort hvad der er afgjørende for det Christelig-Religieuse findes ikke i de opbyggelige Taler, om hvilke Nogle, efter hvad Magisteren sagde, meente at man godt kunde kalde dem Prædikener, medens Andre indvendte, at de vare ikke rigtige Prædikener. Humor er, naar den benytter de christelige Bestemmelser, en falsk Gjengivelse af den christelige Sandhed, da Humor ikke er væsentlig forskjellig fra Ironie, men væsentlig forskjellig fra Christendom, og væsentligen ikke anderledes forskjellig fra Christendom end Ironie er det. Den er kun tilsyneladende forskjellig fra Ironie ved tilsyneladende at have \tlg{tilegnet} sig hele det Christelige, uden dog paa en afgjørende Maade at have \tlg{tilegnet} sig det (men det Christelige ligger netop i Afgjørelsen og Afgjortheden), hvorimod det for Ironien Væsentlige: Erindringens Tilbagetagen ud af Timeligheden ind i det Evige, atter er det Væsentlige for Humor. Tilsyneladende giver Humor det at existere større Betydning end Ironien gjør, men dog er Immanentsen übergreifend, og det Mere eller Mindre er en forsvindende Qvantiteren mod det Christeliges qvalitative Afgjorthed. Humor bliver derfor den sidste terminus a quo i Forhold til at bestemme det Christelige. Humor er, naar den benytter de christelige Bestemmelser (Synd, Synds-Forladelse, Forsoning, Gud i Tiden og saa videre), ikke Christendom, men en hedensk Speculation, der har faaet alt det Christelige at vide. Den kan komme det Christelige skuffende nær, men der hvor Afgjørelsen fanger, der hvor Existentsen fanger den Existerende ligesom Bordet fanger, saa han maa blive i Existentsen, medens Erindringens og Immanentsens Bro bagved er afhugget; der hvor Afgjørelsen bliver i Øieblikket, og Bevægelsen fremad mod Forholdet til den evige Sandhed, der blev til i Tiden: der er Humoren ikke med. Paa samme Maade bedrager den moderne Speculation, ja man kan end ikke sige den bedrager, thi der er snart ikke Nogen at bedrage, og Speculationen gjør det bona fide. Speculationen gjør det Kunststykke, at forstaae hele Christendommen, men vel at mærke, den forstaaer den ikke christeligt, men speculativt, hvilket netop er Misforstaaelsen, da Christendommen er lige Modsætningen til Speculation. – Magister Kierkegaard har formodentligen nok vidst, hvad han gjorde, da han kaldte de opbyggelige Taler: opbyggelige Taler, samt hvorfor han afholdt sig fra Brugen af christelig-dogmatiske Bestemmelser, fra at nævne Christi Navn og saa videre, hvad man ellers i vor Tid gjør frisk væk, medens Categorierne, Tankerne, det Dialektiske i Fremstillingen kun er Immanentsens. Som de pseudonyme Skrifter, foruden hvad de directe ere, indirecte ere en Polemik mod Speculationen, saaledes er disse Taler det ogsaa, ikke derved, at de ikke ere speculative, thi de ere netop speculative, men derved, at de ikke ere Prædikener. Havde Forfatteren kaldt dem Prædikener, havde han været en Sluddermads. De ere opbyggelige Taler, Forfatteren gjentager ordret i Forordet »at han ikke er Lærer« og Talerne ikke »til Opbyggelse«, ved hvilken Bestemmelse deres teleologiske Betydning allerede i Forordet er humoristisk tilbagekaldt. De ere »ikke Prædikener«; Prædikenen svarer nemlig til det Christelige, og til Prædikenen svarer en Præst, og en Præst er væsentligen hvad han er ved Ordinationen, og Ordinationen er en Lærers paradoxe Forvandling i Tiden, hvorved han i Tiden bliver noget Andet end hvad der vilde være den immanente Udvikling af Genie, Talent, Gave og saa videre Fra Evighed er dog vel Ingen ordineret, eller saasnart han fødes, istand til at erindre sig som ordineret. Paa den anden Side er Ordinationen en caracter indelebilis. Hvad vil dette sige Andet, end at atter her Tiden bliver afgjørende for det Evige, hvorved Erindringens immanente Tilbagetagen i det Evige er forhindret. Ved Ordinationen staaer atter det christelige Notabene. Om det er rigtigt, om Speculationen og Humor ikke har Ret, er et ganske andet Spørgsmaal; men om Speculationen end havde nok saa meget Ret, den har aldrig Ret i at udgive sig for Christendom.
\end{quote}
\dcite{AE:246  (SKS 7, 246)}
% AE:246  (SKS 7, 246)

\begin{quote}
Hvad er Daaben uden \tlg{Tilegnelse}? Ja, den er Muligheden af, at det døbte Barn kan blive en Christen, hverken mere eller mindre. Parallellen vilde være: som man maa være født, være bleven til, for at blive Menneske, thi et Barn er det endnu ikke: saaledes maa man være døbt for at blive en Christen. For den Ældre, der ikke er døbt som Barn, gjelder det, at han bliver Christen ved Daaben, fordi han i Daaben kan have Troens \tlg{Tilegnelse}. Tag \tlg{Tilegnelsen} bort fra det Christelige, hvad er saa Luthers Fortjeneste? Men luk ham op, og mærk i hver Linie \tlg{Tilegnelsens} stærke Pulsslag, mærk den i hele Stilens skjælvende Fremadskynden, der bestandig ligesom har hiin Forfærdelsens Tordenveir bag efter sig, som dræbte Alexius og skabte Luther. Havde Papismen ikke Objektivitet og objektive Bestemmelser og det Objektive, det Objektive, det Objektive i Overflod? Hvad manglede den? \tlg{Tilegnelse}, Inderlighed. »Aber unsere spitzfindigen Sophisten sagen in diesen Sacramenten nichts von dem Glauben, sondern plappern nur fleißig von den wirklichen Kräften der Sacramente (det Objektive), denn sie lernen immerdar, und kommen doch nimmer zu der Erkentniß der Wahrheit« (Von der babylonischen Gefangenschaft, Gerlachs lille Udgave 4de Bind pagina 195). Men ad den Vei maatte de jo dog komme til Sandheden, hvis Objektiviteten var Sandheden. Lad det saa ti Gange være sandt, at Christendommen ikke ligger i Differentsen, lad det være Jordlivets saligste Trøst, at det er Christendommens hellige Menneskelighed, at den kan \tlg{tilegnes} af Alle – men skal og bør dog dette forstaaes saaledes, at Enhver er Christen uden videre ved at være døbt da han var fjorten Dage gammel? At være Christen er ingen Mageligheds-Sag, den Eenfoldige skal lige saa vel som den Vise være saa god at existere deri – saa bliver det at være Christen noget Andet end at have en Døbe-Attest liggende i Skuffen, at producere den, naar man skal være Student, naar man vil holde Bryllup; noget Andet end at gaae med en Døbe-Attest i Vestelommen hele Livet igjennem. Men det at være Christen er efterhaanden blevet Noget, man saadan er uden videre, og i Retning af Ansvar snarere Noget, der vedkommer Ens Forældre end En selv: at de dog ikke have forsømt at lade En døbe. Deraf fremkommer det besynderlige Phænomen, hvilket maaskee endda ikke er saa sjeldent i Christenheden, at en Mand, der for sit Vedkommende har tænkt at hans Forældre vel har sørget for at han blev døbt, og dermed ladet den Sag være afgjort, naar han selv bliver Fader, saa vaagner ganske rigtigt den Bekymring, at faae sit Barn døbt, saaledes at Bekymringen for at blive Christen er gaaet over fra Individet selv til Værgen. I Qvalitet af Værge bekymrer Faderen sig om, at Barnet bliver døbt, maaskee ogsaa i Betragtning af alle de Politie-Ubehageligheder og Bryderier som Barnet vil blive udsat for, naar det ikke er døbt. Og Evigheden hisset og Dommens høitidelige Alvor (hvor det vel at mærke skal afgjøres, om jeg var en Christen, ikke om jeg i Qvalitet af Værge har sørget for at mine Børn bleve døbte) forvandles til en Gadescene eller en Scene i et Pas-Contoir, hvor de Afdøde komme rendende med deres Attester – fra Klokkeren. Lad det saa ti Gange være sandt, at Daaben er et guddommeligt Regjerings-Pas for Evigheden; men naar Letsindighed og Verdslighed vil bruge den som Passeerseddel, er den det saa ogsaa? Daaben er dog vel ikke den Lap Papir, som Klokkeren udsteder – og stundom skriver feil paa; Daaben er dog vel ikke blot det udvortes Faktum, at man blev døbt d. 7de Sept. Klokken 11. At Tiden, Existentsen i Tiden bliver afgjørende for en evig Salighed er overhovedet saa paradox, at Hedenskabet ikke kan tænke det; men at faae det Hele afgjort i en Alder af fjorten Dage d. 7de Sept. i Løbet af fem Minutter, synes dog næstendeels lidt for Meget af det Paradoxe. Det manglede blot, at man ogsaa blev viet paa Vuggen til Den og Den, indskrevet i den og den borgerlige Stilling og saa videre, saa vilde man i den unge Alder af fjorten Dage have afgjort Alt for hele sit Liv – med mindre den senere Afgjørelse blev den, at man gjorde det om, hvad man vel vilde finde Umagen værd, i Forhold til det projekterede Ægteskab, men maaskee ikke i Forhold til Christendommen. See, engang i Verden var det saaledes, at naar Alt brast for et Menneske, var der dog det Haab tilbage at blive Christen; nu er man det, og paa saa mange Maader fristet til at glemme – at blive det.
\end{quote}
\dcite{AE:334  (SKS 7, 334)}
% AE:334  (SKS 7, 334)

\begin{quote}
Jeg skal nu korteligen antyde de Opfattelser af Skyld og de dertil svarende Opfattelser af Fyldestgjørelse, der ere lavere end den skjulte Inderligheds evige Erindren af Skylden. Da jeg i den foregaaende § har været saa udførlig, kan jeg her fatte mig desto kortere; thi hvad der i foregaaende § er blevet viist som lavere, maa her vise sig igjen. Som allevegne respekteres ogsaa her kun Categorien, og tager jeg derfor med hvad der, skjøndt ofte kaldet christeligt, naar det henføres til Categorien viser sig ikke at være det. At en Præst, selv en Silke-Præst, at en tituleret og med virkelige Christne rangerende døbt Christen laver Noget sammen, kan da vel ikke gjøre dette Noget til Christendom, saa lidet som det ligefrem følger, fordi en Doctor skriver Noget sammen paa Recept-Papir, at dette Noget derfor er Medicin – det kan jo ogsaa være Pøit. Der er Intet saaledes Nyt i Christendommen, at det jo tilsyneladende har været før i Verden, og dog er Alt Nyt. Dersom En nu bruger Christendommens og Christi Navn, men Categorierne (trods Udtrykkene) ere Intet mindre end christelige, er det saa Christendom? Eller dersom En (confer første Afsnit Capitel II) foredrager, at et Menneske ikke maa have nogen Discipel, og en Anden sætter sig ned som Tilhænger af denne Lære, er saa ikke Misforstaaelsen mellem dem trods alle Tilhængerens Forsikkringer om Beundring og om hvor ganske han har \tlg{tilegnet} sig – Misforstaaelsen? Kjendet paa Christendommen er Paradoxet, det absolute Paradox. Saasnart en saakaldet christelig Speculation hæver Paradoxet, og gjør denne Bestemmelse til et Moment, saa er alle Sphærer forvirrede.
\end{quote}
\dcite{AE:490  (SKS 7, 490)}
% AE:490  (SKS 7, 490)

\begin{quote}
Distinktionen mellem det Pathetiske og det Dialektiske maa dog nærmere bestemmes, thi Religieusiteten A er ingenlunde udialektisk, men den er ikke paradox dialektisk. Religieusiteten A er Inderliggjørelsens Dialektik; den er Forholdet til en evig Salighed ikke betinget ved et Noget, men er Forholdets dialektiske Inderliggjørelse, altsaa kun betinget ved Inderliggjørelsen, der er dialektisk. Religieusiteten B derimod, som den fremtidigen vil blive kaldet, eller den paradoxe Religieusitet, som den er bleven kaldet, eller den Religieusitet, der har det Dialektiske paa andet Sted, gjør Betingelser, saaledes at Betingelserne ikke ere Inderliggjørelsens dialektiske Fordybelser, men et bestemt Noget, der nærmere bestemmer den evige Salighed (medens i A Inderliggjørelsens nærmere Bestemmelse er den eneste nærmere Bestemmelse), ikke idet den nærmere bestemmer Individets \tlg{Tilegnelse} af den, men nærmere bestemmer den evige Salighed dog ikke som Opgave for Tænkning, men netop paradox som frastødende til ny Pathos.
\end{quote}
\dcite{AE:506  (SKS 7, 506)}
% AE:506  (SKS 7, 506)

\begin{quote}
Anmærkning. Dette er det Paradox-Religieuse, Troens Sphære. Det lader sig Altsammen troe – mod Forstand. Vil Nogen bilde sig ind, at han forstaaer det, saa kan han være sikker paa at han misforstaaer det. Den der forstaaer det ligefrem (i Modsætning til at forstaae, at det ikke kan forstaaes), han vil forvexle Christendom med en eller anden Hedenskabets Analogie (hvis Analogie er Bedragets til den faktiske Virkelighed); eller med den til Grund for alle Hedenskabets illusoriske Analogier (der ikke have Guds væsentlige Usynlighed som høiere dialektisk Mellembestemmelse, men bedrages af den æsthetisk-ligefremme Kjendelighed, confer andet Afsnit Capitel 2 Tillæg) liggende Mulighed. Eller han vil forvexle Christendommen med Noget, der dog er opkommet i Menneskets det vil sige Menneskehedens Hjerte, forvexle den med Menneske-Naturens Idee og glemme den qvalitative Forskjel, der accentuerer det absolut forskjellige Udgangspunkt: hvad der kommer fra Gud, og hvad der kommer fra Mennesket; han vil ved en Misforstaaelse istedetfor at bruge Analogien til fra den at bestemme Paradoxet (Christendommens Nyhed er ikke den ligefremme Nyhed, og netop derfor den paradoxe, cfr. det Foregaaende) omvendt tilbagekalde Paradoxet ved Hjælp af Analogien, der dog kun er Bedragets Analogie, hvis Anvendelse derfor er Analogiens Tilbagekaldelse, ikke Paradoxets. Han vil misforstaaende forstaae Christendommen som en Mulighed, og glemme at hvad der er muligt i Mulighedens Phantasie-Medium, muligt i Illusion, eller muligt i den rene Tænkens phantastiske Medium (og dette er det for al speculativ Tale om en evig Gudvorden Tilgrundliggende, at Scenen flyttes over i Mulighedens Medium), i Virkelighedens Medium maa blive det absolute Paradox. Han vil misforstaaende glemme, at det at forstaae kun gjelder i Forhold til Det, hvis Mulighed er høiere end dets Virkelighed, medens her lige omvendt Virkeligheden er det Høieste, det Paradoxe; thi Christendommen som Projekt er ikke vanskelig at forstaae, Vanskeligheden og Paradoxet er, at det er virkeligt. Derfor blev der i andet Afsnit Capitel 3 viist, at Troen er en ganske egen Sphære, der paradox fra det Æsthethiske og Metaphysiske accentuerer Virkelighed, og paradox fra det Ethiske en Andens Virkelighed, ikke den egne; derfor er en religieus Digter en tvivlsom Bestemmelse i Forhold til det Paradox-Religieuse, fordi, æsthetisk, Mulighed er høiere end Virkelighed, og det Digteriske netop ligger i Phantasie-Anskuelsens Idealitet, hvorfor man ikke sjeldent seer Psalmer, der, skjøndt rørende og barnlige og poetiske ved Phantasie-Anstrøget indtil Grændsen af det Phantastiske categorisk betragtede ikke ere christelige, categorisk betragtede ved det poetisk seet Deilige: Lyseblaae, og Ding Dang Klokke Klang befordre det Mythiske langt bedre end nogen Fritænker, thi Fritænkeren udsiger, at Christendommen er Mythe, den naiv-orthodoxe Digter afskyer dette og hævder Christendommens historiske Virkelighed – i phantastiske Vers. Den der forstaaer Paradoxet (i Betydning af at forstaae det ligefrem), han vil misforstaaende glemme, at hvad han eengang i Troens afgjørende Lidenskab greb som det absolute Paradox (ikke som det relative, thi saa var \tlg{Tilegnelsen} ikke Troen), altsaa som Det, der absolut ikke var hans egne Tanker, det kan aldrig blive hans Tanker (i ligefrem Forstand), uden at forvandle Troen til en Illusion, hvorved han da altsaa i et senere Moment kommer til at indsee, at det var en Skuffelse med at han absolut troede, det ikke var hans egne Tanker. I Troen kan han derimod meget godt vedblive at bevare sit Forhold til det absolut Paradoxe. Men indenfor Troens Sphære kan aldrig det Moment indtræde, at han forstaaer det Paradoxe (i ligefrem Betydning): thi skeer det, saa gaaer hele Troens Sphære ud som en Misforstaaelse. Virkeligheden, det vil sige at det og det virkeligen er skeet, er Troens Gjenstand og dog vel ikke noget Menneskes eller Menneskehedens egne Tanker, thi Tanken er da i det Høieste Muligheden, men Muligheden som Forstaaelse er netop den Forstaaelse, ved hvilken det Tilbageskridt gjøres, at Troen hører op. Den, der forstaaer Paradoxet, vil misforstaaende glemme, at Christendommen er det absolute Paradox (ligesom dens Nyhed den paradoxe Nyhed), netop fordi den tilintetgjør en Mulighed (Hedenskabets Analogier, en evig Gudvorden) som et Sandsebedrag, og gjør det til Virkelighed, og netop dette er det Paradoxe, ikke det Fremmede, Usædvanlige ligefrem (æsthetisk), men det tilsyneladende Bekjendte, og dog det absolute Fremmede, der netop som Virkelighed gjør Tilsyneladelsen til Bedraget. Den der forstaaer Paradoxet vil glemme at han ved at forstaae (Muligheden) er gaaet tilbage til det Gamle og gaaet glip af Christendommen. I Mulighedens Phantasie-Medium kan Gud meget godt for Indbildningen smelte sammen med Mennesket, men i Virkeligheden med det enkelte Menneske er netop Paradoxet.
\end{quote}
\dcite{AE:528  (SKS 7, 528)}
% AE:528  (SKS 7, 528)

\begin{quote}
Naar der tales til et Barn om Christendommen, og Barnet ikke voldeligen i overført Forstand mishandles: saa \tlg{tilegner} det sig alt det Milde, det Barnlige, det Elskelige, det Himmelske; det lever sammen med det lille Jesu-Barn, og med Englene, og med de hellige Tre-Konger, det seer Stjernen i den mørke Nat, det reiser den lange Vei, nu er det i Stalden, Forundring over Forundring, altid seer det Himlen aaben, med hele Phantasiens Inderlighed længes det efter disse Billeder – nu, og lad os saa ikke glemme Pebernødderne og al den anden Herlighed, der falder af i den Anledning; thi lad os fremfor Alt ikke blive gamle Labaner, der lyve om Barndommen, der tillyve sig selv dens Overspændthed og fralyve Barndommen dens Virkelighed. I Sandhed, det maatte være en Døgenicht, der ikke finder det Barnlige rørende, og yndigt, og livsaligt; en Humorist skulde man dog vel heller ikke ville mistænke for Miskjendelse af Barndommens Realitet, han Erindringernes ulykkelige-lykkelige Elsker. Men det er visseligen ogsaa en blind Veileder, der paa nogensomhelst Maade vil sige, at dette er den afgjørende Opfattelse af Christendommen, der blev Jøder en Forargelse og Græker en Daarskab. Christus bliver Gudebarnet, eller for det lidt ældre Barn den venlige Skikkelse med det milde Aasyn (den mythiske Commensurabilitet), ikke Paradoxet, hvem Ingen kunde see Noget paa (ligefrem forstaaet), end ikke Døberen Johannes (confer Joh. Ev. 1, 31. 33.), end ikke Disciplene, før de bleve gjorte opmærksomme (Joh. Ev. 1, 36. 42.); hvad Esaias har forudsagt i 53, 2. 3. 4, især i 4. Barne-Opfattelsen af Christus er væsentligen Phantasie-Anskuelsens, og Phantasie-Anskuelsens Idee er Commensurabiliteten, og Commensurabiliteten er væsentligen Hedenskab, den være forresten Magt, Herlighed, Skjønhed, eller indenfor en lidt humoristisk Modsigelse, der dog ikke er et virkeligt Skjul, men et let gjennemskueligt Incognito. Commensurabiliteten er den ligefremme Kjendelighed. Tjenerens Skikkelse er da Incognitoet, men det milde Aasyn er den ligefremme Kjendelighed. Som allevegne saa ogsaa her, der er en vis Orthodoxie, der naar den skal til at slaae rigtigt stort paa ved store Festligheder og afgjørende Leiligheder, bona fide tager lidt Hedenskab til Hjælp – og saa lykkes det allerbedst. En Præst holder sig maaskee jævnt hen til daglig Brug i de strenge og rigtige orthodoxe Bestemmelser, men hvad skeer, en Søndag skal han gribe sig an. For ret at vise hvor levende Christus staaer for ham, skal han aabne os et Indblik i sin Sjel. Nu er det rigtigt. Christus er Troens Gjenstand, men Tro er intet mindre end Phantasie-Anskuelse, og Phantasie-Anskuelse er just heller ikke noget Høiere end Troen. Nu gaaer det da løs: det milde Aasyn, den venlige Skikkelse, Sorgen i Øiet og saa videre Der er slet ikke noget comisk i, at en Mand lærer Hedenskab istedenfor Christendom, men der er noget comisk i, at en Orthodox, naar han ved de store Høitideligheder trækker alle Registere ud, tager feil og trækker Hedenskabets Skuffe ud, uden at mærke det. Dersom Organisten til daglig Brug vilde spille en Valts, vilde han vel blive afsat; men dersom en Organist, der ellers ganske rigtigt foredrog Psalme-Melodier, paa de store Høitider vilde i Betragtning af, at han ledsages af Basuner, spille en Valts – for ret at festligholde Dagen: saa vilde dette jo være comisk. Og dog finder man hos Orthodoxe lidt af dette sentimentale og kjælne Hedenskab, ikke til dagligt Brug, men netop ved de store Festligheder, naar de ret aabne deres Hjerter, og man finder det igjen gjerne i Talens sidste Deel. Den ligefremme Kjendelighed er Hedenskab; alle høitidelige Forsikkringer om, at det jo er Christus og at han er den sande Gud, hjælpe ikke, saasnart det dog ender med den ligefremme Kjendelighed. En mythologisk Figur er ligefrem kjendelig. Gjør man den Orthodoxe denne Indvending, saa bliver han rasende og farer op: ja, men Christus er jo den sande Gud, og altsaa dog vel ingen mythologisk Figur........ det kan man see paa hans blide Aasyn. Men kan man see det paa ham, saa er han eo ipso en mythologisk Figur. Man vil let see, at Pladsen bliver for Troen; thi tag den ligefremme Kjendelighed bort, saa er Troen paa rette Sted. Netop Forstandens og Phantasie-Anskuelsens Korsfæstelse, der ikke kan faae den ligefremme Kjendelighed, er netop Kjendet. Men det er lettere at skulke sig fra Forfærdelsen og skulke sig til lidt Hedenskab, der er gjort ukjendeligt ved det besynderlige Sammenhæng, at det nemlig tjener som sidste og høieste Forklaring i et Foredrag, der maaskee begyndte i ganske rigtige orthodoxe Bestemmelser. Dersom en Orthodox i et fortroligt Øieblik vilde betroe En, at han dog egentligen ikke havde Troen: nu ja, deri er Intet latterligt; men naar en Orthodox i saligt Sværmerie selv forundret næsten over sin Tales høie Sving, ganske aabner sig for En i Fortrolighed, og han er saa uheldig at tage Feil af Directionen, saa han stiger fra det Høiere til det Lavere, saa er det vanskeligere at lade være at smile.
\end{quote}
\dcite{AE:544  (SKS 7, 544)}
% AE:544  (SKS 7, 544)

\begin{quote}
Barne-Alderen (ligefrem forstaaet) er da ikke den sande i Forhold til det at blive Christen. Det er omvendt, den ældre Alder, Modenhedens Alder er Tiden, hvor det skal afgjøres, om et Menneske vil være det eller ikke. Barndommens Religieusitet er det universelle, abstrakte og dog phantasie-inderlige Grundlag for al senere Religieusitet, det at blive Christen er en Afgjørelse, der tilhører en langt senere Alder. Barnets Receptivitet er saa aldeles uden Afgjørelse, at man jo ogsaa siger: man kan bilde et Barn Alt ind. Det forstaaer sig, den Ældre bærer Ansvaret for hvad han tillader sig at bilde Barnet ind, men vist og sandt er det. At Barnet er døbt, kan da heller ikke gjøre det ældre i Forstand, eller modne det til Afgjørelser. Et Jødebarn, et Hedningebarn opdraget fra det Første af ømme christelige Pleieforældre, der behandle det lige saa kjerligt, som Forældre behandle deres eget: vil \tlg{tilegne} sig den samme Christendom som det døbte Barn.
\end{quote}
\dcite{AE:546  (SKS 7, 546)}
% AE:546  (SKS 7, 546)

\begin{quote}
2) Man siger: nei, ikke enhver Antagelse af den christne Lære gjør En til Christen. Det det især kommer an paa er \tlg{Tilegnelsen}, at man \tlg{tilegner} sig og fastholder denne Lære ganske anderledes end noget Andet, at man vil leve i den og døe i den, vove sit Liv for den og saa videre
\end{quote}
\dcite{AE:552.4  (SKS 7, 552)}
% AE:552.4  (SKS 7, 552)

\begin{quote}
Dette seer ud som det var Noget. Den Categorie: ganske anderledes er imidlertid en temmelig maadelig Categorie, og den hele Formel, der gjør et Forsøg paa at bestemme det at være Christen noget mere subjektivt, er hverken det Ene eller det Andet, undgaaer paa en Maade Vanskeligheden med Approximationens Distraction og Svig, men mangler categorisk Bestemmelse. Den \tlg{Tilegnelsens} Pathos, om hvilken her er Tale, er Umiddelbarhedens; man kan lige saa godt sige, at en begeistret Elsker forholder sig saaledes til sin Elskov: fastholder og \tlg{tilegner} sig den ganske anderledes end noget Andet, vil leve i den og døe i den, vove Alt for den. Forsaavidt er der altsaa ingen væsentlig Forskjel mellem en Elsker og en Christen i Henseende til Inderlighed, og man maa da igjen recurrere til det Hvad, som er Læren, og saa gaae vi igjen ind under Nummer 1.
\end{quote}
\dcite{AE:553  (SKS 7, 553)}
% AE:553  (SKS 7, 553)

\begin{quote}
Det gjelder nemlig om at selve \tlg{Tilegnelsens} Pathos i den Troende bestemmes saaledes, at den ikke kan forvexles med nogen anden Pathos. Deri har nemlig den mere subjektive Opfattelse Ret, at det er \tlg{Tilegnelsen}, der gjør Udslaget, men den har Uret i Bestemmelsen af \tlg{Tilegnelsen}, som ingen specifik Forskjel har fra al anden umiddelbar Pathos.
\end{quote}
\dcite{AE:553.2  (SKS 7, 553)}
% AE:553.2  (SKS 7, 553)

\begin{quote}
Dette skeer heller ikke, naar man bestemmer \tlg{Tilegnelsen} som Tro, men strax igjen giver Troen Farten og Retningen hen til Forstaaelsen, saa Troen bliver en midlertidig Funktion, hvorved man midlertidig fastholder hvad der skal blive Gjenstand for Forstaaelsen, en midlertidig Funktion, hvormed fattige Folk og dumme Mennesker maae nøies, medens Privat-Docenter og gode Hoveder gaae videre. Kjendet paa at være Christen (Troen) er \tlg{Tilegnelsen}, men saaledes, at den specifikt ikke er forskjellig fra anden intellectuel \tlg{Tilegnelse}, hvor en foreløbig Antagelse er en midlertidig Funktion i Forhold til Forstaaelsen. Tro bliver ikke det Specifike for Forholdet til Christendommen, og det vil atter blive det Hvad der troes, som afgjør, om En er Christen eller ikke. Men herved er Sagen igjen ført tilbage under Nummer 1.
\end{quote}
\dcite{AE:553.3  (SKS 7, 553)}
% AE:553.3  (SKS 7, 553)

\begin{quote}
Den \tlg{Tilegnelse} nemlig, ved hvilken en Christen er Christen, maa være saa specifik, at den ikke kan forvexles med nogen anden.
\end{quote}
\dcite{AE:553.4  (SKS 7, 553)}
% AE:553.4  (SKS 7, 553)

\begin{quote}
3) Man bestemmer det at blive og at være Christen objektivt ikke ved Lærens Hvad, ei heller subjektivt ved \tlg{Tilegnelsen}, ikke med hvad der er foregaaet i Individet, men med hvad der er foregaaet med Individet: at det er døbt. Forsaavidt man til Daaben føier Antagelsen af Troesbekjendelsen, saa vil herved dog intet Afgjørende vindes, men Bestemmelsen vil vakle mellem at accentuere Hvad (Aproximationens Vei) og ubestemt at tale om Antagelse og Antagelse, og \tlg{Tilegnelse} og saa videre uden specifik Bestemmethed.
\end{quote}
\dcite{AE:554  (SKS 7, 554)}
% AE:554  (SKS 7, 554)

\begin{quote}
Vil derimod Nogen sige, at han jo i Daaben modtog Aanden, og af dens Vidnen med hans Aand veed han at han er døbt: saa er Slutningen lige vendt om, han slutter fra Aandens Vidnesbyrd i ham til at han maa være døbt, ikke fra det at være døbt til at have Aanden. Men skal Slutningen gjøres saaledes, saa er ganske rigtigt Daaben ikke Kjendet, men Inderligheden, og saa maa igjen fordres den Inderlighedens og \tlg{Tilegnelsens} specifike Bestemmelse, hvorved Aandens Vidnesbyrd i en Christen er forskjellig fra al anden (en universellere bestemmet) Aandsvirksomhed i et Menneske.
\end{quote}
\dcite{AE:554.3  (SKS 7, 554)}
% AE:554.3  (SKS 7, 554)

\begin{quote}
Afgjørelsen ligger i Subjektet, \tlg{Tilegnelsen} er den paradoxe Inderlighed, der er specifik fra al anden Inderlighed. Det at være Christen bestemmes ikke ved Christendommens Hvad, men ved den Christnes Hvorledes. Dette Hvorledes kan kun passe til Eet, til det absolute Paradox. Der er derfor ingen ubestemt Talen om at det at være Christen er at antage og antage og ganske anderledes antage, at \tlg{tilegne} sig, at troe, at \tlg{tilegne} sig i Troen ganske anderledes (lutter rhetoriske og fingerede Bestemmelser); men det at troe er specifik bestemmet forskjelligt fra al anden \tlg{Tilegnelse} og Inderlighed. Tro er den objektive Uvished med det Absurdes Frastød, fastholdt i Inderlighedens Lidenskab, der netop er Inderlighedens Forhold potentseret til sit Høieste. Denne Formel passer kun paa den Troende, paa ingen Anden, ikke paa en Elsker, eller en Begeistret, eller en Tænker, men ene og alene paa den Troende, der forholder sig til det absolute Paradox.
\end{quote}
\dcite{AE:555  (SKS 7, 555)}
% AE:555  (SKS 7, 555)

\begin{quote}
Min kjære Læser! At jeg selv skal sige det: jeg er Intet mindre end en Satans Karl i Philosophien, kaldet til at skabe en ny Retning; jeg er et stakkels enkelt existerende Menneske med sunde, naturlige Evner, ikke uden en vis dialektisk Færdighed og heller ei ganske blottet for Studium. Men jeg er forsøgt i Livets casibus, og beraaber mig trøstigt paa mine Lidelser, ikke i apostolisk Forstand som en Æressag, thi de have kun altfor ofte været selvforskyldte Straffe, men dog beraaber jeg mig paa dem, som paa mine Læremestere, og med mere Pathos end naar Stygotius beraaber sig paa alle de Universiteter, ved hvilke han har studeret og disputeret. Jeg trodser paa en vis Oprigtighed, der forbyder mig at eftersnakke hvad jeg ikke kan forstaae, og byder mig, hvad der i Forhold til Hegel længe har voldet mig Smerte i min Forladthed, at forsage at beraabe mig paa ham uden i enkelte Partier, hvilket da er det Samme som at maatte give Afkald paa den Kjendelighed, man vinder ved Tilknytningen, medens jeg bliver, hvad jeg selv indrømmer er uendeligt Lidet, et forsvindende ukjendeligt Atom, som ethvert enkelt Menneske er det; en Oprigtighed, der da igjen trøster og væbner mig med en ualmindeligere Sands for det Comiske og en vis Evne til at gjøre latterligt hvad latterligt er; thi besynderligt nok, hvad der ikke er latterligt, kan jeg heller ikke gjøre latterligt, dertil hører formodentligen andre Evner. Jeg er, saaledes forstaaer jeg mig selv, netop saa meget udviklet ved Selvtænkning, saa meget dannet ved Læsning, saa meget orienteret i mig selv ved at existere, at jeg er istand til at være en Lærling, en Lærende, hvilket allerede er en Opgave. Mere udgiver jeg mig ikke for, end at være duelig til at kunne begynde i en høiere Forstand at lære. Dersom der da blandt Os fandtes Læremesteren! Jeg taler ikke om Læremesteren i Oldtids-Videnskaben, thi en saadan have vi, var det dette jeg skulde lære, da var jeg hjulpen, saasnart jeg havde erhvervet de fornødne Forkundskaber for at kunne begynde at lære; jeg taler ikke om Læremesteren i historisk Philosophie, hvortil jeg vel mangler Forkundskaberne, dersom vi ellers har Læremesteren; jeg taler ikke om Læremesteren i det religieuse Foredrags vanskelige Kunst, thi en saadan Udmærket have vi, og jeg veed, jeg har bestræbt mig for at benytte den alvorlige Veiledning, det veed jeg, om ikke af \tlg{Tilegnelsens} Udbytte, at jeg ikke underfundig skal tillyve mig selv Noget, eller maale hans Betydning med min Tilfældighed, saa veed jeg det af den Ærefrygt, jeg har bevaret for den Høiærværdige; jeg taler ikke om Læremesteren i Poesiens skjønne Kunst og dennes Hemmeligheder med Sproget og Smagen, thi en saadan Indviet have vi, det veed jeg, jeg haaber, jeg skal aldrig glemme hverken ham eller hvad jeg skylder ham. Nei, den Læremester, om hvilken jeg taler og paa en anden Maade, tvetydigt og tvivlende, er Læremesteren i den tvetydige Kunst at tænke over Existents og at existere. Altsaa, dersom han fandtes: saa tør jeg indestaae for, at der sgu nok skulde komme Noget ud af det, naar han paa Prent vilde tage sig af min Underviisning, og til den Ende gaae langsomt og stykkeviis frem, tilladende, som det bør ved en god Underviisning, at jeg spørger mig for, og at jeg forhindrer at Noget forlades førend jeg ganske har forstaaet. Dette kan jeg nemlig ikke antage, at en saadan Læremester skulde mene, at han ikke havde Andet at gjøre, end hvad en maadelig Religionslærer gjør i en Almueskole: sætte mig et vist §-Pensum for hver Dag, som jeg skulde kunne udenad til den næste.
\end{quote}
\dcite{AE:565  (SKS 7, 565)}
% AE:565  (SKS 7, 565)

\begin{quote}
Overhovedet kjender man strax den uendelige Reflexion, i hvilken først Subjektiviteten kan blive bekymret om sin evige Salighed, paa Eet: at den overalt har Dialektiken med sig. Det være et Ord, en Sætning, en Bog, en Mand, et Samfund, det være hvad det være vil, saasnart det saaledes skal være Grændse, at Grændsen selv ikke er dialektisk, er det Overtro og Indskrænkethed. Der er altid en saadan saavel magelig som ogsaa bekymret Trang i et Menneske til at faae noget rigtigt fast, der kan udelukke Dialektiken, men dette er Feighed og Svigagtighed mod Guddommen. Selv det Visseste af Alt: en Aabenbaring, bliver eo ipso dialektisk, idet jeg skal \tlg{tilegne} mig den; selv det fasteste af Alt, den uendelige negative Beslutning, der er Individualitetens uendelige Form for Guds Væren i ham, bliver strax dialektisk. Saasnart jeg tager det Dialektiske bort, er jeg overtroisk og bedrager Gud for Øieblikkets anstrængede Erhverven af det eengang Erhvervede. Derimod er det langt mageligere at være objektiv og overtroisk, og broutende deraf, og proklamerende Tankeløsheden.
\end{quote}
\dcite{AE:573.10  (SKS 7, 573)}
% AE:573.10  (SKS 7, 573)

\begin{quote}
Naar man siger, at deri er det Betryggende mod al Anfægtelse i Tanken om Daaben, at Gud i den gjør Noget ved os, saa er det naturligviis kun en Illusion, der ved en saadan Bestemmelse mener at holde Dialektiken borte, thi Dialektiken kommer strax med Inderliggjørelsen af denne Tanke, med \tlg{Tilegnelsen}. Derpaa har ethvert Genie, det største der nogensinde har levet, udelukkende at anvende al sin Kraft: paa Inderliggjørelsen i sig selv. Men Anfægtelsen ønsker man eengang for alle at være fri for, og i Anfægtelsens Øieblik retter Troen sig derfor ikke mod Gud, men Troen bliver Troen paa, at man virkelig er døbt. Hvis her nu ikke skjulte sig meget Spilfægterie, saa vilde der forlængst være fremkommet psychologisk mærkelige Tilfælde af Bekymring om med Vished at faae at vide, at man er døbt. Det skulde blot være 10,000 Rigsbankdaler det gjaldt om, saa vilde man neppe lade det staae hen med den Art af Vished, som vi Alle har om at være døbte.
\end{quote}
\dcite{AE:573.15  (SKS 7, 573)}
% AE:573.15  (SKS 7, 573)

\subsection{LA \normalfont(SKS 8)}

\begin{quote}
Man hører ikke sjeldent i vor Tid Ynglinge tale vise Ord; naar man hører nærmere til, hører man jo rigtignok, at det er Ynglinge; thi hvad der ikke undgaaer den Vise, at Alt, hvad han siger, angaaer ham selv (saa egoistisk er han!), det undgaaer Ynglingen (saa begeistret er han!). Hvad den Vise gjør, at han forstaaer Alt om sig selv og bestandigt forstaaer sig selv under, det gjør Ynglingen ikke, han gjør Fordringen og underforstaaer ikke sig selv, neppe noget Menneske – saadanne Fordringer gjør han. Man skulde troe, det var en beruset Guddom der talte, saaledes rutter han med Mennesker, saaledes svimler han i en phantastisk Forestilling om en Den, der skulde formaae at \tlg{tilegne} sig det Uhyre, der præsteredes ved Generationernes umenneskeligt forcerede Anstrængelse; mindst af Alt skulde man troe, det var et enkelt Menneske, der talte. Det forstaaer sig han taler rigtignok i Tidens Navn om Tidens Fordring; men det er jo netop Modsigelsen, der burde standse Ynglingens eventyrlige Grusomhed, thi saa grusom har da aldrig Nogen været, ingen Røver-Capitain i Romanen, som Tidens Fordring er det i en Ynglings Mund; saa grusom var end ikke Heliogabal – mod Strudsene, thi han lod dog ikke dræbe flere end han kunde spise. Knappere og knappere tilmaales Diæterne for de Udmærkede, hurtigere og hurtigere vækkes der Mistanke om Utroskab mod Tidens Fordring, nye og nye Kræfter fordrer Tiden i Ynglingens Mund. Og Den, der formodentligen ganske uopfordret gjør Fordringen, Den, der altsaa paa en eller anden Maade ved Utroskab mod sig selv og mod det at være Menneske, har \tlg{tilegnet} sig Fordringen, som fordrer Troskab, han denne uhyre Creditor, der slet ikke kan forklare, hvorledes han selv er kommen i Besiddelse af Fordringen, men blot gjør den gjeldende: han drømmer ikke om, at med den Dom, han dømmer, ja med en strængere Dom og et knappere Maal vil han blive dømt. Inderlighed til at ville lære af Livet finder man sjeldnere, men desto hyppigere Lyst og Tilbøielighed og gjensidig Fremskynden til at blive bedragen af Livet. Uforfærdet synes man ikke socratisk at frygte for at blive bedragen; thi Gudens Stemme er altid en Hvidsken og Tidens Fordring i et tusindtunget Rygtes Form er ikke et almægtigt Bliv der skaber store Mænd, men et Røre i Affaldet, der skaber Forvirrede, et Abracadabra, der frembringer, som al Frembringelse er det, i Lighed med sig selv. Endnu mindre synes man socratisk at frygte meest for at blive bedragen af sig selv, og allermindst synes man opmærksom paa, at hvis af alle Bedragne de ere de Ulykkeligste, der bedrage sig selv, saa er blandt disse de igjen de Ulykkeligste, der i Modsætning til de fromt Bedragne ere formasteligen bedragne af sig selv. Langsom til at høre, hurtig til at dømme, efterkommer man kun halvt det Socratiske: at slutte beskedent fra det Lidet, man forstaaer; thi man slutter overmodigt – fra det Lidet, man forstaaer. Det Øieblikkelige, en glimrende Begyndelse og en ny Tidsregnings Dateren fra denne er det Lidet man forstaaer – dersom det da ellers er muligt at forstaae det Øieblikkelige og forstaae Begyndelsen, da jo det Øieblikkelige mangler det Evige, og Begyndelsen mangler Slutningen førend Meningen er ude, og det burde man dog oppebie, ikke saa meget for Meningens som for Forstaaelsens Skyld. Det Næste altsaa; nu det behøves ikke, naar man taler i Tidens Navn, thi Tiden, saaledes abstrakt forstaaet, er en høist ligegyldig Magt, den kommer saa vist nok godt fra det, selv om ingen Yngling er bekymret paa dens Vegne. Men de enkelte Mennesker, der jo hver skulle have sin Tid? For de enkelte Mennesker kommer der jo et Næste, Eftersætningen efter den glimrende Forsætning – Ungdommens lykkelige Dage, da man selv var Tidens Fordring. Altsaa dette Næste; og jo knappere Tiden tilmaales de Ældre, desto hurtigere kommer jo dette Næste over En, og desto længere bliver Eftersætningen, og desto kortere bliver den smule Forsætning, som man har forstaaet, hvis den ikke bliver til slet Intet, da det bliver umuligt at forstaae den sammen med Eftersætningen. I Sandhed en glimrende Begyndelse med den nye Tidsregning! Det kan man kalde ikke at begynde med Intet men hellere at Begyndelsen endte med at blive til Intet. Ikke Pandoras Æske kunde indeholde saa mange Ulykker og saa megen Elendighed, som der er skjult i det lille Ord: Tidens Fordring. Hver Den, der har leflet med dette Ord, han tage Skade for Hjemgjeld; han kan jo ikke klage over Tidens Fordring, naar han selv har udfordret den! Og dog vil man ikke forstaae det Næste. Bekymret for sig selv at lade sig veilede, i Inderlighed og med Beundring at glædes ved den Ældre, naar han bliver sig selv tro, at opbygges ved de 50 Aars trofaste Tjeneste, at forstaae langsomt, at lære af den Høiærværdige, af hvem man lærer noget Andet end af Øieblikkets Eminentser: til saaledes at lære det Menneskelige og at forsage det Umenneskelige, dertil synes man ikke at have Taalmodighed. Men skulde det ogsaa kunne være Tidens Fordring? Ja, er det ikke Tidens Fordring, saa er det dog den Magts, der er høiere end Tidens, thi det er Evighedens; er det ikke Menneskenes, saa er det Gudens; er det ikke Ynglingens, saa er det den Gamle af Dages! Og er det vel Troskab kun umenneskeligt at fordre Troskab; og er det vel Troskab at love meget Mere end noget Menneske kan holde; og ligger det Comiske ikke stundom nær nok, naar Striden er mellem Tvende, af hvilke den Ene kun umenneskeligt fordrede Troskab uden at forpligte sig selv og derved lære hvad menneskelig Troskab er, den Anden forpligtede sig selv saaledes, at Troskab blev umulig, og Utroskab er de Stridendes fælles Lighed i Forandringen!
\end{quote}
\dcite{LA:14  (SKS 8, 14)}
% LA:14  (SKS 8, 14)

\begin{quote}
Tro har Forfatteren været mod sig selv; vilde Nogen daarligen sige, at han har været tro over Lidet, fordi han ikke har havt travlt med at forandre Opgaven, fordi han ikke har vanket om i alle Regioner, fordi han ikke har gøglet med en prunkende Mangfoldighed: da vilde jeg til denne vistnok fingerede Indvending svare: deri ligger netop Forfatterens store Rigdom. Livs-Anskuelsen, der skaberisk bærer disse Fortællinger, er den samme, medens en sindrig Opfindsomhed, en riig Erfarings erhvervede Forraad, en frugtbar Stemnings vegetative Frodighed tjenende frembringe Forandringen indenfor den skaberiske Gjentagelse. Uroen er væsentligen den samme, Beroligelsen væsentligen den samme, der kommes i dem alle væsentligen fra det Samme til det Samme; Spliden, der sættes, er i Spændstighed væsentligen den samme, Freden og Hvilen igjen den samme det vil sige Livs-Anskuelsen er den samme. Vil man tage en af disse Fortællinger i sin Totalitet som given, og nu betragte de andre, og vil man, som man jo gjør, kalde Digterens Frembringelse en Skabelse: saa vilde jeg helst sammenligne denne hele Række af Frembringelser i Forhold til hiin ene med: Opholdelsen. Men er Gud, med hvem Digteren sammenlignes naar han kaldes skabende, er han mindre beundringsværdig i Opholdelsen end i Skabelsen? Og saaledes ogsaa med hiin Forfatters fortsatte Skabelse. Ja dersom det ikke var saaledes med denne Forfatter, vilde hans egne Frembringelser indeholde en Modsigelse mod sig selv. Livs-Anskuelsen, der er indeholdt i Fortællingen, maa atter indeholdes i Frembringelsen, ellers viser dette Misforhold netop, at en Forfatter ingen Livs-Anskuelse har, trods al sin formeentlige Skaben; hvorimod denne Forfatter i Inderlighedens Troskab reproducerer sin egen Oprindelighed i Gjentagelsen. Vilde man med en psychologisk Gradestok ligesom gradere Kraften, saa vil man finde denne væsentligen at være den samme i alle; vilde man psychologisk ligesom lokke Lidenskabens Fart, vil man finde denne væsentligen at være den samme i alle. I alle Fortællinger den samme Nærhed ved det daglige Livs Virkelighed, og den samme Fjernhed i Opløftelse; den samme Nærhed ved Striden, den samme Fjernhed i Forstaaelsen. Saaledes er denne Forfatter bestandig Virkeligheden lige nær og lige fjern, han er den i Opfattelsen nærmere end en Digter og i den forklarede Gjengivelse mindre fjern end en Digter; thi Forfatterens Livs-Anskuelse ligger i Confiniet mellem det Æsthetiske hen til det Religieuse. Der, hvor Poesien væsentligen hører op, begynder denne Forfatter. Thi Poesien forsoner ikke væsentligen med Virkeligheden, den forsoner gjennem Phantasien med Phantasiens Idealitet, men denne Forsoning er netop i det virkelige Individ den nye Splid med Virkeligheden. Et Kjendskab til Virkelighedens Smerte, som det ikke kan komme frem i Poesien, (fordi den mindre intensiv end Fortvivlelsen har Tidens og Virkelighedens Dialektik) finder sin Forsoning i en Beroligelse, som Poesien (stricte sic dicta) ikke kan \tlg{tilegne} sig, netop fordi det er en Forsoning med Virkelighed. Men i Fortællingerne opgiver Forfatteren sig aldrig selv en saadan Virkelighedens Smerte, at den kun kunde finde sin Beroligelse i afgjort religieuse Bestemmelser og i det Religieuses Idealitet. Enhver Livs-Anskuelse veed Udveien, og er kjendelig paa, hvilken Udvei den veed. Digteren veed Phantasiens Udvei; denne Forfatter veed Virkelighedens; den Religieuse veed det Religieuses. Livs-Anskuelsen er Udveien, og Fortællingen er Veien. Paa Kategoriernes Forhold seer man let, at Ingen kan saaledes forsone med Virkeligheden som denne Forfatter, netop fordi hans Udvei er Virkeligheden. Frembringelsens Mesterskab ligger i, at dette intetsteds læres, men bestandigt skeer, bliver til, at der i Begivenheds og Personligheds Gjennemsigtighed vindes den Ro, i hvilken Fortællingen ender. – At forøvrigt denne Forfatter maatte udmærke sig ved Troskab mod sig selv, er ikke psychologisk vanskeligt at forklare, uden at dog denne ophøiede Lethed gjør Troskaben mindre beundringsværdig. Hverdags-Historien (eller Hverdags-Historierne) viser sig let som en anden Modenheds Frugt, og netop derfor kan der ikke være Tale om de Kriser og afgjørende Forandringer, hvilke en Forfatter, der begynder med en første Modenhed, maa opleve. Den Livs-Anskuelse, der bærer Hverdags-Historien (hiin første som har givet dem alle Navn), maa være modnet i Forfatteren førend Frembringelsen. Hans Frembringelse er ikke Moment i hans egen Udvikling, men idet denne er modnet, afsætter den som sin Frugt en Frembringelse i Inderlighed. Det er ikke det Geniale, ikke det Talentfulde, ikke Virtuositeten, der constituerer Frembringelsen, saa Produktiviteten væsentligen vilde forsvinde med Dettes Forsvinden, nei Frembringelsen selv, Muligheden af at kunne frembringe noget Saadant er snarere Lønnen, som Guden skjenkede Forfatteren, da han anden Gang modnet vandt i en Livs-Anskuelse et Evigt. Derfor kan han være veiledende, fordi han ikke er en Forfatter, der søger sig selv, men En, der har fundet sig selv, førend han blev Forfatter. Og hvad enten man nu kan finde Hvile i den samme Livs-Anskuelse eller ikke: her er dog Fred, og en stille Aands Uforkrænkelighed; saa lidet som her er et Sønderrivelsens sørgelige Budskab (og Gud veed hvor Nogen kan mene, at dette skal forkyndes!) saa lidet er her nogen Vigtighed, der spænder en nysgjerrig Forventning efter, hvad det nu er, der skal tilfredsstille Tidens Fordring, hvad den urolige Forfatter nu behageligst antager at være det Lyksaliggjørende. Nei, Hverdags-Historien er ikke blot en fuldendt Fortælling, men er Fuldendelsens Noveller. Glad gaaer man til Læsningen, man veed forud, at der er en Livs-Anskuelses Tilforladelighed, man er ikke uroliget af en vildfarende Stjernes Omtumlinger, men under Beskyttelse af en ledende Stjerne. Forud kjender man Livs-Anskuelsen, frydes nu ved, at den atter bliver til, forandret, men i Gjentagelsen, en Forandring, der ikke nysgjerrigt er for Andre ved Maskerade-Dragternes Foranledning til at nappe Nyt, men er for Inderligheden. Snart i tyve Aar er dette skeet saaledes, medens det udvortes Væsen ogsaa er det samme; den samme næsten qvindelige Resignation, der dog indgyder Ærbødighed, den samme Beskedenhed i Fremtræden, den samme Afsideshed fra Allarmen og Tidens Fordringer, den samme Hjemlighed og trofaste Vedhængen ved en dansk Læseverden; thi ogsaa dette er saa skjønt, at Hverdags-Historien saa godt som ikke er oversat i noget fremmed Sprog: en Nøisomhedens og Selvbevidsthedens Glæde over det lille Danmark! – O, skjønt saaledes at være Forfatter, skjønt at være den Navnkundige, der som Udgiver har lagt sit berømte Navn til; som Fortællingerne ere det, saa er ogsaa dette Forhold en Forening i Troskab. Udgiveren er ikke bleven berømt ved dem, thi han var det, men Fortællingerne, der selv have erhvervet sig et Navn, næst Udgiverens et af de betydningsfuldeste, vedblive dog i Navnløshedens Hengivelse at søge Tilhold hos den Navnkundige.
\end{quote}
\dcite{LA:18  (SKS 8, 18)}
% LA:18  (SKS 8, 18)

\begin{quote}
Vil man sige om Digteren (og her er ikke Tale om comparative Bestemmelser mellem de enkelte Digtere, men om Qvalitets-Prædikatet for det væsentligen at være Digter), at han river hen, at han begeistrer, fordi han væsentligen virker ved Phantasien, saa vilde jeg sige om Hverdags-Historiens Forfatter: han overtaler; thi at han fængsler ved Læsningen, underholder under Læsningen, er ikke saa charakteristiske Bestemmelser, da det Samme jo kan passe paa Digteren og paa mangen anden Novellist. Han overtaler, og ogsaa dette er en vanskelig, men tillige en velgjørende Kunst. Vil man sige, at Digteren (stricte sic dictus) kun har de høie eller de dybe Toner, og derfor egentligen ikke kan tale med Folk, saadanne som disse ere i Virkeligheden: da gjelder det om Hverdags-Historien, at hans Livs-Anskuelse netop har Overtalelsens Mellemtone, og Livs-Anskuelsens Udtalelse netop sin Fuldendthed heri. Det er ikke stærke, begeistrede og ophidsede Tilraab til at stride, til at nyde, til at trodse Livets Modstand, eller Fortvivlelsens Skrig, som her høres, eller saadanne Optrin, som her skildres, eller i saadan Afgjørelse, her endes. Det er det hyggelige Kabinets indbydende Fortrolighed, der her aabner sin Helligdom for En, men fra hvilken igjen Lidenskabens Heftighed og Afgjørelsens Livsfare og Anstrængelsens Yderste er udelukket, fordi Sligt hverken kan rummes eller kan taales der. At ville lade sig sige er Betingelsen for at \tlg{tilegne} sig Overtalelsens Veiledning; denne nedspændte Eftergivenhed Betingelsen for, at Overtalelsen atter kan aflokke Forstemtheden ny Velklang. Det bliver Alt nok godt igjen. Ved hvilke Midler? Ja spørg sømmeligen; thi om Du end spørger rigtigt, vil Du, hvis Du spørger heftigt have gjort Dig det umuligt at \tlg{tilegne} Dig Besvarelsen. Ved hvilke Midler? Ved Forstandighed til at afvinde Lidelsen en mildere Side, ved Taalmodighed, der venter et Smiil af Lykken igjen, ved venlig Deeltagelse af kjærlige Mennesker, ved Resignationen, der giver Afkald ikke paa Alt, men paa det Høieste, og ved Nøisomhed forvandler det Næste derefter til næsten lige saa godt som det Høieste. Og Alt dette foredrages ikke, det skeer: netop derfor er Overtalelsen saa stor, naar man hengiver sig i den. Ingen Taler kan overtale saaledes, netop fordi han har Hensigten, og Betragtningen altid føder Tvivlen. Overtalelsen her derimod er ikke et Mellemværende mellem tvende Personer, men er Veien i Livs-Anskuelsen, og Novellen fører En ind i den Verden, som Anskuelsen skaberisk bærer. Men see, denne Verden er igjen netop Virkeligheden; Anskuelsen har derfor ikke bedraget En, men netop overtalt En til at blive, hvor man er.
\end{quote}
\dcite{LA:22  (SKS 8, 22)}
% LA:22  (SKS 8, 22)

\begin{quote}
Publikum er ikke et Folk, ikke en Generation, ikke en Samtid, ikke en Menighed, ikke et Selskab, ikke disse bestemte Mennesker, thi alt Sligt er kun ved Concretionen det som det er; ja ikke en eneste af dem, der høre til Publikum, har noget væsentligt Engagement; i nogle Timer af Dagen hører han maaskee med til Publikum, nemlig i de Timer, i hvilke han Ingenting er, thi i de Timer, i hvilke han er det Bestemte han er, hører han ikke til Publikum. Dannet af saadanne Enere, af de Enkelte i de Øieblikke, hvor de Ingenting ere, er Publikum noget uhyre Noget, det abstrakte Øde og Tomme, som er Alle og Ingen. Men af den samme Grund kan Enhver anmasse sig at have et Publikum, og ligesom den romerske Kirke chimairisk udvidede sig ved at udnævne Biskopper in partibus infidelium: saaledes er Publikum Noget, som Enhver, selv en fuld Matros, der foreviser »en Perspektivkasse«, kan \tlg{tilegne} sig, og den fulde Matros har dialektisk consequent absolut samme Ret dertil som den meest Udmærkede, absolut Ret til at sætte alle disse mange, mange Nuller foran sit Eettal. Publikum er Alt og Intet, er den farligste af alle Magter og den meest intetsigende; man kan tale til en heel Nation i Publikums Navn, og dog er Publikum mindre end et eneste nok saa ringe virkeligt Menneske. Bestemmelsen Publikum er det Reflexionens Blendværk, der gøglende har gjort Individerne indbildske, fordi Enhver kan anmasse sig dette Uhyre, i Sammenligning med hvilket Virkelighedens Concretioner synes fattige, Publikum er Forstands-Tidens Eventyr, der gjør de Enkelte phantastisk til mere end at være Konge over et Folk; men Publikum er atter den grusomme Abstraktion, ved hvilken Individerne skal religieust opdrages – eller gaae under.
\end{quote}
\dcite{LA:89  (SKS 8, 89)}
% LA:89  (SKS 8, 89)

\subsection{OTA \normalfont(SKS 8)}

\begin{quote}
Skjøndt denne lille Bog (Leilighedstale kan den kaldes, uden dog at have Leiligheden, som gjør Taleren og gjør ham myndig, eller Leiligheden, der gjør Læseren og gjør ham lærende) i Virkelighedens Forhold er som en Indbildning og liig en Drøm om Dagen: er den dog ikke uden Tillid og ikke uden Haab om Opfyldelsen. Den søger hiin Enkelte, til hvem den ganske giver sig hen, af hvem den ønsker at modtages som var den opkommen i hans eget Hjerte; hiin Enkelte, hvem jeg med Glæde og Taknemlighed kalder min Læser; hiin Enkelte, der i Villighed læser langsomt, læser gjentagent, og som læser høit – for sin egen Skyld. Finder den ham, da er Forstaaelsen fuldkommen i Adskillelsens Fjernhed, naar han beholder den og Forstaaelsen for sig selv i \tlg{Tilegnelsens} Inderlighed.
\end{quote}
\dcite{OTA:121.2  (SKS 8, 121)}
% OTA:121.2  (SKS 8, 121)

\begin{quote}
Lever Du nu saaledes min Tilhører, at Du er Dig tydeligt og evigt bevidst at være en Enkelt? Dette var Talens Spørgsmaal, eller derom spørger Du Dig selv, hvis Du ret levende tænker dens Anledning. Talen skal ikke kunne sige Dig, hvad da ogsaa kun virker forstyrrende, om Mange eller ikke have den Overbeviisning, at et Menneske bør leve i denne Bevidsthed. Men dette gjør jo hverken fra eller til, hvad enten Alle, hver især, have denne Overbeviisning eller Ingen. Den Talende skal heller ikke søge at vinde Dig for denne Overbeviisning, dersom den ellers er hans, – saaledes trænger han ikke til Dig, saa lidet som Du trænger til ham, hvis Du har denne Overbeviisning: det Eviges ophøiede Alvor vil hverken Stemmefleerhedens eller Veltalenhedens Anbefaling. Kun Eet tør Talen ikke love Dig – den ønsker jo heller ei at fornærme; jordisk Fordeel tør den ikke love Dig hvis Du antager og i \tlg{Tilegnelse} fastholder denne Overbeviisning. Tvertimod, den vil, naar den skal fastholdes, gjøre Dig Livet anstrænget, mangen Gang maaskee møisomt; den vil, naar den skal fastholdes, maaskee udsætte Dig for Andres Spot, at jeg ikke skal tale om, hvis Fastholdelsen skulde ville kræve endnu større Offre af Dig. Det forstaaer sig, forstyrre Dig gjør Spotten ikke, dersom Du ellers holder Overbeviisningen fast; Spotten vil endog tjene Dig til det Gode, ogsaa saaledes, at den yderligere forvisser Dig om, at Du er paa den rette Vei. Thi Mængdens Dom har sin Betydning; man skal ikke hovmodigt forblive uvidende om den, nei man skal være opmærksom; naar man saa blot passer paa at gjøre det Modsatte, da træffer man gjerne det Rigtige; eller hvis man oprindeligen gjør det Modsatte, og det da træffer sig saa heldigt, at Mængdens Dom yttrer sig derimod, saa kan man være temmelig sikker paa, at man har grebet det Rette: man har da ikke blot ved sig selv ret inderligen overveiet og prøvet sin Overbeviisning, men man har tillige den Fordeel, endnu engang at faae gjort Prøve paa den – ved Hjælp af Spotten, der muligen saarer, men netop ved Saaret viser, at man er paa den rette Vei, paa Ærens og Seirens, som Krigerens Saar, naar det er paa Brystet, hvor baade Saaret og Hæderstegnet skal bæres. Du har vel lagt Mærke til, at blandt Skoledrengene indbyrdes ansees Den for den Raskeste, der ikke er bange for sin Fader, der tør sige til de Andre: »troer I, at jeg er bange for ham.« Mærke de derimod paa en af deres Jævnlige, at han ganske egentligen og bogstaveligen forstaaet er bange for sin Fader: saa ville de gjerne spotte lidt over ham. Ak, i Menneskenes bange sammenløbne Mængde (thi hvorfor løber man vel sammen i Mængde uden fordi man er bange!) der er det ogsaa rask ikke at være bange, end ikke for Gud; og mærker man, at der er en Enkelt udenfor, som ganske egentligen og bogstaveligen forstaaet er bange – ikke for Mængden, men for Gud: saa vil man gjerne spotte lidt over ham. Man besmykker Spotten og siger: et Menneske skal elske Gud. Ja vist, Gud veed, det er den høieste Trøst for et Menneske, at Gud er Kjærlighed, og at han tør elske Gud; men lad os ikke blive fornemme, og daarligen, ja bespotteligen afskaffe det fra Fædrene Overleverede, af Gud selv Opfundne: at man ganske egentligen og bogstaveligen forstaaet skal være bange for Gud. Og det er Den, som er sig bevidst at være en Enkelt, og derved sig sit evige Ansvar bevidst for Gud, thi han veed, at kunde han end ved Udflugternes og Undskyldningernes Hjælp komme godt gjennem Livet, vandt han end ad denne løvtagte Vei den hele Verden: saa er der dog et Sted hisset, hvor der ikke er mere Udflugt, end der er Skygge i den brændende Ørken. – Videre skal Talen ikke sige herom, den spørger Dig kun gjentagent, lever Du nu saaledes, at Du er Dig bevidst at være en Enkelt, og derved Dig Dit evige Ansvar for Gud bevidst; lever Du saaledes, at denne Bevidsthed kan faae Tid og Stilhed og Raaderum til at gjennemtrænge Dine Livsforhold? Det forlanges ikke, at Du skal trække Dig tilbage fra Livet, fra en hæderlig Gjerning, fra et lykkeligt Husliv; tvertimod skal jo netop hiin Bevidsthed bære og forklare og gjennemlyse Din Færd i Livets Forhold. Du skal ikke trække Dig tilbage og sidde grublende over Dit evige Regnskab, hvorved Du kun forskylder et nyt Ansvar; Dine Pligter og Dine Opgaver vil Du bedre og bedre faae Tid til, medens Bekymringen for Dit evige Ansvar vil forhindre Dig i, at være travl og travlt deeltagende i alt Muligt – en Virksomhed der bedst kan kaldes Tidsspilde.
\end{quote}
\dcite{OTA:235  (SKS 8, 235)}
% OTA:235  (SKS 8, 235)

\begin{quote}
Skjøndt denne lille Bog er uden Lærerens Myndighed, til Overflod, ubetydelig, som Lilien og Fuglen – o, at den var saaledes! – haaber den dog, ved at finde det Eneste den søger: et godt Sted, at finde \tlg{Tilegnelsens} Betydning for hiin Enkelte, hvem jeg med Glæde og Taknemmelighed kalder min Læser.
\end{quote}
\dcite{OTA:257.2  (SKS 8, 257)}
% OTA:257.2  (SKS 8, 257)

\begin{quote}
Thi at finde Hvile det er at dannes for det Evige. Hovedsageligen er der kun Eet, hvori der er Hvile at finde: at lade Gud raade i Alt; hvad et Menneske mere faaer at vide, det angaaer, hvorledes Gud har villet raade. At der er en Forsoning for den Sønderknusede, deri er der Hvile, men han kan ikke finde Hvile i denne evige Tanke, dersom han ikke først hviler i Lydighedens: at Gud skal raade i Alt, thi Forsoningen er jo netop Guds Raad til Menneskets Frelse. At der er gjort Fyldest for Skylden, deri er der Hvile for den Angrende, men han kan ikke finde Hvile i denne evige Tanke, dersom han ikke først hviler i den: at lade Gud raade i Alt, thi Fyldestgjørelsen er jo netop Guds Raad fra Evighed. At Gud vil tage Dig til Naade, deri er der Hvile, men Du kan ikke finde Hvile i denne evige Tanke, hvis Du ikke hviler i den: at Gud skal raade i Alt. Ellers blev jo Guds Naade Din Fortjeneste, istedenfor, at det er Gud, som giver baade at ville og at udrette, giver Væxten og giver Fuldkommelsen, hvad al Din egne Stræben ikke formaaer: den lægger ingen Alen til Din Væxt, ikke en Tomme; hvad al Din egne Selvbekymring ikke formaaer: den bedrøver kun Aanden og sinker Væxten. Men Troen og Troens Lydighed i Lidelser elsker Væxten op, fordi al Troens Arbeiden gaaer ud paa at skaffe det Egne, det Selviske bort, for at Gud ret kan komme til og for da at lade ham raade i Alt. Jo mere der lides, naar der tillige læres af Lidelsen, jo mere tages alt det Selviske bort, det udryddes, og Lydigheden træder isteden som den modtagelige Jordbund, hvor det Evige kan slaae Rødder. Du kan ikke tage det Evige, Du kan kun \tlg{tilegne} Dig det; men Du kan ikke \tlg{tilegne} Dig hvad der er Dit Eget, kun hvad der er en Andens; og dette kan Du atter ikke paa tilladelig Maade \tlg{tilegne} Dig, hvis han ikke vil give Dig det; vil han derimod give Dig det, da er \tlg{Tilegnelsen} Inderliggjørelsen; men i Forhold til Gud og det Evige er \tlg{Tilegnelsen} Lydighed, og i Lydigheden er der Hvile. Der er Hvile i det Evige, dette er den evige Sandhed, men det Evige kan kun hvile i Lydigheden, dette er den evige Sandhed for Dig.
\end{quote}
\dcite{OTA:357  (SKS 8, 357)}
% OTA:357  (SKS 8, 357)

\begin{quote}
Dersom vi tænke os en Yngling, der er vel underviist i Sandhed, da kunne vi jo ingenlunde negte, at han veed det Sande, og dog vil det vel gaae ham, som det gik Andre før ham, at han, naar han er bleven Ældre, har faaet ganske Andet at vide, medens han dog endnu kun veed det Sande. Ynglingen veed vistnok det Sande, men han er uden Kjendskab til, uden Erfaring om de Virkelighedens Forhold, den hele Omgivelse, indenfor hvilken Sandheden skal træde frem. Sjelden er det, at et Menneske i denne Henseende fra sin tidligste Tid er hjulpen ved en mørkere Anelse; som oftest har Ynglingen, hvad der er elskeligt, en Godtroenhed, der rigtignok stundom igjen bliver en Ynglings Fordærvelse. Ynglingens Lærebegjerlighed \tlg{tilegner} sig saa gjerne og saa villigt Sandheden, som den meddeles ham; hans uerfarne men skjønne Indbildning danner ham da et Billede, som han kalder Verden, hvor nu, hvad han har lært, udfolder sig for ham som paa en Skueplads. Det passer i Ynglingens ufordærvede Tanke saa ganske for hinanden: Sandheden, som han har lært den i dens reneste Skikkelse, og Verden, det er den Indbildningens Skueplads, som han selv danner. Saaledes maa Forholdet være mellem Sandheden og Verden, tænker Ynglingen, og saa gaaer han tillidsfuld ud i Virkelighedens Verden.
\end{quote}
\dcite{OTA:420  (SKS 8, 420)}
% OTA:420  (SKS 8, 420)

\subsection{KG \normalfont(SKS 9)}

\begin{quote}
Ja, han var Kjerlighed, og hans Kjerlighed var Lovens Fylde. »Ingen kunde overbevise ham om nogen Synd«, end ikke Loven, der veed Alt med Samvittigheden; »der var og ikke Sviig i hans Mund«, men Alt i ham var Sandhed; der var i hans Kjerlighed ikke et Øiebliks, ikke en Følelses, ikke et Forsæts Afstand fra Lovens Fordring til dens Opfyldelse; han sagde ikke som hiin ene Broder nei, ei heller ja som den anden Broder, thi hans Spise var at gjøre Faderens Villie: saaledes var han Eet med Faderen, Eet med hver een Fordring i Loven, saa det at fuldkomme den var Trangen i ham, hans eneste Livsfornødenhed. – Kjerligheden i ham var idel Handlen; der var intet Øieblik, ikke eet eneste i hans Liv, hvor Kjerligheden i ham kun var en Følelses Uvirksomhed, der søger efter Ord, medens den lader Tiden gaae, eller en Stemning, der er sig selv en Tilfredsstillelse, dvæler ved sig selv, medens der ingen Opgave er, nei hans Kjerlighed var idel Handlen; selv naar han græd, udfyldte dette ikke Tiden, thi vidste Jerusalem end ikke, hvad der tjente til dens Fred, han vidste det, vidste de ved Lazari Grav Sørgende ikke hvad der skulde skee, han vidste dog hvad han vilde gjøre. – Hans Kjerlighed var heel tilstede i det Mindste som i det Største, den samlede sig ikke stærkere i enkelte store Øieblikke, som vare det daglige Livs Timer uden Lovens Fordringer; den var ligelig i ethvert Øieblik, ikke større da han udaandede paa Korset end da han lod sig føde; det var den samme Kjerlighed der sagde »Maria har valgt den bedre Deel«, og den samme Kjerlighed, der med et Blik straffede – eller tilgav Petrus; det var den samme Kjerlighed, da han modtog Disciplene, der glade vendte hjem fra i hans Navn at have gjort Undergjerninger, og den samme Kjerlighed da han fandt dem sovende. – Der var i hans Kjerlighed ingen Fordring til noget andet Menneske, noget andet Menneskes Tid, Kraft, Bistand, Tjeneste, Gjenkjerlighed, thi hvad Christus fordrede af ham var ene det andet Menneskes Gavn, og han fordrede det alene for det andet Menneskes Skyld; intet Menneske levede med ham, der elskede sig selv saa høit som Christus elskede ham. Der var i hans Kjerlighed ingen tingende, eftergivende, partisk Overeenskomst med noget Menneske udenfor den Overeenskomst, som i ham var med Lovens uendelige Fordring; der var i Christi Kjerlighed ingen Fritagelse for ham fordret, ikke den ringeste, ikke en Hviids. – Hans Kjerlighed gjorde ingen Forskjel, ikke den ømmeste mellem Moderen og andre Mennesker, thi han pegede paa sine Disciple og sagde »disse er min Moder;« og atter Discipelens Forskjel gjorde hans Kjerlighed ikke, thi hans eneste Ønske var, at Enhver vilde blive hans Discipel, og dette ønskede han for Enhvers egen Skyld; og atter mellem Disciplene gjorde hans Kjerlighed ingen Forskjel, thi hans guddommelige-menneskelige Kjerlighed var netop den ligelige til alle Mennesker, at ville frelse dem alle, og den ligelige til alle dem, der vilde lade sig frelse. – Hans Liv var idel Kjerlighed og dog var hele dette hans Liv kun een eneste Arbeidsdag, han hvilede ikke før den Nat kom, da han ikke mere kunde arbeide; før den Tid vexlede hans Arbeide ikke med Dagens og Nattens Skifte, thi arbeidede han ikke, saa vaagede han i Bønnen. – Saaledes var han Lovens Fylde. Og som Løn derfor fordrede han Intet, thi hans eneste Fordring, hans eneste Hensigt med hele sit Liv fra Fødselen til Døden var uskyldigt at offre sig selv – hvad end ikke Loven, naar den fordrer Sit til det Yderste, turde fordre. – Saaledes var han Lovens Fylde; han havde ligesom kun een Medvider, der nogenlunde var istand til at følge ham, en Medvider, som var opmærksom og søvnløs nok til at forske efter, det var Loven selv, der fulgte ham Skridt for Skridt, Time for Time med sin uendelige Fordring; men han var Lovens Fylde. – Hvor fattigt aldrig at have elsket, o, men selv det Menneske, som blev rigest ved sin Kjerlighed, hvor er dog hans hele Rigdom kun Armod mod denne Fylde! Og dog, ikke saaledes, lader os aldrig glemme, at der er en evig Forskjel mellem Christus og enhver Christen; er Loven end afskaffet, her staaer den endnu ved Magt, og befæster et evigt svælgende Dyb mellem Gud-Mennesket og ethvert andet Menneske, der end ikke kan fatte, men kun troe, hvad den guddommelige Lov maatte indrømme, at han var Lovens Fylde. Enhver Christen troer det og \tlg{tilegner} sig det troende, men Ingen har vidst det, uden Loven og Han som var Lovens Fylde. Thi at hvad der i et Menneske svagt nok er tilstede i hans stærkeste Øieblik, at det langt stærkere og dog ligeligt var tilstede i ethvert Øieblik, det kan et Menneske ikkun forstaae i det stærkeste Øieblik, men Øieblikket efter kan han ikke forstaae det, og derfor maa han troe og holde sig til Troen, at hans Liv ikke skal blive forvirret ved i eet Øieblik at forstaae og i mange andre Øieblikke ikke at forstaae.
\end{quote}
\dcite{KG:105  (SKS 9, 105)}
% KG:105  (SKS 9, 105)

\begin{quote}
Dersom man med eet eneste Ord skulde angive og betegne den Seier, Christendommen har vundet over Verden, eller endnu rigtigere den Seier, hvorved den har mere end overvundet Verden (da Christendommen jo aldrig har villet seire verdsligt), den Uendelighedens Forandring, som Christendommen tilsigter, hvorved i Sandhed Alt er blevet som det var og dog i Uendelighedens Forstand Alt nyt (thi Christendommen har aldrig været Ven af Nyheds-Kram) – saa veed jeg intet kortere, men heller intet mere afgjørende end dette: den har gjort ethvert menneskeligt Forhold mellem Menneske og Menneske til et Samvittigheds-Forhold. Christendommen har ikke villet styrte Regjeringer fra Thronen, for at sætte sig selv paa Thronen, den har i udvortes Forstand aldrig stridt om Pladsen i den Verden, af hvilken den ikke er, (thi i Hjertets Rum, hvis den der finder Plads, tager den dog ingen Plads i Verden), og dog har den uendeligt forandret alt Det, som den lod og som den lader bestaae. Som nemlig Blodet banker i hver en Nerve, saaledes vil Christendommen med Samvittighedsforholdet gjennemtrænge Alt. Forandringen er ikke i det Udvortes, ikke i det Tilsyneladende, og dog er Forandringen uendelig, som hvis et Menneske, istedenfor Blod i sine Aarer havde hiin guddommelige Vædske, om hvilken Hedenskabet drømte – saaledes vil Christendommen indaande det evige Liv, det Guddommelige i Menneskeslægten. Derfor har man sagt, at de Christne var et Folk af Præster, og derfor kan man, naar man betænker Samvittigheds-Forholdet, sige, at det er et Folk af Konger. Thi tag det ringeste, det meest overseete Tjeneste-Menneske, tænk Dig hvad vi kalde en ret eenfoldig fattig stakkels Arbeidskone, der erhverver sit Udkomme ved det ringeste Arbeide: hun har, christeligt forstaaet, Lov til, ja vi bede hende ret indstændigt i Christendommens Navn, at hun vil gjøre det, hun har Lov til, medens hun udfører sit Arbeide, talende med sig selv og med Gud, hvad der ingenlunde sinker Arbeidet, hun har Lov til at sige »jeg gjør dette Arbeide for Daglønnen, men at jeg gjør det saa omhyggeligt som jeg gjør, det gjør jeg – for Samvittighedens Skyld.« Ak, verdsligt er der kun eet Menneske, eet eneste, der ingen anden Forpligtelse erkjender, end Samvittighedens: det er Kongen. Og dog har hiin ringe Kone, christelig forstaaet, Lov til, kongeligt, at sige til sig selv for Gud »jeg gjør det for Samvittighedens Skyld!« Bliver Konen misfornøiet, fordi intet Menneske vil høre paa denne Tale, saa viser det blot, at hun ikke er christeligt sindet, ellers synes jeg dog, at det kan være nok, at Gud har tilladt mig at tale saaledes med ham – begjerligt at forlange Yttrings-Frihed i denne Henseende er en stor Daarskab mod sig selv; thi der er visse Ting, og deriblandt i Særdeleshed Inderlighedens Hemmeligheder, som tabe ved at offentliggjøres, og som ere ganske tabte, naar Offentliggjørelsen er bleven En det vigtigste, ja der er Hemmeligheder, som i saa Fald ikke blot ere tabte men ligefrem blevne Meningsløshed. Christendommens guddommelige Mening er i Fortrolighed at sige til ethvert Menneske: »hav ikke travlt med at forandre Verdens Skikkelse, eller Dit Vilkaar, som hvis Du, for at blive ved Exemplet, istedenfor at være en stakkels Arbeidskone, maaskee kunde drive det til at blive kaldet Madamme, o nei, \tlg{tilegn} Dig det Christelige, og da skal det vise Dig et Punkt udenfor Verden, ved Hjælp af dette skal Du bevæge Himmel og Jord, ja, Du skal gjøre det endnu Vidunderligere, Du skal bevæge Himmel og Jord saa stille, saa let, at ingen mærker det.«
\end{quote}
\dcite{KG:138  (SKS 9, 138)}
% KG:138  (SKS 9, 138)

\begin{quote}
See derfor ønske vi at ende denne som al vor Tale, der efter den os forundte Evne priser det Christelige, med denne lidet indsmigrende Anbefaling: vogt Dig at begynde paa at gjøre det – dersom det ikke i Sandhed er Din Alvor at ville i Sandhed fornegte Dig selv. Vi have for alvorlig en Forestilling om det Christelige til at ville lokke Nogen, vi ville hellere næsten advare. Den, der i Sandhed vil \tlg{tilegne} sig det Christelige, han vil dog komme til indvortes at opleve ganske andre Forfærdelser end Forfærdelsens Smule Skuespil i en Tale; han maa udefter være ganske anderledes besluttet end han kan blive det ved Hjælp af en Smule Veltalenheds sminkede Usandhed. Vi overlade til Enhver at prøve, om denne vor alvorlige Forestilling skulde synes kold, trøstesløs, uden Begeistring. Forsaavidt En skulde tale om sit eget Forhold til Verden, da vilde dette være en anden Sag, da er det Pligt at tale saa mildt, saa undskyldende som muligt, og selv naar han gjør det, er det Pligt at blive i Kjerlighedens Gjeld. Men naar vi tale veiledende, tør vi ikke fortie, hvad der maaskee er lidet skikket til at indynde Talen i en sværmersk Ynglings higende Forestilling. Vi tør heller ikke anprise, smilende at ville opløfte sig over Verdens Modstand og Daarskab; thi selv om det lader sig gjøre, som det er gjort i Hedenskabet, det lader sig og kun gjøre i Hedenskabet, fordi Hedningen ikke har det Christeliges sande, alvorlige, evigt bekymrede Forestilling om det Sande; thi for denne er det ingenlunde latterligt, at Andre mangle den. Christeligt forstaaet er Verdens væsentlige Daarskab ingenlunde latterlig, selv hvor latterlig den er; thi naar der er en Salighed at vinde eller tabe, saa er det hverken en Spøg, om jeg vinder den, ei heller latterligt, om Nogen forskjertser den.
\end{quote}
\dcite{KG:196  (SKS 9, 196)}
% KG:196  (SKS 9, 196)

\begin{quote}
Overalt hvor det Christelige er, er Forargelsens Mulighed, men Forargelse er den høieste Fare. Enhver, der i Sandhed har \tlg{tilegnet} sig det Christelige eller Noget af det Christelige, han har ogsaa maattet gaae Forargelsens Mulighed saaledes forbi, at han har seet den, og med den for Øie – valgt det Christelige. Skal der tales om det Christelige, maa Talen bestandigt holde Forargelsens Mulighed aaben, men altsaa kan den aldrig komme til ligefremt at anbefale Christendommen, saa Forskjellen mellem Talerne blot vilde være den, at een gjorde det i stærkere, en anden i svagere, en tredie i de allerstærkeste Lovprisningens Udtryk. Christendommen kan kun anprises, idet paa ethvert Punkt Faren i eet væk gjøres aabenbar: hvorledes det Christelige er den blot menneskelige Forestilling Daarskab og Forargelse. Men ved at gjøre dette tydeligt og aabenbart, advares der jo. Saa alvorlig er Christendommen. Hvad der trænger til Menneskenes Bifald gjør sig strax lækkert for dem, men Christendommen er saaledes sikker paa sig selv, og veed med saadan Alvor og Strenghed, at det er Menneskene, der trænge til den, at den just derfor ikke anbefaler sig ligefrem, men først skrækker Menneskene op – som Christus anbefalede sig til Apostlene ved itide at forudsige dem, at de for hans Skyld vilde blive forhadte, ja at Den, som slog dem ihjel, vilde mene at gjøre Gud en Tjeneste.
\end{quote}
\dcite{KG:198  (SKS 9, 198)}
% KG:198  (SKS 9, 198)

\begin{quote}
Men medens man saaledes ved at tage Forargelsens Mulighed bort har faaet hele Verden christnet, skeer bestandigt det Besynderlige: at Verden forarges paa den virkelige Christen. Her kommer Forargelsen, hvis Mulighed nu engang er uadskillelig fra det Christelige. Kun er Forvirringen sørgeligere end nogensinde; thi engang forargedes Verden paa Christendommen – det var der Mening i; men nu har Verden faaet den Indbildning, at den er Christen, at den har \tlg{tilegnet} sig Christendommen uden at mærke Noget til Forargelsens Mulighed – og forarges saa paa den virkelige Christen. Sandelig, et saadant Sandsebedrag er vanskeligt at komme ud af. Vee, de hurtige Penne og de travle Tunger, vee den hele Travlhed, som fordi den hverken veed det Ene eller det Andet, derfor saa uendelig let kan forlige baade det Ene og det Andet!
\end{quote}
\dcite{KG:200  (SKS 9, 200)}
% KG:200  (SKS 9, 200)

\begin{quote}
Et af de overførte Udtryk, som den hellige Skrift hyppigst bruger, eller et af de Ord, som den hellige Skrift hyppigst bruger overført er: at opbygge. Og det er allerede – ja det er saa opbyggeligt at see, hvorledes den hellige Skrift ikke trættes ved dette simple Ord, hvorledes den ikke aandrigt eftertragter Afvexling og nye Vendinger, men tvertimod, hvad der er Aandens sande Væsen, fornyer Tanken i det samme Ord. Og det er – ja det er saa opbyggeligt at see, hvorledes Skriften med det simple Ord faaer betegnet det Høieste og paa den inderligste Maade. Det er fast som hiint Bespisningens Mirakel med det lidet Forraad, der dog ved Velsignelsen strakte rigeligen til og efterlod Overflod. Og det er – ja det er opbyggeligt, om det lykkes Nogen, istedenfor at have travlt med at gjøre nye Opdagelser, som travlt skulle fortrænge det Gamle, ved ydmygt at nøies med Skriftordet, ved taknemligt og inderligt at \tlg{tilegne} sig det fra Fædrene Overleverede, at stifte et nyt Bekjendtskab med det – gamle Bekjendte. Som Børn have vi vel alle ofte leget Fremmed: sandeligen, dette er just Alvor, aandeligt forstaaet at kunne vedblive denne i Alvor opbyggelige Spøg, at lege Fremmed med det gamle Bekjendte.
\end{quote}
\dcite{KG:213.2  (SKS 9, 213)}
% KG:213.2  (SKS 9, 213)

\begin{quote}
Kjerlighed haaber Alt – og bliver dog aldrig til Skamme. Vi tale i Forhold til Haab og Forventning om at blive til Skamme; vi mene da, at En bliver til Skamme, naar hans Haab eller Forventning ikke gaaer i Opfyldelse. Hvori skal nu Skammen ligge? Vel deri, at Ens beregnende Klogskab ikke har regnet rigtigt, at det (til Ens Skam) bliver aabenbart, at man uforstandigt har forregnet sig. Men Herre Gud, den Skam var da ikke saa farlig; det er dog vel nok ogsaa egentligen kun saa i Verdens Øine, hvis Begreb om Ære og Skam man imidlertid ikke skal sætte Ære i at \tlg{tilegne} sig. Thi det, Verden beundrer meest og ene giver Æren, er Klogskab, eller at handle klogt; men det at handle klogt er just det Foragteligste af Alt. Om et Menneske er klog, kan han i en vis Forstand ikke selv gjøre for; at han udvikler sin Klogskab, skal han heller ikke skamme sig ved; men vel desto mere ved at handle klogt. Og vist er det (hvad især gjøres fornødent at sige i disse kloge Tider, hvor Klogskab egentligen er bleven Det, som skal ved Christendommens Hjælp overvindes, ligesom engang Raahed og Vildhed), dersom Menneskene ikke lære at foragte det, at handle klogt, lige saa dybt, som man foragter det, at stjæle, at sige falsk Vidnesbyrd: saa afskaffer man tilsidst ganske det Evige, og dermed Alt, hvad der er helligt og Ære værd – thi det, at handle klogt, er just med sit hele Liv at aflægge falskt Vidnesbyrd mod det Evige, er just at stjæle sin Tilværelse fra Gud. Thi det at handle klogt er Halvhed, hvorved man unegteligt kommer videst i Verden, vinder Verdens Goder og Fordele, Verdens Ære, fordi Verden og Verdens Fordeel er, evigt forstaaet, Halvhed. Men heller aldrig har hverken det Evige, eller den hellige Skrift lært noget Menneske at stræbe efter at komme vidt eller videst i Verden, tvertimod advarer den, ikke at komme for vidt i Verden, for om muligt at bevare sig reen fra Verdens Besmittelse. Men er Forholdet dette, saa synes det ikke anpriseligt at tragte efter at komme videst eller vidt i Verden.
\end{quote}
\dcite{KG:260  (SKS 9, 260)}
% KG:260  (SKS 9, 260)

\begin{quote}
Forunderlige Ihukommelse, som den Kjerlige erhverver til Tak for al sin Arbeiden! Han kan paa en Maade pakke hele sit Liv sammen i en Tankestreg. Han kan sige: jeg har arbeidet trods Nogen, aarle og silde, men hvad jeg har udrettet – en Tankestreg! (Kunde det nemlig sees ligefrem hvad han havde udrettet, saa havde han arbeidet mindre kjerligt.) Jeg har lidt, tungt som noget Menneske, inderligt som kun Kjerlighed kan lide; men hvad jeg har gavnet – en Tankestreg! Jeg har forkyndt det Sande, klart og vel gjennemtænkt trods Nogen; men hvo der har \tlg{tilegnet} sig det – en Tankestreg! Havde han nemlig ikke været den Kjerlige, saa havde han, mindre vel gjennemtænkt, skreget det Sande ligefrem ud, og saa havde han strax havt Tilhængere, der havde \tlg{tilegnet} sig det Sande – og hilst ham som Mester.
\end{quote}
\dcite{KG:277  (SKS 9, 277)}
% KG:277  (SKS 9, 277)

\begin{quote}
Nei, Lige for Lige! vi sige det sandeligen ikke, som var det vor Mening, at et Menneske saa dog til syvende og sidst fortjente Naaden. O, det Første Du lærer af i Alt at forholde Dig til Gud, er netop, at Du slet ingen Fortjeneste har. Prøv det blot at sige til Evigheden »jeg har fortjent«, saa svarer Evigheden »Du har fortjent....«. Vil Du have Fortjeneste og have fortjent Noget, saa er Straf det Eneste; vil Du ikke troende \tlg{tilegne} Dig en Andens Fortjeneste, saa faaer Du efter Fortjeneste. – Vi sige dette heller ikke som var det vor Mening, at det var bedre, om Nogen Dag ud og Dag ind sad i Dødsens Angest for at lytte efter Evighedens Gjentagelse, vi sige end ikke, at det var bedre end den Smaalighed, der i disse Tider benytter Guds Kjerlighed til at sælge Aflad fra enhver farefuldere og mere anstrenget Stræben. Nei, men som det veltugtede Barn har en uforglemmelig Forestilling om Strengheden, saaledes maa ogsaa det Menneske, som forholder sig til Guds Kjerlighed, hvis han ikke »kjerlingeagtigt« (1 Tim. 4, 7.) eller letsindigt skal tage den forfængelig, have en uforglemmelig Frygt og Bæven, skjøndt han hviler i Guds Kjerlighed. En Saadan vil vist ogsaa undgaae at tale til Gud om Andres Uret mod ham, om Splinten i Broderens Øie; thi en Saadan vil helst tale til Gud ene om Naaden, for at ikke dette skjebnesvangre Ord »Ret« skal forspilde ham Alt, ved hvad han selv fremkaldte, det strenge Lige for Lige.
\end{quote}
\dcite{KG:378  (SKS 9, 378)}
% KG:378  (SKS 9, 378)

\subsection{CT \normalfont(SKS 10)}

\begin{quote}
Den ringe Christen, der for Gud er sig selv, er som Christen til for sit Forbillede. Han troer, at Gud har levet paa Jorden, at Han har ladet sig føde i ringe og fattige Kaar, ja i Vanære, og derpaa som Barn levet sammen med den simple Mand, der kaldtes Hans Fader, og den foragtede Jomfrue, der var Hans Moder. At Han derpaa vandrede omkring i en Tjeners ringe Skikkelse, ikke til at adskille fra andre ringe Mennesker end ikke ved sin paafaldende Ringhed, indtil Han endte i den yderste Elendighed, korsfæstet som en Forbryder – og saa rigtignok efterlod sig et Navn; men den ringe Christens Begjæring er jo blot at turde i Liv og Død \tlg{tilegne} sig Hans Navn, eller Navn efter Ham. Den ringe Christen troer, som det er fortalt, at Han til sine Disciple valgte ringe Mennesker af den simpleste Stand, og at Han til sin Omgang søgte dem, som Verden forskjød og foragtede; at Han gjennem sit Livs forskjellige Omvexlinger, da Menneskene vilde høit ophøie Ham, og da de vilde fornedre Ham om muligt end mere, end Han selv havde fornedret sig, blev de ringe Mennesker tro, til hvem Han var knyttet ved nærmere Forbindelser, de ringe Mennesker tro, hvem Han havde knyttet til sig, de foragtede Mennesker tro, hvem man just udstødte af Synagogen, fordi Han havde hjulpet dem. Den ringe Christen troer, at dette ringe Menneske eller at dette Hans Liv i Ringhed har viist, hvad et ringe Menneske har at betyde, ak, og hvad et, menneskelig talt, fornemt Menneske dog egentligen har at betyde, hvor uendelig Meget det kan betyde at være et ringe Menneske, og hvor uendeligt Lidt det at være et fornemt Menneske, hvis man ikke er Andet. Den ringe Christen troer, at dette Forbillede just er til for ham, han, som jo er et ringe Menneske, kjempende maaskee med Armod og trange Kaar, eller det endnu Ringere, foragtet og forstødt. Han tilstaaer vel, at han jo ikke er i det Tilfælde selv at have valgt denne overseete eller foragtede Ringhed, og forsaavidt ikke ligner Forbilledet; men dog fortrøster han sig til, at Forbilledet er til for ham, Forbilledet, der ved Hjælp af Ringheden barmhjertigt ligesom paanøder sig ham, som vilde det sige »ringe Menneske, kan Du ikke see, at dette Forbillede er for Dig«. Vel har han ikke med egne Øine selv seet Forbilledet, men han troer, at Han har været til. Der havde jo i en vis Forstand heller ikke været Noget at see – uden Ringheden (thi Herligheden, den maa troes); og Ringheden kan han da nok gjøre sig en Forestilling om. Han har ikke med egne Øine seet Forbilledet, han gjør heller ikke noget Forsøg paa at lade Sandsen danne sig et saadant Billede. Dog seer han ofte Forbilledet. Thi hver Gang han i Troens Glæde over dette Forbilledes Herlighed ganske forglemmer sin Armod, sin Ringhed, sin Foragtethed, da seer han Forbilledet – og da seer han selv tilnærmelsesviis ud som Forbilledet. Dersom nemlig i et saadant saligt Øieblik, naar han er fortabt i sit Forbillede, et andet Menneske seer paa ham, saa seer dette andet Menneske kun et ringe Menneske for sig: saaledes var det just ogsaa med Forbilledet, man saae kun det ringe Menneske. Han troer og haaber stedse mere og mere at skulle nærme sig til Lighed med dette Forbillede, der først hisset skal vise sig i sin Herlighed; thi her paa Jorden kan det kun være i Ringhed og kun sees i Ringheden. Han troer, at dette Forbillede, hvis han stadigt kæmper for at ligne det, bringer ham anden Gang og i endnu nøiere Slægtskab med Gud, at han ikke blot har Gud til Skaber, som alle Skabninger have, men har Gud til Broder.
\end{quote}
\dcite{CT:54  (SKS 10, 54)}
% CT:54  (SKS 10, 54)

\begin{quote}
Og naar Du da saa ret er bleven fattig, saa har Du ogsaa mere og mere \tlg{tilegnet} Dig Aandens Goder. Da skal Du ogsaa være istand til i endnu fuldere Maade at gjøre Andre rige, ved at meddele Aandens Goder, ved at meddele hvad der i sig selv er Meddelelse. Jo fattigere Du bliver, desto sjeldnere i Dit Liv vil det Øieblik være, hvor Du selvisk er beskjeftiget med Dig selv, eller med hvad der i sig selv er selvisk, det Jordiske, som drager et Menneskes Tanke saaledes til sig, at han saa længe ikke er til for Andre. Men jo sjeldnere saadanne Øieblikke blive, desto stadigere vil da Din Dags Timer være udfyldte med at erhverve og besidde Aandens Goder (og glem det dog ikke, o glem det ikke, ogsaa naar Du gjør det, gjør Du Andre rige!) eller med ligefrem at meddele Aandens Goder, og derved at gjøre Andre rige.
\end{quote}
\dcite{CT:133  (SKS 10, 133)}
% CT:133  (SKS 10, 133)

\begin{quote}
De høre Hans Røst. Og idag er det ganske særligen, er det ene og alene Hans Røst, der skal høres. Alt, hvad der ellers foretages, er blot for at samle Sindets Opmærksomhed paa, at det er Hans Røst, der skal høres. Idag prædikes der ikke. En Skriftetale er ingen Prædiken; den vil ikke belære Dig, eiheller indskjærpe Dig de gamle bekjendte Lærdomme; den vil blot standse Dig paa Veien til Alteret, for at Du ved den Talendes Stemme ved Dig selv skrifter for Gud i Løndom. Thi Du skal af en Skriftetale ikke lære hvad det er at skrifte, det var ogsaa for sildigt, men ved den skrifter Du for Gud. Idag prædikes der ikke. Hvad vi her tale i det foreskrevne korte Øieblik er atter ingen Prædiken; og naar vi have sagt Amen, da er ikke som ellers Gudstjenesten væsentligen forbi, men da begynder den væsentlige. Vor Tale vil derfor blot et Øieblik standse Dig paa Veien til Alteret; thi idag samler Gudstjenesten sig ikke som ellers om Prædikestolen men om Alteret. Og ved Alteret gjælder det fremfor Alt om at høre Hans Røst. Vel skal nemlig en Prædiken ogsaa vidne om Ham, forkynde Hans Ord og Hans Lære; men derfor er en Prædiken dog ikke Hans Røst. Ved Alteret derimod er det Hans Røst, Du skal høre. Om et andet Menneske sagde Dig, hvad der siges ved Alteret, hvis alle Mennesker ville forene sig om at sige Dig det – dersom Du ikke hører Hans Røst, da gik Du forgjeves til Alters. Naar der ved Alteret af Herrens Tjener nøiagtigt siges ethvert Ord, som det er overleveret fra Fædrene; naar Du nøiagtigt hører ethvert Ord, saa der ikke undgaaer Dig det Mindste, ikke en Tøddel – dersom Du ikke hører Hans Røst, at det er Ham, som siger det, da gik Du forgjeves til Alters. Om Du troende \tlg{tilegner} Dig hvert Ord, der siges; om Du alvorligt beslutter at tage Dig det til Hjerte og at indrette dit Liv i Overeensstemmelse dermed – dersom Du ikke hører Hans Røst: saa gik Du forgjeves til Alters. Det maa være Hans Røst Du hører, naar han siger: kommer hid alle I, som arbeide og ere besværede, altsaa Hans Røst, som indbyder Dig; og det maa være Hans Røst, Du hører, naar Han siger: dette er mit Legeme. Thi ved Alteret er der ingen Tale om Ham; der er Han selv personligt tilstede, der er det Ham, som taler – hvis ikke, saa er Du ikke ved Alteret. Sandseligt forstaaet kan man nemlig pege paa Alteret og sige: »der er det«; men aandeligt forstaaet er det dog egentligen kun der, hvis Du der hører Hans Røst.
\end{quote}
\dcite{CT:290  (SKS 10, 290)}
% CT:290  (SKS 10, 290)

\begin{quote}
Stor er Du, o Gud; skjøndt vi kun kjende Dig som i en mørk Tale og som i et Speil, vi tilbede dog undrende Din Storhed – hvor meget mere skulle vi engang prise den, idet vi lære den fuldeligere at kjende! Naar jeg under Himmelens Hvælving staaer omgiven af Skabningens Under, da priser jeg rørt og tilbedende Din Storhed, Du, som let bærer Stjernerne i det Uendelige, og faderligt bekymrer Dig om Spurven. Men naar vi ere forsamlede her i Dit hellige Huus, da ere vi jo ogsaa overalt omgivne af hvad der i end dybere Forstand minder om Din Storhed. Thi stor er Du, Verdens Skaber og Opholder; men da Du, o Gud, tilgav Verdens Synd og forligte Dig med den faldne Slægt, ak, da var Du jo dog end større i Din ubegribelige Forbarmelse! Hvorledes skulle vi da ikke troende prise og takke og tilbede Dig her i Dit hellige Huus, hvor Alt minder os derom, særligen dem, som idag ere forsamlede for at annamme Syndernes Forladelse, og for paa ny at \tlg{tilegne} sig Forligelsen med Dig i Christo!
\end{quote}
\dcite{CT:311.2  (SKS 10, 311)}
% CT:311.2  (SKS 10, 311)

\subsection{TSA \normalfont(SKS 11)}

\begin{quote}
Hvad er da Myndighed? Er Myndighed Lærens Dybsind, dens Fortrinlighed, dens Aandrighed? Ingenlunde. Dersom Myndighed saaledes blot skulde betegne, i en anden Potens eller redupliceret, at Læren er dybsindig: saa er der netop ingen Myndighed; thi hvis da en Lærende heelt og fuldt ved Forstaaelsen \tlg{tilegnede} sig denne Lære, saa blev der jo ingen Forskjel tilbage mellem Læreren og den Lærende. Myndighed er derimod Noget, som bliver uforandret, som man ikke kan erhverve ved paa det fuldeste at have forstaaet Læren. Myndighed er en specifik Qvalitet, der træder til andetstedsfra og netop qvalitativt gjør sig gjældende, naar Udsagnets eller Gjerningens Indhold æsthetisk er sat i Indifferents. Lad os tage et Exempel, saa simpelt som muligt, hvor dog Forholdet viser sig. Naar Den, der har Myndigheden til at sige det, siger til et Menneske: gaa! og naar Den, der ikke har Myndigheden, siger: gaa! saa er jo Udsagnet (gaa!) og dets Indhold identisk, æsthetisk vurderet er det om man saa vil lige godt sagt, men Myndigheden gjør Forskjellen. Dersom Myndigheden ikke er det Andet (το ετεϱον), dersom den paa nogen Maade blot skulde betegne en Potenseren indenfor Identiteten, saa er der netop ingen Myndighed. Dersom saaledes en Lærer er sig begeistret bevidst, at han selv existerende udtrykker og har udtrykt, med Opoffrelse af Alt, den Lære han forkynder: saa kan denne Bevidsthed vel give ham en vis og en fast Aand, men den giver ham ikke Myndighed. Hans Liv som Beviis for Lærens Rigtighed er ikke det Andet (το ετεϱον) men er en simpel Fordoblelse. Det at han lever efter Læren beviser ikke, at den er rigtig; men fordi han er selv overbeviist om Lærens Rigtighed, derfor lever han efter den. Derimod, hvad enten en Politie-Officiant for Exempel er en Slyngel, eller han er en retskaffen Mand, saasnart han er i Function, har han Myndighed.
\end{quote}
\dcite{TSA:103  (SKS 11, 103)}
% TSA:103  (SKS 11, 103)

\begin{quote}
Sagen er imidlertid denne. Tvivlen og Vantroen, der gjør Troen forfængelig, har blandt Andet ogsaa gjort Menneskene generede ved at lyde, ved at bøie sig under Myndighed. Denne Oprørskhed sniger sig ind selv i de Bedres Tankegang, dem maaskee ubevidst, og saa begynder alt dette Forskruede, som i Grunden er Forræderie, om det Dybe og det Dybe, og det Vidunderlig-deilige, som man kan skimte og saa videre Skulde man derfor med eet bestemt Prædikat betegne det christelig-religieuse Foredrag, som det nu høres og læses, da maatte man sige, det er affecteret. Naar man ellers taler om en Præsts Affectation, tænker man maaskee paa, at han pynter og sminker sig, eller at han sødligt smægter med Stemmen, eller at han norsk snurrer paa R og rynker Brynet, eller at han anstrænger sig i Kraftstillinger og i Opvækkelsens Spring og saa videre Dog alt Sligt er mindre vigtigt, om det end altid var ønskeligt, at det ikke var. Men det Fordærvelige er, naar Prædike-Foredragets Tankegang er affecteret, naar dets Rettroenhed naaes ved at lægge Accenten paa et reent forkeert Sted, naar det i Grunden opfordrer til at troe paa Christus, prædiker Troen paa Ham paa Grund af Det, som slet ikke kan blive Troens Gjenstand. Dersom en Søn vilde sige »jeg lyder min Fader, ikke fordi han er min Fader, men fordi han er Genie, eller fordi hans Befalinger altid ere dybsindige og aandrige«: saa er denne sønlige Lydighed affecteret. Sønnen accentuerer noget reent Forkeert, accentuerer det Aandrige, det Dybsindige ved en Befaling, medens en Befaling netop er ligegyldig mod denne Bestemmelse. Sønnen vil lyde i Kraft af Faderens Dybsindighed og Aandrighed; og i Kraft deraf kan han netop ikke lyde ham, thi hans critiske Forhold i Retning af, om nu Befalingen er dybsindig og aandrig, underminerer Lydigheden. Og saaledes er det ogsaa Affectation, naar der er saa megen Tale om at \tlg{tilegne} sig Christendommen og troe Christus, formedelst Lærens Dybsindighed og Dybsindighed. Man \tlg{tilegner} sig Rettroenheden ved at accentuere noget aldeles Forkeert. Den hele moderne Spekulation er derfor affecteret ved at have afskaffet Lydigheden paa den ene Side og Myndigheden paa den anden Side, og ved saa desuagtet at ville være rettroende. En Præst, der er aldeles correct i sit Foredrag, maa tale saaledes, idet han anfører et Ord af Christus: »Dette Ord er af Den, hvem, ifølge hans eget Udsagn, al Magt er given i Himlen og paa Jorden. Du maa nu, min Tilhører, overveie med Dig selv, om Du vil bøie Dig under denne Myndighed eller ikke, antage og troe Ordet eller ikke. Men vil Du det ikke, da gaae for Gud i Himlens Skyld ikke hen og antag Ordet, fordi det er aandrigt eller dybsindigt, eller vidunderlig-deiligt, thi dette er Gudsbespottelse, det er at ville criticere Gud.« Saasnart nemlig Myndighedens, den specifik-paradoxe Myndigheds Dominant er sat i, saa er alle Forhold qvalitativt forandrede, saa er den Art \tlg{Tilegnelse}, som ellers er tilladelig og ønskelig, en Brøde og Formastelighed.
\end{quote}
\dcite{TSA:108  (SKS 11, 108)}
% TSA:108  (SKS 11, 108)

\subsection{IC \normalfont(SKS 12)}

\begin{quote}
Det er den fornedrede Jesus Christus, der har sagt hine Ord. Og intet Ord af Christus, ikke eet eneste har Du Lov at \tlg{tilegne} Dig, Du har ikke mindste Deel i ham, ikke i fjerneste Maade Samfund med ham, hvis Du ikke er blevet saaledes samtidig med ham i hans Fornedrelse, at Du ganske som de Samtidige har maattet blive opmærksom paa den Formaning af ham: salig Den, som ikke forarges paa mig! Du har ikke Lov at tage Christi Ord hen og lyve ham bort; Du har ikke Lov at tage Christi Ord hen og phantastisk gjøre ham til Noget ved Hjælp af Historiens Snaksomhed, der, naar den snakker om ham, ganske bogstavelig ikke veed, hvad den snakker om.
\end{quote}
\dcite{IC:51  (SKS 12, 51)}
% IC:51  (SKS 12, 51)

\begin{quote}
Men som man da i Christenheden har forvirret Alt, saa ogsaa Dette, og derved ganske rigtigt opnaaet, at Christenheden er blevet Hedenskab. Der er i Christenheden en evindelig Prædiken om, hvad der saa skeete efter Christi Død, hvorledes han seirede, og hans Lære seierrigt erobrede hele Verden, kort man hører lutter Prædikener, der mere passende kunde ende med Hurra, end med Amen. Nei, Christi Liv her paa Jorden det er Paradigmet; det er i Lighed dermed jeg og enhver Christen skal stræbe at danne mit Liv, og dette er Prædikenens væsentlige Gjenstand, dertil skal den tjene, til at holde mig i Aande, naar jeg vil blive sløv, og til at styrke, naar man bliver forsagt. Saaledes er han jo Paradigmet i Samtidighedens Situation; i den var der ingen Snakken om hvad der saa skeete bag efter. Men Christenheden har afskaffet Christus, derimod vil den – arve ham, hans store Navn, profitere de uhyre Følger af hans Liv, ja det er ikke langt fra, at den vil \tlg{tilegne} sig dem som dens egne Meriter, og bilde os ind, at Christenheden er Christus. Istedetfor at enhver Generation har at begynde forfra med Christus, og saa at fremstille hans Liv som Paradigma, saa har Christenheden taget sig den Frihed at lægge det hele Forhold reent historisk ud, begynde med at lade ham være død – og saa triumpheres der! Christenheden tiltager siden den Tid Aar for Aar i Mængde – hvad Under vel, der pleier jo de Fleste hellere end gjerne at ville være med, hvor det gjælder om hverken mere eller mindre end at triumphere og om at ride Herredage ind. Og derfor er i Christenheden det at være Christen noget fra det i Samtidighedens Situation at være Christen saa Forskjelligt som Hedenskab og Christendom.
\end{quote}
\dcite{IC:115.3  (SKS 12, 115)}
% IC:115.3  (SKS 12, 115)

\begin{quote}
Man kunde gjøre opmærksom paa, at det, selv om intet Andet var til Hinder, ikke saa let lader sig gjøre, fordi Christi Liv i een Forstand staaer udenfor det ligefremme Forhold til hver Enkelt i Slægten, at han, som Gud-Mennesket, skjøndt sandt Menneske, dog er saaledes ueensartet med det enkelte Menneske, at det ikke saa ganske ligefrem lader sig gjøre uden videre, i en Slags fræk Nærgaaenhed, saadan at ville gaae paa Parti med ham. Man kunde gjøre opmærksom paa, at denne Christi (Gud-Menneskets) Ueensartethed med alle enkelte Mennesker ogsaa er udtrykt ved Læren om hans Gjenkomst. Thi det er ikke med ham som med et andet Menneske, der engang har levet, vundet maaskee en eller anden stor Seier, hvis Følger vi uden videre \tlg{tilegne} os, medens der Intet høres fra ham, endnu mindre, at han kunde komme igjen for at gjøre Regnskabet op med os, for, ved at fordre Sit eller sig selv igjen af os, at dømme os. Anderledes med Christus. Han levede her paa Jorden, dette hans Liv er Forbilledet. Derpaa indgaaer han til Høiheden, og saa siger han ligesom til Slægten: begynder nu I – og hvad er det de skulle begynde paa? paa at leve i Overeensstemmelse med Forbilledet – men, tilføier han, engang ved Tidens Ende kommer jeg igjen. Denne Form af Tilværelse gjør altsaa, om jeg saa tør sige, hele Kirkens Tilværelse her paa Jorden til en Parenthes eller et Parenthetisk i Christi Liv; med Christi Opfart til Høiheden begynder Parenthesens Indhold, og med hans Gjenkomst sluttes den. Det løber altsaa ikke her ud i Eet, som ellers i det historiske Forhold mellem en Enkelt og de Andre, der uden videre tage sig hans Seier til Indtægt; thi hverken er en saadan Enkelt Forbilledet, ei heller er en saadan Enkelt Den der vil komme igjen. Kun Christus er Den, der kan gjøre sit Liv til Prøven for alle Mennesker. Idet han farer op til Himlen, saa begynder Examens Tiden; den har nu varet i 1800 Aar, skal maaskee vare i 18000. Men (og dette hører just med til at det Mellemliggende er Examen) han kommer igjen. Og naar dette er saaledes, saa er al ligefrem Tilslutning til ham for uden videre at tage sig hans Seier til Indtægt umuligere end i Forhold til noget andet Menneske.
\end{quote}
\dcite{IC:199  (SKS 12, 199)}
% IC:199  (SKS 12, 199)

\begin{quote}
Dog dette ville vi her ikke videre opholde os ved, vi ville heller gjøre en anden Betragtning gjældende: er »Sandhed« noget Saadant, at i Forhold til denne det lader sig tænke, at En uden videre kan \tlg{tilegne} sig den ved Hjælp af en Anden? Uden videre, det er, uden selv at ville paa en lignende Maade udvikles, forsøges, kæmpe, lide som Den, der for ham erhvervede Sandheden? er det ikke lige saa umuligt som at sove eller drømme sig til Sandheden, lige saa umuligt uden videre at \tlg{tilegne} sig den, hvor lysvaagen man end er, eller er man lysvaagen, er dette ikke blot en Indbildning, naar man ikke forstaaer eller ikke vil forstaae, at der i Forhold til Sandheden ingen Forkortelse gives, som udelader det at erhverve, og i Forhold til det at erhverve fra Slægt til Slægt ingen væsentlig Forkortning, saa enhver Slægt og Enhver i Slægten væsentligen maa begynde forfra?
\end{quote}
\dcite{IC:200  (SKS 12, 200)}
% IC:200  (SKS 12, 200)

\begin{quote}
Og derfor er, christelig forstaaet, Sandheden naturligviis ikke det, at vide Sandheden, men at være Sandheden. Til Trods for hele den nyeste Philosophi er der en uendelig Forskjel paa dette, hvilket bedst sees af Christi Forhold til Pilatus; thi Christus kunde ikke, kunde kun usandt svare paa det Spørgsmaal, hvad er Sandhed, just fordi han ikke var Den, der vidste hvad Sandhed er, men var Sandheden. Ikke som vidste han ikke, hvad Sandhed er; men naar man er Sandheden, og naar Fordringen er at være Sandhed, er det at vide Sandheden en Usandhed. Thi det at vide Sandheden er Noget, som ganske af sig selv følger med det at være Sandheden, ikke omvendt; og just derfor bliver det Usandhed, naar det at vide Sandheden skilles fra det at være Sandheden, eller naar det at vide Sandheden gjøres til Eet med det at være den, da det forholder sig omvendt: det at være Sandheden er Eet med det at vide Sandheden, og Christus havde aldrig vidst Sandheden, dersom han ikke havde været den; og intet Menneske veed mere af Sandheden end hvad han er af Sandheden. Ja vide Sandheden kan man egentligen ikke; thi veed man Sandheden, maa man jo vide, at Sandheden er at være Sandheden, og saa veed man jo i sin Viden af Sandheden, at det at vide Sandheden er en Usandhed. Vil man sige, at man ved at vide Sandheden er Sandheden, saa siger man det jo selv, at Sandheden er at være Sandheden, idet man siger, at vide Sandheden er at være Sandheden, da man i andet Fald maatte sige: Sandheden er at vide Sandheden, ellers kommer Spørgsmaalet om Sandhed kun igjen, saa Spørgsmaalet ikke er besvaret, men kun den afgjørende Besvarelse udsat, idet man jo igjen maa kunne vide, om man er Sandheden eller ikke. Det vil sige: Viden forholder sig til Sandheden, men saa længe er jeg usandt udenfor mig selv; i mig, det er, naar jeg sandt er i mig (ikke usandt udenfor mig selv), er Sandheden, hvis den er der, en Væren, et Liv. Derfor hedder det: »det er det evige Liv at kjende den eneste sande Gud, og Den, hvem han udsendte«, Sandheden. Det vil sige, kun da kjender jeg i Sandhed Sandheden, naar den bliver et Liv i mig. Derfor sammenligner Christus Sandheden med Spise, og det at \tlg{tilegne} sig den med det at spise; thi ligesom, legemligt, Spisen ved at \tlg{tilegnes} (assimileres) bliver det Livet Opholdende, saaledes er ogsaa, aandeligt, Sandheden baade det Livet Givende og det Livet Opholdende, er Livet. Og derfor seer man, hvilken uhyre Vildfarelse det er, noget nær den størst mulige, at docere Christendom; og hvor forandret Christendommen er blevet ved denne idelige Doceren, seer man deraf, at nu alle Udtryk ere dannede i Retning af, at Sandheden er Erkjendelse, Viden (man taler nu bestandigt om at begribe, speculere, betragte og saa videre), hvorimod i den oprindelige Christendom alle Udtryk ere dannede i Retning af, at Sandheden er en Væren.
\end{quote}
\dcite{IC:203  (SKS 12, 203)}
% IC:203  (SKS 12, 203)

\subsection{OI6 \normalfont(SKS 13)}

\begin{quote}
4) At vi, at »Christenhed« slet ikke kan \tlg{tilegne} sig Christi Forjættelser; thi vi, »Christenhed« er ikke der, hvor Christus og det nye Testamente fordrer, at man skal være for at være Christen.
\end{quote}
\dcite{OI6:255.7  (SKS 13, 255)}
% OI6:255.7  (SKS 13, 255)

\begin{quote}
At vi, at »Christenhed« slet ikke kan \tlg{tilegne} sig Christi Forjættelser; thi vi, »Christenhed«, ere ikke der, hvor Christus og det nye Testamente fordrer, at man skal være for at være Christen.
\end{quote}
\dcite{OI6:268  (SKS 13, 268)}
% OI6:268  (SKS 13, 268)

\subsection{OIA \normalfont(SKS 15)}

\begin{quote}
Dog til Sagen. I denne Comoedie optræder nemlig Skuespildirecteur Dalby i flere Roller, blandt andet ogsaa som Haarskjærer. Det er kun denne Præstation, jeg har med at gjøre. Denne Person skal nu være en vrøvlevorren Hegelianer. Dersom jeg nu vilde sige til Andersen, at jeg aldrig har udgivet mig for nogen Hegelianer, og det forsaavidt var taabeligt af Andersen, at tage Sætninger af min lille Piece og lægge en Hegelianer i Munden, saa vilde jeg næsten forekomme mig selv afsindig. Thi enten maatte jeg med det Ord »Hegelianer« forbinde Forestillingen om en Mand, der med Alvor og Energie havde omfattet denne Tænknings Verdens-Anskuelse, \tlg{tilegnet} sig den, fundet Hvile i den, og nu med en vis sand Stolthed sagde om sig selv: ogsaa jeg har haft den Ære at tjene under Hegel, – og i dette Tilfælde vilde det være et Raserie af mig at sige det i en Samtale med Andersen, da han formodentlig ikke vilde kunne knytte en fornuftig Tanke dertil, ligesom jeg idetmindste vilde tage i Betænkning at bruge et saa betydningsfuldt Prædikat om mig selv, om jeg end var mig bevidst, at jeg havde stræbt at gjøre mig fortrolig med Hegels Philosophie. Eller jeg vilde ved en Hegelianer forstaae et Menneske, der, flygtigt berørt af denne Tænkning, nu bedaarede sig selv i et Resultat, han ikke eiede, – og da vilde det være ikke mindre daarligt at sige dette til Andersen, saafremt man ellers vil give mig Ret i, at den som ikke veed, hvad en Hegelianer er i Sandhed, heller ikke veed hvad en Hegelianer er i sin Usandhed, det vil sige: ingen konkret Forestilling kan have derom.
\end{quote}
\dcite{OIA:9.2  (SKS 15, 9)}
% OIA:9.2  (SKS 15, 9)

\subsection{JC \normalfont(SKS 15)}

\begin{quote}
Forsaavidt der havde været en tidligere Philosophie havde vel ogsaa den enkelte Philosoph benyttet Forgængerne, seet, at han kunde \tlg{tilegne} sig Dette, berigtige Hiint og saa videre, men det var vel ikke faldet ham ind, at ville gjennemskue den evige Nødvendighed, med hvilken den ene Philosoph fremgik af den Anden, han selv af sine Forgængere i en evig Continuitæt. Kunde det end lykkes Tænkningen med Hensyn til det Forbigangne at ahne en saadan indre Nødvendighed, hvorved da ogsaa blev at bemærke, at jo mere fjernet det var, desto større blev Muligheden af en Skuffelse, med Hensyn til det Nærværende syntes det ham en Umulighed. Dette fik da ikke Lov at blive et Nærværende af Iver for, at det jo før, jo hellere skulde blive et Forbigangent, men paa den Maade blev det ingen af Delene. Herover blev han sig selv klar ved at betragte det personlige Liv. Naar Een skuer tilbage over sit Liv, da kan især den tidligere Deel deraf vel vise sig som gjennemtrængt af en Nødvendighed. Hvis derimod Een idet han begynder et bestemt Tids-Afsnit først vil blive sig dette bevidst i sin evige Gyldighed som Moment i sit Liv, saa vil han netop forhindre det fra at komme til at faae Betydning, idet han vil ophæve det, førend det har været, idet han vil, at det, der er et Nærværende, i samme Moment skal vise sig for ham som et Forbigangent.
\end{quote}
\dcite{JC:36  (SKS 15, 36)}
% JC:36  (SKS 15, 36)

\begin{quote}
Skjøndt Johannes vel kunde fastholde dette, saa var hans Sindsstemning af den Beskaffenhed, at han ikke havde Mod til saaledes at være consequent ligeoverfor lovpriste Sandheder. Var det end inconsequent af Geniet at ville fordre det af noget Menneske, saa var det dog consequent af en Stakkels Student at gjøre det. Han indsaae vel det Ufuldkomne i den Maade, paa hvilken han \tlg{tilegnede} sig Sandheden, men derfor vilde han dog ikke slippe Sætningen. Han stræbte da atter at tænke over den, for at see, hvorledes han kunde komme i Forhold til den. Det var endnu ikke Sætningen selv han vilde gjennemtænke; thi han maatte jo først vide, om det kunde lykkes ham at komme i Forhold til den. Han spurgte derfor ikke saaledes: om Tvivlen som Begyndelsen er en Deel af Philosophien eller den er hele Philosophien? Hvis den er en Deel, hvilken da den anden Deel er? Om det er Visheden? Om disse Dele ere i al Evighed adskilte? Hvorledes der kan være tale om et Hele, naar dets Dele udelukke hinanden? Hvad nemlig Epicur paa en sophistisk Maade havde gjort gjeldende med Hensyn til Frygten for Døden, at man ikke skal bekymre sig om den; thi naar jeg er, er Døden ikke, naar Døden er, er jeg ikke, det forekom ham her at have sin Sandhed. Om der da var Noget, der forenede disse tvende Dele til et Hele? Saaledes spurgte han ikke, men derimod om den Enkeltes Forhold til hiin Sætning.
\end{quote}
\dcite{JC:39.2  (SKS 15, 39)}
% JC:39.2  (SKS 15, 39)

\begin{quote}
Om Sætningen da slet ikke lod sig modtage, men blot fremsætte; Enhver modtog den saaledes, at i det Øieblik, han fremsatte den, var det ligegyldigt af hvem, han havde modtaget den, eller om han havde modtaget den, da han først havde modtaget den, idet han selv fremsatte den? – Om den lod sig modtage, om den Enkelte ved en Anden kunde modtage den, om den skulde troes? Idet jeg nemlig troende modtager en Sætning, da er jeg ikke strax istand til at fatte den eller at udføre den, men dog modtager jeg den, idet jeg troer den, der fremsætter den. – Om maaskee Sætningen var af den Beskaffenhed, at der hos den, der skulde fremsætte den, fordredes Myndighed, hos den, der skulde modtage den, Tillid og Hengivelse? – Om den skulde troes saaledes, at den Enkelte ikke gjorde hvad Sætningen sagde, men troede, at den Anden havde gjort det? Om maaskee en enkelt Philosoph havde tvivlet for Alle, ligesom Christus havde lidt for Alle, og man nu blot skulde troe dette, ikke selv tvivle? I saa Fald var jo Sætningen ikke ganske rigtigt fremsat; thi Philosophien begyndte da ikke for den Enkelte med Tvivl, men med Tro paa, at Philosophen N. N. havde tvivlet for ham. – Om Sætningen troende skulde \tlg{tilegnes} saaledes, at den Enkelte gjorde, hvad den sagde? Om da den, der havde fremsat den, havde saa fuldstændigt tvivlet om Alt, at den Enkelte blot repeterede hans Tvivl, og han saaledes i Troen paa Fremsætteren gjorde Tvivlens Bevægelser, saaledes som denne foreskrev det? Om der ved enhver Enkelt bestandig kom et nyt Moment af Tvivl til for den næste? Om man med Hensyn til det, der havde været den tidligere Enkelte muligt at tvivle om, om man med Hensyn til det skulde troe, at han havde tvivlet tilgavns, eller man atter skulde tvivle?
\end{quote}
\dcite{JC:46  (SKS 15, 46)}
% JC:46  (SKS 15, 46)

\subsection{BOA \normalfont(SKS 15)}

\begin{quote}
Saasnart man derimod forvexler Christendommen og det Christelige, saasnart man begynder at tælle Aarene: saa begynder man at ville forandre det Usandsynlige til det Sandsynlige. Man siger: nu har Christendommen bestaaet (thi det Christelige er jo et Faktum for 1800 Aar siden) i 300 Aar, nu i 700, nu i 1800: ja saa maa den vist være det Sande. Ved denne Fremgangsmaade vinder man at forvirre Alt. Afgjørelsen (den at blive Christen) bliver for den Enkelte let en reen Ubetydelighed; falder det ham allerede let nok at antage Skik og Brug i den By, hvor han lever, fordi de Mange gjøre det: skulde det saa ikke falde ganske af sig selv at han blev Christen med – naar Christendommen har bestaaet i 1800 Aar! Paa den anden Side bliver det Christelige svækket, gjort til en Ubetydelighed ved Hjælp af Afstanden, ved Hjælp af de 1800 Aar. Hvad der hvis det skete samtidigen vilde komme som Forfærdelse over Menneskene, vilde oprøre dem Tilværelsen fra Grunden, hvad de, hvis det skete samtidigen, enten vilde forarges paa vilde hade og forfølge og om muligt faae udryddet, eller troende \tlg{tilegne} det sig: det synes man nok man saadan kan antage og troe (det vil sige saadan lade staae hen) siden det er 1800 Aar siden. Thi Samtidigheden er Spændingen, der ikke tillader En at lade det staae hen, men tvinger En til enten at forarges eller tro; Afstanden derimod er det Aflad, der begunstiger Dorskhed i den Grad, at det troende at antage Noget saadan bliver identisk med at lade det staae hen. Hvoraf kommer det vel, at ingen Samtid kan komme ud af det med Sandhedens Vidner, og neppe er Manden død saa kan Alle saa ypperligt komme ud af det med ham? Det kommer der af, at saa længe han lever og de Samtidige med ham i Samtidighedens Situation: saa føle de Braaden af hans Existents, han tvinger dem til en mere anstrænget Afgjørelse. Men naar han er død, saa kan man godt være gode Venner med ham og beundre ham det vil sige saadan tankeløst og mageligt lade det Hele staae hen. Hvorfor mon Socrates sammenlignede sig selv med en Bremse, uden fordi han forstod, at hans Leven med de Samtidige var en Braad. Da han var død, forgudede de ham. Naar en Mand oplever en lille Begivenhed i sit Liv: saa lærer han Noget deraf, og hvorfor? Fordi Begivenheden ret kommer ham paa Livet. Derimod kan den samme Mand sidde i Theateret og see Tragediens store Optrin, han kan læse i Aviser om det Overordentlige, han kan høre Præsten: og det gjør Alt ikke ret noget Indtryk, og hvorfor? Fordi han ikke bliver samtidig med det, fordi han i de to første Tilfælde mangler Phantasie, i det sidste indre Erfaring for ret at blive samtidig med det Fremstillede; fordi han tænker som saa, det er nok mange Aar siden det skete.
\end{quote}
\dcite{BOA:161  (SKS 15, 161)}
% BOA:161  (SKS 15, 161)

\begin{quote}
Dog er den Samtidighed, om hvilken her er Tale, ikke en Apostels Samtidighed, forsaavidt han blev kaldet ved en Aabenbaring, men er blot den Samtidighed, som enhver Medlevende havde, den Mulighed i Samtidighedens Spænding at maatte forarges eller gribe Troen. Og til den Ende er det netop nødvendigt at der luftes ud saaledes, at det, som eengang, bliver muligt, at et Menneske for Alvor kan forarges, eller troende \tlg{tilegne} sig det Christelige, at det ikke bliver med det Christelige som med en Retssag, naar den har staaet hen i umindelige Tider: at man hverken veed ud eller ind paa Grund af den megen Viden. Samtidighedens Situation er det Spændende, som giver Kategorierne qvalitativ Elasticitet; og det maa da være en stor Dosmer, der ikke veed, hvilken uendelig Forskjel det gjør, naar man for sin egen Skyld i Samtidighedens Situation betænker Noget, og naar man saadan tænker over Noget i den Indbildning, at det er 1800 Aar siden; i den Indbildning, ja i den Indbildning, thi netop fordi det Christelige er det qvalitative Paradox, er det en Indbildning, at 1800 Aar er længere siden end igaar.
\end{quote}
\dcite{BOA:166.2  (SKS 15, 166)}
% BOA:166.2  (SKS 15, 166)

\begin{quote}
Hvad er da Myndighed? Er Myndighed Lærens Dybsind, dens Fortrinlighed, dens Aandrighed? Ingenlunde. Dersom Myndighed saaledes blot skulde betegne i en anden Potens eller reduppliceret, at Læren er dybsindig: saa er der netop ingen Myndighed; thi hvis da en Lærende heelt og fuldt ved Forstaaelsen \tlg{tilegnede} sig denne Lære, saa blev der jo ingen Forskjel tilbage mellem Læreren og den Lærende. Myndighed er derimod Noget, som bliver uforandret, som man ikke kan erhverve ved paa det fuldeste at have forstaaet Læren. Myndighed er en specifik Qvalitet der træder til andetstedsfra og netop qvalitativt gjør sig gjeldende, naar Udsagnets eller Gjerningens Indhold æsthetisk er sat i Indifferents. Lad os tage et Exempel, saa simpelt som muligt, hvor dog Forholdet viser sig. Naar Den der har Myndigheden til at sige det, siger til et Menneske: gaae!; og naar Den der ikke har Myndigheden siger: gaae!: saa er jo Udsagnet (gaae!) og dets Indhold identisk, æsthetisk vurderet er det om man saa vil lige godt sagt, men Myndigheden gjør Forskjellen. Dersom Myndigheden ikke er det Andet (το ετεϱον), dersom den paa nogen Maade blot skulde betegne en Potenseren indenfor Identiteten saa er der netop ingen Myndighed. Dersom saaledes en Lærer er sig begeistret bevidst at han selv existerende udtrykker og har udtrykt med Opoffrelse af Alt den Lære han forkynder: saa kan denne Bevidsthed vel give ham en vis og en fast Aand, men det giver ham ikke Myndighed. Hans Liv som Beviis for Lærens Rigtighed er ikke det Andet (το ετεϱον) men er en simpel Fordoblelse. Det at han lever efter Læren beviser ikke at den er rigtig; men fordi han selv er overbeviist om Lærens Rigtighed derfor lever han efter den. Derimod, hvad enten en Politie-Officiant for Exempel er en Slyngel eller han er en retskaffen Mand, saasnart han er i Function har han Myndighed.
\end{quote}
\dcite{BOA:219  (SKS 15, 219)}
% BOA:219  (SKS 15, 219)

\begin{quote}
Sagen er imidlertid denne. Tvivlen og Vantroen, der gjør Troen forfængelig, har blandt Andet ogsaa gjort Menneskene generede ved at lyde, ved at bøie sig under Myndighed. Denne Oprørskhed sniger sig ind selv i de Bedres Tankegang, dem maaskee ubevidst, og saa begynder alt dette Forskruede, som i Grunden er Forræderie, om det Dybe og det Dybe, og det Vidunderlig-Deilige, som man kan skimte og saa videre. Skulde man derfor med eet bestemt Prædikat betegne det christelig-religieuse Foredrag som det nu høres og læses: da maa man sige, det er affecteret. Naar man ellers taler om en Præsts Affectation, tænker man maaskee paa, at han pynter og sminker sig, eller at han sødligt smægter med Stemmen, eller at han norsk snurrer paa R og rynker Brynet, eller at han anstrænger sig i Kraftstillinger og i Opvækkelsens Spring og saa videre. Dog alt Sligt er mindre vigtigt, om det end altid var ønskeligt, at det ikke var. Men det Fordærvelige er, naar Prædike-Foredragets Tankegang er affecteret, naar dets Rettroenhed naaes ved at lægge Accenten paa et reent forkeert Sted, naar det i Grunden opfordrer til at tro paa Christus prædiker Troen paa ham paa Grund af Det, som slet ikke kan blive Troens Gjenstand. Dersom en Søn vilde sige »jeg lyder min Fader, ikke fordi han er min Fader, men fordi han er Genie, eller fordi hans Befalinger altid ere dybsindige og aandrige«: saa er denne sønlige Lydighed affecteret. Sønnen accentuerer noget reent Forkeert, accentuerer det Aandrige det Dybsindige ved en Befaling, medens en Befaling netop er ligegyldig mod denne Bestemmelse. Sønnen vil lyde i Kraft af Faderens Dybsindighed og Aandrighed; og i Kraft deraf kan han netop ikke lyde ham, thi hans critiske Forhold i Retning af om nu Befalingen er dybsindig og aandrig, underminerer Lydigheden. Og saaledes er det ogsaa Affectation naar der er saa megen Tale om at \tlg{tilegne} sig Christendommen og tro Christus formedelst Lærens Dybsindighed og Dybsindighed. Man tillyver sig Rettroenheden ved at accentuere noget aldeles Forkeert. Den hele moderne Spekulation er derfor affecteret ved at have afskaffet Lydigheden paa den ene Side og Myndigheden paa den anden Side, og ved saa desuagtet at ville være rettroende. En Præst der er aldeles correct i sit Foredrag maa tale saaledes, idet han anfører et Ord af Christus: »dette Ord er af Den, hvem, ifølge hans eget Udsagn, al Magt er given i Himlen og paa Jorden. Du maa nu, min Tilhører, overveie med Dig selv, om Du vil bøie Dig under denne Myndighed eller ikke, antage og tro Ordet eller ikke. Men vil Du det ikke, da gaae for Gud i Himlens Skyld ikke hen og antag Ordet fordi det er aandrigt eller dybsindigt eller vidunderlig-deiligt, thi dette er Gudsbespottelse, det er at ville criticere Gud.« Saasnart nemlig Myndighedens, den specifik-paradoxe Myndigheds Dominant er sat i, saa er alle Forhold qvalitativt forandrede, saa er den Art \tlg{Tilegnelse}, som ellers er tilladelig og ønskelig, en Brøde og Formastelighed.
\end{quote}
\dcite{BOA:224  (SKS 15, 224)}
% BOA:224  (SKS 15, 224)

\begin{quote}
Magister Adler er nu i den Alder, da man ordentligviis føler en Trang til at slutte sine Læreaar af, for derpaa at lære Andre. Der pleier da ogsaa i den Alder, der er Modenhedens critiske Alder, at indtræde en Trang til ret dybt at besinde sig paa sit eget Liv. Endnu engang, inden man gjør Overgangen og gaaer videre i Livet, vender man tilbage til sine første Erindringer, til Opdragelsens uforglemmelige Indtryk. Man prøver, hvorledes man nu forholder sig til hvad man den Gang barnligt forstod og \tlg{tilegnede} sig, man prøver om man er i Overeensstemmelse med sig selv, om og hvorvidt man forstaaer sig selv i at forstaae sit Første, og Bekymringen er, at Ens Liv i dybere Forstand maatte være et personligt Liv i væsentligt Sammenhæng med sig selv. Een og Anden staaer vel her paa Skilleveien, om han skal slippe det Første, hugge Broen af, og holde sig til det senere Lærte, eller om han skal søge tilbage til Barndommen og lære omvendt; thi som Barn lærte han jo af Ældre, og nu skulde han altsaa selv Ældre lære af et Barn, lære af sin Barndom. Havde det nu været Magister As Tilfælde, at han, idet han saaledes i sig selv vendte tilbage til sig selv, var stødt paa een eller anden væsentlig christelig Erindring, eet eller andet afgjørende Indtryk af det Christelige: da vil Sagen være bleven en anden, og itide noget alvorligere. Men hans Livs-Udvikling var netop en saadan, at det maa falde ganske naturligt, at han culminerede i at blive Hegelianer, aldeles i samme Forstand som den hegelske Philosophie var Menneskehedens høieste Udvikling. At den hegelske Philosophie derved igjen var Christendommens Culmination, kunde Magister A. ingen Betænkelighed have ved at antage; det fulgte vel snarere ganske af sig selv, og deraf da igjen, at han som Hegelianer var en ualmindelig udviklet og dannet Christen, vel skikket til at lære den hegelske Philosophie fra sig, og altsaa ogsaa Christendom, da disse Bestemmelser vare identiske.
\end{quote}
\dcite{BOA:252  (SKS 15, 252)}
% BOA:252  (SKS 15, 252)

\begin{quote}
Jeg vil tænke mig – thi det er min Sjel imod at blande mig ind i noget andet Menneskes Guds-Forhold, endog blot saaledes, at jeg var vidende om, at der levede en saadan Mand – jeg vil tænke mig en Familiefader. Han er dygtig i sin Livsstilling, behagelig i Omgang, ikke uvittig, vel lidt af Alle, gjestfrie og selskabelig i sit Huus; han er ikke nogen Religionsspotter (dette var om det end er farefuldt maaskee endda bedre for Børnene, thi det giver Elasticitet); han er heller ei i strængeste Forstand Indifferentist: men han er saadan Christen, han vilde finde det underligt, søgt og paafaldende den Sag betræffende udvortes eller indvortes at gjøre videre Ophævelser; han er saadan Christen, og er gjennem Læsning af eet og andet nyere Skrift i Forstaaelse med den Betragtning, at Christendommen er Aandens høieste Dannelse, samt at enhver Dannet er Christen. I sit Husliv, hvad enten der er Selskab eller ikke, faaer han aldrig nogen Anledning til at yttre sig religieust om religieuse Anliggender. Hænder det en sjelden Gang, at een eller anden excentrisk Religieus tildrager sig nogen Opmærksomhed, og bliver Dagens Gjenstand, og derved ogsaa en Gjenstand for Omtale i hiin Familiefaders Huus, da er Dommen over en Saadan og over hans Adfærd ikke nogen religieus men en æsthetisk, der stempler en saadan Tone og et saadant Væsen som anstødeligt, som Noget, der ikke kan taales i de Dannedes Kreds. – Denne Mands Hustru er en elskværdig Kone, fri for alle moderne uqvindelige Nykker, hjertensgod som Moder og Hustrue. Hun har ogsaa i et enkelt Øieblik af sit Liv følt en Trang til en dybere religieus Besindelse. Men da saadanne Anliggender og saadanne Bekymringer aldrig har været bragt paa Bane mellem hende og hendes Mand, og da hun har en maaskee overdreven Respekt for hendes hende overlegne Mands Fordringer i Henseende til Dannelse: saa frygter hun, at det skulde forraade en Mangel paa Dannelse, eller støde, om hun vilde tale om religieuse Gjenstande. Derfor tier hun, og ordner sig med qvindelig Hengivenhed ganske ind under hendes Mand, \tlg{tilegner} sig klædeligt, hvad der mandligt anstaaer ham saa godt; og disse Tvende harmonere som sjeldent et Ægtepar i Christenheden.
\end{quote}
\dcite{BOA:288  (SKS 15, 288)}
% BOA:288  (SKS 15, 288)

\begin{quote}
Under disse Forældres Øine voxe Børnene op; der spares Intet paa at udvikle deres Aand og berige den med Kundskaber, og medens Børnene saaledes voxe i Kundskab og Indsigt \tlg{tilegne} de sig ubevidst Forældrenes dannede Væsen, saa denne Familie virkelig er et yderst behageligt Huus at komme i. Det følger ganske af sig selv at Børnene ere Christne, de ere jo fødte af christne Forældre, og saa er det jo lige saa naturligt som at En er Jøde, der er født af jødiske Forældre, – hvor i al Verden skulde Børnene falde paa Andet. At der er Jøder, Hedninger, Muhamedanere, Fetischdyrkere til, veed de jo rigtignok af Historien og den videnskabelige Religions-Underviisning; men veed det ogsaa som noget dem aldeles Uvedkommende. Lad os saa tage det ældste af Børnene, han er nu i den Alder at han skal confirmeres. Det følger ganske af sig selv, at den Unge svarer »ja« paa de Spørgsmaal som Præsten forelægger ham; hvor i al Verden skulde der vel falde den Unge Andet ind, har man nogensinde hørt, at Nogen har svaret nei, eller er der vel ymtet et eneste Ord til den Unge om, at han ogsaa kan svare nei. Derimod har man maaskee gjort ham opmærksom paa, ikke at svare for høit og heller ikke for sagte, men gjøre det paa en klædelig og beleven Maade; han erindrer maaskee at have hørt Faderen sige, at det var noget Affecteret noget med dette hviskende »Ja« paa Kirkegulvet. I den Grad følger det af sig selv at han skal svare ja, i den Grad, at han istedenfor at gjøres opmærksom paa Svaret, er bleven gjort opmærksom paa Formalitetens reent æsthetiske Side. Saa svarer han da »ja« – hverken for høit eller for sagte, men med den frimodige og dog beskedne Sømmelighed, der klæder et ungt Menneske vel. Faderen er paa Confirmationsdagen noget alvorligere end sædvanligt, dog har hans Alvorlighed mere en festlig end en religieus Charakteer, og den harmonerer derfor ganske med den Oprømthed, der indtræder efter at de fra Kirken er kommen hjem, hvor Faderen ikke blot er behagelig som ellers, men anvender sine Gaver for ret at gjøre denne Dag til en Høitidsdag. Moderen er bevæget, hun har endog grædt i Kirken; men moderlig Ømhed, og verdslig Bekymring for Barnets fremtidige Skjebne er dog vel ikke forbeholdt Christendommen. Den Unge faaer derfor vel paa Confirmations-Dagen et høitideligere Indtryk af Faderen, et rørende Indtryk af Moderen, han vil med Taknemlighed og Glæde mindes denne Dag som en skjøn Erindring; men han faaer ikke noget afgjørende christeligt Indtryk. Den Unge faaer Indtrykket af at denne Dag maa være et betydningsfuldere Moment i hans Liv – om ikke saa betydningsfuldt som for Exempel det at blive Student og immatriculeres.
\end{quote}
\dcite{BOA:289  (SKS 15, 289)}
% BOA:289  (SKS 15, 289)

\subsection{SFV \normalfont(SKS 16)}

\begin{quote}
»Dog fandt han ogsaa her i Verden hvad han søgte: var ingen Anden det, han var selv »hiin Enkelte«, og blev det mere og mere. Det var Christendommens Sag han tjente, hans Liv fra Barndommen af forunderligt dertil anlagt. Saa fuldkommede han Reflexionens Gjerning, at sætte Christendommen, det at blive Christen, heelt og holdent ind i Reflexion. Hans Hjertes Reenhed var: kun at ville Eet; hvad der i levende Live var de Samtidiges Anklage mod ham, at han ikke vilde slaae af, ikke give efter, netop det Samme er Eftertidens Lovtale over ham, at han ikke slog af, ikke gav efter. Men det storartede Foretagende bedaarede ham ikke; medens han dialektisk qua Forfatter holdt Overskuelsen over det Hele, forstod han christeligt, at det Hele for ham betydede, selv at opdrages i Christendom. Den dialektiske Bygning, han fuldførte, hvis enkelte Dele allerede ere Værker, kunde han ikke \tlg{tilegne} noget Menneske, endnu mindre vilde han \tlg{tilegne} sig den selv; skulde han have \tlg{tilegnet} Nogen den, var det blevet Styrelsen, hvem den dog alligevel Dag for Dag, Aar efter Aar var \tlg{tilegnet} af Forfatteren, der, historisk, døde af en dødelig Sygdom, men digterisk døde af Længsel efter Evigheden, for uafbrudt ikke at bestille Andet end takke Gud.«
\end{quote}
\dcite{SFV:76.2  (SKS 16, 76)}
% SFV:76.2  (SKS 16, 76)

\begin{quote}
Kjere! Modtag denne \tlg{Tilegnelse}; den gives hen ligesom iblinde, men derfor ogsaa uforstyrret af nogen Hensyn, i Oprigtighed! Hvo Du er, veed jeg ikke; hvor Du er, veed jeg ikke; hvilket Dit Navn er, veed jeg ikke. Dog er Du mit Haab, min Glæde, min Stolthed, i Uvidenheden min Ære.
\end{quote}
\dcite{SFV:85.4  (SKS 16, 85)}
% SFV:85.4  (SKS 16, 85)

\section{Journals, notebooks and papers}

\subsection{1833}

\begin{quote}
Andet Hovedafsnits

\emph{ 1ste Afdeling. }

om Troen.

§ 56.

Mennesket bliver deelagtig i Frelsen ved Christus ved Troen eller ved Hengiven af det hele menneskelige Væsen i al dets Tænken, Følen og Villen til Christus som den, i hvem den guddommelige Sandheds og Naades Fylde er aabenbaret. Denne Stemning udgaar fra Erkjendelsen saavel af egen Syndighed og egen aandelig Trang, som af den høiere Tilfredsstillelse i Aabenbaringens Sandheder, og Ordet »Tro« betegner derfor ofte i Skriften den blotte Fasthed i Overbeviisningen, enten i Almindelighed eller med Hensyn til Christi Lære og hans Person.

Den almindelige Sætning: Mc: 16, 16. Joh: 3, 16. 6, 40. Act: 10, 43. Rom: 3, 25.28. Gal: 2, 16. Eph: 2, 8. Gjenstanden for Troen angives Rom: 3, 25. 4, 24.

Naar Frelsen ved Christus tilregnes Mennesket umiddelbart, uden at Troen nævnes, maa den dog i Tanken suppleres: Rom: 5, 10. 2 Cor: 5, 21.

§ 57.

Ved Tro ikke ved Gjerninger bliver Mennesket retfærdiggjort for Gud det vil sige Guds Naade kan ei erhverves ved menneskelige Gjerninger, men grundet i Guds eget Væsen, staaer Adgang til den aaben for enhver, som ved Christus har Fortrøstning til Faderens evige Kjærlighed. Den paulinske Lære herom er indeholdt i Christi Lære om Guds Natur og om Troens Virksomhed. Det er kun ved Misforstaaelse, at den er bleven anseet for stridende mod Apostlens Jacobs Lære om den indbyrdes Uadskillelighed af Tro og Gjerninger, og at den høie Trøst som indeholdes i den er bleven skadelig for den christne Dyd.

Om Betydningen af διϰαιουςϑαι og διϰαιοσυνη. \emph{διϰαιουν } betyder ei blot at gjøre retfærdig; men og erklære retfærdig. Tob: 3, 2. Sir: 10, 28. – Luc: 7, 35. – διϰαιουν εαυτον Luc: 10, 29. 16, 15. Saaledes og naar det hedder χϱιςτος eller ϑεος διϰαιουται Luc: 7, 29. 1 Tim: 3, 16.

Allene hos Paulus findes denne Sætning, hvorfor det kunde synes, at den ikke er oprindelig christelig, at han forsaavidt er afvegen fra den christne Lære, som han har lagt mere i den, end der oprindelig tilhørte den. – Men de christne Grundideer, som Paulus fuldstændigen har udviklet, ere: \emph{1.} at ingen Gjerning i og for sig selv betragtet kan ansees for Udtrykket af det dydige Sindelag, og altsaa ei Gjenstand for Guds Velbehag. \emph{2.} at Gud er ophøiet over enhver menneskelig Indvirkning, saa at Menneskene ligesaalidet ved deres Synder kan udelukke sig fra Guds Naade, som de ved gode Gjerninger kunne vinde den. \emph{3} at Troen paa Christus aabner Adgang til Guds Naade.

§ 58.

Ligesom den sande christne Tro er forbunden med Anger og en Omdannelse af Sjælens Virksomhed i Tanke og Villie, saaledes yttrer den sig endvidere som helliggørende Kraft, den henvender Menneskets Stræben til det høie Samfund med Gud og Christus, udvikler af Kjærlighed den rene christelige Dyd, af Haabet den sande Styrke til Opoffrelse for det Gode. Et saadant Liv er den synlige Forherligelse af Troen, Christendommens Øiemed og Veien for Mennesket til Guds Velbehag og den evige Salighed. –

Angeren maa med Hensyn til det forbigangne Liv være forbunden med Troen; men om denne er før hiin eller omvendt kan ei altid afgjøres. Den levende Erkjendelse og Følelse af Synden maa føre til Tro, og paa den anden Side maa Troen vække Angeren, saa at Mennesket, idet det bliver sig Guds Naade bevidst erkjender sin egen Uværdighed. Luc. 15. 2 Cor: 7, 10. – Ligesom Troen ei kan tænkes virksom uden Anger, saaledes er Omvendelsen ogsaa forbunden med den sande Tro. Tænker man sig nemlig et Menneske, i hvis Sind Kjærligheden er traadt i Stedet for Frygt, da vil hans Virksomhed, som tilforn var rettet udad, blive vendt indad, saa at frivillig Lydighed vil afløse den tvungne, forfængelig Stolthed og Hovmod af fortjenstlige Gjerninger vil afløses af den med Ydmyghed ledsagede Bevidsthed, at Ingen kan opfylde den guddommelige Lov. Denne Omdannelse er fremsat under forskjellige Billeder: συναποϑνηςϰειν τῳ χϱιςτῳ, αναγενναςϑαι, συνεγειϱειν συν τω χϱ., παλιγγενεσια, αναϰαινωςις του πνευματος, ανανεουςϑαι τῳ πνευματι, μεταμοϱφουςϑαι, αποδυεςϑαι, αποϑεςϑαι, σταυϱουν τον παλαιον ανϑϱωπον, ενδυεςϑαι τον ϰαινον ανϑ:, σταυϱουν την σαϱϰα, τας πϱαξεις του σωματος, ἡ ϰαινη ϰτιςις, αποϰυειν, γενηϑηναι ανωϑεν. Joh: 3, 3.5. Rom: 6, 6.7. 7, 6. 8, 13. 12, 2. 2 Cor: 3, 18. 5, 17. Gal: 4, 19. 5, 24. 6, 15. Eph: 2, 5.15. 4, 22. Coll: 3, 5.9. Tit: 3, 5. επιστϱεφεςϑαι απο σϰοτους εις φως, εγειϱεςϑαι εξ υπνου, ανοιγειν τους οφϑαλμους, αποϑεςϑαι τα οπλα, eller εϱγα του σϰοτους, ενδυεςϑαι τα οπλα του φωτος. Mth: 18, 3. Act: 14, 15. 26, 18. Rom: 13, 12. Eph: 5, 8. 1 Thess: 5, 4.6. Col: 1, 13.

I Skriften hedder det, at ved Troen meddeles Mennesket den hellige Aand, som tager Menneskets Væsen i Besiddelse og bliver det virksomme Princip i det. το πνευμα αγιον της πιςτεως, δια της πιςτεως. 2 Cor: 4, 13. Gal. 3, 2. 5, 14. Eph: 1, 13. Ved Troen boer Guds Aand i Mennesket 1 Cor: 3, 16. 6, 19. Ved den hellige Aand indvies Mennesket til sand Gudsdyrkelse, saa at dets hele Liv bliver en fortsat aandelig Gudsdyrkelse. 1 Cor: 3, 16. 6, 19.20. 2 Cor: 6, 16. Eph: 1, 19. Rom: 8, 14. Eensbetydende hermed er, at Mennesket skulle stræbe at \tlg{tilegne} sig Christi Aand. Rom: 13, 14. Gal: 3, 27. Christus uddannes i de Christne: Gal: 4, 19. Rom: 6, 5. 8, 29. Gal: 2, 20. Eph: 3, 17.

Kjærligheden indeholder ikke blot Begrebet af en Følelse; men beskrives som den høieste moralske Fuldkommenhed. Mth: 22, 40. Joh: 15, 12. Rom: 13, 9.10. 1 Cor: 13. Gal. 5, 14. 1 Tim: 1, 5. Jac: 2, 8.

Kjærlighedens Væsen bestaaer deri, at Mennesket omfatter det, som ligger udenfor det, og gjør det til Gjenstand for samme Omhu, som det, der er i det selv, eller at Mennesket bliver Eet med Gud. Det Modsatte er det egoistiske, som alene har Agt paa sig selv. Egoisme udspringer af sandselig Attraaer, der stræbe at \tlg{tilegne} sig det Gode for sig alene. – Endvidere skildres Bestandigheden i det Gode, Haabet, som en umiddelbar Følge af Troen: Mth: 10, 22.28. Joh: 16, 33. Rom: 5, 2-5; 8, 25. 12, 12. 15, 13. 1 Cor: 13, 13. Col: 1, 23. Phil: 4, 13. Heb: 10, 35. 12, 1.2.

Den christne Dyd fremstilles som sidste Øiemeed for Christendommen og Særkjende for dens Bekjendere.

Derpaa skulle nemlig Christus Disciple kjendes, at de elske hverandre, at de opfylde Guds Bud, at de efterfølge Christus, at de vare hellige som Gud, at de levede ikke for sig; men for ham: Mth: 7, 21. 12, 50. Joh: 15, 14-17. Rom. 6, 13-22. 2 Cor: 5, 14. Gal: 1, 4. Eph: 1, 4. 2, 10. Tit. 2, 14. 1 Joh: 1, 7. Dette oplyses indirecte, idet de beskrives som Fornægtere af Christus, hvis Vandel ei svare til deres Overbeviisning: αϱνειςϑαι τον χϱιςτον τοις εϱγοις. Tit: 1, 16. βλαςφημειν το ονομα του ϑεου: Rom: 2, 24. 1 Tim: 6, 1. Tit: 2, 5. εχοντες μοϱφωςιν ευσεβειας. 2 Tim: 3, 5. αγαπαν γλωςςῃ 1 Joh: 3, 18. – Jac. 2, 16. Da Opnaaelsen af Christendommens Øiemeed er Opfyldelsen af den guddommelige Verdensplan, opfordres enhver Christen til at bidrage Sit til Guds Riges Fremme. Rom: 2, 6. Mth: 25, 34. Rom: 8, 1.13. 1 Cor: 6, 9. 2 Cor: 9, 7. Eph: 5, 5. 2 Tim: 4, 7.8. Heb: 12, 14. Jac: 1, 27. Joh: 3, 21. 4, 17.

§ 59.

Det var en Følge af at christelig Tro mere og mere forvexledes med blind Underkastelse af de christne Dogmedecreter, at Læren om Troens Værd og Kraft tabte sin Betydning i den christne Kirke. Heraf fulgte et Tilbagefald til den gamle Verdens Overtro om Gjerningers og Bodsøvelsers Kraft, og det laa i det hierarchiske Systems Interesse at bestyrke denne Betydning. I det catholske ved Reformationen fremsatte Lærebegreb, hvor Retfærdiggjørelse som moralsk Forandring i Mennesket selv falder sammen med Helliggjørelse, medens Troen alene betragtes som den foreløbige Betingelse for denne, og Gjerninger som et Middel til at forøge Retfærdigheden og fortjene Saligheden, kunde Afvigelse fra den evangeliske Lære synes mindre væsentlig, naar denne ikke faldt saa meget stærkere i Øinene ved de practiske Misbrug, som have deres naturlige Støttepunct i denne Theorie. –

\emph{Tridentinske Concilium 6te }\emph{Sessio:} justificatio non est sola peccatorum, sed et santificatio et renovatio interioris hominis, renovamur spiritu mentis nostræ, et non modo reputamur, sed vere justi nominamur et sumus, justitiam in nos recipientes unusquisque suam secundum naturam, quam spiritus sanctus partitur singulis prouti vult et secundum propriam cujusque dispositionem et cooperationem. –

Hvad Forholdet mellem Tro og Retfærdiggjørelse angaaer, hedder det: disponuntur homines ad justitiam, dum fidem ex auditu concipientes libere moventur in Deum, credentes, vera esse, quæ divinitus revelata et promissa sunt, et dum a divinæ justitiæ timore ad considerandam Dei misericordiam se convertendo in spem eriguntur, et moventur adversus peccatum per poenitentiam.« Conciliet har erklæret sig mod at Troen skulde være den tilstrækkelige Betingelse, men kun den foreløbige: »fides, nisi ad eam spes accedat et caritas, non unit perfecte cum Christo, si quis dixerit, sola fide impium justificari, et nulla ex parte necesse esse, eum suæ voluntatis motu præparari et disponi, anathema sit.

Om Forholdet, hvori Gjerninger staae til Tro og Retfærdiggjørelse, hedder det: »licet in hac vita quantumvis sancti et justi in levia saltem et quotidiana peccata cadant, non propterea desinunt esse justi. Si quis dixerit, justitiam acceptam non conservari, aut etiam non augeri per bona opera, sed ipsa opera fructus solum, et signa justificationis esse, vel justificatum bonis operibus non vere mereri augmentum gratiæ (vitam æternam et ipsius vitæ æternæ consecutionem) atque gloriæ augmentum, anathema sit.«

§ 60.

I Modsætning til Forestillingerne, af hvilke de fordærvelige Følger vare øiensynlige i den catholske Kirke, blev det Reformatorernes vigtigste og fortjenstfuldeste Værk at kalde den bibelske Lære til Liv om Retfærdiggjørelse som Menneskenes Naadestand, grundet ei i menneskelige Gjerningers Fortjeneste, men i Guds uendelige Naade og Christus Fortjeneste, eller fra Menneskenes Side ved Troen eller den klare og faste Fortrøstning til Forløsningen ved Christus, og om de uudeblivelige Følger eller Virkninger af denne Tro, frivillig Lydighed mod Guds Lov og Flittighed til gode Gjerninger, der saaledes ere Troens Prøvesteen, og med al deres Ufuldkommenhed velbehagelige for Gud. Den rigtige Opfattelse og Anvendelse af denne Lære betrygges kun ved oplyst christen Forstand og levende religieus Følelse. De Misforstaaelser, der yttrede sig allerede i Reformationperioden, vise, hvor skadelig her enhver polemisk Eensidighed virker; og den sørgelige Maade, paa hvilken denne Lære ofte bruges i den evangeliske Kirke til at stille Gjerningerne i Modsætning til Troen, og fremkalde et tomt Gudfrygtigheds Skin, opfordrer til en samvittighedsfuld Omhyggelighed, navnlig i den populaire Udvikling. –

I deres Opposition imod Catholicismen fremhævede Reformatorerne især den practiske Side: Confessio Augustana »olim vexabantur conscientiæ doctrina operum, non audiebant ex evangelio consolationem, quosdam conscientia expulit in desertum, in monasteria, sperantes, ibi se gratiam merituros esse per vitam monasticam. Alii alia excogitaverunt opera ad promerendam gratiam et satisfaciendum pro peccatis; ideo magnopere fuit opus, hanc doctrinam de fide in Christum tradere et renovare, ne deesset consolatio pavidis conscientiis. Tota hæc doctrina ad illud certamen perterrefactæ conscientiæ referenda est, neque sine illo certamine intelligi potest.«

§ 61.

For at anskueliggjøre den fremskridende Udvikling af den Christnes nye aandelige Liv, efter psychologiske Love, har man i det theologiske System benyttet de billedlige Udtryk i Skriften til at abstrahere forskjellige Acter af en saa kaldet ordo elleroeconomia salutis. Denne bestemte Adskillelse og Gradation har imidlertid bidraget til at indtvinge det religieuse Liv i unaturlige Baand, hvor man saaledes, som i de pietistiske Skoler, ei har skjelnet mellem den vidensk: Abstraction og Anvendelse paa det virkelige Liv. –

En Antydning heraf findes i Luthers liden Cathecismus Spiritus sanctus me per evangelium vocavit, suis donis illuminavit, in recta fide sanctificavit et conservavit.« I senere dogmatiske Skrifter blev Terminologien bestemt. Det første Moment er 1) \emph{vocatio.} 2) \emph{illuminatio } Joh: 14, 21. 1 Cor: 2, 13. Eph: 1, 17. Denne illuminatio har man kaldet mediata, for at adskille den fra immediata eller Inspiration. 3.) \emph{conversio } eller\emph{poenitentia } \emph{. } Conf: Aug: »Constat poenitentia duabus partibus, altera est contritio, sive terrores incussi conscientiæ, agnito peccato, altera est fides, quæ concipitur ex evangelio seu absolutione, et credit propter Christum remitti peccata, et consolatur conscientiam, et ex terroribus liberat. Deinde sequi debent bona opera, quæ sunt fructus poenitentiæ, hoc est, mutatio totius vitæ ac morum in melius.« – Det 4) fierde Moment er \emph{santificatio } αγιαςμος. 5) unio mystica. Joh: 14, 23. 1 Joh: 4, 12. Eph: 3, 17.

Andet Hovedafsnits

\emph{ 2den Afdeling. }

Om Naaden.

§ 62.

Ligesom vi prise i Christi Sendelse til Verden, Fuldendelsen af Guds faderlige Naade mod det menneskelige Kjøn, saaledes er det Forhold, i hvilket det enkelte Menneske træder til Christum, til de ved Christus. skjænkede aandelige Goder, den ved Christus. forhvervede Salighed, et Værk af Guds almægtigt virkende Naade. Skriften betegner dette Forhold ved Navnene: Kaldelse (som nærmest udtrykker Meddelelsen af Christendommens Evangelium) og Udvælgelse (som udtrykker den troende Annammelse af Evangeliet, og den derved bevirkede Tilstand af Fromhed og Lydighed, Glæde og Fred med Gud). De billedlige Udtryk, ved hvilke Gud eller den hellig Aand siges at virke, fuldføre og bevare det Gode i Mennesket, indeholde en Advarsel mod at ville udgrunde Virksomhedens Natur og Maade, medens vi saavel ved bestemte Yttringer, som ved Christendommens hele Aand, paa eengang henvises til i Ydmyghed at erkjende det Gode i os, virket ved den Høiestes Velsignelse, og opfordres til at nytte Naadens Vink og Tilskyndelse ved Anvendelsen af egen Kraft og Evne.

ϰλητος bruges om de Christne: Mth: 20, 16. 22, 14. I Brevene bliver ϰαλειν, ϰληςις, ϰλητος brugt enten absolut om de Christne, eller forbundet med andre Ord for Exempel til Guds Rige εις την ϰοινωνιαν του Christus:, εις ζωην αιωνιον. Rom: 8, 30. 9, 24. 1 Cor: 1, 9. 1 Thess: 2, 12. Eph: 4, 4. 1 Tim: 6, 12. 1 Pet: 1, 15. Heb: 3, 1. I en mere udstrakt Betydning forekommer ϰληςις om den Plads, det enkelte Menneske indtager i det hele Samfund: 1 Cor: 7, 15.24. Ordet εϰλεγεσϑαι har samme Betydning som det hebr: \hebr{בחר}‎. ϰαλειν og εϰλ: bruges vel undertiden som synonyma; men ϰληςις betegner dog nærmest Kundgjørelsen eller Meddelelsen af Evangeliet εϰλεγεςϑαι den Indgang, som Evangeliet har fundet: 1 Cor: 1, 27. sq. Jac: 2, 5. Eph: 1, 4. 2 Thess: 2, 13. Mth: 24, 22.24.31. Rom: 8, 33. Col: 3, 12.

Naar saaledes de Christne siges kaldede og udvalgte, hvorledes beskrives da Bistanden af den guddommelige Naade, ved hvilken Kaldelsen siges at være bevirket? Paa mange Steder bliver Virkningerne betingede ved den christne Lære, saa at denne fremstilles som Vehiklet, Organet, hvorved Gud virker paa Menneskene Her bliver saaledes Troen og Omvendelsen, som Virkninger af Evangeliet, ikkun indirecte henførte til Guds Naade. \emph{Rom: 8, 2. 10, 14.17.} Jac: 1, 18.21. 1 Tim: 4, 16. 1 Pet: 1, 23. 2 Pet: 1, 3. Her omtales jo altsaa kun den Guds Virksomhed, som virker gjenem den christne Lære, der jo ikke virker paa anden Maade end enhver anden Lære (i det mindste er der ingen specifisk Forskjel); men der udfordres en vis Receptivitet, Evne og Villie hos Menneskene Disse Betingelser tilskrives atter i den hellige Skrift Gud eller Guds Naade. Gud gjør ligesom Mennesket selv disponeret til at modtage Læren. Saaledes hedder det i Skriften, at Sands for Sandheden, som maa gaa forud for Omvendelsen og Troen, skjænkes af Gud: \emph{1 Cor: 3, 6.} 1 Cor: 1, 4.5. 2 Cor: 4, 6. \emph{Eph: 1, 17.} 2 Tim: 2, 25. Fremdeles siges Troen at bevirkes ved Gud, uden hvis Virksomhed Ingen siges at kunne komme til Christum. Joh: 6, 44. Endelig siges ogsaa Troens Frugt at være bevirket ved Gud. \emph{Rom: 15, 13. 1 Cor: 1, 8.} 2 Cor: 1, 21. \emph{Eph: 3, 16.} Phil. 1, 6. 2, 13. \emph{1 Thess. 3, 12.} 2 Thess. 2, 17. 1 Pet. 1, 5. 5, 10.

§ 63.

Paa Spørgsmaalet, hvorledes det kan forenes med Evangeliets Bestemmelse, at det ikke er bleven alle Mennesker meddelt, henviser Skriften til den samme urandsagelige Viisdom, som først efter Aartusinders Forløb har ladet Lyset opgaae for Jordens Slægter, og som i Tidens Fylde vil føre Alle til saliggjørende Erkjendelse af Sandheden. Paa samme Maade maatte Svaret falde ud, naar spurgtes om Grunden til, at Evangeliet saa ofte er bleven forkyndt, og endnu ofte forkyndes uden at finde Indgang eller yttre sin Kraft. Ved her at gaae ud over de nærmere mellemliggende Aarsager paanøder sig Ideen om en Eenhed af den guddommelige Naade og menneskelige Frihed, som lige nødvendig og ubegribelig, og idet Paulus udleder Christendommens Virkninger paa det ene Menneske fremfor paa det andet af den guddommelige Styrelse (Udvælgelse, Anordning, Forudbestemmelse), og paa samme Tid udhæver den menneskelige Selvbestemmelses Evne paa det stærkeste, antyder han denne Eenhed, men tillige Mysteriets Ufattelighed for den reflecterende Tænkning, og afskjærer Veien for enhver Slutning, som udledes af eensidig Forfølgelse af den ene eller den anden Sætning. –

Hensigten af Christi. Sendelse er at frelse Alle: 1 Tim: 2, 4. Joh: 3, 16. Rom: 3, 29.30. 10, 12. 2 Cor: 5, 14. 1 Joh: 2, 2. Tit: 2, 11. Hermed synes nu de i § anførte Sætninger at stride. Christendommen omtales som μυστηϱιον σεςιγημενον, αποϰεϰϱυμμενον πϱο χϱονων αιωνιων, φανεϱωϑεν εν τω πληϱωματι του χϱονου: Rom: 16, 25-26. 1 Cor: 2, 7. Gal. 4, 4. Eph: 3, 9. 2 Tim: 1, 9.10. Tit. 1, 2.3. De to i § omtalte Phænomener fremstilles som grundede i Guds urandsagelige Villie: Eph: 1, 9. Eph: 3, 3. Col. 1, 27.

Paulus fremstiller den Virkning Evangeliet yttrer paa de Enkelte som Virkning af Guds Forudviden: πϱογινωσϰειν, πϱοὁϱιζειν, ταττειν, εϰλεγεςϑαι, αιϱειν. Rom: 8, 29.30. Eph. 1, 4.5.11. 2 Thess. 2, 13. 1 Pet. 1, 2. Act: 13, 48. Men paa den anden Side fremstilles Troen og Omvendelsen som Virkning af Menneskets egen Selvbestemmelse: Mth: 4, 17. 7, 13. Rom: 12, 1. Eph: 6, 10. Phill: 2, 12.

§ 64.

Skriften beroliger os med Hensyn til Følgerne af Evangeliets endnu stedse indskrænkede Virksomhed; thi naar den fremstiller os den guddommelige Naade ei som bevirket ved Christendommen; men denne som den fulde Aabenbarelse af den evige under forskjellige Former virkende Naade –, da kan den ei tænkes indskrænket til dem, der til en vis Tid ere kaldede til Andeel i Christendommen. Det følger endvidere, at den guddommelige Raadslutning til Frelsens Fuldendelse ved Christendommen maa tænkes i sig selv een, almindelig, omfattende Alle. Udtryk som Udvælgelse, Forskydelse kunne derfor kun betegne Overgangstilstande, en Adskillelse ifølge de forskjellige Forhold til Evangeliet, som til en given Tid finder Sted, men som tilligemed sine Følger vil ophøre. –

\emph{Om Hedningernes Fordømmelse.}

§ 65.

Den Augustinske Theorie om Menneskenaturens fuldkomne Fordærvelse og Kraftløshed maatte ved en umiddelbar Consequents føre til stræng Modsætning af Natur og Naade og derved til en Række af Misforstaaelser og Overdrivelser. Naar al Frihed til det Gode tænktes uddød, maatte 1) Omvendelse og Tro tænkes udledt af en overnaturlig \emph{uimodstaaeligt} virkende Naade, som 2) \emph{var indskrænket til de Christne;} og Grunden til denne 3.) \emph{Naades Omfang og Virksomhed maatte søges i en evig og ubetinget Forudbestemmelse,} ved hvilken nogle bleve udfriede af den almindelige Fordømmelse, og ved Guds Naade førte til Tro og Salighed. Den strenge Consequents i denne Lærebygning, i Forbindelse med den Maade paa hvilken den henviser Mennesket til at søge sin Beroligelse alene i Fortrøstning til Gud, bragte til at oversee den Eensidighed, som sætter den i Strid med Skriften og det christne Livs Fordringer. Efter at have været i Middelalderen Gjenstand for flere Stridigheder, blev Augustins System i det Væsentlige optaget som privat Mening af \emph{samtlige Reformatorer;} men kun i en Deel af den reformerte Kirkes herskende Lære som Følge af Calvins Anseelse. I den catholske Kirke blev den strænge Augustinianisme under overveiende Modstand fornyet af Jansenister. –

\emph{Gotschalck } gik videre og fremsatte Læren om prædestinatio gemina electorum ad requiem – reproborum ad mortem (Imod ham Rabanus Maurus og Hinkmarus Rhemensis.)

§ 66.

Ligesom indtil Augustins Tid Forestillingerne om en fri Samvirken fra Menneskets Side med Guds Naade og om den guddommelige Forudbestemmelse, som grundet i Guds Forudviden, vare almindeligen antagne; saaledes er den græske Kirke stedse bleven disse tro. I den latinske Kirke bleve disse Forestillinger efter at Pelagianismen var banlyst, Grundvolden for det semipelagianske System, der i Middelalderen kan ansees som det herskende, og omsider i den catholske Kirke blev sanctioneret som Kirkelære, og opretholdt i flere dogmatiske Controverser. I den lutherske Kirke fandt Synergismen Forsvar hos Melanchton og hans Skole, og den augsburgske Conf: indskrænker sig uden skarpere positive Bestemmelser til at vise Nødvendigheden af Guds Naade til Omvendelse. Derimod urgerer Concordieformelen med Augustin Menneskets Naturs fuldkomne Kraftløshed; men forkaster derhos Prædestinationen, idet den fremstiller Naaden vel som almindelig; men tillige som modstaaelig fra Menneskets Side. Den er ogsaa heri bleven Grundvold for den senere lutherske Dogmatik, som forgjeves har søgt ved subtile Distinctioner at dække Systemets Inconsequents. I den reformerte Kirke har det strænge Calvinske Lærebegreb fundet Modstand, saavel i flere enkelte symbolske Bøger, som i den hollandske Kirke hos Arminianerne, og for størstedelen i den nyere Theologie. Derimod have de protest: Kirker været enige i at afvise de mystiske sværmærske Forestillinger om overnaturlige og umiddelbare Naadevirkninger, ligesom det ogsaa indtil de nyere Tider almindeligt lærtes, at den saliggjørende Naade var indskrænket til Christendommens Bekjendere.

Confessio Aug: 18de Artikel: humana voluntas non habet vim sine spiritu sancto efficiendæ justitiæ Dei seu justitiæ spiritualis; sed hæc fit in cordibus, cum per verbum spiritus sanctus concipitur. – »Damnant Pelagianos, et alios, qui docent, quod sine spiritu sancto solis naturæ viribus possimus Deum diligere super omnia, item præcepta facere quoad substantiam actuum.« 5te Artikel: \emph{Spiritus sanctus fidem efficit,} ubi et quando visum est Deo, in iis, qui audiunt evangelium.« – 20de Articel \emph{»per fidem efficitur spiritus sanctus«.}

Om det kirkelige Samfund

\emph{ 1ste Afsnit. }

\emph{ Om Kirkens Væsen og Øiemeed. }

§ 67.

Gjenstanden for Christus Virksomhed er ei det enkelte Menneske; men Menneskeslægten, hvis Vee og Vel er ethvert enkelt Menneskes Og Samfundet, som ved Christum er stiftet paa Jorden er i sin Natur og sin Bestemmelse angivet ved Christus. Liv og ved hans Lære. Den guddommelige Aabenbaring forener Menneskene ikke til fælleds Bestræbelse for timelige Formaal; men for at forene dem med Gud ved fælleds Sandheds Erkjendelse, fælleds Tro og Gudsfrygt. Dette Øiemeed er uafhængigt af de ydre Betingelser for enhver anden Samvirken, og saaledes er den Christne. Kirke ligesom Christus. Lære i Omfang ophøiet over enhver Nationalitets Indskrænkning, og i sin Tilværelse over enhver Omvexling af de menneskelige Ting. Dens Eiendommelighed betegnes derfor ved Navnet Guds Rige. –

Χstendommens Universalisme angives: Joh: 3, 16. 8, 12. Act: 10, 34. Eph: 2, 16. Col. 1, 19. Mth: 28, 19. – Kirken Vedvaren: Mth: 16, 18. 28, 20. 1 Cor: 15, 25.

§ 68.

Den christne Kirke forudsætter til Opnaaelse af sit Øiemed Eenhed i Troen, uden hvilken den aandelige Samvirken falder bort; og Frihed for den Enkelte til selvstændig Udvikling af egen Overbeviisning. I begge Retninger møder Kirken Modstand; men hvor Christus er i Sandhed Kirkens Hoved og Troens Grundvold, der kan Kraften ei fattes til at overvinde denne Modstand; thi af dette Forhold udvikler sig den christne Almeenaand, som bevarer Troens Eenhed i Fredens Baand, medens den beskjærmer dens Frihed mod Tvangens Baand. Denne Aand, tænkt som det Alt gjenemtrængende, i Alle virkende Princip betegner Kirkens Ideal, fra hvilken den synlige Kirke til enhver Tid er forskjellig, som Samfundet, der sammenholdes ved Evangeliet og Sacramenterne, og ved disse arbeider under guddommelig Styrelse til sit Maal. –

De Christne opmuntres til Eenhed: ομοφϱονειν, το ἑν φϱονειν. Rom: 15, 5. Phil: 2, 2. 1 Pet: 3, 8. Denne Eenhed betegnes netop som Eenhed i Troen: Eph. 4, 5.13. Gal. 3, 28. Spørges endvidere, hvad forstaaes ved denne Troes Eenhed, da svares herpaa i Joh. 17, 3. Heb. 6, 1. Paa den anden Side omtale ofte Apostlene Forskjel i Aandsgaver som nødvendige og gavnlige: Rom: 12, 3-6. Eph: 4, 7. 1 Cor: 12, 4. Ligeledes at Christus er Kirkens eneste Herre og Mester, saa at Menneskene ei skulle erkjende noget andet Herredømme i Troens Anliggender: Mth: 23, 8. Joh: 13, 13. Saaledes fremkommer den christne Frihed: 2 Cor: 3, 17. Tydeligst bliver dette Forhold af Eenhed og Frihed antydet i Lignelsen af Kirken som et Legeme: 1 Cor: 12. Col. 1, 18. Eph: 1, 22. 4, 15.16. 5, 23.30.

§ 69.

Allerede tidlig blev Kirkens Idee forvansket ved misforstaaet Stræben efter Eenhed, som bragte Traditionen i Skriftens Sted, Bekjendelsen i Troens, og fuldendte Hierarchiets Organisation. Saaledes udviklede Catholicismens Idee sig til Begrebet af Kirken som Samfundet, hvis Eenhed, Almindelighed, Hellighed og apostoliske Oprindelse er umiddelbar givet i den historiske Forbindelse med Christus og hermed tillige Ufeilbarhed i enhver dogmatisk Bestemmelse og kirkelig Indretning. Denne Miskjendelse af Forholdet mellem den stræbende og fuldendte Kirke maatte ophæve al christen Frihed og føre til et System af Troestvang. –

§ 70.

Den protestantiske Kirke sætter Særkjendet for Christi Kirke i Overeensstemmelse med Christi Aand. Den skjelner mellem den usynlige og synlige Kirke, som imellem Opnaaelsen af det ophøiede Øiemeed og Stræben derefter, og erkjender, at til ingen Tid i Samfundet saa lidet som hos den Enkelte aandelig Eensidighed og Feil kan undgaaes; men som det af Christus indstiftede Middel til Opnaaelsen af dette Øiemeed, er de Troendes Forening, sin Ufuldkommenhed uagtet, at ansee som den ene sande Kirke og Veien til at nærme sig Maalet er alle Enkeltes frie og selvstændige Samvirken. –

2det Afsnit.

Om Kirkens Virksomhed.

§ 71.

Kirkens Øiemed tilstæder ingen Anden Art af Virksomhed end aandelig Indvirken efter Christi og Apostlernes Forbillede. Midlerne hertil ere givne i Ordet og Sacramenterne, og Virksomheden af disse til at befordre Guds Riges Seier over dets Fjender, bliver efter Skriftens Lære betrygget ved det vedvarende Samfund med Christus., som Menighedens Herre og Beskytter. Sacramenterne virke ei umiddelbar ex opere operato; men forsaavidt de nyttes og anvendes med troende Sind.

\emph{Anmærkning } Ved denne Tro paa Christus. udelukkes den catholske Kirkes Paakaldelse af Helgenes Forbønner, hvorimod en Forbindelse af Kirkens Fortid med dens Nutid ved Beskuelse, Ihukommelse og Hylding af høie Forbilleder i Tro og christen Dyd kun kan tjene til at oplive og lede Udviklingen af et sandt christent Liv. –

Disse Midler ere: verbum divinum, sacramentum, potestas clavium. Disse Midler staa nu ogsaa i Forhold til Øiemedet; thi Christus. staaer i vedvarende Forbindelse med Kirken Dette har Christus. selv forjættet, Sætningen har Apostlene udført ved at fremstille Christus. som Menighedens Herre, Konge, de Troendes Hyrde: 1 Cor. 15, 25. Eb: 13, 20. 1 Pet. 2, 25. 5, 4. – Christus. fremstilles som siddende ved Guds Høire: Mc. 16, 19. Eph. 1, 20. Eb. 8, 1. Rom: 8, 34. 1 Pet. 3, 22.

§ 72.

Christi. Virksomhed bliver afbildet og fortsat, idet Kirken gjør den fri Forkyndelse og fleersidige Udvikling af det guddommelige Ord, efter den Anviisning, som er indeholdt i den hellige Skrift, til Middelpunct for sin Virksomhed. –

§ 73.

Saavel ved Christendommens hele Tendents som ved Skriftens Yttringer om et Præstedømme, fælleds for alle Christne, er ethvert Hierarchie udelukt af Kirken, som uforeneligt med de Troendes Forhold til Christus., hvorimod Nødvendigheden af en bestemt ordnet Virksomhed med Hensyn til Ordets Forkyndelse, og en Overdragelse af denne Virksomhed til enkelte Medlemmer efter Aandsgavernes Forskjellighed, ved Siden af den fri almindelige Meddelelse, er grundet i Samfundets Natur, og udtrykt i Christus. og Apostlenes Foranstaltninger. Det fulgde af de senere kirkelige Forhold, at en bestemt Udvikling og Dannelse maatte blive en nødvendig Betingelse for denne Virksomhed, og saaledes fremstod en geistlig Stand, som Organet, ved hvilket Christus. skal virke ved sit Ords Kraft. –

\emph{Anmærkning } Den catholske Lære om Ordinationens Natur og Geistlighedens Forhold til Christus og Menigheden er umiddelbart grundet i Læren om Kirkens Væsen. –

§ 74.

Som synlig Fremstilling af Evangeliet, der har Auctoritet og Virksomhed tilfælleds med det guddommelige Ord forudsætte Sacramenterne Indstiftelse af Christus. Den protestantiske Kirke forkaster derfor den catholske Kirkelære om 7 Sacramenter, ligesom den erklærer sig mod enhver Betragtningsmaade af de tvende Sacramenter enten som magisk virkende Tryllemidler eller som blot billedlige Betegnelser. I typologisk Betydning have Dogmatikerne antaget Sacramenter i det Gamle Testamente –

\emph{[i margen:]} Man har anseet denne Lære for skadelig for den christne Dyd. Men det beroer paa en Misforstaaelse af Ordet »Tro« og Gjerning. Det første har man taget blot i intellectuel Forstand; det andet har man brugt blot om de empiriske Phænomener. \emph{Augustinus:} inseparabilis est bona vita a fide, quæ per dilectionem operatur, im̄o ea ipsa est bona vita.« – Denne Frygt har sin Grund i at man overfører en menneskelig Maalestok paa Guddommelige Da vi nemlig ikke kan skue ind i Hjerterne maa vi see paa Gjerninger.

\emph{[i margen:]} I confessio augustana 20de Artikel bruges justificatio og remissio peccatorum. I apologia confessionis hedder det: consequi remissionem peccatorum et justificari. I Concordieformelen: vocabulum justificationis significat justum pronuntiare a peccatis et æternis peccatorum suppliciis absolvere. – I confessio augustana deus non propter nostra merita sed propter Christum justificat hos, qui credunt se propter Christus in gratiam recipi. – Troen har til sit Objekt remitti nostra peccata propter Christum.

\emph{[i margen:]} necessitas meriti og necessitas debiti. og det er mærkeligt, at formula concordiæ søgte at sætte begge Dele i Indifferents ved at erklære sig baade mod Gjerningers Nødvendighed og deres Skadelighed. –

\emph{[i margen:]} Schleiermacher gjør Gjenfødelse til Udgangspunkt for Inddelingen, den bestaaer af to Dele: Omvendelse med Hensyn til Livet, Retfærdiggjørelse med Hensyn til Gud.

\emph{[i margen:]} ved en hellig Kaldelse. 2 Tim: 1, 9.

\emph{[i margen:]} \emph{Melanchton } forandrede sig Noget. I Begyndelsen var han en Prædestinatianer, han siger i loci theol. 1ste Udgave: quandoquidem omnia, quæ eveniunt, necessario juxta prædestinationem divinam eveniunt, nulla est voluntatis nostræ libertas. – \emph{Calvin:} »æternum dei decretum, quo deus apud se constitutum habuit, quid de unoquoque homine fieri vellet. Non enim pari conditione creantur omnes, sed aliis vita æterna, aliis damnatio æterna præordinatur.« – »non ideo deum elegisse quosdam, quod prævidit illos futuros esse sanctos, sed sanctos factos esse, quia elegit illos.« –

\emph{[i margen:]} \emph{Supralapsarii } – \emph{Infralapsarii. }

\emph{[i margen:]} \emph{Michael Bajus.} fordømt ved en Bulle af Pius V af Aar 1567. – \emph{Ludvig Molina } – \emph{Leonhard Lessz. } Clemens d. 8de nedsatte congregationes de auxiliis gratiæ. 1597-1611. \emph{Cornelius Jansen } Biskop i Ypern udgav: Augustinus sive doctrina Augustini de humanæ naturæ sanitate, ægritudine, medicina adversus Pelagianos et Massilienses. – \emph{Paschasius Quesnel } udarbeidede en fransk Oversættelse af Bibelen med moralske Noter. –

\emph{[i margen:]} \emph{Synergisme} (Melanchton?). 1555. Pfeffinger imod ham: N. Amsdorff, og de Jenaiske Flaccius, Museus, Wigandt.

\emph{[i margen:]} I \emph{de dogmatiske Lærebøger} findes en Mængde Definitioner og Distinctioner. gratia præveniens; gratia operans; gratia cooperans. (universalis; resistibilis; amissibilis) – benevolentia universalis sive prædestinatio late sic dicta sive idealis og benevolentia specialis sive prædestinatio stricte sic dicta sive realis. –

\emph{[i margen:]} \emph{Universalisterne.}hypothetici; (Johan Cameron og Moses Amurand; Johan Dallæus og Ludvig Capellus) absoluti.

\emph{[i margen:]} Anabaptister. Schwenkfeldianer. Weigelianer. Böhmister. Quækere og Methodister. (1 Joh: 1, 6. 2 Pet. 1, 4. Eb: 3, 14. 6, 4. 1 Joh. 2, 20.27.)

\emph{[i margen:]} Concordieformelen lærer: 1) at den Hellig-Aand virker i Mennesket det hele Omvendelses Værk: Spiritus sanctus operatur in nobis illud velle et perficere. Den forkaster den Lære: liberum arbitrium Deo occurrere aliquo modo, etsi parum et languide ad conversionem suam conferre, eam adjuvare, cooperari, sese ad gratiam præparare et applicare, eam apprehendere, amplecti, evangelio credere. (Synergismen) 2) at Naaden er ved Christus tilbudt Alle uden Undtagelse: Christus. omnes peccatores ad se vocat et serio vult, ut omnes homines ad se veniant, et sibi consuli et subveniri sinant. 3) at der findes Forskjel mellem de Udvalgte og de Forskudte, eftersom den tilbudte Naade er bortstødt eller modtaget: ut Deus in æterno suo consilio ordinavit, ut Spiritus sanctus electos per verbum vocet, illuminet et convertat atque omnes illos, qui Christum. vera fide amplectantur, justificet et in eos æternam salutem conferat, ita in eodem suo consilio decrevit, quod eos, qui per verbum vocati illud repudiant et Spiritui sancto resistunt et obstinati in ea contumacia perseverant, indurare, reprobare et æternæ damnationi devovere velit. –

\emph{[i margen:]} Det er aabenbart, at medens Tridentiner Concilium bliver staaende paa en Halvei mellem Semipelagianisme og Augustinisme, saa fordyber Concordieformelen sig i et Vildniß, hvoraf det ikke kommer ud. –
\end{quote}
\dcite{Not1:8  (1833; Pap. IIC35; SKS 19, 59)}
% Not1:8  (1833; Pap. IIC35; SKS 19, 59)

\subsection{1835}

\begin{quote}
\emph{Kjøbenhavn,} den 1ste Juni 1835.

De veed, med hvor megen Begeistring jeg i sin Tid hørte Dem tale; hvor enthusiastisk jeg var for Deres Beskrivelse af Deres Ophold i Brasilien, og i den Henseende igjen ikke saa meget med Hensyn til den Masse af enkelte Iagttagelser, hvormed De havde beriget Dem selv og Deres Videnskab, som med Hensyn til det Indtryk, Deres første Udtræden i den vidunderlige Natur giorde paa Dem; Deres paradisiske Lykke og Glæde. Noget Sligt maa altid virke sympathetisk paa ethvert Menneske, som dog ellers har nogen Følelse og Varme, om han endog mener at finde sin Tilfredshed, sin Virksomhed i en ganske anden Sphære, men især paa den Yngre, der endnu kun drømmer om sin Bestemmelse. Vor første Ungdom staaer som Blomsten i Morgendæmring med en deilig Dugdraabe i sit Bæger, hvori alle Omgivelser harmonisk-melancholsk speile sig. Men snart hæver Solen sig paa Horizonten, og Dugdraaben udtørres; med den forsvinder Livets Drømmerier, og nu gjælder det, om Mennesket er istand til, for atter at tage et Billede fra Blomsterne, at udvikle – ved sin egen Kraft ligesom en Nereum, – en Draabe, der kan staae som Frugten af hans Liv. Dertil hører først og fremmest, at man kommer til at staae i den Jordbund, hvor man egenlig hører hjemme; men den er ikke altid saa let at finde. Der gives i den Henseende lykkelige Naturer, som have en saa afgiort Tilbøielighed til en vis Retning, at de troligen arbeide frem ad den engang saaledes anviste Vei, uden at nogensinde den Tanke faaer nogen Magt over dem, at det maaskee egenlig var en ganske anden Vei, de skulde betræde. Der gives Andre, som saa aldeles lade sig styre af deres Omgivelser, at det aldrig gaaer ret op for dem, hvor de egenlig stræbe hen. Ligesom den foregaaende Classe har sit indvortes categoriske Imperativ, saaledes anerkjender denne sidste et udvortes categorisk Imperativ. Men hvor Faa ere ikke de første, og til de sidstes Classe ønsker jeg ikke at høre. Større er deres Tal, som i Livet faae at prøve, hvad egenlig denne Hegelske Dialectik vil sige. Det er nu forresten ganske i sin Orden, at Vinen gjærer, inden den bliver klar; men ubehagelig er desuagtet ofte denne Tilstand i dens enkelte Momenter, ihvorvel den naturligviis betragtet i sin Helhed har sin Behagelighed, forsaavidt den indenfor den almindelige Tvivlen dog har sine relative Resultater. Især har den stor Betydning for den, der igjennem den er kommet paa det Rene med sin Bestemmelse, ikke blot for den paafølgende Ro i Modsætning til den foregaaende Storm: men fordi man da \emph{har Livet} i en ganske anden Betydning end før. Det er dette Faustiske Element, som for en Deel mere eller mindre gjør sig gjældende i enhver intellectuel Udvikling, hvorfor det altid har forekommet mig, at man burde indrømme Ideen til \emph{Faust } en Verdens-Betydning. Saaledes som vore Forfædre havde en Gudinde for Længselen, saaledes, mener jeg, staaer Faust som den personificerede Tvivl. Og mere skal han ikke være, og det er vistnok en Synden mod Ideen, naar Goethe har ladet Faust omvende sig, ligesaavel som naar Merimée har ladet Don Juan omvende sig. Man indvende mig ikke, at Faust dog i det Øieblik, han henvendte sig til Djævelen, gjorde et positivt Skridt, thi heri synes mig netop at ligge Noget af det Dybeste i Sagnet om Faust. Han netop hengav sig til ham for at faae Oplysning og havde den altsaa ikke iforveien, og netop fordi han hengav sig til Djævelen, forøgedes Tvivlen (ligesom en Syg, der falder i en Qvaksalvers Hænder, gjerne bliver endnu værre); thi vel lod Mephistopheles ham see igjennem sine Briller ind i Menneskene og ind i Jordens forborgne Gjemmer, men Faust maatte bestandig nære Tvivl til ham, fordi han aldrig kunde oplyse ham om det allerdybeste i intellectuel Henseende. Til Gud kunde han ifølge sin Idee aldrig komme til at henvende sig, forsaavidt han i det Øieblik, han gjorde det, maatte sige sig selv, at her var i Sandhed Oplysning at finde; men i samme Øieblik havde han jo fornægtet sin Charakteer som Tvivler.

Men en saadan Tvivlen kan ogsaa vise sig i andre Sphærer. Om Mennesket ogsaa er kommet paa det Rene med sig selv om enkelte af saadanne Hovedpuncter, saa gives der i Livet andre ogsaa betydelige Spørgsmaal. Ethvert Menneske ønsker naturligviis at virke efter sine Evner i Verden, men deraf følger da igjen, at han ønsker at uddanne sine Evner i en bestemt Retning, i den nemlig, som er meest passende for hans Individualitet. Men hvilken er denne? Her staaer jeg ved et stort Spørgsmaalstegn. Her staaer jeg som Hercules men ikke paa Skilleveien – nei, her viser sig en langt større Mangfoldighed af Veie, og desto vanskeligere er det altsaa at gribe den rette. Det er maaskee netop en Ulykke ved min Existents, at jeg interesserer mig for altfor Meget og ikke afgjort for Noget; mine Interesser staae ikke alle subordinerede een, men alle staae coordinerede.

Jeg vil forsøge at vise, hvorledes Sagerne staae for mig.

1.      \emph{Naturvidenskaberne.}      Seer jeg for det første hen til hele denne Retning (idet jeg derunder indbefatter alle dem, der søge at forstaaeliggjøre og tyde Naturens Runer, ligefra den, der beregner Stjernernes Fart og saa at sige standser dem deri for nærmere at undersøge dem, til den, der beskriver et enkelt Dyrs Physiologie; fra den, der fra Bjergenes Høider overskuer Jordfladen, til den, der stiger ned i Afgrundens Dybder; fra den, der forfølger Menneskelegemets Dannelse igjennem de utallige Nuancer, til den, der iagttager Indvoldsormene), da har jeg naturligviis paa denne Vei som paa enhver anden (dog fornemlig paa denne) seet Exempler paa Folk, der have skabt sig et Navn i Literaturen ved en uhyre Samlerflid. En stor Mængde af Enkeltheder kjende de, og de have opdaget mange nye; men heller ikke mere. De have blot leveret et Substrat for Andres Tænkning og Bearbeidelse. Og disse Mennesker staae nu der tilfredse med deres Enkeltheder, og dog forekomme de mig at staae ligesom den rige Bonde i Evangeliet; en stor Mængde have de samlet i Laden, men Videnskaben kan sige til dem: »Imorgen vil jeg kræve Dit Liv«, forsaavidt det er den, der afgjør, hvad Betydning hvert enkelt Resultat skal have i det Hele. Forsaavidt nu der var et Slags ubevidst Liv i en saadan Mands Viden, forsaavidt kan Videnskaberne siges at kræve hans Liv; forsaavidt det ikke var, er hans Virksomhed ligesom Menneskets, der ved sit døde Legemes Hensmulren bidrager til Jordens Vedligeholdelse. Anderledes har det naturligviis været med andre Phænomener, med saadanne Naturforskere, som ved deres Speculation have fundet eller stræbt at finde hiint archimediske Punct, som er intetsteds i Verden, og nu derfra have betragtet det Hele og seet Enkelthederne i det rette Lys. Og med Hensyn til saadanne skal jeg ikke nægte, at de have gjort et høist velgjørende Indtryk paa mig. Den Ro, den Harmoni, den Glæde, som man finder hos dem, finder man ellers sjelden. Vi have her i Byen 3 værdige Repræsentanter: en Ørsted, hvis Ansigt bestandig har forekommet mig som en Klangfigur, som Naturen netop har anstrøget paa den rette Maade; en Schouw, der afgiver et Studium for en Maler, der vilde male Adam, hvor han giver alle Dyrene Navne; og endelig en Horneman, der fortrolig med enhver Plante staaer som en Patriarch i Naturen. I den Henseende erindrer jeg ogsaa med Glæde det Indtryk, som \emph{De} gjorde hos mig, der stod som Repræsentant for en stor Natur, der ogsaa skulde have sin Stemme med paa Rigsdagen. Begeistret har jeg været for Naturvidenskaberne og er det endnu; men dog synes mig, at jeg ikke vil gjøre mig dem til mit Hovedstudium. Mig har altid Livet i Kraft af Fornuft og Frihed interesseret meest, at klare og løse Livets Gaade har bestandig været mit Ønske. De 40 Aar i Ørkenen, inden jeg naaede Videnskabernes forjættede Land, forekomme mig for kostbare og det saa meget mere, som jeg troer, at Naturen ogsaa har en Side at betragtes fra, hvortil ikke hører Indsigt i Videnskabens Hemmeligheder. Hvad enten saa jeg i det enkelte Blomster skuer hele Verden, eller jeg lytter til de mange Vink, som Naturen frembyder med Hensyn til Menneskelivet; eller jeg beundrer hine dristige Frihaandstegninger paa Firmamentet; eller jeg for Exempel ved hine Naturlyd paa Ceylon erindres om hine Toner i den aandelige Verden, eller ved Trækfuglenes Bortgang erindres om de dybere Længsler i Menneskets Bryst.

2.      \emph{Theologien.}      Det synes at være det, jeg nærmest har grebet; men ogsaa her møder store Vanskeligheder. Her staaer i Christendommen selv saa store Modsætninger, at det idetmindste meget forhindrer det frie Blik. Orthodoxien er jeg saa at sige voxet op i, som De nok veed; men saasnart jeg selv begyndte at tænke, begyndte efterhaanden den uhyre Colos at vakle. Jeg kalder den med Flid en uhyre Colos, thi den har virkelig i det Hele taget megen Consequents, og de enkelte Dele ere i de mange henrundne Aarhundreder smeltede saa tæt sammen, at det er vanskeligt at komme den tillivs. Jeg kunde nu vel være enig med den i enkelte Puncter, men disse bleve da at betragte som de Spirer, man ofte kan finde i Klipperifter. Paa den anden Side kunde jeg vel ogsaa see det Skjæve, der laae i mange enkelte Puncter, men Hovedbasis maatte jeg for en Tid lade staae in dubio. I det Øieblik \emph{den} forandredes, antog naturligviis det Hele en ganske anden Skikkelse, og saaledes henledes min Opmærksomhed paa et andet Phænomen: Rationalismen, som i det Hele taget gjør en temmelig maadelig Figur. Forsaavidt nemlig Fornuften consequent forfølger sig selv og – ved at gjøre sig Rede for Forholdet mellem Gud og Verden – nu igjen kommer til at see Mennesket i dets dybeste og inderligste Forhold til Gud og i den Henseende ogsaa fra sit Standpunct af betragter Christendommen som den, der i saa mange Aarhundreder har tilfredsstillet Menneskets dybeste Trang, forsaavidt er der Intet at indvende imod den, men det bliver da heller ikke Rationalisme, thi Rationalisme faaer sin egenlige Farve af Christendommen og staaer altsaa paa et ganske andet Gebet og danner ikke et System men en Noahs Ark (for at benytte et af Professor Heiberg ved en anden Leilighed anvendt Udtryk), hvori de rene og urene Dyr ligge ved Siden af hinanden. Den gjør omtrent samme Indtryk som vor Borgervæbning i ældre Tid vilde gjøre ved Siden af Potsdammer-Garden. Den søger derfor væsentlig at knytte sig til Christendommen, baserer sine Udviklinger paa Skriften og skikker en Legion af Skriftsteder forud for hvert enkelt Punct; men selve Udviklingen er ikke gjennemtrængt deraf. De bære dem ad som Cambyses, der paa hans Tog mod Ægypten skikkede de hellige Høns og Katte iforveien, men de ere derfor ogsaa parate til ligesom den romerske Consul at kaste de hellige Høns overbord, naar de ikke vil æde. Feilen ligger altsaa deri, at de, naar de finde Eenhed med Skriften, basere sig derpaa, men ellers ikke, og saaledes staae paa to fremmedartede Standpuncter.

Nonnulla desunt.

– Hvad smaa Ubehageligheder angaaer, vil jeg kun bemærke, at jeg er ifærd med at læse til theologisk Attestats, en Beskjæftigelse, som slet ikke interesserer og som derfor heller ikke gaaer synderlig rask fra Haanden. Jeg har altid holdt mere af det frie, maaskee derfor ogsaa lidt ubestemte Studium end af den Beværtning ved sluttet Bord, hvor man veed forud, hvilke Gjæster man træffer, og hvilke Retter Mad man faaer hver Dag i Ugen. Da det imidlertid engang er en Nødvendighed, og man knap faaer Lov at komme ind paa de videnskabelige Fælleder uden at være indbrændt, og jeg saavel anseer det med Hensyn til min nærværende Sindstilstand for noget Gavnligt for mig, som ogsaa veed, at jeg dermed kan gjøre Fader en stor Glæde (han mener nemlig, at det egenlige Canaan ligger paa den anden Side af theologisk Attestats, men bestiger derhos ligesom Moses fordum Thabor og beretter, at jeg aldrig kommer derind – dog vil jeg haabe, at det ikke ogsaa denne Gang gaaer i Opfyldelse), saa faaer jeg vel at gribe mig an. Hvor lykkelig er ikke De, som i Brasilien har fundet en uhyre Mark for Deres Iagttagelser, hvor der ved hvert Skridt frembyder sig nye Mærkeligheder, hvor den øvrige lærde Republiques Skrig ikke forstyrrer Deres Ro. Thi mig forekommer den lærde theologiske Verden som Strandveien en Søndag Eftermiddag i Dyrehaugstiden1 – de storme hverandre forbi, huje og skrige, lee og gjøre Nar af hverandre, kjøre Hestene ihjel, vælte og blive overkjørte, og naar de saa endelig komme overstøvede og forpustede til Bakken – ja saa see de paa hverandre – og vende hjem.

Hvad Deres Tilbagekomst angaaer, saa vilde det være barnagtigt af mig at paaskynde den, ligesaa barnagtigt som naar Achilles Moder søgte at skjule ham for at undgaae den hurtige, hæderlige Død. – – Lev vel!

Hvor forvirrende er ikke ofte Betragtningen af Livet, naar det viser sig for os i dets hele Rigdom, naar vi studse over Forskjelligheden i Evner og Anlæg, lige fra den, der saa inderlig har levet sig ind i Guddommen, at han ligesom Johannes fordum kan siges at hvile ved Guddommens Bryst, til den, der i dyrisk Raahed misforstaaer og vil misforstaae enhver dybere Bevægelse i Menneskelivet; fra den, der med Lynceus Blik skuer ind i Historiens Gang og næsten vover at stille dens Viser, til den, for hvem endog den simpleste Ting falder vanskelig; – eller vi lægge Mærke til Uligheden i Stand og Stilling og snart med misundeligt Øie føle Savnet af, hvad der er skjænket Andre, snart med taknemlig Veemod see, hvor Meget der er givet os, som er nægtet Andre; – – og nu en kold Philosophi vil forklare os det Hele af en Præexistents og ikke see det som et uendeligt Livsmalerie med sit brogede Farvespil og sine utallige Nuanceringer.

\emph{Gilleleie d. 1. Aug. 1835.}

Saaledes som jeg har søgt at vise det paa de foregaaende Sider, stode virkelig Sagerne for mig. Naar jeg nu derimod søger at bringe mig selv paa det Rene om mit Liv, forekommer det mig anderledes. Ligesom det varer længe inden Barnet lærer at skjelne sig selv fra Gjenstandene, og det saaledes i lang Tid saa lidet udsondrer sig fra sine Omgivelser, at det med Fremhævelse af den lidende Side siger for Exempel: »\emph{mig} slaaer Hesten«, saaledes gjentager det samme Phænomen sig i en høiere aandelig Sphære. Jeg troede derfor, at jeg muligen vilde komme mere til Ro ved at gribe et andet Fag, ved at rette mine Kræfter paa et andet Maal. For en Tid vilde det vel og have lykkedes mig at fordrive en vis Uro, men forstærket vilde den vistnok være vendt tilbage som Feberanfaldet efter Nydelsen af koldt Vand. Det, der egenlig mangler mig, er at komme paa det Rene med mig selv om, \emph{hvad jeg skal gjøre} 2, ikke om hvad jeg skal erkjende, uden forsaavidt en Erkjenden maa gaae forud for enhver Handlen. Det kommer an paa at forstaae min Bestemmelse, at see, hvad Guddommen egenlig vil, at \emph{jeg} skal gjøre; det gjælder om at finde en Sandhed, som er Sandhed \emph{for mig} 3, at finde \emph{den Idee, for hvilken jeg vil leve og døe.} Og hvad nyttede det mig dertil, om jeg udfandt en saakaldet objectiv Sandhed; om jeg arbeidede mig igjennem Philosophernes Systemer og kunde, naar forlangtes, holde Revue over dem; at jeg kunde paavise Inconsequentser indenfor hver enkelt Kreds; – hvad nyttede det mig dertil, at jeg kunde udvikle en Statstheorie og af de mangestedsfra hentede Enkeltheder combinere en Totalitet, construere en Verden, hvori jeg saa igjen ikke levede, men som jeg blot holdt frem til Skue for Andre; – hvad nyttede det mig at kunne udvikle Christendommens Betydning, at kunne forklare mange enkelte Phænomener, naar den for \emph{mig selv} og \emph{mit} Liv \emph{ikke} havde nogen dybere Betydning? Og jo mere jeg kunde det, jo mere jeg saae Andre \tlg{tilegne} sig mine Tankefostre, desto sørgeligere blev min Stilling, ei ulig de Forældres, der af Armod maa skikke deres Børn ud i Verden og overlade dem til Andres Pleie. Hvad nyttede det mig, at Sandheden stod der for mig kold og nøgen, ligegyldig ved, om jeg anerkjendte den eller ikke, snarere bevirkende en ængstelig Gysen end en tillidsfuld Hengivelse? Vel vil jeg ikke nægte, at jeg endnu antager et \emph{Erkjendelsens Imperativ} \emph{;} og at der igjennem det ogsaa lader sig virke paa Menneskene, \emph{men da maa den levende optages i mig,} og \emph{det er det,} jeg nu anerkjender for Hovedsagen. Det er derefter min Sjæl tørster, som Africas Ørkener efter Vand. Det er det, der mangler, og derfor staaer jeg som en Mand, der samlede Indbo og havde leiet Værelser, men havde endnu ikke fundet den Elskede, der skulde dele Livets Med- og Modgang med ham. Men for saaledes at finde hiin Idee eller rettere sagt finde mig selv, nytter det mig ikke at styrte mig endnu mere ind i Verden. Og det var netop det, jeg før gjorde. Derfor troede jeg, at det vilde være godt at kaste mig ind i \emph{Jurisprudentsen,} for at jeg kunde udvikle min Skarpsindighed i de mange Livets Forviklinger. Her tilbød sig netop en stor Masse af Enkeltheder, hvori jeg kunde fortabe mig; her kunde jeg maaskee fra de givne Facta udconstruere en Totalitet, en Organisme af Tyve-Liv, forfølge det i hele dets dunkle Side (ogsaa her er en vis Associations-Aand i høi Grad mærkelig). Derfor kunde jeg ønske at blive \emph{Acteur,} for at jeg ved at sætte mig ind i en Andens Rolle kunde faae, saa at sige, et Surrogat for mit eget Liv, og ved den udvortes Afvexling finde en vis Adspredelse. Det var det, der manglede mig, at føre \emph{et fuldkommen menneskeligt Liv} og ikke blot et Erkjendelsens4, saa jeg derved kommer til ikke at basere mine Tanke-Udviklinger paa – ja paa Noget, man kalder Objectivt, – Noget, som dog i ethvert Tilfælde ikke er mit eget, men paa Noget, som hænger sammen med min Existents's dybeste Rod5, hvorigjennem jeg saa at sige er indvoxet i det Guddommelige, hænger fast dermed, om saa hele Verden styrter sammen. See det er \emph{det, jeg mangler,} og \emph{derhen stræber jeg.} Med Glæde og inderlig Styrkelse betragter jeg derfor de store Mænd, som saaledes have fundet hiin Ædelsteen, for hvilken de sælge Alt, endog deres Liv6, hvad enten jeg seer dem med Kraft gribe ind i Livet, med sikkre Skridt, uden at vakle, at gaae frem paa deres bestemte Bane; eller jeg opdager dem afsides fra den alfare Vei, fordybede i sig selv og i Arbeiden for deres ophøiede Maal. Med Ærbødighed betragter jeg endog de Afveie, der ogsaa her ligge saa nær. Det er denne Menneskets indvortes Handlen, denne Menneskets Guds-Side det kommer an paa, ikke en Masse af Erkjendelser; thi de ville vel følge og ville da ikke vise sig som tilfældige Agregater eller som en Række af Enkeltheder, den ene ved Siden af den anden, uden et System, uden et Brændpunct, hvori alle Radier samles. Et saadant Brændpunct har jeg vel ogsaa søgt. Saavel paa Forlystelsernes bundløse Hav har jeg forgjeves søgt en Ankerplads som i Erkjendelsens Dybder. Jeg har følt den næsten uimodstaaelige Magt, hvormed den ene Forlystelse rækker den anden Haanden; jeg har følt den Art af uægte Begeistring, som den er istand til at fremkalde; jeg har ogsaa følt den Kjedsomhed, den Sønderrevethed, der følger derpaa. Jeg har smagt Frugterne af Kundskabens Træ og tidt og ofte glædet mig ved deres Velsmag. Men denne Glæde var kun i Erkjendelsens Øieblik og efterlod ikke noget dybere Mærke i mig selv. Det forekommer mig, at jeg ikke har drukket af Viisdommens Bæger men er faldet ned deri. Jeg har søgt at finde dette Princip for mit Liv ved Resignation, ved at forestille mig selv, at da Alt gik efter urandsagelige Love, det ikke kunde være anderledes, ved at sløve min Ærgjerrighed og Forfængeligheds Følehorn. Fordi jeg ikke kunde faae Alt til at passe efter mit Hoved, trak jeg mig ud deraf med Bevidsthed om min egen Dygtighed, omtrent som naar en aflægs Præst resignerer med Pension. Hvad fandt jeg? Ikke mit Jeg; thi det søgte jeg netop at finde paa hine Veie (jeg tænkte mig, om jeg saa maa sige, min Sjæl som indesluttet i en Kasse med Springlaas for, og som nu de udvortes Omgivelser ved et Tryk paa Fjedrene skulde faae til at springe op). – Det var altsaa det første, der skulde afgjøres, det var hiin Søgen og Finden af Himmerigets Rige. Ligesaalidet som et Verdenslegeme, naar det tænktes at skulle danne sig, først vilde fatte Bestemmelse om, hvorledes dets Overflade skulde være, til hvilke Legemer det skulde vende den lyse, til hvilke den mørke Side, men først vilde lade Centrifugal- og Centripetal-Kræfternes Harmoni realisere dets Existents og lade Resten komme af sig selv – ligesaalidet nytter det Mennesket først at ville bestemme det Udvortes og siden det Grundconstituerende. Man maa først lære at kjende sig selv, inden man kjender noget Andet (γνωϑι σεαυτον). Først naar Mennesket saaledes inderlig har forstaaet \emph{sig selv} og nu seer sin Gang hen af sin Bane, først da faaer hans Liv Ro og Betydning, først da bliver han fri for hiin besværlige, uheldsvangre Reisekammerat – hiin Livs-Ironi7, der viser sig i Erkjendelsens Sphære og byder den sande Erkjenden at begynde med en Ikke-Erkjenden (Socrates) 8, ligesom Gud skabte Verden af Intet. Men fornemlig hører den hjemme paa Sædelighedens Farvande for dem, som endnu ikke ere komne ind under Dydens Passat. Her tumler den Mennesket paa det rædsomste omkring; snart lader den Mennesket føle sig lykkelig og tilfreds ved Forsættet om at vandre frem ad den rigtige Vei, snart styrter den ham ned i Fortvivlelsens Afgrunde. Ofte dysser den Mennesket i Søvn ved den Tanke: »det kan nu eengang ikke være anderledes«, for pludselig at vække ham til et strængt Forhør. Ofte lader den ligesom Glemselens Slør falde over det Forbigangne, for da atter igjen at lade hver enkelt Ubetydelighed træde frem i et levende Lys. Naar Mennesket her kjæmper sig frem paa den rette Vei, glæder sig over at have beseiret Fristelsernes Magt, kommer maaskee næsten i samme Moment ovenpaa den fuldkomneste Seir en tilsyneladende ringe udvortes Omstændighed, der støder ham lig Sisyphus ned fra Klippens Spidse. Ofte naar Mennesket har concentreret sin Kraft paa Noget, møder en lille udvortes Omstændighed, der tilintetgjør det Hele. (Jeg vilde sige: som en Mand, der kjed af Livet vilde styrte sig i Themsen, og som nu et Myggestik standser netop i det afgjørende Øieblik.) Ofte lader den Mennesket ligesom den Brystsvage befinde sig allerbedst9, naar det er allerværst. Forgjeves søger han at arbeide imod; han har ikke Kraft nok dertil, intet hjælper det ham, at han saa ofte gjennemgaaer det Samme; den Art af Øvelse, som derved opnaaes, kommer det ei an paa. Ligesaa lidet som den, der har nok saa megen Øvelse i at svømme, vil kunne holde sig oppe i et Uveir, men kun den, der inderlig er overbeviist om og har erfaret, at han virkelig er lettere end Vandet, ligesaalidet kan det Menneske, der mangler det indre Holdningspunct, holde sig oppe i Livets Storme. – Kun naar Mennesket saaledes har forstaaet sig selv, er han istand til at hævde sig en selvstændig Existents og saaledes undgaae at opgive sit eget Jeg. Hvor ofte see vi ikke – (i en Tid, hvor vi ved vore Lovtaler ophøie hiin græske Historieskriver, fordi han vidste at adoptere en fremmed Stiil til den meest skuffende Lighed med dens Ophavsmand, istedetfor at han burde dadles, da den første Priis for en Skribent altid er at have en egen Stiil det vil sige en ved sin Individualitet modificeret Udtryks- og Fremstillings-Maade) – hvor ofte see vi ikke Folk, der enten af aandelig Ladhed leve af de Smuler, der falde fra Andres Bord, eller af mere egoistiske Grunde søge at leve sig ind i Andre, for da tilsidst ligesom en Løgner ved den hyppige Gjentagelse af sine Historier selv at troe dem. Uagtet jeg endnu langt fra ikke saaledes inderlig er kommen til at forstaae mig selv, har jeg med dyb Agtelse for dens Betydning søgt at frede om min Individualitet – dyrket den ubekjendte Guddom. Med en utidig Ængstelighed har jeg søgt at undgaae at komme i for nær Berøring med de Phænomener, hvis Tiltrækningskraft maaskee vilde udøve for stor Magt over mig. Jeg har søgt at \tlg{tilegne} mig meget af dem, studeret deres Individualitet og Betydning i Menneskelivet, men dog tillige vogtet mig for ikke som Myggen at komme Lyset for nær. Ved Omgang med de sædvanlige Mennesker har jeg kun havt Lidet at vinde eller tabe. Deels har deres hele Syslen – det saakaldte practiske Liv10 – ikkun lidet interesseret mig; deels har den Kulde og den Mangel paa al Sympathie, hvormed de betragte de aandelige og dybere Bevægelser i Mennesket, fjernet mig endmere fra dem. Mine Omgangsvenner have med faa Undtagelser ikke udøvet nogen synderlig Magt over mig. Et Liv, der ikke er kommen paa det Rene med sig selv, maa nødvendigviis vise en ujævn Side-Overflade frem; her ere de nu kun blevne staaende ved enkelte Facta og deres tilsyneladende Disharmoni, thi til at søge at opløse den i en høiere Harmoni eller til at indsee Nødvendigheden deraf, dertil havde de ikke Interesse nok for mig. Deres Dom over mig var derfor altid eensidig, og jeg har afvexlende enten lagt for megen eller for liden Vægt paa deres Udsagn. Ogsaa deres Indflydelse og de deraf mulige Misvisninger paa mit Livs Compas har jeg nu unddraget mig. Og saaledes staaer jeg da atter paa det Punct, at jeg maa begynde paa en anden Maade. Jeg vil nu forsøge rolig at fæste Blikket paa mig selv og begynde inderlig at handle; thi kun derved bliver jeg istand til, ligesom Barnet ved sin første med Bevidsthed foretagne Handling kalder sig: »jeg«, saaledes istand til i dybere Betydning at kalde mig: »jeg«.

Dog dertil udfordres Udholdenhed, og man kan ikke strax høste, hvor man har saaet. Jeg vil erindre hiin Philosophs Methode: at lade sine Disciple tie i tre Aar, saa kommer det vel. Ligesom man ikke begynder en Fest med Solens Opgang men med dens Nedgang, saaledes maa man ogsaa i den aandelige Verden først arbeide i nogen Tid frem, førend Solen ret kan skinne for os og gaae op i al sin Herlighed; thi vel hedder det: at Gud lader sin Sol opgaae over Gode og Onde og lader regne over Retfærdige og Uretfærdige, men ikke saa i den aandelige Verden. Saa være da Loddet kastet – jeg gaaer over Rubicon! Vel fører denne Vei mig \emph{til Kamp;} men jeg vil ikke forsage. Jeg vil ikke sørge over den forbigangne Tid – thi hvortil Sorg? Med Kraft vil jeg arbeide mig frem og ikke spilde Tiden med at sørge ligesom den, der i en Hængesæk først vilde beregne, hvor dybt han var sunken, uden at erindre, at han i den Tid, han anvender hertil, synker endnu dybere. Jeg vil ile frem paa den fundne Vei og tilraabe Enhver, som jeg møder: ikke som Loths Hustru at see sig tilbage, men huske paa, at det er en Bakke, vi stræbe op af.

Det gaaer ikke an, saaledes at haste frem som Moralisterne, endogsaa forkastende Anger, mene; – see vi ikke i den physiske Verden Taagen som en stille Bøn at hæve sig op fra Jorden for bønhørt at vende tilbage som forfriskende Dug –?

Der er noget besynderligt ironisk i Kjøbenhavnernes Dyrehaugstoure; de søge at ryste Byens spidsborgerlige Støv af dem, flye dem selv – og gjenfinde dem selv paa Bakken.

d. 14. Jan. 37.

Hvor ofte hænder det ikke, at man, naar man allerbedst troer at have grebet sig selv, befindes at have favnet Skyen istedetfor Juno.

Da først faaer Mennesket en indre Erfaring; men hos hvor Mange er ikke Livets forskjellige Indtryk lig de Figurer, Havet tegner i Sandet for strax at udslette dem sporløst igjen.

Deraf lader sig ogsaa forklare et ikke ualmindeligt Phænomen, en vis Gjerrighed1 med Ideer. Netop fordi Livet ikke er sundt men Erkjendelsen for prædominerende, opfattes Ideerne ikke som de naturlige Blomster paa Livets Træ og fastholdes ikke som saadanne og som de, der alene have deres Betydning som saadanne – men som enkelte Lysglimt, som om Livet blev rigt ved en Mængde saadanne saa at sige udvortes Ideer (sit venia verbo; – aphoristisk). De forglemme, at det gaaer med Ideerne som med Thors Hammer – den kommer igjen hvor man kaster den hen, om end i en modificeret Skikkelse.

Et lignende Phænomen er den feilagtige Anskuelse man har af Erkjendelsen og dens Resultater, idet man taler om de objective Resultater og erindrer ikke, at netop den egenlige Philosoph er i høieste Grad subjectiv. Jeg behøver blot at nævne Fichte. Saaledes behandler man ogsaa Vittigheden, seer ikke den som den af hele Forfatterens Individualitet og Omgivelser nødvendig fremspringende Minerva, derfor i en vis Henseende som noget lyrisk – [deraf ogsaa den Rødme, der gjerne ledsager en vis Art af Vittigheder, som netop tyder paa, at den kom naturlig frem, nyfødt; d. 20. Sept. 36.], – men som Blomster, man kan afplukke og gjemme til eget Brug. (Forglemmigeien har sin Plads paa Marken skjult og beskedent men bliver styg i en Hauge.)

Hvor nær ligger Mennesket ellers trods al sin Viden ved Galskab? Hvad er Sandhed andet end en Leven for en Idee? Alt maa til syvende og sidst basere sig paa et Postulat; men i det Øieblik det saaledes ikke længere staaer udenfor ham, men han lever deri, først da er det ophørt at være ham et Postulat. (Dialectik – Disput.)

Saaledes vil det blive os let, naar vi først eengang af Ariadne (Kjærlighed) have modtaget hiint Nøgle Garn, da at gaae igjennem alle Labyrinthens (Livets) Irrgange og dræbe Uhyret. Men hvor Mange styrte sig ikke ind i Livet (Labyrinthen) uden at have iagttaget hiin Forsigtighed (de \emph{unge} Piger og Drengebørn, der hvert Aar offredes til Minotaurus) –?

Vel vedbliver den da ogsaa i en vis Betydning, men han er istand til at taale disse Livets Kastevinde; thi jo mere Mennesket lever for en Idee, desto lettere kommer han ogsaa til at sidde paa Forundringsstolen for hele Verden. – Ofte kan ogsaa en underlig Ængstelighed paanøde sig for, at man, naar man allermeest troer at have forstaaet sig selv, dog egenlig kun har lært en Andens Liv udenad1.

En eiendommelig Art af Ironi findes ogsaa fremstillet i en arabisk Fortælling »Morad der Buckelige« (der findes i Moden Zeitung, »Bilder Magazin« Nr. 40. 1835): en Mand, der kommer i Besiddelse af en Ring, som skaffer ham Alt hvad han ønsker men dog bestandig med et »men« derved, for Exempel naar han ønsker sig i Sikkerhed, befinder han sig i Fængsel etcetera (Denne Fortælling findes i Riises Bibliothek for Ungdommen, 2den Aargangs 6te Hefte, 1836 pagina 453.) – Jeg har ogsaa hørt eller læst et Sted om en Mand, der udenfor et Theater hørte nogle Toner af en saa skjøn og henrivende Sopran, at han strax forliebede sig i Stemmen; han iler ind i Theatret og møder en tyk og fed Mand, der paa hans Spørgsmaal, hvo det var, der sang saa smukt, svarede: »det var mig« – han var en Castrat.

Ordsproget siger ogsaa: »Af Børn og Afsindige skal man høre Sandheden«. Og her er vistnok ikke Tale om at have Sandheden ifølge Præmisser og Conclusioner, men hvor ofte har ikke et Barns eller en Afsindigs Tale nedtordnet den Mand, som Skarpsindigheden Intet kunde udrette med –?

»Es ist, wie mit den anmuthigen Morgentraümen, aus deren einschläferndem Wirbel man nur mit Gewalt sich herausziehen kann, wenn man nicht in immer drückender Müdigkeit gerathen, und so in krankhafter Erschöpfung nachher den ganzen Tag hinschleppen will.« Novalis Schriften, Berlin 1826. 1ster Theil pagina 107.

Dette Liv, som er temmelig gjennemgribende i hele Tidsalderen, viser sig ogsaa i det Store. Medens de ældre Tider opførte Værker, som Betragteren maatte forstumme ved, bygge de nu en Bro under Themsen (Nytte og Fordeel). Ja, Barnet spørger næsten førend det faaer Tid til at beundre en Plantes Skjønhed eller en eller anden Dyreart: hvortil er den gavnlig?
\end{quote}
\dcite{AA:12  (1835; Pap. IA72; SKS 17, 19)}
% AA:12  (1835; Pap. IA72; SKS 17, 19)

\subsection{1836}

\begin{quote}
Forelæsninger af Molbech. 2d Deel. Kbh. 1832.

pagina 15. »At indskrænke den elegiske Digtnings Grændser alene til Kjærlighedsforholdet, mene vi derfor, at være ligesaa urigtigt som at antage, at det alene var den ulykkelige Kjærlighed, hvis Klagelyd skulde høres i Elegien....... I Følelsens hemmelige Dyb ligge Glædens og Sorgens Strænge hinanden saa nær, at de sidste kun altfor let gjenklinge, idet de første røres. Digteren kan midt i sin lykkelige Elskovs Himmel fornemme Ahnelsen om den jordiske Glædes forfængelige Natur, eller han kan midt i Besiddelsens Lyksalighed føle Virkningen af en lønlig Frygt for at tabe den. –

pagina 91. Naar vi gjenemløbe det hele Forraad af lyriske Digte, som nu i hans Værker (Baggesens) ligge samlede for os, da ville vi finde: at den allerstørste Deel, saavel ved deres Gjenstand som i Udførelsen dreie sig om ganske subjective Forhold, hvori atter deels mere eller mindre flygtige erotiske Stemninger spille Hovedrollen; deels det elegante og vittige Galanterie – eller hiint den moderne Tids i lette Former uddannede Forhold til det andet Kjøn, som er en paa Livets Overflade spillende Levning af Middelalderens Ridderaand; et Slags halv ironisk Spøgen med de Situationer, hvori Tilbøiligheden kan prøve, hvorvidt den i dette Tilfælde tør gaae, uden at blive alvorlig, uden at overskride Grændsen, hvor den yttrer sig som Lidenskab.« –

Til Oplysning af Begrebet af det »Romantiske« findes i den 21de Forelæsning et og Andet dog ikke Nyt. pagina 181 »da nu dette (den herskende Drift til at drage det Evige og Uendelige, det hvis Maal og Skranker ingen jordisk Flugt kan naae, ned i Phænomenernes Verden.) selv for den mægtigste Phantasie, den dybeste Fornuft og den meest brændende Kjærligheds Begeistring, dog evig bliver uopnaaeligt: saa maa det, vi kalde det Romantiske, det evigt Higende, det Længselsfulde, den uendelige den savnende Salighed i Følelsen, det Ahnende og Overjordiske i Phantasien (Grundvolden ogsaa for det Eventyrliges hele Verden) det Mystiske og Dybe i Tanken, der ligesom paa eengang søger at identificere sig med Følelsen og Phantasien – dette, siger jeg, maa tjene til at udfylde den brede Kløft mellem Idee og Væsen, imellem det Evige, det Guddommelige, det Overnaturlige, som menneskelig Stræben søger at optage i sig, eller at binde i Konstens forskjellige Skikkelser eller Form-Riger, og disse Skikkelser selv. Det er saaledes hverken det sentimentale, eller det ridderlige, det eventyrlige Element, som udgjør det væsentlige eller nødvendige Grundstof for det Romantiske – det er snarere det Uendelige, den af sandselige Skranker ubundne Frihed i Phantasiens Virken, i Anskuelsen af det Ideelle, i Følelsens Fylde og Dybde, i Tænkningens mod Ideer rettede Kraft, hvori vi maae søge hiin Grundbetingelse for det Romantiske og da tillige for en stor og betydende Deel af den nyere Tids Konst. Det Romantiske, siger Jean Paul, er det Skjønne uden Begrændsning, eller det skjønne Uendelige ligesom der gives et sublimt Uendeligt.« Jean P. ligner det Romantiske med en Egns Belysning af Maaneskinnet, eller med Tonebølgerne i Efterlyden af en Klokkes Ringen af en anslaaet Stræng – en bævende Lyd, der ligesom svømmer bort i fjernere og fjernere Afstand, og endelig taber sig i os selv, og endnu lyder i vort Indre, skjøndt det udenfor os er stille. – fremdeles »er al Digtning et Slags Sandsigerkonst, saa er den romantiske Digtning en Ahnelse af et større Tilkommende.« Man har derfor kaldt det Romantiske Ahnelsens Poesie. – – Men betragte vi nu særskilt de forskjellige Konster i deres Heelhed med Hensyn til deres Slægtskab til det Romantiske eller til den Egenskab at modtage Charakteren af det aandeligt Ubegrændsede; saa ville vi finde, at de i Graden af dette Slægtskab fremtræde i en Orden, hvori Tonekonsten, der overgaaer enhver anden Konst i at gjøre det Uendelige, Uudtømmelige, Uudgrundelige i Sjælen, men her kun gjenem Følelsen, umiddelbar anskueligt, indtager den første Plads; dernæst Digterkonst; Malerkonst; Plastik...... Musiken som den meest romantiske af alle skjønne Konster er tillige den nyeste af alle; ja vi kunne sige, at den egentlig musicalske Composition, Opdagelsen af Harmoniens Theorie og practiske Anvendelse i Musiken, aldeles tilhører vor Tid (det nyere Malerie udmærker sig ved Kjendskab til det Lysdunkle, og det fuldstændige Colorit) – derimod tilhører netop Plastiken den gamle Tid. – pagina 189. Her kan det være fornødent først at møde den Erindring, at det Romantiske i dets Almindelighed ei er nogen udelukkende Eiendom for det christelige Europa; der gives i den orientalske, (indiske), nordiske, celtiske, Mythologie meget romantisk Stof. – pagina 198. En anden Bemærkning, som vi her tilføie hiin, (at en med det Antike nær beslægtet Digteraand endog kan tilhøre vor moderne Tid) er denne: at det Romantiske, hvilket vi modsætte det Antike, ingenlunde maa forvexles med en udelukkende aandelig eller mystisk religieus Stræben efter Forening med det Høiere, Oversandselige. Det rom: har nemlig foruden dets ideelle Grundvæsen eller Grundkraft, der meddeler det en overveiende aandelig, phantasierig, følelsesfuld Charakter, en mod Sandseverdenen henvendt Side, hvorved det bliver i Stand til at udvikle og fremstille dets eiendommelige Konstverden. Denne Anvendelse af det romantiske Stof skeer poetisk paa meget forskjellige Maader; men overhovedet kunne de siges at have den Egenskab tilfælles, at derved søges og fremstilles en Skjønhed i det Mangfoldige en Forbindelse af dettes meest brogede, underfulde, eventyrlige og phantastiske Skikkelser og Billeder til et for den modtagende Indbildningskraft og beskuende Konstsands fatteligt Heelt. Denne Mangfoldighed, som et herskende Træk i det romantiske Konstskjønne, synes ligesaa eiendommelig for dette, som det Stores, det Sublimes Enkelthed, Eenhed og Simplicitet er charakteristisk for den antike Konst. Men under meget forskjellige Yttringer finde vi den romantiske Mangfoldighed i den nyere Tids Digterværker. Saaledes for Exempel fremtræder den i Middelalderen for det meste kun som rige Materialier, eller ufuldendte Elementer til poetiske Konstværker, hvori mangen Gang den herligste, den ædelste og meest livfulde Aand; den reneste religieuse Tro, den meest glødende Følelse rører sig; uden at den af Mangel paa Dannelse og Konstfærdighed har formaaet at forme og behandle det overordentlig rige poetiske Stof, som den har været istand til at \tlg{tilegne} sig. Middelalderens Poesie, der overhovedet har en fremherskende episk Charakteer, er saaledes uhyre indholdsrig, men tidt ikke mindre formløs. – En anden, sildigere Yttring af det Romantiske i Poesien, var mindre ægte og mindre ophøiet end hiin, hvor Aanden dvæler i det Vidunderliges Mysterier og i en levende Tro paa det Hellige og Oversandselige; men holdt sig nærmere til Sandseverdenen og til Hverdagslivet og kaldte det Underfulde fra Troens og Ahnelsens religieuse Sphære ned paa en lavere, mere broget og phantastisk Skueplads, hvor det udsvævende Eventyrlige, Tryllerie, Magie og Hexekonster drive deres Spil; hvor Feer, Gnomer, Alfer og en Mængde andre overnaturlige Væsener og dunkle Magter gribe ind i Menneskets Skjæbne og i Begivenhedernes Gang.

24. Martz. 1836.

\emph{[i margen:]} pagina 84. Kommer Forfatteren en passant til at tale om Hymnen i den gamle græske og romerske Kirke, og citerer i den Anledning Herder Briefe zur Beförd: der Humanität 7te Samml: S: 21 følgende Preisschrift uber die Wirk: der Dichtkunst. Samtliche Wercke zur sch: Lit: IX Seite 419 følgende« For mig kunde det Hele maaskee have Interesse, maa eftersees. –
\end{quote}
\dcite{BB:1  (1836-03-24; Pap. IC88; SKS 17, 59)}
% BB:1  (1836-03-24; Pap. IC88; SKS 17, 59)

\begin{quote}
\emph{Ueber Goethe's Faust. Vorlesungen von Doctor: K. E. Schubarth. Berlin 1830.}

Vi forbigaae indtil videre Dedicationen, Forspillet paa Theatret og Prologen i Himlen.

\emph{1ste Scene.} S: leverer et temmeligt prosaisk resumè af Göethes Ord; og knytter derpaa en Deel weitschweifige Bemærkninger dertil, berigede med Citater af andre Skrifter af Goethe. Derpaa viser han, hvorledes der i Fausts Bevidsthed var indtraadt en \emph{Strid mellem Liv og Viden,} der betingede, at han greb til Magie, der netop har meget tillokkende med sig paa Grund af den muntre Side, den giver Livet, idet den ved Erkjendelsens Udvidelse vil sikkre sig Rigdom, al Slags forhøiet Livsnydelse, langt Liv og saa videre. Magien er derfor ikke i og for sig forkastelig, da den stræber efter Erkjendelse; men forkastelig, fordi den søger det Sande i et underordnet Øiemeed. Havde nu Fausts Videnskabelighed faaet et haardt Stød ved denne Strid, saa kom nu ogsaa den med hans tantaliske Erkjendelsens Stræben forbundne \emph{Utaalmodighed,} han vil favne Alt. Han har følt, at hans Samtids Viden ikke vil tilfredsstille ham, han har ahnet en langt livligere, aandeligere Behandling af Videnskaben; men nu styrter han løs, drømmer sig noget om et ubetinget Herredømme over Naturen som Guds Afbillede. Endelig kommer dertil \emph{hans religieuse Strid med sig selv, Forholdet mellem Auctoritet og egen Skuen»;} Indem er erklärt, das Wunder sei des Glaubens liebstes Kind, theilt er dem Glauben jene Art von Beschränkheit zu, welche von phantastischer Leichtglaubigkeit, oder von dumpfem, kindischen, schwachmüthigen Erstaunen und thörichtem Gefangennehmen des eignen Verstandes sich schwer abgrenzen läßt, anstatt daß der wahre Glaube nicht vor dem Wunder erstaunt, was dem Verstande nur ein Widerspruch ist, sondern das wahrhafte Wunder, in dem findet, was bei hellster Begreiflichkeit und Anschaulichkeit dog unergründlich bleibt und so als dieß wahre Wunder tausendfältig sich überall wiederholt.« Nu følger Citater af Goethes Værker a, som ere meget interessante; men som paa en Maade bortlede Opmærksomheden fra Digtet af. –

Den følgende Scene behandler S: paa samme Maade, en resumè, og en Deel Prædikebemærkninger. Hvad Factumet med Hunden angaaer, da anfører han et Citat af Goethes naturvidenskabelige Værker, hvor han omtaler dette Sted, og gjør opmærksom paa, at, Lys og Dunkel kræve strax deres Modsætning; Han (Goethe) fortæller, at Stedet i Faust var skrevet længe før han gjorde den Opdagelse: »aus dichterischer Ahnung und in halben Bewußtsein, als bei gemäßigtem Licht, vor meinem Fenster auf der Straße ein schwarzer Pudel vorbeilief, der einen hellen Lichtschein nach sich zog: das undeutliche, im Auge gebliebene Bild seiner vorübereilenden Gestalt.«b Rigtigt bemærker nu S:, at der i Faust Gemyt er indtraadt en hidtil ukjendt Ro (da han er kommen hjem fra Spadseretouren), som dog imidlertid kun er et Havblik, og nu gjør Faust det sidste endelige Forsøg ved at læse og fortolke Bibelen for sig, men her strander han ganske, og den fromme Stemning forsvinder og han træder op som Aandebesværger, og nu, da han saaledes ogsaa har opgivet Christenheden, træder Mephistopheles op. Derpaa følger nu hos S: en lang Sludder om det Onde og dets Betydning i Existentsen saavel i physisk som i moralsk Henseende, rigtig Lærebogs-Snak, hvor vi faae en Deel at høre om Judas Ischarioth og om Bonaparte etcetera, han mener, det er det der omtrent kommer ud, at Djævelen hemmer og satiriserer enhver overdreven Stræben, og er altsaa paa en Maade det jeg vilde kalde Ironien. Hvorvidt dette nu stemmer med Göethes Anskuelse af Mephistopheles er et andet Spørgsmaal. Interessantere er derimod for at faae et Blik ind i Goethes Individualitet, hvad han anfører af dennes egne Skrifter om det Indtryk Jordskiælvet i Lissabon gjorde paa ham, og et Citat han anfører af ham: »Daß gewönlich, wenn unser Seelenconcert am geistigsten gestimmt sei, die rohen, kreischenden Töne des Weltwesens am gewaltsamsten und ungestümsten einfallen, und der in Geheim immer fort waltende Contrast, auf einmal hervortretend, nur desto empfindlicher wirke«. – Med Hensyn til det Geisterchor (pagina 82 i min Udgave af Faust) findes følgende Yttring hos S:, som jeg ikke kan forstaae, medens jeg heller ikke selv forstaaer Stedet hos Goethe: »Ein unsichtbares Geisterchor läßt sich ironisch über diese Zerstörung der schönen Welt vernehmen und fordert unsern Helden auf sie prächtiger in seinem Busen aufzubauen, um als Halbgott den neuen Lebenslauf zu beginnen.« Nu vil S. søge at udvikle, hvorledes den uhyre Modsætning der findes i Livet, og som Digteren har concentreret i sin Mephisto, hvorledes og i hvilket Forhold den vil udøve sin Magt over Faust, saaledes, at den, om end i det Enkelte ødelæggende og forstyrrende dog i det Hele vil vise sig som heldbringende. Faust maatte naturligviis ved sin excentriske Stræben være hjemfalden til en Nemesis, og denne Nemesis bliver af en storartet Natur i Forhold til en saadan Character. Mephistopheles paatager sig nemlig fra det Øieblik af, hvor Faust troer uden Haab at maatte fortvivle, at bibringe ham Behag til denne Verden, at overbevise ham om, at den af ham udskieldte Verden ikke er at foragte. Først efter al Nydelse, efter en virkelig iforveien indtraadt Tilfredsstillelse og Mættelse skal F: fortvivle, idet han tilsidst maa erkjende, at Djævelen dog har bedraget ham, at han i egen Lidenskab og Overilelse for tidligt og for snævert har afstukket Grændserne for sin Herlighed. F: skal tilsidst i Følelsen af forspildt Menneske Ære gaae under, han som i Overilelse viste de største Goder, der stod hans Frihed og Besindighed aaben, fra sig, og kun \tlg{tilegnede} sig Existentsens blendende og til Nyden skikkede Side. Nu behandler Mephisto sin Mand med uhyre Ironie, og stiller sig imod en saadan fortvivlende Daare i den avantage, at han ikke engang behøver hine Fornufts og Aands høieste Gaver, for at lade ham finde Livet behageligt. F: fordømte daarligen Livet og nu paatager Mephisto sig at bibringe ham Behag til Livet, idet han berøver ham alt Fornuftigt, ja endog hans Samvittighedsfred. Spørge vi nu, hvad det er for en Sidevei, Mephisto vil føre F. paa, da er det Adspredelsens Vei, hvor Mennesket ganske føler sig som Menneske, idet han fri for Anstrængelser kun tager Livet som en let Nydelse. Naar vi nu fremdeles spørge, hvorledes Mephisto vel har turdet paatage sig at styrte en saa høit stræbende Aand, da kan det kun skee ved et psychologisk Mesterstykke. Han styrter ham ikke umiddelbart i Fornøielse, det vilde under F. nuværende Stemning være uden Følge. Han maa først see Andre nyde og paa et lavere Trin, hvor han selv holdes i en derover svævende Tilstand og halvt ironisk nyder det, ved at vide sig bedre end de Andre. Naar nu dette Syn af den lavere Classes Nyden – forbunden med en halv Mystification af dem fra Mephist: Side, som tjener F. til Moroe, har lidt efter lidt gjort ham fortrolig med det farlige Element, da maa Mephisto paaføre ham en nærmere menneskelig Interesse, dog saaledes, at den kun overrasker ham i sin Nyhed, ikke tilfredsstiller ham, og dertil vil vel Kjærlighed være tjenligst. Kun maa denne menneskelige Interesse flyde af sin sandeste Kilde det vil sige af en virkelig, ufordærvet, quindelig Sjæl. Faust vil snart blive kjed af dette Forhold og Mephisto behøver blot at passe paa, at F. ikke for tidlig bliver sig den sig indledende Katastrophe bevidst, men maa adsprede ham paa andre Maader. Han har kun at sørge for, at F. derved maa forekomme sig selv skyldig, og at da den uendelige Elendighed paa eengang træder op for ham. Ved Synet deraf vil han ikke aflive sig selv; thi dertil er han for svag og slap til; men hans Samvittighed er besmittet, og hvad bliver nu tilovers for ham, efterat han allerede halvironisk er vant til at see al Nydelsessyge, Raahed og Yppighed, bliver trykket af en hemmelig Skyld, end med al Kulde og Ligegyldighed at gaae videre paa Ironiens Vei, og nu under Foragt for sig selv og hele Verden at naae Lystens fulde Maal og selv bestige Lystens Throne og Middelpunct og paa denne Maade at blive et eneste Exempel paa en Virtuos i al jordisk meest forfinede Adspredelse. F. Nydelse maa man heller ikke tænke sig af den aller crasseste Art. Man tænke sig ham med alle Lysters Trylle og Komandostab sættende sig ud over Tid og Rum, nydende netop det Aandrige Pikante, følende sig snart agtet snart frygtet i sin verdenshistoriske Situation, og man kan vel tænke sig, at han vilde ligge under for den Fristelse ikke at ville forlade Verden mere – Hvad nu den Discipel angaaer, der kommer for at belæres af Faust, da bemærker S: rigtigt: der staaer en Faust i Miniatur for os, eller med andre Ord hiin Grundkarakteer i Menneskeslægten stiller sig frem i et ungt Menneske, og saaledes opfyldes allerede her hiint Ord af Mephisto:

»Und immer zirkulirt ein neues frisches Blut.

So geht es fort, man möchte rasend werden«.

Thi neppe er Mephisto i Begreb med at skaffe den gamle Daare i Faust tilside, før der allerede voxer ham en ny i den skjæggeløse Moders Søn.

\emph{De følgende Scener } udføres saa godt som ikke af S:, i Anledning af Branders Sang om Rotten og Mephistos Sang om Loppen, anstiller han en Slags Sammenligning, som falder ud til Fordeel for Mephistos, da han mener, at denne er finere og mere satirisk. I Scenen i Hexekjøkenet (Angaaende Stykket pagina 122 i min Udgave Das ist die Welt etcetera) »Sehr erbauliche Betrachtungen von einem Meerkater über Glanz, Vergänglichkeit und Holheit irdischer Dinge. Er ist leben geblieben, der alte Practicus, und hat Alles durchgekostet; allein er traut dem Sohne nicht dieselbe Geschicklichkeit zu, daher vermacht er ihm als Testament den moralischen Rath, sich der Weltlust zu enthalten, um nicht vorzeitig sterben zu müssen.« En saa barock Spøg som det synes, at Digteren indfører Marrekaten og dens Familie i Hexekjøkenet, saa mener dog S:, at deri ligger en dyb sædelig Betydning, at nemlig Marrekaten, der som bekjendt forener de fleste Dyr-Svagheder i sig, staaer som Repræsentant for det meest Dyriske i Msknaturen. Og Digteren lægger dem ogsaa alt i Munden, hvad der udgjør den lave Msknaturs Bestræbelser og Ønsker; lav Begjerlighed, Spil, Havesyge, Lumskhed, tyrrannisk Undertrykkelse, oprørsk Sind, tøilesløs Frækhed i Tale og Handlen. –

Nu følger Scenen med Margrethe, hun er, som S: rigtig bemærker, »leider keine Kokette, und die unerfahrne Unschuld unterliegt solcher frechen Zudringlichkeit noch schneller als die ausgelernte Buhlkunst. Denn eben weil sie die Unschuld ist, und es nur ehrlich meint, so stellt sie sich selbst in ihrer guten Meinung der Schamlosigkeit zum Anwald auf, und sucht günstig ins Sittige zu deuten, was nur Frechheit ist und bleibt.« Sandseligheden overrumpler nu F:, der kun var vant til at leve blandt Bøger. Lidenskabens første Opblussen søger Mephisto at opflamme end mere ved at gjøre Forhindringer. Nu træder da Margrethe op, der, efter Mephistos egen Tilstaaelse, er et uskyldigt Barn, som Intet har at skrifte. Her anstiller nu Sch: en meget utidig Undersøgelse om Mefistofeles mener det oprigtigt eller ei, hvilket første han erklærer sig for, og faaer da heri et nyt Bidrag til sin Djævletheorie. Smukt er derimod hvad S: siger angaaende Scenen (pagina 141 i min Udgave »Es ist so schwül, so dumpfig hie«) »in dem Weben sympathetischer Gefühle, das uns der Dichter schildert, dürfen wir uns nicht wundern, daß Margrete bei der Rückkehr eine geheimnißvolle, unerklärliche Ahnung von der stattgehabten Anwesenheit der beiden Fremden hat. Es ist zart vom Dichter gewesen, diesen Geist uns fühlen zu lassen, der dem Menschen beigegeben ist, ihn in den bedeutendsten Fällen des Lebens umweht, und ihn als feinstes Gefühl das Geschick seiner fernen Zukunft vorempfinden läßt. So schaudert Margarete ahnungsvoll beim Eintrit ins Zimmer, dessen Luft ihr schwül vorkom̄t, obwohl es draussen nicht warm ist. Sie fühlt sich beklommen und allein, und wünscht der abwesenden Mutter baldige Ankunft.« I Anledning af Scenen mellem Mephisto og Marthe Schwerdtlein (pagina 150 i min Udgave) mener S:, at man ikke kan andet end tænke sig en Sammenligning mellem den Art af Kjærlighed, som fandt Sted mellem Marthe og hendes nu formeentlig afdøde Mand, og den, som finder Sted mellem F: og Grethe: den Art af Kjærlighed som uden mindste Spor af Betydning lever i visse Former, og den der i opblussende Varme skrider ud over disse; og S: mener, at denne »köstliche Vexirscene« Mellem Mephisto og Marthe tjener til at mildne den tragiske Katastrophe med Margrethe.

\emph{S: bemærker med Hensyn til (pagina 182 i min Udgave)} Margrethes Angst for Mephisto, at denne Mands Physiognomie nu først falder hende paa. Hun saae ham jo allerede tidligere hos Fru Marthe og ved deres første Sammenkomst, uden en lignende Bekymring. Er der maaskee med Margr: Indre, som før var reent og uskyldigt foregaaet den Forandring, at det Onde, dets Tilvær, det endog skjulte Udtryk i en Tredies Mine ikke længer er hende ligegyldigt. Det er desværre sandt, at Mennesket ved Forstyrrelsen af hiin rene Følelse taber hiin Uskyld, der lader ham see dristigt og frimodigt lige ud, og at da ganske ligegyldige Ting udvortes fra kun synes ham en Anklage, en Bebreidelse. »Man kann sagen, wenn Mephistopheles für Gretchen früher nicht da war, so fängt er jetzt an für sie Wirklichkeit, Bedeutung und Kraft zu gewinnen. Auch spricht Mephisto das von seiner Seite durch den wegwerfenden Titel »Grasaffe«, den er ihr giebt, aus, während er sie früher »ein gutes unschuldiges Ding« nannte. Es ist überhaupt diese Scene ungemein merkwürdig, wegen der Kräfte und Mächte, die dabei wirksam sind. In das Heiligste, Himmlische, was hier diese beiden freundlichst aufgeschlossenen Gemüther beschäftigt, wirft sich ein satanischer Zug, ein geheimes, ahnungvolles Grauen vor der Hölle.«

\emph{Hvad Samtalen mellem de to Elskende angaaende Religionen betræffer,} da bemærker Sch:, at den aabenbarer Forskjellen mellem begge Kjøns religieuse Trang, den mere almindelig religieuse Anskuelse og den mere individuelle. Derpaa gaaer han over til at betragte Christendommen i Almindelighed, hvorledes de forskjellige Hovedretninger i aandelig Henseende, der har betinget de forskjellige Religioner, nu gjentagende sig i Christendommen betinger de forskjellige Confessioner. Dog synes det mig vedkommer denne hele Undersøgelse slet ikke Stykket især i en saadan Vidtløftighed. – En Charakter, som Gud, Natur i deres Umaadelighed ikke tilfredsstille, er vel so sjelden, at vi nok kunne tiltroe ham en Løben Storm paa Hexebjerget. Imidlertid vil det være smaaligt og modsigende, hvis han ikke formaaede at vikle sig ud af disse Tryllebjergets Irgange, og vedblev i sin gamle Natur, der hverken finder Ro i Godt eller Ondt. Derimod er en saadan Gang til Bloxbjerget med største Nødvendighed begrundet i en saadan Grundcharakter. Det er den samme Kunstnødvendighed, som før lod os see F: i Auerbachs Kjelder, i Hexekjøkenet. De samme Motiver, der overhovedet lode F. slutte Forbund med det onde, benægtende Princip, føre ham op paa al Forvirretheds Top. Det laae nemlig i Djævelens Plan, at bestikke F: ved en ironisk Tilfredsstillelse, idet han paa denne Top af al Forvirring vedblev at bevare en Art Storhed og Overlegenhed. I denne Anskuelse af den store Masse af Ondt bringer Satan ham paa det farlige Punct, at han tilsidst seer sig sat ud over al Forskjel mellem Godt og Ondt. Naar nu det er det egentlig Absolute eller Absoluteste, hvorefter Naturen stræber, saa see vi, at en Bestigelse af Bloksbjerg, for at komme til Følelsen af en saadan Hæven sig ud over sig selv, var uundgaaelig nødvendig for F:. Hiin Ubetingethed, hiin Absoluthed er det nemlig, hvoraf vor Mand lider. Mephisto øiner her det Punct, at han ved at lade Forskjellen mellem Godt og Ondt forsvinde, ved at overskride den, kan gjøre F. modtagelig for hiin høieste, ubetingede Alhed, der forsaavidt da al Forskjel deri forsvinder, skuffende ham bringer ham nær til en vis Art Guds-Lighed. Vel er det ikke andet end den ubetingede Tomhed og Intethed, der som Alhed bliver tilbage efter Ophævelse af al Forskjel. Herved vil nu Mephisto baade tilfredsstille og aandelig tilintetgjøre F., idet han ad denne Bro fører ham over paa den fuldkomne Vilkaarligheds Gebeet.

Men den egentlig Satans Streg, som Mephisto spiller med F. ved denne Bloksbjergs Scene, er hiin hemmelighedsfulde Vision af Gretes Skikkelse, som bliver synlig for F. paa Bloksbjerg. Derved griber han ham dybt ind i hans inderste Samvittighed og knytter ham med i Begyndelsen lette, lidt efter lidt centnersvære Baand til den uoverskuelige Skare af Ondt. Først advarer han F: paa Skrømt for ikke al for nøie at betragte den Idol, derpaa gjør han ham latterlig for denne Carricatur af hans Samvittighed. Er det paa denne Maade lykkedes Mephisto, at vildlede F. Samvittigheds Stemme; saa vil Synet af den uendelige Elendighed kun foraarsage ham en kort og udvortes Rystelse. Han vil meget mere nedbøiet ved den store hemmelighedsfulde Skyld tillige være helbredet for al anden Tanke, tilmed naar han efter hans Sinds og Fornuftigheds Storm ved Opvaagnen mærker, at man endnu beholder Kræfter nok til at leve, og at, hvor Uskyld afslutter Livets Kreds, dets Kræfter pulsere derud over. Ved saaledes at skride ud over de enge Grændser, som Fornuft, Samvittighed sætte for Mennesket, vil F. virkelig føle sig Mephisto i Sandhed forbunden, da han herved aabner ham den grændseløse, umaadelige Udsigt, som alene kan tilfredsstille ham. Thi det er F. Ulykke, at han ikke kan finde sig ind i de almindelig Livsindskrænkninger. Valpurgisnatsdrøm tjener til Udstaffering af den ideale Region paa Bloksbjerg. \emph{En Bemærkning som findes hos S. pagina 327 og 28 vil jeg afskrive:} Indem übrigens der Blocksberg nicht eine Zeitepoche, sondern Aeltestes und Jüngstes in sich befaßt, fängt hier jene Aufhebung des gewönlichen Zeitbegriffs schon sehr bedeutend an, an die wir uns von nun an immer mehr werden gewöhnen müßen. In dem folgende Verlaufe nämlich der Schicksale unseres Helden, wie sie uns der zweite Theil vorlegt, verschwinden Zeit und Oertlichkeit immer mehr, und F: wandelt sich aus einer wirklichen, historischen Person, als die er uns der Hauptsache nach im ersten Theil in Ton und Haltung erscheint, in eine mythisch allegorische und symbolische Persönlichkeit um. Alle Schicksale unsers Helden verlieren den individuellen Bezug, und nähern sich einem Allgemeinen, ja Allgemeinsten der Weltgeschichte«. Idet nu S: til Slutning, før han gaaer over til den anden Deel af Goethes F:, i den 10de Forelæsning tager den hele Udvikling af Mephistos Plan op, omtaler han og paa sit rette Sted F: Kjærlighedshistorie, og viser, hvorledes herfra udviklede sig en af de dybeste »Entzweiungen« i hans Indre, som han ikke kunde blive klar over, om den var ham paaført mere ved egen Brøde, eller ved et ulykkeligt Sam̄enstød af Omstændigheder som for Exempel. ved Valentins Mord, som styrtede ham i en Feil og det af den Art, at den berøvede ham alt Mod til med Uskyldighedens Eftertryk at fortsætte sin Klage over Nægtelsen af hans høieste Ønsker mod de høiere Magter. For at nu ikke Indtrykket af Gretes Elendighed ganske skulde gjenemløbe den ideelle Region, og maaskee gestalte sig til den dybe og inderlige Anger, som vi tilsidst see Margrethe grebet af, saa er vor Mands Aand ved Bestigelsen af Brocken allerede indtaget af et saadant Indtryk af det Onde i Anskuelsen af dets Masse og Magt, at han gjerne vil tage denne Anskuelse til Hjælp, for at forestille sig denne Skabnings Elendighed som et enkelt Tilfælde og et Minimum af hiin Slethed, som hører hjemme paa Jorden, forsaavidt den udgaaer fra Menneskets Personlighed og Villie. Kan han desuagtet ikke ganske overvinde Braadden, saa er dog saa meget vundet, at han i alle lignende tilkommende Tilfælde vil have mindre Mod til skarpt at undersøge, om høieste Fordringer derved blive tilfredsstillede eller ikke. Han vil være desto tilbøiligere til at optage ethvert Nydelsens Tumleliv, for dermed at overdøve hans dybeste, saarede Indre. Alt dette indrømmet, saa seer man, at vor Mand for at unddrage sig de ideelle Bebreidelser, er nødsaget til ja er paa bedste Vei til, umærkelig og næsten mod sin Villie at komme til en real Tilfredsstillelse.

Faust 2den Deel.

Her vil jeg nu blot erindre til egen Efterretning, at Schubarth ikke har den fuldstændige Udgave af F:, som fra Digterens Haand først blev færdig 1831; men maa nøies med de adspredte Fragmenter. Han gjør opmærksom paa, at man har fundet en Fortsættelse af F: deels unødvendig deels umulig, og han mener, at begge Dele paa den bedste Maade ere modbeviste af Digteren. Han erindrer om, at hvad der var udsagt i Prologen kun for en Deel var fuldført »Es irrt der Mensch, so lang er strebt«; men den anden Deel er det ikke »Ein guter Mensch in seinem dunkeln Drange, Ist sich des rechten Weges wohl bewußt«. Ogsaa Mefistofeles har kun bragt sin Mand til Anskuelse af Nydelse, og sat ham i den Nødvendighed i Adspredelse at finde et Hjælpemiddel mod hans hemmelige Halvskyld. Men endnu er der ikke forekommet Noget i den første Deel, hvorved F: selv søger at sætte sig i Middel- og Brændpunctet af Nydelse paa Jorden i den meest concentrerede og fine Betydning. Havde man holdt fast derpaa, mener S:, da havde man omgaaet slige Vildfarelser. Han anfører derpaa en Yttring af Goethe i Kunst und Alterthum i Anledning af Mellemspillet »Helena«: »Darüber aber mußte ich mich wundern, daß diejenigen, welche eine Fortsetzung und Ergänzung meines Fragmentes unternahmen, nicht auf den so nahe liegenden Gedanken gekommen sind, es müsse die Bearbeitung eines zweiten Theils sich nothwendig aus der bisherigen kümmerlichen Sphäre ganz erheben, und einen solchen Mann in höhern Regionen durch würdigere Verhältnisse durchführen.« Til at udvikle og fremstille F: i dette andet Stadium (Nydelse) bruges, mener S:, Scenen ved F: Opvaagnen, ved Keiser-Hoffet og den store Hoffest. Nu følger Fragmentet Helena. I hvilken Henseende han anfører et Sted af Goethe, hvor han gjør opmærksom paa, at efter Sagnet forlangte F: at see Helena. S. vedbliver efter en Exposition af det af Goethe i dette Fragment Givne: Naar vi nu overveie, hvorledes H: og F: her træffe sammen, saa maae vi tilstaae os, at det vel ikke kunde skee paa en mere passende Maade, end idet F: viste sig som Middelpunct i den uhyre Bevægelse af Norden mod Syden, som fra Folkevandringen af gjenem hele Middelalderen til den nyeste Tid har vedligeholdt sig deels igjenem hele Folkemasser deels i Enkelte, som Udtryk for Nordens brændende Længsel og Attraa efter Sydens Skjønheder i ligesaa eventyrlige som yndige Skikkelser. Kun som en saadan Folkefyrste, staaende i Midten af denne Bevægelses meest romantiske Vending, turde F. vove, at anholde om den verdensberømte antike Skjønheds Gunst og kun saaledes som Romantikens Ungdoms-Fyrste viser han sig ikke uværdig til at træde ved Siden af Fylden af al Skjønhed.

Til Slutning gjør S. et Forsøg paa, at construere Finalen til hele Værket. Den er der imidlertid ikke noget synderligt ved. Den Slutnings Bemærkning derimod, hvormed hele hans Skrift ender, er ret god: »Wenn das Streben nach dem Absoluten sich in der neuern und neuesten Menschheit auf mancherlei Weise hervorgethan, und auf dem philosophischen Wege der ernsthaftesten Behandlung unterworfen zu werden, nicht für unwerth befunden worden; so ist es merkwürdig, wie der Dichter dasselbe als einen Wahn von sich zu weisen scheint, dem auf ernstem Wege positiv durchaus nichts abzugewinnen sei. Auf einem scherzhaften Wege hingegen in gründlich, tüchtiger, verneinender Behandlung gewähre es die glänzendsten Vortheile und Befriedigung. Es verleihe nämlich einem, in seiner allgemein menschlichen Begrenztheit sich unbehaglich fühlenden Charakter eben das Recht, sich über alle Grenzen und Schranken, die dem Menschen nothwendiger und zufälliger Weise gezogen sind, hinaus zu versetzen, um mit der Phantasie in jener ungebundenen, dog gefälligen Willkürlichkeit und Zügellosigkeit den nie endenden Wettstreit zu beginnen. Statt daß also das Absolute in die philosophische Region aufzunehmen sei, wo es stets starr, trocken, todt, ungnießbar verbleibe, und verrückte Combinationen veranlasse, gehöre der Begriff desselben recht eigentlich der Poesie an, die ihm allein Giltigkeit zu verschaffen und ihn durch das grenzenlose Schwärmen der Einbildungskraft, dem er entspricht, erst lebendig, wirksam, wahr, so wie ergötzlich und heiter zu machen im Stande sei.« –

\emph{Jeg vil nu gaae tilbage og optage S: Indledning.} Han bemærker, at Goethe i dette Værk fornemlig har brugt Humor: »Eine humoristische Behandlungsart und die damit verknüpfte Ansicht tritt unstreitig überall ein, wo viele Dinge, werthvolle und unwerthe, unserer Anschauung theils zugleich, theils bald nachzeitig entgegen gebracht werden, so daß der wohlbekannte Unterschied derselben zu schwinden scheint, und für den Augenblick durch den Bezug auf etwas noch Höheres, gleichsam als auf ein Absolutes, das bald deutlicher hervortritt, bald mehr im Hintergrunde gehalten wird, seine gewönliche Giltigkeit verliert und in dem bekannten Resultate null wird.« Han anfører som Exempel, det at staae paa et høit Bjerg og derfra skue ned paa den under sig liggende Dal. »Was nun aber dieses Höhere gerade vorstellt, worauf der sonst bekannte Unterschied und Werth der Dinge als verschwindend bezogen wird, bildet den verschiedenen Charakter und die abweichende Beschaffenheit des Humors, der damit sich eben so in die höchsten Regionen erheben, als bei einem sehr trivialen niedern Fluge beharren kann, so daß es nichts Gemeines giebt, was fratzenhaft ausgedrückt, nicht humoristisch aussehn und wirken könnte«. »Jedenfalls wird eine gewisse Breite des Weltzustandes, eine gewisse Mannigfaltigkeit der Gegenstände und Situationen vorausgesetzt, wenn der Humor sich entwicklen soll. Daher es erklärlich ist, warum diese Lebensansicht und Behandlungsart der Dinge bei den Alten weniger angetroffen wird, und erst in der neuern Welt sich entwickelt hat«. Derpaa gjør han en Distinction mellem naiv og realistisk og en mere ideell, phantastisk, længselsfuld Humor. »Der Unterschied beider aussert seine Wirksamkeit vorzüglich in Absicht auf jenes Absolute, den sonstigen Unterschied der Dinge theils aufhebende, theils neu Contrastirende. Wird namlich dieses in irgend etwas Gegenwartiges, noch Erreichbares, Diesseitiges gesetzt, so entsteht jener naive Humor; wird aber dabei in die Ferne, in ein Jenseits, in überirdische Regionen hinaufgegangen, so entspringt der sentimentale Humor.« for Exempel. Jean Paul. Han viser, hvorledes i Digtet Faust Digteren bestandig gaaer ud fra den Tanke, at Livet i dets laveste, meest narragtige og forkerte Manifestationer, dog bestandig er en uskateerlig Gave. Det gjør hans Humor naiv. – Fra pagina 80-84 leverer S: selv en Oversigt over Indholdet af den 2den Forelæsning. En vis eensidig Forkjærlighed for et enkelt Værk af en i Sandhed stor Digter bør man ikke taale, dennes Storhed er at søge i Totaliteten af alle hans Værker. Nu søger S: at møde adskillige Indvendinger mod Værket, saaledes af A: W: Schlegel i »Vorlesungen über dramatische Kunst und Litteratur«, som mener at Digtets Form gaaer langt over den theatralske Fremstilling. Han viser, at nogle blot have opfattet Digtet som en Klage over at Livets høieste Nydelser bleve ham nægtede, og at Lord Byron fra dette Synspunct har reproduceret Stoffet og Indholdet i F. –

I 3die Forelæsning søger S. at udvikle Dedicationen, Forspillet og Prologen. Her findes Intet mærkeligt. pagina 103 findes et nyt Bidrag til hans Djævlelære, som jeg alt i det Foregaaende har udviklet. pagina 106 ligeledes et Bidrag. Han gjør opmærksom paa, at andre store Digtere for Exempel. Klopstock, Milton og endog Lord Byron i hans Kain har opfattet Djævelen fra en anden Side. Goethe mener han nærmer sig mere til Opfattelsen af ham i Hiobs Bog. Resultatet bliver da omtrent, at Djævelen har paataget sig at være Profos ved den guddommelig Justits, en Forretning som jeg egentlig ikke seer der ligger Noget Djævleligt i, især da denne Profos mener det godt med Menneskene

d: 2 Sept: 36.

\emph{[i margen:]} Eet af dem kan jeg dog ei lade være at afskrive: »Man hat oft gesagt und mit Recht, der Unglaube sei ein umgekehrter Aberglaube, und an dem letzten möchte gerade unser Zeit vorzüglich leiden. Eine edle That wird dem Eigennutz, eine heroische Handlung der Eitelkeit, das unlaügbar poetische Product einem fieberhaften Zustande zugeschrieben; ja was wunderlicher ist, das allervorzüglichste, was hervortritt, das allermerkwürdigste, was begegnet, wird so lange als nur möglich ist, verneint«. –

\emph{[i margen:]} Hermed maa sammenlignes hvad Falck siger om Goethes Antipathie mod Hunde.

\emph{[i margen:]} Sin Theorie om Djævelen (en Theorie, som idetmindste ikke synes mig af stor Vigtighed med Hensyn til Digtet) som den, der gaaer Haand i Haand med Vorherre, som ikke er andet end den til Bevidsthed bragte Ironie, synes Schubarth at ville bevise pagina 223. ved at gjøre opmærksom paa den som det synes ham ellers uforklarlige Omstændighed, at han skjøndt han søger at forføre Faust, dog yttrer stor Agtelse for Fornuften (i Monologen, medens Faust gaaer bort og Discipelen melder sig »Verachte nur Vernunft« etcetera), og dog igjen kritiserer den menneskelige Viden paa en meget ufordeelagtig Maade. Men herved maa mærkes, at det for det første maa oplyses i hvilken Betydning Mephisto roser Fornuft, da det jo ogsaa kunde være uden i Samklang med Gud; og dernæst dersom Mephisto virkelig som det synes af mange Yttringer hos Schub:, efter hans Mening, vil det Gode, saa gjør han jo i sin Apologie for det Onde i Verden enten Djævelen til Vorherre eller Vorherre til Djævelen; eller hvilket synes at fremgaae af andre Steder hos Schub:, Djævelen ikke saa meget vil det Gode, som forhindre det Excentriske i enhver Retning, fordi det er ham ubehageligt; thi da gjør han Djævelen til en Epicuræer i stor Stiil. Han har ikke indseet, at om end det Onde i Verden lader sig see som Middel i Guds Haand, det dog ikke er det paa Djævelens om end grandieuse dog ikke alvidende Standpunct.

\emph{[i margen:]} S. K. d. 27 Aug. 36.
\end{quote}
\dcite{BB:7  (1836; Pap. IC96; SKS 17, 77)}
% BB:7  (1836; Pap. IC96; SKS 17, 77)

\begin{quote}
Forholdet mellem Loven som det blot udvortes, objectivt Givne og dens Optagelse i Individet kan man for Exempel see ved en Læges diætetiske Forskrifter, deels er det næsten en Umulighed at følge dem da han egentlig burde staae for i hvert Moment at sige, hvad man burde gjøre, deels om man endogsaa kunde følge dem bringer de ingen »Salighed«, da det egentlig kommer an paa, at man dybere har \tlg{tilegnet} sig det diætetiske og netop nu ved en sikker Takt afgjør det enkelte Tilfælde.

April 1836.
\end{quote}
\dcite{Papir 48:75  (1836-04; Pap. IA152; SKS 27, 107)}
% Papir 48:75  (1836-04; Pap. IA152; SKS 27, 107)

\subsection{1837}

\begin{quote}
Dersom En efter at have læst Afhandlingen vil sige, at jeg vel taler om \emph{den Kunst at fortælle Historier } men i hele Afhandlingen snarere synes at ivre \emph{derimod,} saa vil jeg ikke \emph{ubetinget} indrømme det, da jeg dog kun har ivret mod Misbrug, og tillige gjøre opmærksom paa, at jeg har taget det Ord: »fortælle Historier« i en vidtløftigere Betydning om alt det, hvormed man beskjæftiger Børns Aand uden for de egenlige Læretimer, og som man ikke ligefrem kan kalde Leg, i hvilken Henseende dog vistnok det at fortælle Historier spiller en Hovedrolle.

At saa Mangfoldige give sig af med \emph{at fortælle Børn Historier} er en naturlig Følge af, at der er en saa stor Mængde Børn og af den i dem saa dybt grundede Lyst til at høre fortælle, og dog er der saa Faa, der egenlig have Talent dertil; som en Følge deraf foraarsages der megen Skade derved. Der gives to Maader at fortælle Børn Historier paa, som er at anbefale; men mellem disse to ligger ogsaa en Mangfoldighed af Afveie.

Den \emph{første} er den, som Ammer – og hvad man kan sætte i Categorie dermed – ubevidst følge. Ved dem gaaer en heel phantastisk Verden op for Barnet, og Ammerne ere inderlig overbeviste om og troe paa Sandheden af deres Historier1, hvilket nødvendigviis maa bibringe Barnet en gavnlig Ro, hvor phantastisk forøvrigt Indholdet selv kan være; først naar Barnet selv kommer paa Spor efter, at Vedkommende ikke troer paa sine Historier, først da virke de skadeligt – dog ikke ved Indholdet selv men ved deres Usandhed med Hensyn til Fortælleren – paa Grund af den Mistillid og Mistroiskhed, Barnet efterhaanden udvikler hos sig.

Den \emph{anden} Maade lader sig kun realisere af den, der til fuldkommen Klarhed har reproduceret Barndomslivet, – veed, hvad det forlanger, – veed, hvad der er godt for det, og nu fra sit ophøiede Standpunct byder Børnene en Aandsføde, der er dem tjenlig, – veed at \emph{kunne være} Barn, medens Ammerne igrunden \emph{ere} Børn – (at Børn have Leilighed til at nyde godt af begge Maader, er saare gavnligt, og man troe ingenlunde, at det andet Standpunct aldrig vil anerkjende det første. Nei, tvertimod hvad der altid er Tilfældet med Halvstuderede: at bortskjære Udviklings-Veien, saa yder den, der har en moden Livsanskuelse, hiin sin Anerkjendelse).

Her er nu ikke en lang Forberedelse. Manden kommer hjem fra det travle Contoir, skifter Strømper, faaer sig en Pibe, kysser Mutter paa Kinden og siger: »Naa, min søde Glut,« (det er for at vænne Børnene til kjærlig Omgang) og nu indtræder der da en Begivenhed, som man seer afbildet paa de fleste Børnebøger: Onkel Frands, der fortæller Historier, hvortil Børnene have glædet sig hele Formiddagen, og lille Frits og Marie komme løbende og klappe i Hænderne: »Onkel2 Frands fortæller«. Moderen grupperer sig mellem Børnene med den Mindste paa Armen og siger: hører nu kjønt efter hvad Jeres kjære Fader fortæller!

Dette med Hensyn til Fortællingens \emph{Ramme;} nu til vor \emph{Fortæller.} Al almindelig Beskjæftigelse for Børn udenfor de egenlige Underviisningstimer – og ogsaa i disse saa meget som muligt – bør være \emph{socratisk;} man maa vække Lysten til \emph{at spørge} hos dem, istedetfor at et fornuftigt Spørgsmaal, som maaskee gaaer uden for Onkel Frands' Kreds af Viden eller paa anden Maade er ham ubeleiligt, bliver afviist med de Ord: »den dumme Dreng, kan han ikke tie stille medens jeg fortæller?« – og Moderen for at forhindre alvorligere Optrin forsikkrer, »at han aldrig vil gjøre det mere«. Det, det beroer paa, er at bringe \emph{det Poetiske paa alle Maader i Forhold til deres Liv,} at udøve en Tryllemagt, ved den meest uventede Leilighed pludselig at lade et Glimt see og igjen forsvinde; man skal ikke henlægge det Poetiske til visse Timer og visse Dage. Omkring et saadant Menneske springe Børn ikke som bengelagtige Kalve med dinglende Been og klappe i Hænderne, fordi de \emph{skulle} høre en Historie; \emph{ham} nærme de sig med et aabent, frimodigt, tillidsfuldt Væsen, betroe sig til ham, indvie ogsaa ham med i mangen lille Hemmelighed, fortælle ham deres Lege, og han veed at gaae ind deri, veed ogsaa at give Legen en alvorligere Side. Børnene falde ham aldrig til Besvær, plage ham aldrig, dertil have de for megen Agtelse og Respect for ham3. Han veed, hvad de bestille i Skolen; han læser ikke Lectier med dem, men i Stilhed erkyndiger han sig om hvad de læse, sætter sig ind deri, ikke for at prøve om de kan det, ikke for at tage et enkelt Parti og dramatisere det for dem, ikke for derved at kunne give dem Leilighed til, naar der er Selskab, at glimre; – men for pludselig at lade et Glimt deraf træde frem, paa en individuel Maade sætte det i Forhold til det, som de nu ellers netop beskjæftige sig med, dog aldeles en passant, saa Barnets Sjæl derved electriseres og føler ligesom Allestedsnærværelsen af noget Poetisk, som vel er ham kjært, men som han dog ikke tør træde for nær4. \emph{Derved} næres hos Børnene en bestandig aandelig Bevægelighed, en permanent Opmærksomhed paa hvad de høre og see, en Opmærksomhed, man ellers maa fremkogle paa \emph{udvortes} Maade ved for Exempel at lade Børnene komme ind i et meget oplyst Værelse fra et mindre stærkt oplyst, hvor Onkel Frands sidder – ved at kjede dem hele Dagen med Fortællingen om »hvor rart det er at høre Onkel Frands fortælle« og saa videre

Men uagtet den Klarhed, her hersker, kan der dog let indtræde en vis Sentimentalitet, idet man glemmer, at Manddommen har, hvad Barndommen lovede; man synes, især naar man har med meget opvakte Børn at gjøre, at den dog lovede noget Mere, thi derved griber man ængstende ind i deres Liv, Noget, der virkelig kan komme af denne Grund og ikke altid af trivielt Flæberi. Disse idelige Forsikkringer: »I ere lykkelige, men naar I blive ældre5, saa kommer der Sorger« og saa videre virke \emph{skadeligt,} da de, forsaavidt de slaae Rødder hos Barnet, bibringe det en underlig Angst for, hvor længe det dog endnu kan blive ved at være lykkeligt (og derved \emph{ere} de allerede ulykkelige); – eller forsaavidt denne idelige Jeremiade intet Indtryk gjør, skader den naturligviis som al anden utidig Snak. – Denne Ubestemthed kunde synes at stride imod en vistnok meget rigtig Fordring til Strenghed og skarp Begrændsning: dette skal nærmest repræsenteres i Skolen (hiint er jo Adspredelsestimer) i Personligheden selv. Den, der i sin Barndom aldrig har været under \emph{Evangeliet} men kun under Loven, bliver aldrig fri6 – kanskee det er Uret, men der er noget Nobelt deri; medens jo mere \emph{Loven} er udviklet, desto mere fremspire smaa Drillerier, og Intet er saaledes istand til som den at frembringe Kleinmodighed. Der ligger en Magt i Øiet til at fremlokke det Godes Spire og knuse det Onde – men den misforstaaede Strenghed og Tugt, en Datter af Magelighed, vil næsten lade den ene Generation tage Hævn over den anden for de Prygl, den selv har faaet, og de Mishandlinger den selv har lidt, ved at behandle den følgende ligesaa.

– Men skal man da \emph{ikke} fortælle? Jo, Mythologie og gode Eventyr er hvad Barnet behøver – eller: man lade det selv læse dem og fortælle dem og saa socratisk berigtige dem (ved at spørge efterhaanden berigtige dem, saaledes at Barnet nu slet ikke under Læremesterens Tvang corrigeres, men tvertimod synes at berigtige Andre – og den, der ellers forstaaer at behandle Børn, vil vist ikke være udsat for, at det udarter til Hovmod). Men for alting skee det som Impromptu, ikke til bestemt Tid og Sted; Børn skulle tidligen erfare, at Glæden er en lykkelig Constellation, som man maa nyde med Taknemlighed, men ogsaa vide at afbryde itide; og for Alting glemme man ikke \emph{Pointet } i Historien. – (En Afvei, jeg strax her kan berøre, skjøndt den siden kommer igjen, er det: idelig og saagodtsom hele Dagen at fortælle intetsigende trivielle Historier og derved forarbeide disse Romanlæsere, der hver Dag sluge det ene Bind ovenpaa det andet uden et bestemt Indtryk). Derhos fremkalde man en vis Productivitet (tegne eller paa anden Maade) ved selv idet man fortæller paa mangfoldige Maader at sætte det i Forhold til det, Børnene ellers bevæge og røre sig i.

Nu opkommer det Spørgsmaal: \emph{hvad Betydning} har egenlig Barndommen, er den et blot Trin, der kun har sin Betydning i den Omstændighed, at det paa en Maade betinger de følgende Stadier; eller har det selvstændigt Værdi –? Det Sidste have nu Nogle udvidet i den Grad, at de antoge, at Barndommen i Grunden er det Høieste, Mennesket kommer til, og udover det udarter det mere; det Første har havt den Følge, at man deels søgte blot at faae Tiden til at gaae7, – og kunde man ligesom ved Fierkreaturer ved at slutte dem inde i Mørke fede dem saa stærkt som ellers ikke i et heelt Aar, saa vilde man vist gjøre Alt dertil; deels søgte at gjøre denne »Barndommens kjedsommelige Tid« nyttig og fornemlig sørge for deres physiske Velvære. Den høieste Opdragelses-Maxime paa dette Standpunct lyder saaledes: »Den som ikke spiser Formaden op, faaer ingen Eftermad«. – (Hvor tidt forbittres ikke Børns, det er da især Pigebørns Liv derved, at de idelig maa høre, at man slet ingen Gavn har af dem – og saa videre).

\emph{Afveiene} opstaae idet man kommer ud over Ammernes Standpunct og nu ikke fuldender Løbet, men bliver staaende paa Halvveien.

\emph{Første Stadium:} De, der efter at være komne udover det umiddelbare Standpunct, nu istedetfor, som naturligt var, i den modnere Alder at optage Barndommen forklaret i sig, istedet derfor ere forfaldne til »at være Børn« (jævnfør Foryngelses\emph{drik}), disse lange Labaner, som ere saa uskyldige og saa naive, som vilde give Meget til, at deres Skjæg aldrig blev saa stærkt, at de behøvede at rage det af, for altid at kunne være dunglatte, barhalsede Ynglinger, – der i den Grad ere blevne Børn igjen, at de tale som Børn, \tlg{tilegne} sig alle Barnesprogets Vendinger, og som forlængst vilde have bevirket, at vi alle kom til at tale som Børn og skrive som Børn tale: en Carricatur, der vel ogsaa vil komme, naar først den modsatte, som nu er saa hyppig: at Børn ville være gamle Folk, er overlevet. Det er et tragicomisk Syn at see disse lange barnagtige Gliedermænd springe omkring paa Gulvet og ride paa Kjephest med de søde Smaae og høre deres matte Fortællinger om »den uskyldige og lykkelige Barndom«8. – (Jævnfør deres Sammenstød med halvvoxne Pigebørn, der ville være voxne: de parodiere hinanden).

Deres Fortællinger »for Børn og barnlige Sjæle« (poetisk Skyllevand). Findes hiin Feil som oftest hos Yngre, saa findes en lignende Afvei hos Ældre, som »\emph{nedlade}« sig til Børn i den Overbeviisning, at Barndomslivet er saa tomt og indholdsløst i sig selv, at \emph{de} ligesom ville indblæse det nogen Fylde. I Grunden maa nemlig begge forudsætte Barndommens Tomhed, thi ellers vilde den første ikke indlade sig paa at byde den noget saa Modbydeligt, som en god Natur strax maa secernere; eller den anden indlade sig paa at blæse Livets Aande ind i den. – Man tilintetgjøre heller ikke hele Indtrykket ved efter at have fortalt Noget at ende med: men I begriber da nok, at det \emph{kun} var et Eventyr? – Noget, der ogsaa kommer igjen i en senere Tid hos Folk, der slet ikke have Sands for det Poetiske og derfor fordærve Indtrykket af enhver Anecdote og saa videre ved at anstille Undersøgelse om den factiske Sandhed.

Den phantastiske og eensidige Retning, som Fortællingen har taget. Man fandt det urimeligt og skadeligt for Fremtiden at overfylde Børnenes Phantasie med slige Historier, hvorimod det var ret godt at fortælle Noget for at udfylde Tiden og more dem, og nu begynder da, eftersom det jo var blot til Moro og man egenlig ikke gad anvende nogen Tid til at forberede sig9, hiint uendelige Historie-Vrøvl om den Hund og den Kat og saa videre, i den rædsomste Monotonie, som imidlertid Børnene, da de engang ere forvænte, idelig forlange flere og flere Udgaver af, og som derfor med en eller anden vigtig Forandring (for Exempel at engang var det en rød Hund, en anden Gang en sort) stereotyperede vende tilbage10.

Ogsaa det fandt man imidlertid var galt, da jo dog den Tid kunde anvendes bedre, kunde bruges om end under Form af Spøg og Leg til noget Bedre, og heraf udspandt sig tvende Veie: \emph{enten} at danne dem i, som man kalder det, moralsk Henseende – \emph{eller} bibringe dem en nyttig Viden. Følgerne af at gaae den sidste Vei vil jeg dvæle lidt ved. Nu kom der ligesom ved et Trylleslag en Landeplage af naturhistoriske, ikke Lærebøger men: Læsebøger og allehaande Billedbøger for at bibringe Børnene Gloser af levende Sprog, og Onkel Frands fortalte sine Reiser i Africa og benævnede Dyr og Planter med Navne af Systemer, og Forældre og Andre spurgte: hvad en Næse heed paa Fransk? og saa videre, eller man lærte dem at klimpre et enkelt Stykke paa Fortepiano, – og vil man end ved Sligt forhindre Børn fra at være forlegne \emph{ved} at træde frem, saa skal man dog heller ikke gjøre dem forlegne \emph{for at} træde frem. – Heraf udviklede sig nu en reen atomistisk Viden, som ikke traadte i noget dybere Forhold til Børnene og deres Existents, som \emph{ikke} blev \emph{\tlg{tilegnet}} paa en \emph{sjælelig} Maade og derved berøvedes al mulig Maalestok, og som en Følge deraf forfaldt de til at antage sig selv for store Naturforskere og Sprogmestere; naar først Enkeltheder skal afgjøre Sagen, er det naturligviis aldeles tilfældigt, hvor mange og hvor faa der høre til Mesterskabet. Deraf Coquetteriet, deraf de travle Marther, der glemme det ene Fornødne. Det er ikke om slig atomistisk Viden det gjælder, at hvad man i Ungdommen nemmer, man i Alderdommen ei glemmer. –

Med Hensyn til den \emph{Maade,} paa hvilken jeg troer det er nødvendigt ved al Underviisning og al Opdragelse at lade Barnet \emph{afføde i al Stilhed Livet i sig,} finder jeg nu, ved i disse Dage at læse Steffens »4 Nordmænd«, en god Bemærkning. Jeg har desværre kun den danske Oversættelse, den Steenske Udgave; Stedet findes 2den Deel, Side 250, 51, 52.

Jeg erindrer et Exempel paa, hvorledes i et saadant Liv Alt bliver affødet, Alt, hvad de læste om i Classikerne, afspeiledes; da de nu læste om Ostracisme, indførte de strax den i deres Leg – og saa videre

Og saa disse Børnebøger for: »artige, flittige, lydige, elskværdige, uskyldige, ufordærvede« Børn, – hvor man altsaa ved at forære dem et Exemplar siger dem, at \emph{de ere det,} da det jo ellers var en Misforstaaelse at give dem Bogen11.

1837.

»Ammestue-Historier« – i det Udtryk er der vist ligesaa meget taget Hensyn til Fortællingsmaaden som til Indholdet.

Desværre er det ikke uden Grund, at det altid er \emph{en Onkel,} der optræder som virksom, thi Forældrenes Virksomhed indskrænker sig som oftest til paa den maanedlige Regnskabsdag at optræde enten som Profos eller som Præmie-Uddeler for ædel Daad, – begge Dele med en nøiagtig og punktlig Bogholder-Samvittighed. Naar der derfor var Noget ved Onklerne, var der unægtelig Leilighed nok for \emph{dem} til Virksomhed.

Man skal ogsaa \emph{selv lære af Børn,} af deres vidunderlige Genialitet, som man ogsaa for en Deel maa lade raade og ikke som visse Selvkloge hovmesterere, og derhos erindre Christi Ord, da han var tolv Aar gammel: »Vide I ikke, at det stedse bør mig at være i min Faders Gjerning«? – (Noget Lignende forekommer det mig, at jeg har læst i en af Mynsters Prædikener.) – Man skal heller ikke strax være ved Haanden med det prosaiske Riis som Skolemesteren i »Alferne«, fordi der rører sig noget Dybere i Børnene – \emph{derved} undgaaer man blandt Andet (o guddommelige Nemesis!) at falde 1400 Alen ned under Jorden og blive til et – Mulæsel.

Børn interessere sig ikke saa meget for den \emph{græske} Mythologie, idetmindste ikke for det, der i en modnere Alder ansees for det Herligste (dog Herkules vel nok – NB. \emph{Underværkerne).}

Og Mange begynde saa tidligt dermed, medens Børnene endnu ere ganske smaa, saa der vist stundom falder et saadant Barn ind det, som Abraham af St. Clara fortæller om et lille Barn, at det, lige da det blev født, saae det Jammerlige i Verden saa grant, at det løb ind i Modersliv igjen. – Er det at \emph{styrke} Børn for Livet? Er det ikke at enervere hele deres Liv ved at berøve dem \emph{Enthusiasmens}perpetuum-mobile?

En \emph{Stat} bliver paa en Maade ufri derved, at den giver sig Loven.

Dette ligger i Tidsalderens Hastværk, som igrunden miskjender enhver Alder, fordi den troer, at den ene Alder kun er til for den anden.

Jævnfør\emph{Hamann } \emph{:} »Fünf Hirtenbriefe, das Schuldrama betreffend« i »Sämtl. W.«, 2te Theil, Seite 412 og følgende, – men hans altfor polemiske Ironie gaaer her ogsaa for vidt; saaledes vil han i Grunden have, at man skal lære Alt af Børnene i den strængeste Forstand, hvortil ogsaa hans Motto tyder: »Es ist ein Knabe hie, der hat fünf Gerstenbrod«, hvori der aabenbart ligger altfor Meget. Men det ligger nu igjen i hele hans Retning, thi det er vistnok ikke fordi han selv troer det, men for at ydmyge Verden; det er noget Andet med Socrates – hvilket Hamann ogsaa forlanger, – at spørge som Barn, men det er denne besynderlige Polemik, der gjør, at han heller vilde høre Viisdom af Bileams Æsel end af den viseste Mand, heller af en Pharisæer mod sin Villie end af en Apostel eller Engel (som han selv siger etsteds). Hans Polemik gaaer forvidt og involverer undertiden, saa synes det mig, noget Blasphemisk, Noget, hvorved han synes ligesom at ville »friste Gud«. – Forøvrigt er der naturligviis fortræffelige Ting ogsaa i disse 5 Breve.

Disse kloge Folk, der mene, at det ingen Kunst er at tale med Børn, – til dem vil jeg sige med Hamann: »Kindern zu antworten ist in der That ein Examen rigorosum; auch Kindern durch Fragen anzuholen und zu witzigen ist ein Meisterstück, weil eben Unwissenheit der große Sophist bleibt, der so viele Narren zu starken Geistern krönt – et addit cornua pauperi.« (Horats Od. III, 21).

Engang imellem erindre Vedkommende tilfældigviis en mere eventyrlig Historie fra deres Barndom, men den fortælle de for – saa snart de ere færdige og det Spørgsmaal opkommer: er der saadanne Havfruer til? – at svare: nei, det er noget Folk bilder sig ind. – Eller er da Eventyret noget \emph{saa} Intetsigende, at man strax maa tilintetgjøre Historien og dens Indtryk, at man strax vilde slaae den glimrende Sæbeboble itu for at vise, at den hele Herlighed dog ikke var andet end Sæbevand? Barndommen kræver Eventyret, og det er allerede Beviis nok for dets Gehalt. – Her fremkommer \emph{nu} det Spørgsmaal, hvorvidt man \emph{selv skal troe} paa disse Historier. Jeg troer, at naar vor Fortæller gjør det, idetmindste \emph{det} Spørgsmaal ikke vil fremkomme fra Børnene: er det nu sandt? thi Historien maa udøve en paa engang saa overvældende og saa beroligende Virkning, at det aldrig falder dem ind. Og det: ikke at fortælle Børn slige Phantasien beskjæftigende Eventyr og Sagn, giver netop Plads for en Ængstelse, som ikke modereret ved slige Fortællinger kommer desto stærkere igjen (jævnfør »die Verlobung«, Novelle von Tieck, Dresden 1823, Seite 63 nederst, 64 og 65. – Jævnfør ogsaa den jævne og simple Fortale til »Nordiske Kjæmpehistorier«, udgivet af Rafn, 2det Bind (Nota Bene Fortalen er naturligviis ikke af Rafn), især Slutningen af den, Side 9: »kan det nu vel hænde sig, at En og Anden, som hører disse Fortællinger, vil finde, at Sagaens vældige Bedrifter og Stor-Værker ikke passe med hans Færd og desaarsag forsmaae dem« – rigtigt: \emph{hinc} illæ \emph{lacrymæ} \emph{!})

Jeg seer nu \emph{en god Titel} i den sidste Messecatalog: Blumauer: »die kleinen Enkel auf dem Schooße der erzählenden Großmutter, ein Gegenstück der kleinen Enkel am Knie des erzählenden Großvaters«. (Messecatalog for Juli 1836 – Januar 1837; Side 27, øverst).
\end{quote}
\dcite{BB:37  (1837; Pap. IIA12; SKS 17, 123)}
% BB:37  (1837; Pap. IIA12; SKS 17, 123)

\begin{quote}
Da den valgte Titel paa Stykket synes at indeholde et utidigt Coquetterie skal det hedde.

Den altomfattende Debat af Alt imod Alt

eller

Jo galere jo bedre.

Af Een endnu levendes Papirer

mod hans Villie udgivet

af

S. Kierkegaard.

Striden mellem den gamle og den nye Sæbekielder

heroisk-patriotisk-cosmopolitisk-philantropisk-fatalistisk Drama

i flere Optrin

Samme Stykke er i Begyndelsen meget lystigt, i Fremgangen meget bedrøveligt, men dog i Udgangen meget glædeligt.

\emph{Personer.}

Willibald et ungt Menneske

Eccho hans Ven

Hr v. Springgaasen en Philosoph

Hre Hastværksen et foreløbigt Genie

Hr. Phrase, en Avanturier, Medlem af flere lærde Selskaber og Medarbeider i mangfoldige Journaler.

Hr. Ole Wadt virkelig Krigsassessor forhenværende Skriverlærer.

En Flue, der flere Aar har vidst at overvintre hos salig Hegel, og som har været saa heldig under Affattelsen af hans Værk: Phænomenologie des Geistes flere Gange at have siddet paa hans udødelige Næse.

En Ditto den Omskrevnes Sødskendebarn en Hegelianer.

Et Horn, Organ for Folke-Meningen, som man snart drikker patriotisk af, snart blæser patriotisk i, som Enhver, naar han seer sit Snit, blæser et Stykke paa.

En Bugtaler. En Kjæmper for Orthographien.

En Fodgjænger.

Polyteknikere.

Grosserere.

1ste Act.

Man seer Willibald efter en lille Monolog beslutte sig til at gaae hen i et Theeselskab, hvor han træffer Echo, som henriver hele Selskabet ved sin Aandrighed, sin Vittighed, sine af Willibald laante Indfald. Willibald vil gaae, Vertinden holder ham tilbage, endelig slipper han, iler hjem paa sit Værelse

Scene

\emph{Willibald.}

(Sidder paa sit Arbeids-Værelse i sin Sopha, Tobakspiben i Munden, omgivet af en stor Mængde opslagne Bøger og Lapper Papiir, paa hvilke han har noteret Eet og Andet.)

læser i Peter Schlemihl.

Hvad det dog er for et besynderligt Værk den Schlemihl... eller er det et Værk... er det ikke mig selv... ja rigtigt, en af Chamissos Phantasier er jeg... en Skygge selv der derfor ikke kan kaste nogen Skygge... jeg synker sammen som Auroras Mand og kan snart gjøre mig Haab om som en Naturmærkværdighed under en Glasklokke at pryde et Naturalie-Cabinet.... O da skal jeg engang under en høi- og vellærd Professors Forelæsninger sprenge Glasset og mit Livs Hyperbel overgaae alle hans Decimal- og Allgebra -Beregninger.... da skal mit sidste ildsprudende gigantiske Suk idetmindste indjage Spidsborgerne saa megen Skræk, at jeg skal blive fri for al condolence.... (idet han puster en Sky Tobak ud) og disse Taagemasser det er altsaa det Rige, jeg hører hjemme i. o! see hvor de fortætte og ligesom blive til Skikkelser... Nu vel da skal jeg selv være en Skygge, saa vil jeg idetmindste ogsaa digte en ny, jeg vil skabe en.... (med stor Pathos) der vorde et Menneske........ (idet samme antager Skyen Skikkelse af Echo) hvad seer jeg er det ikke min Plageaand – mit andet jeg..... men er jeg da jeg... ligemeget (griber sin Sabel) herfrem Du alle mine Ords samvittighedsfuldeste Bogholder.... mød og med Dig den hele Skare af Eftersnakkere – gid alle Eders Hoveder sad paa een Hals (hugger til; men der er Intet)... der er Intet... hvad om jeg da gjorde det samme Experiment med mig selv.... (vender Sablen mod sig selv, idet samme kommer han meget stærkt til at hoste) skulde jeg muligen have faaet den Stump Pennefjær i Halsen jeg før legede med..... jeg føler allerede, hvorledes den glider ned i Luftrøret (hoster tager Lyset og seer sig i Speilet)... jeg er usædvanlig bleeg (hoster) (man hører udenfor Een rømme sig – det banker)... det er min Ven naar jeg hoster rømmer han sig.... jeg vil vedde paa at han ogsaa har sunket en Pennefjær... herein! Echo træder ind med en ærbødig Compliment

\emph{W.} Wo komst du her geritten? (omfavner meget venskabeligt)

\emph{E.} For at vedblive som Du har begyndt: Wir satteln nur um Mitternacht.

\emph{W.} Bravo... men desto værre er saa din Tid nu forbi; thi Klokken er lige paa Slaget 1.

\emph{E.} Oh! i saa fortræffeligt Selskab..... (bliver Sablen vaer som Willibald endnu holder i Haanden) hvad seer jeg, hvortil dette Vaaben.

\emph{W.} Jeg var just ifærd med paa Opfordring af en medlidende Djævel at slaae en Græshoppe ihjel, for at den dog kunde falde for dens eneste Vens Haand (lægger Sablen).

\emph{E.} Jeg forstaaer dig ikke

\emph{W.} Godt sagt.

\emph{E.} Forklar Dig, Du er i en ophidset Tilstand... hvorfor forlod Du Selskabet saa tidligt... alle stormede ind paa mig og spurgte mig om Grunden – hvad skulde jeg svaret.

\emph{W.}.. det veed jeg ikke... hvad de ønskede.

\emph{E:} O let Dit Hjerte for mig... mig, Du har gaaet i Skole med; mig som har været indviet i saa mange af dine Planer.

\emph{W.} (afbryder ham) ja Dig, Du som har faaet Prygl tillige med mig af Rector.. Dig min bedste min eneste Ven.

\emph{E.} (vedbliver) mig, til hvem du ved en saa dyb Sympathie er knyttet, jeg føler det Samme... Du kjeder Dig i Selskaber det gjør jeg ogsaa.

\emph{W.} Oh nei jeg synes de ere moersomme.

\emph{E.} det forstaaer sig mange af dem.... Du søger Eensomhed

\emph{W.} det gjør Du ogsaa

\emph{E.} O hvor det er sørgeligt at blive misforstaaet ikke at turde aabne sit hele Hjerte ja Misforstaaelse, ja Misforstaaelse

\emph{W.} (med et ironisk Smiil) ja der gives dog Tilfælde, hvor Vedkommende kan være meget tjent dermed.

\emph{E.} Det nægter jeg ikke – til en vis Grad.

......

nonnulla desunt

\emph{W.} Jeg staaer netop nu og undersøger Lungen og Hjertets Beskaffenhed, mit Aandedrag slaaer Triller, jeg føler nøiagtigt, rledes en Stump Pennefjær jeg har sjunket stiger langsomt ned for at gjøre det af med mig

\emph{E.} Du skulde spise en Bid Brød eller drikke Noget, jeg fik forleden ogsaa en Krumme i den gale Hals

\emph{W.} Jeg spytter Blod – iil efter en Læge

(E. iler i største Hast bort)

\emph{W.} gaae med Gud – (idet E. lukker Døren) gaae Fanden i Vold. endelig blev jeg da af med ham, o hvorfor er jeg dog bleven et selskabeligt Dyr, et Menneske, hvorfor ikke en Ugle eller en Rørdromme, saa var jeg dog frie for den værste Plage – for Venner.

han iler barhovedet ud af Døren stærkere og stærkere, og man hører Intet uden langt borte en svag Musik af en Tromme, Bækner og en forstemt Clarinet der foredrager det bekjendte Thema af Murmesteren: Nei ei forsage Nei ei forsage altid Venner er os nær herover bliver W. end mere fortvivlet og styrter i Raserie afsted.

\emph{E.} kommer igjen med Lægen.

Hans Forundring over ikke at finde W. afløses snart af en vis høitidelig Stemning fremkaldt ved Betragtningen af det nu af sit belivende Princip forladte Værelse.

\emph{E.} Jeg omgaaes ham meget, et høist besynderligt Menneske, en Original, fuld af de kourieuseste Indfald, som flagre om, da han ligesom hiin Gudinde skriver skriver paa Blade, som han lader henveires af Vinden, saaledes har jeg engang tilladt mig at sige i en Novelle, om hvilken en Recenscent har været saa god at sige, at jeg vist vilde forundre mig over paa mine gamle Dage, at jeg kunde have saa stor Genialitet i min Ungdom.

\emph{Doctor. } Oh hvis De til den Tid ikke skulde have mistet aldeles deres Hukommelse, saa vil jo den nok kunde give dem den behørige Oplysning.

\emph{E.} Det er et Menneske, som jeg holder Øie med ligesom Politiet med mistænkelige Personer, som jeg har udtrykt mig i en Novelle og i mit dobbelte Bogholderie over Indfald er en egen Bog indrømmet ham. Naar jeg ikke vilde skade mig selv fornemligt mit Honorar ligesaa meget som jeg kunde gavne det Almene skulde jeg udgive et eget Værk om Indretningen af et sligt Bogholderie. Dem skal jeg imidlertid meddele det.

\emph{Doctor.} Betænk dog, at det jo er det første Skridt

\emph{E.} Frygt kun ikke jeg er overbeviist om at deres Praxis vil forhindre dem i at gjøre Brug deraf. Det har forøvrigt kostet mig megen Tid inden jeg fandt det Rette saavel med Hensyn til Beregningen af Rotationen i denne litteraire Vexeldrift som den chemiske Undersøgelse af den Mergel jeg bruger og dens Forhold til min egen Avlings Beskaffenhed

\emph{Doctor.} Hvor gjerne jeg end vilde høre derpaa, saa vil de dog i den samme Grund, som betingede deres Velvillie til at meddele mig den see en tilstrækkelig Grund og derfor.... maaskee vil et senere mere favorisabelt Øieblik atter bringe os sammen, og da vil de maaskee have den Godhed,.... med mindre de i denne Begunstigelse af Tiden skulde see en Grund til Taushed. altsaa farvel.

\emph{E.} Hvo veed, om det dog ikke var en lykkelig Styrelse, der standsede mig første Gang jeg var ifærd med at forraade min Hemmelighed, ihvorvel jeg dog ikke skal nægte at jeg nok vilde ønske, at Een eller Anden kunde tale til Folk om den store Hemmelighed, jeg er i Besiddelse af.

Vi vende imidlertid tilbage til Willibald. Fortvivlet som vi have seet ham over det venskabelige sidste Afskedsquad styrtede han i Haab om dog engang at undgaae Venner, videre. Dette skulde imidlertid kun altfor snart hende ham. I sin Forvirring var han løbet paa en Mand der som han nu opdagede gik fordybet i en Samtale med tvende Andre. Efter at have gjort sin Undskyldning fik han til Svar, at han ikke behøvede at gjøre nogen Undskyldning; thi han var en stor Synder det vidste han vel og det skulde kun glæde ham, om han maatte vederfares den Lykke at lide Noget for Christi Skyld. – Neppe har Willibald forladt dem, før han i største Iil uden noget Maal skynder sig videre, idet han declamerer høit for sig: Før var min Ære som et speilklart Skjold af slebent Staal,.... nu staaer en Plet af blodig Rust derpaa..... Denne Talen høit er nok til at gjøre en Politie-Agent, der var udsendt for at anholde nogle Opvakte, opmærksom paa ham, og denne Tale om »blodig Rust« mere end tilstrækkelig til at gjøre ham i høi Grad mistænkt. Imidlertid var W. kommen ham af Syne, og da det trods Politiets ivrigste Bestræbelser ikke er lykkedes at paagribe ham, saa indseer Læserne Nødvendigheden af at han ikke mere er at finde paa Jorden. –

2 d Act.

En phantastisk Egn. Prytaneum – hvori de nævnte Personer underholdes paa offentlig Bekostning.

Alt er arrangeret trekantet – der spilles 3 Kort og saa videre

1st Scene.

Ole Wadt.          Holla Hastværksen.

O. Wadt. Som jeg siger jeg misbilliger ingenlunde deres hele Retning, jeg er tvertimod meget villig til at anerkjende deres udødelige Fortjenester af fædrelandske Anliggender; men hvad Stilen, hvad Udtrykket angaaer, da er altid Noget der støder – Deres Pen er ikke blød nok om jeg saa maa sige.

\emph{H.} De mener altsaa, at det kommer af at jeg bruger Staal-Penne.

\emph{W.} Meget rigtigt! De har der gjort en langt dybere Bemærkning, end De selv maaskee troede. Der er Intet der i den Grad fordærver Haanden og Hjertet som Staalpenne. Hvorledes vil vel et Kjærlighedsbrev blive, der skal skrives med Staal.

\emph{H.} Det har vistnok en dyb praktisk Betydning, en symbolsk Charakter, at Staalet saaledes er bleven overført fra Landser og Spyd paa Penne.

\emph{W.} Jeg føler destoværre allerede, at Tidens Sværd er gaaet gjennem mit Hjerte, at Samtiden har tabt det bløde bøielige, elastiske, Yndige, Fortiden glædede sig ved.

\emph{H.} Skal jeg nu plages med den gamle Historie om Fortiden. Det var en idyllisk Uskyldighedsstand, men nu ere vi Mænd. Nu maae vi med Alvor gribe ind i Tingene, med Staalhandsker væbne os

\emph{W.} De mener med Staalpenne.

\emph{H.} O gaae Fanden i Vold med deres Gaasefjær, De har ingen høihjertede Følelser, der er en Sjæl i dem som i en Gaasefjær.

(gaaer bort i Vrede)

2d Scene.

De Forrige.          Phrase.          v. Springgaasen.

(Hr. v. Springgaasen en lille uanseelig Mand, hvis ene Been var et godt qvarteer mindre end det andet, og for at anskueliggjøre hans philosophiske Ideer, pleiede han, efter først at have hævet sig paa det længste Been, at forlade dette ilusoriske Standpunkt som han pleiede, at udtrykke sig, for at vinde den dybere Realitæt).

Phrase standser Hastværksen.

\emph{P.} Lad mig et Øieblik standsere deres bevingede og hastende Fjed Hr Hastværksen. Et Anliggende af stor Vigtighed har i længere Tid fyldt min Sjæl, et Anliggende, ved hvis Udførelse jeg haaber at turde regne paa deres gunstige Medvirkning. Det store Forraad af Kundskaber vi hver for sig ere i Besiddelse af bør vi ingenlunde stræbe at isolere, men ved Haand i Haand at arbeide gjøre Realisation af grandiøse Formaal, dog ikke det alene vi bør ogsaa bestræbe os for at gjøre Videnskabens store Resultater tilgjængelige for Folket, vor Tids Udvikling bør vinde i Extensitæt, hvad den taber i Intensitæt.

\emph{v. S.} Ja det er godt nok med det popoulaire; men min Tvivl er ingenlunde popoulair; det er ikke en Tvivlen om Dit og Dat, om dette eller hiint, nei det er en uendelig Tvivl, ja jeg ængstedes undertiden af en sand videnskabelig Tvivl om jeg dog ogsaa har tvivlet nok; thi Tvivl er det Specifique i den nyere Philosophie, der in parenthesi sagt, begyndte med Cartesius, der sagde de omnibus disputandum est, hvorved han aldeles tilintetgjorde den Sætning, der tilforn gjaldt for en Grundsætning: de gustibus non est disputandum. See Slige store videnskabelige Problemer vilde umulig kunne meddeles Menig-Mand.

\emph{P.} Min Mening er jo heller ikke, at man egentlig skulde skrive for Bønder; Nei for den dannede Mellemstand, for Grosserere, Polyteknikere for Capitalister, og naar man lempede Stilen lidt mere...

\emph{Wadt } om Forladelse, det er netop hvad jeg har udviklet for Hr Hastværksen, ja Stilen, Skrivemaaden, det er det Det kommer an paa. Man tage Noget af det Skarpe, Kantede, det altfor Spidse bort fra den moderne Philosophies Mangesidethed, man afrunde Formerne lidt mere, og der er ingen Tvivl om, at det vil lykkes, at vi vil komme til at skrive den ny Udvikling fra den Dato.

\emph{Hastv:} Philosophie mig her og Philosophie mig der. Det er ikke Philosophien det kommer an paa. Det er de praktiske Spørgsmaal, Livs-Spørgsmaalene – kort sagt Livet

\emph{v. Sp.} ja hvad er Livet

\emph{P.} Livet er et af sig selv udgaaende til sig selv tilbagevendende.

\emph{v. Sp.} etcetera etc. det var kun for at vise Hr Hastværksen at det er saare vanskeligt at blive popoulair at jeg gjorde ham den Indvending. Det gjælder ogsaa her det den nyere Philosophies dybsindige Fordring δοσ μοι που στω. Men hvor skal jeg faae Fodfæste i den vulgaire Resonnaiments-Sphære. δοσ μοι που στω.

\emph{H.} Ja det maa vistnok altid falde dem vanskeligt at faae Fodfæste, og det er omtrent Tilfældet med alle Philosopher, der ere saa slet til Fods som de Hr Springgaasen.

\emph{v. S.} Det er et lavt Angreb.

\emph{Wadt } Ja Formen mangler ham. Det Stødende træder saa stærk frem i enhver Yttring af ham. –

Nonnulla desunt.

3die Scene.

Willibald.

(Noget forvirret i sit Udseende, han seer sig forundret omkring, bliver Indskriftet »Prytaneum« vaer, kaster sig til Jorden af Glæde og kysser den. udbryder i Glæde over at være befriet fra hele det Liv, i hvilket han hidtil havde trællet, og over at see sig hensat i en Egn, hvor Viisdom nødvendigviis maa boe. Han faaer Øie paa Springgaasen gaaer ham med dyb Ærbødighed imøde)

\emph{Willibald.} (nærmer sig til Spr:) Uden egentlig at vide, hvor jeg er kommen hen, er det mig dog altid en Trøst at vi, at jeg har forladt alle mine Plagers Hjem. Den udvortes Omgivelse, det hele Totalindtryk har vakt hos mig en glad Forestilling om, en lykkelig Ahnelse, om at her maa Viisdommen være at finde, at jeg her maa heldbredes for den rædsomme Relativitæt, jeg hidtil har ligget under for.

\emph{v. Sp.} Kjære! Hvad der fattes dem indseer jeg tilfulde. Det er det Faustiske, det er det, som den nyere Philosophie der in parenthesi sagt begyndte med Cartesius, i høi Grad har lidt af. (under denne Replique ere de øvrige Medlemmer af Prytaneum i Samtale traadte ind paa Scenen) (idet Sp. henvender sig til de øvrige Medlemmer)... da jeg just er ifærd med at give en ganske kort Fremstilling af den nyere Philosophie siden Cartesius, saa kan jeg maaskee tjene de øvrige Tilstædeværende ved at tale offentligt om det, saa Alle kunne høre det.

\emph{Phrase } Er det den samme Udvikling, som jeg før allerede har \tlg{tilegnet} mig en stor Deel af; O da vilde de idetmindste gjøre mig en stor Tjeneste ved en saadan Gjentagelse, at sætte mig istand til fra Speculations Himmelfartsbjerg at overskue de store Ideers historiske Incarnationer.

\emph{v. Sp.} O hvilken Glæde, at have en slig Discipel, hvilken Utrættelighed – Ja snart vil Du om end ikke virke saa vidt som Din Mester dog med Hæder kunne beklæde en Docents Post i Nordlandene og adsprede det hyperboræiske Mørke.

Mine Herrer

Det var altsaa med den i sit Liv forfulgte (om han egentlig blev forfulgt veed jeg ikke, men da det er flere Aarhundreder siden han levede, saa er han jo hjemfalden til Mythologien og i dens Lys seet er han nødvendigviis bleven forfulgt) efter sin Død forglemte; men nu for evig udødeliggjorte Cartesius, fra hvem hele den nyere Philosophie daterer sig. Jeg veed vel, at der i vort Selskab desto værre ere indsnegne sig nogle vankundige Mennesker, der, i jordisk Travlhed for de meest underordnede Interesser, antage at Verden begyndte ifjor, og ikke med os, mine Herrer, høitideligholde den store nu verdenshistorisk indtraadte Anskuelsens Bededag, paa hvilken Mennesket det vil sige Philosophen efter fra den umiddelbar paradisiske Uskyldighedstilstand gjennem Livets Dialektik at være vendte tilbage til det Umiddelbare, netop af den Grund ikke bør arbeide men i intuitiv Nyden gjenopleve det allerede oplevede Verdens-Drama. Dog nok Intet skal forstyrre mig i min contemplative Rolighed. Jeg seer den modsatte Anskuelse som et forsvindende Moment, om ikke af anden Grund, saa dog fordi mit jo ellers ikke vilde være det faste. Det var da Cartesius, der udtalte de mærkværdige, evig uforglemmelige Ord: cogito ergo sum og de omnibus disputandum est – Ord, der egentlig i enhver velordnet speculativ Stat burde læres ved Confirmations-Underviisningen, og som idetmindste ingen theologiske Candidat burde være uvidende om, da ingen speculativ Sjælesørger kan uden dem gjøre sig Haab om at røgte sit vanskelige Kald med Held.

Ja store Tanke – saavidt vil det engang komme, at man vil ansee disse Ord, jeg gjentager dem, cogito ergo sum og de omnibus disp: est for Statens videnskabelige Parole, for et Palladium, der vil fjerne Alt Kjætterie, for Ord, der lig Ordet Adam, vil erindre os om vort intellectuelle Livs Skabelse.

Nonnulla desunt.

\emph{Præsidenten} Da det strider mod vore Love som ogsaa mod den dramatiske Velanstændighed, vi hidtil have iagttaget i vort Selskab at holde slige lange Monologer, maa jeg paa Embeds Vegne afbryde dem.

\emph{v. S:} det skulde gjøre mig ondt om min Veltalenhed paa nogen Maade skulde have henrevet mig til at overskride den i vort Selskab vedtagne Convenients. Og hvis det virkelig forholder sig saaledes, da maa Grunden ene være at søge i det friere Foredrag jeg har overladt mig til med Hensyn til vor Katechumen, hvem jeg ønskede at sætte ind paa det rette Standpunkt ogsaa med Hensyn til vort Selskabs Beskaffenhed. Jeg tør med den største Rolighed kalde flere af de tilstædeværende Herrer til Vidne paa, at det korte Foredrag, jeg ellers pleier at holde over den nyere Philosophie siden Cartesius eller rettere sagt over den nyere Philosophie, paa ingen Maade vilde være for langvarigt endog i det bedst indrettede Drama. Ja Hr Præsident jeg vil vedde paa, at det ikke varer over 1½ Minut, da jeg netop har indrettet det med Hensyn til vort Selskabs Tarv.

\emph{Præsident} Jeg nødsages atter til at kalde dem til Orden og paabyde Taushed.

\emph{v. S.} det var Cartesius, der sagde cogito ergo sum og de omnibus dubitandum.

\emph{Præsident} silentium.

\emph{v. S.} Spinoza gjennemførte nu dette Standpunkt reent objektivt, saaledes at al Tilværelse blev det Absolutes Undulationer.

\emph{Præsident} Pedeller træde frem.

\emph{v. S.} denne Objektivitæt bortdestilleredes imidlertid ganske i den critiske Udvikling, og medens Kant kun til en vis Grad gjennemførte denne Skepsis, var det Fichte forbeholdt at skue denne Medusa under Øie i Criticismen og Abstraktionens Nat.

\emph{Præsident} Lægger Haand paa ham og fører ham bort.

\emph{v. S.} da jeg seer man vil bruge Magt, kan jeg ikke faae udført det om Sleiermacher, men Hegel var det, der speculativt concentrerede de forhenværende Systemer og med ham har derfor Erkjendelsen naaet sit egentlige dogmatiske Høidepunkt.

(Pedellerne gjøre Mine til at ville gribe ham)

Nu er jeg færdig og med Hegel er Verdenshistorien forbi, bortfører mig kun; thi nu er der intet andet end Mythologie tilbage og jeg bliver selv en mythologisk Person.

\emph{Phrase } Det er et aldeles eensidigt Standpunkt (rømmer sig) Mine Herrer jeg er kommen ud over Hegel, hvorhen kan jeg endnu ikke sige saa ganske nøie; men jeg er kommen ud over ham.

\emph{v. S.} Hvad maa jeg høre Slange! Judas! slipper mig eller skal da altid den evige Idee ligge under for Masse.

\emph{Præsid:} bortfører ham til Carceret.

(de føre ham bort.)

\emph{Phrase } Jeg gjentager det mine Herrer, jeg er kommen ud over Hegel. Det var nemlig den nyere Philosophie der begyndte med Cartesius, som sagde: cogito ergo sum og de omnibus dubitandum est.

\emph{Præsidenten} (afbryder ham) Da Hr v. Springgaasen har givet Anledning til nogle ubehagelige Optrin, og derved foraarsaget en Spænding i Gemytterne, nødsages jeg til at hæve Forsamlingen.. Enhver gaae til sit – Det, der imidlertid smerter mig meest er, at vi ikke strax kunne faae Leilighed til at høre Hr v. S.s Mening om et meget difficilt Spørgsmaal, der er opkommen i vort Selskab. Det vil imidlertid blive foretaget imorgen paa Generalforsamlingen, cui vos ut frequentes adsitis etiam atque etiam rogamus.      missa est ecclesia. –

Imidlertid vidste man i Prytaneum slet ikke hvad man skulde gjøre med Willibald, der ikkun lidet havde fundet sig opbygget og tilfredsstillet ved v. Sp. philosophiske Forelæsninger. Man havde endelig besluttet at henvise ham til den videnskabelige Anstalt, Prytaneet havde anlagt – den \emph{verdenshistoriske Høiskole.} Denne var dog endnu ikke færdig, og kun Atriet lod sig bruge, men dette var saa stort, at der paa eengang docerede 4 Professorer, uden at forstyrre hinanden, ja det var saa stort, at ikke engang Tilhørerne kunde høre, hvad Docenterne udviklede, ihvorvel disse idelig aftørrede Sveden af deres ved Anstrængelser blødgjorte Pander. 2 af disse 4 Professorer sagde ordret det Samme, og naar de vare færdige vendte de sig om med en Mine som om Ingen i Verden kunde sige noget Saadant.

Her var nu Willibald ved personlig Omgang efterhaanden bleven vundet for de Anskuelser, der gik og gjaldt i Prytanæet, og havde forsaavidt allerede igjen fortrudt et Indfald han i sin Forbittrelse paa v. S. havde indsendt til Præsident, et Spørgsmaal nemlig, hvoraf det kom, at Solen i Prytaneet slet ikke forandrede sin og som en Følge deraf Belysningen bestandig blev den samme, et Spørgsmaal, som nu i høi Grad foruroligede Prytanæet, og gav Anledning til den Generalforsamling, der nu skulde afholdes.

Generalforsamling.

Præsidenten v. S. Phrase. Ole Wadt. Hola Hastv: Polyteknikere, Philologer etcetera

Præsidenten fremsætter Sagen, som allerede foreløbig var bekjendt og haabede derfor at Debatterne kunde blive jævnere

\emph{Hola Hastv.} jeg udbeder mig Ordet. Det Phænomen paa hvis Forklaring Opmærksomheden nu skal henføres er af stor Vigtighed, skjøndt det for enhver, der med vaagent Øie har forfulgt den nyere Tids Kæmpeskridt til at emancipere Videnskaberne og macadamisere dem, er let at forklare det er Morgenbelysning, det er det høitidelige Morgengry, det er Solens Kamp med Mørkets sidste Anstrængelser det er som en Digter har sagt Prytanæets Mai og Prytanæets Morgen. Naar disse først ere overvundne da vil vi faae Øinene op for Resultaterne af vor Tids store Fødselsveer, da ville vi bebude et Gylden-Aar, den rette Nyaars-Tid, da al den gamle Suurdeig og al det gamle Skolefuxerie, al Jesuitisme og Papisme er udfeiet.

\emph{v. S.} Jeg miskjender ingenlunde den ærede Talers Bestræbelser for at gjøre Udsigterne i Fremtiden saa lyse som mulige; men jeg vil sige at han har en Tilbøilighed til at haste alt for meget frem til altfor meget at gaae gerade aus; jeg derimod, jeg føler den dialektiske Puls i mig, jeg bevæger mig i den sande speculative Siksak. Saaledes som man seer at den gemene Fisk, den gemene Fugl, det gemene Dyr i sin Locoumotion altid gaaer slet og ret lige ud efter Næsen, hvorimod den ædle Rovfugl, den værdige Rovfisk, det stolte Rovdyr sees i Spring at bemægte sig sit Bytte, saaledes gaaer det med den ægte speculative Bevægelse, naar man vil gjenemføre dette ægte speculative Billedes dybe Betydning. Udsigter kan der derfor i Opfattelsen af Verdensforholdene ikke blive Tale om idetmindste ikke gerade aus som Hr Hastværksen s fortryllende Udskuelser, snarere maatte man kalde det der bliver for den speculative Betragtning Indsigter, ligesom en Slange, der bider sig selv i Halen (dette uudtømmelige Billede paa Speculationen) jo ikke kan siges at see ud men at see ind, saaledes at den tilsidst seer bagfra ind i sit eget Øie, seer saa at sige ud og ind i sit eget Øie og af sit eget Øie paa eengang.

\emph{Hastv:} Længe nok har jeg hørt paa Hr. v. S: Vidtløftigheder, der dog ikke ere andet end Hegels berømte Tanke-perpetuum-mobile. Men vi der arbeide for Livet vi der ere dimitterede fra Skolen for Livet kunne i Sandhed ikke lade os nøie med at kjøre den ene Tour efter den anden om Hesten uden at komme videre paa Livets Vei, og vi kunne ikke lade os spise af med Hr v. S.s uudtømmelige eller som det snarere synes uopslidelige Slange-Billeder, og med alle de Indfald, der kun kan opkomme i en Mands Hjerne der er født i Steenbukkens Tegn, som Hr v. S. formodentlig er, ja vi springe vistnok med Rette over alle hans Spring.

\emph{v. S.} Uden at lade mig forstyrre af Hr Hastværksens Misforstand, hvilken berøver mig en Deel af min Glæde over at see, at han dog ellers har opfattet en Idee rigtig Overgangs-Kategorien, hvorved man gaaer ud over sin Forgænger, skal jeg gaae over til at løse det forelagde Problem. Phænomenet lader sig nemlig forklare deraf at det er Aftenbelysning. Philosophien er nemlig Livets Aften, og med Hegel, der speculativt concentrerede de foregaaende rationelle Systemer er den verdenshistorisk indtraadt.

\emph{Phrase } Jeg er kommen ud over Hegel.

\emph{Hastværks.} Jeg fordrer Ballotation

\emph{Præsident} skal det være med Kugler eller aabenlyst

\emph{v. S.} jeg forlanger Ordet. Det forekommer mig at være en Urimelighed at ville afgjøre et sligt Spørgsmaal ved Ballotation, og det er min uforgribelige Mening, at den Endelighed i Discussion man opnaaer ved Ballotation egentlig ikke er nogen Endelighed, men snarere den slette Uendelighed.

(Flere i Munden paa hverandre)

A. Det er en Sag af yderste Vigtighed

B. Det er et Livs-Spørgsmaal

C. det er et Princip-Spørgsmaal

D. det er et Livsprincip.

\emph{En Polytekniker} Staten er et galvanisk Apparat

\emph{v.S.} Staten er en Organisme.

\emph{Phrase } Hr Præsident jeg opfordrer dem til at erklære, om det er mig, som taler, ellers forlanger jeg Ordet.

\emph{Hastv:} Jeg forlanger der skal balloteres paa, om der skal balloteres. Jeg kæmper for Friheden, vi ville ikke længer lade os trykke af disse tyranniske Philosopher.

\emph{Ole Wadt.} Venner hvor sørgeligt vilde det ikke være om et sligt Spørgsmaal skulde tilintetgjøre den gode Forstaaelse i vort ældgamle Samfund.

\emph{Præsident} Vanskeligheden bliver om vore Love tillade Ballotation i et sligt Tilfælde

\emph{En Philolog} Jeg andrager paa, at der maa nedsættes en Comite af Folk, der ere kyndige i Alderdomsvidenskaben for ved sund Critik at udfinde Lovens Mening.

\emph{Præsidenten } Den ærede Talers Forslag synes af den Grund overflødigt, at Lovene kun ere 1 Aar gamle.

Midt under denne støiende Forhandling træder Willibald op, der imidlertid i Prytanæets verdenshistoriske Høiskole var vundet for Selskabets Ideer, han forsikkrer, at det slet ikke var Tilfælde, at Solen den physiske nemlig virkelig forandrede sin Stilling, men at han kun ved det Hele havde villet antyde den poetiske, philosophiske, kosmopolitiske Evighed, der allerede var i aandelig Forstand indtraadt i Prytanæet.

Dette beroliger Gemytterne, og Forsamlingen hæves.

3die Akt.

Willibald.

(Han spadserer i en phantastisk Egn i Nærheden af Prytanæet)

(med stigende Pathos) Du Uendelige, Du.. hvad skal jeg kalde Dig, hvorledes skal jeg nævne Dig, Du uendelige Nævner for alle menneskelige Tællere, Du absolute Aand, som nu ikke længer er en Hemmelighed for mig; men hvis forborgne Dybder jeg nu kan lodde; ja her er godt at være her er Viisdommens Hjemstavn her, her hvor jeg fandt min udødelige Lærer v. Springgaasen, ja nu er Lyset gaaet op for mig over Alt. (idet han siger dette flyver en Flue surrende forbi ham og foredrager nogle hegelske Sætninger; og Hornet lader sig høre med nogle politiske Theoremer) ogsaa dette, kun det manglede, visseligen nu er Verdenshistorien forbi; thi nu kan Naturen fastholde Begrebet.

Ole Wadt træder op.

\emph{O. W. } Dyrebare unge Ven De bør gjøre Alt for at tilveiebringe Eenhed i Selskabet; thi den sidste Begivenhed har efterladt en Spænding i Gemytterne, som ængster mig meget.

\emph{W.} Jeg iler Alt vil jeg gjøre, Alt, hvad der staaer i min Magt.

(De gaaer bort Arm i Arm)

Willibald; v. Sp:; Hol. Hastv.; Ole Wadt; Phrase etcetera

\emph{W.} Skjøndt jeg troer, at det egentlige Stridspunkt er hævet ved at jeg har taget mit Andragende tilbage, saa mener jeg dog at vi paa en eller anden Maade bør antyde at Freden igjen er tilveiebragt, jeg mener, at vi bør begynde en ny Tidsregning og til den Ende ogsaa give vort Selskab et nyt Navn, under hvilket det dog iøvrigt forbliver det samme. Jeg foreslaaer derfor at vi for Fremtiden kalde det: det Ny- og det Gamle-Prytaneum, skrevet Nota Bene med Hyphen.

\emph{Hastv.} Jeg for min Deel interesserer mig slet ikke for slige logiske Bestemmelser, kun forsaavidt man dog derved antyder en ny Udvikling bifalder jeg det.

\emph{v. S.} Forslaget selv forekommer mig yderst speculativt, man kunde vel troe, at man ved at udslette den gamle Indskrift »Prytaneum« og derpaa paa ny skrive »Prytaneum« udrettede det Samme; men det vilde dog kun være at vende tilbage til det Umiddelbare, hvor de dialektiske Modsætninger endnu ikke havde udviklet sig og speculativ gjennemtrængt hverandre; men dette er ikke rigtigt. Hertil kommer, at den hele Begivenhed kaster et mærkværdigt Lys over en Mythe, som vist er alle de Tilstædeværende bekjendt Striden nemlig mellem den gamle og den ny Sæbekjelder, og heraf indlyser tillige den speculative Betydning af Mythe overhovedet, at den indeholder en Anticipation af Historien, et Tilløb saa at sige til at blive Historie.

\emph{O. Wadt } Hvormeget jeg end billiger det af Hr v. S. Udtalte, mener jeg dog derfor, at vi tillige paa en mere udvortes Maade bør opbevare Erindringen om denne uforglemmelige Dag ved at opreise en Mindestøtte, et Monument. At Belønninger saavel Pengepræmier som offentlig Roes gavner meget har jeg i sin Tid erfaret som Skriverlærer i Efterslægten.

Derpaa bliver der opreist et Monument, ved hvilken Leilighed der udbringes flere enthusiastiske Skaaler især af Willibald.

Ende. –

\emph{[i margen:]} dediceret til de 7 Gale i Europa, som ingen By har villet vedkjende sig. –

\emph{[i margen:]} Med et Titelkobber som forestiller Luther, der sidder i et Hæsseltræe og skjærer Kjeppe for Folk, der gjøre unyttige Spørgsmaal; man seer nogle af dem ligge paa Jorden, andre vil man finde rundt om i Bogen; Uerfarne ville maaskee tage Feil deraf og ansee dem for Tankestreger.

\emph{[i margen:]} \tlg{Tilegnet} de fire gale Brødre i Claudius Verset aftrykkes.

\emph{[i margen:]} En Slutnings-Vignet skulde forestille Zacchæus i Morbærtræet. –

\emph{[i margen:]} Holla (hans Fornavn)

\emph{[i margen:]} underholdes paa offentlig Bekostning i Prytaneum.

\emph{[i margen:]} Sallust. Jugurtha Capitel IV: profecto existimabunt me magis merito quam ignavia judicium animi mutavisse majusque commodum ex otio meo quam ex aliorum negotiis reipublicæ venturum. –

\emph{[i margen:]} For at imidlertid denne Piece skal være nyttig for Noget, følge et kort Indbegreb af Conversationsthemata belieblich udsatte zum Gebrauche für Jedermann, og en Fortegnelse paa de Uquemsord man kan bruge uden at gjøres ansvarlig efter Trykkefrihedsforordning af 1799

\emph{[i margen:]} Efter Skrift da jeg imidlertid seer, at en Floskelskriver ved Kjøbenhavnspostens Contoir stedse kommer mig i Forkjøbet saa maa jeg med Smerte saa nødes jeg til at udelade det for ikke at blive unyttig

\emph{[i margen:]} Man kunde ogsaa lade Willibald synge: Sjelden Penge, Prygl desfleer etcetera og Echo synge: O fatal, hvilken Qual, maa man ei blive gal. (denne sidste Arie har nemlig i Musik og Idee en Lighed med den første, forsaavidt en Modehandlerindes Fortvivlelse af Auber kan sammenlignes med Leporellos af Mozart.).

\emph{[i margen:]} \emph{W.} alle Romaner lyve

\emph{[i margen:]} \emph{E.} Det har jeg ogsaa lagt Mærke til. Saaledes læste jeg forleden en Roman, der begyndte med de Ord: det gulttavlede Sted paa Hjørnet af Kronprindsens Gade og st. Kjøbmagergade er vist ofte faldet de Forbigaaende paa.« Forunderligt, tænkte jeg, det har du aldrig lagt Mærke til. Øieblikkelig iler jeg derhen for ved Beskuelsen deraf end mere at sætte mig ind i Romanen og om muligt ved at titte ind af Vinduet lære den beskrevne Familie at kjende. Men trods al Anstrængelse var det ikke muligt at finde et sligt Sted.

\emph{[i margen:]} \emph{W.} Besynderligt men har Du ogsaa lagt Mærke til hvad Aar Romanen var fra eller i ethvert Tilfælde, i hvilket Aar Romanen foregaaer, maaskee Stedet er bleven malet om.

\emph{[i margen:]} \emph{E.} Skulde det være muligt, ja det skal jeg dog rigtig erkyndige mig om.

\emph{[i margen:]} \emph{W.} Men sæt nu det aldrig havde været guult, saa er det dog besynderligt, at Du som selv skriver Romaner ikke har bemærket det Kneb og ikke selv brugt det.

\emph{[i margen:]} \emph{E.} Jo vist har jeg brugt det; men det kunde aldrig falde mig ind, at Andre brugte det Samme.

\emph{[i margen:]} W. Saa kommer Du da til næste Gang Du vil gjøre Brug af Sligt, at drage den Reflexion ind med i Romanen, at Du meget godt veed, der gives nogle Forfattere, som bruge dette Kneb, for at skuffe Læserne; men at det, Du siger er saa sandt, at de gjerne maa henvende sig paa Husleie-Contoiret og erkyndige dem, om der ikke findes et sligt Huus paa det betegnede Sted.

\emph{[i margen:]} stor Veddestrid mellem 3 Opvakte, om hvo der var den største Synder. –

\emph{[i margen:]} \emph{Nota Bene } I denne Bog vil for det Følgende kun Tankegangen blive at angive, det Øvrige vil efterhaanden blive at udføre paa Lapper Papir.

\emph{[i margen:]} Naar v. S. ret udfoldede sin Skepsis pleiede han betydningsfuldt at lægge sin Finger paa Næsen, for, som Hastværksen bemærkede, dog at have eet fast Punkt under den uendelige Tvivl. –

\emph{[i margen:]} (denne Tale skulde naturligviis udføres langt videre; men jeg vil nu blot antyde).
\end{quote}
\dcite{DD:208  (1837; Pap. IIB3; SKS 17, 280)}
% DD:208  (1837; Pap. IIB3; SKS 17, 280)

\begin{quote}
3die Foredrag.

d. 22 Nov.

Phænomenologie er altsaa Fremstilling af den videnskabelige Udvikling forsaavidt Selvbevidstheden er dens Maal. Forholdet mellem Subjekt og Objekt, men Objektet her er Christus. Nu kunde man vel her fra Troens Standpunkt sige: ja Christus er Maalet, hvortil Alt tenderer (her er den absolute Identitæt af Aabenbaring og Selvbevidsthed.) Men denne Christus maa dog igjen optages, oversættes i os. Her møder strax Vanskelighed: hvilken er den absolute Christus, ligesom Kunsten giver os forskjellige Billeder af Christus, saaledes ogsaa Erkjendelsen. Her ligger jo ogsaa Differentsen mellem Katholiker og Protestanter. De ere enige om det egentlig factiske, om Troes-Artiklerne, ogsaa om tredie om »Aanden«, naar den kun behandles som den tredie Person i Guddommen; men saasnart Spørgsmaalet bliver om \tlg{Tilegnelsen} divergere de, men Aanden er jo netop det subjektive Princip, hvorved det Objektive oversættes i den concrete Bevidsthed – det Samme viser sig ogsaa i det Mindre, en Christus for Anselms Aand er forskjellig fra en Böhmes. (den første mere blot historisk; den anden langt ideellere; langt mere opfattet i sin uendelige historiske og ideelle Ubiquitet.).

En theologisk Phænomenologie er altsaa Udvikling af den theologiske Viden, forsaavidt Selvbevidstheden er dens Maal.

(Phænomenologie kunde ogsaa tages i en vidtløftigere Forstand og omfatte hele Historien, da det ikke blot er den intellectuelle; men ogsaa den sædelige Bevidsthed for Exempel der gjør sig gjeldende. Hegel kaldte sin første Udgave af Phænomenologien sin første Verdens Omsegling. – senere har han dog alene holdt sig til videnskabelige Bevidsthed.)

den er ikke Khistorie, der udvikler den chr Friheds concrete Fremtrædelser, ikke Dogmehistorie, der beskjæftiger sig med Indholdet selv; ikke Symbolik, der udvikler Controverserne, uagtet begge sidste basere sig paa den.

\emph{ Quicquid cognoscitur cognoscitur per modum cognoscentis. }

(Deraf Betydningen af de forskjellige religieuse Standpunkter, Selvbevidsthedens nødvendige Stadier. – \emph{Betydningen af Accomodation – Antropomorphismer.} Antropomorphismen har i Grunden samme Dialektik som den theologiske Phænomenologie. Til enhver Antropomorphisme svarer en Theomorphisme.)

Daubs sidste Skrift udvikler fornemlig Supranaturalisme og Rationalisme og Pietisme.

Phænomenologie

1ste Stadium.

den apostoliske Kirke, hvor Tro og Troens Gjenstand falde sammen, dens paradisiske Tilstand. – Dog paa det følgende Udviklingstrin som paa ethvert senere have vi kun Christus reflecteret gjenem deres Apostlenes Selvbevidsthed i vor. Nu opkommer der Tvivl om man ogsaa har den rigtig, og det eneste Middel man har til at sikkre sig, er, da Skriften paa dette Standpunkt, som underkastet Fortolkning, er utilstrækkelig, en levende uafbrudt Communication med den apostoliske Kirke, der i ethvert Øieblik umiddelbart og reflexionsfrit kan besvare ethvert Spørgsmaal:

\emph{det er den katholske Kirke}

Men nu opkommer der Tvivl om Kirken virkelig har Objektet (Reformation) og jo mere Kirken udvikler Bestemmelserne om Kirken, desto mere fjerner den sig fra sit Objekt.

2det Stadium

\emph{Reformationen,} som kommer under Veier med, at Kirken ikke har Objektet. Da nu hiin levende Communication er afbrudt saa bliver det naturligt at tye tilbage til det Oprindelige. – Bibelen – Men hvorledes skal den fortolkes? af sig selv svarer Reformatorerne og glemme at de selv bragte et heelt dogmatisk System med dem, og at de saaledes bevæge sig i en Cirkel Troen faaes af Bibelen og Bibelen skal fortolkes af Troen som faaes af Bibelen; og at saaledes Døren er aabnet for enhver Fortolkning, og saaledes involverer Reformationen Rationalismen. –

Forsøget med det apostoliske Symbolum hjælper ikke, da man om man end kunde tilveiebringe en historisk Vished, der vilde vise sig en Mangfoldighed af Fortolkninger, og man saaledes kun tilbage foran den historiske Adskillelse.

3die Stadium.

Christendommen søges i Selvbevidstheden.

\emph{[i margen:]} (Baader.)

\emph{[i margen:]} (Baader.)

\emph{[i margen:]} Jeg formoder at det skeer derved, at Kirken kommer til at reflectere over sig selv og derved taber, som det altid skeer foreløbigt ved enhver Reflexion, sit Objekt.

\emph{[i margen:]} K.

\emph{[i margen:]} Bibelen maa fortolkes og ethvert Fortolknings-Princip har en Forudsætning. Erdmann.
\end{quote}
\dcite{Not4:5  (1837; Pap. IIC14; SKS 19, 128)}
% Not4:5  (1837; Pap. IIC14; SKS 19, 128)

\begin{quote}
\emph{Om Forholdet mellem Kant og Fichte.}

1) betræffende Deductionen af begge disse Standpunkter da forekommer den mig at løbe ud paa det Samme. Fornuften har nemlig nu fundet et Standpunkt, hvorved det undgik den Vanskelighed, at Sandhed, da den var et Objekt, derved blev altereret i det Øieblik: Subjektet vilde \tlg{tilegne} sig den. Denne Sphære er den reen praktiske. Den har det ikke her med Objekter at gjøre (confer pagina 235 nederst) uden forsaavidt disse skulle vorde frembragte. Den har derfor intet Dogme, ingen Lære, men kun Postulat, at troe paa Gud er at realisere ham. Men saaledes forholder det sig jo ogsaa med Fichte, og hvad Erd. i Anmærkninger nærmere udvikler om at Kant ikke altid strængt fastholdt sit Standpunkt; men undertiden tog det Ord »Postulat«, i en Betydning, hvorved det nærmede sig til at blive det samme som et Axiom, saa er der jo ikke taget Hensyn dertil i Deductionen.

2) Men hvorfra faaer K og F. et Objekt, uden at dettes Betragtning standser Handling (og gjelder da ikke om denne Betragtning af Objektet her det som K. engang for Alle har gjort gjeldende om alle Objekter.); men da nu den uendelige Aproximation gjenem Handling netop er Sandhed, saa faaer den jo Sandheden gjenem en Usandhed; og hvorfra faaer de overhovedet disse Ideer om Gud etcetera (om man end og fastholder Idee i den Kantianske Betydning saa det ikke er en theoretisk Erkjendelse; men et »regulativt Princip«.) Disse kunne jo i det høieste staae som Resultater af – den uendelige Udvikling, hvilket ikke vil sige stort – det er da et Indhold, som ingen Væren har, hvorom jeg altsaa ikke kan sige stort, deels fordi jeg da maatte synke ned i en theoretisk Betragtnings Snarer, deels fordi jeg ikke har Tid dertil.

Det Objektive kommer jo bagefter confer pagina 244. nederst og 245 øverst, men kommer det bag efter, og tillige som det der i samme Øieblik skal fortrænges af et nyt subjektivt Maal, saa indseer jeg ikke hvorledes man kan kome til om det endog blot hedder at postulere det; thi end ikke det vil der være Tid til. og ethvert Øieblik der anvendes til at postulere det er jo i Begreb med at blive af theoretisk Art. –

Et Bidrag til at forstaae dette Forhold er pagina 245 og 46., hvor han siger: das nämlich, was bis dahin Object gewesen war (der religiöse Inhalt, Gott) ist itzt moralische Weltordnung das heißt seyende Vernünftigkeit, also die Substants des wollenden Subjectes das heißt mit ihm identisch«. Idet nemlig i de Ord: Gud etcetera hos Kant endnu optraadte en om end nok saa lille Rest af theoretisk Betragtning, hvorimod den reen praktiske, som den ene gjeldende, Side træder stærkere frem i de Benævnelser: moralische Weltordnung.

Det er særdeles rigtigt, hvad han siger om Schl. pagina 251: Der Schleiermachersche Standpunkt hat daher das große Verdienst, daß er das Gefühl der bloßen Accidentalität des Einzelwesens geltend gemacht hat. Eben aber, weil es dies \emph{specifische} Gefühl ist, welches er geltend machte, treffen ihn die Einwände, welche die Gefühlstheologie überhaupt treffen, nicht.

Fremdeles pagina 253. »Es braucht aber der große Unterschied zwischen dem mystischen Gefühl der Einheit mit Gott, wo das Einzelwesen als Einzelwesen sich der Wahrheit theilhaft weiß, und dem schlechthinnigen Abhängigkeitsgefühl, wo es seine Einzelheit als verschwindend weiß, nicht mehr besonders hervorgehoben zu werden.

Fremdeles pagina 266 nederst og 67 øverst Schleiermacher konnte auf einem Standpunkt, wo Alles als Seyn gefaßt wird, auch die Verschiedenheit nur als eine Seyende fassen; daher unterscheiden sich bei ihm die Religionen nicht nur als Stufen sondern auch als Arten. Auf einen Standpunkte, wo dagegen die Wahrheit als sich Entwickelndes gewußt wird, werden die verschiedenen Religionen zu verschiedenen nothwendigen Entwiklungsstufen der Religion überhaupt. –

d. 12 Dec. 37.
\end{quote}
\dcite{Not4:45  (1837; Pap. IIC49; SKS 19, 168)}
% Not4:45  (1837; Pap. IIC49; SKS 19, 168)

\subsection{1838}

\begin{quote}
Christi Forsonings objektive Realitæt ogsaa uafhængig af den sig denne \tlg{tilegnende} Subjektivitæt er meget klarlig antydet i Historien om de 10 Spedalske, der jo alle bleve helbredede, medens det dog kun om den 10de, der taknemligt vendte tilbage at bringe Gud Æren, hedder: din Tro haver frelst Dig. –

d. 18 Sept.

\emph{[i margen:]} Hvad var det da der havde frelst de Andre?
\end{quote}
\dcite{DD:144  (1838-09-18; Pap. IIA263; SKS 17, 262)}
% DD:144  (1838-09-18; Pap. IIA263; SKS 17, 262)

\begin{quote}
Saaledes gaaer det ofte med den egentlig Geniale i hans \tlg{Tilegnelse} af Virkelighedens Indtryk, som Een, der sover, og hører Brandallarm, og henfører nu Alt i sin drømmende Verden, medens det Faktiske deri slet ikke træder frem for ham som saadan. Og der ligger noget ret Betydningsfuldt med Hensyn til Forholdet mellem det Poetiske og det Faktiske (hvilket sidste jo ofte kan i høi Grad være det første) naar man siger jeg veed ikke om jeg har seet det eller jeg har drømt det –

d. 19 Sept.
\end{quote}
\dcite{DD:145  (1838-09-19; Pap. IIA264; SKS 17, 263)}
% DD:145  (1838-09-19; Pap. IIA264; SKS 17, 263)

\begin{quote}
Speculativ Dogmatik.

af

Marthensen.

Den speculative Dogmatik viser os Troens sidste Grunde. Nødvendigheden heraf vil indlyse ved at gaae ind paa Christendommen Især i vor Tid vil den være nødvendig; thi tilforn repræsenterede Theologerne Intelligentsen, deres Erkjendelse var Tidens høieste Erkjendelse. Nu synes Philosopherne at have erholdt denne Plads. Disse ere gaaede ud ikke fra Religionen eller Christendommens Standpunkt, men fra Interessen for Tanken for Aanden, og ikke fra nogen enkelt Form for denne. Deres Princip er: hvad der skal gjælde som Auctoritæt for Mennesket maa kunne bestaae for Tankens Prøve. Dette indførte et Brud mellem Religion og Dannelse. Nogle bleve Vantroe, Andre vilde fremmane Luthers Skygge, Catholikerne tyede til Middelalderen til Kirkens Idee; Andre til Apostlenes og de apostoliske Fædres Tid, da Troen kun var indesluttet i Hjertet. Men en saadan Stræben er forgjæves. Christendommens Indhold er bestandig det samme, men enhver Tidsalder har sit Problem, og Continuitæten i disses Række maa erkjendes. Philosophien har tilintetgjort de tidligere Skikkelser af det christne Liv, men den vil frembringe nye. Tidsalderens dybeste Problem er Religionens Forbindelse med Philosophien. Derfor maa Theologen blive speculativ Theolog for at blive Geistlig i sand Betydning.

§ 1.

Den christne Dogmatik er Videnskaben om de christne Dogmer eller Troeslærdomme, og har som saadan de Troendes Samfund eller Kirken til sin nødvendige Forudsætning. Den christne Troe vinder sin faste Skikkelse som \emph{Dogma,} naar den i Troen umiddelbart indeholdte guddommelige Sandhed hæves til almindelig kirkelig Lovbestemmelse. Denne Troens historiske Udvikling til det kirkelige Dogma, saavelsom de under denne Udvikling indtraadte Differentser og Modsætninger i den almindelige Kirkes Bevidsthed fremstilles i den christne \emph{Dogmehistorie} og \emph{Symbolik.} \emph{Dogmatiken} derimod bliver ei staaende ved at erkjende Dogmet i dets relativ historiske Former, men stræber at erkjende det i dets absolute Sandhed. Denne sin Opgave løser den først, naar den som speculativ Videnskab erkjender Dogmet som Idee. –

§ 2.

Som sepculativ Videnskab er den christne Dogmatik i sit Princip ikke forskjellig fra den christne Religions-Philosophie. Ikke desto mindre ere disse to Videnskaber forskjellige, idet den christne Religions-Philosophie væsentligt har til Opgave at erkjende Christendommen som den absolute Religion i Systemet af Religionens forskjellige historiske Former; hvorimod det er Dogmatikens Opgave, at erkjende den christne Religion i det udviklede System af dens Dogma. Religions-Philosophiens Gebeet er saaledes Religionens hele Rige, medens Dogmatiken, der nærmest udgaaer fra den kirkelig-videnskabelige Interesse finder sit eiendommelige Gebeet indenfor Kredsen af de kirkelige Dogmer. Efter Videnskabens nærværende Standpunkt kan Dogmatiken ikke gjenemføres uden religieus-philosophiske Betragtninger, idet den nemlig deels maa erkjende enhver dogmatisk Bestemmelse som en Bestemmelse i Religionens Idee; deels maa mediere sig Dogmets concrete Erkjendelse ved Overvindelse af Standpunkter, der fra de forchristelige Religioner ere gaaede over i Betragtningen af Christendommen –

§ 3.

Den christne Dogmatik, der er den concreteste Erkjendelse af Religionen, inddeles ved en almindelig religions-philosophisk Betragtning over Religionens Væsen og historiske Udvikling. Religionens Væsen som en bestemt Form af Menneskets Evigheds-Bevidsthed lader sig kun begrebsmæssig opfatte, naar den erkjendes i dens qualitative Forskjel fra de andre Bevidsthedens Former, i hvilke Mennesket træder i Forhold til den absolute Idee. Den absolute Idee er ogsaa Gjenstand for Kunsten, og i dens høieste Bestemmelse som Gudsidee udgjør den Philosophiens væsenlige Indhold. Kunsten \emph{kan være} og Philosophien \emph{er} Gudsbevidsthed, hvilket netop ogsaa er Religionens almindelige Bestemmelse. Men dens qualitative Forskjel fra Kunsten og Philosophien bestaaer nærmest deri, at disse Bevidsthedens Former indeholde det Absolute i dets uendelige Objektivitæt, hvorimod Religionen er Gudsideens uendeligt subjektive og reelle Existents i Mennesket. Religionen er ikke blot Bevidstheden om og Erkjendelsen af det Grundforhold, der bestaaer mellem Gud og Mennesket, men den er dette aandelige Grundforhold selv som reelt existerende. Idet Mennesket saaledes i Religionen ikke blot staaer i et ideelt men existentielt Forhold til Guddommen er Religionen at betragte som den menneskelige Naturs Væsen eller som det, der gjør Mennesket til Menneske –

§ 4.

Menneskets religieuse Grundforhold til Guddommen lader sig i Almindelighed betegne som et gjenem et uendeligt Modsætnings og Afhængigheds Forhold medieret Identitæts og Friheds Forhold. Men deri ligger at Menneskets reale Eenhed med Gud kun er mulig derved, at Gud efter sit Væsens evige Sandhed objektiverer sig for Mennesket i Religionen det vil sige at han aabenbarer sig for ham. Var Gud ikke nærværende i Religionen i sin uendelige Objektivitæt, kunde han ikke vides af Mennesket, da vilde Menneskets subjektive Forhold til Gud være et usandt og ufrit Forhold; thi kun \emph{den sande Tanke om Gud} befrier Mennesket ikke blot fra hans Afhængighed af sig selv, men ogsaa fra den \emph{blinde} ufrie Afhængighed af Gud. Men den Tanke, i hvilken Gud i Religionen er aabenbaret for Mennesket, er efter sin Oprindelse forskjellig fra Speculationens sande Tanke om Gud, fordi den er et integrerende og evigt \emph{Moment} i hiint Menneskets existentielle Grundforhold til Gud, der udgjør Religionens Væsen. Ligesom derfor Menneskets Liv har sin evige Forudsætning og Mulighed i hiint oprindelige Grundforhold, saaledes har al menneskelig Reflexion og Speculation sin evige Forudsætning og Mulighed i Aabenbaringen. –

§ 5.

Religion og Aabenbaring er kun mulig under Forudsætning af Ideen om Guds absolute Personlighed, i hvilken Idee Gud ikke blot erkjendes som den absolute \emph{Substants,} men som det uendelig frie \emph{Subjekt} \emph{,} der i sin absolute Substantialitæt veed og vil sig selv og staaer i et frit skaberisk Forhold til Verden som det fra ham selv realt Forskjellige. Kun den sig selv aabenbare Gud formaaer at aabenbare sig for Mennesket, og kun den sin egen evige Tilværelse uendelig affirmerende og villende Gud formaaer ogsaa at ville denne sin Tilværelse i Mennesket, eller at træde i religieust Forhold til ham. Var Gud ikke personlig, da var Menneskets Personlighed en uopløst og uopløselig Modsigelse; thi den menneskelige Viden finder kun sit sande Indhold i den Viden, hvoraf den selv vides og den menneskelige Villie finder kun sit sande Formaal i den Villie, af hvilken den selv er villet. Religionens og Aabenbaringens System lader sig derfor betegne som Personlighedens System. Negationen af den guddommelige Personlighed, der involverer Negationen af den menneskelige ophæver Religions og Aabenbarings Begrebet i dets dybeste Rod.

§ 6.

Ligesaa nødvendig den Erkjendelse er at Menneskets Væsen maa sættes i hans personlige Eenhed med Gud, ligesaa nødvendig er den Erkjendelse, at dette hans Væsen ikke er \emph{umiddelbart virkeligt,} men først gjenem en uendelig Mediation kan naae det. Denne Mediation er Historien og ikke blot Menneskets religieuse Bevidsthed, men ogsaa den guddommelige Aabenbaring følger den historiske Udviklings Love. Paa Grund af det nødvendige Forhold, der finder Sted mellem den Erkjendende og det Erkjendte, mellem Subjekt og Objekt, er nemlig Guds Aabenbaring for Mennesket betinget ved den menneskelige Selvbevidstheds Maal. Saalænge som derfor Mennesket befinder sig paa et relativt Standpunkt af sin Udvikling, kan Gud kun manifestere et enkelt abstrakt Moment (Phase), der corresponderer med Menneskets endelige Bevidsthed. Men i Tidens Fylde da den menneskelige Selvbevidstheds Væsen bliver virkeligt, og \emph{det sande Menneske aabenbares,} aabenbarer ogsaa den sande Gud sig efter sit Væsens uendelige Indhold.

\emph{Anmærkning } I det indre Forhold, der finder Sted mellem Guds Aabenbaring og den menneskelige Selvbevidsthed have Begreberne om Accomodation (συγϰαταβασις) og Anthropomorphismer deres Grund. Disse Begreber have derfor ikke blot subjektiv; men objektiv Realitæt.

§ 7.

Hedenskabet, Jødedommen og Christendommen ere de Grundformer, i hvilke Religionens og Aabenbaringens Idee historisk har realiseret sig. Idet vi ville erkjende disse Religioners Væsen og indbyrdes Forhold maae de betragtes efter den religieuse Idees almindelige Grundbestemmelser. Disse ere: 1) den almindelige Gudsidee, det Grundforhold mellem Gud og Verden, der udgjør Religionens metaphysiske Grundlag 2) Ideen om Guds Apparents i Endeligheden eller den objektive Anskuelse af Forsoningsideen. 3) Ideen om Menneskets \tlg{Tilegnelse} af Forsoningen, der objektiveres i den religieuse Cultus. Idet den religionsphilosophiske Betragtning erkjender enhver af disse Bestemmelser i deres Forhold til Personlighedens Idee begriber den den Hedenske og Jødedommen som de elementariske Forudsætninger og Gjenemgangspunkter for Christendommen, idet den viser, at disse Systemer kun indeholde abstrakte Momenter af Guds og Menneskets Personlighed, hvorimod Christendommen er Personlighedens virkelige System. –

§ 8.

Hedenskabet er med Rette bleven betragtet som det \emph{naturlige} Menneskes aandelige Standpunkt, forsaavidt som nemlig den menneskelige Selvbevidsthed deels ikke har befriet sig fra den objektive Naturmagt, deels endnu er fængslet i sin egen subjektive Naturlighed, i Individualitætens og Folkeaandens Skranker. Men saa længe som Mennesket ikke har fundet sit eget Jeg og erkjendt sin egen Aands qualitative Forskjel fra Naturen, saa længe formaaer han heller ikke at erkjende den guddommelige Aands qualitative Forskjel fra Verden. Derfor kan Gud kun aabenbare sig for den hedenske Bevidsthed efter sin i Verden manifesterede Substantialitæt, men ikke efter sin uendelige Subjektivitæt og Personlighed. Idet den guddommelige Substants reflecteres i de naturlige Folkeaanders Bevidsthed, opfattes den ikke i sin rene Almindelighed i og for sig selv, men anskues kun i sin umiddelbare Eenhed med den enkelte Folkeaand. Da nu ethvert Folk repræsenterer et Trin af den menneskelige Bevidstheds Udvikling, individualiseres den guddommelige Substants i denne concrete Form og ethvert Folk maa altsaa anskue Guddommen under Skikkelse af den Idee, der udgjør dets eget aandelige Væsen, og afpræger sig i den folkelige Individualitæt. De ethniske Religioner eller Folkereligioner fremstille saaledes den i Mangfoldigheden af sine Momenter adsplittede guddommelige Substants. Deres qualitative Forskjel bestemmes ved den concrete Form, i hvilken Substantsen fremstiller sig, og de høiere Former adskille sig fra de lavere ved den Stræben, de udtrykke efter at udvikle den i Substantsen forborgne Subjektivitæt, en Stræben, der i Hedenskabet ikke naaer sit Maal. –

§ 9.

Da den hedenske Bevidsthed kun kan opfatte Guddommen gjennem Naturanskuelsens Medium er Ideen om Guds Apparents, der i den sande Religion kun er Moment i Guds Aabenbaring, den altomfattende Bestemmelse. Hermed er Polytheismen givet, fordi Naturanskuelsen, der mangler Aandens Eenhed, kun kan fremstille Ideen i den ydre Mangfoldighed af dens Momenter. De ethniske Religioner bevæge sig derfor i en Kreds af Theophanier og da den qualitative Forskjel mellem Naturen og Aanden er ubekjendt, anskues Guddommen i de forskjelligste Ktisiomorphismer snart i naturlig snart i menneskelig Skikkelse. I ingen af disse Theophanier fremtræder Guddommen som en virkelig Person, men Gudskikkelserne ere kun endelige particulaire Individualitæter eller almindelige Symboler og overfladiske Personificationer. Først i den helleniske Anskuelse, hvor den hedenske Natursymbolik hæves til Mythologie, opgaaer Bevidstheden om, at den \emph{menneskelige Skikkelse} er den væsenlige Form for Guddommens Apparents. De græske Guder ere ei abstrakte Personificationer men virkelig levende Individualitæter, der ere befriede fra den umiddelbare Naturlighed og forklarede i Ideen; men deels have de Naturgrund til deres Forudsætning, deels er deres Idee endnu ikke den virkelige Personlighed; thi det ethiske Moment, der er Personlighedens dybeste Kjerne, er tilbagetrængt af det æsthetiske, og Kunstreligionens Anthropomorphisme ere forkrænkelig, fordi den heelt er behæftet med Ktisiomorphisme. –

§ 10.

I den hedenske Cultus kan den religieuse Bevidsthed ikke træde i et selvstændigt og direkte Forhold til Guddommen, men kun forsaavidt som den er et Led i Naturlivets eller Folkelivets almindelige Sammenhæng. Forsoningen mangler den dybere Aandelighed, da Folket befinder sig i en umiddelbar naturlig Eenhed med sine Guder og i sine religieuse Feste kun bringer sig denne Eenhed til Bevidsthed ved en forhøiet Livsnydelse. Erkjendelsen af Menneskets Skilsmisse fra Guddommen, der ikke kan udeblive, udvikler sig ikke af Samvittighedsforholdet, men af Erfaringen om de verdslige Tings Forgængelighed. Derfor indeholder Sonoffret kun Resignation paa den verdslige Besiddelse og Nydelse ikke paa den verdslige Villie. Da den sidste Grund til Livets Nød og Ulykke ikke søges i Menneskets syndige Villie, men udenfor den, er Forestillingen om en dunkel Skjæbne, hvilken Menneskelivet er underkastet og uadskillelig knyttet til, den hedenske Bevidsthed og udbreder et Træk af uopløst Sorg over Hedenskabets livsglade Anskuelser.

§ 11.

I Jødedommen er den qualitative Forskjel mellem Natur og Aand opgaaet for den menneskelige Bevidsthed og Gud aabenbares efter sin Subjektivitæt det vil sige \emph{som Gud.} Grundforholdet mellem Gud og Verden er ikke længere Substantialitætsforholdet men Causalitætsforholdet, der fremtræder som Modsætningsforholdet mellem Skaberen og Skabningen. I denne Religion fremtræder først Aabenbaringen \emph{som saadan,} fordi kun Subjektet virkelig kan siges at aabenbare sig. Men denne Aabenbaring er kun relativ og Guds Personlighed er kun abstrakt og particulair ikke absolut, den guddommelige Subjektivitæt er nemlig ei substantiel, i Dogmet om Skaberen fremtræder det ethiske Moment adskilt fra det metaphysiske og Ideen om Guds Eenhed er kun opfattet negativt i Modsætning til Polytheismen. Endelig er Gud kun dette Folks Gud, en Bestemmelse, der giver dette Folk sin store verdenshistoriske Betydning men ogsaa udfolder Spiren til dets Undergang. –

§ 12.

Da den jødiske Gudsidee indeholder den abstrakte Negation af det Naturlige og creaturlige Element, findes der i denne Religion ingen Apparents af Gud. Gud er kun nærværende for Folket i Loven og i sine Førelser med Folket, der ledsages af Tegn og Under, men selv lader han sig aldrig tilsyne. –

§ 13.

Da Creaturligheden staaer som en uoverstigelig Skranke for den jødiske Bevidsthed er Charakteren af den religieuse Cultus den absolute Afhængighed og Herrens Frygt. Den guddommelige Aabenbaring har kun objektiv og positiv Tilværelse i den theokratiske Forfatning men ingen subjektiv Tilværelse og Udvikling i det menneskelige Individ. Den menneskelige Individualitæt er endnu ei befriet til sand Personlighed, thi den troende Israelit medierer sig kun sin Bevidsthed om Gud gjennem sin Folkebevidsthed og saaledes staaer Jødedommen selv i Folkereligionernes Kreds. Istedetfor den fri Individualitæts personlige Forhold til Gud hersker her Legalitæt og udvortes Offercultus. Subjektivitætens dybere Princip udvikles kun gjennem Propheterne, der ved en dybere Betragtning af Syndighed og en idealere Opfattelse af Forjættelserne bragte Universalismen nærmere.

§ 14.

I den christne Gudsidee, der finder sin dybere Explication i Treenighedens Dogma, er Tanken om den subjektløse Substants og det substantsløse Subjekt negeret. Idet nemlig den uendelige Personlighed her tænkes at fuldbyrde sig gjennem en evig immanent Genesis og Proces, er den tænkt som den absolute Gjennemtrængelse af det Substantielle og det Subjektive det vil sige som den i Sandhed aandelige Person. I Bestemmelsen af Guds Forhold til Verden negerer Christendommen paa den ene Side Akosmismen, idet Christendommen ikke betragter Verden og Endeligheden som et Skin, men tilskriver Creaturligheden en relativ Virkelighed og Selvstændighed udenfor Gud; paa den anden Side negeres den abstrakte Dualisme idet Skabningen ikke blot er af Gud, som dens absolute \emph{Aarsag,} men tillige bestaaer ved Gud, idet Gud i sin uendelige Personlighed har gjort sig til Skabelsens substantielle \emph{Grund} og i reel Nærværelse væsentlig og virkelig er tilstæde i sin Verden, og ikke blot er Skabningen af Gud og ved Gud, men ogsaa \emph{til} Gud, idet han medierer sig med sig selv som dens absolute \emph{Formaal.} Da Gud saaledes er Begyndelsen, Midten og Enden og som Fader, Søn og Aand ikke blot har aabenbaret sin Villie men sit Væsen, staaer Gud gjennem denne sin væsentlige Aabenbaring ikke blot i Relation til Verden, men ogsaa i en uendelig Relation til sig selv.

§ 15.

Den christne Idee om Guds Apparents er indeholdt i Incarnationens Dogma eller Læren om Guds reelle \emph{historiske} Menneskevordelse i Christo. I denne Menneskevordelse er Dualismen mellem Gud og Skabningen faktisk ophævet, da den guddommelige Natur saaledes har forenet sig med den menneskelige, at Gud ikke blot er aabenbaret \emph{formedelst} eller \emph{i} men \emph{som} denne historiske Personlighed. Incarnationens Dogma udelukker saaledes den hedenske Betragtning af Aabenbaringen, ifølge hvilken denne Aabenbaring kun skal have været en mythisk Apparents eller Theophanie. Læren om en mythisk Christus, der som uhistorisk tillige vilde være upersonlig, har nemlig sin dybere Grund i Akosmismen, der er uforenelig med Christendommen Paa den anden Side negerer Incarnationens Dogma den jødiske Betragtning, der kun formaaer at opfatte Christus gjennem Endelighedens og Creaturlighedens Kategorier. –

§ 16.

Incarnationen kunde først finde Sted i Tidens Fylde det vil sige Christi Komme er medieret ved Menneskeslægtens historiske Udvikling og kunde ikke indtræde førend alle relative og creaturlige Former af den religieuse Bevidsthed vare udtømte. Men da Christus ikke lader sig betragte som Menneskeslægtens Produkt er hans hele Personlighed, hans Liv og Gjerninger et Under og Miraklet er uadskilleligt fra Ideen om den historiske Christus:. Miraklet er nemlig intet uden den faktiske Løsning og Ophævelse af Antinomien mellem det Aandelige og det Naturlige, det Ethiske og det Physiske, det Guddommelige og det Creaturlige. Idet den umiddelbare naturlige Tingenes Orden herved ophæves, forstyrres den ikke, men fuldkommes til det, som er Maalet for al creaturlig Udvikling. Forholdet mellem Christus og Menneskeslægten er saaledes et uendeligt Modsætnings og Objektivitæts-Forhold, hvilket i sin største Dybde er udtalt i Læren om Christus som Verdens Frelser, hvorfor ogsaa denne Lære er Middelpunktet i den christne Prædiken og Forkyndelse. –

§ 17.

Om den christne Cultus.

Som den troende Bevidsthed, der bevæger sig mellem Synd og Naade, paa den ene Side staaer i et uendeligt Modsætningsforhold til Frelseren, saaledes skal paa den anden Side denne Kirkens Afhængighed af sin Stifter gaae over i et Identitæts- og Friheds-Forhold, idet Guds Rige, som den udviklede Christus skal være det samme i \emph{Mangfoldighed,} hvad Christus er i \emph{Eenhed.} Dette skeer derved, at Gud som Hellig-Aand selv er tilstæde i sin Menighed. Som Kirken paa den ene Side ved Aanden først bliver en \emph{almindelig} Kirke, saaledes er paa den anden Side Aandens Virksomhed en fortsat Personifications og Individuations-Proces, der medieres ved Ordet og Sacramentet. Det christne Dogma om Aanden udelukker ligesaavel den dualistiske Forestilling, ifølge hvilken Aanden kun er et relativt creaturligt Princip, som den pantheistiske, der identificerer Guds Aand med Verdens og Menneskehedens Aand. –

§ 18.

Troen paa Gud Fader Søn og Hellig-Aand, der efter sin substantielle Gehalt er nedlagt i den hellige Skrift, er som almindelig kirkelig Bekjendelse udtalt i Symbolum apostolicum. Dettes Indhold er i Tidernes Løb nærmere blevet udviklet og bestemt i Modsætning til de hæretiske Systemer. Den christne Dogmatik bør bringe sig det Hæretiske til Bevidsthed og mediere sig Dogmernes Erkjendelse ved at eftervise de hæretiske Standpunkters kirkelige og speculative Utilstrækkelighed. Alle hæretiske Systemer lade sig henføre til de i det foregaaende betegnede dualistiske og pantheistiske Principer og maa betragtes som mislykkede Forsøg paa at conciliere Jødedom og Hedenskab med den christne Bevidsthed.

§ 19.

Catholicismen og Protestantismen ere den christne Bevidstheds verdenshistoriske Former og have deres nødvendige Grund i den menneskelige Aands Udvikling. Differentsen mellem Catholicismen og Protestantismen angaaer ikke den christne Sandheds substantielle Gehalt, der eenstemmig anerkjendes af begge Confessioner, den angaaer derimod Forholdet mellem Aabenbaringens Objekt og det menneskelige Subjekt, \tlg{Tilegnelsen} og Mediationen af dette uendelige Indhold. Differentsen opstaaer derfor, idet \emph{Dogmet om Aanden} kommer til at explicere sig i sin hele Dybde. Medens begge Standpunkter indrømme, at den guddommelige Aabenbaring kun kan \tlg{tilegnes} ved den guddommelige Aand afvige de deri, at i Catholicismen denne Aand fixeres til Kirkens ydre phænomene Objektivitæt, medens den i Protestantismen væsentlig aabenbarer sig i den fri menneskelige Selvbevidstheds eget Element. Denne Modsætning, der lader sig gjenemføre i alle Controverspunkter fremtræder klarest i Controversen om det christne \emph{Erkjendelsesprincip.}

§ 20.

Den christne Bevidsthed, der havde forladt den apostoliske Umiddelbarhed og søgte et Erkjendelsens-Princip for at vedligeholde sin Continuitæt med den apostoliske Bevidsthed, og sikkre sig Identitæten med sit uendelige Objekt, fandt dette i den kirkelige Tradition, hvilket Princip videre maatte udvikle sig i det catholske Dogma om Conciliernes og Pavernes Ufeilbarhed. Dette Standpunkts Endelighed blev efterviist af den reformatoriske Bevidsthed, der med Tilsidesættelse af al kirkelig Auktoritæt satte den hellige Skrift som Norm for Erkjendelsen. Men en dybere Opfattelse af Protestantismens Væsen maatte føre til den Indsigt, at kun Aanden selv i den frie Selvbevidstheds Form er det sande Princip for Erkjendelsen, at Skriftens uendelige Indhold kun \tlg{tilegnes} naar det medieres gjennem den fri Tanke. I dette Frihedens og Selvbevidsthedens Princip, hvorved den protestantiske Aand sætter sig i direkte Forhold til Sandheden, have Protestantismens kirkelige og videnskabelige Modsætninger sin Rod, og heri har Nødvendigheden af en speculativ Dogmatik sin Grund.

\emph{Anmærkning } Naar Catholicismen som en Følge af sin empiriske og endelige Charakteer ikke kan frikjendes for at indeholde et Element af den pelagianske Hæresis, saa vil Protestantismen paa Grund af sin ideale og speculative Charakteer nærmest være udsat for den gnostiske Hæresis.

§ 21.

Den speculative Dogmatik er hverken supranaturalistisk eller rationalistisk, den er hverken udelukkende bibelsk eller symbolsk Dogmatik eller en blot Exposition og Beskrivelse af den christne Bevidstheds umiddelbare Indhold, ligesaa lidt er den et System af rene Fornuftsandheder i deres Abstraktion fra det Positive. Den er den levende Conjunktion, Union og Mediation af alle disse Momenter og derved er den deres rette \emph{Begreb.} Idet Dogmatiken er en philosophisk Videnskab er den tillige i Ordets strængeste Forstand en positiv Videnskab, saavist som den i Dogmet indeholdte Sandhed er en ved guddommelig Auctoritæt givet positiv Sandhed det vil sige en saadan, som den speculative Tanke vel kan mediere sig men ikke producere og a priori construere. De positive Momenter, hvorigjennem Dogmatiken medierer sig sin speculative Indsigt ere Skriften, Kirken og den levende Tro.

§ 22.

Da det Dogmatiske i Christendommen er uadskilleligt fra det Historiske, hvis Middelpunkt er Incarnationens Faktum, maa der om dette Faktum existere et infallibelt historisk og religieust Vidnedsbyrd, uden hvilket ingen historisk Kirke var mulig. Dette infallible og oprindelige Vidnedsbyrd om Christus er nedlagt i den hellige Skrift. I Dogmet om dennes Inspiration er indeholdt at dette Vidnedsbyrd ikke blot er Menneskers men Aandens eget. Derfor erkjender den protestantiske Dogmatik Skriften som Erkjendelsens absolute Norm og Canon. Den bør derfor deels direkte deels indirekte eftervise de speculative Ideers Tilstædeværelse i Skriften, ligesom paa den anden Side den speculative Tænkning ideligt paa nye maa styrke sig og forynge sig ved den bibelske Anskuelses rige Fylde.

§ 23.

Den protestantiske Dogmatik er ikke blot bibelsk men ogsaa kirkelig og den beviser denne sin kirkelige Charakteer derved, at den deels er conform med den Protestantismens særegne Lærebegreb, der er nedlagt i Kirkens Bekjendelsesskrifter, deels med den christne Kirkes almindelige Troesbekjendelse, der er indeholdt i det apostolskeSymbolum. I de kirkelige Symboler har Dogmatiken ikke blot en Norm for Erkjendelsen (norma normata ikke norma normans), men den har deri tillige et virkeligt Erkjendelsesmiddel, idet de symbolske Omrids af den christne Lære tillige ere at betragte som Leed i den speculative Begrebsformation. Men ikke blot Kirken og Skriften ere nødvendige positive Midler for den dogmatiske Erkjendelse, men ogsaa den levende Tro; thi hvad Anskuelsen af Incarnationens Faktum er i objektiv Henseende, det samme er Gjenfødelsen i subjektiv Henseende, hvorfor ogsaa den christne Speculation til alle Tider har vedkjendt sig den Sætning: credam ut intelligam. –
\end{quote}
\dcite{KK:11  (1838; Pap. IIC26; SKS 18, 374)}
% KK:11  (1838; Pap. IIC26; SKS 18, 374)

\subsection{1839}

\begin{quote}
Det er et i Sandhed opbyggende Syn, at see den instinktmæssige Sikkerhed hvormed simple Folk \tlg{tilegne} sig Guds Ord, hvorofte de gaae med sand Velsignelse fra en Prædiken, som de dog langtfra have forstaaet; som Himlens Fugle saae de ikke og høste de ikke, og dog nærer den himmelske Fader dem. –

d. 7 April 39.
\end{quote}
\dcite{EE:44  (1839-04-07; Pap. IIA393; SKS 18, 22)}
% EE:44  (1839-04-07; Pap. IIA393; SKS 18, 22)

\begin{quote}
Den protestantiske Lære om Nadveren danner Modsætning til Socinianerne og Andre (til hvem Zwingli nærmer sig) idet den ikke antager det for et blot Erindrings-Maaltid eller et udvortes Tegn; den danner Modsætning til Katholisismen idet den ikke antager en Forandring af Substantsen, og endnu mindre denne Forandring som en efter Consecrationen indhærerende habitus (thi Læren om Ubiquitæten er fremsat for at forklare den \emph{reelle} og \emph{substantielle} Tilstædeværelse i og under Nydelsen og er ikke nogen rummelig Bestemmelse, eller timelig med Hensyn til en Vedvaren.) Men ihvorvel Luther saaledes urgerer \tlg{Tilegnelsens}-Moment, har han dog ogsaa den substantielle og reale Tilstædeværelse mere for Øie end Calvin, der ved at indskrænke Tilstædeværelsen til de Troende, eensidigt udhæver det subjektive Moment. – I den augsburgske Confession er det udtrykt: docent, quod corpus et sanguis Christi vere adsint et distribuantur vescentibus in coena domini. Confessio variata nærmer sig Calvin: quod \emph{cum} pane et vino. De smalkaldiske Artikler udhæve, at det virkelig er tilstæde baade for Fromme og Ufromme. Formula concordiae Vi troe, at Christi Legeme og Blod er virkeligt og substantielt tilstæde og at det uddeles og modtages tilligemed Brødet og Vinen ligesom Ordene selv lyde efter Bogstaven. Forsaavidt Sagen igjen skulde optages til endelig Afgjørelse har man gjort opmærksom paa, at 1) enten maatte de exegetiske Grunde være afgjørende (men her er den Vanskelighed, at Ordet εστι slet ikke kan tilhøre Christus, saa fremt han har talt syrochaldaisk, at hvis εστι stærkt urgeres nærmer man sig Kath:) 2) eller den eller den anden Opfattelse maa tillægges afgjørende Vægt til at forklare Sacramentets Hemmelighed. Men her gjælder det, at begge Momenter ere nødvendige, baade Tilstædeværelsen og Nydelsen, og disse to Sider blive i de forskjellige Opfattelser blot mere eller mindre urgerede. eller 3) eller den eller den anden Opfattelse maatte tilkjendes afgjørende Indflydelse paa Sacramentets Virksomhed. Men herom gjælder det Samme.
\end{quote}
\dcite{Papir 260:260:3  (1839; Pap. IIC33; SKS 27, 221)}
% Papir 260:260:3  (1839; Pap. IIC33; SKS 27, 221)

\subsection{1840}

\begin{quote}
\emph{Hvad det er at ville gjøre Erfaring}

\emph{over Evangeliet om Nicodemus.}

\emph{Bøn}

... og naar Du benaader os med en af de bedre Timer, da vil vi bede Dig, at den ikke maa være een af de mange flygtige Stemninger, der blot urolige vort Sind i forskjellig Retning men at enhver saadan Time maa i sig indeholde en Forjættelse om at den ikke skal være den eneste, for at Sorgen ikke saaledes skal overvælde os, at naar Besøgelsens Time kommer vi da ikke skulde kunne samle os dertil; eller Glæden saaledes henrive os, at Velsignelsens Øieblik svandt os ubemærket og ubenyttet hen. Vi vide vel at det ikke er givet os skrøbelige Mennesker at have det Guddommelige bestandig iboende i os, men vi vide jo dog ogsaa, at det er derefter vi stræbe, og at Du ikke vil lade Dig uden Vidnedsbyrd, og at Din Naade ikke er en Kilde, der udtørres fordi vi øse af den, men som kun sprudler desto rigeligere og springer desto høiere jo dybere vi grave. Vi formaae ikke at bestemme Tid og Sted, men lær Du os at berede vort Sind, saa det i det beleilige Øieblik ikke maa være udygtigt og uværdigt til at modtage ikke koldt og ufrugtbar. –

at gjøre Erfaring ja det er en stor Tanke at have levet meget, at eie en uudtømmelig aandelig Skat som aldrig kan tages fra os ja det er et stort Maal. Det er dette Ønske, der begeistrer den Yngre, naar hans Hjerte svulmer af ubestemt Længsel, han ønsker at Intet maatte undgaae hans Blik, Intet maatte være uprøvet af ham, at han maatte forsøge sig i Alt, vel ofte forbunden med det forfængelige Ønske, med Rette at kunne sige ogsaa dette har jeg fornummet. Og det er ikke blot det Glædelige men ogsaa det Sørgelige ønsker han at have være prøvet i, for senere med Glæde at erindre sig om, at ogsaa Sorgens Bæger blev ham rakt, at han har smagt dets Bitterhed, og dog holdt sig opreist.

Men her er det Feilagtige allerede at ville have oplevet Alt, før han ret er begyndt at leve.

Og nu den Ældre ham forudsætte vi jo om at han har oplevet det og om den ungdommelige Higen end ofte indbilder sig at det, den Enkelte har oplevet, er Intet mod det, han selv skal opleve, om han end med Stolthed tænker sig selv som den bedagede Mand, der er langt rigere paa Erfaring, saa har han dog en vis Ærbødighed for det den Ældre eier fordi det dog er erfaret; thi hvor sørrgeligt vilde det ikke være, at tænke sig en Olding, der sit hele Liv igjennem blot havde været samtidig med sit eget Liv, vidste om det som han veed om enhver anden udvortes Begivenhed, vidste Time og Klokkeslet, men den indre \tlg{Tilegnelse} deraf, – Vi see saaledes at hverken den urolige Higen ud i Verden, eller den blotte Leven i Verden er nok for at gjøre Erfaring. I en vis Forstand gjør ogsaa de vel Erfaring, men da der ikke i dem udvikler sig den Magt, der overvinder det Hele og overskuer, blive de aldrig \emph{klog af Erfaring.} Der til hører den villies-Akt, som vil gjøre i Erfaring, \emph{som i Sorgens Øieblik} vil \emph{erindre sig den Glæde, den har oplevet, i Glædens Time ikke glemme den Sorg, der truer}
\end{quote}
\dcite{HH:19  (1840; Pap. IIIC13; SKS 18, 136)}
% HH:19  (1840; Pap. IIIC13; SKS 18, 136)

\begin{quote}
Dersom det stod sig rigtigt med Philosophernes Forudsætningsløshed maatte den ogsaa gjøre Rede for Sproget og dets hele Betydning og Forhold til Speculationen, thi heri har jo Speculationen et Medium den ikke har givet sig selv; og hvad Bevidsthedens evige Hemmelighed er for Speculationen som Eenheden af Natur-Bestemmelse og Friheds-Bestemmelse, det samme er Sproget deels det oprindelig Givne, deels det frit sig udviklende. Og ligesaalidt som Individet, hvor frit det end udvikler sig, nogensinde kan komme til det Punkt, at det bliver absolut uafhængig, da tvertimod den sande Frihed bestaaer i frit at \tlg{tilegne} sig det Givne og altsaa i igjennem Frihed at være absolut afhængig, saaledes ogsaa med Sproget, og vi møde jo vel ogsaa stundom denne misforstaaede Tilbøilighed til ikke at ville tage Sproget som det frit \tlg{tilegnede} Givne, men at give sig det selv, hvad enten den nu viser sig i de høieste Regioner, hvor det da gjerne ender meda Taushed, eller i den personlige Isoleren i et Rotwelsk Galimathias. Saaledes lader maaskee og Historien om den babyloniske Sprogforvirring sig forklare, at det var det arbitraire Forsøg paa at constituere et vilkaarligt dannet Fælleds-Sprog, hvilket Forsøg, da det manglede det Alt sammenholdende Fælleds netop maatte adsplittes i de løseste Differentser, thi her gjælder det totum est parte sua prius, hvilket man ikke forstod.

d. 18 Juli 40.

\emph{[i margen:]} Negationen af Sproget.
\end{quote}
\dcite{Papir 264:264:11  (1840-07-18; Pap. IIIA11; SKS 27, 238)}
% Papir 264:264:11  (1840-07-18; Pap. IIIA11; SKS 27, 238)

\subsection{1841}

\begin{quote}
Marheinecke.

at den er for alle Mennesker. Forsoningen forudsætter et Widerspruch mod Gud i alle Mennesker. ϰοςμος betyder i det Nye Testamente den fordærvede Verden. Jeget trækker sig tilbage fra Gud, sucht sich selbst, Selbstsucht og bliver Entzweiung og Entzweiung med Gud. Fordærvelsens Overgang fra Ich til Verden er medieret gjennem Kategorier Einzelnes, Besonderes, Allgemeines. Det Onde er Alles Krig mod Alle. Fordærvelsen viser sig som \emph{Thun} \emph{,} hvori Verdens Gesinnung udtaler sig. Det Negative er væsentlig Løgn, og Verden mangler Sandhed. Dette er den objektiv Trang. Fordærvelsen viser sig som \emph{Liden} \emph{,} den er Hemnung, Vernichtung des Lebens, der Schmerz. Verden behøver Sandhed og Frihed. Fordærvelsen er \emph{Skyld} \emph{,} Identitæten af Thun og Leiden, er Skyld-Last, som hviler paa den, ikke som om den var lagt paa den udvortes

\emph{Forsoningens Nødvendighed fra Guds Side.} Fordærvelsen kan altsaa ikke være i det guddommelige Væsen selv. Gud er ikke der Ungrund til det Onde, heller ikke blot Abfall, der næsten bliver et Zufall, og ikke der Fall, Syndefaldet, det er Verden der fri adskiller sig fra Gud og Verden er den af Gud forstødte. Dog er den ikke kommen ud af alt Forhold til Gud. Gud siger: fiat justitia, men tilføier ikke pereat mundus. Verden maa ikke tænkes i Gud. Guds Retfærdighed vil bortskaffe Synden ved Straf. Uden Guds Retfærdighed var Synden heller ei Uretfærdighed. Tillige er Retfærdigheden Kjærlighed og denne er Retfærdighedens Form. Indvendige Gud har straffet en Uskyldig Svar 1) det er frit af den Uskyldige paataget sig 2) det er Gud selv. Som den forsonende manifesterer Gud først sin evige Kjærlighed. Gud er den forsonende og Forsoningen selv, der er ligesaa vel Guds evige That som Skabelse og Opholdelse.

\emph{Muligheden.} Det der fjerner Gud fra Mennesket, kan ikke være det der forener ham med Gud. Vel siger man, at Synden er Aarsag til Forsoningen. Det Onde kan ikke være Grund til det Gode. Det Onde maa ikke være absolut uden al Godt, det Gode er Trangen til Forsoningen. Mennesket er Syndens Træld, men ikke bestemt til saadan Trælddom. Til dette Gode kan Forsoningen anknüpfen. – Troen paa Gud, hvorledes den end viser sig, har denne Bevidsthed, at Forsoningen ikke er en Fiction, men at Gud ifølge sin Natur er den evig forsonende Kjærlighed. Ingen Religion har manglet ganske denne Tro. Offeret beroer paa den Tanke, at Gud selv er Forsoneren. Heri er Troen paa 1) Guds Retfærdighed, hvorved han ophæver alt Endeligt som saadant, og 2) paa Guds Kjærlighed, der hæver det op til Gud. Forsoningen gaaer i alle Skikkelser ud fra Gud. At Guds Kjærlighed er virkende er Troens Indhold, ved Troen er først Mennesket der versönliche, der versönbare, men denne Troe skylder han Naaden. – Det guddommelige Væsen er det allgenugsame og allgenügende, som tilfredsstiller Alle dem, der ønske at komme ud af Widerspruch, det er sig ikke alene genug, det maa sig selv genugthun, dette er den guddommelige Kjærlighed og heri ligger en Historie, som er alle Religioners Historie. Forsoningen er evig, men det er dog egentlig kun Ideen og Muligheden, det vil sige naar Forsoningen er kan det kun skee ved Gud. Den faktiske Forsoning. Indeholdt Ideen ikke mere end Mulighed saa var den en svag og magtesløs, men er den til alle Tider mere eller mindre realiserende, saa har den ogsaa i Tidens Fylde viist sin fuldk. Realitæt. Uden Ideen lader den faktiske Forsoning sig ikke forklare – Troen paa Forsoning har ikke til Gjenstand et Menneske, der selv er Synder, men den forsonende Kjærlighed, det guddommelige Væsen selv, dog ikke hiinsides Virkeligheden, Verden og Menneskeslægt. De andre Religioner mangle Identitæten af Forsoningen og Forsoneren (Præst – Offer som to forskjellige). Midler var og Moses, men viiste dog kun Muligheden. Skal det gaae over i absolut Virkelighed, saa maatte Gud selv blive Menneske, og gjorde og fuldbragte hvad Mennesket skulde, og det med Gud forenede Menneske paatog sig, hvad der ikke passede sig for Gud, den dybeste Fornedrelse. Satisfactio er saaledes i dobbelt Henseende vicaria. Det Store i Christendommen med Hensyn til Forsoning er at Idee og Realitæt fuldkommen ere sammensluttede, og kan ikke rives fra hinanden.

\emph{Virkeligheden.} 1) Begrebet Forsoning 2) Momenterne 3) der Zweck.

1) \emph{Begrebet} \emph{.} Guds Menneskevordelse i Christo er en Antydning af Forsoningen. Hverken Gud eller Menneske kunde tilveiebringe Forsoningen, er ist an sich die Versohnung. Den virkelige Forsoning er en Thun og Vollbringen. I christi Menneskelige Natur kom Forsoning til Virkelighed. Christus staaer i Forhold til Verdens Synd, selv skyldfri. Den menneskelige Natur har vel Trang til Forsoning, men kan ikke forsone. Der behøves et reent Offer. Christus bærer al Verdens Synd. a) Aufnehmen der Sünden in sich b) das Vertilgen derselben. Dog maae disse to forstaaes dialektisk overgaaende i hinanden. Hans Liden var ikke en Liden i Phantasie. Verdens Synd sluttes inde i en Uskyldig og bæres af ham. Synden har dog bestandig en udvortes Side, og forholder det ikke sig ligesaa udvortes som alle andre Offere? I. Nye Testamente forestilles Christus som Præst og Offer. Det synes som den guddommelige Retfærdighed tager en Uskyldig for de Skyldige. Det er ikke det for alt Godt fremmede Onde som Forsoneren optager i sig, har det intet Forhold til det Gode saa er det et reent negativ og nedsjunken til physiske Bestemelser. Moralsk Ondt er det kun ved Bevidsthed om det Gode. Kun saadanne Synder har han baaret, hvori der endnu er et Glimt af Forsoningstrang. Hans Bæren af Verdens Synd vinder saaledes i Intensitæt, hvorledes kan den gewissensfrie bære den? (thi Christus har ikke Samvittighed, er større). Bevidstheden om Synd, jo dybere den er, desto mere er den gjennemtrængt af det Gode, og dens Ophævelse mulig.

Det Individuelle og Universelle. er Christus blot det enkelte Menneske, saa har Menneskeslægten ikke været indvirksom; har han blot baaret Verdens Synder ikke sine egne, saa er der intet Forhold. Det Første er Rationalisme; det andet Supranaturalisme. Christus opoffrer sig vel for sit Folk, men det er kun en relativ Opoffrelse. Af Analogier kunde man sige, at noget Lignende er mulig for hele Menneskeheden. I Christus er Middelpunkt af Verdenhist. han er et enkelt Menneske tillige hele Artens Liv, som Centrum bestemer han hele Peripherien i sig, som einzelne Menneske er han det almene, som det almene det enkelte, det er den hele Menneskehed der i ham bringer sig Gud til Offer. Her er han Analogie til Adam den anden Adam.

Det egentlig Forsonende i Christo er deels den grændseløse Kjærlighed og Selvfornegtelse i at underkaste sig Faderens Villie. Verdens Selbstsucht staaer i hans Person ligeoverfor den uendelige Kjærlighed, og dette Forhold er et indre, og derved hiin ophævet ved denne. Idet Verdens Lidelse gaaer over i Christi Lidelse, forvandles den fra Nødvendighed til Frihed. Formedelst den uendelige Kjærlighed træder istedetfor Skyld den uendelige Frihed. – Guds Forhold til Forsoneren er ikke blot et trøstende og styrkende, det er et udvortes Forhold, den væsentlig Eenhed komer i Betragtning, han adskiller sig som Menneskets Søn, beder til Gud, det er ikke Faderen, men Sønnen der lider; men det er i begge Naturers Eenhed, og den guddommelige bliver ikke uberørt deraf. Det er Christi Fortjeneste at begive sig i den yderste Endelighed, hvori han næsten forsvinder, men derved opløftes Endeligheden. I Christi Deeltagelse i alle menneskelige Affectioner ligger, at det ikke kan blive ved Døden, men ved Opstandelsen.

Momenterne.

De faktiske Zustände, hvori hans Gesinnung viser sig, an sich ere de satte ved Forsoningens Begreb; die Erscheinungsweise er en udvortes. Ein \emph{Thun} \emph{,} i Modsætning til Verdens onde Thun, en Thun des Einzelnes bezieht sich mittelst des Besonderen auf das Allgemeine. Hans Reflexion i sig er ikke egoistisk. Hans Thun er væsentlig aandelig og saaledes uendelig. Som uendelig svarer det til Ideen, og Forsoningens absolute Idee, der er virksom i ham, udelukker al anden menneskelig Thun ved Siden af ham som forsonende, alt tidligere og sildigere forekommende bezieht sich auf Christum. Denne Thun kaldes i det Nye Testamente Kjærlighed, den er den uendelige Kjærlighed. Identisk med hans Kjærlighed strømmer gjennem hans Thun til Verden uendelig Kjærlighed.

\emph{Leiden.} er aandelig og uendelig, det er det deprimerede Thun, man maa ingenlunde forstaae det som en reen Passivitæt. Denne Liden fremstiller sig først som \emph{legemlig} (Korsfæstelse etc). Den sandselige Affektion, det legemlige Blod er ikke det Forsonende, det er kun en Liden i Endelighed, men det har Betydning som det Aandeliges Erscheinung. \emph{Sjælelidelser.,} Sjæleangst, er vel indvortes men dog endnu sinnlich. Dødsfrygten som almindelig menneskelig har han gjort til sin. \emph{Aandelig Lidelse.,} en Liden i det Godes og Sandes Ideer. Al Liden af endelig Natur er optaget i denne aandelige Liden. Den ene Røver, er den Verden durch die er leidet, den anden er den Verden, für die er leidet, dette er den bodfærdige Røver. Døden er den yderste Widerspruch, som den yderste Liden, bliver Gud vel ikke uden Andeel, men at lide og døe er ikke Guds Bestemmelse. Forsoningen er Aandens Tilbagevenden fra dens Modsættelse med sig selv i sin Eenhed med sig selv. Aandens Rolighed i denne Bevægelse gjennem sin Modsætning, er det høie, og deri befinder han sig allerede i sin Ophøielse. Døden er Begyndelsen til et nyt Liv, det er Lidelsens Tilbagevenden i Thun. Eenheden af Thun og Liden er \emph{Lydighed} \emph{,} det er Begrebets Totalitæt, Thun og Leiden er kun dens Momenter er dens Erscheinung. Kirken indbefatter Alt under obedientia som den deler i activa og passiva. Den gjør det udvortes som om Handlen og Liden traadte tilfældig til. I Lydighed er forenet Nødvendighed og Frihed, Müssen og Wollen. Nødvendighed er Loven, Friheden ophæver den i sig, forvandler det til Villen.

Disses Eenhed i Lyden opløses i den nye Dogmatik. Rationalismen lægger meest Vægt paa Thun, og hans Liden har kun moralsk Nødvendighed, for ikke at blive Forræder mod den erkjendt Sandhed. Lidelse stod ikke i hans Villie, Døden ikke en fri paatagen. Døden var den Declaration, at han hellere vilde døe end blive Sandheden utro, og den var Segl paa hans Lære. Christus adlød sin Pligt. Døden beholder saaledes blot et udvortes Forhold, og den bliver ikke absolut nødvendig. Alt kunde ogsaa være opnaaet uden Død, og Døden som saadan beviser ikke Sandheden. – Supranaturalisme negerer det Frivillige, mors voluntaria, saaledes aner den ikke tænker at den ophørte at være den nødvendige Det Vilkaarlige har kun Skin af Müssen og er ingen Lyden. – Eenheden af begge er den uendelige Lydighed. Christus er ikke en privat Person, og hans Død ikke blot den endelig Død. Han bliver et Offer for sit Folk, der i Dommen giver sig Skin af Retfærdighed; men dette er kun en endelig Lydighed. Det Nødvendige er dermed her det fri. Lydigheden er den Frihed der er sat sig selv. Han har ophævet Loven, det er opfyldt den. Denne Magt er vel over og uden for ham, da det er Faderens Villie. Den Villie gjør Sønnen, dette er Nødvendigheden og Loven; men Faderens Villie er ogsaa Sønnens. Forskjellen er vel sat men i samme Øieblik hævet. Hans Lydighed er uendelig fri og uendelig nødvendig. Forsoningen er den absolute Lyden. For Forsoningens Skyld er Lyden geleistet, idet den er geleistet er Forsoningen fuldbragt.

\emph{ Der Zweck der Versöhnung }

Strauss har ophævet den historiske Proces, idet han sætter Forsoningens Idee, uden en Forsoner, i Totalitæten af Menneskeheden. Det er kun den subjektive Side, \tlg{Tilegnelsen}. Uden den objektive Forsoning vilde Bevidstheden ikke være i Mennesket. Feilen er at Forsoningen an sich realiserer sig umiddelbart uden Midler, Midleren bliver selv her blot en Tanke, der vel ikke har Kraft til at fuldbringe Forsoningen i et enkelt Menneske, men vel i alle Mennesker Men Mennesket er ikke forsonet, men skal vorde det. Forsoningen er et Hiinsides som ved Christus bliver et Dennesidigt. Christus faaer ellers blot Betydning som Exempel. Alle Menneskers Forsoning er betinget ved Christi uendelige Lydighed. Lydigheden i de enkelte Mennesker skal til sin Grund have den Lydighed, som Christus har viist sin Fader. Dette Zweck er endnu ikke opnaaet i Christi Død, og forsaavidt er Forsoningen ikke fuldbragt, thi Menneskeheden er endnu ikke den ved Gud virkelig forsonede, forsaavidt er Christus Død en vicarierende, han offrer i sin Død Faderen den hele Menneskehed Realisationen er Troens Værk. De foregaaende Offere i Verden vare de uegentlige Offere i Forhold til det egentlige som er Christus. Som nu Alt her tenderede til Christi Offer, saaledes har Troen i omvendt Retning sin Stræben tilbage til Christi Offer. Lydighed er Gud kjærere end Offer. I Christo faldt begge Dele sammen. Troen er den subjektive \tlg{Tilegnelse}. Troen er fri: »Lader Eder forsone med Gud«. Det skeer derved, at Christi Livs forsonende Kjendsgjerninger, der som saadanne staae uden for os, gaaer ind i os. Dette anseer Strauss for en rethorisk Anklingen an der Spekulation og mystisk. Men dette er Tilfælde med al objektiv Sandhed naar den gaaer ind i den subjektive Leven. Troen lader endnu dunkle Dele tilbage. Først derved at den hele Menneskehed er den virkelig forsonede er Forsoningens Zweck opnaaet, uden det er Forsoningen blot Mulighed. Dog maa man ikke forstaae det som om Alt blot kom an paa det Indvortes. Det er Mysticisme. Denne Maxime er usædelig, da den antager Forbedring mulig, uden at Mennesket er bragt tilbage i et Forhold til Gud. Mysticisme og Rationalisme behandler Christi Lidelse og Død blot som udvortes, ikke som aandeligt Faktum. Ligesaa utilfredstillende er Christi Fortjeneste, naar man kun udvortes trøster sig ved det, beroliger sig ved Troen om Forsoningen, uden Gjenfødelse. Saa er Christi Fortjeneste Sjælen et Fremmedt. Det er ikke i den sande Tro, den slutter den Troende sammen med den, paa hvem den troer. Troen er Tilregnelse af Christi Fortjeneste, og det skeer af fri Naade, ikke merito. Troen er ogsaa Handlen, som indeslutter Christi Retfærdighed i sig, kun Forskjellen bliver mellem Christi uendelige Lydighed og Menneskets endelige; naar Christus lyder, saa vil det sige, han forsoner; naar det hedder om Menneskene, de lyde, saa betyder det, at de ere mulige til Forsoning, eller forsonede.

\emph{3. Christi kongelige Embede.}

Det er Sandheden af hans prophetiske og præstelige Embede. Derved gaae de over fra Timelighed til Evighed. Mig er givet al Magt, »er givet« udtrykker Afhængigheden, »al Magt« udtrykker Identitæten. hans kongelige Embede er hans evige Udøvelse af hans prophetiske og præstelige Embede. Aandens Rige concret tænkt er Kirken. Her er Christus i sin forklarede Menneskelighed Konge.
\end{quote}
\dcite{Not10:8  (1841; Pap. IIIC26; SKS 19, 288)}
% Not10:8  (1841; Pap. IIIC26; SKS 19, 288)

\begin{quote}
III.                         III

Læren om Aanden.

1. \emph{Treenigheden} \emph{.}

\emph{Den bibelske Lære.} i det gamle Testamente dunkel og skjult. Dog tiltroer Kirken det gamle Testamente denne Lære, dog kommer først Læren i Christendommen En stræng Jøde kan ikke have Treenighed. Det er sandt, hvad Calixt sagde, at Treenigheden ikke lader sig demonstrere af det gamle Testamente Det er deri kun som Hemmelighed, det er derfor deri og dog ikke. Det Skjulte er an sich men det er. Det har Udtrykkene Fader, Søn, Aand, og saaledes en Selvadskillen i Gud, der dog ikke udelukker Eenheden. Alle Stæder derfor man har anført, synes man mere at have lagt denne Lære ind end draget den ud. Ψ 33. Det ypperstepræstelig Tre Gange helligt. – i det Nye Testamente Adskillelsen af Fader og Søn er fremtrædende. Aand bruges ogsaa πνευμα αγιον. Mth. 28, 29. 2 Cor. 11, 13. Paraklet, betyder ikke Hjælp men Hjælperen. Gaverne komme fra Aanden, χαϱιςματα ere altsaa adskilte fra Aanden, og Aanden adskilles igjen fra ὁ ϰυϱιος og ὁ ϑεος. 1 Cor. XII, 5-6. Aandens Virksomhed betegnes ved personlige Egenskaber, see, høre Joh. 16, 13. Stundom adskilles Aanden ikke fra Gud, hvor han nævnes i Forbindelse med Fader og Søn. alvidende 1 Cor. 2. Joh. 16, 13 Almagt 1 Cor XII, 4-11. – inspiration Joh. 14, 15. 1 Pet. 1, 21. Menneskets moralske Forbedring Joh. 3. Apostelen sværger ved ham som ved Christo. Synden ved den er den største. Mth. 12. Kaldes Gud Act: 5, 3.4. – ogsaa Eenheden af Fader Søn og Aand Mth: 28, 18-20. – Bretschneider vil at Døbeformelen ingen dogmatisk Interesse har. –––– \emph{Den kirkelige Forestilling,} den \tlg{tilegner} sig alle de Spoer, der tidligere i Verden maae være givet til en Guds Treenighed. Intet Folk er ganske uden Spor. Inder, Chineser, Perser have Antydninger dertil. Den platoniske Forestilling. Trinitæten gaaer i Almindelighed over hos Hedninge i Fleerguderie. Hos Plato er der Eenhed, han kalder den πατηϱ, λογος, νους, σοφια, og Verdensjæl. I Timæus adskiller han de andre Principer for Gud. Spørgsmaalet er om han hypostaserede de guddommelige Væsener. I Neo-Platonismen komer disse Ideer i Forbindelse med Dogmet., og Kirkefædrene \tlg{tilegnede} sig det Sande deri som det christelige. Ogsaa de alexandrinske Jøder havde en bestemt Art af Treenighedslære, Philo: 1) høieste Gud, 2) λογος, δευτεϱος ϑεος, 3) πνευμα Kabbalisterne havde den ogsaa midt i den strængeste jødiske Monotheisme. – Symbolum apostolicum – Formelerne variere deri, at de ere udførligere og kortere. Fast Skikkelse vandt Formelen først i det apostoliske Symbolum, som efter sin Redaction først falder i det 4d Aarhundrede –

Paulus Samosatenus og Sabellius, den absolute Identitæt, uden al Unterschied. Einerleiheit. – Photinus og Arius. her er wesentlicher Unterschied. Sønnen er en Creatur. – de anerkjende ikke Guds absolute Spiritualisme. Macedonius. Synoden i Nicæa og Constantinopel. nicænoconstantinopolitanum. – Den evangeliske Kirke. slutter sig til de romerske Symboler i dette Punkt. Fordømme Alle der nægte Trinitæten især gamle og ny Samosaterne, Socinianerne. – Den protestantisk Dogmatik. Det indre Forhold mellem de 3 Personer, characteres ellerproprietates personales. Faderens αγηνεσια, generatio activa. Sønnens γενησις, generatio passiva. fælleds for begge er spiratio, Aanden er processio e patre filioque. (Tilsætningen filioque). Den protestantiske Kirke forsvarer den.

Dogmets Begreb.

\emph{Troens Forhold til Læren.} Kirken betegner denne Lære som Hemmelighed. Følelse og Forstand finder liden Næring, Forstanden Modsigelse deri. Det er vel et Mysterium, men ikke absolut ukjendelig. Den spekulative Erkjendelses Opgave var at indsee den Modsigelse, der blev for Forstanden mellem Eenhed og Fleerhed. Eenheden gaaer vel ud i Trehed; men Trehed ikke tilbage i Eenhed. Abstrakt opfattet er Trinitæten ikke Troens Gjenstand. Han troer paa den objektive Personalitæt, først i Reflexionen forlanger han at begribe Eenheden i Flerheden. Herved er Viden og Speculation indtraadt. – Gud er ikke i almindelig Forstand alle Menneskers Fader, han er egentlig kun Fader forsaavidt han er Sønnens Fader og derved Fader til alle. Sønnen er egentlig en negativ Bestemmelse med Hensyn til alle andre Mennesker, de ere ikke i den Forstand Søn som andre Mennesker Supranaturalisme holder især paa Sønnen og Aanden, Rationalismen paa Faderen. Troen tilfredsstiller sig med den hellige Aand, uden at bekymre sig om hans Forhold til Fader og Søn. Gud som Aand er det Standpunkt, hvorfra Videnskaben seer Treenigheden.

\emph{Indvendinger mod denne Lære.} Reflexionen mener vel at man maa tænke Gud som den treenige, men siger ikke, at Gud er den treenige. Der er en Spaltung mellem Seyn og Denken. Det er den reflekterende Bevidsthed, der gjør Indvendingerne. Eenheden holder man fast paa, i Modsætning til Polytheismen. Vi troe alle paa een Gud. eller mere indifferent vi troe alle paa en Gud. Det er rigtigt, forsaavidt der ingen Religion har været saa slet, at den jo har havt en Gud. Spørgsmaalet er, hvilken Gud er den, Du troer paa. Troen vil give Ideen et Indhold. Paa denne Forskjel maa Tænkningen indlade sig. Vi troe ikke alle paa een Gud, om end paa en Gud. Den abstrakte Forstand udvortes og negativt, Gud er ikke Verden og saa videre. Man sætter al Forskjel udenfor Gud og saaledes mener man kun at beholde Eenheden.

Forstanden opfatter Eenheden numerisk, og seer ingen Grund til at blive staaende ved 3, man kan faae en Uendelighed, en Endeløshed. – Augustin siger om Personerne, at de ere in se invicem.

Forstanden stiller Personerne op efter hverandre, og faaer derved, at Sønnen og Aanden ere Faderen underdanige og mindre end Faderen. Alle Forestillinger om Tids Succession maa fjernes.

Dogmet selv.

Hvor forskjellige Formler man end har brugt, saa er der dog visse Bestemelser der er absolut nødvendige Substantialitæt, Subjektivitæt, Identitæt. – Lessing siger, at det fuldkomne Væsen ikke har kunnet beskjæftige sig med Andet end Betragtningen af det Fuldkomneste, her er jo aabenbart Forskjellen og Eenheden, denne er nemlig Betragtningen. Videre siger han, at Gud skabte sig et Væsen som ingen af de Fuldkommenheder manglede, som han selv havde. Her er Feilen at sige, at han skabte. »Dette Væsen er Gud selv og ikke forskjellig fra ham, fordi saa snart man tænker Gud, tænker man det og man kan ikke tænke Gud uden dette, man kan kalde det et Billede, men det er det fuldkommen lige Billede. Dette er Aanden, der er Harmonien.« Schelling bragte Philosophien dette Dogme. »Det Absolute bliver sig selv objektivt i et Modbillede, som tillige er det Selv. Gud overfører hele sin Væsenhed paa det, hvori han bliver objektiv.« Hegel. –

Seer man i alt Sligt kun Construktioner, det er en subjektiv Bestemmelse, saa bliver Alt til Intet.

Dogmet har disse Momenter 1) Gud er absolut Substants. Seyn som saadan., Den abstrakte Væsenhed; man kan ligesaa godt sige: das Seyn ist Gott. han erkjendes som værende in sich er causa sui. Spinoza. Gud er Fader. Her er allerede Personligheden udtalt, Spinoza identificerer Substants med Natur. Han har ikke sin Grund i en anden; thi saa var Sønnen. Gud er Fader uden nogensinde at have været Søn. 2) den absolut Subjekt. für sich. I Naturen er det Materielle og saa videre Naturen selv ikke für sich. Seyn eller Wesen für sich er kun mulig i Tænken, er Bevidsthed. Tænken er kun für sich. Tænken er kun for Tænken. Den Tænken som er Gud, er denkendes Seyn, Realitæt og Idealitæt, Substants og Subjekt. Fixerer man Forskjellen mellem Seyn og Tænken, saa er der ingen Sandhed deri, saa er det et Gud uværdigt Tænken. I den christne Religion er Guds für sich Seyn hypostaseret og forestillet i Tanken om Sønnen. Det guddommelige Væsen aabenbarer sig som Fader i Sønnen. An sich er Gud den uudforskelige. 3) Forholdet mellem Fader og Søn er et negativt, Faderen er ikke Sønnen, Sønnen ikke Faderen. Det maa negeres. Her er Læren om Treenigheden, der manifesterer sig for enhver, der veed at Gud er Aand. er ist der Geist von Vater und Sohn. I disse to er Gud sig selv en anden, forholder sig til sig selv. Han er Identitæten af Differents og Identitæt. Faderen er Aand, Sønnen er Aand, men det er ikke en forskjellig Aand. Aand som Fader er Hellighed, Sønnen er Sandhed, Aanden er Kjærlighed, som forener begge. Kjærlighed er netop det, at den ene er i den anden. Aanden som Kjærlighed hæver Faderens og Sønnens Personligheder evig op i sig, det vil sige, det at de ere Personligheder, ligger netop i, at de ere Aander.

2. Læren om Naaden.

\emph{1. Kaldelsen.} Udvælgelsen. Forudbestemmelse. Augustinus. – Den nyere Dogmatik. Saavel Universalister som Particularister beraabe sig paa Skriften. Bretschneider mener, at Skriften har en dobbelt Læretypus. Saaledes er da den hellige Skrift bragt i Modsigelse med sig selv. Knapp erklærer den bibelske Modsigelse for tilsyneladende og søger at hæve den ved Sprogfortolkning.

\emph{Dogmets Begreb.}

Guds Regjering af Verden er netop, at han som Aand er i Verden. Kaldelsen er den guddommelige Bestemmelse med Mennesket, men denne Bestemmelse er an sich, er endnu umiddelbare, den ene kan blive en Werden für sich, og deri en Werden für Gott, endelig disse to i Eenhed.

Werden für sich er ikke uden det umiddelbar Forhold til Gud, men dette Forhold er det ansichseiende. Menneskets Vorden er ved Natur og Historie bestemt., det er Aandens endelige Liv, eller den endelige Aand. Men deri ligger et Forhold til Gud, umiddelbart og forborgent; tager man Bevidstheden som Jeg, saa har det an sich tomme Jeg sit Indhold i Natur og Historie, og forholder sig efter den anden Side negativ. I dette Jeg er Naturen bleven menneskelig, kommen til sig selv. For Mennesket har Naturen umiddelbar Virkelighed. Naturen faaer en Beziehung auf sich selbst og en Richtung zu sich selbst. Denne Neigung til sig selv er Abneigung mod det, som gaaer ud over den. Det er Selbstsucht. Vil ikke forklares til Aand. Mennesket som han udgaar fra Naturen har samme Retning som Naturen, han indskrænker sig til Verden i Tid og Rum. – Det naturlige Menneske er tillige den Sandheden manglende. Den abstrakte Tænken. Saaledes betegnes Mennesket i Skriften som der selische. –

\emph{Werden für Gott.}

\emph{Identitæt.} Den guddommelige Bestemmelse, at Menneskets Werden für sich ikke er Negationen af dets Werden für Gott; thi denne Negation er netop Fordærvelsen. 1) Det er Aandens Naadeværk, at den guddommelige Kaldelse er, at det identiske Forhold an sich er tilstæde. Denne Mulighed er ikke Menneskets men Guds Værk, gratia præveniens. 2) i Menneskene er det den Fähigkeit at kunne træde i dette Forhold, hvis Mulighed er den guddommelige Naade. Dog er denne Befähigung ikke Menneskets Thun, men er atter Naaden; thi det fornuftige sædelige Væsen kan ligesaa lidet bringe sig i et Forhold til Gud som Gud i et Forhold til sig. Men er dette absolute Forhold an sich, saa er derved Muligheden givet. Denne Mulighed viser sig i Mennesket deri, at han er Fornuft og Frihed. Skjøndt alle Mennesker durch die Gnade berufen sind, saa bringen sie sich nicht selten drum. Her er Modsigelse mellem Naadens Almindelighed og Erfaring, at ikke Alle blive salige. 3) Eenheden heraf er det sande Naadevalg. Den abstrakte Forstand udretter her Intet Speculationen begriber det. Den guddommelige Raadslutning er ubetinget og dog betinget. Raadslutningen har et Forhold til Mennesket og til Indhold den guddommelige Bestemmelse om hans Salighed eller Usalighed. Ogsaa Pelagianerne indrømmede, at Frihed og Fornuft vare Naadegaver; men da Mennesket derved blev salig, saa blev det muligt ved Naaden; men Fornuft og Frihed ere blot Mulighed. Det hjælper ikke at gjøre det betingede Naadevalg gjeldende som begrundet i Forudvidenhed. Det er en aldeles ugyldig Adskillelse af Wollen og Wissen i Gud. Guds Forudbestemmelse er ikke afhængig af hans Forudvidenhed, de ere eet og deri har Calvin Ret. – Den guddommelige Raadslutning er vel ikke afhængig af Tid og Endelighed Denne Raadslutning er villende, vælgende, og det er Naadevalget; men Naadevalgets Sandhed ligger i, at Alle ere valgte. Det egentlige Naadevalg beroer paa en Adskillelse mellem Gode og Onde, og den er ikke i Guds evige Raadslutning. – Modsigelsen hæves, Naadens Almindelighed og Naadens Ubetingethed maa beholdes som begrundede i Gud fra al Abstraktion fra Forskjellen mellem Gode og Onde. est homo a deo creatus prædestinatus. Mennesket som saadan i sin Umiddelbarhed er den til hvem Naadens ubetingede Decret kommer. Den ubetingede og almindelige Naade forvirkliger sig og hæver sig positivt. Selv Particularismen kan ikke nægte, at Mennesket maa holde fast paa Naaden (stille halten), dette er dog vel Thun. Idet det Gode er det an sich er det givet, idet det træder til Mennesket er det modtaget. I Forhold til dets an sich er Mennesket ikke virksom; men i Modtagelsen deraf er han activ. Det ene Forhold negeres ved det andet og saaledes det modsigende Sande; og med Hensyn til Mennesket er det den mod det Sande og Gode rettede Drift, for tillige at komme ud af Modsigelsen. –. Mennesket antager Naaden, idet han anerkjender sin Bestemmelse, han antager den ikke, idet han negerer sin Bestemmelse. Fra Guds Side er Raadslutning ubetinget, fra Menneskets Side er Udvælgelsen den betingede, uden at derfor Naaden ophører at være den ubetingede. At Naaden er en betinget har sin Grund i den selv som den ubetingede. Den abstrakte Modsætning mellem guddommelig og menneskelige Frihed er det der forhindrer at man ikke kan løse Modsætning, som om den guddommelige Frihed var en Hemmung for den menneskelige; den guddommelige er netop derved den ubetingede, at den ikke hemmer Friheden uden for sig, men sætter og begrunder den. Den moralske Friheds Sandhed er den ubetingede Frihed.

\emph{Omvendelsen.}

\emph{Skriftens Lære.} – \emph{Kirkenslære.} i det 16d Aarhundrede bragtes enkelte ny Bestemmelser til, som ikke altid stemme med Skriften. alle Aandens Naadevirkninger til Omvendelse vare middelbare og knyttede til Skriften.. Troen \emph{sola} justificans. 1) interna, ikke den historiske Tro blot. 2) viva. 3) salvifica.

\emph{Dogmets Begreb.} Kaldelsen er alle Naadevirkningers Begyndelse. Den er den menneskelige Bestemmelse tænkt uden Opnaaelsen af Maalet. Dens Maal er Menneskets Omvendelse. Muligheden er indeholdt i Naaden, som er aabenbaret i Christo, og deri at Mennesket midt i den almindelige Fordærvelse dog er forbleven det personlige Væsen. Virkeligheden kommer frem ved at lade sig omvende og ved at omvende sig, og Omvendelsen skeer naar Kaldelsen finder en den adæquat Tilstand hos Mennesket, og begges Virksomhed falder sammen. At Mennesket omvender sig \emph{til} Gud \emph{uden} Gud er en Modsigelse. Men Gud omvender heller Ingen mod sin Villie. Denne Bevægelse maa mediere sig gjennem sine enkelte Momenter. Mediationen er Oplysningen, Helliggjørelsen og Begnadigung. Disse Udtryk antyde allerede selv begge Momenter.

\emph{Die Erleuchtung.} 1) Wahrheit als Wesenheit Gottes. Gud er Lyset og boer i et utilgjængeligt Lys 2) Menneskelige Væsens Forhold til guddommelige Sandhed. Den menneskelige Forstand er af Naturen Dunkel og Mørke 3) Den menneskelige Forstands Opklarelse ved den gud. Sandhed.

1) Die Wahrheit Gottes. Det ligger i at Gud er Aand, som Aand er han villende, som villende er han tænkende. Tænken i Identitæt med Væsen er Sandhed. Aandens Væsen som er Gud er altsaa Sandhed. – Jegheden har al Sandhed til Princip, dog gaaer Fornuften feil i sine Slutninger. Subjektet er ved Spørgsmaalet om Sandheden nødt til at abstrahere fra sig selv. – Guds Forstand er en infinitiv. Men Gud er som Sandhed ogsaa Kjærlighed og som saadan Meddelelse. – Læren om Gud som det utilgængelige Lys maa ikke eensidig fixeres. Dets Aabenbarelse er Sønnen. Ogsaa Aanden bliver kaldet et Lys.

2) Der forudsættes Mørke i den menneskelige Aand. – Subjektet bliver sig bevidst at det mangler Sandheden, derved er Mangelen hævet ud over sig selv og er Trang. Det indseer han, naar han reflekterer paa sig selv og paa Verden. Mennesket har saaledes den absolute Sandheds Idee som Ahnelse og Postulat, men er dog allerede som saadant ude over sig selv som sinnlich og verständig. I Mangelen er han Mørke; men i Bevidstheden herom eller i Trangen er han det ikke mere. Han har den Bevidsthed, at han behøver en høiere Hjælp. Mennesket er Sandheden fähig. an sich er denne Fähigkeit den guddommelige Kaldelse. Denne Tanke om det guddommeligt Sande er ved det guddommelige Sande selv i ham, er ikke mere den blot menneskelige Tanke.

3) Menneskets Oplysning ved den guddommelige Sandhed. Denne Sandhed er i Mennesket derved, at Mennesket er i Gud., og det i Gud værende Menneske er Gud selv.

\emph{ Helliggjørelse }

1.) Gott als die Heiligkeit. Alle Religioner have Forestilling om det Hellige, men blot udvortes, som man kan see af Mythologierne. – Denne Tanke viser sig i den rene Abstraktion i den rene Tanke om ham. Den jødiske Religion danner Modsætningen til Hedenskabet. Især kommer Retfærdigheden tilsyne; Kjærligheden mangler. – I Christendommen er Helligheden ikke seet fra Lovens Standpunkt.

2) Menneskets Forhold til den guddommelige Hellighed. Det Hellige som det Uskyldige forestilles i Skriften som det Rene. Efter sin Natur er Mennesket den ikke hellige, efter sin Bevidsthed den unheilige. – Mennesket har Trang til Helliggjørelse. Dette er den bevidste Trang, og Helliggjørelsens første Rørelse

3) Menneskets Helliggjørelse ved Gud. Er ligesaa meget guddommelig som menneskelig Akt. – Her farer den nyere Theologie vild især siden Kant. Aandens Andeel er aldeles udvortes – det Hellige er ikke det Naturlige, thi i Naturen forekommer den ikke; den er det Historiske og er deels Hellighed, deels Helliggjørelse. Det viser sig i de Enkeltes Liv. De af Christo underrettede Apostle gik ud i Verden, udbredte Evangeliet i den. Til Fromhed hører Begeistring. I denne er Neid og Gleichgültichkeit ikke. Deraf opstaaer Nødvendigheden af at sætte Alt i at fornægte sig selv, og bringe det subjektiv Hellige til Objektivitæt. Martyriet. I Menigheden træder den Enkelte tilbage og er kun i Samfund med Andre. De Helliges Samfund. Det viser sig i Valget af Midler til ethvert Formaal. Der Zweck ist rein geistig so auch die Mittel. Hat kein Mittel, das ein nur sittliches wäre, und auf der Sinnlichkeit abzweckte. – ihr Gottesdienst. Der öffentliche, ist nicht eine Zwangsanstalt.. – Die Sittlichkeit, des Einzelnen in seinem Verhaltniße zu sich selbst, und im Verhaltniß zu Anderen, die Sittlichkeit des Volks. das Gesetz wird heilig gehalten, ist die Anerkenntniß seiner heiligen Ursprung.

3) \emph{Die Begnadigung} \emph{.}

I det gamle Testamente er nærmest Retfærdiggjørelsen i judiciel Forstand – Mennesket kan ikke yde Andet end Troen. Det andet Moment er Fornyelsen. Catholikerne adskille ikke begge Dele. protestantismen fastholder dog ogsaa et nyt Liv som den nødvendige Følge.

\emph{III. Friheden.}

Den bibelske Lære. subjektiv – objektiv – universell.
\end{quote}
\dcite{Not10:9  (1841; SKS 19, 294)}
% Not10:9  (1841; SKS 19, 294)

\begin{quote}
Avling frembringer det Eensartede, Skabelse det Modsatte (ex nihilo)

Sabellius.                            Arius.

Nicæa.

Consubstantialt kan kun være hvad der adskiller sig.

Dogmets Begreb.

1) \emph{Das sichoffenbarseyn Gottes} \emph{.}

Guds Væren som Tænken, Væren alene giver ikke Begreb. Spinoza selv erklærer Guds Væren som en saadan, das an ihm das Denken hat, altsaa Prædikat, Accidents. Der er sat en Tohed, hvortil Gud sætter sig selv.

2) \emph{das Grundseyn Gottes} \emph{– causa sui. } Rationalismen indseer vel, at Gud er Grund til Alt, dette Alt er Verden. Spinoza lader Gud ikke bestemme sig hverken til Grund for Verden eller for sig. Grunden er kun den, hvoraf noget fremgaaer. Denne Bestemmelse af Grundseyn betegner negativt at Gud er sin egen Grund, affirmativt existentia dei som exitus. Denne Existents er i Forhold til Verden Præexistents, Guds ewige Thun – Jacob Bohme: Gott hat aus Einem zwei gemacht und ist doch ein geblieben. – λογος har 2 Sider 1) Gud er Grund ratio 2) Fornuft ratio. –

λογος har auf ewige Weise sin Grund i Gud, og gaaer ligeledes auf ewiger Weise aus. Saaledes er Gud den absolute Intelligents, Fornuft. Bliver man ved Substantsen, saa staaer Guds Offenbarseyn i det Andet uden denne Selvadskillelse ikke til at indsee.

3.) \emph{Gud: Fader og Søn} \emph{.}

Hvad er Ordet uden Tanke, og Tanke uden Ordet. Saaledes er Gud den gud. Fornuft eller λογος. – Forestilling om Gud som Søn er ikke blot skjøn, men umiddelbar i Begrebet. Troen taler i Forestillinger, men det store er at den har Begrebet i sig. Gud er αγενητος som Grund, γενετος som den der gaaer ud. Den negative Forestilling, at den ene ikke er den anden, hænger i de jordiske Forestillinger. Men man maa fastholde at Gud Fader er evig Fader, og Gud Søn evig Søn. Guds Avling maa heller ei tænkes som en Liden, eller som havende Ophør. Faderens Daseyn er ikke Sønnens Existents og omvendt. Det er to Personer (ikke Individualitæter). Person er væsentlig Bevidsthed. – Det at være Guds Søn er endnu ikke at være Menneske, end ikke Menneske ϰατ' εξοχην. Foreningen af de to Naturer forudsætter Adskillelsen, og som Guds Søn er han ene at betragte i det guddommelige Væsen. Forholdet mellem Fader og Søn er absolut Forhold. Sønnen erkjender sig i Faderen og erkjender dermed Faderen, ligesom Faderen. Gud er ikke ewiger Weise Sønnen, men er auf ewiger Weise ikke uden Sønnen. Gud er derfor ikke Fader blot i den Forstand, at han er alle Menneskers Fader. Det er den rationalistiske Polytheisme, der overseer Guds Forhold til sig selv. Kun i Sønnen kan man komme til Faderen.

Guds middelbare Aabenbarelse

1. \emph{Skabelse}

a) Skabelsen selv b) Mennesket i Guds Billede c) Tabet af dette Billede eller det Ondes Oprindelse.

a) Skabelsen selv.

α) \emph{Den bibelske Lære} \emph{.} – Rationalismen siger at ved Udtrykket δι' αυτου, εν αυτῳ om Sønnen betegnes han som Werkzeug. – Supranaturalismen angriber ikke Udtrykkene men begriber dem ikke. – Skabelse ex nihilo. – i Modsætning til Avling der er Meddelelse af Væsen.

β) \emph{Den kirkelige Lære} \emph{;} den udelukker 1) alle Forestillinger, hvori Gud ikke adskilles fra Verden, men opfattes som Einerleiheit: Emanations Systemer. – Arianisme Gud har ved Avling af Sønnen indavlet sit Væsen i Verden. Pantheisme. 2) De Forestillinger, der adskille Gud fra Verden: Materialisme, Hylosoisme; Pelagianisme. Den kirkelige Lære forener sig med den bibelske, der ligesaa meget forudsætter det indre og immanente Forhold som Adskillelsen, ligesaa meget Guds Immanents som Transcendents. I de smalkaldiske Artikler er tilføiet til Troes-Artiklens Indhold, at den \emph{treenige} Gud har skabt Verden.

4) \emph{Dogmets Begreb}.

Den udvortes Verden er Indbegreb af Alt hvad da ist., die erscheinende Welt og Gjenstand for sandselig Iagttagelse, hvilken bestandig med det Einzelnes at gjøre; fører Tanken det Einzelne tilbage til det Almene kan det blive staaende ved Atomet eller ved Materien som Væsen; den fixerer Verden ligeoverfor Bevidstheden. Her er altsaa ikke Tale om Skabelse, men Atomernes concursus, for Religionen har den kun som Resultat Trøstesløshed. Aanden svæver ikke over Vandet, men Vandet beholder Overhaand. Natur (nasci) betegner blot Produktivitæt, Tanken er det supranaturam. – I Idealismen er kun den tænkte Verden den sande, Atomet er den en Tanketing, Materie er Abstraktion, ikke i og for sig. I begge Systemer er Væren og Tænken ikke blot quantitativ overveiende, men hinanden modsatte, og derved berøvede Sandhed. – Identitæten lærer at Væren er Tænken, saaledes, at Væren er det som Væren i sig reflekterte Tænken. Den Aand er Verdens Princip saa vel som Væren som som Tænken. Man maa ikke blive staaende ved Skabelse som blot Skabelse af Natur. Naturen har sin Sandhed i Aanden.

Men er Selvbevidstheden Naturens Forklarelse, saa spørges, hvad er Selvbevidsthedens Sandhed, det er den anden Verden, hvis Rige er Menigheden. Menneskets naturlige Fødsel har blot Betydning af Forudsætning, som den af Gud skabte bliver han først erkjendt i Troen og som gjenfødt.

Den guddommelige Skabelse. Gud staaer igjennem Verden i Forhold til sig. \emph{Dette anerkjendes ikke:} i \emph{Pantheismen} \emph{,} Gud er ikke Skaber thi han er Verden selv.: \emph{i dens kosmologiske} Betragtning, Verden er Virkning af Aarsag og den Aarsag er Gud; her er et Spring fra det Endelige (Verden) til det Uendelige, og det skeer ved en Slutning, man bliver staaende ved Verden som Erscheinung, Gud bliver udenfor Verden. Dette Causalitætsforhold er lige Modsætning til Pantheismens Substantialitæt. Herpaa beroer Pelagianismen idet det regnes med til Guds Naade at have lagt saadanne Kræfter i Verden, hvorved den kan blive fuldkommen. \emph{Den teleologiske.} her er det Zweck og Mittel hiint i Tanken, dette det, hvorved den realiserer sig; men den falder tilbage i den kosmologiske, idet denne Zweckbestimmung bliver udvortes ikke bestemt efter Væsen og an sich, den ender subjektivt; den veed ikke om det bestemte Zweck er an sich og objektiv er tilstæde.

Den christelige Lære indeholder 3 Bestemmelser.

1) Skabelsen ved Guds Søn. alle opera dei ad extra sættes i den samme Virksomhed af Fader Søn og Aand. \emph{Faderens} det vil sige Skabelsen grundet i Guds Almagt men det kunde jo blive ved en Kunnen. Et Guds Forhold til sig selv negativ ikke det Skabte til den Skabende, og gaaer ud over Substantialitæts Forholdet (der da egentlig blot er den abstrakte Identitæt); thi Gud er forskjellig fra det Skabte som den der har Grunden til sig i sig. Tanken om Faderen viser hen til Sønnen og det er Guds Forhold til sig uden alt Forhold til Verden. Deri at Gud er aus sich erkjendes at Verden er af Gud, thi idet Gud er aus sich som Søn har han bestemt sig som Grund. Gud er ikke Gnostikernes αβυϑος eller Urgrund. I Symbolet gaar ogsaa Benævnelsen Fader før Skaberen, var han blot Fader som Skaber, maatte dette Prædikat komme først. \emph{Sønnen} med Sønnens Avling er Verdens Skabelse sat; ikke som om Sønnen er Verden og ellers Intet εν αυτῳ εϰϑιστεν τα παντα Col: 1, 16. Men i Sønnens Avling er Verdens Mulighed sat, om end ikke dens Nødvendighed. Dens Virkelighed har den ved at være skabt af Sønnen. I Sønnen elsker Faderen Verden og deri ligger at Sønnen kan erklære den menneskelige Natur for sin.

\emph{Ved Aanden:} deels forsaavidt Aanden er Eenheden af Faderens og Sønnens Villie og Væsen, deels forsaavidt den er Overgangen fra den indre Virksomhed til Verden.

\emph{Skabelsen af Intet.}

\emph{negativt:} intet Chaos. blot Formation. Tohu vabohu er oftere bleven forklaret som et saadant forhaandenværende Materiale. han har ikke skabt den af Noget. durch kein Ding ist Gott bedingt. Das abstrakte Seyn als identisch mit Nichts. Gnosticismen lader nu Alt fremgaae af Gud selv, der seinen unendlichen Inhalt heraussetzt. Denne Gnosticisme der er uden Rückkehr er Modsætningen til Spinozismen, hvor alt uden Oprindelse synker i Gud. Den Lære, at Verden er af Gud lader sig godt forene med den om Intet. Disse to Sætninger ere ikke identiske ikke heller hinanden udelukkende hinanden, men bestaae med hinanden. \emph{positivt } Verden er ikke begrebet som Produkt, men som Creatur. I dens Nichtseyn trænger den til Tænken. Den endelige Aand fremkalder Verden af sit Intet og lader den fremtræde for Bevidstheden, Gud skaber og denne Nichtss Seyn er Verden, saaledes viser sig bestandig det Nichts, hvoraf det er fremgaaet og som tillige er indgaaet i Verden, da det er dens Nichtigkeit. Gud sagde: det vorde. Der læres derfor heller ei, at Mennesket er skabt af Intet, saaledes som Naturen, men at det er skabt af Guds Væsen.

\emph{Verdens Nødvendighed.}

ikke i udvortes Forstand; eller som om han trængte til den, som om Gud var in der Noth (Nothwendigkeit); eller Gud i Realisationen af Verden først realiserede sit Begreb, som om Gud først havde haft Verden først i Ideen, som om denne igjen i Gud var noget afmægtig, eller den fører til Pantheismen, Identitæt af Gud og Verden, Expansion. – Verden er skabt i og med Tiden, om end Skabelsen er Realisation af en Raadslutning fra evige Tider. Har Gud fra Evighed elsket Verden, saa maa han have skabt den fra Evighed. Verden er altsaa fra Evighed; den er egentlig sat i Avlingen af Sønnen, og med Sønnen følger først Skabelsen, og denne Skabelse er ikke fra Evighed; men er skabt i Tiden, og saaledes bør begge Udtryk fastholdes. I Sønnen er det Verden i sin Uendelighed her i sin Endelighed. – Friheden i Verdens Skabelse er bleven gjort gjeldende af Strauss i hans Rechtsphilosophie, og i en Afhandling af Julius Müller i Studien und Critiken 1835. – Identitæten af Frihed og Nødvendighed er udtrykt i Skriftens Lære at Gud har skabt Verden af Intet. – Verden er ikke blot Daseyn men og Bewußtseyn, men Betingelsen for denne er hiin og dette er Naturens Betydning. I Verden har Gud Historie. – Verdens Nødvendighed overhovedet indeholder eo ipso i sig, at det er den bedste Verden, ikke ved menneskelige Kraft, men derved, at den er skabt af Gud.

\emph{Zweck der Weltschöpfung.}

I det teleologiske Beviis træder Zweck og Mittel ud fra hinanden. Hvilket er Verdens Zweck i Forhold til hvilken Alt er Middel. I den endelige Betragtning bliver Zweck og Mittel relativ. Naturen har vel Zweck, men formaaer ikke at foresætte sig det. I den teleolog. Betragtning er Gud nødvendig sat som »Geist«. Spørges der om et Zweck for das Weltall, saa er det ikke mere sat som Natur, men som Aand.

\emph{Guds Zweck med Verden. Gud kan ikke have andet til Zweck end sig selv,} Verden er ikke skabt ihretvillen, men Gotteswillen. Naar den i sandhed er für Gott saa er den heller ei für sich. Adskiller den sit für sich, saa er den ikke mere i Sandhed für sich, men gaaer sin Nichtsseyn imøde. – Gud har altsaa ingen Trang til Verden, som et Moment i hans Fuldkommenhed. Men saa kunde den jo være bleven uskabt? Nei. Feilen er at man ser Gott blot som den sich genug seiende, ikke den sich genugthuende. – Naar Verden ikke har andet Zweck, saa er det satte Zweck ogsaa det opnaaede, og har ingen Intervall mellem Tænken og Handlen. Saaledes er Guds Zweck uendeligt Zweck. Dermed lader sig godt forene hermed; thi dens Opstaaen i Tiden er dens Endelighed, dens werden und geworden seyn. Denne Forening af det Timelige og det Evige, viser igjen at Gud har intet andet Zweck med Verden end sig selv. Den guddommelige Kjærlighed og Hellighed realiserer sig i den menneskelige Fornuft og Frihed. Dieß vermittelt sig dette gjennem Verden, som altsaa er Middel. Men ogsaa heri har Gud sig selv til Zweck. Verdens Zweck er det samme som Guds Zweck med Verden. Naar Verdens Zweck er opnaaet saa er det i Guds Forherligelse.

\emph{ Vom göttlichen Ebenbilde }

\emph{Den bibelske Lære; den kirkelige Lære; den skolastiske. – Dogmets Begreb.}

Som Skabelsens Middelpunkt og som Nødvendighed i den dialektiske Bevægelse dertil. 1) Menneskets Skabningens Middelpunkt. Det synes egoistisk. \emph{Den materielle} Verden med Solen, \emph{den intellectuelle } med Englene. Expeditioner til Stjernerne. Man har begaaet den Feil at bedømme og betragte disse Himmellegemer ganske i Analogie med Jorden, istedetfor at betragte dem i og for sig, deres Eiendommelighed, der ofte for Exempel ved Sol og Maane forbyder at antage en Befolkning. Det Store i Guds Frembringelse ligger ikke i det Extensive, som om Guds Billedet i Verden successivt udviklede sig paa forskjellige Kloder, men i den Intensitæt der er Aand. For Mennesket er det Substantielle det Accidentelle. – \emph{Den intellectuelle.} Englene. – Det Speculative deri; Dogmatiken har den Opgave at udfinde Grunden til dens Antagelse. Om man her end erhvervede Sandsynlighed, saa er det dog ikke nok til et Dogma, men vel nok til ikke at afskaffe dem i det religieuse Foredrag. – Den intellectuelle Verden kan ikke gaae over i Erscheinung, den Materielle bliver altid i Erscheinung, det bliver nu det høieste Væsen, den forener begge Dele det vil sige Mennesket

\emph{Dette Menneskets Fortrin for Alt viser sig ogsaa ved Betragtning af Naturen.} Alt viser hen paa en Uranfänglikeit af Mennesket, der kun lader sig bevare ved Avling. – Nedstammelse fra et Par. Auktoktonlæren begynder med en større Usandsynlighed end Bibelen. Strauss antager den imidlertid.

\_\_\_\_\_\_\_\_\_\_.

Et Guds Billede kan ikke være i udvortes og endelig Forstand. I dets Sandhed og Uendelighed beroer det paa Identitæt; uden dette kan han ikke være Ebenbild; men dette Ebenbild er Sønnen, der er i absolut Identitæt med Gud, Homoousie, og ikke blot Homoiousie; liig kun i Betydning af Identitæt. Forholdet er et umiddelbar, umiddelbar Lighed, Urbild og Nachbild er ikke wesentlich traadte ud fra hinanden. Er Sønnen Billede saa er den anden det ogsaa nemlig Vorbild für das Nachbild. – ved Menneskets Ebenbild kommer Differentsen i Betragtning. (Han er ikke selv Gud, ligesom Sønnen.) Pantheismen er udelukket; thi der er et Forhold det vil sige Forskjel og Sammenhæng. Her er Homoiousie. Ved Christus bliver det tabte Billede restaureret, men altsaa er det alle angeschaffen. – Differents og Identitæt gaaer sammen, mellem Guds evige Søn i sin Lighed med Gud og hans Menneskevordelse i Christo ligger Adam som Menneske I den anden Adam er Forskjellen hævet, den var sat i den første Adam. Idet Syndens Mulighed er sat er Menneskets Historie først sat som Begyndelse. – Hvad indeholdes nu i det første Menneskes Lighed med Gud. han er forskjellig fra Alt hvad der er ringere end han og fra Gud som er over ham. – Gemyt er Eenheden af Tænken og Villen som Mulighed. Derved er han forskjellig fra Dyret og hele Naturen – i Modsætning til Gud er hans Aand den endelige. positiv i ham er den im Uendliche endliche Geist ist. saa at Mennesket kan optage det der er hans Grund det vil sige Gud i sig. I Fornuften ligger Guds-Ligheden. Dernæst en Delagtighed i den guddl. Natur, (ikke ein Theilseyn Gottes numerisk forstaaet,) en Handlen og Villen, en Magt til at sætte det dem Grund entsprechende, idet Gud i sin Frihed sætter det som af Gud er sat i ham som hans Væsen er han fri. –

Det er Menneskets Bestemmelse der skal realiseres, og den reale Mulighed at være det fornuftige og fri Væsen. Denne Eenhed af hiin Bestemelse og denne Mulighed er den i Gemyttet forenede Guddommelighed. Det er Uskyld der svarer til Guds Hellighed. Men i Unschuld er der en Beziehung auf Schuld i det mindste ϰατα δυναμιν, i Hellighed er det ikke saa.

\emph{Det dialektiske i denne Lære.} 1) er indeholdt hvad Mennesket er ved Gud med al Abstraktion fra, hvad han er ved sig. – Men for at Ligheden skal tilveiebringes maa han være adskilt fra ham. Mennesket er bestemt til i sig at optage alle de evige Ideer, Gud haver i Skabelsen aabenbaret; han maa reflektere paa hvad Gud har gjort ham til. Hvad Mennesket er lader sig nærmest først erkjende af hvad Gud er, Anthropologie forudsætter Theologie, er Gud Skaber saa svarer hertil Menneskets uendelige Bestemmelse. 2) Den indeholder hvad Mennesket er an sich med Abstraktion fra Alt hvad han er für sich. Læren om Uskyld indeholder at han an sich er fornuftig og fri. 3) Det guddommeligt Ebenbilds Gegenwart ist Erinnerung. Det er det, der ligger i Læren om at Guds Billede er tabt og dog tilstæde. Epiphanius forkaster Origenes Lære, at Gudsbillede er tabt, og lærer efter Skriften, at det er bleven i alle Mennesker. Andre mente, at det vel var tabt, men scintillulæ blevne tilbage. Rationalisterne gjør dette Guds Billede til en Abstraktion, i den ene eller den anden Form, som et Ideal, hvorefter der skal stræbes. Først naar det fastholdes som Erindring, først derved bliver det muligt, at Mennesket kan opfatte sin Tilstand ikke blot som ikke-hellig, men som syndig. Som Erindringen er det eine Thatsache des Bewußtseyns. Det er det, der gaaer igjennem alle Folkeslags Bevidsthed. Overeensstemmelse er her ikke udvortes, men indvortes, det er det samme Menneske, den samme Idee, kun rationelt uddannet. – Schelling meente, at den første Tilstand war en høi Grad af Cultur, i hans »academische Studium – Hegel Relig: Ph. 1ste Bind 306. Det fuldkomne Menneske er tillige det ufuldkomneste, thi Alt ligger indesluttet i det; om det end har det Fortrin, at det endnu ikke er forvirret af Culturens Mangfoldighed. Rabbinerne troede, at de første Mennesker havde været i Besiddelse af mathematiske Kundskaber, Astronomie, Lærdom.

\emph{Om det Ondes Oprindelse.}

Den bibelske Lære. Det Onde sættes i Legemet. De onde Tanker blive henførte til Djævelen. Den stiller den herskende Synd i Forbindelse med den første Synd. Aristoteles kalder det Onde συγγενες. Plato siger, at Børnene ikke ere gode φυσει. ogsaa Cicero i Tusculanerne 3, 1. – Udtrykket Arvesynd forekomer vel ikke i Bibelen, men Tanken er der. – Den kirkelige – den symbolske – den supranaturalistiske – den rationalistiske.

\emph{Den symbolske} knytter sig især til AugustinSpeculationer. defectus et carentia justitiæ originalis, concupiscentia, lærer 1) vitiositas et corruptio er universalis. 2) naturalis, congenita, insita. Eph: 2, 3. Ψ 58, 4. 3) dog er denne vitiositas ikke substantialis., imod Flaccius: peccatum originale ipsam esse hominis substantiam. – accidens est quod non per se substitit, sed in aliqua substantia est et ab ea discerni potest.

\emph{Med Hensyn til Forplantningen } er denne vitiositas 1) physica. derved gaaer den tilbage til det første Menneske 2) moralis Enhvers Indvilligelse i den nedarvede Hang. Alle virkelige Synder ere Virkeliggjørelse af Arvesynden, der er Disposition til Alt. 3) per se damnabilis. nærmest den naturlige Straf, at Synden straffer sig selv for Exempel Synd fødes af Synd. Knapp vil m: Hensyn til Døden skille mellem Straf og Følge; men det er en Adskillelse af de øvrige Mennesker fra det første, men det er ikke sandt, der er, som Tertullian siger, kun eet Menneske

\emph{Med Hensyn til det Ondes Oprindelse } læres, at den actuelle og oprindelige Synd har et indre Nexus med den første Synd og bliver Alle tilregnede. Ved den første Synd forstaaer den symbolske Lære uden Omsvøb den første Ulydighed. peccatum originis non esse imputationem solam sed calliginem et depravationem in natura.

\emph{Den supranaturalistiske.} Den tager Fortællingen i Genesis som en slet og ret Historie, hvori Adam træder op som en Privatmand. Gerhard kalder Adam pater totius humani generis. Calov: fons caput seminarium totius generis humani Den nyere Supranat: sætter ham udenfor den menneskelige Natur overhovedet. Den sætter Adams Synd uden Forbindelse med Menneskeslægtens. Den lægger Eftertrykket paa den forbudne Frugt, der har været giftig, og hvorved Syndighed og Død kom ind i Verden. Eftertrykket ligger paa det Onde og det Forbudne.

\emph{Begrebet af det Ondes Oprindelse.} Man maa gaae ud fra den Eenhed mellem Gud og Menneske det vil sige Uskyld, der Ausgangspunkt der Schuld ist die Unschuld. Adskillelsen af Gud og Menneske er nødv; men som sat er den Eenhedens Negative. Gud har derfor ikke med Magt undertrykt det. \emph{Det Ondes Mulighed; Virkelighed; Skyld.} Dets Mulighed. Det Ondes Daseyn an sich er bleven for Mennesket, og han ond. Det hat \emph{sich zugetragen;} hans Forfører var en ond Aand, der ogsaa oprindelig var en god Aand, \emph{es ist geschehen,} daß er von Gott abfiel durch Hochmuth; men hvorledes og hvad der forførte ham, hvorledes Mennesket antog hans Fristelse, det forklares ei. Det Onde er vel saaledes der, men hvorledes Mennesket \tlg{tilegnede} sig det, bliver ubegrebet.

Man kunde nu tage Bevidstheden til Standpunkt, formedelst hans Bevidsthed adskilte han Ret og Uret og ved denne Adskillelse blev han istand til at vælge. Denne Anskuelse forudsætter at Godt og Ondt allerede er; den siger ikke, hvori det Slette ligger, eller hvorfra det kom, eller dets Begyndelse. For at begribe denne maa man gaa ud herover. Verden som Natur er bevidstløs og kun daseiende og mangler Fornuft og Frihed. Naturen som realiserende denne Mulighed bliver Bewußtseyn, og er Welt og Natur als Ichheit.. Det Naturlige naar det har naaet Bevidstheden, er ikke blot Tænkende, men sig Tænkende. her har Verden først fuldkommen sich erfaßt. Men begynder i det Fjerne det Onde. Die bewußtseiende Natur ist sich nur ihres Wollen, og som sich Denken er den og sich Wollen. Al Bestræben gaar ud paa at føre Alt hen til sig; saaledes er den selbstsüchtige Natur geworden und zwar als Freiheit, derved er den bleven afhængig af sig selv og har tabt sin Frihed. Trangen er at Naturen bliver befriet fra Ichheit og Selbstsucht. In diesem Ichwerden der Natur, und Natur-Werden des Ichs ist der Egoismus, und darin der Verlust des göttlichen Ebenbildes.

Denne Bevægelse er Bevidsthedens nødvendige Bestemmelse. Tænkningen af Forskjellen mellem Godt og Ondt er ikke det Onde, men er Betingelsen derfor. Det da at det Gode og Ondes Unterschiedslosigkeit gaaer under, er endnu ikke et Tab men Gevinst; før den Tid er her alt af Gud ikke noget ved sig selv. Overgangen til det Onde danner Forestillingen om det Onde, som det der an sich er i Daseyn, dies ist der Gedanke des bösens Gedankens. Djævelens Hypostasering er ikke det man skal blive staaende ved i Dogmatiken, men hvad er Indholdet deri? Det er, at Mennesket gjør sig den subjektive Tanke om det Onde objektiv, og derved indrømmer det en Magt over sig. Naar nu der Gedanke des bösen Gedankens slaar over i den onde Tanke, saa har Djævelen vundet. Det er altsaa ikke det Onde, at tænke Djævelen. – Forskjellen mellem Gud og Menneske bliver nu til Widerspruch og Gegensatz.

\emph{Det Ondes Virkelighed.,} eller wovon es seine Form hat. Den med Tænkningen af det Onde sig identificerende Wollen. Her først er Mennesket i Modsigelse med sig selv. Det bliver virkeligt idet Frihed og Nødvendighed, Frihed og Lov skilles ad og sætter sig mod hinanden, og derved den som Ufrihed sig sættende Frihed. Bevidstheden har nu det til sit Indhold at der er den Loven modstridende Frihed. – Det Onde er derfor ikke virkelig, men i Vorden, i Entstehen og Vergehen, det Væsen der blot er Væsen og som saadant U-Væsen, der ikke gaaer videre til Begrebet; det forvandler Alt Væsentligt, alt virkeligt Væsentligt, til noget Timeligt og Rummeligt og derved til Uvæsen. Den Onde lever væsentlig i Tid og Rum; istedetfor at Tid og Rum blot ere Momenter. Løgn. Det er i sig ikke noget Virkeligt, det er an sich das Nichtige, das Negative des Wirklichen, det er, det er et positivt, idet det beviser sig paa det Gode. Det Onde er en Position i Negationen. – Er Mennesket af Naturen god eller ond? Det er Forskjel mellem an sich og af Naturen. an sich er han god, men ikke som den der kommer fra Naturen, men fra Aanden, fra det Uendelige, det er hans Uskyld; det er ein Gutseyn, der ogsaa blot er an sich. Det Naturlige er Modsætningen af Aanden. Derfor er det rigtigt at sige: nicht gut ist der Mensch von Natur. Det er en Modsigelse at sige, at han er god af Natur eller ved Natur, det er alene muligt ved Aanden. Hjertets Umiddelbarhed er netop det, der maa forsages. Derfor Gjenfødelser nødvendige. Die Verselbstigung beginnt mit seiner Entselbstigung. Daaben: –––– \emph{Die Schuld.} I Modsætningen mellem Frihed og Nødvendighed er Overgangen fra Uskyld til Skyld sat, saaledes at dermed egentlig ogsaa er sat en Fortjeneste. Den første Overgang, confer det Foregaaende, gjør Mennesket blot til den ikke hellige; den anden Overgang, lader ham ophøre at være blot den nicht-heilige, og gjør ham til den unheilige. Han blev forført, saaledes synes han at være uden Skyld; paa den anden Side han lod sig forføre og har altsaa Skyld. 1) \emph{Mennesket lader sig forføre} \emph{,} dette sker in der Bewegung der Wahrnemung in Bewußtseyn. Das Wahrgenommenes ist nicht nur ein Einzelnes, sondern das Allgemeines, als Wahrg. ist es das Sinnliche. Das wo durch das Sinnliche im Denken eingebildet wird ist die Einbildungskraft. Das Wirkliche ist nicht das Sinnliche sondern das freie. For Mennesket i Erscheinungs-Verden, da han ikke kan wahrnehmen det Aandelige, vender han sig bort fra dette til die Sinnen-Welt. Det er ikke blot en Fühlen, Empfinden, men er Gelüst, da han tillige er den villende. I Mennesket kommer det Dyriske ogsaa til Villen, i Dyret er det blot Begjær. Slangen maa derfor \emph{tale} for at vække Lysten., thi ved denne Talen bringes Mennesket til at imaginere sig det. – Fornuft og Frihed blive nu nedsatte til Tjenere for den sandselige Lyst. Fornuft og Frihed blive altsaa ikke virkelige, forblive i hiin Mulighed, og ere derved degraderede. Denne Fornuftens og Friheds Underordnelse under Lyst og Wahrnehmung er Syndefaldet, Bevidstheden derom er Skyld. Ikke altsaa deri, at han er et sandseligt Væsen, eller lever i en sandselig Verden er han syndig.

\emph{Mennesket bliver forført.} Die Schuld ist die Unschuld. Idet han lader sig forføre er han skyldig activ, men da han bliver forført, forholder sig passiv, saa er det ikke absolut Tab af Fornuft og Frihed, er hat eingebüßt an dem Geiste, aber nicht den Geist selbst verloren. Han har tabt dens Reenhed, men ikke tabt den selv.

\emph{Das Wissen der Schuld ist das Gewissen,} i Identitæt af at lade sig forføre og være forført. Denne Identitæt er medieret gjennem Godt og Ondt, Samvittighed er kun i et Væsen, der veed Godt \emph{og} Ondt, Ondt \emph{og} Godt. Ohne Gewissen ist allein Christus und der Teufel. Gud er hellig, men deri ligger intet Forhold til Godt \emph{og} Ondt. Christus blev derfor fristet, men Grunden, hvorfor han ikke blev overvunden af Fristelsen, er ikke, fordi han var uden Samvittighed, men fordi han ligger ud derover. – Den Onde har heller ei Samvittighed.

Das Gewißen ist im Menschen das Wissen des Guten am Bösen, und des Bösen an dem Gute, og er det Punkt, hvorfra han atter kan opdage det tabte Land, og orientere sig paa ny i det Gode. I Samvittighed skeer det at den guddommelige Villie tager Skikkelse af Lov, men denne Form kan igjen afstryges. Christus er derfor Lovens Ende. Fra Samvittigheden udgaar al Omvendelse og Forbedring, Anger tilhører Samvittigheden, der indeholder en Bevidsthed af ikke selv at kunne retfærdiggjøre sig selv, denne er Bevidstheden af Skyld, det, at der dog er Udsigt til en Retfærdiggjørelse er Bevidstheden af Uskyld. Gjør man Mennesket til egentlig Ophav af det Onde, saa er al Retfærdiggjørelse umulig. Mennesket bliver derved selv gjort til Djævel.

Dette maa betragtes subjektivt det vil sige den enkelte Synd, objektiv det vil sige Arvesynd, Eenheden. \emph{subjektiv.} det er den Enkeltes egen Synd. Hans Natur er vel syndig, men han er dog tillige Aand. Flaccius Illyricus. Modsætningen hertil er at udlede det Onde af Opdragelse og Exempel; men Mangel af Opdragelse og Exempel vilde fordærve det endnu meer, og hvor blev Grunden til den forkerte Opdragelse og Exempel at søge, og hvoraf det kom, at Mennesket var saa begjerlig efter at opfatte det Onde og \tlg{tilegne} sig det. Die einzelne Sünde ist die der Person, die allgemeine ist die der Natur. \emph{Arvesynden.} De aktuelle og Arvesynden ere blot formelt forskjellige, de forholde sig som det Almene, der ikke er uden det Besondere, og det Besondere er i det Almene. Alle Mennesker derfor ere Syndere: Synden forplanter sig metaphysischer og intelligibler Weise i Naturen. Enhver bliver født som Synder, Muligheden til alt Ondt er for Haanden. Men Arvesynden er tillige den igjennem hele Livet blivende Mulighed, som end ikke udryddes med Gjenfødelse. Den ældre Dogmatik bestemte den som negativ: defectus justitiæ originalis; positiv: Hang til Synd, Lyst til det Forbudne, concupiscentia., altsaa ikke blot i Legemet, ikke blot i Uvidenhed og Uforstand, ikke blot i Skrøbelighed forraader denne Slethed sig, men der gives Laster der udføres med den største Besindighed. Man har af gode Mennesker saadanne Bekjendelser: dybt i Hjertet skjuler sig en syndig Hang, en saa afgjort Fordærvethed, at et Smiil kan sige til sine Venner, mine Venner jeg har ingen Ven, Skadefryd endog mod den bedste Ven. Enhver har en Priis, hvorfor han hengiver sig. Ogsaa om det Onde gjelder det: homo sum, nil humani a me alienum puto.

\emph{Identitæt.} Tilregnelse af den adamitiske eller første Synd. er nu Følgerne noget absolut forskjelligt fra den første Synd, saa er Døren aabnet for al Tvivl. Den gamle Dogmatik opfattede Adam blot som den Enkelte. Det er af Vigtighed ikke at modsætte det første Menneske og hans Efterkommere; men Adam er Mennesket i Almindelighed., som han faldt saa falder endnu Enhver. Allerede Kant sagde: mutato nomine de te narratur fabula. Bliver dette det substantielle Indhold af hiin Fortælling, saa er det ligegyldigt om det bør ansees for Poem, Mythe, Philosophie eller Historie. Sandheden bekræftiger sig i Enhvers Bevidsthed. Sandheden af Pauli Lære er at vi alle have syndet i Adam. Augustinus bevæger sig alene i denne Læres Begreb, nos omnes fuimus, sumus ille unus; ikke blot ved den naturlige Nedstammen, men og perconcesionem i hans Synd bliver hans Synd til vor. Enhver anden Synd staar i Sammenhæng med den, som Skriften lærer om Syndens Forplantelse. Arvesynden taaler ikke den Forskjel mellem fremmed og egen Synd. Men denne Modsætning er hævet ved det Egne, der gjør Adams Synd til sin egen, og det Andet, der gjør det til Adams Synd. Man siger, at det strider mod Friheden. Strauss »kan end ikke indsee, at den første Synd var et Misbrug af Friheden, da Friheden netop er Mulighed til Godt og Ondt, det første Forbud har ogsaa reent Charakteer af Vilkaar, det gav ingen Grund, men uden al Humanitæt, behager sig i sin Superioritæt og i sin Inferieur.« Her er Frihed og Vilkaar forvexlet; thi Vilkaaret er Godt og Ondt aldeles ligegyldigt. Den historiske Nødvendighed til at synde er ingen Tvang, Enhver bliver sig sin Kraft til Modstand bevidst.

\emph{Tilregnelsen af Adams Synd som Skyld eller die Erbschuld.} Er Adam nu en aldeles Fremmed, saa er der ingen Mening heri. Men det antager Kirken jo ikke. Saa længe det første Menneske er et andet end alle Andre, kan man endnu spørge saaledes. Paulus og Augustinus udhæve die Erbschaft, saa der ingen væsentlig Forskjel bliver mellem det første Menneske og de følgende. Ved denne Forestilling er Adams Fremmethed hævet. men dog bliver det kun en Forestilling, der er udsat for Tvivl og Indvending, som enhver Lære er det, saa længe den bliver i Forestillingens Sphære. Men der er ikke Tale om ein einzelner Vorgang mit einem Individuum, noch von der lauteren Unschuld die Rede sey, men at her blot er fremstillet hvad Mennesket er efter sit Væsen. Som almindelig er den den Enkelte uvedkommende og fremmed, men idet han gjør den til sin, er den ophørt at være almeen, og er dog i en anden Forstand den almene.

2. \emph{Opholdelse}.

\emph{Den bibelske Lære.}\emph{Kirkelæren } er egentlig ikke meget udviklet, da man mente, at den laae i Skabelsen, med mindre man vil hertil regne Skolastikernes concursus dei, en Lære som Baumgarten har optaget i sin Dogmatik: deum concurrere ad materiale non ad formale actionum; denne Lære blev nærmest udviklet med Hensyn til Menneskets Frihed.

\emph{Dogmets Begreb.} negativ, positiv, concret.

\emph{negativ.} Skabelsen er Begyndelsen, Opholdelsen Fortsættelse. Verdens Princip er det sig fortsættende. Fjerner man den Tanke, 1) at Opholdelsen betyder Verdens Evighed eller Alt Enkelts evige Vedvaren; thi da det er opstaaet og maa forgaae. Men Opholdelsen kommer frem igjen i Apokatastasis των παντων. 2) at det ikke saa vel er den endelige og rummelige Verden Gud opholder, men kun som den er i Gud; thi saa behøver den ingen Opholdelse, saa er den uendelig. Der maa altsaa tænkes paa dens Fürsichseyn. Opholdelsen ligger i, at den Fürsichseyende ikke ophører at være for Gud. Men i det Onde er der Kræfter i Bevægelse der arbeide paa dens Undergang, men Guds Opholden arbeider derimod. 3) om en passiv Opholdet-Væren ved Gud. Saaledes som det Uorganiske indeholdes og opholdes i det Organiske, er det ikke Gjenstand for Guds Opholdelse. Organismen og Mechanismen modstride hinanden. Organismens reale Mulighed er Spiren; det Organiske er som ud af sig saaledes ogsaa i sig, dets Adskillende er Glieder, hvoraf ethvert ligesaa meget bliver opholdt ved sig selv som ved det Andet.

\emph{positivt.} 1.) Substantialitæts Forholdet. Det er da egentlig ikke noget Forhold, hvortil altid hører 2. Natur-Udvikling falder sammen med Guds Almagt. Det er Spinozas Pantheisme. Man bliver med Hensyn til Opholdelsen ved Umiddelbarhed. Verdens Opholdelse er Selvopholdelse. Enten sætter den Gud som den der trænger til Opholdelse, eller opgiver Verden som den der trænger til Opholdelse. 2) Causalitæts-Forhold. Opholdelse er egentlig identisk med Skabelse, og det Særegne der bliver tilbage finder sin Plads i Underet. Gud er adskilt fra Verden. En mechanisk Betragtning. Hieronymus ivrer mod denne Lære hos Pelagius og siger: dormitat ergo deus. 3) Immanents. Ikke eensidig, saa Transcendents udelukkes. I Substantialiæt ist Gott als Vater \emph{ver}kannt, im Causalitæt Verhältniß ebenso der Sohn, Sandheden er i Læren om Aanden. Verden hører til Guds middelbare Aabenbarelse. Bliver man staaende ved Identitæten, saa er Verdens Opholdelse Guds Selv-Opholdelse. At Verden formaaer at opholde sig, ligger ud over den i Gud. Denne Verdens Selvopholdelse er deri tillige dens Modsætning, det Levende er det Dødelige, alt er relativt, men har deri sin Relation til det Absolute; deri er Verden den af Gud opholdte.

\emph{Den concrete Bestemmelse.} Verdens Opholdelse ved sig selv og ved Gud, saa at begge Sider anerkjendes. Hvorved opholder Gud Verden: 1) umiddelbart; her blev nu Verden udelukket. Forholdet er vel et indre, men saa, at Verden er og bliver Gud et Andet. Verden er ikke den umiddelbare Aabenbarelse af Gud; thi det er Sønnen. Dette sker i Naturreligionerne især i Brahma. Derved er egentlig Verdens Virkelighed hævet. – ogsaa som Weltseele er Pantheismen med; Gud er draget med ind i den sandselige Existents, og Gud ikke erkjendt som Aand. 2) middelbart, vermittelst des Gesetzes; Loven er det Almene og det Nødvendige, uden hvilket Verden ikke kan opholde sig selv. Verden opholder sig ved sin Lov, Gud opholder Verden ved sin Lov. Loven er die entstandene Form seines ewigen gottlichen Inhalts. Saalænge Verden var i Samklang med Gud, bar Guds Villie ikke Charakteer af Lov. I Tiden faldt Skabningen af fra Gud. Dens Selbstsucht er dens Undergang; men Natur og Frihed er ikke tilintetgjorte, men forstyrrede. Gud overlod ikke Verden til Undergang, som det hedder i en Psalme, Guds Villie antog nu Skikkelse af Lov, der blev sat som en Grændse for det Onde, at ikke Kræfterne skulle hæve sig. Natur-Lov og Sædeligheds-Lov; der fordi de ere i den og i ham, ikke ere af den og af ham. Solon Lycurg Moses have været virksom til Verdens Opholdelse. 3) Religion. Gud erkjendes som Verdens Opholder, og opholder den derved. Men skikket til Bevidsthed er Mennesket alene, og saaledes opholder Mennesket, idet det opholder sig, ogsaa Verden. Ved Guds-Rige opholder Gud Verden, som det hedder, Gud lader for de Frommes Skyld Verden ikke gaae under, men Verden forvandles bestandig til Guds Rige. Religions-Geschichte er Historien om Guds Opholdelse; idet han til alle Tider tog sig af Mennesket, og til sidst antog Menneskets Natur.

3. Forsynet.

Bibelsk Lære. πϱονοια forekomer i Viisdommens Bog, men i ingen canonisk Bog i denne Forstand. Den kirkelige Lære. I Kirken bliver hvad der er berettet i Skriften til Tro. Forseer sig mod 1) Dualisme, Gnosticisme, Emanation. 2) alle Theorier, der sætte Forsynet som identisk med en Nødvendighed, der udelukker al Frihed. Fatalisme, Occasionalisme, Gud er den umiddelbare Aarsag, til Alt og de mellemliggende Aarsager ere blot Anledninger. Malebranche har mystisk gjennemført det, Bayle rationalistisk og indrømmer selv, at der ikke er Tale om Frihed – Kirkelæren indeholder en Mængde tilsidst reent traditionelle Bestemmelser. Med Hensyn til actus: prognosis, πϱοϑεσις, διοιϰησιν, (Udførelsen selv) med Hensyn til Objekt: providentia generalis, specialis, specialissima; med Hensyn til Middel: providentiaordinaria sive media, extraordinaria sive immediata. providentianaturalis og gratiosa. –

\emph{Begrebet af denne Lære.} 1) Det Substantielle 2) Dialektisk 3) det Concrete. – Det Substantielle er Begrebet i Almindelighed, man er med dette Begreb strax i Aandens Verden, dog saaledes, at den har et Forhold til Naturen. Gud som den evige Aand er baade Naturens og Aandens Princip, og Lovens i begge. Disse Love tænker man nærmest paa ved Forsynet. Altsaa er 1) alt Tilfælde udelukket i Aandens Verden. I Naturen er vel Alt bundet til Love saa at Nødvendighed hersker der, saa at selv det Monstrose har sin Lov; men efter sin Erscheinung falder Alt i det ubestemte. Det Tilfældige er det Tankeløse, ikke bestemt ved sig selv. I Naturen er Aanden ikke für sich. Er al Tilfældigt udelukket, saa er Alt bestemt ved Tanke det vil sige Aandens Rige 2) alt Planløst er udelukket. En Plan indeholder en Raadslutning og dens Udførelse. Sligt Dobbelthed findes ikke i Naturen. Det er umiddelbar Identitæt. Naturen handler vel \emph{efter} Tanker men \emph{uden} Tanker. Mennesket kan gjøre det komme ud i denne Dobbelthed, kan faae en Idealitæt uden Realitæt. Naturen har lutter Realisation. Det guddommelige Forsynets Planfuldhed er Naturen fremmed. Guds Plan er heller ei forskjellig fra hans Udførelse, hans Gjerninger ere Raadslutninger, saaledes ere de Tanker og have som saadanne deres Realitæt i Gud. 3) Zwecklosigkeit. I det Organiske er Hensigtsmæssighed ikke blot udvortes, men og indvortes. Men naar Zweck er Begreb saa tyder den hen paa Fornuft og Viisdom, men dette tyder hen paa Aand. Guds Rige er das Reich der weißesten und besten Zwecken. Fra Guds Side ere Zweckene altid opnaaede, men de sætte Skabningens Fornuft og Frihed og herved Forskjel mellem Zweck og Opfyldelse, denne Forskjel er Verdenshistorien.

\emph{Det Dialektiske:} Bevægelsen over forskjellige Standpunkter. 1) i Troen paa det Guddommelige i Naturen, Augurier, Haruspicier. Naturreligionen er dette naturligt, fordi den hviler i den umiddelbare Identitæt af Natur og Aand. I Tilfældet i τυχη regjerer endnu Naturen under Aandens Bestemmelse. Guds Villie adskilles ikke fra det Tilfældige og saaledes er det Overtro. 2) Forstanden adskiller det Naturlige fra det Aandelige; men derved er Mennesket endnu ikke kommen til Guds Erkjendelse, men blot til Negativitæt. Skjebnen er ud over det Naturlige, men det Aandelige er Gjenstand for Wahn og Furcht., og Alt er bestemt med uforand. Vilkaarlighed som i Islam. 3) Gud i sin Vished er som hævet over det Tilfældige og Naturens Magt saa og over alt Vilkaar. Forsynet skeer sin Ret i Overbeviisningen om Guds absolute aandelige Nærværelse, alt vidende, Alt nærværende. Den absolute Tænken og Villen er eet. I Natur og Historie bliver den evige Aand aabenbar i sin evige Verdens-Regjering

Det Concrete, Guds Eenhed med den af ham skabte og opholdte Verden, Eenheden med den i Verden herskende Aand. 1) Den sædelige Verdensorden. Bekymrer Gud sig alene om det Almene? i denne Betydning er det Almene slet Intet. Forstaaer man det derimod rigtigt saa kan man sige at for Gud existerer kun eet Menneske – 2) Resultatet af Menneskers Handlinger det er det guddommelige Forsyns Virksomhed. Gud medierer sine Raadslutninger ved Menneskets Fornuft og Frihed. Ideerne som Guds concrete Tanker er det, hvorved Gud er virksom i Verden, det Onde bliver til folie for det Gode. 3) Troen derpaa, en ved det guddommelige Aand bestandig bevirket, og tiltager i samme Grad som Fornuft og Frihed tiltager i Mennesket Guds Styrelser blive selv Videnskaben hemmelighedsfuld, derimod kan Viden ikke unddrage sig at vide Nødvendigheden af den Tro, der har denne Hemmelighed til Gjenstand. Vil man ikke vide Nødvendigheden af denne Tro, saa slaaer det over i Tvivl. Guds Planer kan vel være forborgne som Paulus siger, men Visheden om Nødvendigheden af den Tro kan den besidde. –

Christologie.

\emph{ Eenhed af guddommelig og menneskelig Natur }

Den bibelske Lære. De messianske Spaadomme. Hoved Momentet i det gamle Testamente er Ideen om guddommelig Aabenbaring, der ikke er forskjellig fra Guds Søn, og i det Nye Testamente Christi Præexistents. Denne Præexistents var heller ei Jøderne ubekjendt, som de chaldæiske Paraphraser og Rabbinerne. – Det Nye Testamente Lære om guddommelig Natur i Christo, om menneskelig, om Eenheden af dem. – Den nyere Exegese anerkjender nu dette bibliske Faktum eller at det indeholdes i det Nye Testamente Selv Rationalisterne om de end mene, at Læren er ubestemt i det Nye Testamente Supranaturalisterne vil heller ei tillade at denne Lære kommer til Klarhed, de fastholde de praktiske Motiver, men forholde sig ligegyldig mod det egentlig Dogmatiske.

Den kirkelige Lære.

M: Chemnitzde duabus naturis er et meget skarpsindigt Skrift.

unio naturarum, communio, communicatio idiomatum unio naturarum er Handling unitio, qua deus assumsit carnem, Virkningen deraf er communio.

unio personalis er ikke naturalis, qua duæ naturæ in unam coalescunt, ikke accidentalis men necessaria æterna, ikke mystica, ikke moralis, ikke sacramentalis. – communicatio idiomatum.

\emph{ Begrebet af dette Dogme. }

Den nyere Tids kristologiske Forsøg kan her danne en Indledning for at befinde sig i den historiske Udvikling.

1. \emph{Den kirkelig-symboliske}-supranaturalistisk bestemer Intet Nyt, men gjentager. Herimod kom Rationalismen og pirrede den op. Men Rationalismens Christologie er en abstrakt sig tilbagetrækken paa en Spidse af subjektive Tænken, at Christi historiske Erscheinung kommer i Fare. Det Menneskelige i Christo overveier aldeles det Guddommelige i Christo. Hovedsagen er at han er det praktiske Ideal, og at han har gjort sig selv dertil.

2. \emph{Den æsthetiske Theologie} \emph{.} bliver i den umiddelbare Følelse. de Wette-Læren om Christi Guddommelige er intet Begreb, men en æsthetisk Idee, den fromme Christen troer og skuer Gud i Christo, men klügelt nicht, man beholder denne Lære som et skjønt Billede, dog ikke som Digt, men som Ergebniß einer geschichtlichen religieusen Erfahrung.« – Schleiermacher. Det der er avanceret i S. er at han ikke opfattede Christi Betydning alene som Lære, men som Liv, dog blev han i denne Henseende staaende ved en Guds Væren i Christo, og ved det Urbildliche, som man ikke engang i Tanken kan komme ud over.

3. \emph{philosophiske Systemer} \emph{.} – Kant – Ideal. Til Gjendrivelse heraf er det nok, at erindre om den Kløft der etableres mellem Ideal og Virkelighed. I Fichte blev denne Kløft metaphysisk hævet, det Almindelige blev udsagt at Mennesket skulde forene sig med Gud, at det paa en ualmindelige Maade er Tilfælde med Christo, kan man ikke paanøde Nogen, hans Fortjeneste bliver det altid at have først fremsat den. Ogsaa hos Schelling bliver Kløften uopfyldt, han kom ud over Ks Ideal og Fichtes Ichheit og fremstillet igjen Incarnationens Idee. Guds Menneskevordelse er fra Evighed, men heri mangler Forholdet til den christne Lære om et bestemt Tids-Moment i et bestemt Menneske, medens hos Schl. Christus bliver Ideens Symbol, og maa erklære Klæren for ikke adæquat; men Symbol er langtfra at være et Dogme og nærved at være Mythus.

Den første Guds Aabenbarelse er Sønnen, den er umiddelbar, den anden Aabenbarelse er Verden. Disse to ere ikke absolut forskjellige eller efter hinanden i Tiden, men han er sig aabenbar og aabenbar for Verden. I sig selv aabenbar er han i Indifferents som Verden aabenbar i Differents, heri ligger Initiativet til den concrete Eenhed, som er Verden som menneskelig Natur i Eenhed med Gud, Gud-Mennesket forsoner Differentser. At dette er en evig Nødvendighed, uden hvilken hverken Gud eller Verden kan erkjendes, er Dogmatikens Opgave, der i dette Afsnit afhandler Anthropologie i eminent Forstand.

1) Eenheden af guddommelig og menneskelig 2) Den historiske Erscheinung 3) Identitæten af begge eller Gud-Mennesket

Før man erkjender Virkeligheden, maa man have Muligheden. 1) \emph{den menneskelige Natur } 2) \emph{den guddommelige } 3) \emph{Eenheden af begge.}

den menneskelige Natur. \emph{som Natur (nasci) begynder med Liv, Individualitæt, Selbstbewußtseyn Ichheit, Geist og deri Personlichkeit.}

\emph{som Individ.} har det vegetative animalske Liv an sich, Fühlen og Empfinden, der indeholder Handlen og Liden uden at være Handlen og Liden og er reen Receptivitæt. – Dog er Mennesket allerede ved sin Indtrædelse i Verden meer, han er an sich Fornuft-Væsen det vil sige der Möglichkeit nach. – Naturlichkeit i Modsætning til Gemütlichkeit, denne rolige Stilhed er ligesaa meget Følen som Tænken. Dyret er kun empfindende og har saaledes med det Enkelte at gjøre, men at sætte det Enkelte som det Almene er at tænke, det er wahrnehmen. Gemütlichkeit har dog sit Udgangspunkt i Sinnlichkeit og man maa ikke gjøre for meget deraf. Al Følelse er sinnlich er Afhængigheds-Følelse, i denne Følelse og dens Gemutlichkeit er Mennesket blot Natur og kan som Kjød og Blod ikke erkjende Guds Væsen. Den egentlige Erkjendelse komer fra Tanken og fra Guds Tanken, men er da ikke nedbøiende men opløftende, thi dens Princip er Aandens Hellighed. Natur og Sinnlichkeit, Gefuhl og Gemüt hæves og indoptages i det Høiere.

\emph{Selbstbewußtseyn.} Det Moment, hvor han fatter sig i sig selv og adskiller sig fra Verden. Som Gemüths-Wesen bliver han i umiddelbar Eenhed med Natur, men idet den komer til Forstand og Villie komer det frem betinget ved Objektiviteten ved Opdragelse og Underviisning, Spontaneitæten reiser sig derved at det Tilsvarende bydes uden fra. Nærmest er det Sproget. Han erkjender det Sandselige, Endelige, betingede. Komer han ogsaa til Forstand om det Gudd, saa er det atter fordi det Lige udenfra kommer ham imøde. I Ichheit er Mennesket vel det virkelige, men ikke det sande det er han først, naar Fornuften er i Forstanden Frihed i Villie, saaledes er han Aand. Individum som har Sjel og Legeme til Momenter. Men hvorledes komer Mennesket til Fornuft og Frihed og til Personlighed? her viser Mennesket ud over sig til den guddommelige Personlighed, den absolute Identitæt af Fornuft og Frihed, den absolute Betingelse for at blive Fornuft og Frihed er Religion. Den menneskelige Natur er ιδιοσυστατος, men og ανυποστατος og bliver det kun ved Optagelsen i den guddommelige Natur.

\emph{Den guddommelige Natur.} Menneskets Natur udvikler sig af Endeligheden fødes, lever, lider Død. Alt dette er negeret ved den guddommelige Natur, og selv Udtrykket Natur, der staar i Forhold til Vorden, er kun Betegnelsen af den guddommelige Substantialitæt, ogsaa Sinnlichkeit og Gemüthlichkeit er negeret i Gud. Kun en sandselig Tænken tilskriver Gud Affecter. Gud er den hellige det vil sige absolut aandelig Kjærlighed. Mennesket kommer til Forstand lærer at forstaae sig paa Verden og gjør den til Gjenstand for sin Villie, men bliver saaledes i Endeligheden. I Gud som den absolute Sandhed er Adskillelsen med Tænken og Gjenstanden er ogsaa negeret, Guds Tanker ere sande ikke fordi de stemme med Tingene, men fordi disse kun ere ved Guds Tanker. – Ogsaa Mennesket er istand til at fatte det absolute Sande, men behøver den guddommelige Oplysning, som naturlig Menneske veed han Intet om Gud, ved Gud kan han blive deelagtig deri – I Mennesket er Fornuften og Frihed Personlighedens Elementer og den sig vidende Frihed er hans Personlighed, Gud er den sande og hellige Aand. Mennesket som Aand er forskjellig fra Gud som Aand. Rom 8, 14. 2 Cor. XII, 4-6. bliver det ved denne gjensidige Negation, hvis Lighed er Forskjelligheden. Antager man nu ikke med Skriften en Mulighed af Forening og Identitæt, saa bliver Christus ubegribelig. Forøvrigt var end ikke Forskjelligheden ikke lod sig tænke, hvis den skulde være absolut.

\emph{Eenhed.} abstrakt erkjendes den i Skjulthed. Mennesket er ein Geist og Gud er der Geist. Mennesket er den i Endliche Unendliche, Gud den i Uendliche Unendliche. Mennesket er indflættet i die Erscheinung; den uendelige Aand er Negation af al Endelighed, forsaavidt er vi blot ved Forskjellen. Hverken den guddommelige eller den menneskelige Aand er uden en Beziehen. Sandheden af deres Forskjellighed er den enes Skjulthed i den anden. Som den skjulte Aand er den an sich; den menneskelige Aand som behaftet med Individualitæt er det Usande deri, men dog som Skriften siger et Tempel, hvori Gud boer, ved Daaben anerkjendes han som saadan. Har den menneskelige og guddommelige saaledes et indre Forhold saa er det mere end den ligegyldige Forskjellighed; men det er endnu umiddelbar, an sich, Mulighed, det er et forborgent Forhold. Colos. 3, 3. – Men Muligheden er der af et nærmere Forhold; Mennesket er endnu ikke kommen til sig, hvorledes skulde han være kommen til Gud.

\emph{Mediationen } ophæver Skjultheden og manifesterer det ene i det Andet. Denne Mediationen er Religion, der bereder Menneskets Guds-Vordelse og Guds Menneske-Vordelse. Begge Dele er den nødvendige indre Bevægelse i Bevidstheden. I Religionen gaaer den menneskelige og guddommelige Aand ud af deres gjensidige Hiinsidetshed, hvilket ikke er nogen Tilintetgjørelse, Gudsbevidstheden er Menneske-Bevidsthedens Sandhed. Mennesket anerkjender den guddommelige Aand, som det han trænger til; Gud fra sin Side viser og aabenbarer sig i sin Sandhed og Hellighed, Guds evige Godhed, hvori han nedlader sig til den menneskelige Natur. Det kan man med Schelling kalde Guds evige Menneskevordelse. I abstrakt og udvortes Forstand indrømmes det og af Rationalisterne, saaledes af Scott i hans Dogmatik; men ikke som immanent. Religionen er ikke blot efter Indhold men ogsaa efter sin Form Guds-Bevidstheden. Bevidstheden har ikke blot Natur og Historie til Indhold. Bevidstheden er ikke blot naturlig, saaledes begynder den, ikke blot menneskelig, og saaledes gaaer det frem, men væsentlig guddommelig

Guds og Menneskets Forhold til hinanden er som medieret altid et relativt, an sich er det et absolut. Forsaavidt de er Aand er der ingen Forskjel; Gud er den i sin Absoluthed med sig identiske, Mennesket er den i sin Ichheit med sig identiske. Hemmeligheden af denne Forening løser sig i Personlighedens Begreb, der hæver begge Naturer til Identitæt, den ophævede Ichheit sættes som Absoluthed, den ophævede Absoluthed som Ichheit. I christne Dogme er det bestemt som Entäußerung Phil:. Ved at ophæve sig som Ichheit har Mennesket overvundet sin Natur uden at opgive den; saaledes og med Gud. Bevægelsen gaaer alene ud fra den guddommelige Side, den menneskelige Natur kan ikke antage den guddommelige, men omvendt. Derfor bliver Christi Menneskevordelse altid analog med Naaden. – Allerede i gamle Testamente hersker den Forestilling, at hvo som saae Gud, maatte forgaae. – Foreningen af den guddommelige og menneskelige Natur maa vorde, er Process, dette er den nyere Philosophies Fremskridt. Dertil er den historiske Udvikling nødvendig, som Paulus siger, da Tidens Fylde kom. Heri ligger Tanken om Menneskets guddommelige Opdragelse. At Menneskets Forhold til ham selv tillige bliver hans Forhold til Gud, det er den guddommelige Opdragelse. Den trinvise Fremskridt viser, hvorledes den menneskelige Aand har maatte kjæmpe sig hen dertil. –

\emph{Die geschichtliche Erscheinung} eller\emph{den historiske Christus } For Troen er Beviset let ført, det staaer skrevet. Vantroen stiller sig derimod. Er det dogmatiske Begreb ikke udviklet, saa kan det Historiske ikke erkjendes. Af Bibelstæder følger nærmest kun, at de sige det. Troen har idet den griber dette deri die Wahrheit an sich. Men det er endnu ikke erkjendt i sin Sandhed. Kritik og Fortolkning er heller ei tilstrækkelig, og det bliver næsten tilfældigt om man bringer den sande dogmatiske Betydning frem eller ikke. Dogmet om Christus som Gud-Mennesket er ikke at aflede fra Dogmet om Forløsning; da hiint tværtimod begrunder dette. –

1) i Christi Syndefrihed. Christus er stillet midt in Verden; men ogsaa hævet høit over Mennesket Guds Søn. Det synes en Modsigelse, men er kun Tanken om Gud. Mennesket udsagt om en bestemt Person. – Mennesket Jesus som det enkelte Subjekt og det Almene dog prædiceret om ham er en lignende Modsigelse. Naar vi sige Jesus Christus saa sætte vi begge Dele i een Person. Det Almene Christus har naaet sin Personlighed i Jesus. – Uden Synd. Det er i Guddommelighed nødvendig indeholdt som Hellighed; som Menneske er det ikke givet efter Natur. Men i Foreningen i Gud-Mennesket er Syndefrihed nødvendig sat, og det der kan holde sig uden Synd og som har været uden Synd, maa ogsaa i sit Princip være sat saaledes. Derfor læres, at Christus er født af den Hellig Aand. Syndens Princip er fjernet fra Christi Fødsel, saaledes er den rene Menneskehed atter fremstillet. – Men med denne Syndefrihed kan den fri Udvikling bestaae. Den Mening maa fjernes, at det blot var en Natur-Nødvendighed, og Muligheden af at synde bør ikke udelukkes. Er Nødvendighed alene Grund uden Frihed er han sat under Mennesket; udelukkes endog Muligheden af at synde, da er han hævet over Menneske altfor høit. Den høieste Frihed kan ikke være et Værk af en Naturproces, paa den anden Side, da Retningen paa sig og den egen Villie er givet, saa fremt den rene Menneskelighed skal fastholdes, saa maa denne i ethvert Moment være bøiet over i den guddommelige Villie.

2) \emph{Die Verheißung und die Erfüllung} \emph{:} Heri adskilles det, der fra Guds Side fra Evighed ere forenede, begge føres tilbage paa Eenheden af den guddommelige Raadslutning, der er evig. Guds evige Raadslutning bestemer ligesaa vel Tiden til die Verheißung som til Erfullung. Heri ligger den successive Realisation af Opfyldelsen. – Prophetierne vise sig ofte i Forhold til Opfyldelsen meget udvortes, som Gaader, eller Tilfældigheder, hvis Opfyldelse ikke har sin Grund i det Forudsagte. De maa opfattes som en sammenhængende Totalitæt, eller bliver die Beziehung ganske vilkaarlig. – Guds Forsæt at sende Frelseren er altid uforanderlig den samme; Forløseren er bestandig i Forventningen, men er dog ikke deri. Man maa ikke sige at med Guds Forsæt at sende Forløseren tillige var sat først at lade det skee i en sildigere Tid. Med Forsættet begynder Udførelsen. Synden og Naadens Modsætning maatte først drives op paa Spidsen før Opfyldelsen kunde komme. Den tilsyneladende Guds Forkjerlighed for eet Folk er kun et andet Udtryk for Guds almindelig Kjærlighed til alle, der idet den viser sig som Resultat maa den og være sat som Princip. Syndens Bevidsthed maatte udvikles, og medens hele Hedenskabets Opmærksomheden var vendt mod Virkeligheden, sukkede Israel under Loven. – \emph{Die Erfüllung } Det gamle Testamente Theophanier, de ere blot forbigaaende Manifestationer, blot Skygger af det Tilkommende. Ogsaa i det gamle Testamente bliver det udvortes Forhold tilbage. At Christus er den forjættede Messias, han alene, det er Christi og Apostlenes udtrykkelige Lære. Man har nægtet det saaledes, at man har sagt, at den Forjættede var en Anden end Christus (det jødiske Folk, en enkelt Prophet etc), eller man har viist at Forjættelsens enkelte Træk ikke passe paa Christus. Begge Dele er en anden Maade at nægte Forjættelsen. – Opfyldelsen maatte indeholde meget Mere end Forjættelsen – Det er Nødvendigheden af at maatte hæve sig over Forjættelsernes historiske Verlauf. Kun af Guds og Menneskets Idee kan erkjendes om Aabenbarelsen er tilsvarende til Ideen, dog er det ikke en udvortes Reflexion, der kunde være vilkaarlig, eller Ideen først var hos os i Tanken, og Virkeligheden senere, men christelig Bevidsthed er det kun idet det veed, at Christus er Christus. Erfaringen kan ikke give nogen Abschluß, thi en Erfaring kan muligviis en anden modsige; Paa dette Erfarings Standpunkt er Jødedommen bleven staaende. Men alle Folks Udvikling stræber hen til denne ene.

\emph{Christus er Verdenshistoriens Middelpunkt.} Forskjellen mellem Christus og andre Verdenshistoriske Personer er kun quantitativ. Men herved stiller man sig dog eensidig paa hans Menneskelighedsligheds Standpunkt. Først naar han er Middelpunkt, da er han den eneste; thi der kan kun være et Middelpunkt. – I alle endelige Religioner er om ikke den Bevidsthed saa dog den Nødvendighed, ihre Erledigung zu finden i Christendommen Jødedomen havde det i de messianske Spaadomme; thi man kan i aandelige Henseende ikke vide sin Grændse uden allerede at være ude over. Judenthum og Romerthum vare de to verdenshistoriske Folk, der gaae til Grunde paa Christendommen Stoltheden havde begge disse tilfælleds, størst næsten Jøderne.

\emph{Christus som Gud-Msk.} Christus er mere end Middelpunkt, han gaaer ud over al Historie. Strausss Lære, man betvivler ikke den guddommelige og menneskelige Naturs Identitæt og Realitæt, men ikke i eet Individuum, det er ikke den Maade, paa hvilken Ideen pleier at realisere sig; men i en Mangfoldighed af Exemplarer, der gjensidig fuldstændiggjøre hinanden. Menneskeheden er Gud-Mennesket – Vi ere gaaede ud fra Eenh. af guddommelig og menneskelig Natur, derfra til den historiske Erscheinung, nu er det blot tilbage at sætte begge Dele som Eet. – Den Strausiske Lære har det Sande, at Christi Erscheinung ikke er ex abrupto at forstaae, men kun i bestandig Sammenhæng med hele Menneskeheden; den mechaniske Anskuelse af Christi Fremtræden som Deus ex machina tilhører Supranaturalisme. Strauss Feil er at han aldrig kommer videre, ikke til den sande guddommelige Meddelelse og den sande menneskelige Optagelse, Mennesket bliver ham dog et Mellemvæsen. Først naar man naar denne Høide, først da viser Betydning sig af de mangfoldige historiske Manifestationer, der gik forud. – Med den Tanke Menneskeheden tænker man vel at staae i Uendelighed fordi man har det abstrakte Vielheit, tager man den dermed concret saa viser det sig utilstrækkelig, og først det Einzelne er det sande Uendelige. – I Christus er guddommelig og menneskelig Natur forenede som aldrig før og aldrig senere, aldrig senere; thi ogsaa den christelige Menighed kan ikke træde paa Christi Plads, da man i saa Fald forvexler Incarnationen med Christi Aands Inhabitation i den Enkelte. Menigheden forvexles med dens Centrum. – Sandheden af Læren om Incarnationen er at Christus er kommen som den enkelte, dette Individuum. Ogsaa den historiske Udvikling har denne Formation med en enkelt Personlighed i Spidsen. Man siger, at Christus dog ikke har kunnet fremstille Alt i sig (Kunstner, Feldtherre etcetera), men den sædelig-aandelige Grund er indeholdt i ham, og denne Intensitæt er det vigtigste. Naar man siger at Menneskeheden er Guds Søn, saa ophæves egentlig Historie; thi Geniernes Cultus ere og blive syndige Mennesker, og ingen af dem kan være Forløser. Christus er den enkelte og det absolutte Individ, \emph{han er Menneskeheden men in der Einzelnheit} 1) som enkelt er enhver for sig i Forskjel fra Andre han fødes, lider døer og saa videre. I denne Henseende kan han sammenlignes med andre udmærkede Individer; thi det har han tilfælleds, at han er dette enkelte Menneske, Endelighedens Lod underkastet. Han tilhører en Familie. (»og han var dem underdanig«), det viser han endog paa Korset. Han tilhører en Stat. Han lever efter Folkets Skikke og Sæder. 2) Den historiske Retfærdighed fordrer nu ikke at oversee die Allgemeinheit i ham, saa Individualitæt, Familiaritæt, Nationalitæt ingen Skranke var. Han tager det op i sig, men gaaer ud derover. Det beskrænkende heraf er i ham absolut negeret. Heri er han forskjellig fra alle verdenshistoriske Charakterer og er Middelpunkt. Han har i Verden ikke für sich Noget i Verden, han er kun for Verden, i Kjærlighed er han uegennyttig, finder ingen Ære, resignerer paa Alt. Dette Negative er positivt Herskab i ham over Følelse, og over alle naturlige Bevægelser, han tillader Sligt i sig, men han er høiere. Selv hans Kjærlighed til Joh: er ingen personlig beskrænket. Heri er han den uforlignelige. Han er Medlem af Familie, men den indskrænkede Familie-Aand er ikke hans. Han ophævede den jødiske Particularisme, han overskrider sin Genealogie »før Abraham er jeg«, han tilhører derfor ikke den jødiske Nation. 3) Christus er Mennesket afseet fra al Familie og Nation. Menneske-Søn. Magten hertil er den guddommelige Ved denne guddommelige Concentration i et Individ gaar det Guddommelig-Menneskelige over i alle Individualitæts og Familiaritæts Indskrænkninger. – Christus behøver ikke Underviisning, hans Væsenhed og sande Natur er Sandheden og Helligheden selv, som kun successivt kommer tilsyne i hans menneskelige Individ:, før hans Fremtræden var den menneskelige Natur endnu skjult i ham. – Die an sich seiende Eenhed af det Gud. og Mennesket viser sig successivt i Udfoldelse i hans Individualitæt. Han bliver miskjendt, det Guddommelige er endnu skjult, først i sin Opstandelse blev hans guddommelige Natur ret klar og aabenbar. Var Christus i levende Livet bleven tilbedet, saa havde man tilbedet et Individuum.

Læren om Christi to Stænder.

\emph{Den bibelske Lære.} – \emph{Den kirkelige Lære.} Den erindrer om 1) at denne Lære kun gjelder om den menneskelige Natur, da Gud hverken kan ophøies eller fornedres. Men man kan dog egentlig kun sige det om den menneskelige Natur i Forhold til Guddommelige Subjektet i begge Stænder er Gud-Mennesket – 2) til Christi Fornedrelse maa ikke regnes den menneskelige Natur som saadan: Det er ikke i og for sig en Degradation for Gud at blive Menneske 3) den Giesenske og Tübingske Strid. – \emph{den speculative Bestemmelse.} Denne Lære er den dialektiske Bevægelse af Eenheden af guddommelig og menneskelig Natur i Christo. unio, communio, communicatio idiomatum. Eenheden af begge Naturer i Christo er en historisk og Eenheden er en vermittelte, som gaar igjennem disse Naturers Forhold endog indtil Widerspruch, saaledes for Exempel at Christus døer, dog saaledes, at Modsigelsen strander paa Eenheden. Den guddommelige Natur er ikke den menneskelige, den menneskelige ikke den guddommelige Eenheden af begge er nu ikke en umiddelbar, men fremgaaet gjennem Modsætningerne. Det Ophøiede kan kun fornedre sig, det Lave ophøie sig; men kun det Ophøiede kan nedlade sig til det, hvis Væsen det er beslægtet med, saaledes kan kun Gud nedlade sig til Menneske ikke til Dyret. Holder man fast ved Eenheden, saa behøver man ikke efter Calvinsk Methode at adskille Naturen ved ethvert Skridt for at forhindre en Forvirring; thi Identitæten slutter Differentsen i sig. Den lutherske Kirke holder sig mer til Identitæten, og forsaavidt kan Striden mellem dem let beroliges. De to Stænder maa ikke tænkes absolut udenfor hinanden eller i Succession, de er i hinanden, det zwiefache er det einfache αδιαιϱετως. Det komer an paa at fatte det Modsatte som Eet. Mennesket kunne Jøderne ikke miskjende i Christo, faster, beder, er bedrøvet, lider og døer; men han siger: Faderen og jeg er eet. Bevægelsen er den aandelige og derved fri, Kjerlighedens Værk, udgaaende fra den guddommelige Natur, og gjennemtrænger Menneske Den guddommelige Natur fornedrer sig, den menneskelige Natur ophøies. Dog maa den menneskelige Natur ikke tænkes blot passiv, thi ellers er Forholdet blot udvortes; Forholdet er ikke ved Noget udvortes medieret, heller ikke Naadens Værk. Formen herfor er Underet. Den gamle Naturs ϰϱυψις og φανεϱωσις er begge Dele Sider deraf. Christus er et Under selv naar han ikke gjør Mirakler. Eenheden af guddommelige og menneskelige Natur er Principet i alt Under. Det var ubegribelig, hvorledes den kunde vise sig uden Under. Det er imidlertid ikke Beviset for de enkelte Mirakler, som Gud-Mennesket kan forøge eller formindske. Han unddrager sig Mængden og vil ikke udføre Mirakler. Han forlanger at man skal troe alle hans Mirakler, og dog vil han tillige, at man skal være ude derover. Videnskaben har ikke at begribe det enkelte Mirakel, men Mirakelets udeblivelses Nødvendighed. Efter Daub er Miraklet Eenheden af det Historiske og det Dogmatiske. Underets Form er menneskelig Thun, dets Indhold er guddommelig Thun. For Gud abstrakt tænkt er Underet ikke, og det er intet Under; for det naturlige Menneske er der heller intet Under; kun i Menneskeheden, som den er vaagnet for Gud, begynder denne Tro som Beundring. Christi Livs-Momenter ere ikke blot Thaten af ham og med ham men er tillige Dogmer. Herefter maa ogsaa bedømes hvad der hører til de to Stænder. Hvad der gjør Opfattelsen af dem vanskelig er deels den sandselige Vished af Øienvidner, deels hver enkelts Isolerethed, der staaer i Modsigelse med deres uendelige Indhold. Det Sandselige kan ikke hjælpe Aanden. Der gives Mange, som have den sandselige Vished og ansee det som Hovedsagen, hvorved de let taber det Aandelige, der dog først gjør dem til Troes-Gjenstande. Det Sandselige kalder man ogsaa det Faktiske og skjøndt for Exempel ved Christi Opstandelse det Sandselige træder tilbage, naar Christus siger: jeg er Opstandelsen. Ved det Sandselige komer Rum og Tid i Betragtning. Der er vel en Succession i disse Tilstande, og derved har en hver især en für sich seyn, før den gaaer over i den anden; dette synes at hæve den dialektiske Bevægelse. Men aldeles udvortes maa denne Succession ikke opfattes ikke som Tal-Succession. Det maa derfor opløftes i Aandens Region. Disse Momenter staae alle paa Underets Gebeet, og ere Gjenstande for Troen. Viden kan nu ikke ville opløse dette. Heller ikke er det nok med Schelling at ansee Alt for Symbol; thi dermed gaar den historiske Christus ikke ud over Osiris, og disse mythologiske Figurer kunne altsaa med Rette paralelliseres dermed. Disse Momenter maa ikke subjektivt udtydes. Videnskaben maa kjende deres Nødvendighed.

\emph{Christi Fornedrelse.}

Christi Fremtræden i Tjener-Skikkelse. Dette maa ikke forstaaes om Menneskevordelse overhovedet, denne er den ud over Tidens Fylde gaaende fra evig sig Inddannelse i Menneskets Væsen. Tjener-Skikkelsen har til sig tilknyttet som endeligt Menneske mange Lidelser, og ender først i Døden; men den er frivillig paataget for at døde Døden, og hans Død er den forløsende Død, og det Modsatte er forenet, Død og nyt Liv. Da hans Fornedrelse ikke blot er Passivitæt, men og That, saa tager han det tilbage, det er Opstandelse. – Hollenfahrt er Eenheden af Død og Opstandelse, han er vel ikke opstanden, men han lever. Hvori hans Virksomhed har bestaaet i de 3 Dage, har Kirken altid overladt til Enhvers fri Mening. – Opstandelsen: En Skin-Død hæver Døden, en virkelig Død hæver Opstandelsen i Rationalismens Tjeneste. Skriften vil begge Dele, en virkelig Død og en virkelig Opstandelse. Himmelfart. Ogsaa disse 2 Momenter følge efter hinanden i Tiden, væsentlig kunde de ikke adskilles. Efter hans Opstandelse kunde kun hans Troende \emph{see} ham. Hans pneumatiske Legemlighed er Begyndelsen til Himmelfarten. Der er en flydende Overgang mellem Opstandelse, Himmelfart og Sidden ved høire Haand. Weisse gjør opmærksom paa at der ingen væsentlig Forskjel er, og at det at han viser sig for 500 Troende er en Nedstigen fra den Værdighed, hvori han forsvandt. Saa at Opstandelsen ikke er Legemets Fremgang af Graven, men Sjælens Opløftelse fra Hades til Guds Høire, som er den efter Legemet dødede efter Aanden levendegjorte. Aandelig vender han tilbage, tilsidst til Dom, hvorved der vendes tilbage til Begyndelse, og enhver anden Messias er udelukket.

\emph{Forløsningen.}

\emph{bibelske Lære.} Prophet ikke blot Underviisning, men ogsaa hans Liv, hvilket som belærende er hans Exempel. – Præst – Konge det vil sige han er Prophet og Præst ewiger Weise. –

\emph{Kirkelige Forestilling.} Den gamle Inddeling i 3 Virksomheder. – satisfactio. –

\emph{Dogmets Begreb.}

Christus Prophet. Identitæten af den guddommelige og menneskelige Natur i Christo er en Viden, det er en Aabenbaring. Først i den christne Aabenbaring har Aanden naaet sit høieste Trin, og altsaa er enhver foregaaende den ikke fremmed. Den enkelte Bevidsthed og subjektiv Aand har først sin Sandhed \emph{i} den almindelige og objektive Aand. Den subjektive Aand har sin nærmeste Bekræftigelse i Folke-Aand. Subjektets Religion har sin Sandhed i Folkets Religion. Isoleret er den Sværmerie. Først er det \emph{Naturreligion,} eller Ideal, Kunstreligion. Bevidstheden er saaledes deels naturlig deels æsthetisk. Det tredie er Bevægelsen ud over begge – Natur-Religion, den formaaer ikke Idealet, den er tabt i Veltseele, det er Fetischisme. – \emph{Kunst-Religion,} æsthetisk Religion. Bevægelsen er fra Idee til Ideal. – Offret er overalt Udtrykket for Endelighedens Bevidsthed. Identitæten af Idealet og Anskuelsen er Kunst-Religion – I Enden udtrykker Folket sin Bevidsthed om at det kun er ved et Høiere – \emph{den aabenbare Religion.} hermed har Religionen hævet sig til det andet Trin, til Reflexionen. Dens abstrakte Tænken viser sig saaledes tydeligt i hiint Forbud, Du skal ingen Billede danne af Gud. 1) den jødiske Religion. Her hører Propheterne hjemme. Den har Hukommelsen til sit Element, gaaer tilbage i det Historiske. Tradition. paa den ene Side hedder det at Gud \emph{har} viist sig for Abraham, Isaak og Jacob, paa den anden Side er Gud for Folket en Hellighed som ikke kan aabenbare sig, hiinsides liggende. Men herved er det jødiske Folk langt nærmere ved den egentlige Guds-Viden. Som Guds udvalgte Folk er det det verdenshistoriske Folk. – Folket kjender kun Gud in abstracto, han er det høieste Væsen, ens supremum. – Den israelitiske Nation ligger imellem Tidens tvende Dimensioner, Fortid og Fremtid, men mangler den nærværende Tid, dets Guds-Bevisthed er den ulykkelige Guds-Bevidsthed – Abstraktionen hører nu op, Identitæten viser sig. 2) \emph{den christelige Religion,} der er ikke noget Trin mere, den er den absolute. Menneskeslægtens Bevidsthed som det i sin Almindelighed Besondere og Einzelne Ichheden hvori al Negativitæt er ophævet er Jesus Christus. – ikke den tidligere Aabenbaring er Grundlag for den senere, men omvendt. Uden det nye Testamente er det gamle Testamente ikke blot ikke forstaaelig, men uforstandigt. – Som før et Folk var det verdenhistoriske, saa er nu et Individ det verdenshistoriske. 3) \emph{den christelige Religion} som Forløsning ved Christus. a) Stiftelsen af den sande Religion. Det er skeet ved Christo, da han gav den et saadan Indhold og en saadan Form hvori den fra nu af kunde være alle Folks Religion, deri er den den hele Verdens Forløsning. Efter sit Væsen er Christendommen ikke opstaaet, den ved Gud lige evige Bevidsthed ved Gud, den lige evige Guds Aabenbaring har hverken Begyndelse eller Ende. Christendommen er saa gammel som Verden, indeholder alle Religioner als seine Negationen an ihm. – Forskjellen mellem den christelige Religion og de tidligere, er alene Formen. De foregaaende Religioner er Christendommens Usandhed. Christus veed Gud saaledes som Gud veed sig selv. Det Menneskelige i Christendommen er ikke som i de tidligere Religioner det Ufuldkomne. Christus er derfor ikke blot Stifter af en Religion, men hans Liv er Hovedsagen. Christi absolute Viden begynder i alle andre Mennesker med Troen. b) Troen. Apostlene virke til at hæve Folkene ind i denne almene Bevidsthed. Men Læren er ikke noget fra Religionen forskjelligt, ligesom det heller ei er forskjelligt fra hans Liv. Menighedens Virksomhed er ikke at lære, men at troe og optage i sig, iføre sig. – Den christelige Religion overstiger den menneskelige Naturs Erkjendelse, det er det naturlige Menneskes. Det første Indtryk er at den er ham et aldeles fremmed Indhold, og dog er den kun for ham. Han er optaget i denne Lære ved Aandens Naade, derved skeer det, at den guddommelige Sandhed ophører at være ham fremmed. Først i Christendommen finder Mennesket sin Selvbevidstheds Gaade løst, og saaledes gjør den Mennesket først i Sandhed fri. – Forsaavidt Troen beroer paa noget Overleveret, er den traditionel, i Videnskb. er den spekulativ, den har den χstlig Sandhed ikke blot til sin Gjenstand, men Visheden, hvorved den adskiller Troen fra Overtro og Vantroe. Med Hensyn til Indholdet er al Forskjel mellem esoterisk og exoterisk Viden hævet.

Forsoning.

Ypperstepræst.

1) Nødvendighed af dette Dogme, 2) Muligheden, 3) Virkeligheden. –

Strauss opfatter den lige fra Begyndelsen abstrakt, lader Idee og Historie falde ud fra hinanden. –

\emph{Nødvendighed.} for Mennesket for Gud. \emph{for Mennesket } Man har gjort opmærksom paa, at Synden til Christi Fremtrædelses Tid var steget saa høit, at den var nødvendig, heri ligger altsaa, at i Almindelighed vilde den ikke have været nødvendig, for Gud selv er Syndens Størrelse den bevægende Aarsag. Historisk lader det sig bevise, at det var Pharisæerne etc der bragte Christus til Døden. Men denne Betragtning er reent udvortes og tilfældig. Men Skriften lærer,
\end{quote}
\dcite{Not9:1  (1841; Pap. IIIC26; SKS 19, 249)}
% Not9:1  (1841; Pap. IIIC26; SKS 19, 249)

\subsection{1843}

\begin{quote}
Min Bestemmelse synes at være, at jeg skal foredrage Sandheden, saavidt jeg opdager den, saaledes, at det skeer ved samtidig at tilintetgjøre al mulig Myndighed. Idet jeg da bliver umyndig, i høieste Grad upaalidelig i Menneskers Øine, siger jeg det Sande og bringer dem derved i den Modsigelse, hvoraf de kun kan hjælpes ud ved selv at \tlg{tilegne} sig Sandheden. Først \emph{dens} Personlighed er modnet, der \tlg{tilegner} sig det Sande, ligegodt enten det er Bileams Æsel der taler, eller en Grinebidder med sin Skoggerlatter, eller en Apostel og en Engel.
\end{quote}
\dcite{JJ:97  (1843; Pap. IVA87; SKS 18, 170)}
% JJ:97  (1843; Pap. IVA87; SKS 18, 170)

\subsection{1845}

\begin{quote}
confer pagina 182. 189.

Allerede her ligger det Dialektiske i Forhold til det Religieuse. Naar for Exempel et Menneske vilde bedende sige til Gud, at han endnu ikke havde ganske \tlg{tilegnet} sig Christendommen, at han havde Tvivl tilbage, og at Gud derfor vilde give ham Tid til at overvinde dem (og saaledes gaaer der jo altid en Tid hen), saa griber strax det Dialektiske ham, thi denne Tids-Anvendelse bliver i samme Øieblik uendelig dialektisk (sæt han døde imorgen, og Christendommen var det ene Saliggjørende, og han udenfor), fordi den evige Afgjørelse fordrer sig selv med en uendelig Hastighed. Dette viser igjen, hvor vanskeligt det er at faae et historisk Udgangspunkt i Tiden for en evig Salighed.
\end{quote}
\dcite{JJ:340  (1845; Pap. VIA60; SKS 18, 251)}
% JJ:340  (1845; Pap. VIA60; SKS 18, 251)

\subsection{1846}

\begin{quote}
\emph{ Af den Bog om Adler }

Det er overhovedet utroligt, hvilken Confusion den hegelske Philosophie har frembragt i det personlige Liv – sørgelige Følge af, at en Philosoph er Heros og reent personligt en Philister og Nittengryn. Eet er bestandigt undgaaet Hegel hvad det var at leve, han veed kun at \emph{gjengive} Livet, og hvis han er Mester heri, saa er det ogsaa vist, at han er den meest skærende Contrast til en Maieutiker.

Man finder hos Philosopher efter Hegel, som have \tlg{tilegnet} sig den hegelske Methode, forbausende Exempler. En saadan Philosoph skriver en ny Bog, og den bliver han sig bevidst som Moment indenfor den Stræben, der begyndte med hans første Bog; men dette er ikke nok, thi sin hele Stræben (som endnu ikke er til) bliver han sig bevidst som Moment i den hele philosophiske Stræben efter Hegel, og den igjen som Moment i Hegel, og Hegel som Moment i den verdenshistoriske Proces fra Oldtiden gjennem China, Persien, Grækenland, Jødedom, Christendom, Middelalder. Skulde det ogsaa være muligt. Skulde det ikke være rigtigt, at en saadan Privat-Docent, naar han saae vi bleve altfor forbausede, sagde som Præsten, der saae at hans Tilhørere bleve altfor rørte: Græder ikke Børn, det turde være Løgn altsammen.

Men denne usalige Lyst til at opfatte som Moment er bleven en fix Idee. Ethikens Anskuelse: at stræbe, og Metaphysikens at opfatte som Moment stride paa Liv og Død med hinanden. Enhver Levende, som ikke er aldeles tankeløs og distrait, maa vælge, men vælger han det Metaphysiske, saa begaaer han egentlig, aandelig forstaaet et Selvmord.
\end{quote}
\dcite{NB:42  (1846; Pap. VII1A153; SKS 20, 44)}
% NB:42  (1846; Pap. VII1A153; SKS 20, 44)

\begin{quote}
Til den Dedication:

»Hiin Enkelte«

i Leiligheds-Talen, skulde egentlig

have været føiet følgende Skrivelse.

Kjære!

Modtag denne \tlg{Tilegnelse}; den gives hen ligesom iblinde, men derfor ogsaa uforstyrret af noget Hensyn, i Oprigtighed! Hvo Du er veed jeg ikke; hvor Du er, veed jeg ikke; hvilket Dit Navn er, veed jeg ikke – jeg veed end ikke, om Du er til, eller om Du maaskee var til, men ikke mere er det, eller hvorvidt muligen Din Tid engang vil komme. Dog er Du mit Haab, min Glæde, min Stolthed, i Uvisheden min Ære; i Uvisheden, thi kjendte jeg Dig personligen med jordisk Vished, da var dette min Skam, min Brøde – og min Ære tabt.

Det trøster mig, at Beleiligheden er der for Dig, hvad jeg under min Arbeiden og ved min Arbeiden redeligen har tilsigtet og eftertragtet. Thi ikke, dersom dette var muligt, ikke hvis det blev Skik og Brug at læse hvad jeg skriver, eller dog at udgive sig for at have læst, i Haab om at vinde Noget i Verden, ikke da var Beleiligheden for \emph{min} Læser, da havde tvertimod Misforstaaelsen seiret ja den havde tillige bedaaret mig til Uredelighed, hvis jeg ikke af yderste Evne havde forhindret, at noget Saadant skete – hvorimod jeg ved at have gjort Alt for at hindre det havde handlet redeligt. Nei, naar det bliver et tvivlsomt Gode at læse hvad jeg skriver (og dersom jeg efter hvad Evne mig er forundt bidrager her til, da handler jeg redeligen); eller endnu bedre, naar det bliver en Daarskab og en Latterlighed; eller endnu bedre naar det bliver en foragtet Sag, saa Ingen tør være det bekjendt: da er Beleiligheden for min Læser. Da søger han Stilhed, da læser han, ikke for min Skyld, ikke for Verdens Skyld, men for sin egen Skyld, da læser han saaledes, at han ikke søger mit Bekjendtskab men undgaaer det – og da er han \emph{min} Læser.

Ofte har jeg tænkt mig i en Præsts Sted; hvis da Mængden stormer til for at høre, hvis Kirkens Hvælving ikke kan rumme den store Skare, der endnu udenfor staaer lyttende – ja, Ære og og Priis være Den, der er saaledes begavet, at han da kan gribes i Stemning, at han da ret kan tale begeistret, begeistret ved Synet af Mængden, i Forvisning om at hvor den er maa Sandheden være, begeistret ved Tanken om, at der dog maa blive Lidt til Nogle, da der er Mange, og Mange der har lidt af Sandheden er jo Sandheden: mig var det en Umulighed! Men dersom det var en Søndag Eftermiddag, dersom Veiret var uhyggeligt og forstemmende, dersom Vinterstormen gjorde Gaderne tomme, og Enhver, der havde den varme Stue lod Gud vente i Kirken til det blev smukkere Veir – dersom der da i den tomme Kirke sad et Par fattige Koner, der ingen Varme havde i Stuen, og altsaa lige saa gjerne som hjemme kunde fryse i Kirken: sandeligen jeg skulde tale baade dem og mig varme! – Jeg har ofte tænkt mig ved en Grav; hvis da alt Udmærket og Herligt var samlet i Følget, hvis Festligheden var udbredt over den talrige Skare – ja Ære og Priis være Den, der var saaledes begavet at han da kunde forøge Festligheden ved bevæget at være de Manges Tolk, være Udtrykket for Sorgens Sandhed: jeg formaaede det ikke! Men dersom det var den fattige Liigvogn, og Ingen fulgte, dersom kun en fattig Kone gik ved Siden, den Dødes efterladte Hustrue, hvem det nu første Gang hændte, at Manden tog ud uden at tage hende med – hvis hun bad mig derom: ved min Ære jeg skulde holde Liigtale trods Nogen. – Ofte har jeg tænkt mig i Dødens Afgjørelse; hvis der da blev Allarm i Leiren, og megen Løben for at spørge til mig: jeg troer ikke jeg kunde døe, jeg troer min gamle Stridbarhed vilde endnu engang vaagne, saa jeg endnu engang maatte ud og stride med Menneskene. Men hvis jeg laae afsides og ene: jeg haaber til Gud, at jeg skulde døe roligt og saligt.

Der er een Anskuelse af Livet, som mener, at der hvor Mængden er, er ogsaa Sandheden, at det er en Trang i Sandheden selv, at den maa have Mængden for sig. Der er en anden Anskuelse af Livet; den mener, at overalt hvor Mængde er er Usandheden, saaa hvis end alle Enkelte, hver for sig i Stilhed havde Sandheden, vilde dog, dersom de kom sammen i Mængde (saaledes dog, at \emph{Mængden} fik nogensomhelst afgjørende, stemmende, larmende, lydelig Betydning) Usandheden strax være tilstædeb. Men Den, som vedkjender sig denne (sidste) Anskuelse, (der da sjeldent foredrages, thi dette forekommer oftere, at en Mand troer, at Mængden er i Usandheden, men naar denc blot vil antage hans Mening saa er alt rigtigt), han tilstaaer jo selv, at han er den Svage og Afmægtige; hvorledes skulde ogsaa en Enkelt kunne staae mod de Mange, der have Magten! Og dette kunde han da ikke ønske, at have Mængden paa sin Sided: det var jo at spotte sig selv.e Men er saaledes denne Anskuelse fra Første af Svaghedens og Afmagtens Indrømmelse, og synes den maaskee derfor saa lidet indbydendef: saa har den det Gode, at den er ligelig, at den Ingen fornærmer, ikke een Eneste, at den ingen Forskjel gjør, ikke paa een Eneste. Mængden dannes jo af Enkelte; det maa altsaa staae i Enhvers Magt at blive hvad han er: en Enkelt; fra at være en Enkelt er Ingen, Ingen, Ingen udelukket, uden den som udelukker sig selv – ved at blive Mange. At blive Mængde, at samle Mængde om sig er derimod Livets Forskjellighed; selv den meest Velmenende der taler derom kan let fornærme en Enkelt. Men saa har Mængden igjen Magt, Indflydelse, Anseelse og Herredømme – det er ogsaa Livets Forskjellighed, der herskende overseer den Enkelte som den Svage og Afmægtige, timeligt-verdsligt overseer den evige Sandhed: den Enkelte.

\emph{[i margen:]} (for et Øieblik at føre Sagen ud paa sin yderste Spidse)

\emph{[i margen:]} confer Hoslagte

\emph{[i margen:]} – Mængden –

\emph{[i margen:]} for at drive sin Anskuelse igjennem, at Mængden er Usandheden

\emph{[i margen:]} Altsaa tilstaaer han jo han maa ligge under.

\emph{[i margen:]} og høres den maaskee derfor saa sjeldent
\end{quote}
\dcite{NB:64  (1846; Pap. VII1A176; SKS 20, 55)}
% NB:64  (1846; Pap. VII1A176; SKS 20, 55)

\begin{quote}
4.

.... Thi saaledes er det jo, Een giver Gud megen Glæde i Livet og gjør ham stum, en Anden nægter Gud megen Glæde i Livet og gjør ham veltalende: er der saa ikke Ligelighed! En gjør Gud stor i Verden og gjør ham misundt, en Anden gjør han ringe i Verden og gjør ham velsignet: er der saa ikke Ligelighed! En giver Gud den Elskede, men hun forstyrrede Forestillingen, en Anden nægter han den Elskede, men lader ham beholde Forestillingen: er der saa ikke Ligelighed! Een giver Gud Ære i Verden og see han \tlg{tilegner} sig den; en Anden gjør Gud ringeagtet i Verden, og see, han den Ringeagtede giver Gud Æren: er der saa ikke Ligelighed! Dog siger maaskee Nogen: »det er ikke saaledes, det hænger ikke rigtigt sammen med denne Ligelighed, selv i den Talendes Stemme er der Veemod med i Spillet.« Ja vist er der Veemod, og det skal der være; thi al Tale om Menneskets Liv paa Jorden, hvis der i den ikke er Veemod med i Spillet: da er den klangløs eller falsktonende. Ja vist er der Veemod i den; thi den Talende har ogsaa drømt sit Ungdommens Eventyr, det gamle, bekjendte, som Børnene fortælle i Aftenstunden, at langt henne i Skoven saae han et gammelt Slot, hvor der boede en Prindsesse – og ganske saaledes fandt han vel ikke Verden, men han fandt netop heller ikke Ligeligheden i Eventyret.
\end{quote}
\dcite{Papir 339:340:5  (1846-05; Pap. VII1A134; SKS 27, 352)}
% Papir 339:340:5  (1846-05; Pap. VII1A134; SKS 27, 352)

\begin{quote}
10

Fader i Himlene! Ubegribelig er Du i Din Skabning, Du boer fjernt i et Lys, som Ingen kan trænge ind til, og kjendes Du end i Dit Forsyn, vort Kjendskab er dog kun svagt og fordunkler Din Klarhed, Du som er ubegribelig ved Din Klarhed. Men endnu ubegribeligere er Du dog i Din Naade og Barmhjertighed. Thi hvad er et Menneske for Dig Du Uendelige, at Du dog kommer ham ihu – men endnu mere, hvad er den faldne Slægts Søn for Dig, Du Hellige, at Du dog vil besøge ham; ja hvad er en Synder, at for hans Skyld Din Søn dog vilde komme til Verden, ikke for at dømme men for at frelse, ikke for at lade sit Opholdssted blive bekjendt, at det Fortabte kunde søge til ham, men for at søge det Fortabte, uden at have en Hvile, som dog Dyret har det, uden at have en Steen at helde sit Hoved til, hungrende i Ørkenen, tørstende paa Korset. Barmhjertige Gud og Fader! hvad formaaer et Menneske at gjøre til Gjengjeld; end ikke at takke Dig formaaer han – uden Dig. Saa lære Du os da den rette Forstaaelses ydmyge Skjønsomhed, at som Den, der sønderknuset sukker under Skylden siger i sin Sorg: det er umuligt; det er umuligt at Gud saaledes kunde forbarme sig – at saaledes Den, der i Troen \tlg{tilegner} sig det, maatte sige i sin Glæde: det er umuligt. Hvis endog blot Døden syntes at ville adskille de Elskende, og de atter gjengaves hinanden, da var dette jo det første Udraab i den visse Gjenforenings Øieblik: det er umuligt. Og Glædens (dette glade) Budskab om Din Barmhjertighed himmelske Fader – ja, om end et Menneske fra sin tidligste Barndom hørte det: var det derfor mindre ubegribeligt! Om end et Menneske dagligen betænkte det: blev det derfor mindre ubegribeligt!1 Var da Din Naades Ubegribelighed som et Menneskes, der var engang men forsvandt ved det nærmere Bekjendtskab; være (var den) som de Elskendes Lykke: engang i fordums Dage ubegribelig men nu ikke mere! O, dorske menneskelige Skjønsomhed, o svigefulde jordiske Viisdom, o indslumrede Troes matte Tanke, o kolde Hjertes usle Glemsomhed – nei Herre, bevar Du hver Din Troende i den rette Forstaaelses ydmyge Skjønsomhed og frels ham fra det Onde!

Skulde da Din Naades Ubegribelighed være som et Menneskes
\end{quote}
\dcite{Papir 339:340:13  (1846-05; Pap. VII1A142; SKS 27, 356)}
% Papir 339:340:13  (1846-05; Pap. VII1A142; SKS 27, 356)

\begin{quote}
Lovtale over Efteraaret.

af 5 Personer

som endes med et Tutti.

\emph{Efteraaret er Skyernes Tid}

Ifølge den nordiske Mythologie dannedes som bekjendt, Skyerne af Kæmpen             s Hjerne. Og sandeligen der gives intet bedre Sindbillede paa Skyerne end Tanker og intet bedre paa Tanker end Skyer – Skyer er altsaa Hjernespind og hvad er Tanker Andet. See, derfor bliver man træt af alt Andet men ikke af Skyerne – om Efteraaret, der igjen er Tænkningens Tid. Jeg har forlængst træt opgivet Menneskeheden skjøndt jeg dog eller netop fordi jeg har studeret den grundigt. Jeg begyndte som Yngling med at elske ældre Mænd – og blev kjed deraf; derpaa elskede jeg som selv lidt Ældre Ynglinge og blev kjed deraf; saa kom de unge Piger – og saa Matronerne, det Bedste unægteligt, men det er idel Forfængelighed selv hvor de som en Matrone har overstaaet Fatalia. Træt vendte jeg mig fra Menneskeheden og humoristisk til hine Skabninger, der vel ikke som Menneskene gjøre Prætension paa at være Skabningens Underværk, men dog, maaskee just derfor, ere meget underholdende, fordi de ikke prætendere saa meget: jeg mener Køerne. Og unægteligt stiftede jeg ogsaa saa vel paa Fælleden som andetsteds mange interessante Bekjendtskaber – jeg glemmer aldrig den Blakkede – og de glade Timer og saa videre – men dog blev jeg kjed deraf.

Men af Skyerne – om Efteraaret aldrig. Man kjender neppe Skyerne igjen saa forandrede ere de, og Den som aldrig saae dem om Efteraaret saae dem aldrig. Om Vinteren er det for koldt at see paa Skyerne – om Sommeren ere de dvaske, søvnige – men om Efteraaret: drømmende. (udføres). – Om Sommeren staae de stille og kjede sig – om Efteraaret ile de som de flygtige Drømmersker. Om Sommeren er det som gad de ikke hænge deroppe – om Efteraaret er det tydeligt at de bære sig selv i en vellystig Svæven. Om Sommeren gider den ene knap gaae af Veien for den anden: om Efteraaret lege de med hinanden i livsalig Tidsfordriv, ile hinanden forbi, mødes, skilles, længes igjen efter hinanden, blande Farverne (som Venner blande Blod) med hinanden, blive til Eet, skjøndt dog de Tvendes Eiendommelighed skinner igjennem.

Stat derfor stille Du som kalder Dig en Tænker, betragt Skyerne – om Efteraaret. Har Du før tænkt over Andet, saa tænk om igjen. Betænk hvad Du kunde ønske at være – et Menneske? Det kunde dog vel neppe falde noget Menneske ind. En Engel? Det er kjedeligt. et Træ? det er for langvarigt og for beroliget. en Ko? Det er for solid en Existents. Nei en Sky – om Efteraaret. Det gad jeg være. Den øvrige Deel af Aaret vilde jeg opholde mig skjult etsteds, eller i Intet, hvilket ogsaa kan udtrykkes saaledes: jeg vil ikke være til – men hvert Efteraar vilde jeg leve den Maaned. eller jeg vilde leve skjult som en Tanke, indtil Efteraaret kom, da jeg blev en Sky. En Sky er i sig selv et ganske anseeligt Stykke (og jeg vilde da ikke være saadan en lille bitte Makrel-Sky), jeg vilde være en stor, velskabt Sky. Men trods Størrelse er Skyen tillige Letheden selv, elastisk, og har tillige, forsaavidt den ikke har musicalsk Sands (hvad den dog har) Farve-Sands: den veed at bade sig i Farver.

Leve derfor Efteraaret. Med dette Bæger hilser jeg \emph{erindrende} Eder I dybsindige og dog saa flygtige Tænkere, I, mine bedste Tanker, hvem man ganske kan \tlg{tilegne} sig uden Plagiat. Naar Efteraaret kommer, kaster jeg mig i en Vogn, trækker Forlæderets Overdyne helt op over Hovedet, hylder mig i en Kappe – saa kun mine Øine blive tilbage, med hvilke jeg griber efter Eder. Naar da Kudsken kjører saa hurtigt som Heste formaae det (ak ak hvad er dog dette mod Skyer) da er det mig, som var jeg ogsaa blevet en Sky.
\end{quote}
\dcite{Papir 341:344:1  (1846; Pap. VII1B207; SKS 27, 368)}
% Papir 341:344:1  (1846; Pap. VII1B207; SKS 27, 368)

\subsection{1847}

\begin{quote}
Fredags-Talerne kan dog ikke dediceres til Mynster. Jeg vilde det saa gjerne, mindende min Fader. Det var ogsaa ingen almindelig Dedication, men just i det Øieblik, da jeg sætter et Slags Punktum om ikke her saa dog i min Stræben og paa den høitideligste Maade, da at \tlg{tilegne} ham det, det var, som jeg ønskede det, igjen en Concentration af Veneration. Men det lader sig ikke gjøre. Min Gang gjennem Livet er for tvivlsom med Hensyn til om jeg skal nyde Ære og Anseelse eller forhaanes og forfølges til at jeg kan dedicere nogen Levende mit Arbeide; og dernæst er der dog alt for betydelige Differentser mellem os. En saadan Dedication kunde baade paa den ene og paa den anden Maade vikle min Sag forkeert ind i Endelighedens Forhold.
\end{quote}
\dcite{NB3:36  (1847; Pap. VIII1A438; SKS 20, 263)}
% NB3:36  (1847; Pap. VIII1A438; SKS 20, 263)

\begin{quote}
Det Dialektiske i at Meddeleren maa have Øine i Nakken, betræffende Meddelelsens virkelige \tlg{Tilegnelse}.

\emph{[i margen:]} Det Svimlende og Uethiske i at have travlt for at meddele, saa man selv glemmer at være, det man lærer. Gud er ikke i Forlegenhed. Denne Forvirring fremkommer ved at betragte Tilværelsen i et phantastisk Lys.
\end{quote}
\dcite{Papir 364:365:23  (1847; Pap. VIII2B81.32; SKS 27, 398)}
% Papir 364:365:23  (1847; Pap. VIII2B81.32; SKS 27, 398)

\subsection{1848}

\begin{quote}
Christendommen skal ikke doceres. Derfor er det ogsaa Christus siger min Lære er Spise – den skal \tlg{tilegnes}, der skal existeres deri.
\end{quote}
\dcite{NB5:101  (1848; Pap. IXA105; SKS 20, 414)}
% NB5:101  (1848; Pap. IXA105; SKS 20, 414)

\begin{quote}
Ganske simpelt historisk betragtet (altsaa ikke som Troens Gjenstand) er dette maaskee det Punkt hvor Christi Heterogenitet fra ethvert Menneske viser sig stærkest: det at Folket han tilhørte gik under. Ellers er dog Forholdet dette at selv den eminenteste i et Folk, om han saa er død som Martyr saa \tlg{tilegner} Folket sig ham. Men dette er vistnok aldrig faldet Jøderne ind med Christus. Altsaa Folket gik under: dette er Udtrykket for at Christus var den Enkelte der er mere end hele Generationen. Og Folket gik ikke under og forsvandt saa af Historien – nei det bliver staaende i Undergangs-Situation, udtrykkende Undergangen: dette er (ad modum militair Honeur) guddommelig Honeur for Christus i Historien. Alle Følgerne af hans Liv ere dialektisk ikke nær saa mærkelige som denne Følge. Det Forunderlige, Gribende i, at Forsynet maa vaage over bestandigt at holde eet Folk paa samme Punkt: Undergangens, og det Aarhundrede efter Aarhundrede, som i et evigt Omqvæd gjentagende Respekt-Udtrykket for Christus.
\end{quote}
\dcite{NB6:94  (1848; Pap. IXA249; SKS 21, 69)}
% NB6:94  (1848; Pap. IXA249; SKS 21, 69)

\subsection{1849}

\begin{quote}
Afladets og Desliges Oprindelse – dets Udartelse er jo noget Andet – var dog som Alt i Middelalderen: Barnlighedens Misforstaaelse. Naar man læste den og den Bøn saa tidt saa fik man saa og saa tidt Aflad. (for Exempel. confer Liguori Betrachtungen und Gebetbuch Aachen 1840. pagina 599. note.) Men gjorde min Fader ikke lige saa med mig, da jeg var Barn, han lovede mig 1rigsdaler for at læse en af Mynsters Prædikener høit for ham, og 4rigsdaler, hvis jeg vilde skrive den Prædiken op, som jeg havde hørt i Kirken. Jeg gjorde det rigtignok ikke; jeg erindrer endogsaa, at jeg foreholdt ham det urigtige i, at ville lokke mig paa den Maade, fordi han vidste jeg vilde gjerne have Penge. Men Feilen laae dog egentlig ikke hos Fader, snarere hos mig, at jeg aldrig har været Barn. I Forhold til et Barn, som jo ikke har bedre Forstand paa det Høieste, er det jo slet ikke urigtigt; man regner paa at Barnet paa den Maade efterhaanden skal komme i de høiere Forestillingers Magt, og at det saa selv i en senere Tid skal have Leilighed til i dybere Forstand at \tlg{tilegne} sig dem.
\end{quote}
\dcite{NB10:59  (1849; Pap. X1A137; SKS 21, 289)}
% NB10:59  (1849; Pap. X1A137; SKS 21, 289)

\begin{quote}
Derved er det, at man har afskaffet Christendommen, at man overalt har trykket Personligheden tilbage. Man synes at frygte at et Jeg skulde være et Slags Tyrannie, og derfor skal ethvert Jeg nivelleret trykkes tilbage bag en Objektivitet. Jeg skal ikke have Lov til at sige: jeg troer der er en Gud – jeg skal sige: dette er Christendommens Lære: jeg troer, men saa er dette jeg et universellere, ikke mit personlige. Overalt Lære, Objektivitet, og overalt forhindres det, at man skulde faae Indtrykket af, at et Menneske umiddelbart forholder sig til Gud. – Bruger en Præst stundom sit Jeg paa Prædikestolen, saa tilgives det, fordi hans Jeg paa Prædikestolen dog ikke antages at være i strengeste Forstand hans personlige Jeg, men et Slags dramatisk jeg, eller et jeg qua Embedsmand.

Som de lærde Commentarer saa ofte lægge sig forstyrrende mellem en Forfatter og de Læsere, han helst ønskede; som Dag-Pressen har lagt sig forstyrrende mellem den egentlig Literatur og Læseren: saaledes har man skudt overalt Objektiviteten forstyrrende mellem Gud og Virkeligheden, praktiseret Gud langt, langt bort. Istedetfor at Læren skal være det Objektive, og saa skal mit jeg personligt \tlg{tilegne} sig den og jeg altsaa tale i første Person i Virkeligheden, vil man at jeg skal skaffe mit Jeg bort og tale objektivt. Thi det der skal være det Regjerende i Verden, skal ikke være Gud, som virker paa Jeget, men skal være et Objektivt, et Abstraktum, til hvilket de enkelte Jeg'er forholde sig som Blade paa et Træ, som Dyr under en Art.
\end{quote}
\dcite{NB13:76  (1849; Pap. X2A145; SKS 22, 319)}
% NB13:76  (1849; Pap. X2A145; SKS 22, 319)

\begin{quote}
Hvad jeg altid har sagt, og Anti-Climacus da især indskærpet: det Christeliges Vanskelighed fremkommer egentlig først, naar det skal sættes sammen med den Enkelte, naar den Enkelte, Du og Jeg, skal gjøre Alvor af at \tlg{tilegne} sig det, at turde sige: det er mig, det angaaer; thi da er det Christelige altfor høit, og Forargelsen ikke til at undgaae.

Ligesom man nu undgaaer den anden det Christeliges Vanskelighed: Samtidigheden og forvandler det til noget Forbigangent, saaledes undgaaer man den Vanskelighed med den Enkelte, ved bestandigt at underskyde et Objektivt. Naar jeg skal sige: Som en Brudgom elsker Christus mig, mig S. A. Kierkegaard, eller mig, H. Martensen, eller mig, J. P. Mynster – ja saa kniber det. Hvad gjør man saa? Saa underskyder man: Kirken; som en Brudgom elsker Christus Kirken, der er Bruden. See, nu gaaer det ganske behageligt. Et saadant uhyre Noget som Kirken, Kirken, der nu har bestaaet i 1800 Aar, bestaaet, som det hedder, af Milioner og Milioner – et saadant uhyre Noget, det synes at være commensurabelt, i almindelig menneskelig Forstand commensurabelt for Christus – og saa er jo Forargelsen gaaet ud.

Og nu tænke man sig hvilken Confusion! Naar jeg foredrager det om den Enkelte, da skal det være subjektivt Liebhaberie – men det andet det er den objektive Christendom! O, I Gavtyve! eller rettere sagt, de vide end ikke selv, hvor snildt deres Taktik er, for at underminere al Christendom

Til min store Glæde læser jeg idag hos Luther (i Prædikenen paa 20 Søndag efter Trinitatis Evangeliet, Kongen, der gjør sin Søns Bryllup) hvor han udvikler, hvorledes »b Kjød og Blods Blindhed og Forstokkelse vældigt strider mod dette om, at Christus som en Brudgom elsker mig og Dig«, jeg læser disse Ordc:.... Verden indrømmer vel ogsaa tilnøds, at Christus er en smuk, ædel, from og trofast Brudgom, og hans Kirke \emph{en herlig, salig Brud.} Men naar Enhver for sig skal troe, at han ogsaa hører Christus til, samt at Christus bærer saadan hjertelig Kjerlighed til ham – da gaaer Alting istaae.«

Gud skee Lov for Luther! Det er dog immer en god Hjælp mod den næsten afsindigt opblæste dogmatiske og objektive Indbildskhed, der ved at gaae videre, afskaffer Christendommen

\emph{[i margen:]} \emph{Det Christeliges Vanskelighed, og vor Tids Underfundighed.}

\emph{[i margen:]} den gamle Adam

\emph{[i margen:]} pagina 564. anden Spalte
\end{quote}
\dcite{NB14:57  (1849; Pap. X2A231; SKS 22, 378)}
% NB14:57  (1849; Pap. X2A231; SKS 22, 378)

\begin{quote}
\emph{Epistelen paa første Søndag i Advent.}

Det Hele er egentlig eet Billede: at staae op – og klæde sig paa (iføre sig Christum).

Livet sammenlignes med en Dag; man vaagner (Natten er forbi – Dagen er kommet). Saa staaer man op, klæder sig paa (iføre sig Christum).

At iføre Christum er saa \emph{deels} med Hensyn til Forsoningen, at \tlg{tilegne} sig hans Fortjeneste (i Parabelen om Kongen som gjorde sin Søns Bryllup, at der var En, som ikke var iført Bryllups-Klædningen) \emph{deels} at eftertragte at ligne ham, fordi han er Forbilledet og Exemplet. Dette er et substantielt Udtryk i Retning af Inderliggjørelse. Ligesom det Udtryk, han bruger om sin Lære, at den er Spise, er det stærkeste Udtryk for \tlg{Tilegnelsen}, saaledes er det at iføre sig Christum det stærkeste Udtryk for at Efterligningen skal være efter den størst mulige Maalestok; der siges ikke om Christus: ham skal Du stræbe at efterligne (naar der tales saaledes, ligger deri indirecte, at de To dog væsentligen blive ulige) nei Du skal iføre Dig ham, iføre Dig ham – som naar En gaaer i laante Klæder, dette er satisfactio vicaria), iføre Dig ham – som naar En aldeles skuffende ligner en Anden, ikke blot stræber at efterligne ham, men \emph{gjengiver} ham. Christus \emph{giver} Dig sin Klædning (satisfactio) og fordrer nu, at Du skal \emph{gjengive} Ham.
\end{quote}
\dcite{NB14:80  (1849; Pap. X2A255; SKS 22, 391)}
% NB14:80  (1849; Pap. X2A255; SKS 22, 391)

\begin{quote}
(Ein, Zwei, Drei

1

Genierne ere som Tordenveir: de gaae mod Vinden; forfærde Menneskene; rense Luften.

Det »Bestaaende« har opfundet adskillige Torden-Afledere mod eller for Genier: det lykkes – saa meget desto værre for det Bestaaende; thi lykkes det een Gang, to Gange, tre Gange – »det næste Tordenveir« bliver desto forfærdeligere.

Af Genierne er der to Arter. Den første Arts Charakteristiske er Bulderet, Lynet derimod er sparsomt, og slaaer sjeldent ned. Den anden Art har en Bestemmelse af Reflexion i sig, ved hvilken de tvinge sig eller tvinge Bulderet tilbage. Men Lynet er saa desto intensivere; med Lynets Hurtighed og Sikkerhed rammes de udseete enkelte Punkter – og dødbringende.

2

»Havde Apostelen Paulus noget Embede?« Nei, Paulus havde ikke noget Embede. »Havde han da et Levebrød?« Nei, han havde ikke noget Levebrød. »Tjente han da paa anden Maade mange Penge?« Nei, han tjente paa ingen Maade Penge. »Var han da idetmindste ikke gift?« Nei, han var ikke gift. »Men saa var Paulus jo ingen alvorlig Mand!« Nei, Paulus var ingen alvorlig Mand

3

Naar en Mand har Tandpine, siger Verden: stakkels Mand! Naar en Mand er i Pengeforlegenhed, siger Verden: stakkels Mand! Naar en Mands Kone bliver ham utro, siger Verden: stakkels Mand! – Naar Gud lader sig føde, bliver Menneske og lider, siger Verden: stakkels Menneske! Naar en Apostel i guddommeligt Ærende har den ære at lide for Sandhed, siger Verden: stakkels Menneske! – – stakkels Verden!!!

4.

I den pragtfulde Domkirke træder en stadselig Overhofprædikant, det dannede Publikums Udvalgte, frem for en udvalgt Kreds af de Udvalgte, og prædiker rørt – jeg siger »rørt« jeg siger ikke »tørt« nei han prædiker rørt over disse Apostelens Ord: Gud har udvalgt det i Verden Ringe og Foragtede – og Ingen leer!

5.

Eet er at profitere (profiteri) en Kunst, en Videnskab, en Tro – et Andet at have Profit deraf.

6.

Naar man spørger en Mand »veed Du Det og Det« og han da \emph{strax} svarer Ja eller Nei, saa er dette Svar et \emph{populairt} Svar, beviser, at han er en Eenfoldig, Seminarist og Deslige. Varer det derimod for Exempel \emph{10 Aar} inden Svaret kommer; kommer det saa i Skikkelse af en udførlig Afhandling, der nøiagtigt som Holophernes siger tager Tempo'et i agt »Ein, Zwei, Drei,«; og er det saa ved Slutningen af den udførlige Afhandling \emph{ikke ganske klart,} om han veed det eller ikke: saa er dette et ægte speculativt Svar, beviser, at den Spurgte er Professor i Speculation, eller idetmindste saa udspeculeret, at han burde være det.

7.

Den Reflexion som findes under Numero 1 af Reflexioner, som ligge paa et løst Papir i Pulten, den er om en Digter der vil forvexles med et Sandheds-Vidne.
\end{quote}
\dcite{Papir 385:397  (1849; Pap. X6B253; SKS 27, 476)}
% Papir 385:397  (1849; Pap. X6B253; SKS 27, 476)

\subsection{1850}

\begin{quote}
Spørgsmaalet i Forhold til Christendom er altid, hvor langt skal et Menneske omdannes i Retning af at blive Aand, for at han tør \tlg{tilegne} sig Naaden.

Som Forholdene nu er, er det for galt; det er egentlig blot raffineret Hedenskab. Man forbliver i alle Sandselighedens og Verdslighedens Bestemmelser, har sit Liv deri; erkjender saa villigt at man er langt tilbage (det kaldes Anger) men bliver paa samme Sted; anbringer saa »Naaden« som – den nye Lap paa det gamle Klædebon.

Paa den anden Side, skal et Menneske afgjørende udvikles til at være »Aand«, for at turde \tlg{tilegne} sig »Naaden«: saa, Gud veed, hvor mange der vil være i hver Generation, der saa virkelig føle nogen Trang til eller have nogen Brug for Christendom Skal jeg virkelig afdøe fra mine kjereste Ønsker, forsage Alt, hvad der glæder det jordiske Menneske saa bliver jeg jo, menneskelig talt, saa ulykkelig som det er muligt – og saa er Spørgsmaalet, om jeg saa er »Aand«, at jeg virkelig har Brug for Christendom

Det er en skammelig Misbrug der drives med den Inddeling: Loven er det Forfærdende – Evangeliet det Beroligende. Nei, Evangeliet selv er og skal være i sit Første Forfærdende. Havde det ikke været saa, hvor i al Verden kunde det da være gaaet Christus som det gik ham, at da Han sagde: kommer hid – da bleve Alle borte, de flyede ham.

Men klart bliver det mig mere og mere, at kun en Apostel kan i strengere Forstand forkynde Christendom, thi kun han har Myndighed til at være saaledes streng. Et Menneske har ikke saaledes Myndighed, han maa derfor slaae af. Kun Den, der selv i strengere Forstand er forvandlet til »Aand« kun han kan ikke mere forstaae, vil ikke forstaae det usalige Vrøvl, den Skrøbelighed, hvori vi Andre ere hildede, som gjør at vi meget for meget skaane os selv, og for tidligt \emph{hvile} i Naaden, og hvile i den fra Stræben, istedetfor at hvile i den til fornyet og fornyet Stræben.
\end{quote}
\dcite{NB15:109  (1850; Pap. X2A445; SKS 23, 76)}
% NB15:109  (1850; Pap. X2A445; SKS 23, 76)

\begin{quote}
»Standsning« som fornøden for at blive opmærksom paa Christendommen og for at blive Christen.

De fleste Mennesker leve fra Vuggen til Graven ustandseligt og ustandsede hen i det Ustandseliges Medium (Timeligheden, den blotte Qvantiteren og saa videre). Saa kommer endelig Døden og standser dem – og nu blive de opmærksomme paa Christendom, fortryde at de ikke tidligere have \tlg{tilegnet} sig den, naae ved Hjælp af denne Fortrydelse et Forhold til Christendom, og døe saa.

I levende Live frygter det naturlige Menneske ikke Døden mere end han frygter Standsningen. Nu, Døden og Standsningen have ogsaa meget tilfælles. Standsningen er som naar Fisken tages op af Vandet og skal aande i Luften. Det naturlige Menneske gyser for dette andet Element, for den uhyre Magt, der boer i »Standsningen«, og om hvilken han vel forstaaer, at faaer den først den mindste Magt over En, er det uberegneligt, efter hvilken grændseløs Maalestok den faaer Magten. Og det Grændseløse frygter atter det naturlige Menneske som han frygter Døden. Det Grændseløse, det Uendelige, det Eviges Stillestaaen i Standsningen det er ligesom at omkomme for Den, hvis Element er »til en vis Grad«.

Christendom forvandlet til Lære kan ypperlig løbe ind med i Timelighedens og Ustandselighedens Travlhed og Qvantiteren: det fører ikke til Christendom Paa den anden Side: »Standsningen« kan ogsaa blive en § med i Christendommens Lære: det hjælper lige saa lidet.

Der er i Fortællingen om Røveren paa Korset noget Typisk fra en Side som man ikke fremhæver. Alle, alle ere de faldne fra, selv Apostelen har fornegtet Christus – den eneste samtidige Christen med Christus var Røveren paa Korset: saa, om jeg saa tør sige, uendelig meget for høi er Christendom for Menneskene, at (naar Situationen er den meest anstrengede nemlig Samtidighed med Christus) i levende Live end ikke Apostelen kan holde ud med Christus. Kun en Røver – en døende Røver, kun ham hjælper Syndens Bevidsthed og Dødens Situation til at holde Christum fast.

Hvad er det der gjør Døden til »Situationen« for det at blive Christen? Det er det Afsluttende og Afsluttede at det nu er forbi; det hjælper til absolut at byde paa det Absolute. Er det en Døende, hvem man ikke tør sige til, at det er Døden, men hvem man holder hen med at han nok kommer sig, saa er Døden heller ikke »Situationen«.

Selv i Forhold til den sandeste og alvorligste Christen, der har levet, vil det dog være saa, det troer jeg, at først i hans Død blev det absolut sandt at han er Christen. Vi forholde os alle saadan til Christendommen; vi have saa at sige budt paa den. Men jo høiere En i levende Live har budt, desto sandere vil ogsaa hans sidste Bud i Døden blive, desto nærmere vil han være ved at naae i det Øieblik at byde ganske absolut. I levende Live har kun En budet absolut paa den, og holdt ud (skjøndt han vedblev at leve) i ethvert Øieblik at udtrykke og staae ved det eet og samme absolute Bud: han, der selv var det Absolute. Vi Mennesker behøve Understøttelse; og i Dødens Øieblik er et Menneske ved Situationen hjulpet til at blive det Sandeste, han kan blive.
\end{quote}
\dcite{NB18:4  (1850; Pap. X3A47; SKS 23, 256)}
% NB18:4  (1850; Pap. X3A47; SKS 23, 256)

\begin{quote}
Reduplicationens Strenghed.

Naar jeg nu er overbeviist om en Sag; naar jeg da begeistret udtaler den og seer hvorledes hele Mængden gaaer ind derpaa: er det saa ikke for meget forlangt af mig at jeg istedetfor udenvidere at glæde mig derover skal først see efter, blive betænkelig, ja vel endog fatte en Mistanke, at jeg er paa en gal Vei, at det er Misforstaaelse, siden det gaaer saa gesvindt med at blive forstaaet, og det Sande bliver saaledes modtaget med aabne Arme?

Dog saa streng er Sandheden.

\emph{[i margen:]} Dog maa man huske paa, at her ogsaa kan fremkomme noget Selvplagersk; det kan jo dog ogsaa være sandt i Nogle, at de virkelig \tlg{tilegne} sig det Foredragede. Luther siger til et christeligt Liv hører: Tro – Kjerligheds-Gjerninger – og saa Forfølgelse for den Tro og Kjerlighed (Stedet er bemærket i mit Exemplar af hans Prædiken.) Han siger ogsaa et andet Sted (der ligesaa er bemærket i mit Exemplar) at hvor Forfølgelse ikke er, der er Feil i Forkyndelsen. Imidlertid mener jeg dog her atter at man bør være forsigtig med det Selvplagerske; thi ellers kunde det jo blive Opgaven i at bringe Læren ind i Verden, i eetvæk at passe paa, at den ikke kom ind men blev stødt tilbage. Her maa man befale sig Gud i Vold og bede ham vaage over En.
\end{quote}
\dcite{NB18:74  (1850; Pap. X3A124; SKS 23, 303)}
% NB18:74  (1850; Pap. X3A124; SKS 23, 303)

\begin{quote}
\emph{Et Memento for geistlige Talere.}

Det burde stadigt indskærpes, hvad Quinctilian etsteds siger, at han har seet Skuespillere, der lode sig saa stærkt henrive af deres tragiske Rolle, at de endnu hjemme græd. Fremdeles siger Quinctilian, at han selv, da han engang havde paataget sig at vække hos Alle en bestemt Lidenskab, havde \tlg{tilegnet} sig den i den Grad, og følt sig i den Grad overrasket, at han ikke blot selv havde grædt derover, men endog var blegnet, og ganske havde befundet sig som En, der ligger under for Smerten.

Man tænke ogsaa her paa, hvad Hamlet siger om Skuespilleren, der taler om Hecuba.

Enhver sand Digter vil have erfaret det Samme.

Men det Farlige ved Taleren er, at han ikke gjør den Distinction, at det muligt kun er digterisk han sætter sig ind deri.
\end{quote}
\dcite{NB21:37  (1850; Pap. X3A479; SKS 24, 33)}
% NB21:37  (1850; Pap. X3A479; SKS 24, 33)

\subsection{1852}

\begin{quote}
\emph{At vi opdrages fra Barn i Christendom }

har dog den gode Side, at vi, hvis vi virkelig skulle blive Christne, komme til at opleve en Analogie til de med Christus Samtidige. De nærede først jordiske Forventninger, – og saa vendte Alt sig om, og nu blev det Alvor med i Aand og Sandhed at blive Christen. Saaledes naar En er fra Barn opdraget i Christendom. Barnet \tlg{tilegner} sig Christendommen som et jordisk Evangelium – og nu kommer Forfærdelsen for det Menneske i en senere Alder, naar han skal have Aands Indtrykket af Christendommen.
\end{quote}
\dcite{NB25:108  (1852; Pap. X4A539; SKS 24, 520)}
% NB25:108  (1852; Pap. X4A539; SKS 24, 520)

\subsection{1853}

\begin{quote}
\emph{Hvorledes »Christenhed« læser det Nye Testamente især Evangelierne.}

Man deler saaledes: ethvert Ord af Christus i Retning af Anstrengelsen, Lidelsen som er forbunden med det at være Christen, derom siger man: dette angaaer ikke os, det er sagt særligen til Disciplene. Ethvert Trøstens, Forjættelsens Ord derimod tager man sig til Indtægt. Om der i det Nye Testamente staaer: Christus sagde til Apostlene, eller særligen til Apostlene, om det staaer eller ikke, gjør »Christenheden« Intet til Sagen. Nei, naar der i een og samme Forbindelse, om hvilken der i det Nye Testamente siges Christus sagde til Apostlene, findes et Ord i Retning af Anstrengelsen og i samme et Forjættelsens Ord, saa vil det Faktiske i »Christenhed« være dette: hiint første Ord høres næsten aldrig – det var jo ogsaa sagt til Apostlene, det andet høres idelig og idelig – men det var jo ogsaa sagt særligen til Apostlene.

Naar Sagen ikke var saa alvorlig, kunde man fristes til at sige, at det er Latterligt, at see saadanne Christne som i »Χsthed« i Gjennemsnit trøste sig ved saadanne Forjættelser og Trøstens Ord. Thi om det nu end ikke var et Snyderie, at \tlg{tilegne} sig disse, de passe dog alligevel ikke for dem. O, thi i Aandens Verden er en evig Conseqvents: saadanne, saa stærke Trøstens Ord passe ogsaa egentlig kun for den tilsvarende Anstrengelse.

»Christenhed« tillyver sig saaledes Trøsts, Begeistrings-Ord – og til syvende og sidst bruge de dem dog alligevel ikke. Thi i Aandens Verden er Alt Sammenhæng: skal jeg i Sandhed benytte en saa uhyre høi Trøst, saa vil Trøsten først bevirke at mit Liv bliver anstrenget. Trøsten vil i saa Fald først give mig Anstrengelse.

Men »Christenhed« lever i Indbildning og Drømme. Og det forstaaer sig: Penge kan man snyde sig til og Magt og saa videre og det er virkelig Penge og Magt man faaer; men Evighedens Trøst kan man ikke snyde sig til.

Dog mundus vult decipi – ikke blot den bliver bedraget, nei, den \emph{vil} bedrages, den bliver rasende, hader, forfølger, ihjelslaaer Den, som ikke vil bedrage den. Derfor vilde »Christenhed« vel blive rasende, om En oplyste, at den Art af Christendom er en Chimaire. Nei, Indbildningen at være Christen, den vil Christenheden have – og saa er det Sted, hvor det egentlig skal komme En til Gode, om man har været en Christen eller ei – Evigheden, hvor ingen Indbildning hjælper.

Dog, man har nu engang faaet Christendommen vendt saaledes, at Menneskene virkelig troe, at Christendommen hører med for ret at nyde dette Liv, man behøver denne Beroligelse Evigheden betræffende for ret at nyde dette Liv, og den skjenker jo Christendommen En. Hvilken uhyre Bagvendthed!
\end{quote}
\dcite{NB27:77  (1853; Pap. X5A78; SKS 25, 191)}
% NB27:77  (1853; Pap. X5A78; SKS 25, 191)

\subsection{1854}

\begin{quote}
\emph{»Det har Christus særligen sagt til Apostlene«} \emph{.}

Saaledes hedder det gjerne betræffende Alt, hvad der af Christi Udsagn ikke convenerer os.

Convenerer det Udsagte os, saa gjør det slet ikke Noget, at det, endog af Christus siges særligen til Disciplene for Exempel de Ord: salige de Øine, som see Det I see.

I Luthers Prædiken over Evangeliet paa 19de Søndag efter Trinitatis lige i Slutningen seer jeg ogsaa, at han uden videre \tlg{tilegner} sig og enhver Christen de Ord af Christus til Apostlene: annammer den Hellig-Aand; hvem I forlade Synder og saa videre.

Men paa den Maade maatte vi jo ogsaa fordre Miraklets Gave.
\end{quote}
\dcite{NB28:80  (1854; Pap. XI1A26; SKS 25, 278)}
% NB28:80  (1854; Pap. XI1A26; SKS 25, 278)

\begin{quote}
\emph{Christelig Revision.}

Hvad i Endelighedens Verden Penge er, er aandeligt, Begreberne. Det er i disse alle Omsætningerne skee.

Naar det nu saaledes gaaer hen fra Slægt til Slægt, at Enhver overtager Begreberne som han fik dem fra den foregaaende Slægt – og saa anvender han sin Dag og sin Tid til at nyde dette Liv, arbeide for de endelige Formaal og saa videre    : saa afstedkommes det kun altfor let, at efterhaanden Begreberne forvanskes, blive noget ganske Andet end de oprindeligen vare, komme til at betyde noget ganske Andet, blive som falske Penge – medens dog ganske roligt alle Omsætninger vedblive at gjøres deri, hvilket imidlertid ikke saaledes berører Menneskenes egoistiske Interesse, som naar falske Penge komme ud, især naar Begrebs-Forfalskningen er just i Retning af den menneskelige Egoisme, saa Den, der egentlig bliver narret er om jeg saa tør sige, den anden Part i Christendommens Anliggende: Gud i Himlene.

Dog til den Forretning: at revidere Begreberne har Ingen Lyst. Enhver forstaaer mere eller mindre tydeligt, at at anbringes paa den Maade, til den Forretning, er omtrent det Samme som at være offret, betyder, at Ens Liv bliver saaledes lagt under Beslag, at han ikke kan komme til, hvad Mennesket naturligt har Lyst til, at beskjeftige sig med de endelige Formaal. Nei, at tage det saa flygtigt som muligt med Begreberne, og saa jo før jo hellere ind i Livets Concrete, eller i ethvert Fald da ikke at tage det for nøie med Begreberne, ikke nøiere, end at man kan med fuld Fart gaae med i Livets Concrete: det er det Menneskelige.

Dog Revision behøves, og med hver Aarti mere og mere.

Saa maa Styrelsen bemægtige sig et Individ, der bliver at bruge dertil.

En saadan Revisor er naturligviis noget ganske Andet end det hele Sludder-Compagnie af Præster og Professorer – dog er han heller ikke en Apostel, snarere lige det Modsatte.

Det Revisoren behøver er netop hvad Apostelen egentlig ikke har Brug for: Intellectualitet, en eminent Intellectualitet, videre et uhyre Kjendskab til alle mulige Gavtyvestreger og Falsationer, fast som var han selv den durchdrevneste af alle Gavtyve – hans Forretning er jo just i Retning af at »kjende« Falsknerierne.

Da denne hele Viden er af saa uhyre tvetydig Art, at der jo ved den kunde afstedkommes den størst mulige Forvirring, bliver saa Revisoren ikke behandlet som Apostelen. Ak, nei Apostelen er betroet Mand, Revisoren er sat under den skarpeste Kontrol. Jeg har bestandigt kun eet Billede, men det er saa betegnende. Tænk, at Banken i London blev opmærksom paa, at der circulerede falske Sedler – og som vare saa godt gjorte, at det var til at fortvivle over, at sikkre sig at kjende dem, og sikkre sig fremtidigen mod Eftergjørelse. Trods hvad Talenter der kunde være mellem Bankens eller Politiets Personale – der var dog kun Een, som just i denne Retning var ubetinget Talent – men det var en af de dømte Personer, en Forbryder. Han bliver saa at bruge, men han bruges ikke som betroet Mand. Han sættes under den frygteligste Kontrol, med Døden over Hovedet maa han sidde og have med alle disse Penge-Masser at gjøre, han visiteres paa Kroppen for hver Gang og saa videre o: s: v:

Saaledes med den christelige Revisor. Har Apostelen den Opgave at forkynde Sandheden, Revisoren har den Opgave: at opdage Falsknerierne at gjøre dem kjendelige og derved umulige; er Apostelens personlige Egenskab den ædle rene Eenfold (hvilket er Betingelsen for at være den Hellig-Aands Redskab) Revisorens er den tvetydige Viden; er Apostelen i een og kun i god Forstand ganske i Styrelsens Magt, Revisoren er i tvetydig Forstand ganske i Styrelsens Magt; har Apostelen med al sin Anstrengelse og Arbeiden dog ingen Fortjeneste for Gud, Revisoren endnu mindre, kunde (hvis det ellers var muligt) umuligt faae den, da han har et Negativt at aftjene, saa han væsentligen er en Poeniterende: men offrede ere de væsentligen begge; og i Naade ere de begge valgte af Styrelsen, thi det er ikke i Unaade at En vælges til Revisor; og som det begynder med Apostelen, saaledes kan naturligviis Revisoren først komme mod Slutningen, da han jo har til Forudsætning Udbredelsen; og har Apostelen sit Navn af at udsendes fordi han fører fra Gud ud efter, saa er Eftersynets Opgave at trænge igjennem Falsknerierne, at føre tilbage til Gud.

Apostle kan der aldrig mere komme igjen, saa maatte ogsaa Christus kunne komme igjen i anden Forstand end ved sin Gjenkomst. Christi Liv paa Jorden er Christendommen. Apostelen betyder: nu er den bragt an, fra nu af har I Mennesker selv at overtage den, men under Ansvar.

Saa overtog Menneskeheden den. Og er det end evig Løgn, at Christendommen er perfectibel, vist er det, Menneskeheden udfoldede en stigende Perfectibilitet i – at forfalske Christendommen.

Ligeoverfor dette Falsknerie selv om Gud vildea kan han ikke bruge en Apostel af den Grund, at Christenheden ved sin Falsknen har bragt sig paa en saadan Afstand fra Gud, at der ikke kan være Tale om den, hvis jeg saa tør sige, tillidsfulde Henvendelse til Menneskene. Nei, som Christenheden er Falsknen, og som Synden nutildags hovedsaglig er Klogskab, saa er ogsaa fra Styrelsens Side – hvem Mennesket ved sin Falsknen fjernede fra sig – alt Mistroiskhed: der kommer ikke glade Udsendinge fra Gud, saa lidt som man hilser Politiet som saadant, nei der kommer kun Kjendere af Uredelighed, og selv disse blive, da de jo dog væsentligen tilhøre den almindelige Uredelighed, behandlede som Tvetydige af Styrelsen.

Christenheden nutildags er glad og fornøiet. Det er ikke saa ganske sjeldent, at man seer Sagen fremstillet saaledes, at der nu vil begynde en ny Epoche, komme nye Apostle – thi Christenheden, ja den har naturligviis gjort sine Ting ypperligt, den har nu saa fuldkomment indøvet og \tlg{tilegnet} sig hvad der ved Apostlene blev anbragt, at nu maa vi gaae videre. Sandheden er den, at Christenheden har gjort det saa skidt og saa gavtyveagtigt som muligt og at forvente nye Apostle (hvis der ellers var Sandhed i denne Tanke) er den meest qvalificerede Uforskammethed.

\emph{[i margen:]} og selv om der ikke var anden Hindring
\end{quote}
\dcite{NB32:117  (1854; Pap. XI2A36; SKS 26, 205)}
% NB32:117  (1854; Pap. XI2A36; SKS 26, 205)

\begin{quote}
\emph{Christendom (Ironie)}

Vilde Nogen sige: »Christendommen er jo af Gud en Ironie over os Mennesker,« saa vil jeg svare: nei, min gode Mand, men vi Mennesker have det i vor Magt at forvandle Christendommen til en Ironie, en bidende Ironie.

Sagen er ganske simpel. I sin Majestætiskhed tager Gud Tonen saa høit, at hvis et Menneske ikke vil slippe sin endelige Forstandighed, ikke vil forsage den jævne, nydelsessyge Middelmaadighed – saa er Det Gud kalder Hjælp, Frelse, Naade og saa videre den meest bidende Ironie, saa man forsaavidt ikke kan fortænke den endelige Forstand i at sige: nei, Tak, den Hjælp, Frelse, Naade vil jeg dog nok helst være fri for.

Saasnart derimod et Menneske vil som Gud vil, vil lade sig belære af Gud, om Forholdet mellem Timelighed og Evighed, vil see, at \tlg{tilegne} sig Guds Forestilling om Evigheden, vil troe Gud – saa er naturligviis Gud idel idel idel Kjerlighed og Forbarmelse.

Men Christenhedens Udvei at lade Gud være et Snøvl, der vrøvler Noget om en Frelse, Naade og saa videre saaledes, at al Verdslighed finder at der dog er god Mening i denne Tale, denne Udvei er Majestæts-Forbrydelsen.

Som jeg andetsteds har bemærket, der gives intet Hedninge-Folk, der har dyrket en saa latterlig og saa modbydelig Guddom, som »Christenhed«, der dyrker og tilbeder – et Vrøvlehoved; saa passede dog selv den raaeste Hedning bedre paa, at hans Gud dog ragede Noget op over det at være Menneske, stod lidt høiere, medens Christenheden har excelleret i at gjøre Gud ganske menneskelig.

Maaskee vil man sige, at dette er en Overdrivelse. Men det er ikke Tilfældet, tværtimod lader Sagen sig forklare deraf, at Christendommen jo i sin Sandhed har det høieste Gudsbegreb, Gud er Aand, Kjerlighedens Majestæt – men saa vil ogsaa just Demoralisationen af dette Høieste give det allerlaveste Guds Begreb.

At sigte »Christenheden« for Løgn, Hyklerie og Deslige er derfor dog egentlig for høit, er at gjøre den til mere end den er. Nei, den meest passende Sigtelse er: det Hele er en Uanstændighed. Det er en Uanstændighed for en voxen Mand at ride paa Kjephest – og saaledes er det virkelig ogsaa en Uanstændighed at kalde noget Saadant, som det man declamerer under Navn af Christendom, at kalde det Christendom, hvor alle det Guddommeliges Mærker længst ere udviskede og hvor alle Begreber længst ere omstemplede til det menneskelige Vrøvl.
\end{quote}
\dcite{NB34:40  (1854; Pap. XI2A183; SKS 26, 352)}
% NB34:40  (1854; Pap. XI2A183; SKS 26, 352)

\begin{quote}
Børn og Fruentimmer; Manden;

Christendommen Alvoren.

Christendommen er Alvoren. Selvfølgeligt er derfor ogsaa Maalestokken lagt an paa: Manden; den christelige Fordring forholder sig til: Manden, endog til Guds Forestilling om hvad det har at betyde; Manden er Mennesket.

Men Enhver der har den mindste Praxis i Livet, i Forretnings-Livet, veed, at altid naar det gaaer løs paa at skulke sig fra Noget, saa bruges denne Taktik: at anbringe Fruentimmer og Børn.

Og her kommer en ny Forklaring af dette ene Samme, hvorledes Christendommen er blevet ved en Bagvendthed, som Menneskene behageligst kalde Perfectibilitet, blevet lige det Modsatte af hvad den er oprindelig.

Christendommen, hedder det, er jo for Mennesker; et Barn er jo ogsaa et Menneske – ergo er Barnet Maalestok for hvad Christendommen er.

Dette er nu i den Grad Løgn, at det meget mere ligefrem er Samvittighedsløshed at opdrage et Barn i Christendom; thi Barnet kan umuligt \tlg{tilegne} sig hvad Christendommen egentlig er (eller kan vel Barnet have nogen Forestilling om Arvesynden, Barnet, som jo opdrages til i Faderen og Moderen at see sine Velgjørere, at takke dem for den store Velgjerning at være blevet til og saa videre?) Altsaa maa man \emph{enten} gjøre Christendommen til noget Andet end den er for at Barnet kan \tlg{tilegne} sig den – og har man vel Lov til at bilde et Barn Noget ind under Navn af at undervise det; eller Barnet forvandler selv Christendommen, fordi det ikke kan det anderledes – men har man Lov til at foranledige et Barn til at sætte sig fast i Indbildninger, og har man Lov til at gjøre det under Navn af at undervise og belære.

Nei, men det Hele med Barnet i Forhold til Christendommen er en Gavtyvestreeg. Aldeles som i Forretningslivet, naar for Exempel Øvrigheden kommer og spørger efter Manden, som da, han seer at skulke bort og anbringer Fruentimmer og Børn, der skal see at røre og formilde: saaledes er Christenhedens senere Historie denne Gavtyvestreeg, at Manden er skulket fra det, har faaet Christendommen gjort til Noget for Børn og Fruentimmer. Og naar saa det er lykkedes, og Religionen er blevet Barne-Slik, saa kommer Manden hjem denne Skulke-Peer, og saa slikker han med, og taler sentimentalt om, at Religionen, Christendommen er især for Børn og Fruentimmer.

Hvilken Uforskammethed og hvilken nederdrægtig Løgn, at Religionen er Noget for Børn og Fruentimmer, formodentlig fordi dens Opgave er for let for den Stærke: Manden. Jo jeg takker, især for saadanne Mænder, som nu leve.

Nei, Religionen, Christendommen er en Idealitet, en Opgave, som den største Idealiteta

\emph{[i margen:]} af det at være Mand skal fast segne under. Saaledes er Christendommen oprindelig. Fra Østen kommer den. Og hvorledes var Forholdet der? Der var Manden Mennesket, Fruentimmer og Børn omtrent et Slags Husdyr. Men som hele den verdenshistoriske Proces er det Kunststykke at forkluddre alle Guds Anlæg og faae lige det Modsatte ud saaledes er det ogsaa forbeholdt den romantiske protestantiske Christenhed at gjøre Børn og Fruentimmer til »Mennesket« og Manden til et Nummer Nul. Jeg imødeseer, at Perfectibiliteten naaer: Fleer-Mænderiet, Modstykket til: Fleer-Koneriet, som der begyndtes med. Men denne Forandring at Børn og Fruentimmer ere blevne »Mennesket« har atter forandret Christendommen til lige det Modsatte af det Oprindelige: den er blevet idel Præsent – den var idel Opgave.
\end{quote}
\dcite{NB34:43  (1854; Pap. XI2A187; SKS 26, 355)}
% NB34:43  (1854; Pap. XI2A187; SKS 26, 355)

\begin{quote}
Manden; Qvinden; Barnet.

Christendommen.

I Grunden er det dog frygteligt, men sandt er det, og Udtrykket for, efter hvilken rædsom Maalestok det er sandt, at Christendommen slet ikke er til.

Saaledes er egentlig Forholdet i Christenhed, især i Protestantisme.

Mændene – og det endda saadanne usle Skrællinger og Spytteklatte, som det der nu kaldes Mænd, sammenlignet med Østens Idealitet af det at være: Mand – Mændene vende sig med en vis Stolthed og Selvfølelse fra Religionen og sige: Religionen (Christendommen) er Noget for Fruentimmer og Børn.

Og see saa hvad Sandheden er, Sandheden er, at Christendommen, som den er i det nye Testamente, er lagt an efter saadanne Proportioner, at den stricte taget ikke kan være Religion for Fruentimmer, høist paa anden Haand, og for Børn umuligt.

Dette er mit Vidnedsbyrd som Psycholog: ingen Qvinde kan udholde en dialektisk Fordoblelse, og alt Christeligt har et Dialektisk i sig.

For at skulle kunne forholde sig til den christelige Opgave fordres der at være Mand, der behøves Mandens Haardhed og Styrke for endog blot at kunne bære Trykket af Opgaven.

Et Gode, der er kjendeligt paa, at det gjør ondt; en Frelse, der er kjendelig paa, at den gjør mig ulykkelig; en Naade, der er kjendelig paa Lidelse og saa videre o: s: v: alt Sligt, og saaledes er alt Christeligt, kan ingen Qvinde bære, hun vil gaae fra Forstanden, hvis hun skal spændes i denne Anstrengelse.

Hvad Barnet angaaer er det naturligviis aldeles Vrøvl at det skal være Christen.

En Qvinde, og fremfor Alt et Barn, forholder sig til, aander i det Ligefremme. Er Noget et Gode, ja saa maa det være kjendeligt paa, at det gjør godt: det hjælper ikke; at ville tvinge en Qvinde (Barnet vil jeg slet ikke tale om) ind i: et Gode, som gjør ondt er at sprænge hende.

Læg blot Mærke til, hvoraf det vel kan komme, at ingen Qvinde kan taale Ironie, at Ironie i Forhold til hendes Lidenskab er at dræbe hende. Kommer det ikke deraf, at hun ikke kan bære et Dialektisk.

Jeg har i denne Henseende sandeligen taget store Philosophicum. Forsøg det: gjør en Pige ulykkelig og siig saa til hende: ja, men det er af Kjerlighed jeg gjør det – og Du sprænger hende, det brister for hendes Forstand. Lemp Dig saa efter hende, siig: jeg er en letsindig Slyngel – saa kan hun bære det, og hun vil kjerligt tilgive Dig. Men hun slap jo ogsaa for den dialektiske Fordoblelse.

Og saaledes med alt det Christelige. Kun Manden har fra Styrelsens Haand Haardførhed til at kunne bære et Dialektisk.

At skulle bære et Dialektisk er den intensiveste Qval er den intensiveste Qval der er muligt. Barnet, den lille Skjelms-Mester, er fuldkommen betrygget derimod, han kan ikke engang komme det saa nær, at han kunde gaae fra Forstanden derover, om Du end vilde tylde det af al Magt i ham. Qvinden kan komme det saa nær, at hun segner under det, eller at Forstanden, for at hjælpe hende ud af det, lister bort det vil sige hun gaaer fra Forstanden.

At skulle bære et Dialektisk er den intensiveste Qval der er muligt. Det sees da ogsaa let, at langt intensivere end det for Exempel at blive ulykkelig, er den Lidelse: at blive ulykkelig og at det tillige skal betyde just Ens Lykke, og saaledes paa ethvert Punkt. Derfor vil ogsaa Enhver (dersom der ellers lever Nogen, som duer til at forstaae sig paa Sligt) der forstaaer sig paa Sligt, naar han tænker denne Figur: en dialektisk Fordoblelse han vil (som man ved Synet af Marter-Redskaber uvilkaarligt ligesom hører en Martrets Skrig) han vil naar han tænker en Qvinde i Forhold dertil, uvilkaarligt høre dette Skrig: o, frels mig, frels min Forstand.

Saa er altsaa Sandheden af dette med Christenheden denne: at dette Høie, som er det Christelige, dette Høie, som ingen Mand end ikke den Gang da det at være Mand var en Idealitet og end ikke den af saadanne der havde meest Idealitet, som end ikke han overtog eller løftede paa, uden at det satte i Knæerne, dette Høie under hvilket (at jeg skal tage det stærkeste Udtryk) under hvilket selv Verdens Frelser segner, at Gud som er Kjerlighed dog kan forlade ham og gjøre det af Kjerlighed: dette Høie har man i Christenhed saa jævnt og hjerteligt sluddret ned i det Gade-Skidt der er den almindelige menneskelige Tanke-Gang, saa dette Høie nu endog er blevet for let og for lidt for den Art Paaklædninger, man nu kalder Mænd, og overlades til Qvinder og Børn, for hvem egentlig Religionen er.

I det nye Testamente er Forholdet dette: Manden er det lagt an paa; Religionen forholder sig til Manden; Qvinden participerer secundairt i Religionen per Manden; selv bære et Dialektisk kan hun ikke, men ved at være Vidne til hvorledes Manden løfter paa den Opgave, faaer hun dog et Indtryk mere end af det ganske Ligefremme. Barnet skjøtter sig selv, til hans Tid kommer. At ville tylde sand Christendom i et Barn er (hvis det ellers var muligt, hvis Barnets Natur ikke gjorde det umuligt at \tlg{tilegne} sig dette) er lige saa bestialsk som at ville (hvilket dog ofte nok gjøres) tylde Brændeviin i Barnet, fordi Forældrene selv drikke Brændeviin og den søde Glut skal have det lige saa godt som Forældrene. Og under Navn af Christendom at ville tylde Noget i Barnet som ikke er Christendom er jo uforsvarligt.

Men, som sagt, Christenheden har faaet Alt sat over i det Ligefremme – og derfor ganske rigtigt er »Barnet« blevet Maalestokken for det at være Christen! Fortræffeligt! Christenheden synes slet ikke at have mærket, at den just ved »Barnet« er stødt paa et ironisk Problem, et Problem, som saa godhedsfuldt er blevet besvaret, det Problem, hvad skal vi gjøre med Barnet, kan Barnet blive Christen – et Problem, til hvis Besvarelse det nye Testamente Intet indeholder, da det antog, at den Christne – ikke gifter sig. Lige saa ironisk som det er hvad Peer Degn siger om Saxo Grammaticus at han berigede det latinske Sprog med en Mængde Udtryk, lige saa ironisk er den Berigelse, hvorved Christenheden har beriget Christendommen: alt det meget Lærde og Dybsindige om Barnet, Barnets Christendom, Barnets Daab, Tro og saa videre Vær i Forhold til et Barn kun Trediemand, og Du skal see, at Alt stiller sig som det nye Testamente tænker sig det. Men da de rare Christne hittede paa, selv paa en Frisk at ville begynde paa, hvad just Christendommen spærrede af for: saa fik Barnet en anden Betydning. Og paa »Barnet« blev saa – ja det er ganske i sin Orden – Christendommen aldeles endevendt, blev lige stik det Modsatte af hvad den er i det nye Testamente, Barne-Slik, saa forsaavidt Mændene som de nutildags ere, havde Ret i at vende sig bort fra dette, betragtende det som Noget der kun var for Qvinder og Børn, Noget der ækler en Mand ligesom Snakken og Jav'et og Temperaturen i en Barselstue.

Nei, lad det blive igjen som engang, lad det sætte Manden i Knæerne at overtage og bære Opgaven, lad Qvinden gyse ved at see, hvor tung den er – og Barnet? ja, lad det blive som engang, lad os blive fri for Børne-Avlen af Christne: saa er det muligt at Christendommen kan sees igjen. Men ellers er det umuligt; og jeg for mit Vedkommende begriber ikke, hvorledes det er muligt, at Nogen med Indtrykket af Christi Liv og hvad Evangelierne forstaae ved at være Christen, og med Indtrykket af Christi Fordring om Efterfølgelse, er kunnet falde paa – at gifte sig.
\end{quote}
\dcite{NB35:5  (1854; Pap. XI2A192; SKS 26, 366)}
% NB35:5  (1854; Pap. XI2A192; SKS 26, 366)

\subsection{Undated}

\begin{quote}
Det er Løverdag, og jeg kommer ikke til Dig. Da dog min Stemme, hvor høit jeg end talte her hjemme, vilde være for svag til at naae Dit Øre, hvis den ikke, som jeg haaber gjenlyder deri, (og skulde Du en enkelt Gang ikke fornemme den, da maa Du troe, det komer deraf, at den slumrer, men den vaagner snart igjen.), saa sender jeg Dig dette lille Billede, der hvor taust det end er, dog visselig ikke blot kan tale men ogsaa tiltale Dig. Det er en gammel Kone, der læser. Hun læser høit det er aabenbart, imidlertid ikke saa meget fordi der er Nogen, der hører derpaa, som fordi det er en Trang i Menneske, naar det skal \tlg{tilegne} sig det Bedste, da ikke blot at ville see det men ogsaa at ville høre det. (Du kan selv gjøre et Forsøg, jeg tvivler paa, at Du kunde være fornøiet med, at Dit Øie flygtig gleed hen over de Ord: Min Regine, eller udholdende stirrede derpaa, jeg i det mindste siger det høit atter og atter.). – For at jeg dog ogsaa kunde have Noget, har jeg beholdt dit Strikketøi. Det sidder jeg og seer paa. Jeg har skudt det lille Bord bort der staaer foran min Sopha, dit Strikketøi ligger ved min Side; jeg læser høit, der bliver Taushed, jeg bøier mig om for at udhvile mit trætte Hoved ved dit Bryst. – Det ringer – Jeg tager mine Briller paa for at lukke op. – Endnu engang vender jeg tilbage, for at underskrive dette Brev. Og hvad jeg modtog af min Fader, det overgiver jeg til Dig
\end{quote}
\dcite{B127:224  (SKS 28, 224)}
% B127:224  (SKS 28, 224)

\begin{quote}
Saa takker jeg Dem ogsaa for Selvprøvelsen, som De vel ikke har sendt mig, men det nytter Dem ikke; thi paa en vis Maade \tlg{tilegner} jeg mig alligevel Alt hvad De leverer. Om Samtiden vil modtage Anbefalingen – ja, jeg maa næsten tvivle endskjøndt jo nok andet Oplag er udkommet, men der kommer ialfald en Fremtid, der vil det. Samme Pastor Boisen, som ovenfor omtaltes, holdt for noget siden her i Byens Klub et Foredrag over det danske Ordsprog: »der slaaes mangt et Slag i Øresund, hvorefter Intet kjendes« og sagde blandt Andet, at der var to Slags Øresund, det store, ad hvilket man seilede ned til Kjøbenhavn, og det lille Øresund, ad hvilket man seilede ned til det menneskelige Hjertes Havn. Men det traf sig undertiden at man igjennem dette Sund ikke kunde naae Havn, og det kan da nu være af forskjellige Grunde, blandt andre ogsaa af den, at der var slet ingen Havn men kun et øde tomt Rum. Siden den Tid siger jeg ofte til mig selv: der slaaes mangt et Slag i Øresund, hvorefter Intet kjendes, men man kan alligevel ikke bare sig for at slaae i Haab om, at det dog maatte træffe et lille Kjende.
\end{quote}
\dcite{B43:98  (SKS 28, 98)}
% B43:98  (SKS 28, 98)

\part*{III.\quad Hans Brøchner}
\addcontentsline{toc}{part}{III.\quad Hans Brøchner}

\subsection{\textit{Om Søren Kierkegaards Virksomhed som religieus Forfatter} (1855)}

\begin{quote}
Det Christelige er altsaa Maalet, Opgaven, og er det fra først af.
»Forfatterskabet, totalt betragtet, er religieust fra Først til Sidst«. —
»Problemet, hele Forfatterskabets Problem er: det at blive Christen«. — »Til
dette Forfatterskab svarer som Ophav En, der qua Forfatter kun har villet
Eet«. — Nærmere er da dette, han har villet, bestemt saaledes: »Uden
Myndighed at gjøre opmærksom paa det Religieuse, det Christelige«. (Om min
Forfattervirksomhed, S. 7, 9, 14). Men de, der skulde gjøres opmærksomme, vare
de, der i »Christenhed« leve i »det Sandsebedrag: at kalde sig Christen, maaskee
indbilde sig at være det, uden at være det«. Opgaven var da ikke nærmest, at
give, men først ligesom at fratage, at skaffe det bort, der begunstigede
Sandsebedraget, og saaledes ved det Maieutiske, ved denne »Jordemoder-Konst«,
som Sokrates kaldte den, at forhjælpe den Enkelte til selv »at føde Sandheden«.
Men Maaden, paa hvilken der skulde virkes mod Sandsebedraget, var betinget ved
dettes Grund. Denne laa dels i Tidens hele Retning: at den, reflecterende og
lidenskabsløs, havde ladet Forstanden nivellere, og ladet »den Enkelte«
forsvinde i »Abstractionens Phantom, Publikum«, at den, »væsenlig forstandig,
maaskee i Gjennemsnit vidende som ingen Generation var det før den«, havde
»glemt at existere og hvad Inderlighed er«, at den, upersonligt og interesseløst
fortabende sig i det Objective, havde glemt \tlg{Tilegnelsen}, glemt Opgaven, at
existere i sin Viden; — dels laa den i den forvirrende Sammenblanding af de
forskjellige Livsanskuelser, der var begrundet i, at Verdsliggjørelsen af
Christendommen og Forglemmelsen af Personlighedens Lidenskab havde ført til
Forglemmelsen af Christendommens ideale Fordring. For at hæve Bedraget maatte da
først disse Livsanskuelser bringes frem i deres udprægede Eiendommelighed, og
bringes frem, ikke som en Gjenstand for Viden, thi saaledes vilde netop Tidens
Retning bort fra Existens og fra Existensens Afgjørelse begunstiges; men i en
Form, der hævder Existensens Betydning, der tvinger Læseren til ved sig selv
personligt at erhverve Forstaaelsen. Han »gjøres opmærksom«, det Øvrige overlades
til ham selv; den indirecte Meddelelse skjuler den Meddelende for ham, gjør denne
ligesom ukjendelig; han skal nu hjælpe sig selv i \tlg{Tilegnelsen}, og i \tlg{Tilegnelsen}
af det, der ved at fremstilles i Existens kun indirecte meddeles, skal han
personligt erfare, hvor den lavere Livsanskuelses Anstødssten ligger, hvor den
støder an i Lidenskab, hvor Existensen taber sit Grundlag, og hvor Fortvivlelsen
fører enten ud i Fortabelsen eller til det Høiere, til »at frelses uendeligt i
Religieusitetens Væsenlighed« (jfr. En literair Anmeldelse. 1846. S. 108 ff.).
Men idet de tidligere Livsstadier skulde fremstilles, ikke docerende, som
Gjenstand for Viden, men i existerende Individualiteter, maatte den første Række
af Skrifter blive pseudonym. Pseudonymiteten havde »sin væsenlige Grund i selve
Frembringelsen, der for Replikens, for den psychologisk varierede
Individualitets Forskjelligheds Skyld digterisk krævede den Hensynsløshed i Godt
og Ondt, i Sønderknuselse og Overgivenhed, i Fortvivlelse og Overmod, i Lidelse
og Jubel osv., der kun er ideelt begrændset af den psychologiske Conseqvens,
hvilken ingen factisk virkelig Person i Virkelighedens sædelige Begrændsning tør
tillade sig, eller kan ville tillade sig« (Anhang til Afsl. Efterskrift).
Pseudonymerne ere saaledes digtede Forfattere, ligesom dramatiske Personer, der
leve i deres Anskuelser, og udtale disse i hele deres ideale Conseqvens.
\end{quote}
\dcite{Brøchner, Om Søren Kierkegaards Virksomhed som religieus Forfatter, pp. 51–53  [brochner/kierkegaard-virksomhed:4]}
% brochner/kierkegaard-virksomhed:4  (Om Søren Kierkegaards Virksomhed som religieus Forfatter; pp. 51–53)

\begin{quote}
I det første af de pseudonyme Skrifter: Enten-Eller, et Livsfragment,
udgivet af Victor Eremita (1843) er nu de to Livsstadier, der ligge forud for
det Religieuse: det Æsthetiske og det Ethiske, fremstillet i existerende
Personligheder. »Det Æsthetiske er det i Mennesket, ved hvilket det umiddelbart
er det, det er; den, der lever i og ved og af og for dette, lever æsthetisk«;
hans Løsen er Nydelsen, hans Maal Fortabelsen. Første Del fremstiller nu en
saadan æsthetisk Individualitet, omtumlet i Muligheder, som hans Phantasi
fremmaner, hans Dialektik atter tilintetgjør, splittet ad i de isolerede
Øieblikkes Lidenskab, uden Existensens Sammenhæng, uden Virkelighed, fangen af
Fortvivlelsen. Den anden Del fremstiller en enkelt Individualitet, der i
Fortvivlelsens Opgiven af det Æsthetiskes Umiddelbarhed har gjort
Uendelighedens Bevægelse, hvorved det Ethiske først kommer tilsyne, der i
Fortvivlelsen har fundet sig selv igjen i sin evige Gyldighed, valgt sig selv,
som i dette Valg er bleven aabenbar i det Almene, og ved dette har faaet
Muligheden til at give sit Liv Continuitet. Værkets Slutning, en opbyggelig
Tale, fører hen til Grændsen af det Religieuse, og ved at bringe det Opbyggelige
ind, accentueres Subjectivitetens, \tlg{Tilegnelsens} Betydning i Forholdet til
Sandheden.
\end{quote}
\dcite{Brøchner, Om Søren Kierkegaards Virksomhed som religieus Forfatter, p. 53  [brochner/kierkegaard-virksomhed:5]}
% brochner/kierkegaard-virksomhed:5  (Om Søren Kierkegaards Virksomhed som religieus Forfatter; p. 53)

\subsection{\textit{Benedict Spinoza. En Monographie} (1857)}

\begin{quote}
Spinozas Liv og Skrifter. — En Tænkers Livshistorie er hans Tankes Historie, thi hans Livs Indhold er
hans Tanke. Hans ydre Liv faaer ikke umiddelbar Interesse som hos en
udadtil handlende Personlighed, hvis Liv afpræger sig i hans Samtid som
det selv bærer Afpræget af denne. Det bliver enten fuldkommen ligegyldigt
eller kun forsaavidt af Betydning, som dets Overensstemmelse med hans
Tænkning afgiver Borgen for dennes Magt til at gjennemtrænge et
Menneskeliv. Interessen bliver en afledet, væsentlig kan den ikke blive,
thi i Tankeindholdets \tlg{Tilegnelse} forsvinder det personlige Forhold til
Tænkeren, og maa forsvinde, hvis \tlg{Tilegnelsen} skal blive fuldstændig og
fri. Man kan derfor hos en Tænker, selv om man fordrer mere end hans
Tankes Genesis, nøies med et let Omrids af hans Liv, naar kun dette Omrids
er bestemt nok til, at hans Tankes Præg bliver synligt derigjennem; man
kan give Afkald paa alle de selv i deres tilsyneladende Tilfældighed
betydningsfulde Enkeltheder, der gjøre Omridset til et udført Billede.
Hvad Spinoza angaaer, saa er en saadan Nøisomhed nødvendig; han har kun
funden een noget udførligere Biograph, den lutherske Præst
Colerus [note: La vie de Spinoza, tirée des écrits de ce fameux
philosophe et du témoignage de plusieurs personnes dignes de foi, qui
l'ont connu particulièrement. Haag 1706. 8. En tidligere hollandsk Udgave
er udkommen i Utrecht.],

 og denne enfoldige, redelige, trods den theologiske Antipathie med den
uforstaaede Philosoph anerkjendende, Fortæller har i Spinozas simple,
udadtil ensformige Liv kun funden ringe Stof til at meddele, og hans
Meddelelser fuldstændiggjøres kun lidet ved de Enkeltheder, der findes i
Fortalen til Spinozas efterladte Skrifter, hos Bayle (i hans
Dictionnaire historique et critique), i »La vie de Spinoza par un de ses
disciples« (sandsynligviis Lægen Lucas), Amsterdam 1719, samt hos
Nicéron, Kortholt og Leibnitz.
\end{quote}
\dcite{Brøchner, Benedict Spinoza. En Monographie, pp. 27–28  [brochner/spinoza:88]}
% brochner/spinoza:88  (Benedict Spinoza. En Monographie; pp. 27–28)

\begin{quote}
Hvad der gjælder med Hensyn til Forseelse og Lydighed, gjælder ogsaa med Hensyn til Retfærdighed og Uretfærdighed. I Naturtilstanden gives der Intet, der kan siges med Rette at være den Enes og ikke den Andens; men Alt tilhører Alle forsaavidt de have Magt til at hævde det for sig. Men i Staten, hvor det ved den fælles Ret bestemmes, hvad der tilhører den Enkelte, der kaldes den retfærdig, der har en bestandig Villie til at give Hver Sit, uretfærdig den, der stræber at \tlg{tilegne} sig hvad der er en Andens.
\end{quote}
\dcite{Brøchner, Benedict Spinoza. En Monographie, p. 138+  [brochner/spinoza:268]}
% brochner/spinoza:268  (Benedict Spinoza. En Monographie; p. 138+)

\begin{quote}
Det er i Tractatus theologico-politicus at Spinoza udvikler sin Opfattelse af Religionens Væsen og dens Forhold til Staten og Philosophien. For ikke at misforstaae de Udviklinger, han giver, er det her, mere end noget andet Sted, nødvendigt at erindre den Betydning, han indrømmer Accommodationen [note: Tract. de intell. emend. pag. 361 anfører han blandt de foreløbige Leveregler: ad captum vulgi loqui et illa omnia operari, quæ nihil impedimenti adferunt, quo minus nostrum scopum (ɔ: at naae det høieste Gode) attingamus. Nam non parum emolumenti ab eo possumus acquirere, modo ipsius captui, quantum fieri potest, concedamus, adde, quod tali modo amicas præbebunt aures ad veritatem audiendam.]. Hensigten med Afhandlingen var, som det allerede blev angivet i Cap. 1, at skaffe Plads for den frie Philosopheren, og, forat opnaae dette, maatte deels Philosophien skilles fra Troen (Religionen), og begges forskjellige Charakteer paavises; deels maatte den religiøse Friheds Uadskillelighed fra Statens Frihed, og Statsmagtens Myndighed i Forhold til det Religiøse godtgjøres (jfr. tract. theol.-polit. præf. samt pag. 30, 160. 166). Men derved maatte alle de herskende Forestillinger om Religionens

 
Væsen og Skriftens Betydning paavirkes, mange Fordomme maatte fjernes, og efter disse Fordommes Natur og den Stands Charakteer [note: „Summum odium, quale theologicum esse solet.“] der forsvarede dem, og i dem sit eget Herredømme, maatte det skee med Varsomhed. Spinoza søger derfor ikke blot at benytte de Tilknytningspunkter, han kunde finde i sine Landsmænds republicanske Frihedsfølelse og i deres historiske Erfaringer, men i sine Udviklinger læmper han sig saameget som muligt efter den herskende Forestilling, bruger i sin Beviisførelse dens Former og dens Vaaben, \tlg{tilegner} sig dens Udtryk, og sørger kun for at lægge et andet Indhold ind i disse, saaledes at de, der vare fortrolige med hans Tanker, kunde gjenfinde dem i de fremmede Former, de, der ikke kjendte dem, forberedes til at opfatte dem. Erindrer man ikke bestandig dette, vil man ikke kunne undgaae at misforstaae Spinozas Skrift, saaledes som næsten alle hans Samtidige gjorde, og man vil da finde Modsigelser deri mod hans philosophiske Grundanskuelse, der ikke lade sig fjerne.
\end{quote}
\dcite{Brøchner, Benedict Spinoza. En Monographie, pp. 146–147  [brochner/spinoza:279]}
% brochner/spinoza:279  (Benedict Spinoza. En Monographie; pp. 146–147)

\begin{quote}
Men om end Spinoza kun fandt faa Tilhængere, saa
blev dog Spinozismen, som Navn for noget Uforstaaet og
Afskyet, snart et Kjætternavn, med hvilket man gjensidigen
stemplede hinanden, og meget Fritænkeri har vel ogsaa givet
sig et fornemmere Udseende ved at \tlg{tilegne} sig dette Navn.
Derfor kan en omtrent samtidig Forfatter: Roellius (de
religione naturali, § 151, pag. 166, citeret af Brucker)

 sige: Spinozam tota armenta in Belgio sequi ducem. Men
af flere af de endnu bevarede Stridsskrifter mod Spinoza,
der, skjøndt angribende ham, paadroge deres Forfattere Beskyldningen
for at være hans hemmelige Tilhængere, seer
man, hvor svævende og ubestemt man opfattede Spinozismens
Charakteer, ligesom ogsaa den hele polemiske Literatur
uden Undtagelse viser, at Spinoza bestandig misforstodes
netop i de væsentligste Punkter: i Opfattelsen af
Endelighedens Forhold til Gud og af det Guddommeliges
Værens Form. Fra Forestillingen om det Endelige som
reelt existerende, og om Gud som vilkaarligt virkende, —
Forestillinger, man ikke kunde opgive uden at opgive den
hele Grundvold, paa hvilken man stod, kæmpede man mod
Spinoza, og drog absurde Consequentser ud af hans Lære.
\end{quote}
\dcite{Brøchner, Benedict Spinoza. En Monographie, pp. 182–183  [brochner/spinoza:338]}
% brochner/spinoza:338  (Benedict Spinoza. En Monographie; pp. 182–183)

\begin{quote}
I Tidsrummet mellem Leibnitz og de sidste Decennier
af forrige Aarhundrede var Spinoza næsten glemt. Wolf
polemiserede imod ham (i sin theologia naturalis), men med
halve Begreber og uden at forstaae ham, og Tidsalderens
Retning var ugunstig for en dybere Tænkning. Mod Aarhundredets
Slutning paaberaabte den franske Materialisme
sig Spinoza som Autoritet. I hans Gudsbegreb saae
den kun Naturens Begreb, i hans Lære om Nødvendigheden
fandt den sin Determinisme, i hans Lære om Aandens Forhold
til Legemet sin Sensualisme. Det idealistiske Grundlag
hos ham oversaae den; den bragte ikke hine Bestemmelser
i Forbindelse med hans Anskuelsers Helhed, men \tlg{tilegnede}
sig desto ivrigere Enkelthederne af hans Polemik
mod de theologiske Forestillinger, og af hans Kritik.
\end{quote}
\dcite{Brøchner, Benedict Spinoza. En Monographie, p. 186  [brochner/spinoza:343]}
% brochner/spinoza:343  (Benedict Spinoza. En Monographie; p. 186)

\subsection{\textit{Problemet om Tro og Viden} (1868)}

\begin{quote}
Problemets nærværende Stilling er for vor Literaturs Vedkommende
betegnet ved Kampen imellem en Form af Theologien, der væsentligen
laaner de Vaaben, med hvilke den vil hævde Foreningen af Tro og Viden,
paa eklektisk Maade fra nyere philosophiske Blandingsformer, og det
antitheologisk-religiøse Standpunkt, der ved den principielle
Modsætning mellem Tro og Viden vil sikre Religionens og Videnskabens
Gyldighed i deres gjensidige Uafhængighed og Muligheden til deres
fredelige Samværen i Bevidstheden, medens den i Kraft af samme
principielle Modsætning frafjender Theologien al Ret til at bære
Videnskabens Navn, hvilket den \tlg{tilegner} sig idet den gjør Fordring paa
at være en „Troesvidenskab“.
\end{quote}
\dcite{Brøchner, Problemet om Tro og Viden, pp. 1–2  [brochner/problemet-tro-viden:3]}
% brochner/problemet-tro-viden:3  (Problemet om Tro og Viden; pp. 1–2)

\begin{quote}
Det Nye Testamentes Udsagn om Forholdet — Hvad der, idet den nye Lære forkyndes, strax maae træde frem, er
Modsætningen mellem denne og det, som den skal afløse. Denne Modsætning er en
Modsætning mellem Indholdet af det Evangelium om Gudsrigets Komme, der
forkyndes og troende \tlg{tilegnes} af Christne, og af den ikke-christne Verdens
Overbeviisning, der har sit Fællesgebet i en Tænkning, i hvilken den græske og
den østerlandske Aand mødtes. Idet hiin Overbeviisning i dets Modsætning til
Evangeliet nærmest repræsenteres af de Dannede, de Vidende, bliver
Modsætningen en Modsætning mellem Verdens Viisdom og Christi for de Enfoldige
aabenbarede Viisdom (Guds Viisdom); men ogsaa denne Modsætningens Form viser
væsentligen, idet »Viisdom« tages i objectiv Betydning, tilbage til
Modsætningen i Indhold. Et nyt Indhold, fordum skjult for Menneskeheden, er
aabenbaret, en Guds Raadslutning, for hvilken Tidens Fylde er kommen: et
μυστήριον. Mysteriets Begreb viser hen til det nye Indhold, der omfatter ogsaa
det i Fremtiden Forborgne (saaledes Jødefolkets Fremtidsskjæbne, Opstandelsen
og Forvandlingen); men det viser tillige hen til det, der kun bliver aabenbart
for dem, i hvem Aanden virker. Idet Begrebet omfatter denne sidste Betydning,
antydes derved Modsætningens Grund, dens Princip, men idet der ikke kan være
Tale om, at den troende, for Viden fremmede, Bevidsthed kan søge Grunden i et
Videns Princip, saa søges Grunden til Modsætningen mellem den ikke-christelige
Bevidsthed og den christelige i en Sindelagets, en Villiens Mangel, der
forsaavidt hiin Bevidsthed har en Kjendskab til dennes Troesindhold, bliver en
Mangel paa Villie til at antage Aabenbaringens Indhold, og denne Mangel paa
Villie sees atter som begrundet i det »Kjødelige« og det »Psychiske«, der
danner Modsætningen til det »Pneumatiske«. Bagved denne Bestemmelse af
Modsætningen ligger den principielle Modsætning, der i sin Sammenhæng med
Problemet først kommer klart og skarpt frem i Scholasticismen: Modsætningen
mellem det Naturlige og det spiritualistisk Overnaturlige. Det Naturliges
Begreb omfatter den af Synden formørkede og fordunklede Tanke, og for denne er
derfor den høiere overnaturlige Sandhed utilgængelig, der kun gjennem en
guddommelig Aabenbaring bliver tilgængelig for Troen.
\end{quote}
\dcite{Brøchner, Problemet om Tro og Viden, pp. 15–16  [brochner/problemet-tro-viden:29]}
% brochner/problemet-tro-viden:29  (Problemet om Tro og Viden; pp. 15–16)

\begin{quote}
Problemets Udvikling i Scholasticismen — Scholasticismen har Dogmets Udfoldelse til en væsentlig afsluttet
Heelhed til Forudsætning; dens Opgave er, ved en reflecterende Tænknings Hjælp at
\tlg{tilegne} Bevidstheden Dogmet ved at finde, om end ei Dogmets ratio, saa dog
Dogmernes rationes. Naar vi derfor skulle betegne dens Tid efter dens Stræben,
betegne vi den som \tlg{Dogmetilegnelsens} Tid. Ligesom den ved denne Bestemmelse
adskiller sig fra Patristikens Tid, saaledes adskiller den sig fra denne ved de
forandrede Tidsbetingelser. Den Afhængighed af Oldtiden, der var begrundet i den
umiddelbare Nationalitetssammenhæng, var ophørt, og Dogmefrembringelsen var
saaledes afsluttet, at et System, omfattende Troens og Videns Totalitet, nu kunde
forsøges. I Kirkens Tjeneste og bøiende sig under dens Autoritet stræber Tanken at
opføre en saadan christelig Lærebygning, i hvilken Troesindholdet og
Tankeindholdet skal sammensmeltes til en Heelhed. Sammenføiningens Art bliver
tydelig naar vi betragte de store Systemer, der frembragtes i Scholasticismens
Blomstringsperiode. I disse see vi den antike Verdensanskuelse, der samlede sig i
Forestillingen om den harmonisk ordnede guddommelige Κόσμος, ligesom omfattet af
det Christeliges Ramme. Anskuelsen af Verdensvirkeligheden, der maatte tages fra
den som historisk Forudsætning givne Philosophie, for hvilken Verden var en
Virkelighed, og som for Oldtiden var det Hele, blev nu, forbunden med det
christelige Troesindhold, til en Deel, indføiet i en større Heelhed. Det fra
Oldtiden optagne Tankeindhold udfyldte kun den midterste Deel af det
Verdensdrama, af hvilket første og sidste Act tilhørte Christendommen. Forud for
det naturlige Universum sattes den naturfrie, almægtige, skabende Gud, efter det
den christelige Himmel, Naadens overnaturlige Rige. Men i denne Sammensætning
blev, efter Scholasticismens egen Indrømmelse, Begyndelsen og Slutningen et for
Tanken Utilgængeligt, altsaa i Væsen uensartet med det, der udfyldte
Mellemrummet; det specifisk christelige Indhold stiller sig altsaa, trods den
forsøgte Sammenføining, som et Ubegribeligt fra Tankeindholdet. Og idet man søger
Grunden til Ubegribeligheden, saa finder man den i Incongruentsen imellem
naturlig Tænkning og overnaturlig Aabenbaring. Troens »høiere Erkjendelse«, der
skylder en ved en guddommelig Undervirkning tilveiebragt Potentsation af de
naturlige Evner sin Tilværelse, har kun Navnet tilfælleds med Videns Erkjendelse,
der hviler paa disse naturlige Evner, som ikke staae i et continuerligt
Udviklingsforhold til det ved den miraculøse Potentsation Satte. Den naturlige
Erkjendelses Betydning for Troesindholdet bliver kun den: i negativ Form at hævde
det, som det Ikke-Umulige, ikke den: positivt at begrunde det. Men trods den
Væsensdifferents, der kun lader Navnets Lighed tilbage, bestemmes dog Forholdet
mellem Tro og Viden som et Forhold mellem to Erkjendelsesformer. Alligevel ligger
der i Modsætningernes Uensartethed en Tilskyndelse til at føre dem tilbage til
forskjellige Aandsfunctioner, og en saadan Tilbageføren finde vi f. Ex. hos
Duns Scotus, naar han sætter Troen som kun blivende til ved en Villiesact og
derfor lader den paa Troen hvilende Theologie, hvis Maal ikke nærmest er
Erkjendelsen, men Saligheden, Nydelsen af Gud, staae i nærmere Forhold til Villien
end til Forstanden, idet den ikke bliver en speculativ, men en praktisk Lære, der
skal gribe vor Villie.
\end{quote}
\dcite{Brøchner, Problemet om Tro og Viden, pp. 18–19  [brochner/problemet-tro-viden:34]}
% brochner/problemet-tro-viden:34  (Problemet om Tro og Viden; pp. 18–19)

\begin{quote}
Hos Kierkegaard er det Religiøse den oprindelige og bevægende personlige
Interesse og det Maal, til hvilket der fra Begyndelsen af sigtes hen; men i
Udviklingen af Standpunktet er Viden Udgangspunktet, og det er gjennem Videns
Selvbegrændsning, gjennem Erkjendelsen af den Skranke, Existentsforholdet sætter
for den existerende Vidende, at Troens og Troesforholdets Nødvendighed
anerkjendes. Videns Gyldighed bliver da ikke tilintetgjort, men anerkjendes som
relativ, indenfor sin bestemte Grændse; Troen bliver, som den for den
Existerende høieste Vishedsform, det absolut Gyldige, det, der ene giver den
Existerende den, ligesom af Uvisheden udrevne, forsonende Vished om det, der
tilfredsstiller Personlighedens dybeste Trang. Viden kan, netop paa Grund af sin
Begrændsning, ikke forstyrre Troens Indhold, ikke heller begribende \tlg{tilegne} sig
det; det Eneste, den i Forhold til Troen kan gjøre, er at blive tjenende, idet
den værner om Paradoxet, bevarer det som saadant i dets Skarphed og
Utilgængelighed. Omvendt lader Troen Viden være uforstyrret paa dens begrændsede
Omraade, og ifølge dette Forhold kunne de to Principer med deres Sphærer faae
deres Plads i samme Bevidsthed uden forstyrrende Sammenstød.
\end{quote}
\dcite{Brøchner, Problemet om Tro og Viden, pp. 26–27  [brochner/problemet-tro-viden:48]}
% brochner/problemet-tro-viden:48  (Problemet om Tro og Viden; pp. 26–27)

\begin{quote}
Hvad først Sphærernes Uendelighed angaaer, der umiddelbart maatte følge som
Consequents af Principernes absolute Gyldighed, og som maatte indeslutte, at al
Virkelighed, Virkeligheden i dens Fuldstændighed, fik sin Plads indenfor hver af
de to Sphærer, saa lader den sig ikke fastholde. Med Hensyn til Videnssphæren,
der synes at faae en saadan Uendelighed derved, at Viden ogsaa har det til
Gjenstand, som den blot afspeiler, vilde Uendeligheden ialfald blive en
Uendelighed med Hensyn til Omfang, ikke til Indhold. Ved selve Adskillelsen
mellem det, der assimileres — det ved Tanken Begrundede og Begrebne
assimileres — og det, der kun afspeiles, bliver der i Vidensindholdet sat en
Grændse, der viser hen til en Væsensbegrændsning i Viden som menneskelig Viden,
selv om denne Begrændsning postuleres hævet i Videns Idee. Men
Væsensbegrændsningen, der umiddelbart træder frem gjennem Vidensindholdets
Begrændsning, maa uundgaaelig atter føre tilbage til en Begrændsning med Hensyn
til Omfang, paa Grund af det dialektiske Forhold mellem Omfang og Indhold og
deres nødvendige gjensidige Bestemmen. Det er umiddelbart indlysende, at det,
der kun afspeiles i Viden, ikke blot naar Viden tages i stræng Forstand som
\tlg{tilegnende} (assimilerende) Viden, og kun i denne Forstand er den egentlig Viden,
kun vides efter sin abstracte Væren; men at det, selv naar Viden tages i udstrakt
Betydning, ogsaa kun vides abstract, fordi den concrete Viden kun er ved
Erkjendelsen af Sammenhængen, men for den afspeilende Viden maa det Afspeiledes
Sammenhæng unddrage sig idet dens Form — paa Grund af den absolute
Uensartethed i Principerne — bliver absolut incommensurabel med Videns. Og
hvad der gjælder med Hensyn til Videnssphærens Uendelighed har sit Tilsvarende
med Hensyn til Troessphærens. Ikke heller denne kan fastholdes i den angivne
Betydning. Indenfor Troen findes der intet Forhold, der svarer til
Afspeilingsforholdet i Viden og ifølge hvilket Videns Indhold kunde optages i
den. Og selv om et saadant Forhold fandt Sted, vilde Troen ikke faae al
Virkelighed til Gjenstand, af den tilsvarende Grund til den, af hvilken Viden
ikke havde den. Troessphæren kan saaledes ikke hævde sin Uendelighed. Kun dette
Resultat søges her; i det Efterfølgende vil det paa sit Sted blive viist, at
Sphærernes factiske Begrændsethed, hvormeget end Selvbegrændsningen accentueres,
dog vil føre til et Sammenstød ved Grændsen med den Sphære, der i denne berøres.
\end{quote}
\dcite{Brøchner, Problemet om Tro og Viden, pp. 27–29  [brochner/problemet-tro-viden:51]}
% brochner/problemet-tro-viden:51  (Problemet om Tro og Viden; pp. 27–29)

\begin{quote}
[label= )]
 Det religiøse Standpunkt, der sætter Troen, som det ved Guds overnaturlige
 Virken betingede \tlg{Tilegnelsesforhold} til den guddommelige Aabenbaring, som
 absolut hævet over den naturlige Viden, der opfattes som den, der ikke blot
 ikke kan erkjende Aabenbaringens Indhold, men maa fornægte det.
 Theologiens og en enkelt philosophisk Forms Standpunkt, der statuerer en
 quantitativ Differents til Fordeel for Troen [note: Til dette Standpunkt
 gaaer i Realiteten den scholastiske Sondring af duæ veritates
 tilbage.]. 
 Den rene Rationalismes Standpunkt, paa hvilket Forskjellen udvises ved
 Forestillingen om en Fornuftreligion, en naturlig Religion, med hvilken
 Christendommen identificeres.
 Mæglingens Standpunkt indenfor den idealistiske Philosophie, paa hvilket
 der statueres en Formdifferents til Fordeel for Philosophien.
 Den realistiske og den consequente idealistiske Philosophies Standpunkt,
 paa hvilket, idet Forholdet bestemmes fra Videns Synspunkt, Viden faaer en
 absolut Overvægt over Troen, og bliver det eneste Sandhedsorgan.
 Sondringen af Sphærerne for at udstille det Religiøse.
 Sondringen af Sphærerne for at tilveiebringe begges Forening i samme
 Bevidsthed.
\end{quote}
\dcite{Brøchner, Problemet om Tro og Viden, pp. 30–31  [brochner/problemet-tro-viden:56]}
% brochner/problemet-tro-viden:56  (Problemet om Tro og Viden; pp. 30–31)

\begin{quote}
Hvad det første Punkt angaaer, saa bestemmes Troen som en Villiens \tlg{Tilegnelse} i
Kjærlighed af den guddommelige Aabenbaring. Villien er altsaa den Menneskeaandens
Function, til hvilken den har sit væsentlige Forhold, men Villien er, ligesom den
i Universet er det Første, det egentlig Constitutive i den menneskelige
Personlighed: Menneskets Grundværen er Villie. Men saaledes bliver Villien det
Høieste, det, der har Principatet i Forhold til de andre Aandsfunctioner, der ere
afhængige af den. At det, der i Troen gribes gjennem Villien, ikke er
commensurabelt med Tanken, med Begrebet, viser derfor kun, at hint er det Høiere
og Rigere.
\end{quote}
\dcite{Brøchner, Problemet om Tro og Viden, p. 46+  [brochner/problemet-tro-viden:84]}
% brochner/problemet-tro-viden:84  (Problemet om Tro og Viden; p. 46+)

\begin{quote}
Den nye Erkjendelsens Form, der sættes som en høiere end den blot logiske Tænkning, skal
være den høiere og gjøre den derpaa byggede positive Videnskab til en høiere, fordi den er
en Erkjendelse af den frie personlige Guds Aabenbaringer. Men det Begreb om en
guddommelig Personlighed, til hvilket der vises hen, er allerede belyst, og det er viist,
at det kun har den endelige Enkeltpersonligheds Indhold, hvorved der bringes en
Endelighedsbestemthed ind, der gjør det af denne Kilde Fremgaaende til et Lavere end det
af Ideen Udledede. Man bliver stillet ligeoverfor det Dilemma: naar Naturen og
Historien skulle erkjendes som Guds Aabenbaringer, saa maae de sees som begrundede i Guds
Væsen, og dette kan da, som al Tilværelses absolute Princip, ikke bestemmes ved
Personlighed; — fastholdes derimod Personlighedens Bestemmelse, saa bliver Gud
det absolut Ubegribelige, og Naturen og Historien kunne ikke af Tanken opfattes som Guds
Aabenbaringer. Personlighed skal netop være det, der ikke gaaer op i Begrebet, det,
der kun bliver Gjenstand for Anskuelse, Følelse, Tilbedelse. Den absolute
Personlighed vilde være den allerpersonligste Personlighed, det absolut singulære Væsen,
der aldeles ikke havde almindelige Prædicater, og derfor ingen Tilknytningspunkter for
Tænkningen. „Vel faaer Gud, naar han bestemmes som personlig, en Bestemmelse, der ogsaa
tilkommer mig, men dens Absoluthed er dens Singularitet, dens Distinction fra mig.
Angaaende en anden endelig Personlighed kan jeg, paa Grund af dens Beslægtethed med mig,
vel gjætte og formode Noget; men om den absolute Person kan jeg, da her endog Basis for
Forudsætningen om en mulig Lighed med mig i Tænken og Handlen falder bort, kun drømme og
phantasere.“ (Feuerbach.) — Hvad der aabenbares af den frie, personlige Gud, skal kun
kunne erkjendes af Mennesket ved Frihed, og dette Friheds-Forhold, der
modsættes Tankenødvendighedsforholdet, skal saa betegne den nye Erkjendelses og den
nye Videnskabs høiere Værd. Men denne frie Erkjendelse, der skal være en høiere
Erkjendelse, maa tænkes at blive til ved et Brud, sat ved Villien, — en virkelig
Begrebsbestemmelse af Forholdet giver Martensen ikke —; saaledes at Tankenødvendigheden
afbrydes ved Villiesacten, der sætter en ny Begyndelse. Villien, der sætter denne,
bliver da en fra Erkjendelsen forskjellig Villie, men for en saadan Villie have vi ingen
anden Betegnelse end Vilkaarlighed, og den Erkjendelse, der skal hvile paa denne
Villiesact, bliver en saadan, der bæres af „Villiens Fjederham“ Phantasien, og savner
ethvert Correctiv, enhver Regulator. Dens „Frihed“ er et Ringere end Tankens Nødvendighed;
thi den Tankenødvendighed, der ved den rationelle Videnskab skal være Tegnet paa dennes
lavere Værd, er netop Udtrykket for Tankens Bestemthed ved sit eget Væsen, sin egen Lov;
den er derfor Særkjendet for al egentlig Videnskab, og det er en
Misforstaaelse, naar den indskrænkes til Videnskabens abstracteste og
indholdsløseste Former. Jo concretere og rigere Videnskabens Indhold er, desto
vanskeligere bliver den Assimilation, der giver Gjenstanden Tankenødvendighedens Form, men
derfor bliver denne Assimilation ligefuldt Opgaven, naar ikke det Concrete, kun successivt
\tlg{Tilegnelige}, forvexles med det ifølge sin irrationelle Natur med Tanken Incommensurable og
for den absolut \tlg{Utilegnelige}. Og paa disse Gebeter faae vi ikke en abstract
Tankenødvendighed, men den concrete Tankes Nødvendighed, der er tvingende for Tanken
fordi den er Tankens eget Væsen.
\end{quote}
\dcite{Brøchner, Problemet om Tro og Viden, pp. 106–108  [brochner/problemet-tro-viden:192]}
% brochner/problemet-tro-viden:192  (Problemet om Tro og Viden; pp. 106–108)

\begin{quote}
Kierkegaards Opfattelse udvikles gjennem en Polemik mod Speculationen,
navnlig saaledes som den fremtræder i den Hegelske Philosophie; men
Polemiken rettes mod Speculationen fordi Kierkegaard i den seer det
stærkeste Udtryk for hele Tidens Misvisning: at have glemt at existere,
glemt hvad Inderlighed er. For den speculative Philosophie var det
Menneskeaandens høieste Opgave at hensætte sig i et Evighedens Medium,
leve et rent Evighedens Liv, og denne Opgave skulde realiseres ved den
rene Tanke, idet Individets Tænken af det Absolute sattes som Eet med
det Absolutes Tænken af sig selv, med den absolute, i Evighedens
Element værende, Tanke. Individets Liv, der fandt sit udtømmende Udtryk
i den logiske Tanke, gjennem dets Tænken af det Absolutes, al Sandhed
og Virkelighed omfattende, evige Livsproces, blev et Moment i denne
evige Proces, og frigjordes fra Endelighedsbestemtheden. Mod denne 
Opfattelse af Menneskelivet og dets Opgave protesterer Kierkegaard.
Mod Speculationens Forflygtigelse af Existentsen, mod dets Glemmen af
Existentsens Krav til det det Ethiskes Fordringer underkastede Individ,
mod dets Overseen af Existentsens Betingelser, mod den illusoriske
Bevægelse i en abstract Evigheds Sphære accentuerer Kierkegaard
Existentsen, samler Alt om det Spørgsmaal: hvorledes udtrykker jeg
existerende Sandheden? Som existerende er Mennesket ikke blot evigt,
hvad Speculationen vilde, men er en Synthese af det Timelige og Evige,
af Endeligt og Uendeligt; som existerende er den efter sit Væsen evige
Aand underkastet Endelighedsbetingelser, er i Vorden; som existerende
skal Tænkeren i al sin Tanke tænke det med, at han er existerende; han
skal sætte al sin Tanke i Vorden. Men Vorden er Alternationen mellem
Væren og Ikke-Væren, ved Vorden bringes i hvert Moment Negationens
Bestemmelse ind i Aandens Væren og Tænken. Heraf følger saa med
Nødvendighed, at Mennesket ved Existentsen er forhindret fra at naae
den positive, betryggede Hvile i Evigheden, thi idet det, som
existerende, ikke kan fastholde det Uendelige og Evige, der er det
eneste Visse, i Uendelighedens Element, men skal fastholde det i
Endelighedens Element, i Existentsen, saa bliver Forholdet ikke den
betryggede Visheds, men det maa udtrykke den Negationens Bestemmelse,
der ligger i Vorden, og det Evige kan derfor af den Existerende, for
hvem det omsætter sig i et Tilkommende, kun haves for Erkjendelsen i
Uvishedens Form. Heri ligger saa, at den objective Viden i ingen af
dens Former kan faae en saadan Sikkerhed, at den kan blive Grundlaget
for den uendelige Afgjørelse for Subjectet, for en Afgjørelse med
Hensyn til det, der angaaer Subjectets dybeste Trang. Den sandselige
Vished er et Bedrag, den historiske Viden er en Approximation, den
mathematiske Erkjendelse er en Abstraction, den speculative Tænkning
ligeledes en Abstraction og i sidste Instants en Hypothese i Forhold
til Tilværelsen; paa ingen af disse Videns Former kan derfor den
Individet uendeligt interesserende Afgjørelse bygges. Men kan Maalet
ikke naaes ad Objectivitetens Vei, saa maa det naaes ad
Subjectivitetens: i et Subjectivitetsforhold maa det Evige gribes, og
dette Subjectivitetsforhold er det, i hvilket Visheden om det Evige med
Troens Lidenskab gribes ud af den objective Uvished. Troen er „den
objective Uvished, fastholdt i den meest lidenskabelige Inderligheds
\tlg{Tilegnelse}“.
\end{quote}
\dcite{Brøchner, Problemet om Tro og Viden, pp. 118–120  [brochner/problemet-tro-viden:212]}
% brochner/problemet-tro-viden:212  (Problemet om Tro og Viden; pp. 118–120)

\begin{quote}
Den speculative Religionsphilosophie havde ladet Religionen være et
Trin i Aandens Erkjendelsesudvikling, der fandt sin Fuldkommelse i
Philosophien, i den speculative Tanke, der med logisk Nødvendighed
udfoldede sig af den lavere Form. Først i Begrebet lod den Troens
Indhold finde dets sande Form; ved at omsættes i den philosophiske
Erkjendelses Form fik Troesindholdet en høiere Sandhed; thi Tankens,
Begrebets Sandhed, er den høieste Sandhed. Hos Kierkegaard sees
Sphærernes Forhold fra et andet Synspunkt. Videnskaben bliver Et,
Livet et Andet; Udfoldelsen af Tanken i Objectivitetens Form, i
Mulighedens Medium, efter Nødvendighedens Lov, skilles fra Sandhedens 
\tlg{Tilegnelse} i Individet i Subjectivitetens Form, i Virkelighedens
Medium, efter Frihedens Lov. Sandhedens personlige Existents bliver
accentueret som det Høieste, Subjectivitetsforholdet bliver denne
Existentses Form, og indenfor de personlige Existentsformer bliver den
religiøse den høieste, og indenfor den atter Troesforholdet til det
for Tanken absolut Utilgængelige: det Paradoxe, Christendommens
specifiske Bestemmelse. Tankens Opgave bliver her ikke at assimilere
Troesgjenstanden, at opløse den i Begrebet, men at opdage Tankens
Modsætning: det absolut Utænkelige, Paradoxet, og at fastholde det
som Paradox forat der saa kan troes i Kraft af det Absurde.
\end{quote}
\dcite{Brøchner, Problemet om Tro og Viden, pp. 123–124  [brochner/problemet-tro-viden:219]}
% brochner/problemet-tro-viden:219  (Problemet om Tro og Viden; pp. 123–124)

\begin{quote}
Den almindelige Theologie fordyber sig objectivt i lærde
Undersøgelser; den vil gjennem disse bringe frem hvad Christendommen
er som objectiv Lære; ad Historiens og Kritikens Vei vil den sikkre
sig Visheden herom, og under denne Stræben, der, efter Sagens Natur,
taber sig i en uendelig Approximation, lader den det være
Forudsætningen, at den Undersøgende er Christen, og at det, det
kommer an paa, er det videnskabelige Resultat, hvis \tlg{Tilegnelse} skal
følge af sig selv. For Kierkegaard er Existentsens uendelige
Afgjørelse Hovedsagen: Approximationen, ved hvilken netop Afgjørelsen
forhindres, er Bedraget; \tlg{Tilegnelsen} er Opgaven, og det Væsentlige
ikke et Hvad, men et Hvorledes. Der spørges ikke om, hvad nu
Christendommen historisk er; men der spørges: hvorledes bliver jeg
Christen? Sandheden er \tlg{Tilegnelsens} Inderlighed, Sandheden er
Subjectiviteten.
\end{quote}
\dcite{Brøchner, Problemet om Tro og Viden, p. 124  [brochner/problemet-tro-viden:220]}
% brochner/problemet-tro-viden:220  (Problemet om Tro og Viden; p. 124)

\begin{quote}
En lignende Uklarhed og Vilkaarlighed træder frem hvor Modsætningen bestemmes ved det
Subjective og det Objective, af hvilke hint forholder sig til Livet, denne til
Videnskaben. Hvor Modsætningen bestemmes ved disse Former (f. Ex. i phil. Propædeutik
§ 15) tages „Subjectiviteten“ og „det Subjective“ i Dobbeltbetydning, saaledes at
det snart kommer til at udtrykke det tilfældigt Individuelle og det Egoistiske, der falder
udenfor det objectivt Gyldige, snart Subjectivitetens almene Form. Ogsaa her er Forholdet
til Kierkegaard bestemmende. For Kierkegaard betegnede den subjective Sandhed den
\tlg{tilegnede} Sandhed. Denne Sandhed er paa det Ethiskes Gebet det Almene, der realiseres i
Existents, og som maa tænkes forat der kan existeres deri. Paa det paradox Religiøses
Gebet har den Sandhed, der skal \tlg{tilegnes}, Paradoxets Bestemmelse. I alle Sphærer sættes
Sandhedens \tlg{Tilegnelse}, fordi Existentsen forhindrer den objective Vished, som betinget ved
en Act, der overalt kan betegnes som en Troens Act. Hos Nielsen, hvor det ikke er
Existentsens Usikkerhed, men den abstracte Modsætning mellem det Almene og det Enkelte, der
betinger Disjunctionen mellem Viden og Existents, bliver det Subjective, i Forhold til
Videns Objectivitet, sat under en ensidig Modsætnings Bestemthed, som Væren i
Individualitetens Modsætning til det Almene, og derfra kommer Dobbeltbetydningen, idet
deels den Kierkegaardske Bestemmelse eftervirker, deels den nye Betydning udelukkende
træder frem. Men idet Subjectivitetens Begreb opfattes i denne Betydning gjennem en ensidig
og abstract Modsætning, oversees det, at Subjectivitetens Form ogsaa rummer det
objectivt Gyldige, at Subjectet optager det som Norm i sig uden at Indholdet forandres, at
den subjective Handling kan have objective Principer, at der er et objectivt Sandt og
Gyldigt i den fornuftbestemte Subjectivitet. Subjectivitetens, Livets, Standpunkt
kommer i hiin Opfattelse til at falde sammen med Egoismens Standpunkt.
\end{quote}
\dcite{Brøchner, Problemet om Tro og Viden, pp. 131–132  [brochner/problemet-tro-viden:233]}
% brochner/problemet-tro-viden:233  (Problemet om Tro og Viden; pp. 131–132)

\begin{quote}
Paa samme Maade som Subjectivitetens Begreb, faaer ogsaa Objectivitetens Begreb en
Dobbeltbetydning. Hvor Subjectivitetens Hovedstadier udvikles [note: Nielsen:
Philosophisk Propædeutik, p. anf. St.] betyder det Objective snart det objective Ydre
eller det blot abstracte Almene, snart det det Subjectives Totalitet omfattende Almene og
Almeengyldige. Og idet den abstracte Modsætning definitivt fastholdes, bliver det
vilkaarligt ignoreret, at paa de forskjellige Stadier Subjectets Bestemthed — hvad der
især fremtræder tydeligt paa det Stadium, der betegnes som den uendelige Subjectivitets,
— netop er i Kraft af hvad der er bestemt som Tilværelsens objective Princip og at den
udtrykker dette. Det bliver overseet, at den subjective Sandhed netop er den \tlg{tilegnede}
og i \tlg{Tilegnelsen} med sig væsentlig lige objective Sandhed.
\end{quote}
\dcite{Brøchner, Problemet om Tro og Viden, pp. 132–133  [brochner/problemet-tro-viden:234]}
% brochner/problemet-tro-viden:234  (Problemet om Tro og Viden; pp. 132–133)

\begin{quote}
Naar vi gaae ud fra den Indrømmelse, at Erkjendelse og Villie maae
sættes som eiendommelige, selvgyldige Yttringsformer af Væsenet, der
ikke kunne reduceres til Modificationer af hinanden, saa bliver det
strax Opgaven at bestemme, Hvormeget der ligger i denne Indrømmelse og
om den maa føre til de postulerede Consequentser:
Subordinationsforholdet og den absolute Uensartethed. Nielsen bestemmer
Villiens Forhold til Selvet i den Sætning: „det i Subjectiviteten
Allersubjectiveste er Selvet; men det Selviske i Selvet er
Selvbestemmen, er Villie“. Ved at bruge Begrebet Selvbestemmen
saaledes, at det paa een Gang bliver det væsentlige Udtryk for
Selvhedsenergien og identificeres med Villien, falder Eftertrykket paa
denne sidste; den \tlg{tilegner} sig — om end Tænkningens, ligesom ogsaa
den kunstneriske Phantasies, oprindelige Eiendommelighed udtales — i
Virkeligheden Principatet, og den nærmeste og umiddelbare Consequents
bliver, at den Sidestillen af Villie og Viden, til hvilken den
principielle Modsætning fører, bliver til Vidensprincipets
Subordination under Villiesprincipet: en Subordination, der under Theoriens
Udvikling træder frem i mange Former. Til denne Opfattelse bidrager ogsaa den ovenfor
paaviste ensidige
Opfattelse af Begrebet Existents (Liv), der identificerer det exclusivt
med det Praktiske, med Villien („Livet er i Villien“). Anderledes
stiller Forholdet sig naar hiin Vilkaarlighed undgaaes, naar vi, i
Overensstemmelse med hvad der paa et tidligere Punkt blev antydet, ved
Opfattelsen af det Menneskelige gaae ud fra den Bestemmelse, der
kategorisk sondrer Menneskelivet fra Livets lavere Stadier: fra
Væsensuniversalitetens Bestemmelse, der udtrykker sig i den
Forsigværen i den frigjorte Almindeligheds Form, der er
Selvbevidsthedens, og naar vi opfatte Væsenets Grundbestemmelse som
Energie, nærmere som Selvhedsenergie. Begge Former: Erkjendelse og
Villie, blive da primitive Udtryk for det i Selvbevidstheden som Enhed
aabenbarede Væsens Energie, for den Selvets Selvbestemmen, gjennem
hvilken det, bestemmende sig, realiserer sig: i begge Functioner hævder
Selvheden og Energien sig under eiendommelige Grundformer. Opgaven
bliver da, nærmere at bestemme Yttringsformernes Eiendommelighed, og
Bestemmelsen af den vil afgjøre Spørgsmaalet om Modsætningsforholdet
og nærmere belyse det postulerede Subordinationsforhold.
Eiendommeligheden bliver at bestemme saaledes, at i Erkjendelsen
Selvhedsenergien nærmest har Almindelighedens og Idealitetens, i
Villien Enkelthedens og Realitetens Form. Men disse som Udgangspunkter
givne Bestemmelser vise hen paa hinanden og bestemme sig henimod
hinanden. Realiteten, der constitueres ved Villien, er en Realisation
af det ideelle Væsens Bestemmelser i Selvet eller udenfor Selvet.
Idealiteten, der virkeliggjøres gjennem Erkjendelsen, er en ved
Selvets Realitetsside betinget Idealiseren af et Realt og er en saadan
selv i det ideelle Selvs Selverkjendelse. Og paa samme Maade er
Almindeligheden, der charakteriserer Erkjendelsen, en Almindelighed, der
stræber henimod Concretionen, Punktualiteten, og hvis sidste Maal er
den concrete Erkjendelse af Tilværelsen, hvilken Erkjendelsen af Selvet
efter dets individualiserede Væsensbestemmelser og væsentlige Forhold
er indbefattet; med andre Ord: det individuelt-almene Væsens
Gjennemsigtighed for sig selv som Led i en Tilværelsens Totalitet.
Omvendt er Enkelthedsbestemmelsen i Villien en saadan, der med
Nødvendighed udfolder det Almene i sig, fordi Villien kun er Villie som
det selvbevidste, ɔ: efter Begrebet almene og universelle, Væsens
Villie. Villien maa udfolde sig fra den ved det Enkelte bestemte, det
Enkelte sættende, Villie til den ved det, Væsenet udtrykkende,
Almeenfornuftige bestemte, det Almene realiserende Villie.
Villiesyttringens Enkelthedsform er her ikke ophævet, men
Enkelthedsbestemtheden er her bleven gjennemsigtig i det Almene. Paa
lignende Maade som med Almindelighedens og Enkelthedens Bestemmelser
forholder det sig med de dermed saa nøie sammenhængende:
Objectivitetens og Subjectivitetens. Erkjendelsen har nærmest
Objectivitetens Præg; gjennem Bestemtheden ved Andetheden, gjennem
Subjectets Forhold til de for alle erkjendende Individer fælleds
almeengyldige Tankeformer og Tankelove, der bestemme det med
Nødvendighed. Men Erkjendelsen fuldender sig kun i den Subjectet
omfattende ligesom om det for sig selv klare Selv concentrerede
Erkjendelse, og den er, som Virkeliggjørelse af en Subjectets
Væsensbestemmelse, som ideel Udvidelse af Selvet, Gjenstand for en
Interesse, der kan have det intensiveste Subjectivitetens Præg, som en
Erkjendelsens Lidenskab, og i den er overhovedet, fordi den er
Væsensyttring, det subjective Væsen virksomt efter sin Heelhed, sin
udeelte Personlighed. Paa den anden Side har Villiens Subjectivitet, jo
mere Villiesrealisationen fjerner sig fra det Instinctmæssige,
Driftagtige og Egoistiske, jo mere den bliver en Virkeliggjørelse af
det Almeenfornuftige, af Væsenets Lov, en Væsensvillen, Objectivitetens,
det Almeengyldiges, Bestemthed i sig.
\end{quote}
\dcite{Brøchner, Problemet om Tro og Viden, pp. 134–136  [brochner/problemet-tro-viden:239]}
% brochner/problemet-tro-viden:239  (Problemet om Tro og Viden; pp. 134–136)

\begin{quote}
Fra det kriticerede Standpunkt vilde det her Udviklede vel blive
indrømmet i den Forstand, at Tænkningens Nødvendighed for Handlingen
anerkjendtes, men denne for Handlingen nødvendige Tænkning vilde saa
blive bestemt som praktisk Tænkning, og mellem denne og den theoretiske
Tænkning vilde der blive sat en diametral Modsætning. Nielsen lægger et
Eftertryk paa denne Bestemmelse: praktisk Tænkning, men den faaer
imidlertid hos ham ingensteds nogen virkelig Begrebsforklaring,
ligesaalidt som det tilsvarende Begreb: praktisk Erkjendelse, og det
maa overhovedet være paafaldende, paa hvilken Maade han, under
Udviklingen af sin Theorie, anvender sine psychologiske Kategorier, uden
skarp Begrebsbestemmelse og uden Sondring af det væsentlig
Forskjellige. Saaledes bruges Bestemmelserne: Bevidsthed, Viden,
Tænkning, Erkjendelse, ganske iflæng. — Praktisk Tænkning maatte
nærmest betyde en Reflexion, der bestemmer Sammenhængen mellem Princip
eller Forudsætninger og Consequentser, mellem Midler og Øiemed, paa det
Praktiskes Gebet, altsaa, efter Forudsætningen, en formel Udviklen af
Consequentser af det, der ikke opløses i Reflexionen, og denne
Reflexion, der dog overhovedet ikke vilde faae nogen Plads paa det
paradox Ethiskes Gebet, maatte, idet Handlingen indtræder, tænkes som
absorberet, usynliggjort i Handlingens Enkelthed. En saadan praktisk
Tænken vilde imidlertid i sin Form ikke adskille sig fra den
theoretiske Tænken, der uddrager Consequentser af et endnu ikke af
Erkjendelsen Assimileret, f. Ex. af et blot empirisk Bestemt, eller
som danner Constructionsbegreber. Bestemmelsen af Udgangspunktet
(Forudsætningen) og af Øiemedet kan tænkes som sat ved en
Driftbestemthed eller ud fra den ethiske Subjectivitets Villiesbestemmen.
I første Tilfælde er den det, der skal gjennemklares ud fra
Væsenserkjendelsen, i sidste Tilfælde har den en saadan
Væsenserkjendelse til Forudsætning, men Væsenserkjendelsen er paa een
Gang theoretisk og praktisk. I den udprægede ethiske Charakteer kan den
praktiske Tænkning antage en saadan Umiddelbarhedens Form, at der
handles tilsyneladende reflexionsløst, ligesom ifølge en
Naturnødvendighed, der har Analogie med Geniets Virken. Men denne
reflexionsløse Handlen er ensartet med den Virken af den udviklede
Erkjendelse, der tilsyneladende overspringer Mellemleddene og
umiddelbart naaer Resultatet, men i Virkeligheden gjør det ved, ifølge
Reflexionens Uddannethed, uendelig hurtigt at gjennemløbe
Mellemleddene. Naar saaledes den praktiske Tænkning i sit Væsen —
thi Væsenet bestemmes ved Formen — bliver liig med den theoretiske,
der forudsættes gjennem og i den praktiske, saa gjælder det Samme om
den „praktiske Erkjendelses“ Forhold til den theoretiske. De kunne
ikke adskilles ved den Tænkningens Form, ved hvilken de komme tilveie;
thi den praktiske Erkjendelse, der individualiserer den almene Lov i
dens Anvendelse paa Subjectet, giver den existentiel Anvendelse, følger
netop de selvsamme Love som den theoretiske combinerende Tænkning.
Differentsen skulde da ligge i, at den praktiske Erkjendelse
(Erkjendelsen „til praktisk Brug“) tages i Kantisk Betydning om det
ud fra en Subjectets Trang Postulerede, der ikke retfærdiggjøres ved
den theoretiske Erkjendelse [note: Af Troens Gud skal der saaledes kun gives en praktisk
Erkjendelse, jfr. Afhandl. om theor. og prakt. Erkjendelse. S. 72.]. Men et saadant
Postulat, der, forat have
Gyldighed, maa have sin Grund i Subjectets Væsensbestemmelse, i det
Aprioriske i det Ethiske, har kun sin Sandhed ved at kunne erkjendes
gjennem Væsenets Selvaabenbarethed i Erkjendelsen. Postulatet er en
Anticipation i Kraft af en i Bevidstheden virkende apriorisk
Bestemmelse, svarende til de Anticipationer, der kunne finde Sted paa
den theoretiske Erkjendelses Gebet, og paa hvilke den almindelige
psychologiske Iagttagelse og Erfaring, f. Ex. af Barnets Anvendelse af
Begreber, og Philosophiens Historie frembyder utallige Exempler. Hos
Nielsen sættes der imidlertid en absolut Modsætning mellem theoretisk 
og praktisk Erkjendelse. Der siges, at Valgfrihed, Valg, Beslutning,
kun kunne være Gjenstand for den sidste, ikke for den første. De skulle
vel falde ind under Tankens Belysning, men i Principet ligefuldt være
Villieskategorier og deres Erkjendelse derfor væsentlig praktisk. Det,
der erkjendes theoretisk, skal ikke kunne erkjendes praktisk, fordi i
enhver afgjørende Erkjendelse Indholdet bliver bestemt ved Maaden og
Maaden ved Indholdet. I disse Udsagn er praktisk Erkjendelse brugt som
enstydig med praktisk Bevidsthed, d. v. s. med den umiddelbare
Bevidsthed om Subjectets praktiske Bestemmethed: en Bevidsthed, der
naturligviis i sin Umiddelbarhed sondrer sig fra den paa Reflexion
hvilende Erkjanden, men i hvilken Forholdet er det samme som overhovedet
i Erfaringsforholdet til Gjenstanden. Indeslutter den praktiske
Bevidsthed ikke blot Bevidstheden om sin Gjenstands Væren, men om sin
Gjenstands Gyldigshed, saa viser dette kun, at den i Gjenstanden
forholder sig til et Apriorisk, men ikke, at dette Aprioriske er et for
den theoretiske Erkjendelse Utilgængeligt. Er den praktiske Bevidstheds
Indhold Væsensbestemmelse i Subjectet, maa det kunne blive
gjennemsigtigt i theoretisk Erkjendelse, begrebet ud fra
Væsenserkjendelsen. I dette Forhold indrømmer Nielsen dog kun en
„Afspeiling“, ikke en virkelig Begriben, en \tlg{tilegnende} Erkjanden af
Gjenstanden efter dens Virkelighedsbetydning. Det Gode, som praktisk
Livsidee, siges at indeholde et Umiddelbarhedens Overskud, en Villiens
Oprindelighed, der er principielt utilgængelig for Videnskaben, og
derfor bliver Forholdet kun Afspeilingens. Men denne Betegnelse:
Afspeiling, er ikke en Kategorie, men kun en uheldig Metaphor, der,
hvorledes vi end see paa Sagen, fører ind i Modsigelser. Ved dette
Billede skal det Forhold udtrykkes, i hvilket det med Viden absolut
Uensartede er værende for Viden. Idet Intet absolut skal falde udenfor
Viden, men der dog statueres et med Viden absolut Uensartet, bliver det
nødvendigt at bestemme Forholdet saaledes, at Viden optager ligesom et
Billede af Tingen, men ikke det Afbildede selv efter dets
Virkelighedsbestemthed; den skal saaledes i Forholdet til det af et
andet Princip bestemte Ethiske, i „theoretisk Afspeiling for sig selv
kunne eftergjøre de Bevægelser, der existentielt, punktuelt, ja
antirationelt foregaae i Villien“; men om en assimilerende Begriben
skal der ikke være Tale. Men denne Opfattelse af Forholdet viser sig
let som uholdbar. Fastholdes det Afspeiledes absolute
Uensartethed, vil end ikke en Afspeiling kunne tænkes som
mulig, thi mellem det absolut Uensartede kan intet
Vexelvirkningsforhold finde Sted. Hvert Princip vilde udfolde sine
Bestemmelser af sig, ligegyldige mod hinanden og uden nogen gjensidig
Bestemmen, og man vilde, paa Grund af den absolute Uensartethed, i den
uafhængige Parallelisme ikke kunne faae Noget, der svarede til Spinozas
idea corporis: en Bestemmelse, der hos ham, idet den gjør det
med Tanken absolut uensartede Legemlige til Tankens Indhold, kun er et
Postulat, da Tankens Begreb Intet indeholder, der viser hen til Legeme.
Afspeilingens Mulighed er altsaa, under den antagne Forudsætning, et
uberettiget Postulat. Men selv om man per impossibile
supponerede, at det med Viden absolut Uensartede kunde være tilstede i
Viden, vilde Forestillingen om en Afspeiling dog kun føre til
allehaande Modsigelser. Den Afspeiling, fra hvilken Udtrykket hentes,
er, som det fremgaaer af den Maade, paa hvilken Betegnelsen bruges, den
normale Speilflades troe Afbilden af det, der afspeiles ɔ: af en
legemlig Gjenstands Overflade. Ved en saadan er der en virkelig Lighed,
en virkelig Tilsvaren mellem Speilbilledet og det umiddelbart for Synet
Værende. Men naar Afspeilingens Bestemmelse anvendes paa Forholdet
mellem Viden og Villie, paa hiins ideelle Gjengiven af dennes
Bestemmelser, saa er der ingen Tilsvaren, idet i den ideelle Gjengiven
af det ideelle Væsentlige dette Væsentlige ikke gjengives efter sin
Væsensbestemthed, og saa maa den Lighed, der skulde berettige
Anvendelsen af Begrebet, forsvinde paa Grund af det optagende Mediums
absolute Formforskjellighed fra det Optagne. Sætter man nemlig den
afspeilende Viden som under Afspeilingen virkende efter sin
Eiendommelighed, saa bliver det „Afspeilede“ i Realiteten ikke
afspeilet, men absolut forvandlet ved det optagende Mediums med det
uensartede Natur. Vilde man undgaae dette ved at lade Viden, trods
Selvvirksomheden, ikke forvandle det Afspeilede, saa vilde det
indeslutte, at den, imod Forudsætningen, var ensartet med det, og da
blev Forholdet Mere end en Afspeiling. Vilde man begrunde
Ikke-Forvandlingen ved, at Viden forholdt sig absolut passiv, saa vilde
en saadan Passivitet være uforenelig med Videns Væsen, og Afspeilingen,
der fra det Afspeiledes Side maatte sees som en Indvirken paa det
Afspeilede, vilde blive utænkelig, da en Indvirkning forudsætter
Tilbagevirkning, altsaa Activitet. Som en Indvirkning fra det
Afspeilede maatte i ethvert Tilfælde Afspeilingens Tilbliven forklares;
thi der kan ikke ved en immanent Udvikling ud fra Viden selv med logisk
Nødvendighed naaes til det Afspeilede, da dette er det absolut
Fremmede, der ikke kan komme frem i Sammenhæng med Videns Bestemmelser.
Men Indvirkningen viser jo gjennem det Vexelvirkningsforhold, den
forudsætter, hen til Ligheden i Væsen, til Væsensenhed; den ophæver
altsaa den Forudsætning, hvorpaa Forestillingen om en blot Afspeiling
hviler. Hvorledes man altsaa end søger at opfatte Forholdet, bliver
Betegnelsen Afspeiling uheldig og ubrugelig til at afgive en kategorisk
Bestemmelse.
\end{quote}
\dcite{Brøchner, Problemet om Tro og Viden, pp. 146–150  [brochner/problemet-tro-viden:249]}
% brochner/problemet-tro-viden:249  (Problemet om Tro og Viden; pp. 146–150)

\begin{quote}
3. Ved Siden af de Beviser for den postulerede absolute
Uensartethed mellem Viden og Villie, der kunne søges i
Kjendsgjerninger, men, som det har viist sig, kun gjennem en uberettiget
Fortolkning af disse, vilde der saa kunne stilles et Beviis, bygget paa
en Analogieslutning fra Forholdet mellem den kunstneriske Phantasies
Virken og Erkjendelsen af dens Productioner, idet den supponerede
absolute Uensartethed i dette, som analogt med Villiens og Videns
forudsatte, Forhold benyttedes til at støtte Paastanden om hiin
postulerede Uensartethed. Der kunde saa siges: Det Skjønne og det Sande ere absolut
uensartede, hvad der sees deraf, at en æsthetisk Analyse aldrig kan give Kunstværket: det
Skjønne kan ikke tænkes saaledes ind i Tanken, at de to Ideer blive identiske. Men hvad der
gjælder om det, gjælder ogsaa om det Gode. Villien, Tanken og Anskuelsen afspeile alle den
hele Idee. — Med Hensyn til denne Parallel vil det først være at undersøge, hvorvidt
Analogieslutningen vilde være berettiget, hvilket afhænger af, om Phantasiens og Villiens
Differents fra Erkjendelsen har en væsentlig Lighed, og dernæst, hvorledes den kunstneriske
Frembringelses Forhold til Erkjendelsen er at opfatte. — Hvad det første Punkt angaaer,
saa viser der sig i Phantasiens og Villiens Forhold til Viden en Forskjel, der maa faae en
afgjørende Betydning. I den kunstnerisk frembringende Phantasie er
Umiddelbarhedsbestemmelsen en anden end i Villien: det er en Naturbestemthedens
Umiddelbarhed, i hvilken Aandsvirksomheden træder frem i den genialt skabende Phantasie; det
er en Aandsbestemthedens Umiddelbarhed, der gjør sig gjældende i Villien som Villie ɔ: i den
over den blotte Drift hævede bevidste Villie. Hvad der i den kunstneriske Genialitets
umiddelbare Forbindelse og Sammensmeltning af Natur og Aand gjør sig gjældende, kan ikke
uden Videre overføres paa Villiesvirksomheden, og vilde derfor, selv om det viste hen til en
absolut Uensartethed mellem Phantasie og Viden, ikke bevise Noget med Hensyn til Villiens og
Videns Forhold. Men at det ikke viser hen til en saadan Uensartethed, bliver tydeligt ved
Betragtningen af det andet ovenfor udhævede Forhold. Med Hensyn til dette betragte vi
særskilt den ved Geniets Naturgave betingede active Frembringen af det Skjønne og den
receptive \tlg{Tilegnelse} af det genialt Frembragte i en ved Anskuelse ledet Erkjenden.
Hvad den geniale Frembringelse angaaer saa forholder det sig med den paa Phantasiens Gebet
som paa Videns: den er betinget ved et af Individets frie Beslutning uafhængigt, oprindeligt
ɔ: forud for Individets Selvbestemmen givet, Anlæg, i hvilket det Aandelige gaaer
i umiddelbar Enhed med Naturbestemtheden, ligesom under Geniets Udvikling den
kunstneriske Reflexion gaaer i umiddelbar Enhed med, ligesom smelter sammen med, den
geniale Frembringen. At der bliver en Nuance i Forholdet mellem Naturbestemtheden og
Aandsbestemtheden i den Genialitet, der forholder sig til Erkjendelse, og i den, der
forholder sig til kunstnerisk Phantasie, faaer i denne Sammenhæng ingen Betydning: her
falder Eftertrykket kun paa Naturanlæggenes Uafhængighed af den bevidste
Virksomhed og den frie Selvbestemmen. Naar der tænkes paa denne Uafhængighed, kan det
naturligviis med rette siges, at „den æsthetiske Analyse aldrig kan give Kunstværket“,
ɔ: at man ikke ved Erkjendelsen kan naae til at give sig Betingelserne for en genial
Produceren paa det kunstneriske Gebet; at man ikke ud fra Erkjendelsen kan naae til
Frembringelsen af det originale og som saadant individuelt betingede Kunstværk. Men dette
beviser Intet med Hensyn til det foreliggende Spørgsmaal. — Hvad Erkjendelsens Forhold
som receptivt opfattende det kunstnerisk Frembragte angaaer, saa bliver \tlg{Tilegnelsen} vel,
som overhovedet \tlg{Tilegnelsen} af det concrete Givne, en, om end ikke ved faste Grændser
bestemt, approximerende \tlg{Tilegnen}, men under denne Approximation bevares Harmonien mellem
Erkjendelsens Fordringer og Kunstværkets Skjønhedslov, og Erkjendelsen stiller sine
Fordringer i Overensstemmelse med en, ikke blot i Erkjendelsen afspeilet, men i
Erkjendelsen indoptagen Skjønhedsidee. Den analyserende Erkjendelses Optagen af
Kunstværket er vel ledet af en Anskuelse, der i Receptivitetens Form danner Parallelen til
den kunstnerisk frembringende Phantasies Productivitet; men denne Anskuelses Umiddelbarhed
omsættes i Reflexion ligesom Erkjendelsen omsættes i Umiddelbarhed, og det ved
Genialitetens aandelige Naturmagt Frembragte bliver, efter at være frembragt, den
erkjendende Aands Eiendom i en endnu høiere Betydning end den concrete Naturfrembringelse
bliver det, — Ikke heller ad denne Vei kan derfor den absolute Uensartethed hævdes.
\end{quote}
\dcite{Brøchner, Problemet om Tro og Viden, pp. 150–153  [brochner/problemet-tro-viden:250]}
% brochner/problemet-tro-viden:250  (Problemet om Tro og Viden; pp. 150–153)

\begin{quote}
I den sidste Form for Dualismen reflecterer sig den dualistiske Modsætning mellem det
Antirationelle og det Rationelle (jfr. Gr. L. II, 278 Anm.), den Modsætning, hvis Bevaring
for Nielsen er det oprindelige Diemed, og idet nu Dualismen i Videns Idee reflecterer sig i
den menneskelige Viden, saa viser den sig her deels directe gjennem Dobbeltheden i Viden som
\tlg{tilegnende} og som blot afspeilende: en modsigende Dobbelthed, der maa medføre en Spaltning i
Vidensindholdet [note: Var „Afspeilingen“ en væsentlig Bestemmelse i Viden, en
væsensnødvendig Videns Form, maatte der ved en immanent Udvikling naaes til det
„Afspeilede“; men dette er en Umulighed paa Grund af dettes absolute Uensartethed med det
begrebne Vidensindhold. Viden maa som Viden udfolde sine Muligheder systematisk, men
Mulighederne vilde, naar de skulde omfatte det absolut Uensartede, umiddelbart modsige den
Enhed, den systematiske Udfoldning forudsætter.]; deels indirecte derigjennem, at Viden, idet
den ifølge sit Princip skal udfolde sig i Selvstændighed, dog skal tillæmpe sine Bestemmelser
efter det Antirationelle. Fordringen af en saadan Tillæmpning ligger i det tidligere anførte
Udsagn, at Videnskaben skal lade Miraklets Mulighed staae som det, der ikke kan modbevises:
et Udsagn, der trodser de Modsigelser, til hvilke Indrømmelsen af denne Mulighed
fører [note: I Forestillingen om Miraklet træder netop den Modsigelsen indesluttende
Dualisme frem; thi naar Miraklet tages i sin egentlige Betydning, betegner det det abstract
Ideelles vilkaarlige Virken med Negation af ethvert Realt, af alle reale Betingelser: den
samme Virken, der er sat i Skabelsesforestillingen. Muligheden af Miraklet i denne Betydning
søger Nielsen at hævde ved sin efter den religiøse Forestilling tillæmpede Bestemmelse af den
absolute Subjectivitet, der frit skal kunne transscendere alle Betingelser. Men derved faaer
man enten umiddelbart den angivne Dualisme mellem et Ideelt og det Reale, eller, naar den
ideelle Subjectivitet dog fra Begyndelsen af tillægges Realitetsbestemmelser, bliver dets
Transscenderen af disse, ved hvilken de negeres, idet nye sættes, kun tænkelig som et Brud,
hvorved der altsaa sættes en Dualisme i selve det absolute Princip og i den reale
Tilværelse. Og Dualismen i Principet bliver saa ikke blot den dualistiske Sondring mellem det
Absolute, bestemt ved de oprindelige og ved de senere satte Realitetsbestemmelser; men mellem
det Absolutets Idealitetsside og dets Realitetsside, hvilken sidste sættes som et Intet, idet
den transscenderes. — Vilde man begrunde Umuligheden af at modbevise Miraklets Mulighed
derved, at Viden, anerkjendende sin Grændse, maa anerkjende Personlighedens Fordringer, og at
Personligheden gjennem Fordringen af en hellig Gud fordrer en almægtig, skabende Gud, saa er
denne Vending allerede belyst i det Foregaaende. Personlighedens Krav er i sin Sandhed
Væsenets Krav, der ikke kan modsige den Væsenet udtrykkende Viden.], og staaer i Strid med
andre bestemte Udsagn [note: F. Ex. Gr. L. II, 190. „En Almagt, der ubetinget skaber —
— af Intet, er ikke blot for os ufattelig, men i Ideen ubegribelig“.] . En saadan
Tillæmpning bliver synlig i Opfattelsen af Magtbegrebet, idet Magtens Begreb hos
Nielsen vakler over mod Almagtens Forestilling, og den bliver synlig idet en
Forestillingssondring fastholdes mellem den absolute, i Villie concentrerede mægtige
Subjectivitet og den reale Gjenstand; thi gjennem denne Forestillingssondring trænger sig
Forestillingen ind om en almægtig ɔ: absolut vilkaarlig, skabende Villie, der i ubetinget
Frihed hæver sig over alle Betingelser: den Forestilling, paa hvilken i sidste Instants
Paastanden om det for Videnskaben Umulige at modbevise Miraklets Mulighed hviler.
\end{quote}
\dcite{Brøchner, Problemet om Tro og Viden, pp. 174–176  [brochner/problemet-tro-viden:285]}
% brochner/problemet-tro-viden:285  (Problemet om Tro og Viden; pp. 174–176)

\subsection{\textit{Et Svar til Professor R. Nielsen} (1868)}

\begin{quote}
var, at det, idet den abstracte Modsætning definitivt fastholdes, vilkaarligt
bliver ignoreret, at paa de forskjellige Subjectivitets-Stadier Subjectets
Bestemthed —

 
hvad der især tydelig fremtræder paa det Stadium, der betegnes som den uendelige
Subjectivitets — netop er i Kraft af hvad der er bestemt som Tilværelsens
objective Princip, og at den udtrykker dette. Det Andet var, at det bliver
overseet, at den subjective Sandhed netop er den \tlg{tilegnede} og i \tlg{Tilegnelsen} med
sig væsentlig lige objective Sandhed. Ved ikke at forstaae, at den sidste Passus
generelt udtrykker Forholdet, og ved, paa en høist vilkaarlig Maade, at referere
den exclusivt til et enkelt af de nævnte Stadier, faaer Prof. N. ud, at jeg
aldeles skal have misforstaaet hans Ord om dette, og at jeg skal have viklet mig
ind i gyselige Modsigelser. Saasnart mine Ord forstaaes rigtigt, seer man, at de
ingen Misforstaaelse af Prof. N. forudsætte, og at de ere fuldkommen correcte.
Det viser sig især tydeligt paa Stoicismens og Hedonismens Standpunkt, at
Subjectets Bestemthed udtrykker Tilværelsens objective Princips Bestemthed. Dette
behøver ingen Forklaring for den, der blot kjender lidt til Philosophiens
Historie, og ikke heller behøver det en Forklaring, at Tilværelsens Princip
betegnes som Tilværelsens objective Princip. Og naar saa den generelle Sætning om
Forholdet mellem den subjective og den objective Sandhed anvendes paa disse to
Systemer, saa bliver altsaa Consequentsen, at den relative Sandhed, som det
enkelte ensidige, men i sin Ensidighed et Gyldigt indesluttende System eier, er
Maalet for den subjective Sandhed hos den, der practisk \tlg{tilegner} sig Systemets
objective Sandhed [note: Paa Prof. N.'s snurrige Argumentation fra det, at en Stoiker kan „sætte sig
 ligesaa sikkert ind i“ den hedoniske Theorie som en Hedoniker og dog ikke
 blive Hedoniker, skal jeg kun henlede Læserens Opmærksomhed: den behøver ingen
 Commentar!].
\end{quote}
\dcite{Brøchner, Et Svar til Professor R. Nielsen, pp. 16–17  [brochner/svar-nielsen:51]}
% brochner/svar-nielsen:51  (Et Svar til Professor R. Nielsen; pp. 16–17)

\begin{quote}
det rette Videnskabelige. Egentlig skal nok Feilen ligge i, at jeg ikke har hævet
mig til Indsigt i „Forholdet mellem System og Schema“. Her skal jeg nu villig
indrømme, — thi Noget bør man jo være saa billig at indrømme — at jeg
virkelig ikke har kunnet naae til at \tlg{tilegne} mig Prof. N.'s Systematik. Men
Prof. N. har ogsaa øiensynligt udviklet det Systematiske til en Høide, der ikke
er kjendt i nogen tidligere Philosophie. Prof. N.'s philosophiske Skrifter ere
nemlig i en saa forbausende Grad gjennemtrængte af det Systematiske, at man kan
skjære dem over i større og mindre Stykker, og naar man saa sætter Stykkerne af
de forskjellige Bøger sammen efter Pagina og indbinder dem, saa er der strax et
System. Der manglede endnu kun, forat Prof. N. kunde naae det Systematiskes
„absolute Grændse“, at han indrettede sine Skrifter saaledes, at de, ligesom
visse af Middelalderens cabbalistiske Formler, kunde læses ligegodt bagfra

 
og forfra. Forat nu Læseren ikke skal troe, at jeg gjør Løier, og forat Prof. N.
ikke atter skal bebreide mig, at jeg taler almindeligt i al Almindelighed om det
Almindelige, saa skal jeg nu give en „Sagudvikling“, saa „punktuel“, at Prof. N.
neppe kan ønske den punktuellere.
\end{quote}
\dcite{Brøchner, Et Svar til Professor R. Nielsen, pp. 23–24  [brochner/svar-nielsen:73]}
% brochner/svar-nielsen:73  (Et Svar til Professor R. Nielsen; pp. 23–24)

\subsection{\textit{Om det Religiøse i dets Enhed med det Humane} (1869)}

\begin{quote}
Begreb i et andet Forhold til Erkjendelse og Tanke end f. Ex. til Villien;
den bliver det abstracte, formelle, Udtryk for Erkjendelsen, det, der
kan udfoldes til en \tlg{tilegnende} Erkjenden, til begribende
Tænkning, til Aandens Gjennemsigtighed for sig i sit Indhold; og det, at den
træder i Forhold til de forskjellige Aandsfunctioner, idet den bliver Formen
for deres Væren for Selvet, viser netop hen til Erkjendelsens ideelle
Principat, og til Muligheden af en Selvets Gjennemsigtighed for sig gjennem
Erkjendelsen i sine forskjellige eiendommelige
Yttringsformer [note: Med Hensyn til den Beskyldning, en Kritiker har gjort
mod mig for Sammenblanding af det væsentlig Forskjellige, bemærker jeg kun,
hvad enhver eftertænksom Læser af det tidligere Skrift vel vil have seet, at
jeg ikke identificerer de umiddelbare Bevidsthedsformer med den
udfoldede videnskabelige Erkjendelse, men kun hævder, at det, der
træder frem i hine Umiddelbarhedens Former, deels er et saadant, der kan
udfoldes i klar Erkjendelse, deels er et resulterende Umiddelbart, og
nærmere et Saadant, der er resulteret af en forudgaaende Erkjendelse.
Ligesaalidt sammenblander jeg det Subjective med det Objective,
men jeg hævder, at hint, som Udtryk for Subjectets Væsen, har den samme
Gjennemskuelighed for Erkjendelsen ɔ: for en theoretisk Erkjendelse, som
hint Umiddelbare.].
\end{quote}
\dcite{Brøchner, Om det Religiøse i dets Enhed med det Humane, p. 6+  [brochner/det-religioese:28]}
% brochner/det-religioese:28  (Om det Religiøse i dets Enhed med det Humane; p. 6+)

\begin{quote}
Bestemmes saa Modsætningen mellem de to Erkjendelsesarter saaledes, at den
theoretiske Erkjendelse er bevæget af Nødvendighed, medens den
praktiske „grunder i Frihed“, idet hiin, som fortabt i det Almene,
tvingende bestemmes af de uimodsigelige Vidensgrunde, denne, som betinget ved
Selvhensyn, hviler paa Samvittighedsgrunde, saa viser ogsaa denne Modsætning
sig som dialektisk. Den theoretiske Erkjendelse er behersket af Nødvendighed
hvor Slutningens Forudsætninger ere fuldt gjennemsigtige, hvor Beviset faaer
Stringents ved klare og fuldstændige Præ-

 misser. Men ligesaa fuldt er den „praktiske Erkjendelse“ bestemt med
Nødvendighed hvor disse Betingelser ere tilstede, hvor Forudsætningerne ere
givne klare og fuldstændige. Hvad der tilveiebringer Skinnet af en abstract
Modsætning, er Forskjellen i Concretion mellem den praktiske
Erkjendelses Gjenstand og visse Sphærer af det Theoretiske. Paa det concrete
Ethiskes Gebet kan der vedblive at være en tilsyneladende uophævelig
Differents netop paa Grund af Præmissernes Ufuldstændighed, paa samme Maade
som indenfor det Theoretiske, paa de concretere Gebeter. Differentsen kan
holde sig, fordi de mangfoldige Forudsætninger ikke ere klarede og derfor
partielt suppleres paa individuel Maade. Saaledes bliver, under samme
Forudsætninger, Nødvendigheden den samme. Og den Nødvendighed, der bestemmer
den theoretiske Erkjendelse, er, naar Erkjendelsen er begribende og som
saadan \tlg{tilegnende}, et Udtryk for selve det erkjendende Væsens
Lovmæssighed, og omsætter sig saaledes i Frihed. I den ufrie
Nødvendigheds Form fremtræder den theoretiske Erkjendelse kun hvor det ved
Sandsningen Givne, i dets Fremmedhed, er bestemmende. Men til dette Stadium
svarer det Stadium i den praktiske „Erkjendelses“ Udvikling, hvor det
Selvhensyn, der bestemmer den, er Hensynet til et uklart Driftagtigt, ved
hvilket Villien bringes ind under Naturnødvendigheden. Kun som Udtryk for et
Væsentligt har Subjectets Trang, der bestemmer hiin „Erkjendelse“,
Berettigelse, men som Udtryk for Væsenet har denne Trang paa een Gang den
sande Nødvendigheds og Frihedens Præg, medens Driftbestemthedens saakaldte
„Frihed“ er ufri Nødvendighed ligesom Sandsningens Nødvendighed. Baade i den
theoretiske og i den praktiske Erkjendelse fremtræde altsaa begge de
\end{quote}
\dcite{Brøchner, Om det Religiøse i dets Enhed med det Humane, pp. 12–13  [brochner/det-religioese:44]}
% brochner/det-religioese:44  (Om det Religiøse i dets Enhed med det Humane; pp. 12–13)

\begin{quote}
Den theoretiske Erkjendelse skal være contemplativ, gaaende op i
Begrundelsen af det Objective; men den er det kun i den ved de tidligere
Bestemmelser begrændsede Betydning. Fra den beskuende, selvforglemmende
Fordybelse i det Almene og Objective søger den tilbage til det Punktuelle og
til den Selverkjendelse, til hvilken en uimodstaaelig indre Trang driver den,
og idet den stræber hen mod sit Maal, er der i Erkjendelsens Kamp for
\tlg{Tilegnelsen} af Virkeligheden, i den sande Erkjendelses Kamp mod det
Usande, et Tilsvarende til Villiens Conflict med Existentsen, til dens Kamp
for sin Realisation, til den gode Villies Kamp mod det Onde.
\end{quote}
\dcite{Brøchner, Om det Religiøse i dets Enhed med det Humane, p. 14+  [brochner/det-religioese:48]}
% brochner/det-religioese:48  (Om det Religiøse i dets Enhed med det Humane; p. 14+)

\begin{quote}
2cm0.4pt — Søge vi, efterat have belyst det almene Forhold mellem det Theoretiske og det
Praktiske, den nærmere Bestemmelse af Erkjendelsens sande Form: af den
Form, der er dens Udviklings Maal, saa bliver den at finde ud fra Bevidsthedens
Grundforhold. I det Forhold, der i Menneskeaanden er givet mellem
Selvbevidsthed og Bevidsthed, ligger

 Nødvendigheden af en Stræben efter at bringe alt Bevidsthedsindhold ind under
Selvbevidsthedens Enhed, at ophæve Bevidsthedsgjenstandens Fremmedhed og at
gjøre den gjennemsigtig for Aanden i en begribende Erkjenden, i hvilken Selvet
har sin ideelle Udfoldelse og sin Selvaabenbarelse, idet det ideelt \tlg{tilegner} sig
Virkeligheden. Først i en saadan begribende Erkjenden, i hvilken
Bevidsthedsindholdet for Aanden danner en Enhed, svarende til Selvbevidsthedens
Enhed, er den Modsigelse hævet, der ligger i Bevidsthedsforholdet: at
Gjenstanden er i min Bevidsthed som min Bevidstheds Bestemmelse, og dog
efter sit Indhold er et for mig Fremmedt, og først i den er ligeledes den
Modsigelse hævet, at Aanden, der i Selvbevidstheden paa uendelig Maade forholder
sig til sig selv og er i Enhed med sig selv, tillige er fremmed for sig selv ved
det fremmede Indhold, Selvet omfatter.
\end{quote}
\dcite{Brøchner, Om det Religiøse i dets Enhed med det Humane, pp. 19–20  [brochner/det-religioese:63]}
% brochner/det-religioese:63  (Om det Religiøse i dets Enhed med det Humane; pp. 19–20)

\begin{quote}
Vi bestemte der det religiøse Forhold som „det absolute
Forhold [note: Hvis det for „omvendte“ Dialektikere skulde synes
modsigende at tale om et absolut Forhold, fordi Forhold udtrykker
Relation, behøve vi vel blot at minde om det Absolutes Forhold til sig
selv, og om Aandens Uendelighedsforhold til sig selv i Selvbevidsthedens
Concretion.] til det Absolute“, altsaa ved et Udtryk, der ikke blot i
Formen nærmer sig til S. Kierkegaards almindelige Betegnelse af den endnu
ikke christelige Religiøsitets Stadium, hvilket han charakteriserer ved „den
absolute Retning mod det absolute τέλος“, men som allerede er brugt af ham
hvor han sammenstiller det religiøse absolute Forhold med de relative
Forhold, der hos ham ogsaa omfatte de ethiske Forhold, i hvilke hiint ikke
skal udtømme sig eller blive concret. Medens vi saaledes \tlg{tilegne} os hans
\end{quote}
\dcite{Brøchner, Om det Religiøse i dets Enhed med det Humane, p. 30+  [brochner/det-religioese:88]}
% brochner/det-religioese:88  (Om det Religiøse i dets Enhed med det Humane; p. 30+)

\begin{quote}
a) Betingelsen for, at Erkjendelsen for Menneskeaanden kan blive Formen for et
absolut Forhold, er den, at den kan bevare Gjenstandens sande Uendelighed i
Forholdet. At den maa kunne bevare Uendelighedsbestemtheden, ligger forsaavidt
allerede i Erkjendelsens Form, som denne er bestemt ved det Almene, og dette
Almene, hvis Væren som Alment for den menneskelige Bevidsthed er betinget ved
Væsenets indre Uendelighed, staaer i Væsenssammenhæng med det Uendelige. Men
forat gjennem denne Erkjendelsens Form det Uendelige kan bevares som det i
Sandhed Uendelige, maa Erkjendelsens Almindelighedsform være i concret Enhed
med Begrebets andre Grundformer. Forat den kan være dette, forudsættes, at
Erkjendelsen i Forholdet til det Absolute ikke staaer i abstract Modsætning til
hvad der ligesaafuldt er en Grundfunction af Aanden, at den ikke bliver en
ensidig, fra det Supplerende og Totaliserende isoleret Form for Aandens
Virksomhed. Som saadan maatte den f. Ex. uundgaaelig opfattes naar den ved en
principiel absolut Uensartethed sondredes fra Villien; thi i denne Sondring
vilde Villien \tlg{tilegne} sig Enkelthedsbestemmelsen som sin Grundform og lade
Erkjendelsens Almindelighedsform blive den abstracte Almindeligheds, medens
Villiens Enkelthedsform, i sin Sondring fra det Almene, netop kun blev den
umiddelbare Naturenkelthedsform. I denne abstracte Modsætning vilde ingen af
Aandsfunctionerne blive opfattet efter sin Sandhed, ɔ: som Yttring af den
udelte, totale Menneskeaand, og ingen af dem vilde komme til i Sandhed at
forholde sig til det Uendelige, fordi det Uendelige ikke i sin sande
Uendelighed vilde kunne være for nogen af dem.
\end{quote}
\dcite{Brøchner, Om det Religiøse i dets Enhed med det Humane, p. 44  [brochner/det-religioese:124]}
% brochner/det-religioese:124  (Om det Religiøse i dets Enhed med det Humane; p. 44)

\begin{quote}
Til Forestillingen om en Skabelse staaer den angivne Opfattelse i det
Forhold, at den fra denne Forestilling \tlg{tilegner} sig den Sandhed, at det
Reale er begrundet i det Ideelle, Aandelige og bestemt ved dette; medens den
seer det Usande i Skabelsesforestillingen deri, at denne ikke lader det Reale
være af væsentlig Betydning for det Ideales Væren, ikke lader det ideelle
Absolute i dets Ubetingethed betinge sig selv gjennem det Reale, at den ikke
anerkjender dette Reales Væsens-Enhed med det Ideale. Den fastholder, at det
Ubetingede kun er ubetinget som det, der gjennem et System af reale
Betingelser betinger sig selv; at det Absolute er principium sui,
causa sui, idet det er Princip for et Andet, og at dette Andet, som
Principets Andet, kun kan være dette ideelle Princips egen Bestemmelse.
\end{quote}
\dcite{Brøchner, Om det Religiøse i dets Enhed med det Humane, p. 71  [brochner/det-religioese:190]}
% brochner/det-religioese:190  (Om det Religiøse i dets Enhed med det Humane; p. 71)

\begin{quote}
Gjennem Gud faaer Mennesket dette, der, ved at sættes i Gud, Aandens væsentlige
Gjenstand, allerede indirecte er udtalt som en væsentlig Bestemmelse ved den
menneskelige Natur, i uendelig rigere Maal tilbage. Religionens Hemmelighed er
den, at Mennesket objectiverer sit Væsen, men atter gjør sig til Object for
dette objectiverede, til et Subject forvandlede Væsen. Vel har Mennesket i
Religionen Gud til Øiemed, men Gud har intet Andet til Øiemed end Menneskets
Vel: den guddommelige Virksomhed bliver Middel for dette, altsaa har Mennesket
i Religionen kun sig selv til Øiemed i og gjennem Gud. Og hvad der, ligesom ved
en religiøs Repulsionskrafts Act, sættes som en Virksomhed af Gud, Menneskets
udsondrede Subjectivitet, fra hvem den religiøse Bevidsthed maa lade enhver
Impuls udgaae, det bliver atter ved en religiøs Attractionskrafts Act, idet den
guddommelige Virksomhed skal virke paa Mennesket, i Mennesket, blive Princip
for Menneskets gode Sindelag, for dets gode Handlinger, gjort til en i sit
Væsen menneskelig Virksomhed. Religionens Udviklingsgang bestaaer derfor
nærmere deri, at Mennesket bestandig \tlg{tilegner} sig Mere af det, hvad det før
tilskrev Gud. Hvad der paa det tidligere Trin var

 Religion, er det ikke længere paa det sildigere: hvad der paa hint var
Atheisme, gjælder paa dette for Religion.
\end{quote}
\dcite{Brøchner, Om det Religiøse i dets Enhed med det Humane, pp. 106–107  [brochner/det-religioese:268]}
% brochner/det-religioese:268  (Om det Religiøse i dets Enhed med det Humane; pp. 106–107)

\begin{quote}
Naar nu Religionens Standpunkt er det praktiske eller subjective, saa ligger
deri, at kun det praktiske, efter sine bevidste, physiske eller
moralske, Øiemed handlende Menneske, der alene betragter Verden i Relation til
disse Øiemed,

 ikke i og for sig, gjælder den for det hele væsentlige Menneske, og det er
i Relation til Mennesket i denne Bestemthed, at Religionens høieste Væsen
bestemmes som et Gemyttets Væsen, som et Væsen, der tilfredsstiller
Gemyttets og Phantasiens Trang. Alt det, der ligger bag den praktiske
Bevidsthed, men er den væsentlige Gjenstand for Theorien — for Theorien i
den oprindeligste og almindeligste Betydning, i Betydning af objectiv Anskuelse
og Erfaring, Fornuft, Videnskab overhovedet — forholder sig ikke væsentlig
til den religiøse Bevidsthed. Religionen mangler, idet den er udelukkende
praktisk, subjectiv, den sande Anskuelse af Universet, af det virkelige
Uendelige, den sande Bevidsthed om Slægten. Og idet nu Naturens og Menneskets
sande almindelige Væsen som et saadant, der ikke kan negeres, gjør sig
gjældende for Aanden, saa søges vel hiin Mangel erstattet i Gud, i et særskilt,
personligt Væsen: Naturens og Menneskets kun for det theoretiske Øie
tilgængelige Væsensbestemmelser lægges uvilkaarligt over i Gud, men deels
saaledes, at de kun udgjøre den religiøse Bevidstheds yderste
Tilknytningspunkt, „Religionens mathematiske Punkt“; deels saaledes, at de
omformes, og i denne Omformning \tlg{tilegnes} af Gemyttet og Phantasien. Idet
Religionen antager Menneskehedens og Naturens tilsyneladende, overfladiske
Væsen for deres sande, indre Væsen, og forestiller sig deres sande esoteriske
Væsen som èt andet Væsen, metamorphoserer den under Personificationen
Menneskets og Naturens sande Væsen til en overmenneskelig og
overnaturlig Personlighed, ubegribelig for Erkjendelsen, virkende gjennem
Underet og gjennem dette tilfredsstillende det religiøse Individs praktiske
Trang i Modsigelse med Fornuftloven.
\end{quote}
\dcite{Brøchner, Om det Religiøse i dets Enhed med det Humane, pp. 107–108  [brochner/det-religioese:272]}
% brochner/det-religioese:272  (Om det Religiøse i dets Enhed med det Humane; pp. 107–108)

\begin{quote}
er der, trods Opfattelsen af den guddommelige Villie som
vilkaarlig, en Mulighed for, at et Tankeindhold kan lægge sig ind i
de religiøse Former, og i dem hævde det Rationelle en væsentlig Betydning,
og dette Tankeindhold, der ligesom fældes ned i hine Former, er deels fremkommet
gjennem en Tankeudvikling, der ikke har sondret sig fra det Religiøse, men er
fremmet ved selve Religionen, deels ifølge en indre Trang \tlg{tilegnet} af den
religiøse Bevidsthed fra en af den relativ uafhængig Tankeudvikling.
\end{quote}
\dcite{Brøchner, Om det Religiøse i dets Enhed med det Humane, p. 123  [brochner/det-religioese:302]}
% brochner/det-religioese:302  (Om det Religiøse i dets Enhed med det Humane; p. 123)

\begin{quote}
Muligheden og Grunden til det saaledes bestemte ethisk Onde søge vi i
selve den Grundbestemmelse i det menneskelige Væsen, ifølge hvilken det, i sin
Universalitet, som ideelt-realt Led i Tilværelsens organiske Heelhed, er en i
væsensgyldig Realitetssondring bestaaende relativ Totalitet, der har Mulighed til
Bevidsthed, om end i den relative Abstractions Form, om sin indre Uendelighed, og
til en fri Forsigværen. Men denne det individuelt Ideelles Mulighed til en
uendelig Væren for sig er en i sig fordobblet Mulighed: den er Muligheden til
den sande udfyldte Forsigværen i Hengivelsen til og \tlg{Tilegnelsen} af det
gjennemvirkende Universelle, og Muligheden til en Afslutten i sig i sin
væsentlige Realitetssondring fra det Absolute, i en indre Uendelighed, der i
Sondringen fra det sande Uendelige bliver en abstract, tom, ved Løsrevetheden
endelig, og derfor selvmodsigende Uendelighed, hvis illusoriske
Udfyldning maa søges i det for den Tilfældige, der kun har Endelighedens
Bestemthed. I hiin første Forsigværen vil Individet sig som indordnet i
Heelheden, underordnet Principet, det vil sig altsaa som dettes eiendommelige
Organ, og kun saaledes kan det, hvad der ligger i Grundforholdet, ville sig efter
sit Væsen. Men idet
\end{quote}
\dcite{Brøchner, Om det Religiøse i dets Enhed med det Humane, p. 151  [brochner/det-religioese:380]}
% brochner/det-religioese:380  (Om det Religiøse i dets Enhed med det Humane; p. 151)

\begin{quote}
Tabet af dets Integritet. Ved begge disse Bestemmelser kan der føres over paa det
Paradoxes Gebet. Opfattes nemlig Forholdet til den guddommelige Villie som
Forhold til en Villie, hvis Lov er den med Erkjendelsens Fornuftlov uensartede, i
sin Vilkaarlighed ubegribelige, saa bliver Synden det Paradoxe, Bevidstheden om
Synden som Synd betinget ved et Troesforhold af paradox Natur, og Syndens
Forsoning bliver betinget ved og Gjenstand for et Troesforhold af samme Art. Og
paa samme Maade forholder det sig med den Opfattelse af Syndens Virkning, der i
Læren om Arvesynden statuerer Menneskenaturens Fordærvelse ved den første
syndige Villiesact. Denne Fordærvelse, saaledes som Kirkelæren opfatter den,
bliver det Ubegribelige. Menneskevæsenet skal tænkes som fordærvet ved en af
selve Væsenet, og, bestemtere, af det endnu i Integritetens Tilstand værende
Væsen, fremgaaende Act. Villiesyttringen, der realiserer en Mulighed, der er
anlagt i Væsenet, skal forvandle Væsenet, ved een Act destruere det. Som Parallel
til dette Ubegribelige staaer saa det Ubegribelige i Restitutionen ved et
guddommeligt Under, ved Guds paradoxe Menneskevordelse, ved Gudmenneskets
Selvopoffrelse, der \tlg{tilegnes} gjennem Troen. I denne Opfattelse unddrager Synden
og Forsoningen sig al Erkjendelse. Men Bestemmelsen af Synd som omfattende
saavel en Strid mod den guddommelige Villie, som en Indvirkning paa
Væsenet, kan fastholdes uden at det Rationelles Gebet forlades. Bestemmelsen
Synd kan fastholdes i Betydningen af Overtrædelse af den guddommelige Villies
Fordring, idet den guddommelige Villies Begreb opfattes paa den Maade, der
blev angiven under Udviklingen af det Absolutes Begreb. Og Begrebet om en
Affection af Væsenets Integritet ved Synden kan fastholdes, idet det opfattes i
Overensstemmelse med hvad der blev udviklet i de foregaaende Afsnit. Ifølge den i
den menneskelige Aand tilstedeværende dobbelte Mulighed: at være og at ville sig
som bevidst indordnet Led i Totaliteten, og at være og at ville sig i Isolationen
efter sin negative Uendelighed,
\end{quote}
\dcite{Brøchner, Om det Religiøse i dets Enhed med det Humane, p. 170  [brochner/det-religioese:424]}
% brochner/det-religioese:424  (Om det Religiøse i dets Enhed med det Humane; p. 170)

\begin{quote}
Til den objective Betingelse for Ophævelsen af Syndestraffen træder saa
den subjective: Troen, som \tlg{Tilegnelse} af den objective
Fyldestgjørelse. Troen som den retfærdiggjørende Tro er Troen paa, at
Christi Død er fyldestgjørende for dette troende Individ. Ved
denne Tro, ikke ved Gjerninger, der, for at have Sandhed, skulle have Troen til
Forudsætning, er det, at Forsoningen vindes.
\end{quote}
\dcite{Brøchner, Om det Religiøse i dets Enhed med det Humane, p. 196+  [brochner/det-religioese:491]}
% brochner/det-religioese:491  (Om det Religiøse i dets Enhed med det Humane; p. 196+)

\subsection{\textit{Bidrag til Opfattelsen af Philosophiens historiske Udvikling} (1869)}

\begin{quote}
Princips Forsigværen gjennem den som Andethed satte An-
dethed. — To Epocher af Philosophien, i deres Princip bestemte ved
væsentlige Ideeforhold og i deres Følge ved en Udviklings Nød-
vendighed, foreligge saaledes historisk som fuldkomment
afsluttede, og ved det Grundproblem, de have stillet,
men ikke løst, vise de hen paa en tredie Hovedform, hvis
Udvikling, der ved selve den abstracte Formel antydes som
betinget af de til det Reale umiddelbart sig forholdende Viden-
skabers \tlg{Tilegnelse} og Assimilation ved den philosophiske Tanke,
tilhører Fremtiden, den hele Fremtid.
\end{quote}
\dcite{Brøchner, Bidrag til Opfattelsen af Philosophiens historiske Udvikling, p. 29+  [brochner/philosophiens-udvikling:87]}
% brochner/philosophiens-udvikling:87  (Bidrag til Opfattelsen af Philosophiens historiske Udvikling; p. 29+)

\begin{quote}
Principet som Stof vil derfor uundgaaeligt skride frem mod
ideellere Bestemmelser: Principer som Vand, Luft, ubegrændset
og ubestemt Stof, ville afløses af Principer som Tal, Væren
( ), begge opfattede som det Stof, hvoraf Tingene bestaae,
som den legemlige Tilværelses Substants, og idet Problemet
om Vorden, om det Mangfoldiges Tilbliven af Principet,
kommer frem som det væsentlige, bliver Vordens Begreb
det, der, medens det, uforklaret, \tlg{tilegner} sig det Sandse-
liges Omraade, luttrer Begrebet Væren til Ideen. Og saa-
snart Ideens Begreb er fundet: det Begreb, til hvilket selve
Stoffet viste hen som til sin Fuldendelse, bliver Ideen sat som
den sande, væsentlige Virkelighed, og Tankens Opgave
bliver: at søge et Udtryk for Enheden af de to Mo-
menter, i hvilket denne Forudsætning er fastholdt.
Men idet denne Opgave kommer frem, saa forudsætter den, at
Tanken bevidst har sondret de to Momenter, der for
Anskuelsen ere forenede i Kunstværkets Enhed, og denne
Enheds Art umuliggjør, som det ovenfor blev paaviist, en ny
Forening ved den græske Tanke, for hvilken Harmoniens
Umiddelbarhed var Enhedens eneste Form. Skjønhedens
Umiddelbarhed taaler ikke Tankens sondrende Magt, derfor
bringer Tanken Splid ind i det, hvorpaa den hviler, og gjen-
nem denne Splid maa denne Tankeform gaae tilgrunde. Mo-
menterne, der i deres Sondring begge tiltvinge sig en Be-
stemmelse af Væren, staae dualistisk imod hinanden, og
den græske Philosophie maa ende i den frugtesløse Bestræbelse
for, efterat Bevidstheden om Momenternes Væsensforskjel er
vakt, atter at finde en Forsoning.
\end{quote}
\dcite{Brøchner, Bidrag til Opfattelsen af Philosophiens historiske Udvikling, p. 36  [brochner/philosophiens-udvikling:116]}
% brochner/philosophiens-udvikling:116  (Bidrag til Opfattelsen af Philosophiens historiske Udvikling; p. 36)

\begin{quote}
En saadan Overgangsperiode var nødvendig ved Mid-
delalderens Slutning før Philosophiens nye Hovedform kunde
fremtræde i sin udviklede Bestemthed. Tænkningen havde i
Scholasticismen, idet dens Gjenstand var det for Tanken Ube-
gribelige, kun vundet en formel Udvikling uden at kunne
udfolde et selvstændigt Tankeindhold. Idet den frigjorde sig
fra Kirkens Herredømme, og, forkyndende en ny Verdens-
anskuelse, vendte sig mod Scholasticismen, maatte den
derfor søge det for denne Kamp nødvendige Indhold, der
maatte være et saadant, i hvilket den befriede menneskelige
Tanke kunde gjenfinde sig og gjenkjende Virkeligheden, hvor
det forelaae som givet, og den fandt det da hos Oldtidens
Philosophie, med hvilken den følte sig i Slægtskab. Den om-
byttede da Scholasticismens tvungne Underkastelse under en
fremmed Autoritet med en selvvalgt Underkastelse under
en beslægtet, i hvilken den nye Tids Tanke, der nu kunde
have Virkelighed og Nærværelse i sit Indhold, kunde vinde
Fylde og Selvstændighed gjennem \tlg{Tilegnelsen} og Udformningen
af dette Virkelighedsindhold, i hvilket den først skulde leve sig
ind. I alle de mangfoldige og forskjelligartede philosophiske
Forsøg, der falde i dette Tidsrum, er det fælleds Særkjende,
ved Siden af Oppositionen mod den tidligere Philosophie, denne

 Afhængighed af Oldtiden, trods Bestræbelsen for og den
efterhaanden fremtrædende Fordring paa Uafhængighed.
\end{quote}
\dcite{Brøchner, Bidrag til Opfattelsen af Philosophiens historiske Udvikling, pp. 54–55  [brochner/philosophiens-udvikling:183]}
% brochner/philosophiens-udvikling:183  (Bidrag til Opfattelsen af Philosophiens historiske Udvikling; pp. 54–55)

\begin{quote}
See vi altsaa hen til de Tiders Eiendommelighed, som
dette Tidsrum forbinder, saa har Scholasticismens, den
«christelige Philosophies», Tid viist sig som ufri, splidagtig og
usammenhængende i sit Tankeliv. Scholasticismen optog sine

 Tankeformer fra en fremmed Livsanskuelse, optog dem der-
for kun som Former, i hvilke den bevægede sig ved en
Forstandsreflexion, der mere og mere fjernede sig fra Vir-
keligheden, og idet den kun kunde anvende dem paa det
over Tanken som urokkelig Autoritet staaende christelige
Dogme som ydre Former, i hvilke det ikke forvandledes til
et Tankeindhold, eller \tlg{tilegnedes} som begrebet af Bevidstheden,
saa blev dens Charakteer Ufrihed, Virkelighedsløshed og
Formalisme. Den opløstes ved Modsigelsen mellem dens
Elementer, ved den tænkende og religiøse Aands Utilfredsstil-
lethed i den, ved Aandens i mangfoldige Phænomener frem-
trædende, gjennem Oldtidens Gjenoplivelse styrkede Stræben
henimod Naturen og mod den frie Selvbevidstheds ideale Virke-
lighed. Trangen til Frihed og Trangen til Virkelighed,
Tankelivets dobbelte Betingelse, førte Bevidstheden bort fra
Kirkens Autoritet og fra det naturløse transscendente Aandelige,
og idet den førte den bort derfra, førte den den hen til Selv-
bevidstheden som til Sandhedens Kilde, og til Naturen som
Tankens Indhold, som det i den objective Virkeligheds Form
existerende Fornuftige og Sande, som Grundlaget for en concret
Udfoldelse af det menneskelige Liv. Ved denne forandrede
Retning vandt Aanden, som det allerede tidligere blev viist,
Betingelsen tilbage for Udviklingen af en virkelig Philosophie,
og denne Retnings methodiske, bevidste Gjennemførelse
\end{quote}
\dcite{Brøchner, Bidrag til Opfattelsen af Philosophiens historiske Udvikling, pp. 55–56  [brochner/philosophiens-udvikling:185]}
% brochner/philosophiens-udvikling:185  (Bidrag til Opfattelsen af Philosophiens historiske Udvikling; pp. 55–56)

\begin{quote}
Idet Giordano Bruno \tlg{tilegner} sig den Opfattelse af
Gud som Modsætningernes Enhed, der ved Periodens
Begyndelse var gjort gjældende af Nicolaus Cusanus, og idet
han nærmere bestemmer Gud som Enheden af de fire aristo-
\end{quote}
\dcite{Brøchner, Bidrag til Opfattelsen af Philosophiens historiske Udvikling, p. 60+  [brochner/philosophiens-udvikling:201]}
% brochner/philosophiens-udvikling:201  (Bidrag til Opfattelsen af Philosophiens historiske Udvikling; p. 60+)

\begin{quote}
Idealismens og Realismens Indvirkning paa hin-
anden i den nyere Philosophie. — Den gjensidige Indvirkning af den nyere Philosophies to
Hovedretninger see vi under Synspunktet af en gjensidig,
ubevidst Paavirken af deres dualistisk sondrede Principer.
Principernes gjensidige Bestemmen viser sig i negativ Form
idet hvert oprindelig hævder sig sine Prædicater exclusivt, og
gjennem denne exclusive \tlg{Tilegnen} af dem indirecte bestemmer
Modsætningens. Men fra denne Bestemmen, der falder sammen
med den i det Foregaaende paaviste Virkning af Dualismen,
see vi i denne Sammenhæng bort, og betragte her den ubevidste
gjensidige Paavirkning, som de to Principer udøve paa hinan-
den under den nærmere Bestemmelse af selve Principet og
under Udfoldelsen af det af Principbestemmelsen Afhængige, og
som fremtræder deri, at de her overføre deres Bestem-
melser paa Modsætningen, hvor de vise sig som In-
consequentser. Denne ubevidste Overføren af Bestemmelser
paa Modsætningen opfatte vi som nødvendig paa Grund af
Ensidigheden i det enkelte Princip, der i denne Ensidig-
hed skal udtrykke den omfattende Totalitet, og derfor
ligesom kræver en Fuldstændiggjørelse. Og gjennem den ud-
vortes, inconsequente Optagen af Modsætningens Bestemmel-
ser see vi den sande concrete Enhed antydet, som Frem-
tidsudviklingen skal virkeliggjøre.
\end{quote}
\dcite{Brøchner, Bidrag til Opfattelsen af Philosophiens historiske Udvikling, p. 170+  [brochner/philosophiens-udvikling:559]}
% brochner/philosophiens-udvikling:559  (Bidrag til Opfattelsen af Philosophiens historiske Udvikling; p. 170+)

\begin{quote}
Absolute, Guddommelige; kun paa dette passer Definitionen
fuldstændigt, paa det endelige Tænkende og det endelige Le-
gemlige kun med den Modification, at det til sin Existents
ikke behøver noget Andet end Gud. Der kommer saaledes,
naar Substantsbegrebet alligevel bibeholdes som Bestemmelse
for det Endelige, en Tvetydighed ind deri. Der tales om
den uendelige Substants og om de endelige Substantser,
og om disse atter snart saaledes, at det Aandelige og det Le-
gemlige som Heelhed betegnes ved dette Navn, snart saa-
ledes, at det individuelle Aandelige og Legemlige \tlg{tilegner}
sig det.
\end{quote}
\dcite{Brøchner, Bidrag til Opfattelsen af Philosophiens historiske Udvikling, p. 173+  [brochner/philosophiens-udvikling:574]}
% brochner/philosophiens-udvikling:574  (Bidrag til Opfattelsen af Philosophiens historiske Udvikling; p. 173+)

\begin{quote}
kjenden af det Objective, til Gjenstand, faaer Charakteren
af kritisk Idealisme, viser Paavirkningen sig paany. Vi
kunne see bort fra den negative Form af den, der er iden-
tisk med Dualismens Indvirkning, og som, naar vi tage
Stadiet som Heelhed, viser sig deri, at i det afsluttende Sy-
stem den Enhed, der naaes, trods Bestræbelsen for at aner-
kjende Naturens Væsentlighed, dog bliver en ensidig Enhed,
idet Aandsbegrebet, der skulde omfatte Naturens Væsen i sig,
ifølge den dualistiske Modsætning, ikke kan \tlg{tilegne} sig det
Naturliges eiendommelige Bestemmelser, men maa skyde dem
fra sig som usande. Men vi fremhæve en positiv Paavirk-
ning gjennem Overførelser fra Modsætningen, saavel i Be-
stemmelsen af Principet, som i det af Principet Afledede.
Hvad Bestemmelsen af Principet angaaer, saa forbigaae vi
Kant fordi den Maade, paa hvilken han unddrager det Ab-
solute fra Erkjendelsen, idethøieste tillader at søge en Paa-
virkning gjennem svage Antydninger. Hos Fichte finde vi
Paavirkningen, i hans ældre Philosophie, idet en uforklaret
Naturbestemmelse, en Naturskranke, trænger ind i Jegets
Activitet, og i hans Philosophies sildigere Form idet Væ-
rens Begreb afløser Jegets som det høieste. Hos Schel-
ling finde vi den i Identitetsphilosophiens Grundbegreb, der,
trods Fordringen af Subjectivitetens Overmagt, dog gaaer til-
bage til den spinozistiske uendelige Væren. Hos Hegel
endelig finde vi den i den tvetydige Sammenblanding af Sub-
jectiviteten og Objectiviteten, der bringer en Værens Reflex
ind i den logiske Idee og ligesom indsmugler Naturen i et
Moment af dens Proces: i Processens Sammenslutning i Umid-
delbarhed. — I det af Principet Afledede viser Paavirk-
ningen sig hos Kant naar Naturcausaliteten bliver Syns-
punktet for Opfattelsen af Aandsphænomenerne, idet det af
Aanden, der bliver Gjenstand for Erkjendelse, gaaer ind
under denne Causalitet, og den viser sig — hvad der umid-
delbart hænger sammen hermed — i den Bestemmelse af
\end{quote}
\dcite{Brøchner, Bidrag til Opfattelsen af Philosophiens historiske Udvikling, pp. 178–179  [brochner/philosophiens-udvikling:591]}
% brochner/philosophiens-udvikling:591  (Bidrag til Opfattelsen af Philosophiens historiske Udvikling; pp. 178–179)

\begin{quote}
Men naar nu en historisk Fremstilling følger denne Me-
thode; naar den saaledes samler Alt under een Tanke; naar
den lader en heel historisk Udfoldelse vise sig som denne
ene Tankes Historie og lader Forholdet til den bestemme
hvad der er væsentligt og hvad der er uvæsentligt; naar den
forklarer Alt ud fra denne Tankes Raaden, og, idet den be-
standig lader det Universelles Magt blive synlig, hyppigt lige-
som underlægger Tænkerne en anden Tanke end den, de be-
vidst forholde sig til, og lader en dobbelt Magt virke i deres
Tænken, — saa udsætter den sig for den Fare, at blive
vilkaarlig, eller idetmindste for at paadrage sig Skinnet
af en Vilkaarlighed, der gjør en subjectiv Opfattelse gjældende
paa den objectives Bekostning, og ikke respecterer den histo-
riske Kjendsgjerning i dens Givethed og Concretion. Opgaven
bliver da, baade i Virkeligheden at undgaae en saadan Vil-
kaarlighed og at frigjøre sig for Skinnet af den. Forat Vil-
kaarligheden kan undgaaes, fordres der foruden de Aandsevner,
der overhovedet ere Betingelsen for en klar \tlg{Tilegnelse} af de
historiske Kjendsgjerninger og en dybere Indtrængen i deres
Sammenhæng, et Studium af de philosophiske Systemer, der
paa een Gang er omfattende, universelt nok til at vise Tan-
kens store Bevægelser, ligesom Acterne i Tankedramaet, og
med det Samme fører saaledes ind i de enkelte store Tænkeres
Eiendommelighed, giver en saadan Fortrolighed med dem, at
man faaer Gehør for deres Tankebevægelses Rhythmus og
derved ledes til den inderlige Forstaaelse. Til dette Øiemed
\end{quote}
\dcite{Brøchner, Bidrag til Opfattelsen af Philosophiens historiske Udvikling, p. 182+  [brochner/philosophiens-udvikling:624]}
% brochner/philosophiens-udvikling:624  (Bidrag til Opfattelsen af Philosophiens historiske Udvikling; p. 182+)

\begin{quote}
gede Tankens blivende Enhed med den kunstneriske An-
skuelse og derigjennem den Begrebsdistinctionens Art, der
kunde betegnes som en Anskuelsesdistinction. I de
nævnte Former af den nyere Philosophie er det Dualismen,
der, idet den sætter Spaltningen mellem Tanken og Tankens
Andet: det Reale, lader dette, som Gjenstand for Tanken, ikke
blive gjennemtrængt og \tlg{tilegnet}, men kun ligesom udenfra
formet af denne, og derfor overføre sin Udvorteshedsson-
dring paa Realbegreberne, og gjennem dem paa Begreberne
overhovedet.
\end{quote}
\dcite{Brøchner, Bidrag til Opfattelsen af Philosophiens historiske Udvikling, p. 230  [brochner/philosophiens-udvikling:773]}
% brochner/philosophiens-udvikling:773  (Bidrag til Opfattelsen af Philosophiens historiske Udvikling; p. 230)

\begin{quote}
Opfattelse af Naturens Væsenløshed som skabt og som Skue-
plads for Underet; i Modsætning til Phantasiens og den i
Phantasien endende metaphysisk-abstracte Naturopfat-
telse, der skjult og indirecte statuerer Væsenløsheden af det
Naturlige som Naturligt, hævder Realismen, gjennem Fordrin-
gen paa en physisk Naturopfattelse, paa en Opfattelse af
Naturen ved den selv, gjennem Accentueringen af de virkende
Aarsager, Naturens Væsentlighed. Idet den opfatter Na-
turen som det Lovbegrundede og i sine egne, uforan-
derlige Love Hvilende, som det, der er Erkjendelsens og
Sandhedens Kilde, og som kun \tlg{tilegnes} ved Erfaringen
\end{quote}
\dcite{Brøchner, Bidrag til Opfattelsen af Philosophiens historiske Udvikling, p. 233  [brochner/philosophiens-udvikling:787]}
% brochner/philosophiens-udvikling:787  (Bidrag til Opfattelsen af Philosophiens historiske Udvikling; p. 233)

\begin{quote}
tikens Udvikling af sit concrete Indhold vil da den fremskridende,
men ikke fuldbyrdede \tlg{Tilegnelse} af det Reale træde

 frem i den Form, at det endnu ikke \tlg{Tilegnede} saaledes i Ud-
viklingen knyttes til det \tlg{Tilegnede}, at de Punkter hæves frem,
i hvilke Idealiteten ligesom lyser igjennem og som derved
pege hen paa den mere omfattende \tlg{Tilegnelse}. Den philoso-
phiske Fremstilling maa da omgive sit formede Indhold med
et endnu ikke i dens Form omdannet Stof, men med et saa-
dant, der røber Muligheden af en \tlg{Tilegnelse}, idet det allerede
ligesom er gjennemskinnet af enkelte Straaler af den forkla-
rede Tanke.
\end{quote}
\dcite{Brøchner, Bidrag til Opfattelsen af Philosophiens historiske Udvikling, pp. 274–275  [brochner/philosophiens-udvikling:927]}
% brochner/philosophiens-udvikling:927  (Bidrag til Opfattelsen af Philosophiens historiske Udvikling; pp. 274–275)

\begin{quote}
bliver den adæquate Realisation af Videnskabens Begreb,
maa den faae Principatet i Forholdet til de særskilte Viden-
skaber, og blive disses Regulator. Den bliver dette ved sin
Bevidsthed om de relative Principers Gyldighed og om Videns
Grundforhold; men idet den controllerer de særskilte Viden-
skabers Begrebsbestemmelser og deres Anvendelse af deres
Methode, paatvinger den dem ikke en for dem fremmed Er-
kjendelsesart, men lader netop deres i sin Eiendommelighed
bevarede Erkjendelsesart, idet den sikkres mod en Overskriden
af sine Grændser, føre til Resultater, som den philosophiske
Erkjendelse \tlg{tilegner} sig og assimilerer gjennem en Methode,
der ikke er fremmed for hine Videnskabers, men som netop
har \tlg{tilegnet} sig denne som Moment. I dette Forhold viser
Philosophien sig i sin indre Enhed med de enkelte Videnska-
ber, og netop idet de udvikle sig i deres hele Eiendommelig-
hed, blive de sig deres fælleds Opgave bevidste og arbeide
bevidst hen mod en Videnskabens Enhed, der realiseres
gjennem den philosophiske Erkjendelse. De reale Videnskabers,
af Philosophien besjælede, totale Organisation som en
saadan Videnskabernes Enhed, i hvilken de enkelte Led have
deres relative Selvstændighed, maa blive Maalet for Philo-
sophiens og de særskilte Videnskabers historiske Udvikling.
\end{quote}
\dcite{Brøchner, Bidrag til Opfattelsen af Philosophiens historiske Udvikling, p. 276  [brochner/philosophiens-udvikling:929]}
% brochner/philosophiens-udvikling:929  (Bidrag til Opfattelsen af Philosophiens historiske Udvikling; p. 276)

\subsection{\textit{Den græske Philosophies Historie i Grundrids} (1873–1874)}

\begin{quote}
bestemte idealt-reale Princip. —

 To Hovedformer af Philosophien foreligge saaledes historisk
for os som heelt afsluttede, og ved det Grundproblem, de have
stillet, men ikke løst, vise de hen paa en tredie Hovedform,
hvis Udvikling, der ved selve den abstracte Formel antydes
som betinget af, at den philosophiske Tanke \tlg{tilegner} og assi-
milerer sig de Videnskaber, der forholde sig umiddelbart til
det Reale, tilhører Fremtiden, den hele Fremtid.
\end{quote}
\dcite{Brøchner, Den græske Philosophies Historie i Grundrids, pp. 9–10  [brochner/philosophiens-historie-1:49]}
% brochner/philosophiens-historie-1:49  (Den græske Philosophies Historie i Grundrids; pp. 9–10)

\begin{quote}
Tal, Væren ( ), begge opfattede som det Stof, hvoraf Tin-
gene bestaae, som den legemlige Tilværelses Substants, og idet
Problemet om det Mangfoldiges Tilbliven af Principet, der
nødvendigen knytter sig til Problemet om Bestemmelsen af
Principet, træder frem og gjør sig gjældende som det væsent-
lige, bliver Vordens Begreb det, der medens det uforklaret
\tlg{tilegner} sig det Sandseliges Omraade, luttrer Begrebet Væren,
der nu søger et andet Omraade, til Ideen. Og saasnart Ideens
Begreb er fundet: det Begreb, til hvilket selve Stoffet viste
hen som til sin Fuldendelse, bliver Ideen sat som den sande,
væsentlige Virkelighed, og Tankens Opgave bliver: at søge et
Udtryk for Enheden af de to Momenter, i hvilket denne Forud-
sætning er fastholdt. Men idet denne Opgave kommer frem,
saa forudsætter den, at Tanken bevidst har sondret de to Mo-
menter, der for Anskuelsen ere forenede i Kunstværkets Enhed,
og denne Enheds Art umuliggjør, som det ovenfor blev berørt,
en ny Forening ved den græske Tanke, for hvilken Harmoniens
Umiddelbarhed var Enhedens eneste Form. Skjønhedens Umid-
delbarhed taaler ikke Tankens sondrende Magt, derfor bringer
Tanken Splid ind i det, hvorpaa den hviler, og gjennem denne
Splid maa denne Tankeform gaae tilgrunde. Momenterne, der
i deres Sondring begge tiltvinge sig en Bestemmelse af Væren,
staae dualistisk imod hinanden, og den græske Philosophie
maa ende i den frugtesløse Bestræbelse for, efterat Bevidst-
heden om Momenternes Væsensforskjel er klaret, atter at finde
en Forsoning.
\end{quote}
\dcite{Brøchner, Den græske Philosophies Historie i Grundrids, p. 20  [brochner/philosophiens-historie-1:89]}
% brochner/philosophiens-historie-1:89  (Den græske Philosophies Historie i Grundrids; p. 20)

\begin{quote}
( ), og kommer til at staae Herskerne nærmere end i
Anthropologien stod, skal leve alene for
Statens Forsvar indadtil og udadtil. Den tredie Stand skal
leve alene for Erhverv, for Tilveiebringelsen af alle Livsfor-
nødenheder, i stræng og tvungen Underordning under de andre.
I den uindskrænket herskende, Ideens Herredømme repræsen-
terende Stand har Statsideen en særskilt Existents, og
det er derfor consequent, at den Enkelte kun kan være Med-
lem af denne Stand ved at opgive de individuelle Interesser
for det Almene. For den herskende Stand ophører Privat-
eiendom og Familieliv, og den samme Lov gjælder for dens
«Hjælpere»: Krigerne. Statsmyndigheden ordner Kjønsforbin-
delserne, fastsætter Antallet af de Børn, der maae leve, og ud-
vælger det blandt de Fødte, \tlg{tilegner} sig fuldstændig Børnene
ligefra Fødselen, bestemmer, efter de Anlæg, de vise, deres
Livsstilling, og leder de til Statens Vogtere Udvalgtes strængt
ordnede Opdragelse og Uddannelse. (De to til tre første Leve-
aar ere for Børnene bestemte til Legemspleie; fra det 3die til
det 6te uddannes de ved Fortælling af Myther; fra det 7de
til det 10de øves de i Gymnastik; fra det 10de til det 13de
i Læsen og Skriven; fra det 14de til det 16de i Kjendskab til
\end{quote}
\dcite{Brøchner, Den græske Philosophies Historie i Grundrids, p. 130+  [brochner/philosophiens-historie-1:434]}
% brochner/philosophiens-historie-1:434  (Den græske Philosophies Historie i Grundrids; p. 130+)

\begin{quote}
I Bestemmelsen af Aristoteles's Forhold til Platon ligger
Bestemmelsen af hans Forhold til den græske Philosophie over-
hovedet. Aristoteles er den græske Philosophies
ideelle Afslutning. Hvad den græske Tænkning stræbte
efter som sit høieste Maal: i Begrebets Form at udtrykke
den Tilværelsens Enhed, der umiddelbart er tilstede for
Anskuelsen i Verdenskunstværket, er hos Aristoteles forsøgt
virkeliggjort med en større Tankens Energie end af nogen
anden græsk Tænker, og ved Bestemmelsen af Formen som
iboende og af Stoffet som det efter Muligheden Værende ere
Principerne satte i et saadant Forhold til hinanden, at hiin
Stræben har et virkeligt Grundlag. Og medens Aristoteles
energisk tilstræber den principielle Enhed, omfatter og \tlg{tilegner}
hans Philosophie sig — netop som Følge af hans Opfattelse
af Grundmodsætningernes Forhold — i saa rigt et Maal som
noget græsk System den reale Virkelighed, medens den tillige
\end{quote}
\dcite{Brøchner, Den græske Philosophies Historie i Grundrids, pp. 138–139  [brochner/philosophiens-historie-1:465]}
% brochner/philosophiens-historie-1:465  (Den græske Philosophies Historie i Grundrids; pp. 138–139)

\begin{quote}
bereder Aristoteles gjennem en Kritik af sine Forgængeres
Bestemmelser af det i Sandhed Værende, og navnlig gjennem
Kritiken af den platoniske Ideelære, der først har angivet
Formprincipet, om den end ikke har bestemt det paa den
rette Maade. Han \tlg{tilegner} sig den platoniske Bestemmelse af
\end{quote}
\dcite{Brøchner, Den græske Philosophies Historie i Grundrids, p. 146+  [brochner/philosophiens-historie-1:505]}
% brochner/philosophiens-historie-1:505  (Den græske Philosophies Historie i Grundrids; p. 146+)

\begin{quote}
f. Physiken. — Vi have i det Foregaaende fremhævet — den store Forskjel, der er mellem Platon og Aristoteles i deres
Forhold til det Physiske. Platon iler bort fra Naturvirkelig-
heden for at fortabe sig i de evige Ideer, Philosophiens eneste
adæquate Gjenstand; Aristoteles drages bestandig hen mod den
Virkelighed, gjennem hvilken Formen aabenbarer sig og som
netop derfor overalt har noget Guddommeligt og Beundrings-
værdigt og gjemmer en uudtømmelig Rigdom for Erkjendelsen
i sig, og han nægter ikke denne naturvidenskabelige Erkjen-
delse Philosophiens Navn, selv om den kun fører til det Sand-
synlige, Mulige og relativt Almindelige, og selv om den maa
standse før den fuldstændig har bemægtiget sig Gjenstanden.
Aristoteles er paa een Gang Philosoph og Naturforsker: han
er Mester i Iagttagelsen af Kjendsgjerningerne, ved sin Stræben
efter at forklare dem af deres eiendommelige physiske Grunde
staaer han høiest i sin Periode, og han har ikke blot \tlg{tilegnet}
sig al sin Tidsalders Kundskab til alle Naturens Gebeter, men
har selv i mange Retninger givet den større Omfang og
Dybde.
\end{quote}
\dcite{Brøchner, Den græske Philosophies Historie i Grundrids, p. 165+  [brochner/philosophiens-historie-1:610]}
% brochner/philosophiens-historie-1:610  (Den græske Philosophies Historie i Grundrids; p. 165+)

\begin{quote}
6. Eklekticismen. — a. Eklekticismens almindelige Characteer. — Indtil
henimod Begyndelsen af det første forchristelige Aar-
hundrede gaae de forskjellige philosophiske Retninger ublan-
dede ved Siden af hverandre, saa begynder der en Sammen-

 smeltning gjennem Mæglings- og Foreningsforsøg, af hvilke
kun Epikuræismen holder sig uberørt. Den ved de forskjellige
Skolers gjensidige Kritik klarede Bevidsthed om de enkelte
Systemers Mangler fører, idet den philosophiske Productivitet
er afmattet, til en eklektisk Vælgen af det i de forskjellige
Systemer tilsyneladende Sikkre, som saa den enkelte Skole
\tlg{tilegner} sig idet den ved en Tillæmpen af sine egne Læresæt-
ninger tilslører Differentsen fra det Fremmede. Ogsaa den
skeptiske Retning maatte tilsidst føre til Eklekticisme, idet
der, efter at den sikkre Viden var opgiven, af Hensyn til
Handlingen gaves det Sandsynlige en Betydning, der, hvor den
praktiske Trang endnu bestemtere traadte i Forgrunden, atter
lod det indtage det Sandes Sted, og idet dette Sandsynlige af
den skeptiske Bevidsthed ikke kunde søges i et enkelt System,
men, efter den subjective Trangs Anvisning, maatte søges i de
forskjellige for Skeptikeren ligeberettigede Systemer. Navnlig
fra Karneades's Efterfølgere i Akademiet udgik en eklektisk
Retning, der støttede sig paa Nødvendigheden af bestemte
Anskuelser som Grundlag for det praktiske Liv. En stor Ind-
flydelse paa Udbredelsen af Eklekticismen fik den græske
\end{quote}
\dcite{Brøchner, Den græske Philosophies Historie i Grundrids, pp. 229–230  [brochner/philosophiens-historie-1:889]}
% brochner/philosophiens-historie-1:889  (Den græske Philosophies Historie i Grundrids; pp. 229–230)

\begin{quote}
Naar vi betegnede Kirkefædrenes Tid som Dogmeudpræg-
ningens, kunne vi betegne Scholasticismens som Dogmebear-
beidelsens, den formelle \tlg{Dogmetilegnelses} Tid. I hint Tids-
afsnit havde den christelige Aand udfoldet Dogmet til en væ-
sentlig afsluttet Heelhed af objective Sandheder; i denne har
den Dogmeudfoldelsen som et tilbagelagt Stadium, og Opgaven
bliver nu den: at \tlg{tilegne} Bevidstheden den objective Sandhed.
Som Middel hertil benyttes Philosophien, der ogsaa i dette
Tidsrum optages fra Oldtiden, navnlig Platons og Aristoteles's
Philosophie. Dens Former anvendes paa den christelige Troes
Indhold, der i Kirkelæren forelaae som et Object med urokke-
\end{quote}
\dcite{Brøchner, Den græske Philosophies Historie i Grundrids, p. 262+  [brochner/philosophiens-historie-1:998]}
% brochner/philosophiens-historie-1:998  (Den græske Philosophies Historie i Grundrids; p. 262+)

\begin{quote}
længe, ved de optagne Former, udspindes Consequentser af de
ubegribelige Forudsætninger, bliver dette dog ikke noget Tanke-
indhold. Idet Tankeformerne blive ydre og fremmede Former
for Troesgjenstanden, kan den ikke gjennem dem \tlg{tilegnes} af
Selvbevidstheden; der kan i det Høieste ved dem tilveiebringes
\end{quote}
\dcite{Brøchner, Den græske Philosophies Historie i Grundrids, p. 272+  [brochner/philosophiens-historie-1:1042]}
% brochner/philosophiens-historie-1:1042  (Den græske Philosophies Historie i Grundrids; p. 272+)

\subsection{\textit{Den nyere Philosophies Historie i Grundrids} (1873–1874)}

\begin{quote}
kunne udfolde et selvstændigt Tankeindhold. Idet den fri-
gjorde sig fra Kirkens Herredømme, og, forkyndende en ny
Verdensanskuelse, vendte sig kæmpende mod Scholasticismen,
maatte den derfor søge det for denne Kamp nødvendige Ind-
hold, der maatte være et saadant, i hvilket den befriede
menneskelige Tanke kunde gjenfinde sig og gjenkjende Virke-
ligheden, hvor det forelaae som givet, og den fandt det da
hos Oldtidens Philosophie, med hvilken den følte sig i Slægt-
skab. Den ombyttede da Scholasticismens tvungne Under-
kastelse under en fremmed Autoritet med en selvvalgt Under-
kastelse under en beslægtet, i hvilken den nye Tids Tanke,
der nu kunde have Virkelighed og Nærværelse i sit Indhold,
kunde vinde Fylde og Selvstændighed gjennem \tlg{Tilegnelsen} og
Udformningen af dette Virkelighedsindhold, i hvilket den først

 skulde leve sig ind. I alle de mangfoldige og forskjelligartede
philosophiske Forsøg, der falde i dette Tidsrum, er det fælleds
Særkjende, ved Siden af Oppositionen mod Middelalderens Phi-
losophie, denne Afhængighed af Oldtiden, trods Bestræbelsen
for og den efterhaanden fremtrædende Fordring paa Uaf-
hængighed.
\end{quote}
\dcite{Brøchner, Den nyere Philosophies Historie i Grundrids, pp. 2–3  [brochner/philosophiens-historie-2:27]}
% brochner/philosophiens-historie-2:27  (Den nyere Philosophies Historie i Grundrids; pp. 2–3)

\begin{quote}
Hos Giordano Bruno (født c. 1548, brændt i Rom — 1600) finder Tidsalderens Tænkning sit klareste og meest cha-
rakteristiske Udtryk. Idet han \tlg{tilegner} sig den Opfattelse af
Gud som Modsætningernes Enhed, der ved Periodens Begyn-
delse var gjort gjældende af Nicolaus Cusanus, og idet han
\end{quote}
\dcite{Brøchner, Den nyere Philosophies Historie i Grundrids, p. 13+  [brochner/philosophiens-historie-2:72]}
% brochner/philosophiens-historie-2:72  (Den nyere Philosophies Historie i Grundrids; p. 13+)

\begin{quote}
ogsaa \tlg{tilegner} sig Substantsens Navn og bringer en Tvetydighed — ind i dens Begreb, og der bliver en ny Dualisme mellem Gud
og den skabte Tilværelse.
\end{quote}
\dcite{Brøchner, Den nyere Philosophies Historie i Grundrids, p. 88  [brochner/philosophiens-historie-2:416]}
% brochner/philosophiens-historie-2:416  (Den nyere Philosophies Historie i Grundrids; p. 88)

\begin{quote}
sin afhængige Væren. Og idet den, efterat have foretaget
denne Udelukkelse, abstract sondrer mellem den umiddelbare
Vished og den middelbare Viden, miskjender den den psycho-
logiske Genesis af de Forestillinger, der forkynde sig med
umiddelbar Vished; den miskjender Muligheden af det fra først
af Middelbares Omsætten i et Umiddelbart gjennem \tlg{Tilegnel}-
sen; den glemmer, at den umiddelbare Vished om det Over-
sandselige er medieret ved en Opløftelse over Sandseligheden,
og at den concrete Enhed, der sammenslutter modsatte Be-
stemmelser i sig, indeslutter dem saaledes, at hver af dem er
betinget, medieret ved den anden, hvilket er det Samme som,
at Enheden gjennem denne Mægling slutter sig sammen med
sig selv i en høiere Umiddelbarhed, paa een Gang altsaa en
Mediation og umiddelbar Relation til sig, og i selve Media-
tionen ophæver Mediationen. Saaledes erkjendes det Absolute
kun som Aand idet det begribes som i sig medierende sig
med sig selv. Den abstracte Sondring mellem den umiddel-
bare Vished og den middelbare Viden fører uundgaaelig til
den Consequents, at Troens Indhold enten kun bliver det Ab-
stracte, Bestemmelsesløse, da enhver concret Bestemmelse
\end{quote}
\dcite{Brøchner, Den nyere Philosophies Historie i Grundrids, p. 190  [brochner/philosophiens-historie-2:1115]}
% brochner/philosophiens-historie-2:1115  (Den nyere Philosophies Historie i Grundrids; p. 190)

\begin{quote}
var han saa Professor i Würzburg, siden Medlem af Viden-
skabernes Akademie i München og Akademiets Generalsecretær;
fra 1820—26 holdt han Forelæsninger i Erlangen, blev, ved
Münchener-Universitetets Oprettelse, atter kaldet dertil og vir-
kede der indtil 1841, da Romantikeren, Kong Friedrich Wil-
helm IV af Preussen drog ham til Berlin for at skulle be-
kæmpe Hegelianismen og begrunde en christelig Philosophie.
I Berlin foredrog han i nogle Aar sin nye Philosophie, der
successivt havde udviklet sig af hans Identitetsphilosophie,
under Paavirkning af de mystiske og theosophiske Systemer,
med hvilke Schelling efterhaanden blev bekjendt, og af hvilke
han paa sin Viis \tlg{tilegnede} sig Elementer. Sine sidste Aar levede
han i fornem Tilbagetrukkenhed, beskæftiget med Udarbei-
delsen af det længe forjættede Aabenbaringssystem. Han
døde 1854.
\end{quote}
\dcite{Brøchner, Den nyere Philosophies Historie i Grundrids, pp. 225–226  [brochner/philosophiens-historie-2:1348]}
% brochner/philosophiens-historie-2:1348  (Den nyere Philosophies Historie i Grundrids; pp. 225–226)

\begin{quote}
værende Aand, der veed sig i Mennesket, hvis Viden — om Gud bliver til dets Viden af sig i Gud, adskilles for
Forestillingen i udenfor hinanden faldende Processer, der hver
for sig fremstille det absolute Indhold. Den guddommelige
Treenigheds evige Proces adskilles fra den Proces, hvis
Momenter ere Skabelsen af Naturen og den endelige Aand,
det Skabtes Brud med det Guddommelige, og Forsoningen ved
det historiske Gudmenneske, der \tlg{tilegnes} gjennem Troen og
Negationen af Endeligheden i den sande Cultus. Idet Christi
Død som hans sandselige Existentses Ophør gjør en aandelig
Opfattelse af ham mulig, og det i ham Anskuede faaer en
aandelig Existents i Menighedens Bevidsthed, forvandles for
Subjectet, der opgiver sin Endelighed, den ydre Forsoning til
en indre, og Tanken erkjender i det, som Forestillingen ad-
skilte, den absolute Aands uadskillelige evige ene Selvaabenbaring.
\end{quote}
\dcite{Brøchner, Den nyere Philosophies Historie i Grundrids, p. 268  [brochner/philosophiens-historie-2:1609]}
% brochner/philosophiens-historie-2:1609  (Den nyere Philosophies Historie i Grundrids; p. 268)

\begin{quote}
denne Processes Trin i sin relative Abstraction giver et prin-
cipielt sandt, ligesom ved sig selv mod sin Integration hen-
stræbende, Udtryk for det Virkelige, der paa hvert Punkt
aabenbarer Principet og Totaliteten. Under Dialektikens Ud-
vikling af sit concrete Indhold maa den fremskridende, men
ikke fuldbyrdede \tlg{Tilegnelse} af det Reale træde frem i den
Form, at det endnu ikke \tlg{Tilegnede} saaledes i Udviklingen
knyttes til det \tlg{Tilegnede}, at de Punkter hæves frem, i hvilke
Idealiteten ligesom lyser igjennem og som derved pege hen
paa den mere omfattende \tlg{Tilegnelse}. Den philosophiske Frem-
stilling vil da omgive sit formede Indhold med et endnu ikke
i Begrebet omsat Stof, men med et saadant, der røber
Muligheden af denne Omsætning, idet det allerede ligesom er
gjennemskinnet af enkelte Straaler af den forklarende Tanke.
\end{quote}
\dcite{Brøchner, Den nyere Philosophies Historie i Grundrids, p. 281  [brochner/philosophiens-historie-2:1672]}
% brochner/philosophiens-historie-2:1672  (Den nyere Philosophies Historie i Grundrids; p. 281)

\begin{quote}
Grundforhold (jfr. 1. Deel, S. 1 f.), men idet den skal control-
lere de særskilte Videnskabers Begrebsbestemmelser og deres
Anvendelser af deres Methode, kan den ikke ville paatvinge
dem en for dem fremmed Erkjendelsesart, men maa netop lade
deres i sin Eiendommelighed bevarede Erkjendelsesart, idet den
sikkres mod en Overskriden af sine Grændser, føre til Resul-
tater, som den philosophiske Erkjendelse \tlg{tilegner} sig og assi-
milerer gjennem en Methode, der ikke er fremmed for hine
Videnskabers, men som netop har \tlg{tilegnet} sig denne som Mo-
ment. I dette Forhold viser Philosophien sig i sin indre
Enhed med de særskilte Videnskaber, og netop idet disse ud-
vikle sig i deres hele Eiendommelighed, blive de sig deres
fælleds Opgave og deres fælleds Midler bevidste, og arbeide
\end{quote}
\dcite{Brøchner, Den nyere Philosophies Historie i Grundrids, p. 282+  [brochner/philosophiens-historie-2:1677]}
% brochner/philosophiens-historie-2:1677  (Den nyere Philosophies Historie i Grundrids; p. 282+)

\part*{IV.\quad Harald Høffding}
\addcontentsline{toc}{part}{IV.\quad Harald Høffding}

\subsection{\textit{Filosofien og Darwinismen} (1874)}

\begin{quote}
Der er altsaa god Grund til at anstille en Prøvelse af Darwinismens
logiske og ethiske Værd for at komme til en Afgjørelse af, hvilken af
disse stridende Domme der er den sande, og paa hvilken Maade vi skulle
stille os til denne mærkelige Anskuelse, som fra Videnskabsmandens
stille Kammer nu er trængt ud paa Gader og Stræder og anvendes som et
mægtigt Vaaben i vor Tids aandelige Kampe. Filosofien har, afseet fra
disse ydre Foranledninger, i Følge sin egen Opgave den Pligt at anstille
en saadan Prøvelse. Det er Filosofiens Opgave at blive sig den
videnskabelige Erkjendelses Væsen og Methode bevidst, at skærpe Indsigten
i de Stadier, en Anskuelse maa gjennemgaa, førend den kan indlemmes i de
videnskabelige Sandheders Rige, og at hævde de nødvendige Betingelser
for Erkjendelsens Sandhed og Gyldighed. Men dens Kald er ikke begrændset
hertil. Idet den undersøger det menneskelige Bevidsthedsliv, dets Love
og dets Elementer, og forfølger dets Udvikling gjennem Historien, søger
den tillige at indordne ethvert Erkjendelsesresultat som Led i Aandens
hele Organisme og maa derfor paa ethvert Punkt prøve, hvorledes det
vundne Resultat forholder sig som Bidrag til en almindelig
Verdensanskuelse. Intet Resultat er i Sandhed \tlg{tilegnet} af Menneskeaanden,
førend denne Prøvelse er anstillet. Og først naar dette er gjort, kan
Resultatet faae berettigede praktiske Følger, idet det tredie Led i den
filosofiske Undersøgelse bliver: hvilken praktisk og ethisk Betydning
den \tlg{tilegnede} Sandhed faaer.
\end{quote}
\dcite{Høffding, Filosofien og Darwinismen  [hoeffding/darwinismen:2]}
% hoeffding/darwinismen:2  (Filosofien og Darwinismen)

\begin{quote}
Det er saa langt fra, at Darwinismen fører til Materialisme, at den
netop maa siges at vise i den helt modsatte Retning. Den lægger ganske
vist en stærk Vægt paa de ydre Vilkaar; men disse Vilkaar virke kun
gjennem Organismens Lempelse efter dem. Og hvad man kalder Tilfælde, er
kun Exemple paa Betingelsernes, de indre og ydre Betingelsers, uendelig
forgrenede og indviklede Sammenhæng, som vi ikke kunne overse; i de
saakaldte tilfældig opstaaede Former have vi netop et Bevis paa Livets
Energi, der veed at \tlg{tilegne} sig endog de fineste Nuancer i de ydre
Vilkaar og anvende dem i sin Tjeneste.
\end{quote}
\dcite{Høffding, Filosofien og Darwinismen  [hoeffding/darwinismen:19]}
% hoeffding/darwinismen:19  (Filosofien og Darwinismen)

\begin{quote}
Jeg har hørt en begavet Dame sige, at dersom Darwin havde Ret, og vi
virkelig nedstammede fra Aberne, vilde Livet tabe alt Værd og Betydning
for hende. Det er maaske dette Punkt, som for de Fleste staaer som det
mest afskrækkende. Men den, som seer, at selv det Højeste og Ædleste dog
har sine bestemte Betingelser, at det ikke kan være uden dem, kun er ved
at \tlg{tilegne} sig dem, — den, som erkjender, at de tilsyneladende absolute
Forskjelle forsvinde i den uendelige Sammenhæng, vil se Sagen paa en
anden Maade. Han vil ogsaa i Menneskeaanden se en Form for den kæmpende
Ide og vil ikke se nogen Nødvendighed for, at det Aandelige ude- eller
ovenfra indblæses i Materien; det Aandelige staaer ikke som et fremmed
Element i Tilværelsen; det fremkommer paa sit bestemte Trin i
Udviklingen, arbejder sig frem af selve Materien og afgiver i sidste
Instans den eneste Forklaring af denne. Det Paafaldende i den
Fremstilling, at Menneskene skulde nedstamme fra Aberne, har sin simple
Grund deri, at man her umiddelbart stiller to Udviklingstrin sammen, der
ligge uendelig langt fra hinanden; man overspringer da netop det
Væsenlige — selve Udviklingen. Hvor stor en Kløft er der ikke mellem
de Menneskeædere, fra hvilke vi maaske nedstamme, og vor Tids
civiliserede Mennesker!
\end{quote}
\dcite{Høffding, Filosofien og Darwinismen  [hoeffding/darwinismen:24]}
% hoeffding/darwinismen:24  (Filosofien og Darwinismen)

\subsection{\textit{Om Realisme i Videnskab og Tro} (1882 / 1884)}

\begin{quote}
Realismen bliver derfor først en ejendommelig Verdensopfattelse, naar
den vover sig ud over det blot Givne eller det hidtil Givne og aabner
Udsigten til en total og gennemført Anvendelse af de naturlige Aarsagers
Princip. Det er dette, som er sket i de store naturvidenskabelige
Hypoteser, hvis Række aabnes af Kants og Laplaces Teori og slutter med
Darwins og Spencers. Disse Hypoteser ere voxede naturligt frem som
Konsekvenser af Videnskabens Udviklingsgang. I dem træder Realismen
tydeligt frem som almindeligt videnskabeligt Princip. Det er ikke altid
sagt, at Detailforskerne blive sig dette Princip bevidst. Arbejdets
Deling paa Videnskabens Omraade har ført det med sig, at den enkelte
Forsker med Lethed og tilsyneladende ogsaa med god Samvittighed kan
grave sig ned paa et enkelt Punkt uden at bekymre sig om, hvor vidt de
Principer og Metoder, han følger, have almindeligere Betydning og
indeholde Konsekvenser af indgribende Betydning for hele
Verdensopfattelsen. Mange Naturforskeres Stilling til Darwinismen finder
sin Forklaring herved. Den, som derimod fæster sit Blik paa
Videnskabernes almindelige Udviklingsgang og sér, hvorledes ethvert
Fremskridt i Forstaaelse er vundet ved Anvendelse af de naturlige
Aarsagers Princip, betragter dette Princip som det, der er kaldet til at
beherske hele vor Naturopfattelse. Det maa udfolde alle sine Konsekvenser;
vi maa \tlg{tilegne} os det; uden om det fører ingen Vej.
\end{quote}
\dcite{Høffding, Om Realisme i Videnskab og Tro, p. 6  [hoeffding/realisme:12]}
% hoeffding/realisme:12  (Om Realisme i Videnskab og Tro; p. 6)

\subsection{\textit{Forholdet mellem Tro og Viden i dets historiske Udvikling} (1885)}

\begin{quote}
II. — Tiden gik, Afslutningen kom ikke. Kirken trøstede sig med, at for Gud er
tusind Aar som én Dag, og én Dag som tusind Aar (2. Peters Brev 3. Kap),
og begyndte snart at falde til Ro i Verden, at indrette sig efter
Forholdene og at tage Stilling over for de forskellige Sider af de
kulturhistoriske Bestræbelser. Over for Videnskaben blev dens Stilling den,
at den ikke længer i Troen fandt al fornøden Erkendelse, men søgte at vise,
at Troens Lærdomme satte Kronen paa al menneskelig Erkendelse. Kirken
\tlg{tilegnede} sig den Videnskab, Grækerne havde udviklet: den græske Astronomi,
Medicin og Filosofi, og anbragte sin Teologi som Spiret paa Videnskabens
Bygning, — paa lignende Maade som den omdannede de græske Templer til
Kirker og f. Eks. indviede Partenon, Athenes Tempel, til den »hellige
Visdom«.
\end{quote}
\dcite{Høffding, Forholdet mellem Tro og Viden i dets historiske Udvikling, pp. 196–197  [hoeffding/forholdet-tro-viden:4]}
% hoeffding/forholdet-tro-viden:4  (Forholdet mellem Tro og Viden i dets historiske Udvikling; pp. 196–197)

\begin{quote}
For Protestantismen er det 18. Aarhundrede, hvad det 14. var for
Katolicismen. Men den Krise, som det bragte, var langt grundigere, bundede
langt dybere end Renaissancebevægelsen. Det er den selvstændige humane
Livsanskuelses Fødselstid. Forløberne vare naturligvis mange — men først
nu forstodes og anerkendtes de. Først nu fandt f. Eks. Spinoza, der i
halvandet Aarhundrede havde ligget »som en død Hund«, en Kreds, der kunde
tænke hans Tanker og \tlg{tilegne} sig dem. Først nu fik den teologiske
Livsanskuelse en jævnbyrdig Modstander. Jeg tænker her ikke paa noget
bestemt, afsluttet System, men paa en almindelig Tanke- og Livsretning, som
i de store Træk var fælles for Mænd som Lessing og Kant, Schiller og Goethe.
— Som man ser, fandt denne nye Livsanskuelse straks mere end én Dante.
\end{quote}
\dcite{Høffding, Forholdet mellem Tro og Viden i dets historiske Udvikling, p. 199+  [hoeffding/forholdet-tro-viden:15]}
% hoeffding/forholdet-tro-viden:15  (Forholdet mellem Tro og Viden i dets historiske Udvikling; p. 199+)

\subsection{\textit{Kontinuiteten i Kants filosofiske Udviklingsgang} (1893)}

\begin{quote}
Herved adskiller Kants Etik sig fra den tidligere rationalistiske Etik,
saaledes som den udformedes af «Intellektualisterne» i det 17.
Aarhundrede og af Price. Naar disse grundede det Etiske paa Fornuft,
mente de dermed Tingenes objektive, fornuftige Sammenhæng, som Mennesket
opfatter ved sin Fornuft og derefter bøjer sig for. Fornuften som
menneskelig Evne fik her, som Fr. Jodl træffende udtrykker
sig [note: Geschichte der Ethik in der neueren Philosophie.
I. München 1882. p. 221.], en blot receptiv Rolle at udføre, nemlig at
\tlg{tilegne} sig den Fornuft, der antoges given i Verdensforholdene, og
anvende den i sin Handlen. Kants praktiske Fornuft er derimod, ligesom
hans teoretiske Fornuft, Udtryk for Menneskets eget inderste Væsen, et
Indbegreb af Love for Menneskets Aandsvirksomhed. Den dybsindige Tanke i
Kants definitive Etik er, at Mennesket ved at blive sig den moralske Lov
bevidst, bliver sig bevidst som Borger i to forskellige Verdener, der
mødes i hans Person: «Den moralske Skullen er Subjektets egen Villen som
Led i en intelligibel Verden og tænkes kun som Skullen, for saa vidt
Subjektet tillige betragter sig som Led af
Sanseverdenen». [note: Grundlegung zur Metaphysik der Sitten.
3. Aufl. Riga 1792. p. 113.] I streng rationalistisk Form har Kant
her paa ny udviklet en Tanke, der i mere empirisk Form kom frem i
«Beobachtungen», naar der appelleredes til den menneskelige Naturs
Værdighed. Det er denne Værdighed, han mente at maatte søge en dybere
Begrundelse for end Psykologien kunde give. Ideen om Menneskets Værdighed
er det kontinuerlige Element i Kants Etik; hvad der varierer, er
Begrundelsen.
\end{quote}
\dcite{Høffding, Kontinuiteten i Kants filosofiske Udviklingsgang, pp. 45–46  [hoeffding/kants-udvikling:97]}
% hoeffding/kants-udvikling:97  (Kontinuiteten i Kants filosofiske Udviklingsgang; pp. 45–46)

\begin{quote}
Betydningen af Kants Syntesebegreb er for det Første den, at han derved
viser os, hvad det er at forstaa en Ting. Hans Erkendelseslære hviler paa
den simple Tanke, at det at forstaa Noget er at se det i saa nøje
Sammenhæng som muligt med alt Andet, vi kende. En forbindende
Aandsvirksomhed maa gøre sig gældende ved al Erkendelse; kun naar vi
kunne anvende den, have vi \tlg{tilegnet} os det givne Stof og føle den
ejendommelige Tilfredsstillelse, som følger med Afhjælpelsen af en Trang.
Hvad vi ikke forstaa, er det Isolerede og Sammenhængsløse. Naturligvis er
der mange Grader mellem de to Yderpunkter: forstaa og ikke forstaa,
Sammenhæng og Sammenhængsløshed, og, som vi have set (ovenfor § 14–15),
begik Kant den Fejl ikke at agte paa disse Gradsforskelle. Men den
Fortjeneste har han, først at have set Erkendelsesproblemet paa denne
simple Maade. Det er et Kolumbusæg, han har fundet og stillet paa
Spidsen. Mærkværdigt nok havde Hume fundet det samme Kolumbusæg, men lod
det trille fra sig. I sin «Treatise» (som Kant ikke kendte) kom Hume
egentligt til samme Grundbegreb, som Kant naaede til i «Kritik der reinen
Vernunft». Hume førte nemlig dér Erkendelsens Problem fra de specielle
Problemer (særligt Problemet om Geometriens reale Gyldighed og om
Aarsagsforholdet) tilbage til den Grundvanskelighed, som for ham —
formedelst den atomistiske Psykologi, han gik ud fra — laa i
Forbindelsen mellem Forestillinger i det Hele. Den store Gaade laa for
ham tilsidst i Enhedsprincipet (the uniting
principle). [note: Treatise of human nature. I, 3, 14 (ed.
Selby-Bigge. Oxford 1888. p. 169). Kfr. Appendix (ibid. p. 636).] 
Kant viser nu, at Forskellen mellem det Forstaaelige og det Gaadefulde
beror paa Forskellen mellem det, der kan sammenfattes i kontinuerlig
Sammenhæng, og det, hvor en saadan Funktion ikke er mulig. Selve den
forbindende Funktion kan man saa aabenbart kun finde gaadefuld, enten
naar man gaar ud fra, at det Sammenhængsløse skulde være det Forstaaelige,
eller naar man spørger om den større Sammenhæng, af hvilken vor
sammenfattende Funktion igen er betinget. Den første Vej fører ud i det
Meningsløse; den anden Vej fører til at stille os det sidste Problem,
Erkendelsesteorien og Psykologien i det Hele kunne staa over for, og som
betegner deres Grænse: nemlig om selve Erkendelsens, eller i det Hele
Bevidsthedens Opstaaen i Tilværelsen. Tanken er sig sit eget sidste
Problem. [note: Smlgn. min Psykologi. 3. Udg. p. 403.] — Det
vilde sikkert have været af stor Betydning for Kants Udvikling, om han
havde kendt Humes Hovedværk paa et tidligt Stadium af sin Bane.
Sammenstødet mellem de to store Tænkere vilde da have angaaet et endnu
mere centralt Spørgsmaal, og Kant vilde maaske være bleven ledet til en
fuldkommen Gennemarbejdelse af det psykologiske Grundlag for sin
Erkendelsesteori, medens han paa den anden Side vilde have kunnet spare
sig et langt skolastisk Arbejde. Ægget vilde tidligere og lettere være
blevet sat paa Spidsen. Det var en for Filosofiens følgende Historie
skæbnesvanger Handling, Hume begik, da han, misfornøjet med den ringe
literære Lykke, hans Ungdomsværk gjorde og fortrædelig over de ubelejlige
Angreb, det var Genstand for fra teologisk Side, fornegtede dette sit
storartede Værk. [note: Eduard Grimm (Zur Geschichte des
Erkenntnissproblems. Leipzig 1890. p. 584) ser i denne Fornegtelse en
Indrømmelse fra Humes Side af det Ensidige i «Treatise». Men at Motivet
afgjort var det ovenfor fremførte, ses tydeligt dels af Humes
Selvbiografi, dels af det Brev, i hvilket Hume opfordrede sin Forlægger
til at lade den Erklæring, i hvilken han fornegtede «Treatise», følge med
«Inquiry». Letters of David Hume to William Strahan. Edited by
C. Birkbeck Hill. Oxford 1888. p. 288 f. — Det var ogsaa af den
sidst omtalte Grund, at Hume ikke lod «Dialogerne» udkomme i sin Levetid.
Ibid. p. 330.]
\end{quote}
\dcite{Høffding, Kontinuiteten i Kants filosofiske Udviklingsgang, pp. 56–57  [hoeffding/kants-udvikling:121]}
% hoeffding/kants-udvikling:121  (Kontinuiteten i Kants filosofiske Udviklingsgang; pp. 56–57)

\subsection{\textit{Religionsfilosofi} (1901)}

\begin{quote}
Den dybeste Begrundelse for de naturlige Aarsagers Princip finder jeg
i en Trang til Kontinuitet, som ligger i vor Bevidstheds Natur, og
som allerede lægger sig for Dagen i den almindelige Trang til at
finde Aarsager. Vor Bevidsthed stræber uvilkaarligt at sammenfatte og
sammenholde sine Elementer og at bevare Enheden med sig selv, ogsaa
naar den staar overfor et nyt Indhold. Det nye Indhold stræber den da
at forme, ændre og lægge til Rette, saaledes at det Nye viser sig at
være en ny Form for et Gammelt. Selvgenkendelse bliver da mulig trods
Indholdets Skiften, og det personlige Livs Enhed afspejler sig i den
Kontinuitet, som derved tilvejebringes. Erkendelsesarbejdet hænger
her nøje sammen med Arbejdet paa Karakterens Udvikling. Om en
Karakter tale vi — i udpræget Forstand — kun, hvor der hersker
Enhed og Kontinuitet i den hele Stræben og ikke atter og atter
springes over til noget Nyt. I Analogi hermed søger Erkendelsen at
finde saa megen Enhed og Kontinuitet som muligt, at paavise det Nye
som Omsætning eller Fortsættelse af det Gamle, og kun hvor dette
lykkes, er en fuld \tlg{Tilegnelse} og Forstaaelse vunden. De naturlige
Aarsagers Princip, der fører til at bygge saa lange og saa
kontinuerlige Rækker som muligt, er derfor ingen Tilfældighed, men
hænger nøje sammen med Personlighedens eget Væsen. Det er derfor
urigtigt, naar Personlighed og Videnskab sættes i skarp Modsætning
til hinanden. Videnskabeligt Arbejde er et Personlighedsarbejde. Og
jo mere dette Arbejde lykkes, des mere nærmer Aarsagsforklaringen sig
til Begrundelsen og Genkendelsen. Forskellighederne i Tilværelsen ere
reducerede til det mindst Mulige, naar vi med ét Blik, “under
Evighedens Synspunkt” overskue store Egne af Tilværelsen.
\end{quote}
\dcite{Høffding, Religionsfilosofi, pp. 24–25  [hoeffding/religionsfilosofi:41]}
% hoeffding/religionsfilosofi:41  (Religionsfilosofi; pp. 24–25)

\subsection{\textit{Vort Hjem. En Indledning} (1903)}

\begin{quote}
I. — Ordet Hjem vækker Forestillinger om Inderlighed, Fylde og Hvile.
Om Inderlighed — thi skønt Hjemmet ude fra set er et Sted som alle andre
Steder, er det dog for os det Sted, hvor vi modtog alt, hvad Skæbnen bragte
os, den Vig, til hvilken Virkningerne fra den ydre Verdens store Hav maatte
forplante sig, for at vi skulde mærke og værdisætte dem; og det er tillige det
Sted, fra hvilket de Virkninger, vi selv maatte have øvet, ere udgaaede. I al
vor Virken og Liden, al vor Given og Modtagen væver derfor Hjemmets Stemning
sig ind og danner en Baggrund, som vi maaske først blive os klart bevidst,
naar den trues af Opløsning eller dog undergaar dybere gaaende Ændring. Og
hermed hænger ogsaa Fylden sammen. Thi det er en stor Kreds af Erindringer og
Vaner, egentligt hele vort Livs Indhold, der hænger sammen med denne centrale
Forestilling. Hjemmet er som Nettet, der ikke kan drages op, uden at et Mylr
af levende Indhold følger med. Hvilen endelig beror baade paa Inderlig-heden
og Fylden. Hvile kan ikke findes i det fremmede og nye, kun i det, der er
kendt og prøvet, og hvis \tlg{Tilegnelse} ingen Anspændelse kræver. Det hjemlige er
lempet efter vort Væsens Former; alle Folder og Fuger i vor Natur finde der et
Leje, som passer til dem. Naar det hjemlige atter dukker op for os, er det,
som om det aldrig havde været borte; let og hurtigt føje vi os ind i de
forlængst tillempede Rammer. Og den Rigdom af Vaner og Erindringer, Hjemmet
bærer, bringer ogsaa Hvile. Paa en skarp Od kan man ikke hvile; der maa være
en Mangfoldighed i Indholdet, for at den stadige Dvælen ved et og det samme
ikke skal trætte og sløve.
\end{quote}
\dcite{Høffding, Vort Hjem. En Indledning, pp. 3–4  [hoeffding/vort-hjem:1]}
% hoeffding/vort-hjem:1  (Vort Hjem. En Indledning; pp. 3–4)

\begin{quote}
Men for at forstaa de Tvivl, der kunne opstaa om Hjemmets Værdi, maa vi mærke
os, at Loven om det omvendte Forhold mellem Styrke og Omfang ikke er den
eneste Lov for menneskelige Interessers Udvikling. Der viser sig nemlig paa
Følelseslivets Omraade ogsaa et omvendt Forhold mellem Styrke og Indhold. En
Følelses Indhold er de Forestillinger, til hvilke den er knyttet, og som giver
den dens ejendommelige Beskaffenhed. Glæde og Sorg nuanceres og varieres efter
det Forestillingsindhold, der betinger dem. Nu er det klart, at med det snævre
Omfang følger ogsaa let et snævert Indhold. I de smaa Kredse antager
Følelseslivet let en rent elementær Karakter, kommer til at mangle bestemt og
mangesidig Udformning. Ved Isolation fra Livet i de større Kredse indtræder da
— selv om Følelsens Styrke bevares, hvilket paa Grund af Vanens og
Ensformighedens Indflydelse ikke altid sker — let en aandelig Fattigdom.
Følelsen bliver da blind. Det Indhold, den har, og som maaske fastholdes med
stor Inderlighed og Lidenskab, stammer væsenligt fra Fortiden. Familiens
Husguder ere de forbigangne Tiders Heroer; men det er ikke sagt, at de egne
sig til at være Førere i de nye Tider. Hjemmets konservative Karakter betinges
herved. Inderligheden, Trygheden og Hvilen bliver da til Træghed og afføder
Uvillie og Frygt over for det nye. Hvad der ellers udgør Hjemmets Styrke,
bliver her dets Svaghed. Alle Tiders Historie frembyder Erfaringer i denne
Retning. Her ville vi kun fremdrage et Par særligt iøjnefaldende Træk. —
Nyligt har en japanesisk Forfatter, Tokiwo Yokoi, i Anledning af Krigen mellem
Kina og Japan anstillet en Sammenligning mellem de to Folks herskende Moral og
har herved lagt stor Vægt paa, at den kinesiske Moral væsenligt er en
Familiemoral: »Det var den sønlige Pietets Princip, som var Grundlaget for
Kinesernes Moral, snarere end Principet om Loyalitet over for Fyrsten.« Denne
Retning slog den kinesiske Etik saa meget mere ind paa, som Dynastierne i Kina
have skiftet meget hyppigt. Familieforholdet blev da det eneste Baand, der
kunde hævde Enheden og Kontinuiteten i Samfundet, og det med saadan Fasthed,
at »disse tre Hundrede Millioner Mennesker have bevaret deres Traditioner og
Skikke uforandrede i tre Tusinde Aar og trods tre og tyve
Dynastiforandringer.« Men de Dyder, der udmærke Kineserne, ere derfor ogsaa
kun det private og huslige Livs Dyder, hvorimod det skorter paa Patriotisme,
paa Sans for og Ærlighed i offentlige Forhold. Deraf forklares Vanskeligheden
ved Optagelse og Gennemførelse af Reformer i Kina, medens Japanerne, der
gennem lange Tider vare øvede i at underordne private og huslige Interesser
under Statens Krav, med største Iver og Opofrelse slog ind paa Fremskridtets
Vej. — Kyndige Iagttagere af Forholdene i Tyrkiet og i Persien have fundet,
at de orientalske Haremer maaske yde den vigtigste Modstand mod en ny Kulturs
Indtrængen. »Medens der,« siger Vambéry (Der Islam des neunzehnten
Jahrhunderts), »i Selamlik (den Del af Huset, der er tilgængelig for Mændene
og den fremmede Verden) hist og her viser sig Spor af Fornyelse og à la
franca saa at sige er blevet et Løsen, er i Haremet, den muhamedanske
Families Helligdom, alt blevet ved det gamle, og der er ikke sket den mindste
Forandring i Møblementet, Sæd, Skik og Husorden, saa lidt som i Dragt,
Talebrug og Tænkemaade Hos nogle forstokkede Ortodokse er Haremet i
Virkeligheden blevet et lykkeligt Tilflugtssted, medens det selvfølgelig er en
vældig Hæmsko for de Forkæmpere for det nye, der mene det alvorligt. Det nye,
der indføres i Selamlik, bliver i Haremet enten ignoreret eller revet ned igen
af ubændig Fanatisme.« — Det mangler ikke i Europa paa Tendenser i lignende
Retning. Man har saaledes — baade i protestantiske og i katolske Lande —
opstillet den Paastand, at Skolen skal være afhængig af Hjemmet. »Hjem, Stat
og Kirke«, sagde Biskop Ketteler, »maa have Skoler, der svare til deres Aand
og deres Fordringer.« Der mangler her den supplerende Sætning, at Skolen — i
videste Forstand: Samfundet til Fremme og \tlg{Tilegnelse} af Videnskaben — atter
skal virke tilbage paa den Aand og de Fordringer, der raade i Hjem, Stat og
Kirke, — at der er en større Horisont, som Hjemmets (og de andre nævnte
Samfunds) Tendens til Afslutning let fører til at begrænse. Hjemmets ofte
lumre og forgemte Luft trænger til at renses gennem friske Strømninger. I
saadanne Tilfælde kæmper man nu egentligt ikke for Hjemmet, men man bruger
Hjemmet som Middel til at bevare Dogmer, som ikke paa anden Maade kunne holdes
oppe. Men det Kulturindhold, paa hvilket Hjemmet skal tære, vil ved denne
Fremgangsmaade blive ringere og vil stivne i døde Former.
\end{quote}
\dcite{Høffding, Vort Hjem. En Indledning, pp. 7–9  [hoeffding/vort-hjem:11]}
% hoeffding/vort-hjem:11  (Vort Hjem. En Indledning; pp. 7–9)

\subsection{\textit{Den menneskelige Tanke} (1910)}

\begin{quote}
Overfor den Bemærkning, at de Begreber, vi i vor Eftertanke bruge til
at karakterisere Sjælelivet, — Syntese, Enhed, Mangfoldighed
o. lign., — finde Anvendelse paa alle Emner, vil jeg hævde, at
Grunden til, at hin Gruppe af Begreber stadigt finde Anvendelse og maa
finde Anvendelse paa alle Emner, netop maa søges i, at de udtrykke de
Former, i hvilke Sjælelivet og særligt Erkendelsen udfolder sig. Kun
ved Anvendelse af disse Former bliver den \tlg{Tilegnelse}, den Optagen i
vort eget Væsens Sammenhæng, hvori al Forstaaelse, al Erkendelse
bestaar, mulig. Intet Under, at alle Erkendelsesprodukter bære deres
Præg! Vor Erkendelse bestaar i en bevidst eller ubevidst Analogi med
vort eget Væsen. —
\end{quote}
\dcite{Høffding, Den menneskelige Tanke, p. 17+  [hoeffding/menneskelige-tanke:39]}
% hoeffding/menneskelige-tanke:39  (Den menneskelige Tanke; p. 17+)

\begin{quote}
1) Formale Forskelligheder — a) Elementære Værdier forudsætte et sanseligt Nærværende som Genstand for
Nydelse og \tlg{Tilegnelse}, som Spore og som Støtte. Ideelle Værdier forudsætte
Erindring eller Forventning, Eftertanke eller Fantasi. Da de ikke støttes af
sanselig Nærværelse, forudsætte de større psykisk Energi for at kunne fastholdes.
De have tillige det Fortrin at være uafhængige af ydre Betingelser, saa at de
kunne staa til Raadighed, saasnart de indre Betingelser ere tilstede. Dermed staa
de over de elementære Værdier.
\end{quote}
\dcite{Høffding, Den menneskelige Tanke, p. 241+  [hoeffding/menneskelige-tanke:580]}
% hoeffding/menneskelige-tanke:580  (Den menneskelige Tanke; p. 241+)

\begin{quote}
Det viste sig i Kategorilæren, at Formerne antog stedse speciellere Karakter
formedelst de Emner, der forelaa. Her ligger det første filosofiske Problem,
bestemt ved Forholdet mellem Form og Emne. Kunne Formerne afledes af Emnerne,
eller omvendt? Med hvilken Ret anvendes Formerne paa Emnerne, hvis de ikke kunne
afledes af disse? Er det muligt at naa en fuldkommen og udtømmende \tlg{Tilegnelse} af
Emnerne ved Hjælp af Tankens Former, saaledes som disse nu engang ere? —
Saadanne Spørgsmaal høre under Erkendelsesproblemet. Det tredie Element,
Motivet eller Værdien, gør sig ganske vist ogsaa her gældende. Der kunde jo
opkastes det Spørgsmaal, om ikke Valget af de Emner, der behandles, og af de
Former, der anvendes, skyldes et Hensyn til Værdi, et Motiv, der kan komme mere
eller mindre til Bevidsthed. Ogsaa dette Spørgsmaal maa drøftes under
Erkendelsesproblemet, selvom det ellers kan neutraliseres ved, at vi foreløbigt
blot forudsætte den intellektuelle Interesse, som umiddelbart er rettet imod
Klarhed over Forholdet mellem Former og Emner.
\end{quote}
\dcite{Høffding, Den menneskelige Tanke, p. 249  [hoeffding/menneskelige-tanke:595]}
% hoeffding/menneskelige-tanke:595  (Den menneskelige Tanke; p. 249)

\begin{quote}
Hvis man savner det æstetiske Problem i denne Opstilling, vil jeg pege hen
paa dets Analogi med det etiske Problem. Æstetisk Værdi beror paa, at et Emne
fremtræder som en karakteristisk Helhed og anskues og fastholdes blot som saadan,
afset fra intellektuel og praktisk Interesse. Æstetisk betragte vi Virkeligheden
som Billede og Billedet som Virkelighed; det er os nok, naar Emnet — det være nu
selve Virkeligheden eller dens Billede — afrunder sig til et ejendommeligt Hele.
Ogsaa det Hæslige og det Disharmoniske kan falde ind herunder; det æstetiske
Problem vilde da angaa Betingelsen for, at saadanne Billeder kunne blive til og
blive Genstand for Arbejde og \tlg{Tilegnelse}. I Etiken drejer det sig derimod om
Arbejdet for selve den virkelige Helhed, der staar som Grundværdien. Her bliver
Forskellen mellem Billede og Virkelighed af afgørende Betydning. Det er
Helhedskategorien, der ligger til Grund for al Vurdering. Vi kunne opfatte Helheden
som Billede (æstetisk Problem), som Grundlag for Bedømmelse af Handlinger og
Bestræbelser (etisk Problem), og som den, hvis Skæbne bestemmer vor Følelse af
Livets og Tilværelsens Værdi (religiøst Problem). Til en fuldstændig Fremstilling
af de filosofiske Problemer vilde kræves, at det æstetiske Problem fik en særegen
Plads. Men det Typiske i vort Forhold til Problemerne vil forhaabentligt kunne
træde frem, selvom jeg holder mig til de Problemer, jeg føler mig bedst forberedt
til at behandle. Iøvrigt maa jeg henvise til min Psykologi (V B, 12; VI C, 9) og
min Etik (XXX, 1–2).
\end{quote}
\dcite{Høffding, Den menneskelige Tanke, pp. 250–251  [hoeffding/menneskelige-tanke:598]}
% hoeffding/menneskelige-tanke:598  (Den menneskelige Tanke; pp. 250–251)

\begin{quote}
Der rører sig en Mangfoldighed af Følelser, Drifter og Tanker, af Evner og Tilskyndelser,
som kan nærme sig mere eller mindre til en kaotisk Forskelsrække, og som det gælder om at
ordne og forme. Og det personlige Liv udfolder sig i en Række af Øjeblikke og gennem
skiftende Situationer. Den successive saavel som den simultane Mangfoldighed skal formes
saaledes, at Helheden hævdes, uden at Fylden lider derunder. I det primitive Aandsliv raade
her Instinkt og Tradition. Den psykiske Energi vækkes let og ledes ligesaa let ad bestemte
Veje. Den primitive Etik er dels biologisk, dels sociologisk. Allerede det organiske Liv
kræver til sin Bestaaen, at der bevares en Sammenhæng mellem Livets Elementer. Her raader
Instinktet. Og da ethvert menneskeligt Individ bliver til indenfor en Gruppe og
uvilkaarligt lever sig ind i dens Ordning og dens psykiske Atmosfære, saa faar den Maade,
paa hvilken det stræber efter at harmonisere sine indre Emner, fra først af en social
Karakter. Som det \tlg{tilegner} sig Sproget, og altsaa former sin Trang til at give sine
Sindsbevægelser, sine Erfaringer og sine Tanker Udtryk efter Omgivelsernes Forbillede,
saaledes er det ogsaa de sociale Forbilleder og de i Gruppen overleverede Bud, der forme
dets Villiesliv. Kun naar den uvilkaarlige Tryghed, som Overleveringer, Autoriteter og
Forbilleder fremvirke, af en eller anden Aarsag ophæves, optræder der etiske Problemer. Da
ere det uvilkaarligt stillede Formaal og det med instinktiv Hengivelse fulgte Forbillede
ikke længere selvfølgelige. Og det bliver i hvert Tilfælde ikke selvfølgeligt, at Maalet
skal naas ad den hidtil fulgte Vej, eller at man skal efterfølge Forbilledet paa den hidtil
antagne Maade. Man opdager Vanskeligheder og Konflikter, hvor Alt før syntes saa naturligt
og saa nødvendigt. Nye Værdier, og dermed nye Formaal, kunne opdages; kunne de forenes med
de overleverede, og kan der overhovedet findes et fælles Synspunkt for dem? — Der kan
her, praktisk set, vise sig Mulighed for en kaotisk Forskelsrække. Mennesket vil dette, men
vil ogsaa hint; Villiens naturlige Enhed, der tilsidst bunder i den psykiske Energis Enhed,
truer med at sprænges. Der maa gennem de mange Formaal og de mange Midler strække sig en
skjult Identitet, eller i hvert Tilfælde en Kontinuitet, hvis den personlige Helhed skal
bestaa.
\end{quote}
\dcite{Høffding, Den menneskelige Tanke, pp. 350–351  [hoeffding/menneskelige-tanke:825]}
% hoeffding/menneskelige-tanke:825  (Den menneskelige Tanke; pp. 350–351)

\subsection{\textit{Erindringer — Uddrag af Kapitel II} (1928)}

\begin{quote}
24. Novbr. 1867. — En Forandring, der er foregaaet med mig i dette Efteraar, bestaar i, at jeg ikke længere kan stemme med Rasmus Nielsen, til hvem jeg, som du véd, før idethele havde sluttet mig. Dels en Vaklen hos ham selv, dels en slaaende Kritik af Professor Brøchner, dels og fremfor Alt en fornyet Prøvelse af Sagen selv (ikke mindst i Samtaler med Poul Kierkegaard) har vist mig det Uholdbare i hans Anskuelse. Dog er jeg derved ikke ført nærmere til det Kirkelige... Tværtimod er det med det samme gaaet op for mig, at jeg paa flere Punkter havde Meninger, som ikke havde nogen sand Rod hos mig, men hang fast som noget Udvortes og Overleveret. Kristendommen vil — det er jeg fuldt og fast overbevist om — altid være mig den højeste Sandhed, den, i hvilken jeg ene kan finde Livsgaaderne løste. Men kun hvad jeg personligt kan erfare og \tlg{tilegne} mig, kan være Sandhed for mig; og den Overbevisning er jeg kommen til, at hvad man mener sig at besidde, som gaar ud over det, beror paa ubevidst Hykleri. Den Anskuelse, hvis Uholdbarhed jeg nu er kommen til at indse, har naturligvis ikke været uden Betydning; det var et Udviklingstrin, som maatte gennemgaas.
\end{quote}
\dcite{Høffding, Erindringer — Uddrag af Kapitel II, pp. 67–69  [hoeffding/erindringer:31]}
% hoeffding/erindringer:31  (Erindringer — Uddrag af Kapitel II; pp. 67–69)

\begin{quote}
Han var en filosofisk Eneboer. Den Kreds af Tanker, han havde \tlg{tilegnet} sig, var hans aandelige Hjem. Han havde ikke megen Evne til at meddele sig i fri Form, skiftende sine Udtryk efter Situationer og Personligheder, og han havde en Gang for alle sammenfattet sine Ideer i bestemte, meget abstrakte Former. Men baade paa Forelæsninger og i Samtale kunde der komme en Vibreren i Stemme og Ansigtsudtryk, et Tegn paa den dybe Stemning, der laa bagved og vidnede om Tankens Sammenhæng med hele Personligheden. Han var en Tankens Ridder, den sidste Repræsentant iblandt os for den Idealisme, der i Tillidt til, at Tilværelsens Inderste kunde finde og havde fun det Udtryk i Tankens Form, tillige i denne Form fandt Udtryk for det personlige Livs højeste Higen. Problemerne ere for dem, der fulgte efter ham paa Filosofiens Omraade, blevne mere sammensatte og mangfoldige. Paa en af sine sidste Forelæsninger erklærede Brøchner sig for Tilhænger af „den idealistiske Filosofi, som Tyskland har Æren af at have frembragt“, og der har neppe været mange, hos hvem denne Filosofis Hovedtænker i den Grad gennemtrængte og adlede Personligheden som hos ham. Billedet af „den stille Professor“, som han er bleven kaldt, har fulgt mig paa min Vej og mindet mig om, hvilken Værdi Tankelivet kan have, ogsaa ud over det rent intellektuelle Omraade, naar der er Alvor i Tankens Arbejde. Hans Venskab var mig en Ære og en Opmuntring, og det var med dyb Bevægelse, jeg tog Afsked med ham paa hans Dødsleje. Det var Dagen, før han døde, jeg saa ham sidste Gang. Han talte med fuldstændig Ro om sin Bortgang, og aftalte med mig Ordningen af forskellige Sager. Da han bød mig Farvel og takkede mig for mit Venskab, var jeg mere betagen end han. Hans Husbestyrerinde fortalte senere, at han havde sagt ganske roligt, da jeg var gaaet: „Tag det Glas bort, ti Dr. Høffdings Taarer ere deri.“
\end{quote}
\dcite{Høffding, Erindringer — Uddrag af Kapitel II, pp. 78–79  [hoeffding/erindringer:36]}
% hoeffding/erindringer:36  (Erindringer — Uddrag af Kapitel II; pp. 78–79)

\begin{quote}
I Fortalen til min Udgave af Brøchners og Molbechs Brevvexling og i „Danske Filosofer“ har jeg givet en Karakteristik af Brøchner som Personlighed og som Tænker, og mit Værk „Den nyere Filosofis Historie“ har jeg \tlg{tilegnet} „Mindet om min Lærer og Ven“.
\end{quote}
\dcite{Høffding, Erindringer — Uddrag af Kapitel II, p. 80  [hoeffding/erindringer:38]}
% hoeffding/erindringer:38  (Erindringer — Uddrag af Kapitel II; p. 80)


\end{document}
